\documentclass[11pt]{article}
\usepackage[a4paper,margin=24mm]{geometry}
\usepackage{amsmath,amssymb,mathtools}
\usepackage{fontspec,unicode-math}
\setmainfont{Cambria}
\setmathfont{Cambria Math}
\setmonofont{Cambria Math}
\newfontfamily\Rsixtyonesymbolfont{Segoe UI Symbol}
\usepackage{longtable,booktabs,array,calc,graphicx}
\usepackage{pdflscape}
\usepackage{fvextra}
\DefineVerbatimEnvironment{verbatim}{Verbatim}{breaklines,breakanywhere,fontsize=\small}
\usepackage[hidelinks,unicode]{hyperref}
\providecommand{\tightlist}{\setlength{\itemsep}{0pt}\setlength{\parskip}{0pt}}
\providecommand{\pandocbounded}[1]{#1}
\setlength{\emergencystretch}{2em}
\title{Derived Categories Source Review and Cross-Chapter Corrections\\Complete Editorial Arguments}
\author{AI editorial work: OpenAI Codex---GPT-6 Astra, Ultra effort}
\date{30 September 2026}
\begin{document}
\maketitle
\noindent
The original source is by the Stacks Project authors. These twenty-five notes
supply the complete arguments for the Derived Categories review and its
receiving passages in More on Algebra, Cohomology of Sheaves, and the
derived categories of schemes, varieties and spaces. The chapter files
retain the source exposition with eighty-nine bounded edits in sixty-four units. The longer arguments, proof
completions, and stronger conclusions here are identified as editorial
work. The thirty-three recorded consequences and the refinement of UNDER-007
are editorial results, with no novelty claim.
The original proof notes and their source locators are included. Statements
about review progress in those notes record their preparation history;
release status is recorded separately in the candidate receipts.
\setcounter{secnumdepth}{1}
\makeatletter
\renewcommand\@pnumwidth{2.5em}
\renewcommand\@tocrmarg{3.5em}
\makeatother
\tableofcontents\clearpage
\section{\texorpdfstring{Derived Categories:\allowbreak{} opening definitions and elementary results}{Derived  opening definitions and elementary results}}\label{proofs-0001-1140-derived-categories-opening-definitions-and-elementary-results}

Authority:\allowbreak{} the Stacks Project authors,\allowbreak{} \href{https://github.com/stacks/stacks-project/blob/a04446e57ec1fbc252a871afcec7752fb2807b14/derived.tex}{Derived Categories,\allowbreak{} frozen source},\allowbreak{} commit \texttt{a04446e5\allowbreak{}7ec1fbc2\allowbreak{}52a871af\allowbreak{}cec7752f\allowbreak{}b2807b14},\allowbreak{} SHA-256 \texttt{EFDA4CB3\allowbreak{}61DD4090\allowbreak{}9D199116\allowbreak{}1EAE716E\allowbreak{}58B6CF9E\allowbreak{}79380971\allowbreak{}502E1DA1\allowbreak{}2396D402}. Line numbers in this note refer to that original source unless explicitly identified as cumulative-source lines. The compared cumulative source is bound separately in \texttt{INPUT\_\allowbreak{}BINDING.\allowbreak{}json} and \texttt{CURRENT\_\allowbreak{}PUBLIC\_\allowbreak{}BINDING.\allowbreak{}json}. These arguments are editorial review,\allowbreak{} not prose attributed to the source authors and not claims of mathematical novelty. No live chapter is edited by this note. Generalizations remain separate from the source text.

\subsection{1. Conventions and the first three reports}\label{proofs-0001-1140-conventions-and-the-first-three-reports}

Write a triangle as \(T=(X,Y,Z,f,g,h)\),\allowbreak{} with \(f:X\to Y\),\allowbreak{} \(g:Y\to Z\),\allowbreak{} \(h:Z\to X[1]\). A morphism \((a,b,c):T\to T'\) satisfies \[
bf=f'a,\qquad cg=g'b,\qquad a[1]h=h'c.
\] Keep the source rotation \[
R(T)=(Y,Z,X[1],g,h,-f[1]).
\] The chosen comparison isomorphisms for repeated translations are those of original lines 82–\allowbreak{}88; no strict equality of unrelated presentations is being imposed. Every iterated-shift expression below uses these comparisons. For complexes the source shift is strict.

Applying TR3 from \((X,X,0,1,0,0)\) to \(T\) gives \(gf=0\). Rotation gives \(hg=0\) and \(f[1]h=0\). If \(a:W\to Y\) satisfies \(ga=0\),\allowbreak{} rotate the diagram from \((W,W,0,1,0,0)\) to \(T\) and apply TR3 to the two supplied arrows \(a,0\). Full faithfulness of the shift supplies \(b:W\to X\) with \(fb=a\). This proves exactness of the representable sequence at \(\operatorname{Hom}(W,Y)\). Rotation proves exactness at its other terms. The dual construction proves the contravariant representable sequence. These are the exactness statements used in the proofs below.

Original lines 699–\allowbreak{}727 give the splitting criterion. In detail,\allowbreak{} if \(h=0\),\allowbreak{} exactness of \(\operatorname{Hom}(Z,Y)\to\operatorname{Hom}(Z,Z)
\to\operatorname{Hom}(Z,X[1])\) supplies \(s:Z\to Y\) with \(gs=1\). Conversely \(gs=1\) implies \(h=hgs=0\). The sequence \[
0\longrightarrow\operatorname{Hom}(W,X)
\xrightarrow{f_*}\operatorname{Hom}(W,Y)
\xrightarrow{g_*}\operatorname{Hom}(W,Z)\longrightarrow0
\] is split by \(s_*\); its left injectivity follows from the preceding zero map induced by \(h[-1]\). Thus \((f,s):X\oplus Z\to Y\) induces a bijection on every representable functor and is an isomorphism.

For the kernel criterion at 729–\allowbreak{}759,\allowbreak{} a monomorphism \(f\) forces \(h[-1]=0\) because \(fh[-1]=0\). It therefore splits by the previous paragraph. If \(k:K\to X\) is a kernel of an arbitrary \(f\),\allowbreak{} the monomorphism \(k\) splits,\allowbreak{} giving \(X\cong K\oplus X'\). The restriction \(f':X'\to Y\) is monic:\allowbreak{} if \(f't=0\),\allowbreak{} its composite with the inclusion into \(X\) factors through \(k\); projection to \(X'\) gives \(t=0\). Hence \(Y\cong X'\oplus Y'\). This is precisely the original projection–\allowbreak{}coprojection decomposition,\allowbreak{} with every object retained. The cokernel argument follows in the opposite category. This proves the mathematical context of \textbf{OCC-01491 /\allowbreak{} P02-E0019}:\allowbreak{} deleting the extra article in ``is the isomorphic'' changes only grammar.

\textbf{OCC-02889 /\allowbreak{} P02-E1426} correctly supplies ``of'' after ``application'' and changes ``are'' to ``is'' at 1068–\allowbreak{}1069. The two applications of TR4 and all their arrows are verified in §3 below. The singular ``conclusion'' may introduce the enumerated assertions; no mathematical amendment is required by this report.

\textbf{OCC-02890 /\allowbreak{} P02-E1427} correctly removes the period between the two displayed triangle morphisms and ``are morphisms of triangles'' at 1080–\allowbreak{}1081. The two morphisms are the compound subject of this one sentence. The proposed operation removes that period without changing either map.

The resolution-only record \textbf{OCC-02894 /\allowbreak{} P02-ER0003} has no numeric range. Its exact text routes it to P02-E0094 at 11522 and 11624–\allowbreak{}11625. It is retained for adjudication with that later report,\allowbreak{} not counted as a new defect or silently discarded.

\subsection{2. Explicit functor maps}\label{proofs-0001-1140-explicit-functor-maps}

At 190–\allowbreak{}200 an exact functor is \((F,\xi)\),\allowbreak{} where \(\xi_X:F(X[1])\xrightarrow{\sim}F(X)[1]\). Thus the fourth arrow of its image triangle is \(\xi_XF(h)\),\allowbreak{} with codomain \(F(X)[1]\); \(F(h)\) alone has codomain \(F(X[1])\).

\textbf{DERIVED-NEW-001.} In the additivity proof,\allowbreak{} 893–\allowbreak{}898,\allowbreak{} retain this comparison for the triangle on the zero object:\allowbreak{} \[
(F(0),F(0),F(0),1_{F(0)},1_{F(0)},\xi_0F(0)).
\] Here the last \(0\) denotes the original zero morphism \(0\to0[1]\). Name \(\xi\) in the proof. The first two arrows still imply \(1_{F(0)}=1_{F(0)}1_{F(0)}=0\). An object whose identity is zero is a zero object:\allowbreak{} every map into or out of it equals its composition with the zero identity. Consequently \(F\) preserves every zero map,\allowbreak{} since such a map factors through \(0\).

For the biproduct \(X\oplus Y\),\allowbreak{} the image of the distinguished split triangle has the two maps \(F(i_X),F(p_Y)\). The map \(F(i_Y)\) is a right inverse of \(F(p_Y)\),\allowbreak{} so the splitting criterion proves that \((F(i_X),F(i_Y)):F(X)\oplus F(Y)\to F(X\oplus Y)\) is an isomorphism. The images of the projections give its inverse,\allowbreak{} because \(F(p_X)F(i_X)=1\),\allowbreak{} \(F(p_Y)F(i_Y)=1\),\allowbreak{} and both cross composites are zero. Hence the biproduct identities and their diagonal and codiagonal maps are preserved. For \(a,b:U\to V\),\allowbreak{} the formula \(a+b=\nabla_V(a\oplus b)\Delta_U\),\allowbreak{} with these same biproduct maps,\allowbreak{} proves \(F(a+b)=F(a)+F(b)\). No additivity was assumed while proving that \(F(0)\) is zero.

\textbf{DERIVED-NEW-002.} At 917–\allowbreak{}920,\allowbreak{} name \((F,\xi)\) and put \(\xi_XF(h)\) in the displayed image triangle. Here is the fully typed reflection argument. Choose a distinguished completion \((X,Y,Z',f,g',h')\). Its image and the assumed image of \(T\) have the same first map. TR3 and representable exactness give an isomorphism \((1,1,c')\) between them. Fullness supplies \(c:Z\to Z'\) with \(F(c)=c'\),\allowbreak{} and also supplies a preimage of \((c')^{-1}\). Faithfulness proves that the two preimages are inverse. The second triangle square gives \(F(cg)=F(g')\),\allowbreak{} hence \(cg=g'\). The third gives \[
\xi_XF(h)=\xi_XF(h')F(c).
\] Cancelling \(\xi_X\) and then using faithfulness gives \(h=h'c\). Thus \((1,1,c)\) is the required isomorphism of original triangles. The mathematical theorem is unchanged; its comparison map is made explicit. The source may intend an implicit identification,\allowbreak{} but its definition allows arbitrary translation isomorphisms rather than strict commutation.

\textbf{DERIVED-NEW-003.} At 1011 the triangle attached to the short exact sequence \(0\to A\xrightarrow{a}B\xrightarrow{b}C\to0\) is \[
(G(A),G(B),G(C),G(a),G(b),\delta),
\] as the earlier definition at 298–\allowbreak{}331 explicitly requires. The arrows \(a,b\) themselves belong to the abelian source category. Replace them by their images in this tuple. Applying \(H\) after translations gives the exact sequence in the source proof,\allowbreak{} with connecting map \(H^n(\delta):H^nG(C)\to H^{n+1}G(A)\),\allowbreak{} using the fixed shift comparisons. Its preceding term \(H^{-1}G(C)\) vanishes by the stated hypothesis,\allowbreak{} so the displayed initial zero is justified. Naturality of \(\delta\) gives naturality of every connecting map.

For completeness,\allowbreak{} \(G\) is additive even though the definition calls it a functor without assuming additivity. Apply it to the zero short exact sequence:\allowbreak{} the resulting triangle has its first two maps equal to \(1_{G(0)}\),\allowbreak{} so \(G(0)=0\) by the previous zero-object argument. Apply it next to a split short exact sequence \(0\to A\to A\oplus B\to B\to0\). Its projection still has the right inverse supplied by \(G\),\allowbreak{} so the same biproduct argument proves that \(G\) is additive. Consequently every \(H^nG\) is additive,\allowbreak{} as required in the receiving abelian-category definition of a delta-functor.

These three findings are explicit-map repairs. They do not assert that the intended additivity,\allowbreak{} reflection,\allowbreak{} or delta-functor results fail.

\subsection{\texorpdfstring{3. The 3×3 proposition:\allowbreak{} all nine squares and the outer signs}{3. The 3×3  all nine squares and the outer signs}}\label{proofs-0001-1140-the-33-proposition-all-nine-squares-and-the-outer-signs}

\textbf{DERIVED-NEW-004} concerns 1048–\allowbreak{}1052. The first three rows and first three columns are distinguished as drawn. The fourth row and fourth column are the unsigned shifts of the first row and column; each becomes distinguished after negating its last arrow. They need not be distinguished with all three unsigned shifted arrows.

We verify the construction at 1055–\allowbreak{}1140 without suppressing its maps. Let the given square have arrows \[
f:X\to Y,\quad u:X\to X',\quad v:Y\to Y',\quad
f':X'\to Y',\qquad vf=f'u=w.
\] Choose the five original distinguished triangles \[
\begin{split}
&(X,Y,Z,f,g,h),\qquad(X',Y',Z',f',g',h'),\\
&(X,X',X'',u,p,q),\qquad(Y,Y',Y'',v,r,s),\\
&(X,Y',A,w,t,e).
\end{split}
\] For the factorization \(w=vf\),\allowbreak{} TR4 supplies \(a:Z\to A\) and \(b:A\to Y''\) with distinguished triangle \((Z,A,Y'',a,b,c)\),\allowbreak{} where \(c=g[1]s\). Its two morphisms of triangles give the exact identities \[
ag=tv,\quad ea=h,\quad bt=r,\quad sb=f[1]e.
\] For \(w=f'u\),\allowbreak{} TR4 supplies \(a':X''\to A\) and \(b':A\to Z'\) with distinguished triangle \((X'',A,Z',a',b',c')\),\allowbreak{} where \(c'=p[1]h'\). This time the identities are \[
a'p=tf',\quad ea'=q,\quad b't=g',\quad h'b'=u[1]e.
\] Set \(k=ba':X''\to Y''\),\allowbreak{} \(z=b'a:Z\to Z'\),\allowbreak{} and choose a distinguished triangle \((X'',Y'',Z'',k,m,n)\). Apply TR4 to \(X''\xrightarrow{a'}A\xrightarrow bY''\),\allowbreak{} using the rotation \((A,Y'',Z[1],b,c,-a[1])\) for the third input triangle. It gives \(\alpha:Z'\to Z''\),\allowbreak{} \(\beta:Z''\to Z[1]\),\allowbreak{} and \[
\alpha b'=mb,\quad n\alpha=c'=p[1]h',\quad
\beta m=c=g[1]s,\quad -a[1]\beta=a'[1]n.
\] The fourth triangle is \((Z',Z'',Z[1],\alpha,\beta,-(b'a)[1])\). Its inverse rotation is exactly \((Z,Z',Z'',z,\alpha,\beta)\),\allowbreak{} distinguished by TR2. Thus the full diagram,\allowbreak{} including the names of every arrow,\allowbreak{} is \[
\begin{array}{ccccccc}
X&\xrightarrow f&Y&\xrightarrow g&Z&\xrightarrow h&X[1]\\
{\scriptstyle u}\downarrow&&{\scriptstyle v}\downarrow&&
 {\scriptstyle z}\downarrow&&{\scriptstyle u[1]}\downarrow\\
X'&\xrightarrow{f'}&Y'&\xrightarrow{g'}&Z'&\xrightarrow{h'}&X'[1]\\
{\scriptstyle p}\downarrow&&{\scriptstyle r}\downarrow&&
 {\scriptstyle\alpha}\downarrow&&{\scriptstyle p[1]}\downarrow\\
X''&\xrightarrow k&Y''&\xrightarrow m&Z''&\xrightarrow n&X''[1]\\
{\scriptstyle q}\downarrow&&{\scriptstyle s}\downarrow&&
 {\scriptstyle\beta}\downarrow&&{\scriptstyle q[1]}\downarrow\\
X[1]&\xrightarrow{f[1]}&Y[1]&\xrightarrow{g[1]}&Z[1]&
 \xrightarrow{h[1]}&X[2].
\end{array}
\] In row order the nine square relations are \[
\begin{array}{lll}
vf=f'u,& zg=g'v,&u[1]h=h'z,\\
kp=rf',&\alpha g'=mr,&p[1]h'=n\alpha,\\
f[1]q=sk,&g[1]s=\beta m,&q[1]n=-h[1]\beta.
\end{array}
\] Indeed,\allowbreak{} the second and third follow by composing the two first-octahedron identities with \(b'\). The fourth is \(ba'p=btf'=rf'\). The fifth is \(\alpha g'=\alpha b't=mbt=mr\). The sixth and eighth are displayed third-octahedron identities. The seventh is \(sba'=f[1]ea'=f[1]q\). For the last retain the actual minus sign:\allowbreak{} \[
q[1]n=e[1]a'[1]n=-e[1]a[1]\beta=-h[1]\beta.
\] This proves commutativity of eight squares and anticommutativity of the lower right square. It also proves the six required distinguished-triangle assertions without assuming anything about the unsigned fourth ones.

Three rotations of \(T\) give \((X[1],Y[1],Z[1],-f[1],-g[1],-h[1])\). The three isomorphisms \((1,-1,1)\) turn this into \((X[1],Y[1],Z[1],f[1],g[1],-h[1])\). This proves the precise statement about the fourth row; apply the same construction to the first column for the fourth column. It is a change of sign on the displayed last arrow,\allowbreak{} not an unrecorded change of shift convention.

\subsubsection{Counterexample to the unsigned fourth-row assertion}\label{proofs-0001-1140-counterexample-to-the-unsigned-fourth-row-assertion}

Use the source\textquotesingle s homotopy category of complexes of abelian groups. Its cone convention is original 2470–\allowbreak{}2490,\allowbreak{} and its distinguished cone triangle is \((K,L,C(f),f,i,-p)\),\allowbreak{} explicitly stated and proved at 2995–\allowbreak{}3026 and 3319–\allowbreak{}3386. In particular the final sign is \textbf{minus} the cone projection; we retain it.

Let \(X=Y=\mathbb Z[0]\) and let \(f=3:X\to Y\). Its cone has \(C^{-1}=\mathbb Z\),\allowbreak{} \(C^0=\mathbb Z\),\allowbreak{} differential \(d_C^{-1}=3\),\allowbreak{} and all other terms and differentials zero. Let \(i:Y\to C\) be identity in degree zero and zero in all other degrees,\allowbreak{} and \(p:C\to X[1]\) be identity in degree minus one and zero elsewhere. Then \[
T=(X,Y,C,3,i,-p)
\] is distinguished. Its correctly signed shift is \((X[1],Y[1],C[1],3[1],i[1],p[1])\). Suppose the unsigned shift \((X[1],Y[1],C[1],3[1],i[1],-p[1])\) were also distinguished. TR3 applied to identity maps on its first two objects would supply \(c:C[1]\to C[1]\) satisfying,\allowbreak{} in the homotopy category,\allowbreak{} \[
ci[1]=i[1],\qquad p[1]c=-p[1].
\]

The shifted cone has terms \(\mathbb Z\) in degrees \(-2,-1\) and differential \(-3\). A chain endomorphism has integer components \(a,b\) with \(-3a=-3b\),\allowbreak{} hence \(a=b=m\). Its only possibly nonzero homotopy component is \(H^{-1}:\mathbb Z\to\mathbb Z\),\allowbreak{} an integer \(t\). Both components of \(dH+Hd\) equal \(-3t\). Thus endomorphism homotopy classes are integers modulo 3.

A chain map \(Y[1]\to C[1]\) has its sole component in degree \(-1\); homotopies change it by \(-3t\). Therefore the first displayed equality requires \(m\equiv1\pmod3\). A chain map \(C[1]\to X[2]\) has its sole component in degree \(-2\); homotopies again change it by \(-3t\). The second displayed equality requires \(m\equiv-1\pmod3\). These requirements contradict one another. Hence the unsigned shifted triangle is not distinguished.

This also defeats an existential reading in which the 3×3 construction may choose a different completion of \(f\):\allowbreak{} every such completion is isomorphic to \(T\) with identity maps on \(X,Y\). Shifting this isomorphism is an isomorphism of the unsigned shifted sextuples. If any of them were distinguished,\allowbreak{} TR1 would make the unsigned shift of \(T\) distinguished,\allowbreak{} contrary to the calculation. For example,\allowbreak{} take as the initial square the identity comparison of \(3:X\to Y\) with itself.

\subsubsection{Receiving uses of the corrected proposition}\label{proofs-0001-1140-receiving-uses-of-the-corrected-proposition}

The targeted search of top-level current chapter TeX found six uses. Each needs the first three rows and columns,\allowbreak{} which remain proved above. This checks propagation of the sign clarification,\allowbreak{} not every independent claim of those chapters.

* Original Derived Categories 1979–\allowbreak{}1995:\allowbreak{} the two first column cones \(X'',Y''\) lie in the chosen full triangulated subcategory; the third row puts \(Z''\) there. The third column then identifies the required denominator. No unsigned fourth triangle is used. * Original Derived Categories 11774–\allowbreak{}11810:\allowbreak{} the first two row cones \(Q_1,Q_2\) lie in the right orthogonal; the third column puts \(Q_3\) there,\allowbreak{} while the first column puts \(A_3\) in the original subcategory. Rotations for the other cases use TR2 with its sign. The fourth row is a display of the shifted maps only. * Current \texttt{cotangent.\allowbreak{}tex} 2831–\allowbreak{}2866 and 4591–\allowbreak{}4631:\allowbreak{} orient the upward diagrams with their bottom row first. The initial square has \(X=0\),\allowbreak{} \(Y=L_{B/R}\otimes_B^{\mathbf L}(A\otimes_RB)\),\allowbreak{} \(X'=L_{A/R}\otimes_A^{\mathbf L}(A\otimes_RB)\),\allowbreak{} \(Y'=L_{A\otimes_RB/R}\) in the ring case. The two first columns and rows are the stated fundamental triangles. The third column and third row supply the two triangles ending in \(E\). For spaces replace these exact entries by \(0,Lq^*L_{Y/B},Lp^*L_{X/B},
L_{X\times_BY/B}\). These are the displayed original objects; neither use requires the fourth row or column to be unsigned distinguished. * Current \texttt{more-\allowbreak{}algebra.\allowbreak{}tex} 30690–\allowbreak{}30709:\allowbreak{} use the square \(F\to M\) with vertical multiplication by the given ring element \(f\). Its row cone is \(M_f\),\allowbreak{} and multiplication on \(M_f\) is invertible,\allowbreak{} so the third column cone is zero. The third row therefore identifies the cones of \(f:F\to F\) and \(f:M\to M\). The fourth-triangle claim plays no role. * Current \texttt{derham.\allowbreak{}tex} 2979–\allowbreak{}3008:\allowbreak{} its displayed two-row comparison is extended by column cones. The argument identifies the middle and right column cones using the preceding comparison result,\allowbreak{} so the third row makes the left column cone zero. Again this only uses the first three triangles in each direction. The separate cohomological comparison result is not re-proved or newly certified by this bounded sign check.

No receiving source edit is required by this sign correction.

\subsection{\texorpdfstring{4. DERIVED-UNDER-001:\allowbreak{} the exact ambiguity of a triangle completion}{4.  the exact ambiguity of a triangle completion}}\label{proofs-0001-1140-derived-under-001-the-exact-ambiguity-of-a-triangle-completion}

The original square-zero assertion at 483–\allowbreak{}511 concerns endomorphisms of one triangle. Its same proof works for composable morphisms between different distinguished triangles. The precise statement is as follows.

For distinguished triangles \(T,T'\),\allowbreak{} define the abelian subgroup \[
I(T,T')=\{k:Y\to Y'\mid kf=0,\ g'k=0\}.
\] By covariant and contravariant representable exactness,\allowbreak{} \[
I(T,T')=
\operatorname{im}\bigl(f'\circ-:\operatorname{Hom}(Y,X')
\to\operatorname{Hom}(Y,Y')\bigr)
\ \cap\
\operatorname{im}\bigl(-\circ g:\operatorname{Hom}(Z,Y')
\to\operatorname{Hom}(Y,Y')\bigr).
\] This is an equality of subgroups of the displayed original Hom group. Indeed,\allowbreak{} \(g'k=0\) is exactly the first-image condition and \(kf=0\) is exactly the second. No choice of factorization is claimed unique.

If \(k\in I(T,T')\) and \(l\in I(T',T'')\),\allowbreak{} choose \(k=f'\alpha\) and \(l=\beta g'\). Then \[
lk=\beta g'f'\alpha=0.
\] These subgroups form an ideal in the additive category of distinguished triangles:\allowbreak{} a map in the ideal is exactly a triple \((0,k,0)\),\allowbreak{} and composing it on either side with any triangle morphism still has first and third components zero. Its square is zero by the displayed calculation. This proves the categorical strengthening with no octahedral axiom.

Given \(a:X\to X'\),\allowbreak{} \(c:Z\to Z'\) satisfying \(a[1]h=h'c\),\allowbreak{} twice rotate the two triangles and use TR3 on \(c:Z\to Z'\) and \(a[1]:X[1]\to X'[1]\). The resulting third map is \(b[1]\) for a unique \(b:Y\to Y'\),\allowbreak{} since the shift is fully faithful. The equations with \(-f[1]\) and \(-g[1]\) give \(bf=f'a\) and \(g'b=cg\). Thus completions exist. The difference of two completions belongs to \(I(T,T')\),\allowbreak{} and adding any element of that subgroup preserves all three triangle equations. The set of completions is therefore a torsor for this explicitly defined group. In particular,\allowbreak{} the exact uniqueness condition is \(I(T,T')=0\); the five vanishing conditions in 558–\allowbreak{}613 are sufficient conditions for that equality,\allowbreak{} not necessary ones.

If \(b_0\) is an isomorphism completion and \(b=b_0+k\) is another,\allowbreak{} the element \(t=b_0^{-1}k\) is in \(I(T,T)\),\allowbreak{} so \(t^2=0\). Consequently the exact inverse is \[
b^{-1}=(1_Y-b_0^{-1}k)b_0^{-1}.
\] Both products with \(b=b_0(1_Y+t)\) equal the appropriate identity. This describes the nonuniqueness of isomorphism completions without claiming functorial choices of cones.

For the original idempotent lift,\allowbreak{} let \(b'\) complete idempotents \(a,c\). The defect \(d=(b')^2-b'\) lies in \(I(T,T)\),\allowbreak{} hence \(d^2=0\),\allowbreak{} and it commutes with \(b'\) because it is a polynomial in \(b'\). Retain the original formula \[
b=b'-(2b'-1)d=3(b')^2-2(b')^3.
\] The correction term remains in the ideal,\allowbreak{} so \((a,b,c)\) is a triangle endomorphism. Expanding the polynomial gives exactly \[
b^2-b=(4(b')^2-4b'-3)d^2=0.
\] Thus the broader ideal description recovers the original formula and all of its integer coefficients rather than replacing it with an abstract existence assertion. Later uses at the categorical idempotent-completion and uniqueness loci remain to be checked in source order; none is pre-certified by its matching label.

\subsection{\texorpdfstring{5. DERIVED-UNDER-002:\allowbreak{} propagate the countable-power refinement}{5.  propagate the countable-power refinement}}\label{proofs-0001-1140-derived-under-002-propagate-the-countable-power-refinement}

The previously published \href{https://github.com/KokunoYumeto/unofficial-stacks-project-ai-drafts/blob/95a51593aceb247ac7a183111678c3a0f8b8f4b5/ai-integrated/candidates/commons/stacks/errata/r60/evidence/review/PROOFS_0001_1900.md}{Homology editorial argument,\allowbreak{} §1} proves HOMOLOGY-UNDER-001. Its receiving lemma here is original 805–\allowbreak{}822. In a pre-triangulated category every epimorphism has a zero cokernel and hence has a kernel by the kernel–\allowbreak{}cokernel equivalence proved in §1. In particular this applies to each split epimorphism used in the countable-product construction.

For a fixed object \(X\) whose countable power \(P=\prod_{n\ge1}X\) exists,\allowbreak{} and a fixed idempotent \(e:X\to X\),\allowbreak{} keep the actual coordinate maps \(p_n:P\to X\) and put \[
p_n\Phi=e p_n+(1_X-e)p_{n+1},\quad
p_1\Psi=p_1,\quad
p_n\Psi=(1_X-e)p_{n-1}+e p_n\quad(n\ge2).
\] For the first coordinate,\allowbreak{} \(p_1\Phi\Psi=e p_1+(1-e)((1-e)p_1+e p_2)=p_1\). For \(n\ge2\),\allowbreak{} \[
p_n\Phi\Psi=e((1-e)p_{n-1}+e p_n)
 +(1-e)((1-e)p_n+e p_{n+1})=p_n.
\] Therefore \(\Phi\Psi=1_P\),\allowbreak{} and \(\Phi\) has a kernel \(k:K\to P\) by the preceding pre-triangulated argument. Set \(i=p_1k\). Multiplying \(p_n\Phi k=0\) by \(e\) and \(1-e\) gives \(e p_nk=0\) and \((1-e)p_{n+1}k=0\). Thus \(p_nk=0\) for \(n\ge2\),\allowbreak{} while \(ei=0\). If \(t:W\to X\) has \(et=0\),\allowbreak{} the product map \(v:W\to P\) with coordinates \((t,0,0,\ldots)\) satisfies \(\Phi v=0\). Its unique factorization through \(k\) gives a unique factorization of \(t\) through \(i\),\allowbreak{} since the remaining coordinates of \(k\) are zero. Hence \(i\) is a kernel of \(e\). Apply the same construction to \(1-e\) to obtain \(j:J\to X\).

Because \(e(1-e)=0\),\allowbreak{} there is \(r:X\to K\) with \(ir=1-e\). Similarly there is \(s:X\to J\) with \(js=e\). Monicity gives \(ri=1_K\),\allowbreak{} \(sj=1_J\),\allowbreak{} \(rj=0\),\allowbreak{} \(si=0\),\allowbreak{} and \(ir+js=1_X\). These are inverse biproduct maps \((i,j):K\oplus J\to X\),\allowbreak{} \((r,s):X\to K\oplus J\),\allowbreak{} with the original \(e\) acting as zero on \(K\) and identity on \(J\).

Thus a countable power of this individual object suffices to split every one of its idempotents. Countable powers for all objects suffice for the category to be Karoubian; products of differing objects are unnecessary. The corresponding result for a countable coproduct follows by applying the proved construction to the opposite category,\allowbreak{} which exchanges all the displayed domains,\allowbreak{} codomains,\allowbreak{} products and kernels.

This also supplies an exact consequence for the already included Verdier lemma \texttt{lemma-\allowbreak{}functorial-\allowbreak{}triangles-\allowbreak{}decomposable} (current lines 824–\allowbreak{}880)\allowbreak{}. Its countable-products-or-coproducts assumption may be replaced by the assumption that the category is Karoubian. Here is the complete argument,\allowbreak{} including the naturality needed by its Yoneda step.

Let \(T\) be its section of the first-arrow functor,\allowbreak{} and retain \(T(f)=(X,Y,Z_f,f,g_f,h_f)\),\allowbreak{} \(d_f=-h_f[-1]:Z_f[-1]\to X\). The subfunctor \(\mathcal K_f(W)=\{v:W\to X\mid fv=0\}\) is the image of the natural surjection \(q_f:\operatorname{Hom}(-,Z_f[-1])\to\mathcal K_f\),\allowbreak{} \(q_f(t)=d_ft\),\allowbreak{} by representable exactness. Write \(T(W\to0)=(W,0,Z_W,0,0,h_W)\); exactness or two rotations gives that \(h_W:Z_W\to W[1]\) is an isomorphism. For \(v\in\mathcal K_f(W)\),\allowbreak{} the square \((v,0)\) yields a triangle map with third arrow \(c_v:Z_W\to Z_f\). Put \(s_W(v)=-c_v[-1]h_W[-1]^{-1}\). Then \(q_fs_W(v)=v\).

If \(u:W'\to W\),\allowbreak{} functoriality of \(T\) gives an arrow \(c_u:Z_{W'}\to Z_W\) with \(c_{vu}=c_vc_u\) and \(h_Wc_u=u[1]h_{W'}\). Shifting that equality and cancelling the two isomorphisms gives \(c_u[-1]h_{W'}[-1]^{-1}=h_W[-1]^{-1}u\),\allowbreak{} so \(s_{W'}(vu)=s_W(v)u\). Hence \(s\) is natural,\allowbreak{} even without assuming the chosen section \(T\) additive.

The natural endomorphism \(sq_f\) of the representable functor is induced by a morphism \(e_f:Z_f[-1]\to Z_f[-1]\) by Yoneda,\allowbreak{} and \((sq_f)^2=sq_f\) gives \(e_f^2=e_f\). Choose a splitting \(e_f=ip\),\allowbreak{} \(pi=1_K\),\allowbreak{} by the Karoubian assumption. Since \(q_fe_f=q_f\),\allowbreak{} the map \(d_fi:K\to X\) is a kernel of \(f\):\allowbreak{} for \(v\in\mathcal K_f(W)\),\allowbreak{} its factor is \(p s_W(v)\),\allowbreak{} and \(d_fip s_W(v)=q_fe_fs_W(v)=v\). If \(d_fit=v\),\allowbreak{} naturality and \(sq_f=e_f\) imply \(it=e_f it=s_W(v)\),\allowbreak{} hence \(t=p s_W(v)\). The kernel criterion of §1 now gives the full projection–\allowbreak{}coprojection decomposition of \(f\).

The earlier Verdier insertion remains byte-for-byte intact. These weaker hypotheses and exact maps are separate editorial consequences,\allowbreak{} with HOMOLOGY-UNDER-001 recorded as the input result. They do not claim a new reading or wholesale acceptance of the other queued Verdier candidates.
\clearpage
\section{\texorpdfstring{Derived Categories:\allowbreak{} localization and the first kernel lemmas}{Derived  localization and the first kernel lemmas}}\label{proofs-1149-1843-derived-categories-localization-and-the-first-kernel-lemmas}

Authority and comparison identities are those of \texttt{INPUT\_\allowbreak{}BINDING.\allowbreak{}json}. All unqualified Derived Categories line numbers refer to the Stacks Project authors\textquotesingle{} frozen source at \texttt{a04446e5\allowbreak{}7ec1fbc2\allowbreak{}52a871af\allowbreak{}cec7752f\allowbreak{}b2807b14}. Reading now reaches original line 1860; the boundedness definition beginning at 1845 continues past that point and is not included in this completed review interval. These are complete editorial arguments for the listed decisions,\allowbreak{} not a replacement edition of the chapter. No novelty claim is made.

\subsection{1. The localization axioms and the earlier factorization repair}\label{proofs-1149-1843-the-localization-axioms-and-the-earlier-factorization-repair}

Keep the source MS5 condition \(f\in S\iff f[1]\in S\),\allowbreak{} its MS6 condition for the third map of triangles,\allowbreak{} and both left and right fraction axioms. MS6 and identity denominators imply that every isomorphism lies in \(S\):\allowbreak{} apply MS6 to the two split triangles \((0,X,X,0,1,0)\),\allowbreak{} \((0,X,Y,0,f,0)\). The third map is forced to be \(f\). Composition with isomorphisms therefore preserves \(S\),\allowbreak{} by MS1,\allowbreak{} and comparison isomorphisms between iterated shifts do not alter membership in \(S\). In particular,\allowbreak{} for \(s\in S\),\allowbreak{} its shift \(s[-1][1]\) is conjugate to \(s\) by the specified natural isomorphisms,\allowbreak{} so \(s[-1][1]\in S\) and MS5 gives \(s[-1]\in S\). Induction yields stability under every integral shift.

For LMS2,\allowbreak{} start with \(f:X\to Y\),\allowbreak{} \(t:X\to X'\) in \(S\),\allowbreak{} and a distinguished \((X,Y,Z,f,g,h)\). Complete \(t[1]h:Z\to X'[1]\),\allowbreak{} rotating as necessary,\allowbreak{} to a triangle \((X',Y',Z,f',g',t[1]h)\). Twice rotate both triangles. Their first two vertical maps are \(1_Z,t[1]\),\allowbreak{} which belong to \(S\); MS6 gives a third map \(s'[1]\in S\). Full faithfulness of the shift and MS5 give \(s':Y\to Y'\) in \(S\) with \(s'f=f't\). The reversed construction proves RMS2. This supplies the exact shift and denominator argument at 1187–\allowbreak{}1253.

For an exact functor \((F,\xi)\),\allowbreak{} the class of arrows inverted by \(F\) satisfies MS1. It satisfies the source\textquotesingle s MS4:\allowbreak{} if \(F(fg)\) and \(F(gh)\) are invertible,\allowbreak{} then \(F(g)\) has a left inverse and a right inverse,\allowbreak{} which coincide. MS5 follows from \[
F(f[1])=\xi_Y^{-1}F(f)[1]\xi_X.
\] For MS6 choose a third map using TR3,\allowbreak{} apply the exact functor with its actual translation comparison,\allowbreak{} and use the two-isomorphisms criterion on the image triangles. The preceding paragraph gives MS2. For LMS3 retain \(a=f-g\),\allowbreak{} a triangle \((Z,X,Q,t,d,h)\),\allowbreak{} and the factorization \(a=id\) supplied by contravariant representable exactness. Completing \(i:Q\to Y\) gives \((Q,Y,W,i,j,k)\). Since \(F(t)\) is invertible,\allowbreak{} \(F(Q)=0\). The image of the latter triangle shows \(F(j)\) invertible. Thus \(j\in S\) and \(ja=jid=0\),\allowbreak{} as required. Reverse all arrows for RMS3.

The same reasoning applies to the arrows inverted by every \(H^n\) for a homological functor \(H\). Shift comparisons identify the condition at \(f[1]\) with the condition at all consecutive \(H^n(f)\). The long exact sequences give MS6 by the five lemma. For LMS3 retain the source notation at 1334–\allowbreak{}1355:\allowbreak{} \[
f,g:X\to Y,\quad a=f-g,\quad t:Z\to X,\quad
(Z,X,Q,t,u,v),\quad i:Q\to Y,\quad (Q,Y,W,i,j,k).
\] Since \(at=0\),\allowbreak{} exactness gives \textbf{\(iu=a\)}. For every integer \(n\),\allowbreak{} the exact segment \[
H^n(Z)\xrightarrow{H^n(t)}H^n(X)\longrightarrow H^n(Q)
\longrightarrow H^{n+1}(Z)\xrightarrow{H^{n+1}(t)}H^{n+1}(X)
\] has both displayed \(t\)-maps invertible,\allowbreak{} so \(H^n(Q)=0\). The long exact sequence for the second triangle now makes each \(H^n(j)\) invertible. Finally \(ja=jiu=0\).

This proves the semantic match of \textbf{OCC-01492 /\allowbreak{} P02-E0020} with \textbf{MC-STK-ERR-0812},\allowbreak{} already present in the current chapter. Its historical proof annotation calls \(g\) undefined. That explanation is inaccurate:\allowbreak{} \(g:X\to Y\) was declared at 1334. It is the wrong map for this composition; \(i\) has domain \(Q\),\allowbreak{} whereas \(u:X\to Q\) is the triangle arrow giving the factorization. The exact prior edit is correct and preserved. This note corrects the explanation without rewriting the immutable earlier receipt or creating a duplicate edit.

\subsection{2. The localization triangle morphism and two earlier copyedits}\label{proofs-1149-1843-the-localization-triangle-morphism-and-two-earlier-copyedits}

For the construction at 1394–\allowbreak{}1442 keep the original common-target diagram and all of its arrows:\allowbreak{} \[
f:X\to Y,\quad f':X'\to Y',\quad f'':X''\to Y'',\quad
k:X\to X'',\quad l:Y\to Y'',\quad s:X'\to X'',\quad t:Y'\to Y'',
\] where \(s,t\in S\),\allowbreak{} \(f''k=lf\),\allowbreak{} \(f''s=tf'\),\allowbreak{} \(a=Q(s)^{-1}Q(k)\),\allowbreak{} and \(b=Q(t)^{-1}Q(l)\). TR3 supplies \((k,l,m)\) to a completion \((X'',Y'',Z'',f'',g'',h'')\). MS6 supplies \((s,t,r)\) from the primed triangle,\allowbreak{} with \(r\in S\). Consequently \[
mg=g''l,\qquad h''m=k[1]h,
\] and \[
rg'=g''t,\qquad h''r=s[1]h'.
\] The first two arrows\textquotesingle{} equation is already \(f''k=lf\). The middle correction \textbf{MC-STK-ERR-0813} labels the fourth vertical arrow \(k[1]:X[1]\to X''[1]\),\allowbreak{} exactly as required; \(g[1]\) has the wrong source and target. This is the semantic match for \textbf{OCC-01493 /\allowbreak{} P02-E0021}.

Put \(c=Q(r)^{-1}Q(m)\). Applying \(Q\) to the displayed identities gives \[
cQ(g)=Q(r)^{-1}Q(g'')Q(l)=Q(g')Q(t)^{-1}Q(l)=Q(g')b.
\] Similarly,\allowbreak{} since the localized shift has identity comparison with \(Q\),\allowbreak{} \[
Q(h')c=Q(s[1])^{-1}Q(h'')Q(m)
=Q(s[1])^{-1}Q(k[1])Q(h)=a[1]Q(h).
\] Along with \(bQ(f)=Q(f')a\),\allowbreak{} these are all three equations of the localized triangle morphism. No map has been identified merely from a matching numerical locus.

\textbf{MC-STK-ERR-0814} at 1445 correctly replaces ``lemma'' by “proposition”:\allowbreak{} the enclosing environment is the proposition beginning at 1359. No received physical report in this interval corresponds to that older copyedit,\allowbreak{} but the unit itself has been reconciled.

\textbf{OCC-01494 /\allowbreak{} P02-E0022} matches \textbf{MC-STK-ERR-0815} at 1521:\allowbreak{} the duplicated ``is'' is already removed. Here is the required mathematical context. For the full pre-triangulated subcategory \(\mathcal D'\),\allowbreak{} the restricted denominators \(S'\) are the arrows inverted by \(Q|_{\mathcal D'}\),\allowbreak{} because \(S\) is saturated. Section 1 therefore gives their multiplicative-system and triangle compatibility properties. Every object \(X\) receives a denominator \(X'\to X\) from \(\mathcal D'\),\allowbreak{} giving essential surjectivity of the localized inclusion. For a roof \(A\leftarrow X\to B\),\allowbreak{} with \(A,B\in\mathcal D'\),\allowbreak{} choose a denominator \(X'\to X\) with \(X'\in\mathcal D'\). Composing gives an equivalent roof wholly inside the full subcategory,\allowbreak{} proving fullness. If two such roofs are equal in the ambient localization,\allowbreak{} choose a common right refinement on which their numerator maps agree. Apply the same density assumption to its source; the resulting refinement lies in \(\mathcal D'\),\allowbreak{} its denominators lie in \(S'\),\allowbreak{} and the numerators remain equal. This proves faithfulness as well as the cofinality implicit in the source\textquotesingle s colimit argument.

\subsection{\texorpdfstring{3. DERIVED-NEW-005:\allowbreak{} fix the comparison in the uniqueness assertion}{3.  fix the comparison in the uniqueness assertion}}\label{proofs-1149-1843-derived-new-005-fix-the-comparison-in-the-uniqueness-assertion}

At 1363–\allowbreak{}1365 the uniqueness assertion requires that the localization functor\textquotesingle s given translation comparison be the identity under \([1]Q=Q[1]\). An exact functor in the chapter\textquotesingle s definition includes that comparison; equality of the two composite functors does not itself force every natural automorphism between them to be the identity. The intended structure in the proof has the identity comparison.

With it fixed,\allowbreak{} the proof is exact. MS5 lets the shift act on a left fraction by \[
\bigl(Q(s)^{-1}Q(f)\bigr)[1]=Q(s[1])^{-1}Q(f[1]).
\] The localization universal property makes this well-defined and unique; the inverse shift descends in the same way. Canonical distinguished triangles are the isomorphic images of distinguished triangles of \(\mathcal D\). If a second pre-triangulated structure on this fixed localized additive category and shift makes \((Q,1)\) exact,\allowbreak{} it contains all these canonical triangles. Any distinguished triangle of the second structure has some first arrow \(a\). A canonical completion of \(a\) exists by the source\textquotesingle s TR1 construction. Both triangles now belong to the second structure and share that first arrow,\allowbreak{} so TR3 and representable exactness give an isomorphism between them with identity first two components. Isomorphism closure of the canonical class puts the original triangle in the canonical class. Thus the two classes agree. This proves uniqueness with the exact comparison specified,\allowbreak{} rather than adding an unproved uniqueness claim.

To see the limitation of leaving the comparison free,\allowbreak{} let \((\mathcal D,[1],\mathcal T)\) be a triangulated category and define \[
\mathcal T^-=
\{(X,Y,Z,f,g,-h):(X,Y,Z,f,g,h)\in\mathcal T\}.
\] This is another triangulation on the same additive category and shift. TR1 follows directly. For TR2,\allowbreak{} the rotation of \((X,Y,Z,f,g,-h)\) is \((Y,Z,X[1],g,-h,-f[1])\). It is in \(\mathcal T^-\) exactly when \((Y,Z,X[1],g,-h,f[1])\) is in \(\mathcal T\). The latter is isomorphic via \((1,1,-1)\) to \((Y,Z,X[1],g,h,-f[1])\); TR2 for \(\mathcal T\) gives the required equivalence in both directions. For TR3,\allowbreak{} negating both boundary maps leaves the third commuting-square equation equivalent to the original one,\allowbreak{} so the same third map works.

For TR4 start with three \(\mathcal T^-\)-triangles with boundary maps \(d_1,d_2,d_3\). Their corresponding \(\mathcal T\)-triangles have boundaries \(-d_1,-d_2,-d_3\). TR4 in \(\mathcal T\) provides \(a,b\) with connecting arrow \(-p_1[1]d_3\). Negating that arrow gives the required \(\mathcal T^-\)-triangle with connecting arrow \(p_1[1]d_3\). The other two triangle-morphism conditions still hold because both boundaries in each equation were negated. This proves all four axioms.

Now \((\operatorname{id},-1_{[1]}):
(\mathcal D,\mathcal T)\to(\mathcal D,\mathcal T^-)\) is an exact functor and commutes strictly with the shift as a functor. Take \(S\) to be all isomorphisms,\allowbreak{} so that the identity is a model of the localization functor. Both structures satisfy an assertion requiring only an unspecified exact structure on that functor. They need not coincide. In the cone example of \texttt{PROOFS\_\allowbreak{}0001\_\allowbreak{}1140.\allowbreak{}md},\allowbreak{} §3,\allowbreak{} the triangle \((\mathbb Z[0],\mathbb Z[0],C(3),3,i,-p)\) belongs to \(\mathcal T\),\allowbreak{} whereas \((\mathbb Z[0],\mathbb Z[0],C(3),3,i,p)\) does not:\allowbreak{} comparison with identity first two maps would require an integer endomorphism class \(m\) of the cone to satisfy both \(m\equiv1\pmod3\) and \(m\equiv-1\pmod3\). The unshifted homotopy differential is \(3\),\allowbreak{} so homotopies change components by \(3t\),\allowbreak{} giving exactly those congruences. This proves distinctness. Transport along the canonical equivalence gives the same example on the particular fraction model of \(S^{-1}\mathcal D\).

This is a concrete consequence of the preceding sign finding and the chapter\textquotesingle s own definition of an exact functor. The repair makes explicit the intended comparison; it does not change the constructed localization.

The four external current-chapter citations to this proposition were checked at \texttt{adequate.\allowbreak{}tex:\allowbreak{}2202–\allowbreak{}2227},\allowbreak{} \texttt{dga.\allowbreak{}tex:\allowbreak{}2568–\allowbreak{}2590},\allowbreak{} \texttt{perfect.\allowbreak{}tex:\allowbreak{}9810–\allowbreak{}9840},\allowbreak{} and \texttt{sdga.\allowbreak{}tex:\allowbreak{}4187–\allowbreak{}4212}. They use existence of the canonical localization or its exact universal factorization,\allowbreak{} which uses precisely \((Q,1)\). None needs uniqueness with a free translation comparison. The Perfect Complexes passage also uses the reflection lemma; its exact functor retains its own specified comparison as in DERIVED-NEW-002. No change to these receiving passages is required. Later internal citations are still checked in chapter order.

The exact structure on the factorization \(F'\) at 1483–\allowbreak{}1488 also has a definite receiving map. Since localization has the same objects,\allowbreak{} put \(\xi'_{Q(X)}=\xi_X\). Naturality on \(Q(f)\) is the original naturality,\allowbreak{} and naturality on inverses \(Q(s)^{-1}\) follows by multiplying that same commuting equation by inverses. It therefore holds on every fraction. The images of canonical distinguished triangles are distinguished,\allowbreak{} so \((F',\xi')\) is exact. This proves the relevant part of the omitted universal-property argument with the translation maps exposed.

\subsection{4. Kernel witnesses and filtered triangle denominators}\label{proofs-1149-1843-kernel-witnesses-and-filtered-triangle-denominators}

There is no new error at 1574–\allowbreak{}1577. Given the zero denominator \(0:Z\to Z'\),\allowbreak{} the source\textquotesingle s split triangle \[
(Z'[-1],Z'[-1]\oplus Z,Z,\iota,\pi,0)
\] rotates \textbf{backwards} to \[
(Z[-1],Z'[-1],Z'[-1]\oplus Z,-0[-1],\iota,\pi).
\] The first arrow is in \(S\) by MS5 and closure under comparison isomorphisms; its minus sign does not change a zero map. The third object is isomorphic,\allowbreak{} by the actual biproduct interchange map,\allowbreak{} to \(Z\oplus Z'[-1]\). This is exactly condition (4)\allowbreak{},\allowbreak{} using \(Z'[-1]\) as the complementary object. A forward rotation alone would expose a shifted copy of \(Z\); the inverse rotation proves the stated unshifted result. Conversely,\allowbreak{} if \(f\in S\) and its triangle cone is \(Z\oplus Z'\),\allowbreak{} exactness of \(Q\) gives \(Q(Z)\oplus Q(Z')=0\). Each summand\textquotesingle s identity factors through the zero object,\allowbreak{} so each summand is zero. For saturated \(S\),\allowbreak{} the two zero-object arrows are denominators exactly when their localized images are invertible,\allowbreak{} giving conditions (5)\allowbreak{} and (6)\allowbreak{}; the triangle \((0,Z,Z,0,1,0)\) gives condition (7)\allowbreak{}.

For the category \(\mathcal I\) at 1595–\allowbreak{}1735,\allowbreak{} an object is the actual denominator triangle map \(s:T\to T'\),\allowbreak{} with all three components in \(S\); a morphism is a triangle map commuting with these structure maps. Every component of such a morphism is in \(S\):\allowbreak{} its image under \(Q\) is the quotient of two invertible structure maps,\allowbreak{} and \(S\) is saturated. Thus subsequent denominator completions used in the proof really have the required membership.

MS2 followed by MS6 completes any prescribed denominator \(X\to X'\) to an object of \(\mathcal I\). Rotation gives the same assertion for either other component. Applied to a target triangle,\allowbreak{} this also completes any subsequent denominator on one component to a triangle map with all three components in \(S\). The latter assertion permits successive replacements of a target triangle; it does not prescribe its two other new objects in advance.

For two objects \(s_1:T\to T_1\),\allowbreak{} \(s_2:T\to T_2\),\allowbreak{} filteredness of each component category \(X/S,Y/S,Z/S\) and these completions give a common later target and three maps \(a,b,c\) from \(T_1\) to that target,\allowbreak{} commuting with the three maps out of \(T\). Each remaining triangle-square defect becomes zero after \(Q\):\allowbreak{} both triples there are the quotient of the same two structure isomorphisms. The fraction equality criterion gives a denominator out of the target component that kills one defect. Complete it to a map of target triangles,\allowbreak{} postcompose the triple,\allowbreak{} and do this for each of its other two components,\allowbreak{} using rotations when the defect lands in a shifted component. Once a square commutes,\allowbreak{} postcomposition by a triangle map preserves its commutativity. After these three steps all squares commute,\allowbreak{} giving the required common object of \(\mathcal I\).

For two parallel arrows in \(\mathcal I\),\allowbreak{} their first components agree after composition with the first denominator of their source. LMS3 gives a target denominator making those components equal. Complete it to a map of triangles and then do the same for the second and third components. Equalities already obtained survive composition,\allowbreak{} so the resulting two triples are equal. The identity of \(T\) supplies a nonempty object. These arguments verify all three filteredness axioms.

Finally let \(u:X\to U\) be an object of \(X/S\),\allowbreak{} with two arrows to the first components of two objects of \(\mathcal I\). Choose a common target in \(\mathcal I\). The two resulting arrows from \(U\) to its first component can be equalized in the filtered category \(X/S\). Complete that equalizing denominator to one more target triangle map. This supplies a common later object with equal induced arrows out of \(U\),\allowbreak{} giving a connected comma category. The earlier completion gives its nonemptiness. This proves cofinality in the source\textquotesingle s exact definition; rotations prove the other two cofinality assertions. No claim that a surjective-on-objects functor alone is cofinal is used.

\subsection{\texorpdfstring{5. DERIVED-NEW-006:\allowbreak{} propagate the image-triangle comparison}{5.  propagate the image-triangle comparison}}\label{proofs-1149-1843-derived-new-006-propagate-the-image-triangle-comparison}

The same typed-boundary repair as DERIVED-NEW-002 applies at 1787 and 1803. Name the exact functor \((F,\xi)\),\allowbreak{} and write the image triangle with boundary \(\xi_XF(h)\). The kernel conclusion stays unchanged:\allowbreak{} naturality and invertibility of \(\xi\) carry zero objects through both shifts; the image triangle with its first two objects zero has third object zero by representable exactness. Fullness of the kernel subcategory and TR3 supply all of its triangle maps; the cone criterion then makes it pre-triangulated,\allowbreak{} and TR4 restricts when available in the ambient category. If an object \(X\oplus Y\) maps to zero,\allowbreak{} additivity of \(F\) gives \(F(X)\oplus F(Y)=0\); projection and inclusion show the identities of both summands are zero. Hence the kernel is saturated and strictly full. The same argument applied to every \(H^n\) proves the following homological-kernel lemma at 1812–\allowbreak{}1843,\allowbreak{} using its long exact sequence instead of an image triangle.

The new operations in this interval are one comparison qualification in the localization proposition and two explicit-map operations in the kernel lemma. Four earlier units remain unchanged; three further physical reports are accounted for by those earlier corrections.
\clearpage
\section{\texorpdfstring{Quotients,\allowbreak{} the retained Verdier insertion,\allowbreak{} and exact adjoints}{ the retained Verdier  and exact adjoints}}\label{proofs-1844-2401-quotients-the-retained-verdier-insertion-and-exact-adjoints}

Authority:\allowbreak{} the Stacks Project authors\textquotesingle{} \texttt{derived.\allowbreak{}tex},\allowbreak{} commit \texttt{a04446e5\allowbreak{}7ec1fbc2\allowbreak{}52a871af\allowbreak{}cec7752f\allowbreak{}b2807b14}; exact source and current comparison are bound in \texttt{INPUT\_\allowbreak{}BINDING.\allowbreak{}json}. Original line numbers are used except where current lines are explicitly indicated. This note preserves the source statements and records editorial arguments and generalizations separately. No novelty is claimed.

\subsection{1. Boundedness and the multiplicative system of a subcategory}\label{proofs-1844-2401-boundedness-and-the-multiplicative-system-of-a-subcategory}

For the bounded-below subcategory at 1845–\allowbreak{}1887,\allowbreak{} choose integers \(a,b\) with \(H^n(X)=0\) for \(n<a\),\allowbreak{} \(H^n(Y)=0\) for \(n<b\). In a distinguished triangle \(X\to Y\to Z\to X[1]\),\allowbreak{} the exact segment \(H^n(Y)\to H^n(Z)\to H^{n+1}(X)\) gives \(H^n(Z)=0\) for \(n<\min(b,a-1)\). The shift moves the bound by its integer exponent,\allowbreak{} so the class is stable under both shifts. Finite direct sums and direct summands retain a bound by the explicit biproduct identity \(H^n(X\oplus Y)\cong H^n(X)\oplus H^n(Y)\). For upper bounds,\allowbreak{} if the two groups vanish respectively for \(n>a\) and \(n>b\),\allowbreak{} then \(H^n(Z)=0\) for \(n>\max(b,a-1)\). Combining these two calculations handles bounded objects. The implication displayed at 1881–\allowbreak{}1886 is valid for every fixed \(n\); it does not assert that bounded-below objects vanish in every degree. No correction of that quantifier is needed.

Let \(\mathcal B\) be the given full triangulated subcategory and let \(S(\mathcal B)\) consist of arrows whose triangle cone is isomorphic to an object of \(\mathcal B\). Cones with the same first arrow are isomorphic by TR3 and representable exactness,\allowbreak{} so the membership test is independent of a chosen cone. The octahedron for \(X\xrightarrow fY\xrightarrow gZ\) supplies the original triangle \(Q_1\to Q_2\to Q_3\to Q_1[1]\); if \(Q_1,Q_3\) are isomorphic to objects of \(\mathcal B\),\allowbreak{} then so is \(Q_2\). This proves closure under composition. For LMS3,\allowbreak{} retain \(a:X\to Y\),\allowbreak{} \(t:Z\to X\),\allowbreak{} \(at=0\),\allowbreak{} the triangle \((Z,X,Q,t,g,h)\),\allowbreak{} and \(i:Q\to Y\) with \(ig=a\). In \((Q,Y,W,i,s,k)\),\allowbreak{} rotation makes the cone of \(s\) equal to \(Q[1]\),\allowbreak{} up to the chosen shift comparison. Thus \(s\in S\) and \(sa=sig=0\). Reversing the arrows proves RMS3.

MS5 uses the correctly signed shift of a triangle. Its cone is still the shifted original cone,\allowbreak{} so membership in \(\mathcal B\),\allowbreak{} up to isomorphism,\allowbreak{} is preserved in both directions. \textbf{OCC-02891 /\allowbreak{} P02-E1428} correctly requests the missing period at 1976; the mathematical sentence and its objects remain unchanged.

MS6 follows from the corrected 3×3 proposition:\allowbreak{} the first two column cones lie in \(\mathcal B\),\allowbreak{} the third row puts the third column cone there,\allowbreak{} and that third column gives the required denominator. Only the first three rows and columns are used. MS2 then follows by the explicit rotated construction in \texttt{PROOFS\_\allowbreak{}1149\_\allowbreak{}1843.\allowbreak{}md},\allowbreak{} §1.

For the saturation equivalence retain the original four objects and maps \(W\xrightarrow hX\xrightarrow gY\xrightarrow fZ\). Assume \(fg,gh\in S\). The five chosen cones are \(Q_1,Q_2,Q_3,Q_{12},Q_{23}\),\allowbreak{} with \(Q_{12},Q_{23}\) in \(\mathcal B\) up to isomorphism. The octahedral triangles have boundary maps \[
s:Q_2\longrightarrow X[1]\longrightarrow Q_1[1],\qquad
t:Q_3\longrightarrow Y[1]\longrightarrow Q_2[1].
\] Their cones are \(Q_{12}[1]\),\allowbreak{} \(Q_{23}[1]\),\allowbreak{} so \(s,t\in S\). Their shifted composite is the full original map \[
s[1]t:Q_3\longrightarrow Y[1]\longrightarrow Q_2[1]
\longrightarrow X[2]\longrightarrow Q_1[2].
\] It is zero because the middle two arrows are consecutive in the shifted triangle for \(g\). It is also in \(S\),\allowbreak{} by the composition and shift properties just proved. A cone of this zero arrow is \(Q_1[2]\oplus Q_3[1]\):\allowbreak{} take the direct sum of the two identity triangles and rotate with their signs. This cone belongs to \(\mathcal B\) up to isomorphism. Saturation of \(\mathcal B\) puts both summands,\allowbreak{} and hence \(Q_1,Q_3\),\allowbreak{} in it up to isomorphism. The triangle \(Q_1\to Q_{12}\to Q_2\to Q_1[1]\) then puts \(Q_2\) there. Thus \(g\in S\); the same cone test also gives \(h,f\in S\). \textbf{OCC-01495 /\allowbreak{} P02-E0023} therefore agrees with the already implemented \textbf{MC-STK-ERR-0816}:\allowbreak{} remove the extra ``is'' from ``It is also follows”,\allowbreak{} without weakening the parenthetical claim.

Conversely,\allowbreak{} if \(S\) is saturated and \(X\oplus Y\) belongs to \(\mathcal B\) up to isomorphism,\allowbreak{} use the three composable arrows \(Y[-1]\to0\to X\to X\oplus Y\). The two adjacent composites have cones isomorphic to \(X\oplus Y\),\allowbreak{} using the displayed split triangle and its inverse rotation. MS4 makes \(0\to X\) a denominator. Its cone is \(X\),\allowbreak{} which therefore belongs to \(\mathcal B\) up to isomorphism. Interchange the two summands for \(Y\). This proves the other implication without assuming strict fullness of the original subcategory.

\subsection{2. Quotient kernel and the retained Verdier quotient results}\label{proofs-1844-2401-quotient-kernel-and-the-retained-verdier-quotient-results}

The canonical quotient functor uses its identity translation comparison,\allowbreak{} as made explicit in DERIVED-NEW-005. If an exact functor annihilates \(\mathcal B\),\allowbreak{} its image of the cone triangle of a denominator has zero cone,\allowbreak{} so it inverts that denominator. For a homological functor whose entire shifted kernel contains \(\mathcal B\),\allowbreak{} the relevant long exact sequence gives the same conclusion. The ordinary fraction factorization is consequently defined,\allowbreak{} and its exact or homological structure descends as in \texttt{PROOFS\_\allowbreak{}1149\_\allowbreak{}1843.\allowbreak{}md},\allowbreak{} §3. No additional triangulation or choice of boundary comparison is hidden in this use.

The kernel of \(Q:\mathcal D\to\mathcal D/\mathcal B\) consists exactly of those \(Z\) for which \(Z\oplus Z'\) is isomorphic to an object of \(\mathcal B\) for some \(Z'\). Indeed,\allowbreak{} the localization kernel lemma supplies a denominator triangle with cone \(Z\oplus Z'\),\allowbreak{} giving the forward implication. Conversely the whole biproduct maps to zero and its inclusion and projection show that \(1_{Q(Z)}=0\). Every strictly full saturated triangulated subcategory containing \(\mathcal B\) contains all such summands. The kernel itself has those three properties,\allowbreak{} proving the minimality.

The operation \(S\mapsto\mathcal B(S)\mapsto S(\mathcal B(S))\) returns exactly the arrows inverted by the localization,\allowbreak{} by its image cone triangle and the zero-cone criterion. This is the saturation of \(S\). The opposite composite returns the strictly full saturated closure of \(\mathcal B\),\allowbreak{} by the preceding paragraph. The older \textbf{MC-STK-ERR-0817} correctly inserts ``a'' at 2204 in the description of this argument; it has no additional physical report in this range.

The existing Verdier insertion at \textbf{current lines 2287–\allowbreak{}2487} was read in full. It is retained unchanged. Here is its source-level comparison and its exact receiving maps,\allowbreak{} rather than a new claim to have reread Verdier\textquotesingle s historical source.

For \(\mathcal B\subset\mathcal A\subset\mathcal D\),\allowbreak{} the cone criterion gives \[
S_{\mathcal A}(\mathcal B)=S_{\mathcal D}(\mathcal B)
\cap\operatorname{Arrows}(\mathcal A).
\] For the reverse inclusion,\allowbreak{} replace the cone by an isomorphic object of \(\mathcal B\); the endpoints and all arrows then belong to \(\mathcal A\) by fullness,\allowbreak{} so its triangle belongs to the subcategory. If \(s:X'\to X\) is a denominator with \(X\in\mathcal A\),\allowbreak{} its cone belongs to \(\mathcal B\) up to isomorphism. Rotate to a triangle beginning with \(X\to B\); a completion of that map inside \(\mathcal A\) shows that \(X'[1]\),\allowbreak{} hence \(X'\),\allowbreak{} is isomorphic to an object of \(\mathcal A\). Precomposing \(s\) with that isomorphism gives an isomorphic denominator whose source belongs to \(\mathcal A\).

The full denominator subcategory is therefore an equivalence of index categories:\allowbreak{} every object has such an isomorphism to an object in it,\allowbreak{} and fullness supplies all arrows between those objects. In particular it is cofinal in the direction used for the right-fraction Hom colimit. Replacing every roof by this isomorphic roof gives the fully faithful map \(\mathcal A/\mathcal B\to\mathcal D/\mathcal B\). The object map is the original inclusion because a fraction localization does not identify its objects. Its full image is triangulated:\allowbreak{} shifts are inherited,\allowbreak{} and the cone of an arrow in that image can be computed in \(\mathcal A/\mathcal B\) and mapped across the exact inclusion.

Let \(q_B,q_A\) denote the two quotient functors from \(\mathcal D\),\allowbreak{} and \(q_{A/B}\) the quotient of \(\mathcal D/\mathcal B\) by the displayed image. The composite \(q_{A/B}q_B\) kills \(\mathcal A\),\allowbreak{} so it induces \[
U:\mathcal D/\mathcal A\longrightarrow
(\mathcal D/\mathcal B)/(\mathcal A/\mathcal B),\qquad
Uq_A=q_{A/B}q_B.
\] Conversely \(q_A\) kills \(\mathcal B\),\allowbreak{} descends across \(q_B\),\allowbreak{} and its descended functor kills \(\mathcal A/\mathcal B\). This gives \(V\) in the reverse direction with \(Vq_{A/B}q_B=q_A\). The two composites agree with identity on the original category. Uniqueness of factorization on arrows and their inverted denominators gives \(VU=1\) and \(UV=1\) in these fixed fraction models. With arbitrary equivalent models these equalities give the corresponding natural isomorphisms. Their shift comparisons are inherited identity comparisons throughout. This proves the claimed exact equivalence.

For the insertion\textquotesingle s subcategory correspondence,\allowbreak{} put \(\mathcal K=\ker Q\). The inverse image of the strictly full closure of \(\mathcal C\) is closed under shifts,\allowbreak{} cones and isomorphisms because \(Q\) is exact. The cone criterion gives its stated inverse-image denominator system. If \(\mathcal A\) is strictly full,\allowbreak{} triangulated,\allowbreak{} and contains \(\mathcal K\),\allowbreak{} and \(Q(X)\cong Q(Y)\) for \(Y\in\mathcal A\),\allowbreak{} represent that isomorphism by \(Y\xleftarrow sZ\xrightarrow fX\). Both \(Q(s)\) and \(Q(f)\) are invertible. Their cones therefore belong to \(\mathcal K\),\allowbreak{} even when the original denominator system was not saturated. The first cone triangle puts \(Z\) in \(\mathcal A\),\allowbreak{} and the second puts \(X\) there. This proves \(Q^{-1}(\overline{Q(\mathcal A)})=\mathcal A\). The reverse composite is equality because \(Q\) has the same objects as the source. Finally every object of the quotient is \(Q(X)\),\allowbreak{} and additivity identifies \(Q(X\oplus Y)\) with \(Q(X)\oplus Q(Y)\). Applying the just-proved inverse-image identity to a direct sum proves saturation in either direction. Passing to strictly full closures proves the insertion\textquotesingle s final assertion for non-strictly-full subcategories. This preserves the exact scope of that assertion; no assumption of splitting every idempotent is added.

\subsection{\texorpdfstring{3. Exact right adjoints,\allowbreak{} with the shift comparison and five-lemma step}{3. Exact right  with the shift comparison and five-lemma step}}\label{proofs-1844-2401-exact-right-adjoints-with-the-shift-comparison-and-five-lemma-step}

Let \((F,\xi):\mathcal D\to\mathcal E\) be exact and let \(G:\mathcal E\to\mathcal D\) be a right adjoint. Write \(\eta_X:X\to GF(X)\),\allowbreak{} \(\epsilon_A:FG(A)\to A\) for its unit and counit. Define \(\lambda_A:G(A)[1]\to G(A[1])\) to be adjoint to \[
F(G(A)[1])\xrightarrow{\xi_{G(A)}}F(G(A))[1]
\xrightarrow{\epsilon_A[1]}A[1].
\] For \(u:A\to B\),\allowbreak{} the adjoints of \(G(u[1])\lambda_A\) and \(\lambda_B G(u)[1]\) agree because \[
u[1]\epsilon_A[1]\xi_{G(A)}
=\epsilon_B[1]F(G(u))[1]\xi_{G(A)}
=\epsilon_B[1]\xi_{G(B)}F(G(u)[1]).
\] Thus \(\lambda\) is natural.

Composition with \(\lambda_A\) on \(\operatorname{Hom}_{\mathcal D}(X[1],-)\) is the composite of the four bijections \[
\begin{split}
\operatorname{Hom}_{\mathcal D}(X[1],G(A)[1])
&\cong\operatorname{Hom}_{\mathcal D}(X,G(A))\\
&\cong\operatorname{Hom}_{\mathcal E}(F(X),A)\\
&\cong\operatorname{Hom}_{\mathcal E}(F(X[1]),A[1])\\
&\cong\operatorname{Hom}_{\mathcal D}(X[1],G(A[1])).
\end{split}
\] The third sends \(h\) to \(h[1]\xi_X\). To verify that these bijections are composition with the defined map,\allowbreak{} write an input as \(g[1]\); its adjoint after composition is \[
\epsilon_A[1]\xi_{G(A)}F(g[1])
=\epsilon_A[1]F(g)[1]\xi_X
=(\epsilon_AF(g))[1]\xi_X,
\] which is the same image. Every test object is isomorphic to some \(X[1]\); consequently \(\lambda_A\) is an isomorphism. Put \(\gamma_A=\lambda_A^{-1}:G(A[1])\to G(A)[1]\).

The adjunction bijection on each Hom group is additive:\allowbreak{} it sends \(v:W\to G(A)\) to \(\epsilon_AF(v)\),\allowbreak{} and \(F\) is additive. Naturality of these group isomorphisms and bilinearity of composition show \(G(a+b)=G(a)+G(b)\). The same argument gives \(G(0)=0\). Thus the underlying right adjoint has the needed additivity.

For a distinguished triangle \((A,B,C,a,b,c)\) of \(\mathcal E\),\allowbreak{} complete \(G(a)\) to \((G(A),G(B),X,G(a),v,w)\) in \(\mathcal D\). Applying \(F\) gives boundary \(\xi_{G(A)}F(w)\). TR3 with the counit supplies \(t:F(X)\to C\) satisfying \[
tF(v)=b\epsilon_B,\qquad ct=\epsilon_A[1]\xi_{G(A)}F(w).
\] Let \(j=G(t)\eta_X:X\to G(C)\). Naturality of \(\eta\) and the triangle identity give \[
jv=G(tF(v))\eta_{G(B)}
=G(b)G(\epsilon_B)\eta_{G(B)}=G(b).
\] The second displayed equation says,\allowbreak{} by adjunction,\allowbreak{} \(G(c)j=\lambda_Aw\),\allowbreak{} or \(\gamma_AG(c)j=w\). Hence \((1,1,j)\) is a morphism from the chosen distinguished triangle to the triangle with image objects and boundary \(\gamma_AG(c)\).

Here is the exactness needed for the source\textquotesingle s five-lemma step,\allowbreak{} without assuming the conclusion. Apply \(\operatorname{Hom}_{\mathcal D}(W,-)\) to the five consecutive terms \[
G(A),\ G(B),\ G(C),\ G(A)[1],\ G(B)[1].
\] Adjunction on the first three terms,\allowbreak{} followed by \(\lambda_A\) and \(\lambda_B\) on the last two,\allowbreak{} identifies this sequence with 
\begingroup\small
\[
\operatorname{Hom}_{\mathcal E}(F(W),A)\to
\operatorname{Hom}_{\mathcal E}(F(W),B)\to
\operatorname{Hom}_{\mathcal E}(F(W),C)\to
\operatorname{Hom}_{\mathcal E}(F(W),A[1])\to
\operatorname{Hom}_{\mathcal E}(F(W),B[1]).
\]
\endgroup
 The sequence is exact by representable exactness in \(\mathcal E\). The boundary identification uses \(\lambda_A\gamma_A=1\); the following arrow uses naturality of \(\lambda\). The other row,\allowbreak{} with \(X\) replacing \(G(C)\),\allowbreak{} is exact because its triangle was chosen distinguished. The four surrounding vertical maps are identities on the original \(G(A),G(B)\) and their shifts. The five lemma therefore makes \(\operatorname{Hom}_{\mathcal D}(W,j)\) bijective for every \(W\),\allowbreak{} and Yoneda makes \(j\) an isomorphism. The image triangle is distinguished,\allowbreak{} proving that \((G,\gamma)\) is exact.

This proves the mathematical context of \textbf{OCC-01496 /\allowbreak{} P02-E0024}. All four erroneous category occurrences at 2361–\allowbreak{}2364 are in \(\mathcal D\),\allowbreak{} not \(\mathcal D'\) (called \(\mathcal E\) above)\allowbreak{}. The already implemented \textbf{MC-STK-ERR-0818},\allowbreak{} \textbf{0819},\allowbreak{} and \textbf{0820} correct respectively the functor,\allowbreak{} the test object,\allowbreak{} and the two Hom groups. The extra exactness argument here is an editorial completion of the omitted details,\allowbreak{} not another source error or a replacement attributed to the authors.

\subsection{\texorpdfstring{4. DERIVED-UNDER-003:\allowbreak{} pre-triangulated adjunctions and the equivalence test}{4.  pre-triangulated adjunctions and the equivalence test}}\label{proofs-1844-2401-derived-under-003-pre-triangulated-adjunctions-and-the-equivalence-test}

The preceding proof uses TR1,\allowbreak{} TR2,\allowbreak{} TR3 and representable exactness; it never uses TR4. Thus it proves the adjoint-exactness result already for \textbf{pre-triangulated} categories. The unit and counit moreover respect the exact structures just constructed. The counit equation is \[
\epsilon_{A[1]}=\epsilon_A[1]\xi_{G(A)}F(\gamma_A),
\] obtained by writing the defining adjunction equation for \(\lambda_A\) and composing with \(F(\gamma_A)\). For the unit the equation is \[
\eta_X[1]=\gamma_{F(X)}G(\xi_X)\eta_{X[1]}.
\] Indeed,\allowbreak{} after composing with \(\lambda_{F(X)}\),\allowbreak{} the two sides have the same adjoint:\allowbreak{} naturality of \(\xi\) and \(\epsilon_{F(X)}F(\eta_X)=1\) make both adjoints equal to \(\xi_X\). Faithfulness of the adjunction bijection proves equality. These are the two explicit compatibility squares of exact functors.

For the next lemma at 2372–\allowbreak{}2401 it consequently suffices to assume that \(F\) is fully faithful and exact between pre-triangulated categories,\allowbreak{} that \(G\) is its right adjoint,\allowbreak{} and that \(\ker G=0\). One need not separately assume \(G\) exact or either TR4 axiom. Full faithfulness makes \(\eta_X\) an isomorphism:\allowbreak{} composition with it on each test Hom group is the fully faithful map induced by \(F\),\allowbreak{} followed by the adjunction bijection. For every \(A\),\allowbreak{} the triangle identity \[
G(\epsilon_A)\eta_{G(A)}=1_{G(A)}
\] then shows that \(G(\epsilon_A)\) is the inverse of the specified isomorphism \(\eta_{G(A)}\). Complete \(\epsilon_A:FG(A)\to A\) to a distinguished triangle with cone \(Y\). Exactness of \((G,\gamma)\) implies \(G(Y)=0\),\allowbreak{} since its first arrow is an isomorphism. The zero-kernel assumption gives \(Y=0\); hence \(\epsilon_A\) is an isomorphism. The unit and counit,\allowbreak{} with the two proved comparison equations,\allowbreak{} exhibit an exact equivalence.

This also proves that the source\textquotesingle s identification of \(GFG(A)\) with \(G(A)\) hides a definite inverse pair,\allowbreak{} rather than merely an object isomorphism from which invertibility of an arbitrary endomorphism would not follow. Its intended arrow is precisely \(G(\epsilon_A)\).

The original theorem statements remain unchanged. The direct uses in \texttt{equiv.\allowbreak{}tex:\allowbreak{}880–\allowbreak{}890} and \texttt{injectives.\allowbreak{}tex:\allowbreak{}2804–\allowbreak{}2825} retain their existing triangulated hypotheses and are covered by the proved stronger result. The former also invokes the left-adjoint version,\allowbreak{} obtained by reversing the adjunction and both categories,\allowbreak{} with the corresponding inverse shift comparisons. The other equivalence-test citations located in DGA,\allowbreak{} sheaf DGA,\allowbreak{} Perfect Complexes,\allowbreak{} algebraic-space chapters and Stacks Sheaves are routing records pending their separate receiving checks; this note does not pretend to have read those proofs.
\clearpage
\section{\texorpdfstring{Cones,\allowbreak{} strictification,\allowbreak{} and the chain-level square-zero calculation}{  and the chain-level square-zero calculation}}\label{proofs-2419-2994-cones-strictification-and-the-chain-level-square-zero-calculation}

Authority and line conventions are those of \texttt{INPUT\_\allowbreak{}BINDING.\allowbreak{}json} and the earlier Derived review notes. Original source reading now reaches line 2994. This note retains all original components,\allowbreak{} degree indices,\allowbreak{} signs and splittings. Additional proofs are editorial; no novelty or newly verified historical-source reading is claimed.

\subsection{1. Cone maps and the missing shift}\label{proofs-2419-2994-cone-maps-and-the-missing-shift}

At 2424 \textbf{MC-STK-ERR-0821} changes ``We set ... be'' to ``We let ... be''. It is correct,\allowbreak{} with no new report in this interval. The categories remain those of cochain complexes and their homotopy classes.

For \(f:K^\bullet\to L^\bullet\),\allowbreak{} retain \[
C(f)^n=L^n\oplus K^{n+1},\qquad
D^n=\begin{pmatrix}d_L^n&f^{n+1}\\0&-d_K^{n+1}\end{pmatrix}.
\] The upper-right entry of \(D^{n+1}D^n\) is \(d_L^{n+1}f^{n+1}-f^{n+2}d_K^{n+1}=0\),\allowbreak{} while the diagonal entries vanish because the two original differentials square to zero. The inclusion \(i^n(l)=(l,0)\) is a chain map. The projection \(p^n(l,k)=k\) is a chain map to \(K[1]\),\allowbreak{} whose differential in degree \(n\) is \(-d_K^{n+1}\).

In the source\textquotesingle s homotopy-commutative square keep \(a:K_1\to K_2\),\allowbreak{} \(b:L_1\to L_2\),\allowbreak{} and \(h^n:K_1^n\to L_2^{n-1}\),\allowbreak{} with \[
b^nf_1^n-f_2^na^n=d_{L_2}^{n-1}h^n+h^{n+1}d_{K_1}^n.
\] The source matrix is \[
c^n=\begin{pmatrix}b^n&h^{n+1}\\0&a^{n+1}\end{pmatrix}.
\] Its chain-map equation compares \[
D_2^nc^n=
\begin{pmatrix}
d_{L_2}^nb^n&d_{L_2}^nh^{n+1}+f_2^{n+1}a^{n+1}\\
0&-d_{K_2}^{n+1}a^{n+1}
\end{pmatrix}
\] with \[
c^{n+1}D_1^n=
\begin{pmatrix}
b^{n+1}d_{L_1}^n&b^{n+1}f_1^{n+1}-h^{n+2}d_{K_1}^{n+1}\\
0&-a^{n+2}d_{K_1}^{n+1}
\end{pmatrix}.
\] The diagonal equations are the chain-map equations for \(a,b\); the upper-right equation is the displayed homotopy identity in degree \(n+1\). This proves the surrounding claim of \textbf{OCC-01497 /\allowbreak{} P02-E0025}. Its ``show'' to ``shows'' copyedit is already present as \textbf{MC-STK-ERR-0822}.

The two further equations are \[
c^ni_1^n=i_2^nb^n,\qquad p_2^nc^n=a^{n+1}p_1^n.
\] Thus the third-square identity is \(p_2c=a[1]p_1\),\allowbreak{} not the unshifted expression at original 2535. This is \textbf{DERIVED-NEW-007}. The codomain of \(p_1\) is \(K_1[1]\),\allowbreak{} so the right factor must be followed by \(a[1]:K_1[1]\to K_2[1]\). The operation inserts that shift and leaves the source\textquotesingle s chosen homotopy and matrix untouched.

At 2493 the source calls the sextuple with \textbf{positive} \(p\) a triangle; it does not yet call it distinguished. That statement is correct by the definition of a triangle. The later distinguished cone triangle uses \(-p\). No correction of the positive \(p\) at 2493 is justified,\allowbreak{} and none is proposed.

For the receiving factorization lemma,\allowbreak{} if \(gf=d_Mh+hd_K\),\allowbreak{} the first map is explicitly \[
u^n=(g^n,h^{n+1}):L^n\oplus K^{n+1}\to M^n.
\] Its chain equation is the same upper-right identity above,\allowbreak{} and \(ui=g\) as chain maps. For the other map retain \[
C(g)[-1]^n=M^{n-1}\oplus L^n,\qquad
d_{C(g)[-1]}^n=
\begin{pmatrix}-d_M^{n-1}&-g^n\\0&d_L^n\end{pmatrix},\qquad
v^n=\begin{pmatrix}-h^n\\f^n\end{pmatrix}.
\] Then the first component of \(d_{C(g)[-1]}^nv^n\) is \(d_M^{n-1}h^n-g^nf^n=-h^{n+1}d_K^n\); the second is \(d_L^nf^n=f^{n+1}d_K^n\). Thus \(v\) is a chain map and \(p[-1]v=f\). Both factorizations hold in the category of complexes,\allowbreak{} with the dependence on the chosen homotopy retained.

\subsection{2. Strictification with exactly one split side}\label{proofs-2419-2994-strictification-with-exactly-one-split-side}

For the square at 2590–\allowbreak{}2615 let \(bf-ga=d_Dh+hd_A\),\allowbreak{} with \(h^n:A^n\to D^{n-1}\). If \(\pi^nf^n=1_{A^n}\),\allowbreak{} put \(H^n=h^n\pi^n:B^n\to D^{n-1}\) and \[
(b')^n=b^n-d_D^{n-1}H^n-H^{n+1}d_B^n.
\] This is a chain map homotopic to \(b\). Since \(f\) is a chain map,\allowbreak{} \[
(b')^nf^n=b^nf^n-d_D^{n-1}h^n
-h^{n+1}\pi^{n+1}f^{n+1}d_A^n=g^na^n.
\] No assertion that \(\pi\) is a chain map is used. If instead \(g^ns^n=1_{D^n}\),\allowbreak{} put \(J^n=s^{n-1}h^n:A^n\to C^{n-1}\) and \[
(a')^n=a^n+d_C^{n-1}J^n+J^{n+1}d_A^n.
\] It is homotopic to \(a\),\allowbreak{} and the chain-map equation for \(g\) gives \(g^n(a')^n=g^na^n+d_D^{n-1}h^n+h^{n+1}d_A^n=b^nf^n\). These are the two exact corrections asserted by the source.

The source\textquotesingle s replacement of \(\alpha:K\to L\) by a split injection has degree \(n\) object \(L^n\oplus K^n\oplus K^{n+1}\),\allowbreak{} differential \[
\begin{pmatrix}d_L^n&0&0\\0&d_K^n&1_{K^{n+1}}\\
0&0&-d_K^{n+1}\end{pmatrix},
\] and maps \(\widetilde\alpha^n=(\alpha^n,1,0)^t\),\allowbreak{} \(\pi^n=(1,0,0)\),\allowbreak{} \(s^n=(1,0,0)^t\). Their displayed compositions are the original \(\alpha\) and identity. The homotopy \(H^n(l,k,k^+)=(0,0,k)\) gives \[
dH(l,k,k^+)=(0,k,-d_K^nk),\qquad
Hd(l,k,k^+)=(0,0,d_K^nk+k^+),
\] so \(dH+Hd=(0,k,k^+)=1-s\pi\),\allowbreak{} with no missing sign.

For the split-surjection replacement retain \(K^n\oplus L^{n-1}\oplus L^n\) and differential \[
\begin{pmatrix}d_K^n&0&0\\0&-d_L^{n-1}&1_{L^n}\\
0&0&d_L^n\end{pmatrix}.
\] The source maps are \(\widetilde\alpha^n=(\alpha^n,0,1)\),\allowbreak{} \(i^n=(1,0,0)^t\),\allowbreak{} \(s^n=(1,0,0)\). Here \(H^n(k,l^-,l)=(0,0,l^-)\),\allowbreak{} and the two homotopy terms are \((0,l^-,d_L^{n-1}l^-)\) and \((0,0,-d_L^{n-1}l^-+l)\). Their sum is \(1-is\). Each construction adds only these specified shifts and finite direct sums,\allowbreak{} so it preserves all three boundedness classes in the statements.

\subsection{3. Changing splittings and the obstruction to changing both squares}\label{proofs-2419-2994-changing-splittings-and-the-obstruction-to-changing-both-squares}

For the original split sequence with maps \(\alpha,\beta\),\allowbreak{} keep \(\pi\alpha=1\),\allowbreak{} \(\beta s=1\),\allowbreak{} \(\pi s=0\),\allowbreak{} \(\beta\alpha=0\),\allowbreak{} and \(\alpha\pi+s\beta=1\) degree by degree. The boundary is \(\delta^n=\pi^{n+1}d_B^ns^n\). If \(s'^n=s^n+\alpha^nh^n\),\allowbreak{} then \(\pi'^n=\pi^n-h^n\beta^n\). Expanding all four terms gives \[
\begin{split}
(\delta')^n={}&\pi^{n+1}d_B^ns^n
+\pi^{n+1}d_B^n\alpha^nh^n
-h^{n+1}\beta^{n+1}d_B^ns^n\\
&-h^{n+1}\beta^{n+1}d_B^n\alpha^nh^n\\
={}&\delta^n+d_A^nh^n-h^{n+1}d_C^n.
\end{split}
\] The fourth term is zero because \(\beta^{n+1}d_B^n\alpha^n=\beta^{n+1}\alpha^{n+1}d_A^n=0\). The family \(-h^n:C^n\to A[1]^{n-1}\) is therefore a homotopy from \(\delta\) to \(\delta'\),\allowbreak{} since the target differential is \(-d_A^n\). This agrees with the earlier Homology result and proves the identity triple is an isomorphism of triangles in the homotopy category. \textbf{OCC-02892 /\allowbreak{} P02-E1429} matches \textbf{MC-STK-ERR-0823}:\allowbreak{} there is one sequence and two choices of splitting,\allowbreak{} so the corrected singular ``a termwise split exact sequence'' is appropriate.

The limitation at 2876–\allowbreak{}2900 is also real. In the cited original example,\allowbreak{} let \(R=k[\epsilon]/(\epsilon^2)\). Keep \(A^1=R\),\allowbreak{} \(C^0=R\),\allowbreak{} and \(B=(R\xrightarrow\epsilon R)\) in degrees zero and one,\allowbreak{} with the source\textquotesingle s termwise identity inclusion and projection. The endpoint maps are \(1+\epsilon\) on \(A\) and \(1\) on \(C\),\allowbreak{} while the middle map is \(1_B\). To make both squares commute as chain maps,\allowbreak{} a replacement middle map would have to be \(1\) in degree zero and \(1+\epsilon\) in degree one. It is a chain map,\allowbreak{} because \(\epsilon^2=0\),\allowbreak{} but it is not homotopic to \(1_B\):\allowbreak{} a homotopy is multiplication by a single \(r\in R\) from \(B^1\) to \(B^0\),\allowbreak{} and its two components \(Hd\) and \(dH\) are both multiplication by \(r\epsilon\). They cannot equal zero in degree zero and \(\epsilon\ne0\) in degree one. Thus the two strictifications cannot in general be imposed on one middle map in the same homotopy class. This verifies the receiving example directly,\allowbreak{} without relying only on its trace assertion or extending the one-square lemma too far.

\subsection{4. DERIVED-NEW-008 and the square-zero consequence without kernels}\label{proofs-2419-2994-derived-new-008-and-the-square-zero-consequence-without-kernels}

The category at 2902–\allowbreak{}2938 is only additive. The image and kernel notation at 2935–\allowbreak{}2936 need not denote existing objects for arbitrary maps \(b^n,(b')^n\). The proof needs only the given split sequence,\allowbreak{} and can state its factorization directly. This is a clarification of the existing argument,\allowbreak{} not a counterexample to the lemma.

First replace \(b\) by a homotopic chain map with \(\beta_2b=0\),\allowbreak{} using the termwise split surjection \(\beta_2\). Independently replace \(b'\) by a homotopic chain map with \(b'\alpha_2=0\),\allowbreak{} using the termwise split injection \(\alpha_2\). Write \(\pi_2,s_2\) for the supplied degreewise splitting maps. For each degree the exact biproduct identity gives \[
b^n=(\alpha_2^n\pi_2^n+s_2^n\beta_2^n)b^n
=\alpha_2^n\pi_2^nb^n.
\] Consequently \[
(b')^nb^n=(b')^n\alpha_2^n\pi_2^nb^n=0.
\] This proves zero composition for the strictified chain maps. Composition of chain maps preserves homotopies,\allowbreak{} so their original homotopy classes also have zero product. Neither an image of \(b^n\) nor a kernel of \((b')^n\) is required. The proposed source operation uses exactly these equalities and the original split sequence.

For the cone comparison at 2940–\allowbreak{}2984,\allowbreak{} an endomorphism \((1,1,c)\) of the cone triangle gives two commuting squares,\allowbreak{} up to homotopy,\allowbreak{} on the termwise split sequence \(L\to C(f)\to K[1]\). Apply the proved chain-level square-zero assertion to \(c-1\) twice:\allowbreak{} \((c-1)^2=0\). Thus \[
c(2-c)=(2-c)c=2c-c^2=1.
\] For a general \((a,b,c)\) with \(a,b\) homotopy equivalences,\allowbreak{} choose representatives of their inverses and construct a cone map \(c'\) using the homotopy-commutative inverse square. Both \(cc'\) and \(c'c\) are then invertible by the preceding calculation. Thus \(c' (cc')^{-1}\) is a right inverse of \(c\),\allowbreak{} and \((c'c)^{-1}c'\) is a left inverse; the two coincide. This proves the cone comparison claim in the source.

This is also a receiving instance of \textbf{DERIVED-UNDER-001},\allowbreak{} the square-zero ideal for triangle maps. The present proof was made directly on complexes so that it is independent of the triangulated structure on the homotopy category,\allowbreak{} which the chapter constructs later using these lemmas. Using the abstract triangulated result here as the foundation would have introduced a circular proof; no such replacement is made.
\clearpage
\section{Split sequences and the homotopy triangulation}\label{proofs-2995-3485-split-sequences-and-the-homotopy-triangulation}

Source:\allowbreak{} the Stacks Project authors\textquotesingle{} \texttt{derived.\allowbreak{}tex} at the authority commit bound in \texttt{INPUT\_\allowbreak{}BINDING.\allowbreak{}json}. Original line numbers are used. The arguments below verify the source\textquotesingle s objects and maps without changing its definition of distinguished triangles. New observations are editorial; no novelty is claimed.

\subsection{1. The signed cone comparison and its two homotopies}\label{proofs-2995-3485-the-signed-cone-comparison-and-its-two-homotopies}

Retain the termwise split sequence with maps \(\alpha,\beta\),\allowbreak{} sections \(s^n\),\allowbreak{} retractions \(\pi^n\),\allowbreak{} and boundary \(\delta^n=\pi^{n+1}d_B^ns^n\). Besides the splitting identities,\allowbreak{} the differential satisfies \[
d_B^ns^n-s^{n+1}d_C^n=\alpha^{n+1}\delta^n,
\qquad
\pi^{n+1}d_B^n-d_A^n\pi^n=\delta^n\beta^n.
\] For the second equality insert \(1_{B^n}=\alpha^n\pi^n+s^n\beta^n\) after \(d_B^n\),\allowbreak{} and use \(d_B^n\alpha^n=\alpha^{n+1}d_A^n\). The first equality follows by inserting the analogous identity before \(d_B^ns^n\) and using \(\beta^{n+1}d_B^ns^n=d_C^n\).

Write \(D\) for the cone differential of \(\alpha\). The source maps are \[
P^n=(\beta^n,0):B^n\oplus A^{n+1}\to C^n,
\qquad S^n=\binom{s^n}{-\delta^n}:C^n\to B^n\oplus A^{n+1}.
\] The displayed identities and \(d_A^{n+1}\delta^n+\delta^{n+1}d_C^n=0\) give \(d_C^nP^n=P^{n+1}D^n\),\allowbreak{} and \(D^nS^n=S^{n+1}d_C^n\). More explicitly the latter has first component \(d_B^ns^n-\alpha^{n+1}\delta^n=s^{n+1}d_C^n\) and second component \(d_A^{n+1}\delta^n=-\delta^{n+1}d_C^n\). Also \(PS=1_C\).

The proposed homotopy is \[
H^n=\begin{pmatrix}0&0\\\pi^n&0\end{pmatrix}:
C(\alpha)^n\longrightarrow C(\alpha)^{n-1}.
\] Its two products add to \[
D^{n-1}H^n+H^{n+1}D^n
=\begin{pmatrix}
\alpha^n\pi^n&0\\
-d_A^n\pi^n+\pi^{n+1}d_B^n&\pi^{n+1}\alpha^{n+1}
\end{pmatrix}
=\begin{pmatrix}1-s^n\beta^n&0\\\delta^n\beta^n&1\end{pmatrix}
=1-S^nP^n.
\] This proves the stated homotopy equivalence. The last triangle square also has a direct homotopy:\allowbreak{} let \(T^n=(\pi^n,0):C(\alpha)^n\to A[1]^{n-1}=A^n\). Then \[
d_{A[1]}^{n-1}T^n+T^{n+1}D^n
=(\delta^n\beta^n,1)=\delta^nP^n+p^n.
\] Thus \(\delta P\) equals \(-p\) in the homotopy category. The other two squares commute as chain maps,\allowbreak{} since \(Pi=\beta\). This proves the precise isomorphism in (1)\allowbreak{}.

For (2)\allowbreak{},\allowbreak{} the split-injection replacement preserves \(K\),\allowbreak{} and its projection \(\widetilde L\to L\) is a homotopy equivalence. The cone construction with its chosen commuting square gives an isomorphism of the positive-projection triangles. Negating both last arrows preserves that same triangle isomorphism,\allowbreak{} so it applies to the signed triangles used in (1)\allowbreak{}. Taking the degreewise complements of the split injection supplies the required quotient complex. The three chain-map equations defining that quotient follow by projecting the original differential,\allowbreak{} since the injected subobject is a subcomplex.

\textbf{OCC-02893 /\allowbreak{} P02-E1430} and French \textbf{DERIVED-013},\allowbreak{} \textbf{DERIVED-014} correctly identify the two missing first arrows in the displayed tuples. The existing \textbf{MC-STK-ERR-0824},\allowbreak{} \textbf{0825} insert \(f\) and \(\widetilde f\). The positive last arrows in those comparison tuples are legitimate; they can be negated on both sides as just proved. They are not another sign error.

For a finite composable sequence the induction at 3141–\allowbreak{}3173 applies the same split-injection replacement to \(B_{n-1}\to A_{n-1}\to A_n\). Its projection commutes with that composite as a chain map,\allowbreak{} and it is the same explicitly contracted homotopy equivalence proved earlier. Finite sums and one shift at each finite stage preserve the three stated boundedness classes.

\subsection{\texorpdfstring{2. DERIVED-NEW-009:\allowbreak{} the object carrying the rotation sign}{2.  the object carrying the rotation sign}}\label{proofs-2995-3485-derived-new-009-the-object-carrying-the-rotation-sign}

The decomposition in the rotation lemma identifies \(B^n\) with \(A^n\oplus C^n\). Its differential is \[
\begin{pmatrix}d_A^n&\delta^n\\0&d_C^n\end{pmatrix}.
\] For the cone of \(\delta[-1]:C[-1]\to A\),\allowbreak{} the lower diagonal is \(-d_{C[-1]}^{n+1}=d_C^n\),\allowbreak{} and the upper right is exactly \(\delta^n\). Thus the source\textquotesingle s comparison \(q:B\to C(\delta[-1])\) is a chain isomorphism with \(q\alpha=i\) and \(pq=\beta\).

The backward-rotated triangle is \[
S=(C[-1],A,B,-\delta[-1],\alpha,\beta),
\] whereas the distinguished cone triangle wanted in the proof is \[
T=(C[-1],A,C(\delta[-1]),\delta[-1],i,-p).
\] The precise isomorphism \(S\to T\) is \((-1_{C[-1]},1_A,q)\). Its three squares are \[
1_A(-\delta[-1])=\delta[-1](-1_{C[-1]}),\quad
q\alpha=i1_A,\quad (-p)q=-\beta=(-1_{C[-1]})[1]\beta.
\] If \((r_X,r_Y,r_Z)\) is the original comparison from \((X,Y,Z,f,g,h)\) to \(S\),\allowbreak{} the resulting comparison to \(T\) is \((-r_X,r_Y,qr_Z)\). Consequently the instruction at original 3350 must say \textbf{\(-1\) on \(X\)},\allowbreak{} not on \(Y\). Changing only the \(Y\)-component would make the second square require \(i=-i\),\allowbreak{} which is false,\allowbreak{} for example for a nonzero direct summand inclusion of complexes of free abelian groups concentrated in degree zero. The theorem remains valid with the corrected component.

For forward rotation the source compares 
\par\noindent \((L,C(f),K[1],i,-p,-f[1])\)\par\noindent
 with 
\par\noindent \((L,C(f),K[1],i,p,f[1])\)\par\noindent
. Here the isomorphism is \((1_L,1_{C(f)},-1_{K[1]})\). The last boundary of the termwise split cone sequence is exactly \(f^{n+1}\) in degree \(n\),\allowbreak{} by its differential matrix. This proves both directions of TR2 with the original convention \(-p\) for the distinguished cone triangle.

\subsection{\texorpdfstring{3. DERIVED-NEW-010:\allowbreak{} the second octahedral square needs a homotopy}{3.  the second octahedral square needs a homotopy}}\label{proofs-2995-3485-derived-new-010-the-second-octahedral-square-needs-a-homotopy}

The definition at 2582–\allowbreak{}2587 requires an actual direct-summand inclusion,\allowbreak{} not merely the existence of an arbitrary left inverse. Therefore the degreewise complements used at 3273–\allowbreak{}3274 do exist. For any other left inverse of such an inclusion its kernel is the graph of a map from the known complement:\allowbreak{} in coordinates a left inverse has the form \((1,t)\),\allowbreak{} with kernel inclusion \(q\mapsto(-tq,q)\). The projection to that complement gives the universal factorization of every map killed by \((1,t)\). No assumption that the whole additive category is Karoubian is needed.

Retain the original \(A,B,C,Q_1,Q_2,Q_3\) and the chosen retractions. In their actual degreewise decompositions write \[
B^n=A^n\oplus Q_1^n,\quad
C^n=A^n\oplus Q_1^n\oplus Q_3^n,\quad
Q_2^n=Q_1^n\oplus Q_3^n.
\] All entries of the differential of \(C\) are retained as \[
d_C^n=\begin{pmatrix}
d_A^n&u^n&v^n\\0&d_{Q_1}^n&w^n\\0&0&d_{Q_3}^n
\end{pmatrix}.
\] The zero lower entries follow from the two inclusions being chain maps. The three off-diagonal identities in \(d_C^{n+1}d_C^n=0\) are,\allowbreak{} without discarding any term,\allowbreak{} \[
d_A^{n+1}u^n+u^{n+1}d_{Q_1}^n=0,
\] \[
d_{Q_1}^{n+1}w^n+w^{n+1}d_{Q_3}^n=0,
\] \[
d_A^{n+1}v^n+u^{n+1}w^n+v^{n+1}d_{Q_3}^n=0.
\] The quotient differentials are the displayed diagonal block and its \(2\times2\) upper triangular submatrix. The three connecting maps are \[
d_1^n=u^n,\qquad d_2^n=(u^n,v^n),\qquad
d_3^n=\binom{v^n}{w^n}:Q_3^n\to B^{n+1}.
\] These identities prove that each is a chain map into the specified shifted target.

Put \(a^n(q_1)=(q_1,0)\) and \(b^n(q_1,q_3)=q_3\). These are chain maps. The first downward triple \((1_A,\beta,a)\) commutes in every square already on complexes,\allowbreak{} including \(d_2a=d_1\). For the second downward triple \((\alpha,1_C,b)\),\allowbreak{} the first two squares commute on complexes,\allowbreak{} but the last square has difference \[
d_3^nb^n-\alpha^{n+1}d_2^n
=\begin{pmatrix}-u^n&0\\0&w^n\end{pmatrix}:
Q_1^n\oplus Q_3^n\to A^{n+1}\oplus Q_1^{n+1}.
\] It need not vanish. Its exact homotopy is \[
H^n:Q_2^n\longrightarrow B[1]^{n-1}=B^n,
\qquad H^n(q_1,q_3)=(0,q_1).
\] Indeed,\allowbreak{} \[
d_{B[1]}^{n-1}H^n+H^{n+1}d_{Q_2}^n
=\begin{pmatrix}-u^n&0\\-d_{Q_1}^n&0\end{pmatrix}
+\begin{pmatrix}0&0\\d_{Q_1}^n&w^n\end{pmatrix}
=d_3^nb^n-\alpha^{n+1}d_2^n.
\] This proves the required morphism of triangles in \(K(\mathcal A)\). It also identifies the defect in the source\textquotesingle s claim at 3291–\allowbreak{}3293 that both downward triples are compatible with the chosen splittings:\allowbreak{} \(\pi_3(a,q_1,q_3)=(a,q_1)\),\allowbreak{} whereas \(\alpha\pi_2(a,q_1,q_3)=(a,0)\). For an explicit nonzero boundary defect take \(Q_3=0\),\allowbreak{} \(Q_1^0=\mathbf Z\),\allowbreak{} \(A^1=\mathbf Z\),\allowbreak{} all other components zero,\allowbreak{} and \(u^0=1\). Then \(C=B\),\allowbreak{} \(\beta=1\),\allowbreak{} and the displayed boundary difference is \(-1\) in degree zero. It is null-homotopic by precisely the displayed \(H\).

The remaining TR4 triangle comes from the termwise split sequence \(Q_1\to Q_2\to Q_3\). Its connecting map is \(w\),\allowbreak{} and \[
p_1[1]d_3=(0,1)\binom v w=w
\] as chain maps. Thus every part of TR4 holds with the original \(a,b,p_1[1]d_3\). The source sentence is repaired by saying that the boundary squares commute up to homotopy; the theorem and its subsequent application require no alteration.

The grammatical subject at 3305–\allowbreak{}3306 is “splittings”,\allowbreak{} so ``provides'' becomes ``provide'' (\textbf{DERIVED-NEW-011})\allowbreak{}. The remark at 3395–\allowbreak{}3396 omits the noun ``complexes'' after ``bounded (above,\allowbreak{} below)\allowbreak{}”; restoring that noun is \textbf{DERIVED-NEW-012}. Neither copyedit changes a mathematical condition.

\subsection{4. Completing the triangulation and its immediate receiving uses}\label{proofs-2995-3485-completing-the-triangulation-and-its-immediate-receiving-uses}

TR1 is supplied by the split identity sequence and the signed cone comparison of §1. TR2 is §2. For TR3 choose representatives of the given morphisms and a homotopy of their square. The cone matrix in \texttt{PROOFS\_\allowbreak{}2419\_\allowbreak{}2994.\allowbreak{}md},\allowbreak{} §1 has \(p'c=a[1]p\),\allowbreak{} so also \((-p')c=a[1](-p)\). It is therefore the required morphism of the distinguished signed cone triangles. This is the direct receiving check for \textbf{DERIVED-NEW-007}.

Only after TR1–\allowbreak{}TR3 have been established does the source use the earlier easier-TR4 lemma. The finite-sequence replacement above produces a strictly commuting diagram of the two original composable maps with homotopy equivalences on objects. Section 3 supplies the three distinguished triangles and their octahedral maps in that replacement. Transport through those object isomorphisms gives TR4 for the original maps. No use of the triangulated structure was made before it had been constructed.

For the bounded subcategories let \(A^n=0\) for \(n<a\) and \(B^n=0\) for \(n<b\). The cone degree \(B^n\oplus A^{n+1}\) vanishes for \(n<\min(b,a-1)\). With upper bounds \(a,b\) it vanishes for \(n>\max(b,a-1)\). These exact bounds and the shift functors prove closure under the required completions. Fullness is the original Hom definition. The alternative proof uses the same split-injection replacement and its degreewise quotient,\allowbreak{} so it has the same boundedness property. \textbf{OCC-01498 /\allowbreak{} P02-E0026} and French \textbf{DERIVED-015} match \textbf{MC-STK-ERR-0826},\allowbreak{} changing ``an termwise'' to ``a termwise''.

An additive functor \(F:\mathcal A\to\mathcal B\) preserves every splitting equation,\allowbreak{} including the full identity \(\alpha\pi+s\beta=1\). It commutes with the cochain shift on these exact chosen complexes because \(F(-d_A^{n+1})=-F(d_A^{n+1})\). Thus its translation comparison is the identity,\allowbreak{} and \[
F(\pi^{n+1}d_B^ns^n)=F(\pi^{n+1})F(d_B^n)F(s^n)
\] is the connecting map of the image split sequence. This proves exactness for all four homotopy categories. The missing-comparison findings earlier in the chapter require no edit here:\allowbreak{} the comparison in this construction is explicitly identity on the same complexes.

Finally,\allowbreak{} for the improvement lemma keep the given boundary \(c:C\to A[1]\),\allowbreak{} its cone \(W\),\allowbreak{} and the source maps \(i:A[1]\to W\),\allowbreak{} \(p:W\to C[1]\). Two backward rotations give exactly \((A,W[-1],C,-i[-1],p[-1],c)\). The first backward rotation has first arrow \(p[-1]\),\allowbreak{} and the second has first arrow \(-i[-1]\); the signs are therefore correct. To apply TR3 with identities on \(A,C\),\allowbreak{} first rotate both triangles twice so that \(c\) is their first arrow,\allowbreak{} complete \((1_C,1_{A[1]})\),\allowbreak{} and rotate back. The third-component isomorphism criterion makes the resulting middle map invertible. Degreewise \(W[-1]^n=A^n\oplus C^n\). With the first map the negative inclusion and the second the projection,\allowbreak{} the retraction is the negative projection to \(A^n\) and the section is the positive inclusion of \(C^n\). All splitting identities hold. The differential\textquotesingle s upper right is \(-c^n\); composing with that negative retraction gives connecting map \(c^n\),\allowbreak{} as required. No further sign repair is needed at 3475–\allowbreak{}3484.

\subsection{5. Report-accounting repair}\label{proofs-2995-3485-report-accounting-repair}

The earlier review files bound the CJK reports and the existing French-origin correction operations,\allowbreak{} but did not yet record the fifteen physical French reports through 3422 as separately adjudicated occurrences. Their full original observations,\allowbreak{} adverse evidence and proposals have now been read,\allowbreak{} not inferred from overlapping loci. DERIVED-001 through DERIVED-015 correspond respectively to MC-STK-ERR-0812 through MC-STK-ERR-0826,\allowbreak{} with the proofs in the earlier notes and this note. DERIVED-001\textquotesingle s explanation that \(g\) is undefined is incorrect:\allowbreak{} it is defined but is the wrong map for the stated composite. Its proposed replacement by \(u\) is correct. The rest of those fifteen proposed changes agree with their bound operations. A supplemental disposition record joins these physical reports to their semantic groups and preserves the earlier checkpoint files. Future source-order selection uses the common routing records for both report schemas,\allowbreak{} avoiding the CJK-only field selection.
\clearpage
\section{\texorpdfstring{Totalization,\allowbreak{} canonical triangles,\allowbreak{} and the vanishing-composition bound}{ canonical  and the vanishing-composition bound}}\label{proofs-3487-4161-totalization-canonical-triangles-and-the-vanishing-composition-bound}

Authority:\allowbreak{} the Stacks Project authors\textquotesingle{} \texttt{derived.\allowbreak{}tex},\allowbreak{} commit \texttt{a04446e5\allowbreak{}7ec1fbc2\allowbreak{}52a871af\allowbreak{}cec7752f\allowbreak{}b2807b14},\allowbreak{} bound in \texttt{INPUT\_\allowbreak{}BINDING.\allowbreak{}json}. Line numbers without a chapter prefix refer to that original file. Other chapter passages are the current public sources bound in the accompanying reading record. This is an editorial proof record. Source statements and their signs remain identifiable; no novelty is claimed.

\subsection{\texorpdfstring{1. Exact totalization,\allowbreak{} preserving both component conventions}{1. Exact  preserving both component conventions}}\label{proofs-3487-4161-exact-totalization-preserving-both-component-conventions}

The source uses commuting double-complex differentials and the full total differential \[
T(A)^n=\bigoplus_{p+q=n}A^{p,q},\qquad
D|_{A^{p,q}}=d_1^{p,q}+(-1)^p d_2^{p,q}.
\] The two diagonal squares of \(D^2\) vanish because \(d_1^2=d_2^2=0\). Its two mixed terms are \((-1)^{p+1}d_2^{p+1,q}d_1^{p,q}\) and \((-1)^pd_1^{p,q+1}d_2^{p,q}\); they cancel by the commuting-square condition. Equality on every summand inclusion proves the equality on the coproduct; no rearrangement of an infinite sum of morphisms is required.

For the first convention,\allowbreak{} with outer degree \(q\),\allowbreak{} the source shift is \([0,1]\). The exact functor comparison \(\xi_A:T(A[0,1])\to T(A)[1]\) multiplies the summand \(A^{p,q+1}\) by \((-1)^p\). This is the inverse of the original \(\gamma\) in Homology,\allowbreak{} and its inverse has the same component sign. On a horizontal differential,\allowbreak{} the source has coefficient \(1\) and the target has coefficient \(-1\),\allowbreak{} balanced by \((-1)^{p+1}=-(-1)^p\). On a vertical differential,\allowbreak{} both sides have coefficient \((-1)^{p+1}\) before applying that same sign. Thus this is an actual chain isomorphism,\allowbreak{} natural in all bigraded chain maps.

If \(h^{p,q}:A^{p,q}\to B^{p,q-1}\) is an outer homotopy,\allowbreak{} its total homotopy is \(H|_{A^{p,q}}=(-1)^ph^{p,q}\). The four terms of \(DH+HD\) on this summand are \[
(-1)^pd_1h+d_2h+(-1)^{p+1}hd_1+hd_2.
\] The first and third cancel because each \(h^q\) is a horizontal chain map. The other two are precisely \(f^{p,q}-g^{p,q}\). For a split sequence of these outer complexes,\allowbreak{} totalizing the actual sections and retractions gives a split sequence degree by degree. The connecting map on \(C^{p,q}\) is \[
\pi^{p+1,q}d_1s^{p,q}
+(-1)^p\pi^{p,q+1}d_2s^{p,q}
=(-1)^p\pi^{p,q+1}d_2s^{p,q}.
\] The omitted first term is zero for a reason:\allowbreak{} the section commutes with \(d_1\),\allowbreak{} and \(\pi s=0\). The remaining term is exactly \(\xi_A T(\delta)\). Hence images of the distinguished split triangles have the required boundary.

For the second convention,\allowbreak{} with outer degree \(p\),\allowbreak{} the shift is \([1,0]\). The comparison \(T(A[1,0])\to T(A)[1]\) is the identity on the reindexed summands. For an outer homotopy \(h^{p,q}:A^{p,q}\to B^{p-1,q}\),\allowbreak{} use \(H=h\). The four terms are \[
d_1h+(-1)^{p-1}d_2h+hd_1+(-1)^phd_2.
\] Here the two vertical terms cancel,\allowbreak{} and the horizontal pair gives \(f-g\). The connecting map of a totalized split sequence is \(\pi^{p+1,q}d_1s^{p,q}\):\allowbreak{} its vertical term vanishes because the sections are vertical chain maps. These are the two exactness claims at 3487–\allowbreak{}3532,\allowbreak{} with the full original comparisons retained. They agree with the previously published Homology review,\allowbreak{} rather than defining a replacement convention.

\subsection{2. DERIVED-NEW-013 and DERIVED-UNDER-004}\label{proofs-3487-4161-derived-new-013-and-derived-under-004}

For fixed \(X^\bullet\in\operatorname{Comp}(\mathcal A)\),\allowbreak{} the variable \(Y^\bullet\) belongs to \(\operatorname{Comp}(\mathcal B)\). Consequently the domain at 3572 is \(\operatorname{Comp}(\mathcal B)\). For fixed \(Y^\bullet\),\allowbreak{} the domain of the functor varying \(X^\bullet\),\allowbreak{} at 3574,\allowbreak{} is \(\operatorname{Comp}(\mathcal A)\). Those two source categories are interchanged in the printed proof. The two theorem assertions at 3555–\allowbreak{}3568 already have the correct domains. \textbf{DERIVED-NEW-013} repairs just the two domains in the proof.

Bilinearity supplies all needed identities:\allowbreak{} \(1_X\otimes h\) is the outer homotopy for the first functor,\allowbreak{} the horizontal chain-map identity follows from functoriality in the other variable,\allowbreak{} and the five biproduct splitting equations are preserved by \(X^p\otimes-\). The second variable-fixed construction uses \(h\otimes1_Y\) and the same equations. Thus both functors into the respective double-complex homotopy categories are exact with their literal outer shifts. Composing with §1 proves the totalized claims. The first totalized comparison has component \((-1)^p\) on \(X^p\otimes Y^{q+1}\); the second has component identity on \(X^{p+1}\otimes Y^q\). The direct receiving tensor lemmas in More on Algebra,\allowbreak{} Cohomology and Cohomology on Sites therefore keep their original conclusions and all flatness-related distinctions.

The same argument yields an editorial extension without requiring all countable coproducts. Let \(\mathcal E\) be an additive category. Inside either of the two homotopy categories of double complexes,\allowbreak{} take the full subcategory on those actual double complexes \(A\) for which every displayed diagonal coproduct \(\bigoplus_{p+q=n}A^{p,q}\) exists. This full subcategory is triangulated,\allowbreak{} and totalization on it is exact with the same comparisons.

Here is the closure and exactness proof. Diagonal sums of \(A[0,1]\) and \(A[1,0]\) are reindexings of the next diagonal sum of \(A\); the inverse shifts work in the same way. For a first-convention cone of \(f:A\to B\),\allowbreak{} the terms are \(B^{p,q}\oplus A^{p,q+1}\),\allowbreak{} so its diagonal sum exists and is canonically the finite biproduct \(T(B)^n\oplus T(A)^{n+1}\). For the second convention they are \(B^{p,q}\oplus A^{p+1,q}\),\allowbreak{} with the same conclusion. These identifications follow from the coproduct universal property,\allowbreak{} without discarding or identifying any labelled summands. The zero object and finite direct sums also remain in the subcategory. The already proved triangulated-subcategory criterion therefore applies.

In the first convention the isomorphism from the totalized outer cone to \(C(T(f))\) is identity on each \(B^{p,q}\) and \((-1)^p\) on each \(A^{p,q+1}\). The totalized outer-cone differential has upper-right entry \((-1)^pf\),\allowbreak{} source-part horizontal differential \(d_1\),\allowbreak{} and source-part vertical differential \(-(-1)^pd_2\). After the stated map,\allowbreak{} these become respectively \(f\),\allowbreak{} \(-d_1\),\allowbreak{} and \(-(-1)^pd_2\),\allowbreak{} the three entries required by the original cone of \(T(f)\). The changing sign between horizontal source indices accounts for \(-d_1\). The inclusion square is identity,\allowbreak{} and the last square is the exact comparison \(\xi_A\) on the source shift,\allowbreak{} with the retained negative projection. In the second convention the comparison is identity on both sets of summands:\allowbreak{} the outer cone already has source-part differential \(-d_1+(-1)^{p}d_2\) at new index \(p\),\allowbreak{} which is \(-d_1-(-1)^{p+1}d_2\) at its old source index. Its upper-right entry is \(f\). Thus totalization preserves distinguished cone triangles and is exact in both cases.

For a fixed tensor factor,\allowbreak{} restrict the variable-complex homotopy category to those objects for which the tensor double complex has these diagonal coproducts. Bilinearity and the same cone and shift calculation prove closure and exactness on that full triangulated subcategory. In particular,\allowbreak{} when both factors are bounded on the same side,\allowbreak{} every diagonal has finite support and finite biproducts suffice. For lower bounds \(p\geq a,q\geq b\),\allowbreak{} a fixed diagonal satisfies \(a\leq p\leq n-b\); for upper bounds it satisfies \(n-b\leq p\leq a\). These are the actual finite index intervals.

This is \textbf{DERIVED-UNDER-004},\allowbreak{} propagating the diagonal-coproduct calculation in \textbf{HOMOLOGY-UNDER-009} to exact functors and their domains. Existence of the sums remains essential. The Homology counterexample in finite-dimensional vector spaces,\allowbreak{} with identity horizontal pairs in bidegrees \((2i,-2i),(2i+1,-2i)\),\allowbreak{} has no degree-zero diagonal coproduct. It is excluded by the precise domain just defined. Also that domain need not be closed under arbitrary outer homotopy equivalences:\allowbreak{} the same example is outer contractible. No claim of strict fullness is made. No source theorem is silently replaced by this editorial extension.

\subsection{\texorpdfstring{3. Cohomology,\allowbreak{} localization and bounded models}{3.  localization and bounded models}}\label{proofs-3487-4161-cohomology-localization-and-bounded-models}

For a termwise split short exact sequence,\allowbreak{} its connecting morphism is \(\delta^n=\pi^{n+1}d_B^ns^n\). On a cycle,\allowbreak{} its image in \(H^{n+1}(A)\) is the snake-lemma boundary obtained by lifting along \(s^n\),\allowbreak{} applying \(d_B^n\),\allowbreak{} and retracting along \(\pi^{n+1}\). Independence follows because a different lift adds an \(A^n\)-component and therefore changes the result by an \(A\)-boundary. The general abelian-category assertion is the same pullback/\allowbreak{}kernel calculation in the referenced Homology lemma,\allowbreak{} not an assumption that a morphism from every test object can be lifted. Consequently \(H^0\) is homological and the source\textquotesingle s long sequence has boundary \(H^i(\delta)\) with its displayed sign.

\textbf{OCC-01499 /\allowbreak{} P02-E0027} and French \textbf{DERIVED-016} correctly request deletion of the unfinished reference placeholder at 3651; \textbf{MC-STK-ERR-0827} already does this,\allowbreak{} retaining both actual Homology references. French \textbf{DERIVED-017} matches \textbf{MC-STK-ERR-0828},\allowbreak{} the insertion of ``by'' at 3702. The functor being named is the unique factorization of the original \(H^0\).

The shifted kernel of \(H^0\) is exactly the acyclic complexes. For a distinguished triangle,\allowbreak{} the long sequence proves closure under cones and both shifts,\allowbreak{} and additivity proves closure under summands. The cone of a map is acyclic precisely when its cohomology maps are isomorphisms. Thus the localization denominators are exactly the quasi-isomorphisms,\allowbreak{} its kernel is exactly the acyclic subcategory,\allowbreak{} and its canonical shift comparison is identity. The quotient and bounded-subcategory constructions established earlier apply with these exact objects.

The size warning in the next remark is substantive. In the cited example,\allowbreak{} kernels and cokernels of an equivariant abelian-group map inherit each ordinal-indexed operator,\allowbreak{} and extending the shorter operator list by zeros gives the necessary common index. Finite biproducts do the same. This verifies the abelian structure used there. For the extensions of \(Z\) by \(Z\),\allowbreak{} an extension isomorphism has underlying matrix \(\begin{psmallmatrix}1&t\\0&1\end{psmallmatrix}\). If the nonzero operator of one extension occurs at ordinal \(\alpha\) and that of another at \(\beta\ne\alpha\),\allowbreak{} intertwining the operator at \(\alpha\) would require the product of that invertible matrix with \(\begin{psmallmatrix}0&1\\0&0\end{psmallmatrix}\) to be zero. It is not zero. Thus the distinct ordinals yield distinct extension classes. The use of their identification with derived Hom is the cited later Ext lemma,\allowbreak{} whose complete source-order review remains later; no new set-sized Hom claim is made for this large category.

For a Grothendieck abelian category the source invokes the specific K-injective existence theorem in Injectives. Its statement and the start of its construction were read here; its entire transfinite proof was not reaudited. The receiving comparison is exact:\allowbreak{} if \(I\) is K-injective and \(s:M\to N\) is a quasi-isomorphism,\allowbreak{} its acyclic cone and all shifts have zero Hom into \(I\). The long Hom sequence makes \(\operatorname{Hom}_K(N,I)\to\operatorname{Hom}_K(M,I)\) bijective. Every roof \(N\leftarrow M\to I\) therefore has a unique representative \(N\to I\); common refinements prove that this identification is independent of the roof. Composing with the selected quasi-isomorphism \(L\to I\) gives exactly the Hom comparison at 3746–\allowbreak{}3747,\allowbreak{} and its right side is a set.

The bounded models are obtained with their actual boundary objects:\allowbreak{} \[
(\tau_{\leq b}K)^i=
\begin{cases}K^i&i<b,\\\ker d_K^b&i=b,\\0&i>b,\end{cases}
\quad
(\tau_{\geq a}K)^i=
\begin{cases}0&i<a,\\\operatorname{coker}d_K^{a-1}&i=a,\\K^i&i>a.
\end{cases}
\] Their differentials are the original ones or their induced kernel and cokernel maps. Choose \(a<b\) outside the nonzero cohomology range. The two truncations commute on each degree,\allowbreak{} including the distinct boundary degrees,\allowbreak{} so their composite is the original \(N\) in the four-corner diagram. All four maps induce the identity on the surviving cohomology and zero-to-zero outside it.

For bounded-below complexes a left fraction \(s^{-1}f\) can have its common target replaced by its lower truncation,\allowbreak{} since that target is quasi-isomorphic to the bounded-below denominator source. This gives fullness. If the fraction is zero,\allowbreak{} a further denominator kills \(f\); lower-truncate its target again to give the same equality inside the bounded category. Additivity reduces equality of arbitrary fractions to this zero test,\allowbreak{} proving faithfulness. For bounded-above complexes use right fractions and upper truncation of their common source. For bounded complexes perform the bounded-below argument inside the bounded-above category; lower truncation preserves that upper bound. This supplies all three equivalences,\allowbreak{} with the kernel and essential surjectivity already verified.

\subsection{4. The canonical delta-functor and truncation triangles}\label{proofs-3487-4161-the-canonical-delta-functor-and-truncation-triangles}

For the original short exact sequence with maps \(a,b\),\allowbreak{} put \(q=(b,0):C(a)\to C\). Its kernel is identified with \(C(1_A)\) by \((x,y)\mapsto(a(x),y)\). The differential under that identification is \(\begin{psmallmatrix}d_A&1\\0&-d_A\end{psmallmatrix}\). The homotopy \(H^n(x,y)=(0,x)\) has \(dH+Hd=(x,y)\),\allowbreak{} so the kernel is contractible. The exact cohomology sequence makes \(q\) a quasi-isomorphism.

For a commuting map of short exact sequences \((f,g,h)\),\allowbreak{} the canonical cone map in degree \(n\) is \(c^n=\begin{psmallmatrix}g^n&0\\0&f^{n+1}\end{psmallmatrix}\). It satisfies \(q'c=hq\),\allowbreak{} \(p'c=f[1]p\). Hence,\allowbreak{} in the derived category,\allowbreak{} \[
(-p'(q')^{-1})h=-p'c q^{-1}=f[1](-pq^{-1}).
\] This proves the actual boundary naturality,\allowbreak{} including the shift. The signed cone triangle and the isomorphism \(q\) prove the distinguishedness. This is a further checked receiving use of DERIVED-NEW-007. All bounded variants preserve the cone\textquotesingle s one shift and finite sum. French \textbf{DERIVED-018} matches \textbf{MC-STK-ERR-0829},\allowbreak{} adding ``above'' at 3968 without changing this construction.

When all three vertical maps are quasi-isomorphisms,\allowbreak{} this same naturality gives the claimed triangle isomorphism. In the split case,\allowbreak{} the earlier homotopy inverse \(S=(s,-\delta)\) satisfies \(qS=1\),\allowbreak{} \(Sq\simeq1\),\allowbreak{} and \(-pS=\delta\). Thus the canonical boundary is exactly the image of the original split boundary. \textbf{OCC-01500 /\allowbreak{} P02-E0028},\allowbreak{} French \textbf{DERIVED-019},\allowbreak{} and French \textbf{DERIVED-020} match \textbf{MC-STK-ERR-0830},\allowbreak{} \textbf{0831},\allowbreak{} the two singular-sequence copyedits.

For the first truncation triangle,\allowbreak{} write \(U=K/\tau_{\leq a}K\). The natural map \(U\to\tau_{\geq a+1}K\) has kernel concentrated in degrees \(a,a+1\),\allowbreak{} with both terms \(\operatorname{im}d_K^a\) and differential identity under the canonical image/\allowbreak{}coimage identification. Hence the kernel is acyclic and the map is a quasi-isomorphism. The canonical delta-functor triangle transports across that precise map. Applying the same construction to \(\tau_{\leq a+1}K\) leaves only the quotient cohomology \(H^{a+1}(K)[-a-1]\). For the lower truncations,\allowbreak{} the kernel of \(\tau_{\geq a}K\to\tau_{\geq a+1}K\) has term \(K^a/\operatorname{im}d^{a-1}\) at \(a\),\allowbreak{} term \(\operatorname{im}d^a\) at \(a+1\),\allowbreak{} and the induced surjection between them. Inclusion of its kernel \(H^a(K)[-a]\) is a quasi-isomorphism. This proves the third triangle and its final shift \([-a+1]\),\allowbreak{} with all original quotient contributions retained.

\subsection{\texorpdfstring{5. DERIVED-UNDER-005:\allowbreak{} an explicit bound and derived morphisms}{5.  an explicit bound and derived morphisms}}\label{proofs-3487-4161-derived-under-005-an-explicit-bound-and-derived-morphisms}

The induction at 4154–\allowbreak{}4160 applies its one-step case after shifting the cohomological cutoff. Its full form is the following. Let \(f_j:K_j\to K_{j+1}\),\allowbreak{} \(0\leq j<n\),\allowbreak{} be morphisms in \(D(\mathcal A)\),\allowbreak{} and fix an integer \(b\). Assume \[
H^i(K_0)=0\quad(i>b),\qquad H^{b-j}(f_j)=0\quad(0\leq j<n).
\] Then the composite factors through \(\tau_{\leq b-n}K_n\to K_n\). The intermediate complexes need not have any boundedness assumption. This includes arbitrary derived morphisms,\allowbreak{} not only the chain maps in the source statement.

We prove the needed one-step assertion without assuming a representative exists on the original pair of complexes. If \(X\) has cohomology zero above \(m\),\allowbreak{} represent \(f:X\to Y\) by a roof \(X\xleftarrow{s}Z\xrightarrow{g}Y\). The map \(\tau_{\leq m}Z\to Z\) is a quasi-isomorphism. The restricted chain map factors through \(\tau_{\leq m}Y\),\allowbreak{} because its degree-\(m\) component carries cycles to cycles. The composite \(\tau_{\leq m}Z\to H^m(Y)[-m]\) equals the canonical projection to \(H^m(Z)[-m]\) followed by \(H^m(g)\). If \(H^m(f)=0\),\allowbreak{} then \(H^m(g)=H^m(f)H^m(s)=0\). Thus that composite is zero as a chain map. Applying the exact representable Hom sequence to the canonical triangle \[
\tau_{\leq m-1}Y\to\tau_{\leq m}Y\to H^m(Y)[-m]
\] gives the required lift \(X\to\tau_{\leq m-1}Y\).

These truncation constructions are natural on the derived category. On complexes their kernel/\allowbreak{}cokernel definitions are functorial and preserve quasi-isomorphisms by the displayed cohomology formulas. They also preserve homotopies:\allowbreak{} at the upper cutoff restrict the last homotopy component to the kernel; at the lower cutoff project its first nonzero target component to the cokernel. The boundary equations are exactly the original \(f-g=dh+hd\),\allowbreak{} restricted or projected. The universal property of localization therefore supplies the induced functors and their canonical maps used above.

Starting with \(u_0:K_0\to\tau_{\leq b}K_0\),\allowbreak{} the inverse of the canonical quasi-isomorphism,\allowbreak{} suppose the first \(j\) maps factor as \(u_j:K_0\to\tau_{\leq b-j}K_j\) followed by the canonical inclusion. Apply the one-step assertion at \(m=b-j\) to \[
\tau_{\leq b-j}K_j\longrightarrow K_j\xrightarrow{f_j}K_{j+1}.
\] Its degree-\(b-j\) cohomology map vanishes by hypothesis. It therefore factors through \(\tau_{\leq b-j-1}K_{j+1}\). Composing that factorization with \(u_j\) constructs \(u_{j+1}\). Induction proves the stated result. For the source\textquotesingle s \(n=2\) step this is the application at cutoff \(-1\),\allowbreak{} equivalently the cutoff-zero case on complexes shifted by \([-1]\). It is an implicit shift in the source,\allowbreak{} not a failure of its conclusion.

The dual assertion says that for arrows \(K_n\to\cdots\to K_0\),\allowbreak{} with \(H^i(K_0)=0\) below \(a\) and \(H^{a+j}(K_{j+1}\to K_j)=0\),\allowbreak{} the composite factors through \(K_n\to\tau_{\geq a+n}K_n\). To obtain it,\allowbreak{} reverse arrows and degrees:\allowbreak{} a cochain complex in \(\mathcal A^{op}\) associated to \(K\) has degree \(i\) object \(K^{-i}\) and differential the opposite of the original map \(d_K^{-i-1}\). Cohomology in degree \(i\) is the opposite of \(H^{-i}(K)\); quasi-isomorphisms and their roofs are reversed,\allowbreak{} and upper truncation becomes lower truncation under this operation. Applying the just-proved factorization with bound \(-a\) and reversing it back gives exactly the claimed arrow and indices.

Now suppose additionally that \(H^i(K_n)=0\) for \(i<a\) and \(n>b-a\). The receiving truncation \(\tau_{\leq b-n}K_n\) is acyclic,\allowbreak{} hence zero in the derived category. The composite is zero. In particular,\allowbreak{} for a complex with cohomology in \([a,b]\),\allowbreak{} a composite of \(b-a+1\) endomorphisms that each act by zero on all cohomology is zero. These endomorphisms form an ideal,\allowbreak{} since cohomology is an additive functor. Its \((b-a+1)\)-st power is therefore zero. If \(u\) is one of them,\allowbreak{} the complete finite formula \[
(1-u)^{-1}=1+u+u^2+\cdots+u^{b-a}
\] follows by multiplying on either side and retaining the terminal term \(-u^{b-a+1}=0\). No global dimension or existence of projective/\allowbreak{}injective resolutions is used for this consequence.

This is \textbf{DERIVED-UNDER-005}. It is a proved corollary and extension of the original factorization argument,\allowbreak{} not a claim of a new theorem in the literature. Its source theorem stays unchanged.

\subsection{6. Propagation and DERIVED-NEW-014}\label{proofs-3487-4161-propagation-and-derived-new-014}

In \texttt{more-\allowbreak{}algebra.\allowbreak{}tex},\allowbreak{} current 23390–\allowbreak{}23402 (original 22711–\allowbreak{}22723)\allowbreak{},\allowbreak{} the earlier \textbf{MC-STK-ERR-0959} corrected both reversed composites to \(\psi\varphi\). The reports \textbf{OCC-01806 /\allowbreak{} P02-E0339},\allowbreak{} \textbf{OCC-02517 /\allowbreak{} P02-E1050},\allowbreak{} and \textbf{OCC-12698 /\allowbreak{} MORE-ALGEBRA-G-016} concern that order and do not ask to correct the receiving truncation. Their original report texts were read here; this is a propagation check,\allowbreak{} not a new whole-chapter report closure.

With the original labels \(K_1\xrightarrow\varphi K_2
\xrightarrow\psi K_3\),\allowbreak{} §5 applies at \(b=0,n=2\). It factors the composite through \textbf{\(\tau_{\leq-2}K_3\)}. The proof currently writes \(K_2\). Since both truncations vanish,\allowbreak{} the displayed existence of some factorization through zero is not a counterexample to the theorem. Nevertheless the cited argument\textquotesingle s receiving object is \(K_3\); \textbf{DERIVED-NEW-014} corrects that index in the proof and retains the earlier order corrections. The theorem itself needs only the source upper bound on \(K_1\),\allowbreak{} the final lower bound on \(K_3\),\allowbreak{} and the two specified vanishing cohomology maps; §5 proves this stronger editorial statement without requiring \(K_2\) to have amplitude \([-1,0]\).

The other direct uses were checked as receiving applications:\allowbreak{}

\begin{itemize}
\tightlist
\item
  In \texttt{cotangent.\allowbreak{}tex:\allowbreak{}2250–\allowbreak{}2288},\allowbreak{} the original \(c=k^2\) is retained. There are \(k\) transitions,\allowbreak{} each lowering the ideal index by \(k\),\allowbreak{} and each kills degrees \(0,-1,\ldots,-k+1\). The composite factors through the target cutoff \(-k\),\allowbreak{} exactly as asserted. For \(k=0\) the assertion is only the existing vanishing in positive cohomological degrees; no zero-fold composition is claimed to vanish.
\item
  In \texttt{local-\allowbreak{}cohomology.\allowbreak{}tex:\allowbreak{}3464–\allowbreak{}3484},\allowbreak{} the truncated dualizing complex has amplitude \([-d+1,0]\). Multiplication by the chosen \(f\) is zero on its cohomology,\allowbreak{} so its \(d\)-th power is zero. For \(d=0\) that truncated object is already zero,\allowbreak{} consistent with the conclusion.
\item
  In \texttt{local-\allowbreak{}cohomology.\allowbreak{}tex:\allowbreak{}4375–\allowbreak{}4393},\allowbreak{} the homological cone of the augmentation of the Koszul complex occupies degrees \(0\) through \(t\). Its cochain amplitude is therefore contained in \([-t,0]\). If the recorded \(a_1^c\) kills its cohomology,\allowbreak{} the explicit sufficient exponent is \(c(t+1)\). This proves the receiving step that enlarges \(c\); it does not claim to optimize the final ideal exponent in the whole lemma.
\item
  In \texttt{local-\allowbreak{}cohomology.\allowbreak{}tex:\allowbreak{}4685–\allowbreak{}4701},\allowbreak{} the \(t\) transitions kill negative degrees \(-t,\ldots,-1\). Apply the dual assertion with \(a=-t\) to the truncated objects to factor through the degree-zero quotient \(M/I^{n+tc'}M\) of the source. The displayed transitions at index \(n+tc'\) and the later index \(n-c\) are retained; no change of the constants or source objects is made.
\end{itemize}

These are bounded checks of the use of the factorization lemma. They do not certify every surrounding argument or report in those chapters. The broader joint synthesis remains after the core review queue.
\clearpage
\section{Filtered derived objects and the domain of a derived functor}\label{proofs-4166-5089-filtered-derived-objects-and-the-domain-of-a-derived-functor}

This review concerns original \texttt{derived.\allowbreak{}tex} lines 4166–\allowbreak{}5089 at commit \texttt{a04446e5\allowbreak{}7ec1fbc2\allowbreak{}52a871af\allowbreak{}cec7752f\allowbreak{}b2807b14}. The immutable original and current chapter snapshots in this intake remain the comparison objects. The current chapter still has the mathematical defects identified below. All constructions retain the source\textquotesingle s decreasing filtrations,\allowbreak{} cochain differentials,\allowbreak{} fraction direction,\allowbreak{} and shift comparisons. Supplemental arguments and strengthenings are editorial; they do not replace the source\textquotesingle s exposition. No novelty is claimed.

\subsection{1. Finite filtrations in each degree and the target of gr}\label{proofs-4166-5089-finite-filtrations-in-each-degree-and-the-target-of-gr}

In an abelian category A,\allowbreak{} an object of Fil\textasciicircum{}f(A)\allowbreak{} is an object M with a decreasing filtration F\^{}p M,\allowbreak{} equal to M for sufficiently small p and to zero for sufficiently large p. Its p-th graded object is F\^{}p M/\allowbreak{}F\textasciicircum{}\{p\allowbreak{}+1\} M. Gr(A)\allowbreak{} is the category of families indexed by p,\allowbreak{} with morphisms and abelian operations computed componentwise. In particular,\allowbreak{} the notation gr(M)\allowbreak{} does not require a countable coproduct in A.

Each gr\textasciicircum{}p,\allowbreak{} gr,\allowbreak{} and the forgetful functor is additive. An additive functor preserves the actual biproducts used in shifts and cones. For a map u of cochain complexes,\allowbreak{} its cone in degree n is K\textasciicircum{}\{n\allowbreak{}+1\} direct-sum L\^{}n with differential (-d\_\allowbreak{}K,\allowbreak{}0;u,\allowbreak{}d\_\allowbreak{}L)\allowbreak{}. Applying the functor gives precisely the cone matrix for its image. The shift comparison is the componentwise identity. Thus all three functors induce exact functors on homotopy categories. The composition of gr with H\^{}0 is homological because cohomology of a cone gives the usual long exact sequence in Gr(A)\allowbreak{}. The same argument applies to gr\^{}p and to forgetting F.

Consequently the complexes with H\^{}n(gr K)\allowbreak{}\allowbreak{}=0 for every n form the strictly full saturated triangulated kernel described by the previously proved homological-kernel construction. A morphism is a filtered quasi-isomorphism exactly when its cone belongs to that kernel:\allowbreak{} apply the long exact graded cohomology sequence in every degree. This proves the source\textquotesingle s statements about FAc,\allowbreak{} FQis,\allowbreak{} the quotient,\allowbreak{} and the factorization of H\^{}0 gr. Kernels,\allowbreak{} cycles,\allowbreak{} and cohomology here are taken in Gr(A)\allowbreak{},\allowbreak{} which is abelian; no abelian structure on Fil\textasciicircum{}f(A)\allowbreak{} is being presumed.

For completeness,\allowbreak{} the finite-filtration exactness argument used to forget F is local in cochain degree. For three finite filtered objects A0 \allowbreak{}-> A1 \allowbreak{}-> A2,\allowbreak{} exactness of every graded sequence implies exactness of the underlying sequence and of the sequences of filtration steps and filtration quotients. To prove this,\allowbreak{} choose a common finite filtration interval for these three objects. Remove its highest nonzero step. The highest-step complex is the corresponding graded complex,\allowbreak{} and the quotient has a shorter filtration. Induction applies to the quotient; the termwise short exact sequence from highest steps to the original complex to the quotient then gives exactness at the middle. That last assertion follows directly by lifting a middle cycle to the left term of the quotient,\allowbreak{} subtracting its lifted boundary,\allowbreak{} and lifting the resulting highest-step cycle. In an arbitrary abelian category these lifts are obtained after epimorphic pullbacks; images descend along these epimorphisms. The complete test-object argument is in \texttt{.\allowbreak{}.\allowbreak{}/\allowbreak{}HOMOLOGY\_\allowbreak{}INTAKE\_\allowbreak{}20260930/\allowbreak{}PROOFS\_\allowbreak{}4681\_\allowbreak{}5648.\allowbreak{}md},\allowbreak{} lines 265–\allowbreak{}281. Applying the same induction to the induced step filtrations and quotient filtrations gives their exactness as well.

Write f:\allowbreak{}A0 \allowbreak{}-> A1 and g:\allowbreak{}A1 \allowbreak{}-> A2. Step exactness yields

\[
 f(F^p A_0)=\ker(g)\cap F^p A_1=f(A_0)\cap F^p A_1.
\]

Quotient exactness yields

\[
 g^{-1}(F^p A_2)=F^p A_1+\ker(g),
\]

which implies g(F\^{}p A1)\allowbreak{}\allowbreak{}=g(A1)\allowbreak{} intersect F\^{}p A2. Thus both maps are strict. These are the two strictness conclusions actually needed later; exactness of the underlying sequence alone would not suffice.

Apply this three-term result to K\^{}\{n-1\} \allowbreak{}-> K\^{}n \allowbreak{}-> K\textasciicircum{}\{n\allowbreak{}+1\} for each n when gr K is acyclic. It proves H\textasciicircum{}n(K)\allowbreak{}\allowbreak{}=0. The common filtration interval may depend on n:\allowbreak{} the source\textquotesingle s degreewise finiteness suffices. Applying this result to a filtered cone proves that forgetting F sends FQis to quasi-isomorphisms. Therefore localization gives exactly

\[
 \operatorname{gr}^p,\operatorname{forget}F:DF(A)\longrightarrow D(A),
 \qquad \operatorname{gr}:DF(A)\longrightarrow D(\operatorname{Gr}(A)).
\]

DERIVED-RECON-022 reconciles French DERIVED-021 with the existing MC-STK-ERR-0832 deletion of the stray words at original 4361. Its mathematical argument is valid for the reason just proved.

DERIVED-NEW-015:\allowbreak{} original 4379,\allowbreak{} still current 4732,\allowbreak{} puts gr(X)\allowbreak{} in D\textasciicircum{}b(A)\allowbreak{},\allowbreak{} although the immediately preceding functor has target D(Gr(A)\allowbreak{})\allowbreak{}. The correctly typed condition is gr(X)\allowbreak{} in D\textasciicircum{}b(Gr(A)\allowbreak{})\allowbreak{}. This is boundedness in the cochain degree,\allowbreak{} uniformly across graded components; it is not a separate,\allowbreak{} component-dependent bound for each p. The analogous plus/\allowbreak{}minus conditions have the same graded target. The following bounded-model proof uses precisely this condition.

\subsection{2. Bounded filtered models and full faithfulness}\label{proofs-4166-5089-bounded-filtered-models-and-full-faithfulness}

Suppose H\^{}n(gr K)\allowbreak{}\allowbreak{}=0 for n\textless a. The three-term result applied with middle term K\^{}\{a-1\} shows that d\^{}\{a-1\} is strict. Put I\allowbreak{}=im(d\textasciicircum{}\{a-1\})\allowbreak{},\allowbreak{} with the induced filtration from K\^{}a. Strictness identifies the quotient filtration on its coimage with this filtration,\allowbreak{} so

\[
 \operatorname{gr}(I)=\operatorname{im}(\operatorname{gr}(d^{a-1})),
 \quad
 \operatorname{gr}(K^a/I)=
 \operatorname{gr}(K^a)/\operatorname{im}(\operatorname{gr}(d^{a-1})).
\]

The second identity follows by the exact graded sequence of an induced subobject and its filtered quotient. Define L\textasciicircum{}n\allowbreak{}=0 for n<a,\allowbreak{} L\textasciicircum{}a\allowbreak{}=K\textasciicircum{}a/\allowbreak{}I,\allowbreak{} and L\textasciicircum{}n\allowbreak{}=K\textasciicircum{}n for n>a,\allowbreak{} with the induced differential at a and the original differentials thereafter. All terms have finite filtrations. The original quotient map K \allowbreak{}-> L is a filtered chain map. Its image under gr is the canonical lower truncation map of gr K; the vanishing assumption makes that map a quasi-isomorphism.

For the other direction suppose H\^{}n(gr K)\allowbreak{}\allowbreak{}=0 for n\textgreater b. The three-term argument with middle term K\textasciicircum{}\{b\allowbreak{}+1\} makes d\^{}b strict. Give J\allowbreak{}=ker(d\textasciicircum{}b)\allowbreak{} the induced filtration in K\^{}b. Strictness gives gr J\allowbreak{}=ker(gr d\textasciicircum{}b)\allowbreak{}. Set M\textasciicircum{}b\allowbreak{}=J,\allowbreak{} M\textasciicircum{}n\allowbreak{}=K\textasciicircum{}n for n<b,\allowbreak{} and M\textasciicircum{}n\allowbreak{}=0 for n\textgreater b. Its differential into J is the factorization of d\^{}\{b-1\}; the inclusion M \allowbreak{}-> K becomes the upper truncation quasi-isomorphism under gr.

If graded cohomology is bounded,\allowbreak{} choose a\textless b enclosing its support; enlarging the chosen interval ensures this strict inequality without changing K. The two constructions act at different degrees. Thus N\allowbreak{}=tau\_\allowbreak{}\{>\allowbreak{}=a\}tau\_\allowbreak{}\{<\allowbreak{}=b\}K\allowbreak{}=tau\_\allowbreak{}\{<\allowbreak{}=b\}tau\_\allowbreak{}\{>\allowbreak{}=a\}K,\allowbreak{} with the quotient term K\textasciicircum{}a/\allowbreak{}im(d\textasciicircum{}\{a-1\})\allowbreak{},\allowbreak{} the kernel term ker(d\textasciicircum{}b)\allowbreak{},\allowbreak{} the original intermediate terms,\allowbreak{} and induced differentials. The quotient and inclusion maps give the source\textquotesingle s actual commutative square

\[
 \begin{array}{ccc}
 K&\longrightarrow&L=\tau_{\ge a}K\\
 \uparrow&&\uparrow\\
 M=\tau_{\le b}K&\longrightarrow&N.
 \end{array}
\]

All four maps are filtered quasi-isomorphisms by the two truncation arguments. This supplies the source\textquotesingle s omitted verification of part (3)\allowbreak{}. In particular N has only finitely many nonzero terms,\allowbreak{} so their finitely many filtration bounds give a common finite filtration interval for N. This is a consequence about the representative just constructed,\allowbreak{} not a claim that the original K has a uniform filtration bound.

The bounded localization assertions follow with the original fraction directions. For bounded-below K,\allowbreak{}L,\allowbreak{} represent a morphism in DF(A)\allowbreak{} as s\^{}\{-1\}f with f:\allowbreak{}K \allowbreak{}-> L\textquotesingle{} and s:\allowbreak{}L \allowbreak{}-> L\textquotesingle{} a filtered quasi-isomorphism. Then gr L\textquotesingle{} has cohomology bounded below. Postcompose both maps with the filtered quasi-isomorphism L\textquotesingle{} \allowbreak{}-> L\textquotesingle\textquotesingle{} just constructed,\allowbreak{} with L\textquotesingle\textquotesingle{} bounded below. This represents the same morphism in the bounded-below localization,\allowbreak{} proving fullness. If a bounded-below fraction becomes zero,\allowbreak{} the fraction equality criterion gives a denominator L\textquotesingle{} \allowbreak{}-> T annihilating f in the homotopy category. Postcompose with T \allowbreak{}-> T\textquotesingle{} for T\textquotesingle{} bounded below. The composite still annihilates f and is a denominator in the bounded-below category,\allowbreak{} proving faithfulness.

For bounded-above objects use right fractions K \textless- K\textquotesingle{} \allowbreak{}-> L and precompose with a bounded-above filtered quasi-isomorphism K\textquotesingle\textquotesingle{} \allowbreak{}-> K\textquotesingle. For a zero fraction precompose its annihilating denominator once more with such a bounded-above replacement. These prove fullness and faithfulness with all arrows reversed. Finally work inside the bounded-above category and apply the bounded-below construction; lower truncation preserves an existing upper bound. The same two fraction arguments give the bounded equivalence. Essential surjectivity is supplied by the explicit models. In all three categories the acyclic kernel is closed under shifts,\allowbreak{} cones,\allowbreak{} and existing direct summands by graded cohomology,\allowbreak{} proving the stated saturation and the identification of the multiplicative systems.

\subsection{\texorpdfstring{3. Derived functor maps,\allowbreak{} denominators,\allowbreak{} and shifts}{3. Derived functor   and shifts}}\label{proofs-4166-5089-derived-functor-maps-denominators-and-shifts}

Keep the source\textquotesingle s exact functor F:\allowbreak{}D \allowbreak{}-> D\textquotesingle{} and saturated multiplicative system S. For X in D and W in D',\allowbreak{} define the abelian group

\[
 T_X(W)=\mathop{\operatorname{colim}}_{s:X\to X'\ \in X/S}
             \operatorname{Hom}_{D'}(W,F(X')).                 \tag{3.1}
\]

This is a colimit of abelian groups,\allowbreak{} not a presumed colimit in D\textquotesingle. The latter need not exist. An essentially constant system with value A has T\_\allowbreak{}X naturally isomorphic to Hom(-,\allowbreak{}A)\allowbreak{}. Conversely if T\_\allowbreak{}X is represented by A,\allowbreak{} the class corresponding to 1\_\allowbreak{}A has a representative alpha:\allowbreak{}A \allowbreak{}-> F(X')\allowbreak{} at some index. For every W its induced map Hom(W,\allowbreak{}A)\allowbreak{} \allowbreak{}-> T\_\allowbreak{}X(W)\allowbreak{} is the representing isomorphism,\allowbreak{} by naturality. This is exactly characterization (4)\allowbreak{} in Categories,\allowbreak{} lemma-characterize-essentially-constant-ind,\allowbreak{} so the system is essentially constant with value A. This also describes why ordinary existence of a colimit in D\textquotesingle{} is a weaker condition.

For f:\allowbreak{}X \allowbreak{}-> Y and an index s:\allowbreak{}X \allowbreak{}-> X',\allowbreak{} MS2 gives t:\allowbreak{}Y \allowbreak{}-> Y\textquotesingle{} in S and f':\allowbreak{}X' \allowbreak{}-> Y\textquotesingle{} with f's\allowbreak{}=tf. Map a class [a:\allowbreak{}W \allowbreak{}-> F(X')\allowbreak{}]\allowbreak{} to [F(f')\allowbreak{}a]\allowbreak{} in T\_\allowbreak{}Y(W)\allowbreak{}. For two such squares,\allowbreak{} first take a common refinement of their Y-indices. Their two maps out of X\textquotesingle{} then agree after precomposition with s. MS3 supplies a further denominator equalizing them. This proves independence. Applying the same argument to a transition X\textquotesingle{} \allowbreak{}-> X\textquotesingle\textquotesingle{} proves compatibility with the colimit. Composing squares proves T(gf)\allowbreak{}\allowbreak{}=T(g)\allowbreak{}T(f)\allowbreak{},\allowbreak{} and the identity square proves T(1)\allowbreak{}\allowbreak{}=1. For f\allowbreak{}+g take common denominator squares for f and g; the sum of the refined maps is a square for f\allowbreak{}+g. Additivity of F gives T(f\allowbreak{}+g)\allowbreak{}\allowbreak{}=T(f)\allowbreak{}\allowbreak{}+T(g)\allowbreak{}.

When T\_\allowbreak{}X and T\_\allowbreak{}Y are represented,\allowbreak{} Yoneda gives the unique RF(f)\allowbreak{} with the source\textquotesingle s commuting squares. Equivalently,\allowbreak{} whenever the two colimits in D\textquotesingle{} exist,\allowbreak{} the same refined-square construction on their cocones proves the first source lemma. No triangle argument is needed for that weaker assertion.

If s:\allowbreak{}X \allowbreak{}-> Y belongs to S,\allowbreak{} composition gives j:\allowbreak{}Y/\allowbreak{}S \allowbreak{}-> X/\allowbreak{}S. The category Y/\allowbreak{}S identifies with the category of objects under s in X/\allowbreak{}S:\allowbreak{} a factorization t\allowbreak{}=u s with t,\allowbreak{}s in S has u in S by saturation. For a filtered category I and an object i,\allowbreak{} the forgetful functor i/\allowbreak{}I \allowbreak{}-> I is cofinal. Indeed any k and i have a common successor. For two successors with maps from k,\allowbreak{} filteredness first gives a common target and then equalizes both pairs of arrows from i and k. This proves nonemptiness and connectedness of the required comma category. It follows that j is cofinal,\allowbreak{} so (3.1)\allowbreak{} is unchanged. Its isomorphism agrees with T(s)\allowbreak{} by the defining squares. Hence representability at X and at Y is equivalent and RF(s)\allowbreak{} is the induced isomorphism.

Shift by n gives an equivalence X/\allowbreak{}S \allowbreak{}-> X[n]\allowbreak{}/\allowbreak{}S because S is closed under both shifts and their inverses. The original comparison for F,\allowbreak{} iterated with its inverse when n<0,\allowbreak{} identifies F(X')\allowbreak{}[n]\allowbreak{} with F(X'[n]\allowbreak{})\allowbreak{}. Applying this to every index and every transition gives the canonical comparison RF(X)\allowbreak{}[n]\allowbreak{} \allowbreak{}-> RF(X[n]\allowbreak{})\allowbreak{}. Its coherence and naturality follow on each cocone map from those of F,\allowbreak{} so they hold after the colimit.

For LF use instead the functor

\[
 U_X(W)=\mathop{\operatorname{colim}}_{(S/X)^{\mathrm{op}}}
                         \operatorname{Hom}_{D'}(F(X'),W).
\]

Its representability is Hom(LF(X)\allowbreak{},\allowbreak{}-)\allowbreak{}. Reverse the MS2 squares and the MS3 equalizations above. Composition with a denominator now identifies the relevant category over an object in the cofiltered category S/\allowbreak{}X; passing to the opposite category gives the same filtered cofinality proof. Yoneda reverses the represented natural transformation,\allowbreak{} producing LF(f)\allowbreak{}:\allowbreak{}LF(X)\allowbreak{} \allowbreak{}-> LF(Y)\allowbreak{}. Shift comparisons are induced by the same iterated comparisons of F. These observations prove both dual source assertions with their actual variances.

\subsection{\texorpdfstring{4. Two out of three,\allowbreak{} including the Hom comparison direction}{4. Two out of  including the Hom comparison direction}}\label{proofs-4166-5089-two-out-of-three-including-the-hom-comparison-direction}

The category I of denominator maps out of a fixed distinguished triangle is filtered and its three component functors are cofinal. The complete source-bound proof is \texttt{PROOFS\_\allowbreak{}1149\_\allowbreak{}1843.\allowbreak{}md},\allowbreak{} section 4,\allowbreak{} lines 268–\allowbreak{}318:\allowbreak{} MS2 and MS6 extend a chosen component denominator; MS3 kills the three square defects in succession without destroying earlier equalities. The same operation equalizes parallel triples and proves connectedness of each component comma category. Thus cofinality here has been proved,\allowbreak{} not inferred merely from surjectivity on objects.

Let X \allowbreak{}-> Y \allowbreak{}-> Z \allowbreak{}-> X[1]\allowbreak{} be distinguished,\allowbreak{} with RF(X)\allowbreak{}\allowbreak{}=A and RF(Y)\allowbreak{}\allowbreak{}=B. Choose a triangle A \allowbreak{}-> B \allowbreak{}-> C \allowbreak{}-> A[1]\allowbreak{} on RF(f)\allowbreak{}. Choose an essential-constancy witness alpha:\allowbreak{}A \allowbreak{}-> F(X')\allowbreak{}. MS2 gives a square to f':\allowbreak{}X' \allowbreak{}-> Y\textquotesingle{} with both structure arrows in S. Refine Y\textquotesingle{} to an index admitting a witness beta:\allowbreak{}B \allowbreak{}-> F(Y')\allowbreak{}. Write the cocone maps a:\allowbreak{}F(X')\allowbreak{} \allowbreak{}-> A and b:\allowbreak{}F(Y')\allowbreak{} \allowbreak{}-> B. They satisfy a alpha\allowbreak{}=1,\allowbreak{} b beta\allowbreak{}=1,\allowbreak{} and b F(f')\allowbreak{}\allowbreak{}=RF(f)\allowbreak{} a. Therefore the square defect

\[
 \delta=F(f')\alpha-\beta RF(f):A\longrightarrow F(Y')
 \quad\text{satisfies}\quad b\delta=0.
\]

Essential constancy gives a transition v:\allowbreak{}Y' \allowbreak{}-> Y\textquotesingle\textquotesingle{} for which F(v)\allowbreak{}\allowbreak{}=F(v)\allowbreak{} beta b. To obtain the same transition on both appearances of Y',\allowbreak{} use the witness factorization and then filteredness to equalize the two parallel index arrows. Consequently F(v)\allowbreak{}delta\allowbreak{}=0. Replace Y\textquotesingle{} by Y\textquotesingle\textquotesingle{} and beta by F(v)\allowbreak{}beta. The first square now commutes. This is the precise step requiring the indexing arrows Y \allowbreak{}-> Y',\allowbreak{} as corrected by MC-STK-ERR-0841 and reconciled in DERIVED-RECON-023.

MS6 completes the source square to a denominator triangle map,\allowbreak{} and TR3 completes (alpha,\allowbreak{}beta)\allowbreak{} to (alpha,\allowbreak{}beta,\allowbreak{}gamma)\allowbreak{} with gamma:\allowbreak{}C \allowbreak{}-> F(Z')\allowbreak{}. Here and below the last arrow of an image triangle is xi\_\allowbreak{}X\textasciicircum{}\{-1\}F(h)\allowbreak{},\allowbreak{} where xi\_\allowbreak{}X:\allowbreak{}F(X)\allowbreak{}[1]\allowbreak{} \allowbreak{}-> F(X[1]\allowbreak{})\allowbreak{} is the exact-functor comparison. This convention makes each last square well typed.

Over the filtered category I of subsequent triangle refinements,\allowbreak{} Hom(W,\allowbreak{}-)\allowbreak{} gives a morphism between the exact five-term sequences


\par\begin{landscape}
\noindent\textbf{Exact comparison diagram.}\par
Editorial counting correction: each row has five terms; all five comparison maps are shown.\par
\begingroup\setlength{\arraycolsep}{2pt}
\[
\begin{array}{ccccccccc}
\operatorname{Hom}(W,A)&\to&\operatorname{Hom}(W,B)&\to&
\operatorname{Hom}(W,C)&\to&\operatorname{Hom}(W,A[1])&\to&\operatorname{Hom}(W,B[1])\\
\downarrow&&\downarrow&&\downarrow&&\downarrow&&\downarrow\\
\operatorname{colim}_I\operatorname{Hom}(W,F(X''))&\to&
\operatorname{colim}_I\operatorname{Hom}(W,F(Y''))&\to&
\operatorname{colim}_I\operatorname{Hom}(W,F(Z''))&\to&
\operatorname{colim}_I\operatorname{Hom}(W,F(X'')[1])&\to&
\operatorname{colim}_I\operatorname{Hom}(W,F(Y'')[1]).
\end{array}
\]
\endgroup\end{landscape}


The bottom sequence is exact:\allowbreak{} an element whose next image is zero has zero image at a later index; exactness there gives a preimage,\allowbreak{} which defines a preimage in the filtered colimit. The four outside vertical arrows are bijective by essential constancy,\allowbreak{} component cofinality,\allowbreak{} and the shift comparison. The five lemma makes the middle arrow bijective. Explicitly that arrow is

\[
 \operatorname{Hom}_{D'}(W,C)\longrightarrow
 \operatorname{colim}_{I}\operatorname{Hom}_{D'}(W,F(Z'')),
 \qquad u\longmapsto[\gamma u].                            \tag{4.1}
\]

Characterization (4)\allowbreak{} of essential constancy now applies exactly to gamma. It makes C a value of RF(Z')\allowbreak{},\allowbreak{} hence of RF(Z)\allowbreak{}. The compatible maps identify the derived triangle with A \allowbreak{}-> B \allowbreak{}-> C \allowbreak{}-> A[1]\allowbreak{}. Rotating the initial triangle,\allowbreak{} including its original minus sign,\allowbreak{} reduces either other pair of defined objects to this case. For LF reverse the triangle-indexing arrows,\allowbreak{} pass to the opposite filtered category,\allowbreak{} and apply Hom(-,\allowbreak{}W)\allowbreak{}; the same five-term argument proves the dual assertion.

DERIVED-NEW-016 is a direction clarification at original 4792–\allowbreak{}4794. The displayed map there points from the colimit to Hom(W,\allowbreak{}C)\allowbreak{},\allowbreak{} whereas the diagram and the cited characterization (4)\allowbreak{} construct (4.1)\allowbreak{}. Its inverse does exist after the five-lemma argument. The correction displays the constructed map in its actual direction; it does not claim that the theorem or the existence of the inverse was false.

\subsection{5. Direct-sum cofinality and DERIVED-UNDER-006}\label{proofs-4166-5089-direct-sum-cofinality-and-derived-under-006}

Put Z\allowbreak{}=X direct-sum Y. The functor J:\allowbreak{}X/\allowbreak{}S times Y/\allowbreak{}S \allowbreak{}-> Z/\allowbreak{}S sends (s,\allowbreak{}t)\allowbreak{} to s direct-sum t. It is defined because Q is additive and S is precisely the arrows Q inverts. The two product projections are cofinal:\allowbreak{} fix an object in the omitted nonempty filtered factor,\allowbreak{} take componentwise common successors,\allowbreak{} and equalize parallel component maps. These operations also prove that the product category is filtered.

Here is the separate proof that J is cofinal. Given u:\allowbreak{}Z \allowbreak{}-> V in S,\allowbreak{} apply MS2 to u and each original projection p\_\allowbreak{}X:\allowbreak{}Z \allowbreak{}-> X and p\_\allowbreak{}Y:\allowbreak{}Z \allowbreak{}-> Y. Obtain s:\allowbreak{}X \allowbreak{}-> X',\allowbreak{} t:\allowbreak{}Y \allowbreak{}-> Y\textquotesingle{} in S and a:\allowbreak{}V \allowbreak{}-> X',\allowbreak{} b:\allowbreak{}V \allowbreak{}-> Y\textquotesingle{} with a u\allowbreak{}=s p\_\allowbreak{}X and b u\allowbreak{}=t p\_\allowbreak{}Y. The column (a,\allowbreak{}b)\allowbreak{}:\allowbreak{}V \allowbreak{}-> X\textquotesingle{} direct-sum Y\textquotesingle{} is a morphism from u to J(s,\allowbreak{}t)\allowbreak{}. This proves nonemptiness of each comma category.

For two such comma objects,\allowbreak{} choose common successors in X/\allowbreak{}S and Y/\allowbreak{}S. Call the maps to the X successor A:\allowbreak{}X'\_\allowbreak{}1 \allowbreak{}-> X'',\allowbreak{} B:\allowbreak{}X'\_\allowbreak{}2 \allowbreak{}-> X'',\allowbreak{} and to the Y successor C:\allowbreak{}Y'\_\allowbreak{}1 \allowbreak{}-> Y'',\allowbreak{} D:\allowbreak{}Y'\_\allowbreak{}2 \allowbreak{}-> Y\textquotesingle\textquotesingle. For their maps a\_\allowbreak{}i:\allowbreak{}V \allowbreak{}-> X'\_\allowbreak{}i and b\_\allowbreak{}i:\allowbreak{}V \allowbreak{}-> Y'\_\allowbreak{}i,\allowbreak{} one has (A a\_\allowbreak{}1)\allowbreak{}u\allowbreak{}=(B a\_\allowbreak{}2)\allowbreak{}u and (C b\_\allowbreak{}1)\allowbreak{}u\allowbreak{}=(D b\_\allowbreak{}2)\allowbreak{}u. MS3 gives a denominator out of X\textquotesingle\textquotesingle{} equalizing the first pair and one out of Y\textquotesingle\textquotesingle{} equalizing the second. Postcomposing the successors gives a common comma object and actual morphisms to it from both original objects. This proves connectedness,\allowbreak{} including the equalities on V rather than only after Q.

Thus,\allowbreak{} naturally in W and with the original inclusion and projection maps,\allowbreak{} (3.1)\allowbreak{} gives

\[
 T_Z(W)=T_X(W)\oplus T_Y(W).                              \tag{5.1}
\]

To see the colimit step explicitly,\allowbreak{} any pair of representatives reaches a common product index; two pairs agree precisely when their two component equalities hold after a common refinement. This proves the bijection with the finite biproduct of the two colimits. If both summands are represented,\allowbreak{} their biproduct represents T\_\allowbreak{}Z,\allowbreak{} proving part (1)\allowbreak{} of lemma-direct-sum-defined and its specified decomposition. The zero and inclusion/\allowbreak{}projection maps agree with these component maps by the functor construction in section 3.

Now assume RF(Z)\allowbreak{}\allowbreak{}=V exists,\allowbreak{} without assuming D\textquotesingle{} Karoubian. Let

\[
 e=RF(i_Xp_X):V\longrightarrow V.
\]

This is defined since it is an endomorphism of Z,\allowbreak{} even before RF(X)\allowbreak{} is known to exist. Functoriality gives e\textasciicircum{}2\allowbreak{}=e. Under (5.1)\allowbreak{},\allowbreak{} Hom(-,\allowbreak{}e)\allowbreak{} is exactly the projection onto T\_\allowbreak{}X. If e\allowbreak{}=j r with r:\allowbreak{}V \allowbreak{}-> A,\allowbreak{} j:\allowbreak{}A \allowbreak{}-> V,\allowbreak{} and rj\allowbreak{}=1\_\allowbreak{}A,\allowbreak{} then Hom(-,\allowbreak{}A)\allowbreak{},\allowbreak{} included by j and projected by r,\allowbreak{} represents T\_\allowbreak{}X. Conversely a representing object for T\_\allowbreak{}X and the inclusion and projection in (5.1)\allowbreak{} give j and r by Yoneda and satisfy those same identities. Therefore

\[
 RF(X)\text{ exists}\quad\Longleftrightarrow\quad e\text{ splits}.
                                                               \tag{5.2}
\]

The corresponding statement for Y uses 1-e. This is the receiving application of HOMOLOGY-UNDER-010 and its individual-witness addendum. That earlier proof also gives an explicit essential-constancy witness:\allowbreak{} if c\_\allowbreak{}i:\allowbreak{}H\_\allowbreak{}i \allowbreak{}-> V and t:\allowbreak{}V \allowbreak{}-> H\_\allowbreak{}i witness the sum diagram,\allowbreak{} then p\_\allowbreak{}i t j:\allowbreak{}A \allowbreak{}-> F(X\_\allowbreak{}i)\allowbreak{} witnesses its X summand; the cocone is r c\_\allowbreak{}i i\_\allowbreak{}i. The identities are

\[
 r c_i i_i p_i t j=r e c_i t j=rj=1_A,
 \qquad F(\beta)=F(\alpha)p_i t j(r c_k i_k)
\]

for the eventual pair of index maps alpha and beta supplied by the sum witness. They follow by restricting that witness\textquotesingle s equality to the summand and using c\_\allowbreak{}k i\_\allowbreak{}k\allowbreak{}=e c\_\allowbreak{}k i\_\allowbreak{}k\allowbreak{}=j r c\_\allowbreak{}k i\_\allowbreak{}k.

In the present triangulated target,\allowbreak{} splitting e also splits 1-e. This requires an argument; it is not automatic in an arbitrary additive category. Complete j:\allowbreak{}A \allowbreak{}-> V to a triangle A \allowbreak{}-> V \allowbreak{}-> C \allowbreak{}-> A[1]\allowbreak{},\allowbreak{} writing its second map q and last map h. Since j[1]\allowbreak{}h\allowbreak{}=0 and rj\allowbreak{}=1,\allowbreak{} one has h\allowbreak{}=r[1]\allowbreak{}j[1]\allowbreak{}h\allowbreak{}=0. Exactness of Hom(C,\allowbreak{}-)\allowbreak{} gives k\_\allowbreak{}0:\allowbreak{}C \allowbreak{}-> V with qk\_\allowbreak{}0\allowbreak{}=1. Set k\allowbreak{}=k\_\allowbreak{}0-jr k\_\allowbreak{}0,\allowbreak{} so qk\allowbreak{}=1 and rk\allowbreak{}=0. The map d\allowbreak{}=1\_\allowbreak{}V-jr-kq has qd\allowbreak{}=0,\allowbreak{} hence factors as jv by exactness of Hom(V,\allowbreak{}-)\allowbreak{}. But rd\allowbreak{}=0,\allowbreak{} so v\allowbreak{}=rjv\allowbreak{}=rd\allowbreak{}=0. Thus

\[
 kq=1_V-e,\qquad qk=1_C,\qquad
 (j,k)(r,q)=1_V,\qquad (r,q)(j,k)=1_{A\oplus C}.          \tag{5.3}
\]

Consequently either condition in (5.2)\allowbreak{},\allowbreak{} or its Y analogue,\allowbreak{} guarantees both values and their biproduct decomposition. The original Karoubian hypothesis ensures every needed e splits,\allowbreak{} but only this specified idempotent is necessary for this pair. This precise criterion is DERIVED-UNDER-006; it is kept editorial. Applying the construction in D\textquotesingle\^{}op gives the same criterion for LF,\allowbreak{} with e\allowbreak{}=LF(i\_\allowbreak{}Xp\_\allowbreak{}X)\allowbreak{} on LF(Z)\allowbreak{} and the representing functors Hom(A,\allowbreak{}-)\allowbreak{}.

DERIVED-RECON-024 reconciles P02-E0030 and French DERIVED-031/\allowbreak{}032 with both existing corrections MC-STK-ERR-0842/\allowbreak{}0843:\allowbreak{} the two restrictions are of F,\allowbreak{} the only functor supplied in the situation. DERIVED-RECON-025 accepts P02-E0031\textquotesingle s missing agreement correction at original 4891,\allowbreak{} changing ``The proof ... are dual'' to ``is dual''.

\subsection{\texorpdfstring{6. The full domain,\allowbreak{} its localization,\allowbreak{} and the zero-object overclaim}{6. The full  its  and the zero-object overclaim}}\label{proofs-4166-5089-the-full-domain-its-localization-and-the-zero-object-overclaim}

Let E be the full subcategory where RF exists. The zero object is in E:\allowbreak{} the object id\_\allowbreak{}0 of 0/\allowbreak{}S is terminal,\allowbreak{} and its image is F(0)\allowbreak{}\allowbreak{}=0. The essential-constancy witness is 0 \allowbreak{}-> F(0)\allowbreak{}; every index has its unique map to this terminal object,\allowbreak{} giving the required eventual equality. The dual initial object of S/\allowbreak{}0 proves the LF assertion as well. Inverting isomorphisms,\allowbreak{} shifts,\allowbreak{} and section 4 prove E is strictly full and triangulated. Sections 3–\allowbreak{}4 give the exact functor RF on it.

If an arrow of S has either endpoint in E,\allowbreak{} section 3 puts the other endpoint in E. This property proves all the multiplicative-system axioms for S\_\allowbreak{}E inside E:\allowbreak{} an MS2 construction starting with objects in E has its new endpoint in E because the new comparison arrow is in S; the same applies to an MS3 equalizing denominator. Identities,\allowbreak{} composition and saturation restrict directly. Shifts preserve E; in MS6 the third vertex belongs to E by its triangle,\allowbreak{} and the third denominator has both endpoints there. Hence S\_\allowbreak{}E is compatible.

For X,\allowbreak{}Y in E a fraction X \allowbreak{}-> T \textless- Y in the larger localization already has T in E,\allowbreak{} since the arrow out of Y is a denominator. The same is true of every later denominator used to compare two such fractions. Thus both the morphisms and their equality relation are exactly those in S\_\allowbreak{}E\textasciicircum{}\{-1\}E. This proves full faithfulness,\allowbreak{} not just faithfulness on original morphisms. Localized distinguished triangles are images of triangles up to isomorphism,\allowbreak{} so this embedding is exact. RF inverts S\_\allowbreak{}E and therefore factors with its specified shift comparison as the stated exact functor on S\_\allowbreak{}E\textasciicircum{}\{-1\}E.

If D\textquotesingle{} is Karoubian,\allowbreak{} every e of section 5 splits,\allowbreak{} so E is closed under existing summands of its objects. More precisely,\allowbreak{} the same conclusion holds whenever just the idempotents RF(i\_\allowbreak{}Xp\_\allowbreak{}X)\allowbreak{} arising from decompositions of objects in E split. Conversely closure under these summands forces these idempotents to split by (5.2)\allowbreak{}. This propagates DERIVED-UNDER-006 to part (8)\allowbreak{} of proposition-derived-functor without replacing that proposition\textquotesingle s original sufficient hypothesis.

DERIVED-NEW-018 concerns original 4982–\allowbreak{}4983. The statement that the canonical map might ``never'' be an isomorphism is literally false:\allowbreak{} F(0)\allowbreak{} \allowbreak{}-> RF(0)\allowbreak{} is always the isomorphism 0 \allowbreak{}-> 0 just proved. Its intended warning is retained by saying the map need not be an isomorphism for a given X. For an explicit example even when RF is everywhere defined,\allowbreak{} take D\allowbreak{}=D'\allowbreak{}=K\textasciicircum{}b(Ab)\allowbreak{},\allowbreak{} F the identity,\allowbreak{} and S all morphisms. MS1–\allowbreak{}MS3 hold (postcompose or precompose with a zero object where necessary)\allowbreak{},\allowbreak{} shift compatibility is automatic,\allowbreak{} and MS6 is TR3 since every completed third map lies in S. The localization is the zero category:\allowbreak{} inverting the zero endomorphism makes each identity zero. Thus S is saturated. For each X,\allowbreak{} X/\allowbreak{}S has terminal object X \allowbreak{}-> 0. Its diagram is essentially constant with value 0,\allowbreak{} using the zero witness and the maps to that terminal index. Hence RF is the zero functor. At X\allowbreak{}=Z[0]\allowbreak{},\allowbreak{} F(X)\allowbreak{} \allowbreak{}-> RF(X)\allowbreak{} is not an isomorphism:\allowbreak{} a homotopy from 1\_\allowbreak{}\{Z[0]\allowbreak{}\} to zero would have to be zero in every degree and cannot equal its nonzero degree-zero identity. At X\allowbreak{}=0 it is an isomorphism. This verifies both the correction and the retained warning with the source\textquotesingle s exact assumptions.

DERIVED-NEW-019 groups three unreported grammar corrections:\allowbreak{} ``fully faithfulness'' becomes ``full faithfulness'' at 4933,\allowbreak{} and the two quoted phrases at 4965/\allowbreak{}4967 gain ``is'' before ``everywhere defined''. These change none of the domains,\allowbreak{} quantifiers or hypotheses.

\subsection{7. Computing objects and the receiving variance correction}\label{proofs-4166-5089-computing-objects-and-the-receiving-variance-correction}

The canonical cocone map eta\_\allowbreak{}X:\allowbreak{}F(X)\allowbreak{} \allowbreak{}-> RF(X)\allowbreak{} is natural in X by section 3 with the identity denominator. It commutes with the shift comparisons because those comparisons are induced at every index. Therefore eta\_\allowbreak{}X is invertible exactly when eta\_\allowbreak{}\{X[n]\allowbreak{}\} is. In a distinguished triangle of defined objects,\allowbreak{} eta gives a morphism of triangles; if two components are invertible,\allowbreak{} the Hom long exact sequences and the five lemma make the third invertible. Rotation gives any pair of components. The LF argument uses the natural cone map LF(X)\allowbreak{} \allowbreak{}-> F(X)\allowbreak{} and the same triangle and shift comparisons.

Suppose Z\allowbreak{}=X direct-sum Y computes RF. Through eta\_\allowbreak{}Z,\allowbreak{} the value RF(Z)\allowbreak{} is F(X)\allowbreak{} direct-sum F(Y)\allowbreak{},\allowbreak{} and the idempotent e of section 5 is exactly F(i\_\allowbreak{}Xp\_\allowbreak{}X)\allowbreak{}. Its splitting is explicitly F(i\_\allowbreak{}X)\allowbreak{},\allowbreak{} F(p\_\allowbreak{}X)\allowbreak{}. Hence RF(X)\allowbreak{} and RF(Y)\allowbreak{} exist. Moreover (5.1)\allowbreak{} identifies

\[
 \operatorname{Hom}(W,F(X\oplus Y))\longrightarrow T_Z(W)
\]

with the direct sum of the two canonical maps Hom(W,\allowbreak{}F(X)\allowbreak{})\allowbreak{} \allowbreak{}-> T\_\allowbreak{}X(W)\allowbreak{} and Hom(W,\allowbreak{}F(Y)\allowbreak{})\allowbreak{} \allowbreak{}-> T\_\allowbreak{}Y(W)\allowbreak{}. A block diagonal map of two abelian groups is bijective only if both component maps are bijective:\allowbreak{} injectivity follows by putting zero in the other component,\allowbreak{} and surjectivity by lifting an element with other component zero and then projecting. Thus eta\_\allowbreak{}X and eta\_\allowbreak{}Y are isomorphisms by the representability criterion,\allowbreak{} with precisely these canonical maps. This proves lemma-direct-sum-computes without a Karoubian assumption and completes the propagation of HOMOLOGY-UNDER-010 through its cofinality step and this receiving lemma.

DERIVED-NEW-017:\allowbreak{} original 5056,\allowbreak{} 5058,\allowbreak{} 5063,\allowbreak{} 5065 instead have Hom(F(-)\allowbreak{},\allowbreak{}W)\allowbreak{} under a colimit indexed by the forward filtered category. For a transition v:\allowbreak{}X' \allowbreak{}-> X'',\allowbreak{} this Hom expression supplies

\[
 \operatorname{Hom}(F(X''),W)\longrightarrow
 \operatorname{Hom}(F(X'),W),
\]

which is in the wrong direction for the displayed diagram. The cited ind-object characterization requires Hom(W,\allowbreak{}F(-)\allowbreak{})\allowbreak{}. Reversing the arguments in all four occurrences gives the two exact comparison maps just proved and restores the forward transition by postcomposition with F(v)\allowbreak{}. It leaves W,\allowbreak{} every denominator,\allowbreak{} both index categories,\allowbreak{} and the right-derived conclusion unchanged. Replacing the colimit by a limit would not be this proof:\allowbreak{} it would impose a different variance and would not give characterization (4)\allowbreak{}.

For LF one genuinely uses Hom(F(-)\allowbreak{},\allowbreak{}W)\allowbreak{},\allowbreak{} but now the index is the opposite of S/\allowbreak{}X,\allowbreak{} as in section 3; the reversed denominators provide the needed forward Hom transitions. This proves the dual result without conflating its indexing category with X/\allowbreak{}S.

Finally lemma-existence-computes follows from the already proved denominator invariance:\allowbreak{} any specified denominator connecting X to a computing object puts X in the domain. This proves the right and left statements at original 5073–\allowbreak{}5089. The next criterion,\allowbreak{} lemma-find-existence-computes,\allowbreak{} is reserved for the next source-ordered review,\allowbreak{} despite its opening having been read through original 5120.

All seven received physical reports in this interval have been read and accounted for. The eight later French records DERIVED-022 through DERIVED-029 have original loci from 5260 through 6197 and remain in their later source positions; their numbering is not source order. The wider mathematical synthesis and literature novelty assessment remain follow-on work after the correction and received-theorem queue.
\clearpage
\section{\texorpdfstring{Computing classes,\allowbreak{} finite acyclic prefixes,\allowbreak{} and a repaired existence proof}{Computing  finite acyclic  and a repaired existence proof}}\label{proofs-5091-5900-computing-classes-finite-acyclic-prefixes-and-a-repaired-existence-proof}

Authority:\allowbreak{} \texttt{derived.\allowbreak{}tex} at \texttt{a04446e5\allowbreak{}7ec1fbc2\allowbreak{}52a871af\allowbreak{}cec7752f\allowbreak{}b2807b14},\allowbreak{} original lines 5091–\allowbreak{}5900. The private original/\allowbreak{}current snapshots fix the comparison. The receiving proof in \texttt{perfect.\allowbreak{}tex} is compared at the same authority and the public R60 source. Corrections are distinguished below from the editorial refinement DERIVED-UNDER-007. No novelty is claimed.

\subsection{1. Cofinal computing classes and the composition comparison}\label{proofs-5091-5900-cofinal-computing-classes-and-the-composition-comparison}

For the source\textquotesingle s class I,\allowbreak{} let C\_\allowbreak{}X be the full subcategory of X/\allowbreak{}S whose targets lie in I. Every object s:\allowbreak{}X \allowbreak{}-> Y has an arrow to an object of C\_\allowbreak{}X:\allowbreak{} choose the asserted denominator t:\allowbreak{}Y \allowbreak{}-> I\_\allowbreak{}Y and compose ts. Two objects of C\_\allowbreak{}X have a common successor in X/\allowbreak{}S; refine its target once more into I. Parallel arrows can first be equalized in X/\allowbreak{}S and then refined into I. Thus C\_\allowbreak{}X is filtered. The same procedure applied to two objects under a fixed s proves that the corresponding comma category is nonempty and connected. This proves cofinality,\allowbreak{} including the connectedness condition.

Every transition u between its targets is in S,\allowbreak{} since Q(u)\allowbreak{} is the quotient of two invertible structure maps and S is saturated. Therefore F(u)\allowbreak{} is an isomorphism by the source\textquotesingle s second hypothesis. Choose one index i. For any j choose a common successor k with arrows a:\allowbreak{}i \allowbreak{}-> k,\allowbreak{} b:\allowbreak{}j \allowbreak{}-> k,\allowbreak{} and set c\_\allowbreak{}j\allowbreak{}=F(a)\allowbreak{}\textasciicircum{}\{-1\}F(b)\allowbreak{}:\allowbreak{}F(X\_\allowbreak{}j)\allowbreak{} \allowbreak{}-> F(X\_\allowbreak{}i)\allowbreak{}. Further common refinements and equalization show this is independent of both choices. These maps form a cocone; c\_\allowbreak{}i\allowbreak{}=1. Their defining equalities give the essential-constancy witness at index i. Thus the source\textquotesingle s word ``constant'' means canonically isomorphic to a constant diagram,\allowbreak{} rather than literal equality of its objects and arrows. When X belongs to I,\allowbreak{} choose i\allowbreak{}=id\_\allowbreak{}X. The cocone identifies the canonical map F(X)\allowbreak{} \allowbreak{}-> RF(X)\allowbreak{} with the identity. This proves the computing assertion as well as existence. In the opposite category,\allowbreak{} the identical common-predecessor and equalization construction proves the left-derived criterion for P.

For composition retain the original Q\_\allowbreak{}A,\allowbreak{} Q\_\allowbreak{}B and F'\allowbreak{}=Q\_\allowbreak{}B F. Let c\_\allowbreak{}\{s\}:\allowbreak{}F'(X\_\allowbreak{}s)\allowbreak{} \allowbreak{}-> RF'(Q\_\allowbreak{}A X)\allowbreak{} be the right-derived cocone,\allowbreak{} and write eta\textasciicircum{}G\_\allowbreak{}Y:\allowbreak{}G(Y)\allowbreak{} \allowbreak{}-> RG(Q\_\allowbreak{}B Y)\allowbreak{}. The maps

\[
 G(F(X_s))\xrightarrow{\eta^G_{F(X_s)}}RG(F'(X_s))
        \xrightarrow{RG(c_s)}RG(RF'(Q_A X))                 \tag{1.1}
\]

form a cocone,\allowbreak{} by naturality of eta\^{}G. The universal property of R(GF)\allowbreak{}(Q\_\allowbreak{}A X)\allowbreak{} therefore gives the source\textquotesingle s comparison t\_\allowbreak{}X. The only interchange of RG and a colimit is justified by essential constancy:\allowbreak{} applying any functor to a witness,\allowbreak{} its section identity and its eventual factorization identity gives another such witness. Hence the image diagram has value RG(RF'(Q\_\allowbreak{}A X)\allowbreak{})\allowbreak{}; no preservation of arbitrary filtered colimits is being assumed.

For f:\allowbreak{}X \allowbreak{}-> Y choose its denominator squares as in the preceding review. Naturality of eta\^{}G and the cocone identities for RF\textquotesingle{} make the composites of (1.1)\allowbreak{} on these squares equal. Equality on every colimit structure map proves naturality of t. Denominators are sent to isomorphisms by all involved derived functors,\allowbreak{} so the equality also holds for every localized fraction. The iterated original shift comparisons commute on each map (1.1)\allowbreak{},\allowbreak{} and thus on t. This supplies the omitted transformation check without identifying Q\_\allowbreak{}A with F or suppressing Q\_\allowbreak{}B.

For the left version,\allowbreak{} use the projections LF'(Q\_\allowbreak{}A X)\allowbreak{} \allowbreak{}-> F'(X\_\allowbreak{}s)\allowbreak{},\allowbreak{} apply LG,\allowbreak{} and follow by the canonical maps LG(Q\_\allowbreak{}B F(X\_\allowbreak{}s)\allowbreak{})\allowbreak{} \allowbreak{}-> G(F(X\_\allowbreak{}s)\allowbreak{})\allowbreak{}. They form a cone,\allowbreak{} inducing LG(LF'(Q\_\allowbreak{}A X)\allowbreak{})\allowbreak{} \allowbreak{}-> L(GF)\allowbreak{}(Q\_\allowbreak{}A X)\allowbreak{}. Applying a functor to the pro-object witness justifies the analogous limit interchange. The reversed denominator squares prove naturality and compatibility with the same shift comparisons.

\subsection{2. Bounded and unbounded definitions at the same object}\label{proofs-5091-5900-bounded-and-unbounded-definitions-at-the-same-object}

For X in K\textasciicircum{}\allowbreak{}+(A)\allowbreak{},\allowbreak{} the category X/\allowbreak{}Qis\textasciicircum{}\allowbreak{}+(A)\allowbreak{} is a full subcategory of X/\allowbreak{}Qis(A)\allowbreak{}. Every index X \allowbreak{}-> X\textquotesingle{} can be followed by the lower truncation quasi-isomorphism X\textquotesingle{} \allowbreak{}-> X'',\allowbreak{} for X\textquotesingle\textquotesingle{} bounded below. Common successors and equalizations in the larger filtered category followed by this construction prove cofinality,\allowbreak{} including connected comma categories. Thus the two indexing diagrams agree after cofinal restriction.

There is also a target-category issue to verify. The bounded diagram takes values in the fully faithfully embedded D\textasciicircum{}\allowbreak{}+(B)\allowbreak{}. If it is essentially constant there,\allowbreak{} its witness remains a witness in D(B)\allowbreak{}. Conversely,\allowbreak{} if its value V exists as an essentially constant object in D(B)\allowbreak{},\allowbreak{} the witness makes V a retract of F(X'')\allowbreak{} for one bounded-below complex X\textquotesingle\textquotesingle. Cohomology in every degree below that complex\textquotesingle s lower bound is then a retract of zero,\allowbreak{} hence zero. Consequently V belongs to D\textasciicircum{}\allowbreak{}+(B)\allowbreak{}. Fullness puts its cocone and witness maps in that category as well. This proves both implications of lemma-irrelevant,\allowbreak{} not just cofinality of the source categories. The identity-denominator maps are preserved,\allowbreak{} so the assertion about computing objects follows. Reversing arrows,\allowbreak{} using upper truncations and D\textasciicircum{}-(B)\allowbreak{},\allowbreak{} proves both left-derived assertions.

DERIVED-RECON-026 reconciles French DERIVED-022 with MC-STK-ERR-0833\textquotesingle s existing indefinite article at original 5260. The underlying replacement argument is the one just proved.

\subsection{3. Resolutions by a prescribed class and the acyclic-class criterion}\label{proofs-5091-5900-resolutions-by-a-prescribed-class-and-the-acyclic-class-criterion}

For the source\textquotesingle s left-resolution construction,\allowbreak{} retain its induction data in degrees j>\allowbreak{}=n. Write Z\_\allowbreak{}P\textasciicircum{}n\allowbreak{}=ker(d\_\allowbreak{}P\textasciicircum{}n)\allowbreak{},\allowbreak{} Z\_\allowbreak{}K\textasciicircum{}n\allowbreak{}=ker(d\_\allowbreak{}K\textasciicircum{}n)\allowbreak{}. The induction hypothesis includes an epimorphism Z\_\allowbreak{}P\textasciicircum{}n \allowbreak{}-> Z\_\allowbreak{}K\textasciicircum{}n. The actual object used at the next step is

\[
 E=K^{n-1}\mathbin{\times}_{K^n}Z_P^n
  =K^{n-1}\mathbin{\times}_{Z_K^n}Z_P^n,                 \tag{3.1}
\]

because d\_\allowbreak{}K\textasciicircum{}\{n-1\} lands in Z\_\allowbreak{}K\textasciicircum{}n. Its projection to K\^{}\{n-1\} is epic by stability of epimorphisms under pullback in an abelian category. Choose the given class epimorphism P\^{}\{n-1\} \allowbreak{}-> E. Its two projections give alpha\^{}\{n-1\} and d\_\allowbreak{}P\textasciicircum{}\{n-1\}. They make the square commute,\allowbreak{} alpha\^{}\{n-1\} epic,\allowbreak{} and d\_\allowbreak{}P\textasciicircum{}n d\_\allowbreak{}P\textasciicircum{}\{n-1\}\allowbreak{}=0.

The image of E \allowbreak{}-> Z\_\allowbreak{}P\textasciicircum{}n is the inverse image of im(d\_\allowbreak{}K\textasciicircum{}\{n-1\})\allowbreak{} under Z\_\allowbreak{}P\textasciicircum{}n \allowbreak{}-> Z\_\allowbreak{}K\textasciicircum{}n. Indeed (3.1)\allowbreak{} is the pullback of d\_\allowbreak{}K\textasciicircum{}\{n-1\}; factor that map as an epimorphism followed by its image inclusion and pull back. The image is therefore exactly the kernel of the epimorphism Z\_\allowbreak{}P\textasciicircum{}n \allowbreak{}-> H\textasciicircum{}n(K)\allowbreak{}. Since P\^{}\{n-1\} \allowbreak{}-> E is epic,\allowbreak{} quotienting by im(d\_\allowbreak{}P\textasciicircum{}\{n-1\})\allowbreak{} makes H\textasciicircum{}n(P)\allowbreak{} \allowbreak{}-> H\textasciicircum{}n(K)\allowbreak{} an isomorphism. The kernel of E \allowbreak{}-> Z\_\allowbreak{}P\textasciicircum{}n is Z\_\allowbreak{}K\textasciicircum{}\{n-1\}. Pulling the epimorphism P\^{}\{n-1\} \allowbreak{}-> E back along this kernel shows that Z\_\allowbreak{}P\textasciicircum{}\{n-1\} \allowbreak{}-> Z\_\allowbreak{}K\textasciicircum{}\{n-1\} is epic. Higher cohomology is unchanged. These are all clauses of IH\_\allowbreak{}\{n-1\}. The initial zero terms at n\allowbreak{}=a\allowbreak{}+1 satisfy the induction hypotheses because K vanishes above a. Successive descending choices yield the entire requested complex; each degree and differential is fixed at a finite stage. Applying this construction to tau\_\allowbreak{}\{<\allowbreak{}=a\}K proves part (2)\allowbreak{},\allowbreak{} with the stated quasi-isomorphism into the original K.

The right-resolution construction is its exact dual. More explicitly,\allowbreak{} replace the pullback by the pushout

\[
 K^{n+1}\mathbin{\amalg}_{K^n}\operatorname{coker}(d_I^{n-1})
 =K^{n+1}\mathbin{\amalg}_{\operatorname{coker}(d_K^{n-1})}
                         \operatorname{coker}(d_I^{n-1}),       \tag{3.2}
\]

and embed (3.2)\allowbreak{} into a chosen class object I\textasciicircum{}\{n\allowbreak{}+1\}. The inductive monomorphism of cokernels makes K\textasciicircum{}\{n\allowbreak{}+1\} \allowbreak{}-> (3.2)\allowbreak{} monic by stability of monomorphisms under pushout. The maps from K\textasciicircum{}\{n\allowbreak{}+1\} and I\^{}n give the next component and differential. The dual kernel/\allowbreak{}image argument above proves the cohomology and cokernel induction conditions. Start with zero terms below a and proceed upward. For cohomology bounded below first use K \allowbreak{}-> tau\_\allowbreak{}\{>\allowbreak{}=a\}K. These constructions require neither projectivity/\allowbreak{}injectivity nor closure of the class under sums.

Now impose the hypotheses of lemma-subcategory-right-acyclics. The class I is nonempty by its embedding property. The sequence 0 \allowbreak{}-> P \allowbreak{}-> P \allowbreak{}-> 0 \allowbreak{}-> 0 for P in I shows 0 is in I. The same cokernel condition with P\allowbreak{}=0 shows closure under isomorphism. Class resolutions form a cofinal full subcategory of denominators:\allowbreak{} resolve the target of each denominator,\allowbreak{} then refine common successors and equalizations as in section 1.

Let A in I and A[0]\allowbreak{} \allowbreak{}-> I\^{}bullet be a quasi-isomorphism with I\^{}n in I and I bounded below. Choose n0\textless0 with I\textasciicircum{}n\allowbreak{}=0 below n0. Put B\textasciicircum{}n\allowbreak{}=im(d\textasciicircum{}\{n-1\})\allowbreak{} and Z\textasciicircum{}n\allowbreak{}=ker(d\textasciicircum{}n)\allowbreak{}. Below degree zero,\allowbreak{} B\textasciicircum{}n\allowbreak{}=Z\textasciicircum{}n. Starting with B\textasciicircum{}\{n0\}\allowbreak{}=0,\allowbreak{} the sequences 0 \allowbreak{}-> B\^{}n \allowbreak{}-> I\^{}n \allowbreak{}-> B\textasciicircum{}\{n\allowbreak{}+1\} \allowbreak{}-> 0 for n0<\allowbreak{}=n<0 and the class hypothesis show that all these B's,\allowbreak{} including B\textasciicircum{}0,\allowbreak{} belong to I,\allowbreak{} and that their F-images are exact sequences. Next 0 \allowbreak{}-> B\^{}0 \allowbreak{}-> I\^{}0 \allowbreak{}-> I\textasciicircum{}0/\allowbreak{}B\textasciicircum{}0 \allowbreak{}-> 0 puts I\textasciicircum{}0/\allowbreak{}B\textasciicircum{}0 in I. The quasi-isomorphism identifies A with Z\textasciicircum{}0/\allowbreak{}B\textasciicircum{}0,\allowbreak{} giving 0 \allowbreak{}-> A \allowbreak{}-> I\textasciicircum{}0/\allowbreak{}B\textasciicircum{}0 \allowbreak{}-> B\^{}1 \allowbreak{}-> 0. Thus B\^{}1 is in I and its F-sequence is exact. Continue with 0 \allowbreak{}-> B\^{}n \allowbreak{}-> I\^{}n \allowbreak{}-> B\textasciicircum{}\{n\allowbreak{}+1\} \allowbreak{}-> 0 for n>\allowbreak{}=1.

Splicing these exact F-sequences shows that H\textasciicircum{}n(F(I)\allowbreak{})\allowbreak{}\allowbreak{}=0 for n!\allowbreak{}=0 and that the actual map F(A)\allowbreak{} \allowbreak{}-> H\textasciicircum{}0(F(I)\allowbreak{})\allowbreak{} is an isomorphism. At degree zero this follows by composing F(I\textasciicircum{}0)\allowbreak{} \allowbreak{}-> F(I\textasciicircum{}0/\allowbreak{}B\textasciicircum{}0)\allowbreak{} with the kernel identification F(A)\allowbreak{}\allowbreak{}=ker(F(I\textasciicircum{}0/\allowbreak{}B\textasciicircum{}0)\allowbreak{} \allowbreak{}-> F(B\textasciicircum{}1)\allowbreak{})\allowbreak{}; it is the map induced by the original A \allowbreak{}-> I\^{}0. Hence F(A)\allowbreak{}[0]\allowbreak{} \allowbreak{}-> F(I)\allowbreak{} is a quasi-isomorphism. The resulting constant cocone proves A computes RF. Reversing all sequences and using the quotient class P proves lemma-subcategory-left-acyclics with its stated hypotheses.

DERIVED-RECON-027 reconciles P02-E0034 and French DERIVED-023 with the existing MC-STK-ERR-0834 singular ``element''. DERIVED-RECON-028 reconciles P02-E0033 and French DERIVED-024 with MC-STK-ERR-0835:\allowbreak{} the proved property is quasi-isomorphism. A termwise isomorphism is not asserted. For example take F the identity on Ab,\allowbreak{} class I all objects,\allowbreak{} and I\textasciicircum{}bullet\allowbreak{}=Z[0]\allowbreak{} direct-sum (Z \allowbreak{}-> Z)\allowbreak{} in degrees 1 and 2 with identity differential. The natural Z[0]\allowbreak{} \allowbreak{}-> I is a quasi-isomorphism and satisfies every hypothesis,\allowbreak{} but is not an isomorphism of complexes.

\subsection{\texorpdfstring{4. Lower bounds,\allowbreak{} higher derived functors,\allowbreak{} and their connecting maps}{4. Lower  higher derived  and their connecting maps}}\label{proofs-5091-5900-lower-bounds-higher-derived-functors-and-their-connecting-maps}

For arbitrary K with H\textasciicircum{}i(K)\allowbreak{}\allowbreak{}=0 when i<a,\allowbreak{} every denominator K \allowbreak{}-> L can be followed by L \allowbreak{}-> tau\_\allowbreak{}\{>\allowbreak{}=a\}L,\allowbreak{} a quasi-isomorphism. The full subcategory of targets zero below a is cofinal by the same successor and equalization argument. If RF(K)\allowbreak{} exists,\allowbreak{} its essential-constancy witness factors its identity through some F(L)\allowbreak{} with L\textasciicircum{}i\allowbreak{}=0 below a. Applying H\^{}i gives a factorization through zero for i<a,\allowbreak{} so those cohomology objects of RF(K)\allowbreak{} are zero.

DERIVED-NEW-020:\allowbreak{} the category at original 5515 is written K/\allowbreak{}Qis\textasciicircum{}\allowbreak{}+(A)\allowbreak{},\allowbreak{} although this lemma allows an arbitrary complex K whose cohomology,\allowbreak{} rather than its terms,\allowbreak{} is bounded below. Such a K need not be an object of K\textasciicircum{}\allowbreak{}+(A)\allowbreak{}. The proof just given works directly in K/\allowbreak{}Qis(A)\allowbreak{}; replacing Qis\textasciicircum{}\allowbreak{}+ by Qis at this single occurrence makes the original proof well typed without replacing K. If K already belongs to K\textasciicircum{}\allowbreak{}+(A)\allowbreak{},\allowbreak{} section 2 gives the compatible bounded version.

For part (2)\allowbreak{} apply RF to the original truncation triangle

\[
 \tau_{\le a}K\longrightarrow K\longrightarrow\tau_{\ge a+1}K
                      \longrightarrow(\tau_{\le a}K)[1].       \tag{4.1}
\]

Existence on the first two vertices gives existence on the third by the previously proved two-out-of-three result. The third image has H\textasciicircum{}i\allowbreak{}=0 for i<a\allowbreak{}+1. The exact sequence with terms H\^{}\{i-1\} of that image,\allowbreak{} H\^{}i of the first two images,\allowbreak{} and H\^{}i of the third proves the stated canonical isomorphism for i<\allowbreak{}=a. No upper-bound conclusion is inferred.

When RF exists on D\textasciicircum{}\allowbreak{}+(A)\allowbreak{},\allowbreak{} define R\^{}iF on A by first placing A in degree zero. The lower-bound result gives R\textasciicircum{}iF\allowbreak{}=0 for i\textless0. For 0 \allowbreak{}-> A \allowbreak{}-> B \allowbreak{}-> C \allowbreak{}-> 0 the canonical connecting triangle and RF give the exact sequence

\[
 0\to R^0F(A)\to R^0F(B)\to R^0F(C)\to R^1F(A)\to\cdots.
\]

Thus R\^{}0F is left exact. If F \allowbreak{}-> R\^{}0F is an isomorphism,\allowbreak{} F is left exact. Conversely,\allowbreak{} for a left exact F,\allowbreak{} restrict denominators A[0]\allowbreak{} \allowbreak{}-> K to complexes zero below zero. Then 0 \allowbreak{}-> A \allowbreak{}-> K\^{}0 \allowbreak{}-> K\^{}1 is exact,\allowbreak{} so F(A)\allowbreak{}\allowbreak{}=ker(F(K\textasciicircum{}0)\allowbreak{} \allowbreak{}-> F(K\textasciicircum{}1)\allowbreak{})\allowbreak{}\allowbreak{}=H\textasciicircum{}0(F(K)\allowbreak{})\allowbreak{}. Applying H\^{}0 to the essential-constancy witness shows this cocone has value R\textasciicircum{}0F(A)\allowbreak{}. Its maps are exactly the canonical comparison,\allowbreak{} proving F\allowbreak{}=R\textasciicircum{}0F naturally. DERIVED-RECON-029 reconciles P02-E0032 and French DERIVED-025 with MC-STK-ERR-0836\textquotesingle s missing outer parenthesis at 5580.

A computes RF exactly when F(A)\allowbreak{}[0]\allowbreak{} \allowbreak{}-> RF(A[0]\allowbreak{})\allowbreak{} induces isomorphisms in all degrees:\allowbreak{} the cone of the comparison is zero exactly when all its cohomology is zero. The negative degrees already vanish. This proves lemma-F-acyclic,\allowbreak{} including the R\^{}0 condition for a general additive F and its removal only when F is left exact.

For the three assertions of lemma-F-acyclic-ses,\allowbreak{} keep the source\textquotesingle s left exactness assumption. If A and C are acyclic,\allowbreak{} the adjacent terms R\textasciicircum{}iF(A)\allowbreak{},\allowbreak{}R\textasciicircum{}iF(C)\allowbreak{} vanish for i\textgreater0 and force R\textasciicircum{}iF(B)\allowbreak{}\allowbreak{}=0. If A and B are acyclic,\allowbreak{} the segment R\textasciicircum{}iF(B)\allowbreak{} \allowbreak{}-> R\textasciicircum{}iF(C)\allowbreak{} \allowbreak{}-> R\textasciicircum{}\{i\allowbreak{}+1\}F(A)\allowbreak{} does the same for C. If B and C are acyclic,\allowbreak{} higher degrees force R\textasciicircum{}iF(A)\allowbreak{}\allowbreak{}=0 for i>\allowbreak{}=2; the additional surjectivity F(B)\allowbreak{} \allowbreak{}-> F(C)\allowbreak{} forces R\textasciicircum{}1F(A)\allowbreak{}\allowbreak{}=0 as well. In all cases R\textasciicircum{}1F(A)\allowbreak{}\allowbreak{}=0 makes 0 \allowbreak{}-> F(A)\allowbreak{} \allowbreak{}-> F(B)\allowbreak{} \allowbreak{}-> F(C)\allowbreak{} \allowbreak{}-> 0 exact. These arguments do not drop left exactness,\allowbreak{} which will matter in section 6.

For completeness the connecting maps of the higher delta-functor are H\^{}i of RF(c)\allowbreak{},\allowbreak{} followed by the inverse shift comparison and the identification H\textasciicircum{}i(RF(A)\allowbreak{}[1]\allowbreak{})\allowbreak{}\allowbreak{}=H\textasciicircum{}\{i\allowbreak{}+1\}(RF(A)\allowbreak{})\allowbreak{},\allowbreak{} where c:\allowbreak{}C[0]\allowbreak{} \allowbreak{}-> A[1]\allowbreak{} is the original canonical connecting morphism. Naturality of c and of the exact-functor comparison gives naturality of these maps. The triangle long exact sequence gives all delta-functor axioms. If every A embeds in an F-acyclic I,\allowbreak{} the induced R\textasciicircum{}iF(A)\allowbreak{} \allowbreak{}-> R\textasciicircum{}iF(I)\allowbreak{} is zero for i\textgreater0. The source\textquotesingle s effacement theorem therefore applies. Its full cokernel construction,\allowbreak{} independence of the embedding,\allowbreak{} naturality,\allowbreak{} connecting-map compatibility,\allowbreak{} and uniqueness are already proved in \texttt{.\allowbreak{}.\allowbreak{}/\allowbreak{}HOMOLOGY\_\allowbreak{}INTAKE\_\allowbreak{}20260930/\allowbreak{}PROOFS\_\allowbreak{}1901\_\allowbreak{}2765.\allowbreak{}md},\allowbreak{} lines 310–\allowbreak{}367. Applied with these R\^{}iF and their just-specified connecting maps,\allowbreak{} it proves universality. No new assumption that F itself is left exact is needed for this universality of the R\^{}iF system.

\subsection{5. Leray\textquotesingle s lemma and DERIVED-UNDER-007}\label{proofs-5091-5900-lerays-lemma-and-derived-under-007}

A bounded complex of right F-acyclic objects computes RF. For a one-term complex this is the definition and shift invariance. For support [a,\allowbreak{}b]\allowbreak{},\allowbreak{} the source\textquotesingle s termwise split sequence

\[
0\to\sigma_{\ge a+1}A^\bullet\to A^\bullet
                      \to\sigma_{\le a}A^\bullet\to0          \tag{5.1}
\]

gives a distinguished triangle already in the homotopy category. F preserves its termwise splittings,\allowbreak{} so this is a triangle to which the earlier computing-object two-out-of-three lemma applies. Induction proves both existence of RF and its canonical comparison isomorphism.

For a bounded-below A\^{}bullet with RF(A)\allowbreak{} defined,\allowbreak{} fix q and put P\allowbreak{}=sigma\_\allowbreak{}\{<\allowbreak{}=q\}A,\allowbreak{} T\allowbreak{}=sigma\_\allowbreak{}\{>\allowbreak{}=q\allowbreak{}+1\}A. There is a termwise split sequence 0 \allowbreak{}-> T \allowbreak{}-> A \allowbreak{}-> P \allowbreak{}-> 0 with its actual inclusion,\allowbreak{} projection,\allowbreak{} and connecting map. Suppose only that A\^{}n is right F-acyclic for n<\allowbreak{}=q. Then P is bounded and computes RF by the preceding argument. RF(T)\allowbreak{} exists by two-out-of-three. The complexes T and F(T)\allowbreak{} are zero below q\allowbreak{}+1,\allowbreak{} so section 4 gives H\textasciicircum{}j(RF(T)\allowbreak{})\allowbreak{}\allowbreak{}=0 for j<\allowbreak{}=q; the same vanishing holds for H\textasciicircum{}j(F(T)\allowbreak{})\allowbreak{} by its terms. The comparison of the two triangle cohomology sequences now proves

\[
 H^j(F(A))\xrightarrow{\sim}H^j(RF(A))\quad(j<q),
 \qquad H^q(F(A))\lhook\joinrel\longrightarrow H^q(RF(A)). \tag{5.2}
\]

There is an exact description of the possible cokernel in degree q. Let beta:\allowbreak{}H\textasciicircum{}q(F(P)\allowbreak{})\allowbreak{} \allowbreak{}-> H\textasciicircum{}q(RF(P)\allowbreak{})\allowbreak{} be the known isomorphism,\allowbreak{} delta\_\allowbreak{}F:\allowbreak{}H\textasciicircum{}q(F(P)\allowbreak{})\allowbreak{} \allowbreak{}-> H\textasciicircum{}\{q\allowbreak{}+1\}(F(T)\allowbreak{})\allowbreak{} the connecting map,\allowbreak{} and gamma\allowbreak{}=H\textasciicircum{}\{q\allowbreak{}+1\}(eta\_\allowbreak{}T)\allowbreak{}:\allowbreak{}H\textasciicircum{}\{q\allowbreak{}+1\}(F(T)\allowbreak{})\allowbreak{} \allowbreak{}-> H\textasciicircum{}\{q\allowbreak{}+1\}(RF(T)\allowbreak{})\allowbreak{}. The other connecting map satisfies delta\_\allowbreak{}R beta\allowbreak{}=gamma delta\_\allowbreak{}F,\allowbreak{} with the signs fixed by the same original split sequence. Set

\[
 O_q=\operatorname{im}(\delta_F)\cap\ker(\gamma)
                  \ \subset H^{q+1}(F(T)).                    \tag{5.3}
\]

Using beta to identify the middle terms,\allowbreak{} H\textasciicircum{}q(F(A)\allowbreak{})\allowbreak{} is ker(delta\_\allowbreak{}F)\allowbreak{} and H\textasciicircum{}q(RF(A)\allowbreak{})\allowbreak{} is ker(gamma delta\_\allowbreak{}F)\allowbreak{}. The map from the latter to O\_\allowbreak{}q is the restriction of delta\_\allowbreak{}F,\allowbreak{} and its kernel is the former. Its image is the intersection in (5.3)\allowbreak{}:\allowbreak{} pulling the epimorphism H\textasciicircum{}q(F(P)\allowbreak{})\allowbreak{} \allowbreak{}-> im(delta\_\allowbreak{}F)\allowbreak{} back along that intersection proves this image assertion in an arbitrary abelian category. Consequently

\[
0\to H^q(F(A))\xrightarrow{H^q(\eta_A)}H^q(RF(A))
                                      \longrightarrow O_q\to0.\tag{5.4}
\]

If F is left exact,\allowbreak{} gamma is an isomorphism and O\_\allowbreak{}q\allowbreak{}=0. To prove this without assuming RF exists on any additional object,\allowbreak{} let T be zero below c\allowbreak{}=q\allowbreak{}+1 and use the cofinal denominators T \allowbreak{}-> L with L zero below c. For each such L,\allowbreak{} left exactness gives H\textasciicircum{}c(F(L)\allowbreak{})\allowbreak{}\allowbreak{}=F(H\textasciicircum{}c(L)\allowbreak{})\allowbreak{},\allowbreak{} and the quasi-isomorphism identifies this with F(H\textasciicircum{}c(T)\allowbreak{})\allowbreak{}. Applying H\^{}c to the essential-constancy witness shows H\textasciicircum{}c(RF(T)\allowbreak{})\allowbreak{} has that value. The same kernel computation gives H\textasciicircum{}c(F(T)\allowbreak{})\allowbreak{}\allowbreak{}=F(H\textasciicircum{}c(T)\allowbreak{})\allowbreak{},\allowbreak{} and the induced map is precisely gamma. Thus (5.2)\allowbreak{} extends to an isomorphism for all j<\allowbreak{}=q when F is left exact. This is an actual comparison proof,\allowbreak{} not a presumption that F commutes with all cohomology.

The finite-prefix statement,\allowbreak{} (5.3)\allowbreak{}–\allowbreak{}(5.4)\allowbreak{},\allowbreak{} and this left-exact refinement are DERIVED-UNDER-007. Their hypotheses preserve the original bounded below A and the given existence of RF(A)\allowbreak{},\allowbreak{} while requiring acyclicity only through the specified degree q. Taking q\allowbreak{}=i\allowbreak{}+1 for every i when all terms are acyclic proves the source\textquotesingle s full Leray lemma exactly. The construction does not erase the differential out of degree q:\allowbreak{} H\textasciicircum{}q(F(A)\allowbreak{})\allowbreak{} still uses F(A\textasciicircum{}\{q\allowbreak{}+1\})\allowbreak{},\allowbreak{} even when that term is not assumed acyclic.

\subsection{\texorpdfstring{6. DERIVED-NEW-021:\allowbreak{} the enough-acyclics proof and its counterexample}{6.  the enough-acyclics proof and its counterexample}}\label{proofs-5091-5900-derived-new-021-the-enough-acyclics-proof-and-its-counterexample}

The proposition\textquotesingle s statement is valid for an arbitrary additive F. The induction at original 5856–\allowbreak{}5877 is not valid at that generality:\allowbreak{} its short exact sequences of image objects are triangles in D(A)\allowbreak{},\allowbreak{} whereas the earlier computing lemma applies to the exact functor K(A)\allowbreak{} \allowbreak{}-> D(B)\allowbreak{} and triangles in K(A)\allowbreak{}. A nonsplit short exact sequence need not give such a triangle there. Applying F to that sequence cannot be justified before the desired exactness has been proved.

Here is an explicit counterexample to the claimed acyclicity of all images,\allowbreak{} including the proposition\textquotesingle s enough-acyclics hypothesis. Fix a prime p. For an abelian group M let D(M)\allowbreak{} be its largest divisible subgroup and set

\[
 F(M)=M[p]/D(M)[p].                                      \tag{6.1}
\]

The sum of all divisible subgroups is divisible:\allowbreak{} any element is a finite sum,\allowbreak{} and its division by any positive integer is obtained by dividing each summand. Thus this sum defines D(M)\allowbreak{}. Every group homomorphism maps divisible subgroups into divisible subgroups,\allowbreak{} so it induces the displayed quotient map. These induced maps preserve composition,\allowbreak{} identities and addition,\allowbreak{} making (6.1)\allowbreak{} an additive functor. In particular D(M direct-sum N)\allowbreak{}\allowbreak{}=D(M)\allowbreak{} direct-sum D(N)\allowbreak{},\allowbreak{} by projection and by inclusion,\allowbreak{} and (6.1)\allowbreak{} preserves that biproduct with its maps.

Every group A embeds in a divisible group,\allowbreak{} without using injective resolutions:\allowbreak{} choose a free abelian group P surjecting onto A with kernel N. The inclusion P \allowbreak{}-> Q tensor\_\allowbreak{}Z P is injective because P is a sum of copies of Z. It induces an injection

\[
 A=P/N\longrightarrow (\mathbf Q\otimes_{\mathbf Z}P)/N.
\]

The target is divisible,\allowbreak{} since a quotient of a divisible group is divisible. Section 3 therefore resolves every bounded-below complex by a bounded-below complex of divisible groups. F vanishes on every divisible group,\allowbreak{} so it sends all these complexes and every transition between them to zero. Section 1,\allowbreak{} applied to this class of complexes,\allowbreak{} proves RF is everywhere defined on D\textasciicircum{}\allowbreak{}+(Ab)\allowbreak{} and is the zero functor. Thus M is right F-acyclic exactly when F(M)\allowbreak{}\allowbreak{}=0. Divisible groups and Z are acyclic,\allowbreak{} whereas F(Z/\allowbreak{}p)\allowbreak{}\allowbreak{}=Z/\allowbreak{}p is nonzero,\allowbreak{} so Z/\allowbreak{}p is not.

Consider the original cochain complex,\allowbreak{} in degrees 0 through 3,\allowbreak{}

\[
 I^\bullet:
 \mathbf Z\xrightarrow{\,p\,}\mathbf Z
 \xrightarrow{m\mapsto m/p+\mathbf Z}\mathbf Q/\mathbf Z
 \xrightarrow{\,p\,}\mathbf Q/\mathbf Z.                   \tag{6.2}
\]

Its first map is injective. The kernel of the middle map is pZ,\allowbreak{} equal to the preceding image. Its image is (p\textasciicircum{}\{-1\}Z)\allowbreak{}/\allowbreak{}Z,\allowbreak{} precisely the kernel of multiplication by p on Q/\allowbreak{}Z. The final map is surjective because Q/\allowbreak{}Z is divisible. Hence (6.2)\allowbreak{} is acyclic. Every term is right F-acyclic,\allowbreak{} and the category has enough such objects,\allowbreak{} as just proved. Nevertheless the source\textquotesingle s J\textasciicircum{}1\allowbreak{}=im(d\textasciicircum{}1)\allowbreak{} is isomorphic to Z/\allowbreak{}p and is not right F-acyclic. Thus the stated image induction fails under its own hypotheses. F(I\textasciicircum{}bullet)\allowbreak{} is still zero and hence acyclic; the counterexample refutes the proof\textquotesingle s intermediate claim,\allowbreak{} not the proposition\textquotesingle s conclusion.

The correct proof is shorter and preserves the claimed generality. An acyclic bounded-below complex I\^{}bullet is isomorphic to zero in D\textasciicircum{}\allowbreak{}+(A)\allowbreak{}. RF is defined at zero with value zero:\allowbreak{} the denominator category has the terminal zero index and F(0)\allowbreak{}\allowbreak{}=0. Denominator invariance therefore makes RF(I\textasciicircum{}bullet)\allowbreak{} defined with value zero. Because all its terms are right F-acyclic,\allowbreak{} the already proved Leray lemma gives the actual isomorphism

\[
 F(I^\bullet)\xrightarrow{\sim}RF(I^\bullet)=0.             \tag{6.3}
\]

This is not circular. Leray\textquotesingle s lemma requires existence only at the one complex in question,\allowbreak{} supplied here by its quasi-isomorphism to zero; it does not use the enough-acyclics proposition.

For a quasi-isomorphism between two bounded-below complexes of acyclic objects,\allowbreak{} its cone is bounded below and acyclic. Its terms are finite biproducts of acyclic objects:\allowbreak{} those are acyclic by the earlier direct-sum computing result applied to the degree-zero complexes. Additivity identifies F of this cone with the cone of the image map,\allowbreak{} including the original minus sign on the shifted first component. Equation (6.3)\allowbreak{} makes that image map a quasi-isomorphism. The source\textquotesingle s class resolutions and section 1 now prove existence everywhere and computation by every such complex. Reversing arrows and using the upper-bounded version of Leray gives the left-derived assertion. This is the proposed replacement of the faulty paragraph.

DERIVED-RECON-030 accepts the still-missing P02-E0035 article at 5826–\allowbreak{}5827,\allowbreak{} changing ``is quotient'' to ``is a quotient''.

\subsection{\texorpdfstring{7. Receiving proofs,\allowbreak{} including DERIVED-NEW-022 in Perfect Complexes}{7. Receiving  including DERIVED-NEW-022 in Perfect Complexes}}\label{proofs-5091-5900-receiving-proofs-including-derived-new-022-in-perfect-complexes}

The repaired proposition retains its entire statement. Two external direct uses were checked for the implication that receives it.

In More Algebra,\allowbreak{} lemma-compute-Rlim,\allowbreak{} current 23649–\allowbreak{}23665,\allowbreak{} every inverse system A embeds in the displayed ML system B\_\allowbreak{}n\allowbreak{}=A\_\allowbreak{}n direct-sum A\_\allowbreak{}\{n-1\} direct-sum ... direct-sum A\_\allowbreak{}1 via (1,\allowbreak{}f\_\allowbreak{}n,\allowbreak{}f\_\allowbreak{}\{n-1\}f\_\allowbreak{}n,\allowbreak{}...,\allowbreak{}f\_\allowbreak{}2...f\_\allowbreak{}n)\allowbreak{}. Its first projection splits the injection at each n,\allowbreak{} and B\textquotesingle s transition projections are surjective,\allowbreak{} so B is ML. The earlier reviewed Mittag–\allowbreak{}Leffler theorem supplies exactness of lim on the specified short exact sequences with ML kernel; quotients of ML systems are ML. Thus the right-acyclic-class criterion supplies the acyclic objects used in the bounded Rlim existence step. The repaired proposition proves exactly that step. The subsequent unbounded Rlim calculations are not certified by this bounded receiving check.

In Etale Cohomology,\allowbreak{} current 8872–\allowbreak{}8889,\allowbreak{} each discrete G-module is a quotient of a sum of Z[G/\allowbreak{}U]\allowbreak{}:\allowbreak{} a vector has open stabilizer U,\allowbreak{} and the basis coset gU maps to its translate by g. The underlying groups of these modules are free,\allowbreak{} hence flat over Z. In a short exact sequence with the last two terms flat,\allowbreak{} its first term is flat and tensoring with R stays exact; the Tor\_\allowbreak{}1 contribution from the last flat term is zero. These are precisely the left-acyclic-class hypotheses for H(M)\allowbreak{}\allowbreak{}=M tensor\_\allowbreak{}Z R. The dual repaired proposition therefore supplies the bounded-above LH and its computation stated there. The later adjunction and Ext comparison remain outside this receiving check.

The internal references at original 7003,\allowbreak{} 10007 and 10051 retain the same needed conclusion:\allowbreak{} enough known acyclic objects give the bounded derived functor,\allowbreak{} and bounded complexes of them compute it. Their separate injective,\allowbreak{} unbounded,\allowbreak{} and spectral-sequence arguments remain in their source-order review positions. Nothing in (6.3)\allowbreak{} requires adding left exactness to those uses.

For DERIVED-UNDER-007,\allowbreak{} the Cohomology remark at current 2195–\allowbreak{}2220 uses K\allowbreak{}=f\_\allowbreak{}*I\textasciicircum{}bullet and the left exact functor Gamma(Y,\allowbreak{}-)\allowbreak{}. It assumes all K\^{}n are Gamma-acyclic. For computing through degree q>\allowbreak{}=0 the proved refinement needs that acyclicity only for n<\allowbreak{}=q,\allowbreak{} while retaining the actual formula

\[
 H^q(\Gamma(Y,K^\bullet))=
 \ker(\Gamma(Y,K^q)\to\Gamma(Y,K^{q+1}))\big/
 \operatorname{im}(\Gamma(Y,K^{q-1})\to\Gamma(Y,K^q)).     \tag{7.1}
\]

The identity Gamma(Y,\allowbreak{}f\_\allowbreak{}*I\textasciicircum{}n)\allowbreak{}\allowbreak{}=Gamma(X,\allowbreak{}I\textasciicircum{}n)\allowbreak{} respects every displayed differential. Thus this application makes a degree-bounded comparison; it neither deletes I\textasciicircum{}\{q\allowbreak{}+1\} nor claims a new full derived-composition theorem. The original all-degree assertion is retained.

DERIVED-NEW-022 is a receiving proof repair in \texttt{perfect.\allowbreak{}tex},\allowbreak{} original and current 529–\allowbreak{}536. The argument passes to the canonical lower truncation tau\_\allowbreak{}\{>\allowbreak{}=-n\}F\textasciicircum{}bullet and then uses the acyclicity of its terms. Its bottom term is F\textasciicircum{}\{-n\}/\allowbreak{}im(d\textasciicircum{}\{-n-1\})\allowbreak{},\allowbreak{} whose acyclicity is not established by the assertion about the original terms. The following argument supplies the needed comparison while keeping all those original terms. The theorem\textquotesingle s conclusion and hypotheses stay unchanged; no geometric counterexample to that theorem is asserted.

Work over the source\textquotesingle s affine S. Keep its exact bound N from lemma-quasi-coherence-direct-image:\allowbreak{} for E with quasi-coherent cohomology and H\textasciicircum{}j(E)\allowbreak{}\allowbreak{}=0 for j>0,\allowbreak{} one has H\textasciicircum{}j(Rf\_\allowbreak{}*E)\allowbreak{}\allowbreak{}=0 for j>\allowbreak{}=N. Choose an integer n>\allowbreak{}=max(1,\allowbreak{}N)\allowbreak{}. Put

\[
 B^\bullet=\sigma_{\ge -n}\mathcal F^\bullet,
 \qquad C^\bullet=\sigma_{\le -n-1}\mathcal F^\bullet.
\]

The termwise split sequence 0 \allowbreak{}-> B \allowbreak{}-> F \allowbreak{}-> C \allowbreak{}-> 0 gives its actual distinguished triangle. Every cohomology sheaf of C is quasi-coherent,\allowbreak{} and it vanishes in degrees greater than -n-1. Apply the stated bound to C[-n-1]\allowbreak{}. Since H\textasciicircum{}j(C[-n-1]\allowbreak{})\allowbreak{}\allowbreak{}=H\textasciicircum{}\{j-n-1\}(C)\allowbreak{},\allowbreak{} this shift has cohomology zero above zero. Exactness and its specified shift comparison give

\[
 H^j(Rf_*C)=0\quad\text{for }j\ge N-n-1.                 \tag{7.2}
\]

The choice n>\allowbreak{}=N makes both H\textasciicircum{}\{-1\}(Rf\_\allowbreak{}*C)\allowbreak{} and H\textasciicircum{}0(Rf\_\allowbreak{}*C)\allowbreak{} zero. Thus H\textasciicircum{}0(Rf\_\allowbreak{}*B)\allowbreak{} \allowbreak{}-> H\textasciicircum{}0(Rf\_\allowbreak{}*F)\allowbreak{} is an isomorphism. The bound n>\allowbreak{}=1 also gives H\textasciicircum{}0(f\_\allowbreak{}*B)\allowbreak{}\allowbreak{}=H\textasciicircum{}0(f\_\allowbreak{}*F)\allowbreak{},\allowbreak{} with exactly the same terms in degrees -1,\allowbreak{}0,\allowbreak{}1 and the same two differentials. B is bounded below and consists of the original right f\_\allowbreak{}*-acyclic terms. Leray\textquotesingle s lemma therefore identifies f\_\allowbreak{}*B with Rf\_\allowbreak{}*B. Naturality with respect to B \allowbreak{}-> F now proves the desired H\^{}0 comparison for the original F. Applying this argument after every shift proves the full theorem.

The exact relation to the source\textquotesingle s original truncation is retained:\allowbreak{} if J\allowbreak{}=im(d\textasciicircum{}\{-n-1\})\allowbreak{},\allowbreak{} there is a short exact sequence of complexes

\[
 0\to J[n]\to\sigma_{\ge -n}\mathcal F^\bullet
               \to\tau_{\ge -n}\mathcal F^\bullet\to0.       \tag{7.3}
\]

In degree -n it is the actual kernel/\allowbreak{}image/\allowbreak{}quotient sequence; all other degrees are identities or zero. The sheaf J is quasi-coherent. The same bound gives H\textasciicircum{}j(Rf\_\allowbreak{}*J[n]\allowbreak{})\allowbreak{}\allowbreak{}=R\textasciicircum{}\{j\allowbreak{}+n\}f\_\allowbreak{}*J\allowbreak{}=0 for j\allowbreak{}+n>\allowbreak{}=N,\allowbreak{} so (7.3)\allowbreak{} makes the two derived truncations agree on H\^{}0 for the chosen n. Their bottom objects are not identified,\allowbreak{} and acyclicity of that quotient was not presumed. Equations (7.2)\allowbreak{}–\allowbreak{}(7.3)\allowbreak{} give the precise bridge and preserve the stated constant N and both truncations.

Finally,\allowbreak{} for the source\textquotesingle s exact-functor lemma at 5883–\allowbreak{}5900,\allowbreak{} exactness of F identifies F of kernels,\allowbreak{} images and cokernels with those of the image maps. Hence H\textasciicircum{}n(F(K)\allowbreak{})\allowbreak{}\allowbreak{}=F(H\textasciicircum{}n(K)\allowbreak{})\allowbreak{},\allowbreak{} naturally with the original differentials. It sends every quasi-isomorphism to a quasi-isomorphism. In the denominator diagram all transitions become isomorphisms in D(B)\allowbreak{},\allowbreak{} and the identity denominator has value F(K)\allowbreak{}. Section 1\textquotesingle s witness proves RF(K)\allowbreak{}\allowbreak{}=F(K)\allowbreak{},\allowbreak{} including for unbounded K. The bounded versions and R\textasciicircum{}iF\allowbreak{}=0 for i!\allowbreak{}=0 follow immediately by evaluating a degree-zero object. The reversed argument also proves the LF version.

These are bounded source and receiving reviews. Root-chapter search hits for other Leray and lower-vanishing uses are routing evidence,\allowbreak{} not claims to have read or accepted those entire proofs. The complete new operations and explicit reference changes are bound in the part07 record. The synthesis across all findings remains later work.
\clearpage
\section{A further degree in the left-exact finite-prefix comparison}\label{proofs-5091-5900-addendum-01-a-further-degree-in-the-left-exact-finite-prefix-comparison}

This strengthens the left-exact clause of DERIVED-UNDER-007 in section 5 of \texttt{PROOFS\_\allowbreak{}5091\_\allowbreak{}5900.\allowbreak{}md}. The earlier statement and its exact defect sequence remain valid. No source-body replacement or new claim ID is introduced,\allowbreak{} and no novelty is asserted.

Let F:\allowbreak{}A \allowbreak{}-> B be additive and left exact,\allowbreak{} let A\^{}bullet be bounded below,\allowbreak{} and suppose RF(A\textasciicircum{}bullet)\allowbreak{} is defined. If A\^{}n is right F-acyclic for n<\allowbreak{}=m,\allowbreak{} then its canonical comparison satisfies

\[
 H^j(F(A^\bullet))\xrightarrow{\sim}H^j(RF(A^\bullet))
       \quad(j\le m+1),
 \qquad
 H^{m+2}(F(A^\bullet))\lhook\joinrel\longrightarrow
                                  H^{m+2}(RF(A^\bullet)).       \tag{1}
\]

Thus for an isomorphism through degree q,\allowbreak{} acyclicity through degree q-1 suffices. This retains every term and differential of A,\allowbreak{} including the term A\textasciicircum{}\{q\allowbreak{}+1\} used to calculate H\^{}q of F(A)\allowbreak{}. No existence of RF on all complexes or on any additional cohomology object is assumed.

First consider a complex T zero below an integer a,\allowbreak{} with RF(T)\allowbreak{} defined. For a quasi-isomorphism s:\allowbreak{}T \allowbreak{}-> L with L also zero below a,\allowbreak{} the original cone C(s)\allowbreak{} has components T\textasciicircum{}\{n\allowbreak{}+1\} direct-sum L\^{}n and differential (-d\_\allowbreak{}T,\allowbreak{}0;s,\allowbreak{}d\_\allowbreak{}L)\allowbreak{}. It is zero below a-1 and acyclic,\allowbreak{} so

\[
 0\longrightarrow C(s)^{a-1}\longrightarrow C(s)^a
                         \longrightarrow C(s)^{a+1}
\]

is exact. Left exactness preserves this exact sequence,\allowbreak{} including its kernel in the middle. Hence H\textasciicircum{}\{a-1\}(F(C(s)\allowbreak{})\allowbreak{})\allowbreak{}\allowbreak{}=H\textasciicircum{}a(F(C(s)\allowbreak{})\allowbreak{})\allowbreak{}\allowbreak{}=0. Additivity identifies F(C(s)\allowbreak{})\allowbreak{} with C(F(s)\allowbreak{})\allowbreak{},\allowbreak{} preserving that entire matrix and its minus sign. The cone long exact sequence consequently makes H\textasciicircum{}a(F(s)\allowbreak{})\allowbreak{} an isomorphism and H\textasciicircum{}\{a\allowbreak{}+1\}(F(s)\allowbreak{})\allowbreak{} a monomorphism.

Targets L zero below a form a cofinal subcategory of T/\allowbreak{}Qis(A)\allowbreak{},\allowbreak{} as proved in section 4 of the main note. It includes the identity of T. Applying H\^{}a to the essential-constancy diagram,\allowbreak{} every transition is an isomorphism by the preceding calculation. Its canonical map from H\textasciicircum{}a(F(T)\allowbreak{})\allowbreak{} to the value H\textasciicircum{}a(RF(T)\allowbreak{})\allowbreak{} is therefore an isomorphism.

For the next degree,\allowbreak{} apply H\textasciicircum{}\{a\allowbreak{}+1\} to the original essential-constancy witness. Write c:\allowbreak{}H\textasciicircum{}\{a\allowbreak{}+1\}(F(T)\allowbreak{})\allowbreak{} \allowbreak{}-> H\textasciicircum{}\{a\allowbreak{}+1\}(RF(T)\allowbreak{})\allowbreak{} for the canonical map. The eventual factorization property at the identity index gives some denominator s:\allowbreak{}T \allowbreak{}-> L and a map v from H\textasciicircum{}\{a\allowbreak{}+1\}(RF(T)\allowbreak{})\allowbreak{} to H\textasciicircum{}\{a\allowbreak{}+1\}(F(L)\allowbreak{})\allowbreak{} such that

\[
 H^{a+1}(F(s))=v c.
\]

The left side is monic by the cone calculation; hence c is monic. This proof uses the witness itself and does not assume that an arbitrary colimit of monomorphisms exists or stays monic in B. We have proved the exact lowest-degree comparison

\[
 H^a(F(T))\simeq H^a(RF(T)),\qquad
 H^{a+1}(F(T))\lhook\joinrel\longrightarrow H^{a+1}(RF(T)). \tag{2}
\]

Now put P\allowbreak{}=sigma\_\allowbreak{}\{<\allowbreak{}=m\}A and T\allowbreak{}=sigma\_\allowbreak{}\{>\allowbreak{}=m\allowbreak{}+1\}A,\allowbreak{} keeping the termwise split sequence 0 \allowbreak{}-> T \allowbreak{}-> A \allowbreak{}-> P \allowbreak{}-> 0. P is bounded and consists of acyclic objects,\allowbreak{} so it computes RF. Two-out-of-three therefore defines RF(T)\allowbreak{}. All cohomology of F(P)\allowbreak{} and RF(P)\allowbreak{} above m is zero. At degree m\allowbreak{}+1 the comparison of the two triangle sequences is a map between the cokernels of

\[
 H^m(F(P))\longrightarrow H^{m+1}(F(T)),
 \qquad
 H^m(RF(P))\longrightarrow H^{m+1}(RF(T)).
\]

Both vertical maps are isomorphisms:\allowbreak{} the first by computation on P,\allowbreak{} the second by (2)\allowbreak{} with a\allowbreak{}=m\allowbreak{}+1. Their cokernels,\allowbreak{} respectively H\textasciicircum{}\{m\allowbreak{}+1\}(F(A)\allowbreak{})\allowbreak{} and H\textasciicircum{}\{m\allowbreak{}+1\}(RF(A)\allowbreak{})\allowbreak{},\allowbreak{} are therefore canonically isomorphic. The lower degrees are isomorphic by the main note\textquotesingle s argument. At degree m\allowbreak{}+2 the same triangle sequences identify H\textasciicircum{}\{m\allowbreak{}+2\}(F(A)\allowbreak{})\allowbreak{} with H\textasciicircum{}\{m\allowbreak{}+2\}(F(T)\allowbreak{})\allowbreak{} and H\textasciicircum{}\{m\allowbreak{}+2\}(RF(A)\allowbreak{})\allowbreak{} with H\textasciicircum{}\{m\allowbreak{}+2\}(RF(T)\allowbreak{})\allowbreak{},\allowbreak{} because both adjacent cohomology objects of P vanish. Under those identifications the comparison is the monomorphism in (2)\allowbreak{}. This proves (1)\allowbreak{} with the actual canonical maps.

For the receiving global-sections calculation in the main note,\allowbreak{} degree q therefore only requires Gamma(Y,\allowbreak{}-)\allowbreak{}-acyclicity of K\^{}n for n<\allowbreak{}=q-1. Its displayed formula (7.1)\allowbreak{},\allowbreak{} including K\textasciicircum{}\{q\allowbreak{}+1\} and both differentials,\allowbreak{} is unchanged. For the repaired H\^{}0 step in Perfect Complexes,\allowbreak{} the bounded-below original-term truncation only needs acyclicity in its negative degrees to obtain that H\^{}0 comparison; the original theorem supplies acyclicity of all terms,\allowbreak{} so its statement and the recorded source repair remain unchanged.
\clearpage
\section{\texorpdfstring{Serre subcategories,\allowbreak{} quotient complexes,\allowbreak{} and simultaneous replacements}{Serre  quotient  and simultaneous replacements}}\label{proofs-5910-6199-serre-subcategories-quotient-complexes-and-simultaneous-replacements}

This is an editorial review of original \texttt{derived.\allowbreak{}tex} lines 5910–\allowbreak{}6199 at authority commit a04446e5\allowbreak{}7ec1fbc2\allowbreak{}52a871af\allowbreak{}cec7752f\allowbreak{}b2807b14. The current chapter retains the previously admitted corrections and Verdier additions. The original statements and constructions remain identifiable. No novelty is asserted. All complexes have the original cochain grading and shifts.

\subsection{1. Weak Serre subcategories and the exact comparison}\label{proofs-5910-6199-weak-serre-subcategories-and-the-exact-comparison}

Let B be the given weak Serre subcategory of the abelian category A. For a complex X,\allowbreak{} write H\textasciicircum{}n(X)\allowbreak{} for its cohomology in A. By the definition and characterization in current Homology 2091–\allowbreak{}2176,\allowbreak{} B is strictly full,\allowbreak{} closed under kernels and cokernels of its internal morphisms and under extensions,\allowbreak{} and its inclusion into A is exact. Finite biproducts belong to B because they are extensions.

The class D\_\allowbreak{}B(A)\allowbreak{} of complexes with all H\^{}n in B is therefore additive and preserved by every integer shift. If X direct-sum Y belongs to it,\allowbreak{} put M\allowbreak{}=H\textasciicircum{}n(X)\allowbreak{} direct-sum H\textasciicircum{}n(Y)\allowbreak{}. The two endomorphisms

\[
 e_X=\begin{pmatrix}1&0\\0&0\end{pmatrix},\qquad
 e_Y=\begin{pmatrix}0&0\\0&1\end{pmatrix}:M\longrightarrow M
\]

are morphisms between objects of B. Their kernels are respectively H\textasciicircum{}n(Y)\allowbreak{} and H\textasciicircum{}n(X)\allowbreak{},\allowbreak{} which therefore belong to B. This proves the actual summand assertion underlying the corrected phrase at original 5948; it does not assume that B is closed under arbitrary A-subobjects.

For a distinguished triangle X \allowbreak{}-> Y \allowbreak{}-> Z \allowbreak{}-> X[1]\allowbreak{},\allowbreak{} the exact five-term segment is

\[
 H^n(X)\longrightarrow H^n(Y)\longrightarrow H^n(Z)
       \longrightarrow H^{n+1}(X)\longrightarrow H^{n+1}(Y).
\]

If X,\allowbreak{}Y lie in D\_\allowbreak{}B(A)\allowbreak{},\allowbreak{} the defining five-term condition puts H\textasciicircum{}n(Z)\allowbreak{} in B. Rotation gives every other two-out-of-three case. Cohomological upper or lower bounds are preserved under finite sums,\allowbreak{} summands,\allowbreak{} shifts and triangles,\allowbreak{} proving all stated bounded versions as well.

Exactness of B \allowbreak{}-> A gives the equality of the corresponding kernels,\allowbreak{} images,\allowbreak{} quotients and hence cohomology for any B-complex. Applying the inclusion to the original cone matrix preserves its entries and signs. It sends quasi-isomorphisms to quasi-isomorphisms and descends to the exact comparison D(B)\allowbreak{} \allowbreak{}-> D\_\allowbreak{}B(A)\allowbreak{},\allowbreak{} with the same assertion for \allowbreak{}+,\allowbreak{} -,\allowbreak{} b. No full faithfulness is inferred at this stage.

DERIVED-RECON-031 reconciles P02-E0036 and French DERIVED-026 with MC-STK-ERR-0837:\allowbreak{} the missing verb is already supplied. DERIVED-RECON-032 reconciles P02-E0037 and French DERIVED-027 with MC-STK-ERR-0838:\allowbreak{} the duplicated words are already removed,\allowbreak{} and the endomorphisms above justify the resulting mathematical sentence.

For a Serre subcategory B,\allowbreak{} let v:\allowbreak{}A \allowbreak{}-> C\allowbreak{}=A/\allowbreak{}B be the exact quotient. Since v(H\textasciicircum{}n(X)\allowbreak{})\allowbreak{}\allowbreak{}=H\textasciicircum{}n(vX)\allowbreak{} and ker(v)\allowbreak{}\allowbreak{}=B,\allowbreak{} the kernel of D(v)\allowbreak{} is exactly D\_\allowbreak{}B(A)\allowbreak{}. This proves the stated factorization through the Verdier quotient,\allowbreak{} rather than identifying its source with its target in advance.

\subsection{2. Lifting a complex from the abelian quotient}\label{proofs-5910-6199-lifting-a-complex-from-the-abelian-quotient}

Let the given C-complex have objects X\^{}i of A and differentials d\^{}i in C. Let S be the morphisms of A with kernel and cokernel in B. The roof calculus and its pushout/\allowbreak{}pullback stability are the ones in current Homology 2235–\allowbreak{}2289. We retain each step of the source\textquotesingle s two-sided construction,\allowbreak{} including the finite sums and intersections.

For i>\allowbreak{}=0,\allowbreak{} write d\textasciicircum{}i\allowbreak{}=v(f\textasciicircum{}i)\allowbreak{}v(s\textasciicircum{}i)\allowbreak{}\textasciicircum{}\{-1\},\allowbreak{} with s\textasciicircum{}i:\allowbreak{}Y\textasciicircum{}i \allowbreak{}-> X\^{}i in S and f\textasciicircum{}i:\allowbreak{}Y\textasciicircum{}i \allowbreak{}-> X\textasciicircum{}\{i\allowbreak{}+1\}. Starting with j\textasciicircum{}0\allowbreak{}=1 on X\textasciicircum{}0,\allowbreak{} form recursively

\[
 (X')^{i+1}
   =\operatorname{Coker}\big((j^is^i,-f^i):
             Y^i\longrightarrow (X')^i\oplus X^{i+1}\big).
                                                        \tag{2.1}
\]

Let a\textasciicircum{}i:\allowbreak{}(X')\allowbreak{}\textasciicircum{}i \allowbreak{}-> (X')\allowbreak{}\textasciicircum{}\{i\allowbreak{}+1\} and j\textasciicircum{}\{i\allowbreak{}+1\}:\allowbreak{}X\textasciicircum{}\{i\allowbreak{}+1\} \allowbreak{}-> (X')\allowbreak{}\textasciicircum{}\{i\allowbreak{}+1\} be the two pushout maps. The minus sign in (2.1)\allowbreak{} gives the exact identity a\^{}i j\^{}i s\textasciicircum{}i\allowbreak{}=j\textasciicircum{}\{i\allowbreak{}+1\} f\^{}i. Since j\^{}i s\^{}i is in S,\allowbreak{} its pushout j\textasciicircum{}\{i\allowbreak{}+1\} is in S. Induction proves this at every finite nonnegative degree. Applying v yields

\[
 v(a^i)v(j^i)=v(j^{i+1})d^i.
\]

Take (X')\allowbreak{}\textasciicircum{}n\allowbreak{}=X\textasciicircum{}n and j\textasciicircum{}n\allowbreak{}=1 for n<\allowbreak{}=0,\allowbreak{} and transport the negative differentials in C. This gives an isomorphism of C-complexes and represents every nonnegative differential by an actual A-map. No infinite pushout or limit has been used.

For the resulting negative differentials,\allowbreak{} write d\textasciicircum{}i\allowbreak{}=v(t\textasciicircum{}\{i\allowbreak{}+1\})\allowbreak{}\textasciicircum{}\{-1\}v(g\textasciicircum{}i)\allowbreak{},\allowbreak{} with t\textasciicircum{}\{i\allowbreak{}+1\}:\allowbreak{}X\textasciicircum{}\{i\allowbreak{}+1\} \allowbreak{}-> Z\textasciicircum{}\{i\allowbreak{}+1\} in S and g\textasciicircum{}i:\allowbreak{}X\textasciicircum{}i \allowbreak{}-> Z\textasciicircum{}\{i\allowbreak{}+1\}. Starting at j\textasciicircum{}0\allowbreak{}=1,\allowbreak{} form,\allowbreak{} by descending induction,\allowbreak{}

\[
 (X')^i=X^i\mathbin{\times}_{Z^{i+1}}(X')^{i+1},\quad
 j^i:(X')^i\longrightarrow X^i,\quad
 a^i:(X')^i\longrightarrow (X')^{i+1},                    \tag{2.2}
\]

where the second map to Z\textasciicircum{}\{i\allowbreak{}+1\} is t\textasciicircum{}\{i\allowbreak{}+1\}j\textasciicircum{}\{i\allowbreak{}+1\}. The pullback identity is g\^{}i j\textasciicircum{}i\allowbreak{}=t\textasciicircum{}\{i\allowbreak{}+1\}j\textasciicircum{}\{i\allowbreak{}+1\}a\textasciicircum{}i. Pullback stability of S makes j\^{}i an S-map,\allowbreak{} and after v the identity becomes d\^{}i v(j\textasciicircum{}i)\allowbreak{}\allowbreak{}=v(j\textasciicircum{}\{i\allowbreak{}+1\})\allowbreak{}v(a\textasciicircum{}i)\allowbreak{}. Keep the already constructed terms and maps in degrees >\allowbreak{}=0. We now have A-maps in every degree,\allowbreak{} whose successive composites vanish after v. These maps need not yet form an A-complex; that is precisely what the next two steps achieve.

DERIVED-RECON-033 /\allowbreak{} P02-E0040 corrects the four diagram labels in original 6063–\allowbreak{}6064 to t\textasciicircum{}\{-2\},\allowbreak{}g\textasciicircum{}\{-2\},\allowbreak{}t\textasciicircum{}\{-1\},\allowbreak{}g\textasciicircum{}\{-1\}. Their domains and codomains in (2.2)\allowbreak{} are respectively X\^{}\{-2\} \allowbreak{}-> Z\textasciicircum{}\{-2\},\allowbreak{} X\^{}\{-2\} \allowbreak{}-> Z\textasciicircum{}\{-1\},\allowbreak{} X\^{}\{-1\} \allowbreak{}-> Z\textasciicircum{}\{-1\},\allowbreak{} X\^{}\{-1\} \allowbreak{}-> Z\^{}0. The existing subscript labels do not match the declared cochain components. All four remain uncorrected in the current chapter and receive separate exact replacements.

To check the final source formulas,\allowbreak{} write a\_\allowbreak{}\{j,\allowbreak{}i\}\allowbreak{}=d\textasciicircum{}\{j-1\}...d\textasciicircum{}i for j\textgreater i and a\_\allowbreak{}\{i,\allowbreak{}i\}\allowbreak{}=1,\allowbreak{} using the A-maps just constructed. Put I\_\allowbreak{}k\allowbreak{}=Im(d\textasciicircum{}\{k-1\}d\textasciicircum{}\{k-2\})\allowbreak{} inside X\^{}k. Exactness of v implies v(I\_\allowbreak{}k)\allowbreak{}\allowbreak{}=0,\allowbreak{} so I\_\allowbreak{}k belongs to B. For n\textgreater0 retain the entire source expression

\[
 N^n=\sum_{0<k\le n}\operatorname{Im}
                  (I_k\longrightarrow X^n),\qquad
 \widetilde X^n=X^n/N^n,                                \tag{2.3}
\]

where the indicated map is the restriction of a\_\allowbreak{}\{n,\allowbreak{}k\}. Set N\textasciicircum{}n\allowbreak{}=0 for n<\allowbreak{}=0. Each N\^{}n belongs to B,\allowbreak{} since the sum is finite and B is Serre. Moreover d\textasciicircum{}n(N\textasciicircum{}n)\allowbreak{} is contained in N\textasciicircum{}\{n\allowbreak{}+1\},\allowbreak{} so all differentials descend. For n>\allowbreak{}=0 the image of d\^{}n d\^{}\{n-1\} is contained in I\_\allowbreak{}\{n\allowbreak{}+1\},\allowbreak{} one of the retained summands in N\textasciicircum{}\{n\allowbreak{}+1\}. Thus these composites are zero on the quotient. Each X\^{}n \allowbreak{}-> tilde X\^{}n becomes an isomorphism under v. This proves the first source replacement with its exact range and without dropping the k\allowbreak{}=1 contribution from X\^{}\{-1\}.

For the remaining step use the new objects tilde X and their induced maps; denote their path composites by tilde a. Set Q\textasciicircum{}n\allowbreak{}=tilde X\^{}n for n>\allowbreak{}=0,\allowbreak{} and for n\textless0 keep the full finite intersection

\[
 Q^n=\bigcap_{n\le k<0}
       (\widetilde a_{k,n})^{-1}
       \ker(\widetilde d^{k+1}\widetilde d^k).             \tag{2.4}
\]

Equivalently,\allowbreak{} Q\^{}n is the kernel of the map to the finite biproduct of the images of tilde a\_\allowbreak{}\{k\allowbreak{}+2,\allowbreak{}n\},\allowbreak{} n<\allowbreak{}=k<0. Each such image belongs to B,\allowbreak{} since v kills the consecutive two-map factor. Consequently tilde X\textasciicircum{}n/\allowbreak{}Q\textasciicircum{}n lies in B and Q\^{}n \allowbreak{}-> tilde X\^{}n becomes an isomorphism under v. If x denotes a generalized section of Q\textasciicircum{}n,\allowbreak{} every condition at degree n\allowbreak{}+1 follows by composing the corresponding condition for x with tilde d\^{}n; categorically,\allowbreak{} the image of tilde d\textasciicircum{}n|\_\allowbreak{}\{Q\textasciicircum{}n\} lies in each of the kernels defining Q\textasciicircum{}\{n\allowbreak{}+1\}. Thus Q is preserved by the differential. For n\textless0 its k\allowbreak{}=n condition says exactly that the two-map composite starting at Q\^{}n is zero. The composites starting in nonnegative degree already vanished by (2.3)\allowbreak{}. Hence Q is an A-complex. The maps in (2.1)\allowbreak{}–\allowbreak{}(2.4)\allowbreak{},\allowbreak{} after v,\allowbreak{} give actual chain isomorphisms connecting vQ to the original C-complex. This proves essential surjectivity with no omitted boundary degree.

\subsection{3. The left adjoint descends after the quotient}\label{proofs-5910-6199-the-left-adjoint-descends-after-the-quotient}

Keep the source hypothesis:\allowbreak{} u:\allowbreak{}C \allowbreak{}-> A is left adjoint to the exact quotient v:\allowbreak{}A \allowbreak{}-> C and vu is naturally isomorphic to Id\_\allowbreak{}C. Let eta:\allowbreak{}Id\_\allowbreak{}C \allowbreak{}-> vu and epsilon:\allowbreak{}uv \allowbreak{}-> Id\_\allowbreak{}A be the actual unit and counit. The latter natural isomorphism need not initially have been specified as eta,\allowbreak{} so we first justify that eta itself is invertible.

Write T\allowbreak{}=vu,\allowbreak{} mu\allowbreak{}=v epsilon u,\allowbreak{} and choose the given natural isomorphism alpha:\allowbreak{}T \allowbreak{}-> Id. Define natural endomorphisms of Id\_\allowbreak{}C by

\[
 e=\alpha\eta,\qquad
 m=\alpha\mu\,T(\alpha^{-1})\alpha^{-1}.                 \tag{3.1}
\]

Naturality of eta and the monad unit identity mu eta\_\allowbreak{}T\allowbreak{}=1 give

\[
 me=\alpha\mu T(\alpha^{-1})\eta
    =\alpha\mu\eta_T\alpha^{-1}=1.
\]

Natural endomorphisms of the identity commute:\allowbreak{} naturality of e at the morphism m\_\allowbreak{}X gives e\_\allowbreak{}X m\_\allowbreak{}X\allowbreak{}=m\_\allowbreak{}X e\_\allowbreak{}X for every X. Thus em\allowbreak{}=1 as well. It follows that eta\allowbreak{}=alpha\textasciicircum{}\{-1\}e is invertible. The adjunction triangle identity now gives the specific formula

\[
 v(\epsilon_X)=\eta_{vX}^{-1}.                            \tag{3.2}
\]

The left adjoint u is additive. One direct verification is that its left-adjoint property preserves zero objects and finite coproducts; in an abelian category these are biproducts,\allowbreak{} and the zero maps,\allowbreak{} projections and diagonals are characterized by their biproduct equations. Hence u preserves addition as defined by diagonal,\allowbreak{} biproduct and codiagonal. It therefore acts termwise on complexes and homotopies and preserves the cone matrix,\allowbreak{} with its minus sign.

It is not necessary,\allowbreak{} and is generally false,\allowbreak{} that u send every C-quasi-isomorphism to an A-quasi-isomorphism. Let

\[
 T_A=D(A)/D_B(A),\qquad q:D(A)\longrightarrow T_A.
\]

If s is a quasi-isomorphism of C-complexes,\allowbreak{} then v u C(s)\allowbreak{} is isomorphic to C(s)\allowbreak{} via eta\^{}\{-1\}. Thus H\^{}n(u C(s)\allowbreak{})\allowbreak{} lies in B for every n. It follows that q(u C(s)\allowbreak{})\allowbreak{}\allowbreak{}=0 and q(u s)\allowbreak{} is invertible. The universal property of localization therefore gives an exact functor

\[
 G:D(C)\longrightarrow T_A,\qquad G(Y)=q(uY),              \tag{3.3}
\]

where uY means the termwise construction before localization. The localization proof is what makes (3.3)\allowbreak{} independent of a chosen complex representative. Naturality of eta and epsilon on chain maps extends to inverted arrows by multiplying their naturality identities by inverses. We obtain natural isomorphisms

\[
 F G\xrightarrow{\eta^{-1}}\operatorname{Id}_{D(C)},\qquad
 G F\xrightarrow{q(\epsilon)}\operatorname{Id}_{T_A},      \tag{3.4}
\]

where F:\allowbreak{}T\_\allowbreak{}A \allowbreak{}-> D(C)\allowbreak{} is induced by v. The second is invertible because its cone has cohomology in B by (3.2)\allowbreak{}. For arbitrary X,\allowbreak{}Y in D(A)\allowbreak{},\allowbreak{} the inverse on Hom of F is the explicit map

\[
 a:vX\longrightarrow vY
 \quad\longmapsto\quad
 q(\epsilon_Y)\,G(a)\,q(\epsilon_X)^{-1}.                 \tag{3.5}
\]

Applying F to (3.5)\allowbreak{},\allowbreak{} formula (3.2)\allowbreak{} and naturality of eta give a. Conversely,\allowbreak{} naturality of q(epsilon)\allowbreak{} shows that (3.5)\allowbreak{} takes F(b)\allowbreak{} back to b for every b:\allowbreak{}qX \allowbreak{}-> qY,\allowbreak{} including roofs. This proves full faithfulness and essential surjectivity of F with the original maps.

For each of \allowbreak{}+,\allowbreak{}-,\allowbreak{}b perform the same construction on complexes whose terms satisfy the indicated bound. Additivity of u preserves their zero terms. An acyclic cone still maps to a complex with B-cohomology,\allowbreak{} so (3.3)\allowbreak{} descends to D\textasciicircum{}*(C)\allowbreak{} \allowbreak{}-> D\textasciicircum{}*(A)\allowbreak{}/\allowbreak{}D\_\allowbreak{}B\textasciicircum{}*(A)\allowbreak{}. All identities (3.2)\allowbreak{}–\allowbreak{}(3.5)\allowbreak{} remain within these bounded categories. One never assumes that u preserves cohomological bounds on an arbitrary representative.

DERIVED-RECON-034 confirms MC-STK-ERR-0839\textquotesingle s parenthesis repair for P02-E0038 and French DERIVED-028. That repair is valid but insufficient to express the scope of (3.5)\allowbreak{}:\allowbreak{} original 6128 and current 6481 quantify only X,\allowbreak{}Y in A. \textbf{DERIVED-NEW-023} changes that domain to D(A)\allowbreak{} and adds the termwise descent explanation after original 6118. This is a proof-domain repair and clarification,\allowbreak{} not a counterexample to the quotient theorem or an assertion that u is exact. The old punctuation correction remains preserved in the new expression.

For a concrete check that the distinction is necessary,\allowbreak{} fix a prime p. Let A be the abelian category of left modules over

\[
 R=\begin{pmatrix}\mathbf Z&0\\\mathbf F_p&\mathbf F_p\end{pmatrix}.
\]

Equivalently its objects are triples (M,\allowbreak{}N,\allowbreak{}phi)\allowbreak{},\allowbreak{} where M is an abelian group,\allowbreak{} N an F\_\allowbreak{}p-vector space and phi:\allowbreak{}M/\allowbreak{}pM \allowbreak{}-> N is F\_\allowbreak{}p-linear. A morphism is a pair (f,\allowbreak{}g)\allowbreak{} satisfying g phi\allowbreak{}=phi' bar f. Kernels and cokernels are computed on the two components with the induced structure map,\allowbreak{} as also follows by applying the two complementary idempotents of R. The functor v(M,\allowbreak{}N,\allowbreak{}phi)\allowbreak{}\allowbreak{}=M is exact. Its left adjoint is u(M)\allowbreak{}\allowbreak{}=(M,\allowbreak{}M/\allowbreak{}pM,\allowbreak{}1)\allowbreak{}:\allowbreak{} a morphism u(M)\allowbreak{} \allowbreak{}-> (M',\allowbreak{}N',\allowbreak{}phi')\allowbreak{} is uniquely determined by f:\allowbreak{}M \allowbreak{}-> M',\allowbreak{} its second component being phi\textquotesingle{} bar f. Thus vu\allowbreak{}=Id,\allowbreak{} and the counit has components (1,\allowbreak{}phi)\allowbreak{}.

The Serre kernel B consists of triples (0,\allowbreak{}N,\allowbreak{}0)\allowbreak{}. Every counit has kernel and cokernel in B; in the localization it identifies every object with u(vX)\allowbreak{}. The same pair-of-components adjunction proves that the induced A/\allowbreak{}B \allowbreak{}-> Ab is an equivalence,\allowbreak{} so this is a concrete instance of the source setup. The injection p:\allowbreak{}Z \allowbreak{}-> Z is sent by u to (p,\allowbreak{}0)\allowbreak{},\allowbreak{} which has the nonzero kernel (0,\allowbreak{}F\_\allowbreak{}p,\allowbreak{}0)\allowbreak{}. More explicitly,\allowbreak{} the acyclic complex Z --p-\allowbreak{}-> Z \allowbreak{}-> Z/\allowbreak{}p in degrees 0,\allowbreak{}1,\allowbreak{}2 is sent to a complex whose first component is that exact complex and whose second is F\_\allowbreak{}p --0-\allowbreak{}-> F\_\allowbreak{}p --1-\allowbreak{}-> F\_\allowbreak{}p. Its degree-zero cohomology is (0,\allowbreak{}F\_\allowbreak{}p,\allowbreak{}0)\allowbreak{},\allowbreak{} and all other cohomology is zero. It is not acyclic in A; its cohomology lies in B and vanishes in the required quotient,\allowbreak{} exactly as (3.3)\allowbreak{} asserts. No contribution has been removed.

\subsection{4. Subcomplex replacements and all derived morphisms}\label{proofs-5910-6199-subcomplex-replacements-and-all-derived-morphisms}

Now assume precisely the hypothesis of original 6143–\allowbreak{}6150:\allowbreak{} B is Serre in A,\allowbreak{} and every epimorphism X \allowbreak{}-> Y with Y in B admits a B-subobject X\textquotesingle{} of X which still maps onto Y. Let X\^{}bullet be bounded above with H\textasciicircum{}i(X)\allowbreak{} in B,\allowbreak{} and retain arbitrary prescribed B-subobjects B\^{}i of X\^{}i.

Choose C\^{}i inside ker(d\textasciicircum{}i)\allowbreak{} in B surjecting onto H\textasciicircum{}i(X)\allowbreak{}. Such a choice is possible by applying the hypothesis to ker(d\textasciicircum{}i)\allowbreak{} \allowbreak{}-> H\textasciicircum{}i(X)\allowbreak{}. Keep the original formula

\[
 D^i=C^i+d^{i-1}(B^{i-1})+B^i.                           \tag{4.1}
\]

Every summand,\allowbreak{} and their finite sum,\allowbreak{} lies in B. The image d\textasciicircum{}i(D\textasciicircum{}i)\allowbreak{} equals d\textasciicircum{}i(B\textasciicircum{}i)\allowbreak{},\allowbreak{} which is contained in D\textasciicircum{}\{i\allowbreak{}+1\}; hence D is a subcomplex. The chosen C\^{}i show that H\textasciicircum{}i(D)\allowbreak{} \allowbreak{}-> H\textasciicircum{}i(X)\allowbreak{} is epic.

Above the upper term bound set E\textasciicircum{}i\allowbreak{}=0. Descending inductively,\allowbreak{} put

\[
 V^{i+1}=(D^{i+1}+E^{i+1})\cap\operatorname{Im}(d^i).
\]

This object lies in B,\allowbreak{} because it is a subobject of a B-object. Apply the hypothesis to the epimorphism (d\textasciicircum{}i)\allowbreak{}\textasciicircum{}\{-1\}(V\textasciicircum{}\{i\allowbreak{}+1\})\allowbreak{} \allowbreak{}-> V\textasciicircum{}\{i\allowbreak{}+1\} to obtain E\^{}i inside that preimage,\allowbreak{} with E\^{}i in B and d\textasciicircum{}i(E\textasciicircum{}i)\allowbreak{}\allowbreak{}=V\textasciicircum{}\{i\allowbreak{}+1\}. Set Y\textasciicircum{}i\allowbreak{}=D\textasciicircum{}i\allowbreak{}+E\textasciicircum{}i. The formulas give an actual subcomplex Y of X,\allowbreak{} containing each prescribed B\textasciicircum{}i,\allowbreak{} with terms in B. Its cohomology remains surjective onto that of X. For injectivity,\allowbreak{} the exact subobject identity is

\[
 \ker(d^i|_{Y^i})\cap\operatorname{Im}(d^{i-1}_X)
   =Y^i\cap\operatorname{Im}(d^{i-1}_X)
   =d^{i-1}(E^{i-1})=\operatorname{Im}(d^{i-1}_Y).        \tag{4.2}
\]

The first equality uses d\^{}i d\textasciicircum{}\{i-1\}\allowbreak{}=0,\allowbreak{} the second is the inductive choice,\allowbreak{} and the last follows from E\^{}\{i-1\} contained in Y\^{}\{i-1\} together with the reverse containment for every Y-boundary. Thus Y \allowbreak{}-> X is a quasi-isomorphism. This verifies the entire subcomplex claim without replacing images by unspecified elements.

For faithfulness first take a chain map f:\allowbreak{}U \allowbreak{}-> V between bounded above B-complexes whose derived image in A is zero. The fraction calculus supplies a quasi-isomorphism s:\allowbreak{}V \allowbreak{}-> X into a bounded above A-complex and a homotopy h\textasciicircum{}i:\allowbreak{}U\textasciicircum{}i \allowbreak{}-> X\^{}\{i-1\} satisfying

\[
 (sf)^i=d_X^{i-1}h^i+h^{i+1}d_U^i.                      \tag{4.3}
\]

Take the prescribed subobjects B\textasciicircum{}i\allowbreak{}=Im(h\textasciicircum{}\{i\allowbreak{}+1\})\allowbreak{}\allowbreak{}+Im(s\textasciicircum{}i)\allowbreak{}. They belong to B. The construction gives a quasi-isomorphism Y \allowbreak{}-> X through which s and h factor. Write s':\allowbreak{}V \allowbreak{}-> Y for the resulting chain map. Since both s and Y \allowbreak{}-> X are quasi-isomorphisms,\allowbreak{} so is s\textquotesingle. Equation (4.3)\allowbreak{} holds in Y,\allowbreak{} by monicity of Y \allowbreak{}-> X,\allowbreak{} and makes s\textquotesingle f null-homotopic there. Hence f is zero in D\textasciicircum{}-(B)\allowbreak{}. For an arbitrary derived arrow represented by f t\textasciicircum{}\{-1\},\allowbreak{} vanishing after inclusion implies that f vanishes after inclusion,\allowbreak{} because t is already a quasi-isomorphism. The preceding argument applies,\allowbreak{} proving faithfulness for every derived arrow.

For fullness,\allowbreak{} represent any A-derived arrow U \allowbreak{}-> V by a right roof

\[
 U\xrightarrow{f}X\xleftarrow{s}V,\qquad s\text{ a quasi-isomorphism}.
                                                               \tag{4.4}
\]

Take B\textasciicircum{}i\allowbreak{}=Im(f\textasciicircum{}i)\allowbreak{}\allowbreak{}+Im(s\textasciicircum{}i)\allowbreak{}. A replacement Y \allowbreak{}-> X containing these images factors the roof into B-complexes. The factor s':\allowbreak{}V \allowbreak{}-> Y is a quasi-isomorphism,\allowbreak{} and (s')\allowbreak{}\textasciicircum{}\{-1\}f' in D\textasciicircum{}-(B)\allowbreak{} maps to the original arrow. This is the exact fullness argument abbreviated in the source. Essential surjectivity follows from the construction with B\textasciicircum{}i\allowbreak{}=0.

DERIVED-RECON-035 verifies French DERIVED-029 and MC-STK-ERR-0840:\allowbreak{} after faithfulness,\allowbreak{} the remaining assertion is indeed Fullness. P02-E0039 also reports original 7396–\allowbreak{}7397. Only its 6197–\allowbreak{}6198 component is reconciled here. The physical report remains partly open until its second context is reviewed; it is not counted among the completely resolved reports at this checkpoint.

\subsection{\texorpdfstring{5. DERIVED-UNDER-008:\allowbreak{} simultaneous factorization and consequences}{5.  simultaneous factorization and consequences}}\label{proofs-5910-6199-derived-under-008-simultaneous-factorization-and-consequences}

Keep all hypotheses of section 4. For a fixed bounded above X with B-cohomology,\allowbreak{} let P\_\allowbreak{}X be its B-term subcomplexes Y for which Y \allowbreak{}-> X is a quasi-isomorphism,\allowbreak{} ordered by inclusion. It is nonempty by section 4. Given Y\_\allowbreak{}1,\allowbreak{}...,\allowbreak{}Y\_\allowbreak{}r,\allowbreak{} apply that construction with B\textasciicircum{}i\allowbreak{}=sum\_\allowbreak{}j Y\_\allowbreak{}j\textasciicircum{}i. It yields Y containing every Y\_\allowbreak{}j. Thus P\_\allowbreak{}X is filtered; as an inclusion poset it has no distinct parallel arrows requiring equalization. Each inclusion Y \allowbreak{}-> Y\textquotesingle{} in P\_\allowbreak{}X is itself a quasi-isomorphism by two-out-of-three.

More strongly,\allowbreak{} any finite family of chain maps from bounded above B-complexes into X factors simultaneously through one Y in P\_\allowbreak{}X. Any specified finite family of homotopies between those maps is retained there as well. In degree i prescribe the sum of all map images from degree i and all homotopy images from degree i\allowbreak{}+1. Each is a quotient of a B-object,\allowbreak{} so this is a B-subobject of X\^{}i. The construction and the monicity argument in (4.3)\allowbreak{} prove the claim with all equations unchanged.

In particular,\allowbreak{} for a fixed bounded above B-complex U,\allowbreak{} the canonical maps induce the exact Hom formula

\[
 \mathop{\rm colim}_{Y\in P_X}
       \operatorname{Hom}_{K^-(B)}(U,Y)
       \xrightarrow{\ \sim\ }
       \operatorname{Hom}_{K^-(A)}(U,X).                  \tag{5.1}
\]

Here the colimit is realized by the group on the right:\allowbreak{} every chain map factors by the image construction,\allowbreak{} and if two factorizations give the same homotopy class in X,\allowbreak{} include both target subcomplexes and the images of one witnessing homotopy in a common replacement. They then agree in its homotopy category. This proves precisely the equality relation in the filtered colimit. Sums are compatible by passing to a common replacement and adding the factor maps,\allowbreak{} so (5.1)\allowbreak{} is an isomorphism of abelian groups,\allowbreak{} not only a bijection. The same argument proves its universal property for compatible families of homomorphisms. Thus no preservation of a previously unconstructed colimit in A is assumed.

The following diagram records the actual simultaneous replacement; the diagonal inclusions and all maps shown are chain maps,\allowbreak{} and the maps from Y\_\allowbreak{}1,\allowbreak{}Y\_\allowbreak{}2,\allowbreak{}Y to X are quasi-isomorphisms.

\[
 \begin{array}{ccccc}
 Y_1&\lhook\joinrel\longrightarrow&Y&\longleftarrow\joinrel\rhook&Y_2\\[-2pt]
 &&\big\downarrow&&\\[-2pt]
 &&X&&
 \end{array}
\]

Both paths to X are the given subcomplex inclusions. The homotopy operators,\allowbreak{} when prescribed,\allowbreak{} factor term by term through this same Y with their original degree -1.

The equivalence of section 4 restricts to

\[
 D^b(B)\xrightarrow{\sim}D^b_B(A).                      \tag{5.2}
\]

Indeed,\allowbreak{} if the A-cohomology of a resulting Y vanishes below a,\allowbreak{} exactness of B \allowbreak{}-> A gives the same vanishing in B. Replace Y,\allowbreak{} within B,\allowbreak{} by its canonical truncation tau\_\allowbreak{}\{>\allowbreak{}=a\}Y. Its degree-a term is coker(d\textasciicircum{}\{a-1\}\_\allowbreak{}Y)\allowbreak{},\allowbreak{} an object of B,\allowbreak{} its higher terms are the original Y\textasciicircum{}i,\allowbreak{} and the map Y \allowbreak{}-> tau\_\allowbreak{}\{>\allowbreak{}=a\}Y is a quasi-isomorphism. The upper term bound is retained. Full faithfulness follows by restricting the already proved D\^{}- equivalence and the bounded embeddings. This proves (5.2)\allowbreak{} without asserting that this canonical truncation is a subcomplex of X.

For M,\allowbreak{}N in B,\allowbreak{} the comparison therefore identifies,\allowbreak{} for every n,\allowbreak{}

\[
 \operatorname{Hom}_{D^b(B)}(M[0],N[n])
   \xrightarrow{\sim}
 \operatorname{Hom}_{D^b(A)}(M[0],N[n]).                 \tag{5.3}
\]

These are the derived-category Ext groups. Compatibility with composition and shift follows from the original exact inclusion; in particular their products are preserved. We do not use an unproved additional identification with a different Ext presentation.

These are editorial consequences of the source construction. The bounded restriction is also used in the existing Perfect proposition below; it is not advertised as a result missing from the corpus. The source body is not rewritten to incorporate (5.1)\allowbreak{}–\allowbreak{}(5.3)\allowbreak{}.

\subsection{6. Checked receiving implications}\label{proofs-5910-6199-checked-receiving-implications}

In current Adequate Modules 2013–\allowbreak{}2031,\allowbreak{} u is the associated adequate module functor,\allowbreak{} v is restriction to the small Zariski site,\allowbreak{} and the unit is an isomorphism. Crucially,\allowbreak{} the source explicitly says u is not left exact in general. The adjunction and unit are described at 1737–\allowbreak{}1765; parasitic objects form the Serre kernel,\allowbreak{} as in 2040–\allowbreak{}2114. Applying (3.2)\allowbreak{} to the actual adjunction shows that the cone of the termwise counit has parasitic cohomology. Equations (3.3)\allowbreak{}–\allowbreak{}(3.5)\allowbreak{} therefore give the inverse and both natural isomorphisms for the equivalence asserted at 2164–\allowbreak{}2179,\allowbreak{} including its bounded versions. The clarification preserves this receiving theorem and does not add the false requirement that its u be exact. Only this dependency implication is reviewed here,\allowbreak{} not every earlier adequate-module lemma.

For current Perfect Complexes 2564–\allowbreak{}2592 set A\allowbreak{}=QCoh(O\_\allowbreak{}X)\allowbreak{},\allowbreak{} B\allowbreak{}=Coh(O\_\allowbreak{}X)\allowbreak{},\allowbreak{} with X Noetherian. The kernels,\allowbreak{} quotients and extensions required for the Serre property are those in current Coherent Sheaves 2448–\allowbreak{}2488:\allowbreak{} on a Noetherian affine,\allowbreak{} submodules and quotients of finite modules are finite,\allowbreak{} and extensions of finite modules are finite. The local exact module-to-quasi-coherent-sheaf correspondence gives the stated sheaf operations.

For an epimorphism G \allowbreak{}-> F from a quasi-coherent sheaf to a coherent one,\allowbreak{} Properties 3009–\allowbreak{}3033 writes G as a directed union of its finite-type quasi-coherent submodules. These are coherent on a Noetherian scheme. Their images form a directed union equal to F. Choose a finite affine cover of X. On each affine,\allowbreak{} finitely many generators of F are contained in one of the image submodules; by directedness and the finiteness of the cover,\allowbreak{} one coherent submodule of G has image all of F. This verifies the epimorphism hypothesis used in section 4 with the receiving objects unchanged.

Consequently the concrete coherent subcomplex construction in (4.1)\allowbreak{}–\allowbreak{}(4.2)\allowbreak{},\allowbreak{} and its simultaneous map/\allowbreak{}homotopy refinement (5.1)\allowbreak{},\allowbreak{} apply to bounded above quasi-coherent complexes with coherent cohomology. The bounded restriction (5.2)\allowbreak{} and groups (5.3)\allowbreak{} follow for Coh inside QCoh. Current Perfect 2595–\allowbreak{}2631 already uses the bounded restriction in its larger comparison with all O\_\allowbreak{}X-modules. That further comparison invokes its separate proposition-Noetherian; we do not claim a new review of that proposition from this bounded dependency check. No Perfect source change follows from UNDER-008.
\clearpage
\section{\texorpdfstring{Injective and projective resolutions:\allowbreak{} signs,\allowbreak{} relative lifts,\allowbreak{} and fixed outer terms}{Injective and projective   relative  and fixed outer terms}}\label{proofs-6209-7089-injective-and-projective-resolutions-signs-relative-lifts-and-fixed-outer-terms}

This editorial review covers original Derived Categories 6209–\allowbreak{}7089 at authority commit a04446e5\allowbreak{}7ec1fbc2\allowbreak{}52a871af\allowbreak{}cec7752f\allowbreak{}b2807b14. Earlier corrections are checked against that source and the frozen current chapter. The two strengthenings below remain separate from the source body. They are consequences of the reviewed constructions; no novelty is asserted. Cohomological shifts have differential d\_\allowbreak{}\{X[1]\allowbreak{}\}\textasciicircum{}n\allowbreak{}=-d\_\allowbreak{}X\textasciicircum{}\{n\allowbreak{}+1\},\allowbreak{} and a homotopy from a map f to zero satisfies f\textasciicircum{}n\allowbreak{}=d\_\allowbreak{}Y\textasciicircum{}\{n-1\}h\textasciicircum{}n\allowbreak{}+h\textasciicircum{}\{n\allowbreak{}+1\}d\_\allowbreak{}X\textasciicircum{}n.

\subsection{1. Existence and the original bounds}\label{proofs-6209-7089-existence-and-the-original-bounds}

An injective resolution of an object A is the source\textquotesingle s complex I with I\textasciicircum{}n\allowbreak{}=0 for n<0,\allowbreak{} injective terms,\allowbreak{} A identified with ker(d\_\allowbreak{}I\textasciicircum{}0)\allowbreak{},\allowbreak{} and H\textasciicircum{}n(I)\allowbreak{}\allowbreak{}=0 for n\textgreater0. Thus A[0]\allowbreak{} \allowbreak{}-> I is a quasi-isomorphism. A resolution of a general complex K requires I to be bounded below,\allowbreak{} so its existence forces H\textasciicircum{}n(K)\allowbreak{}\allowbreak{}=0 in all sufficiently small degrees.

Conversely,\allowbreak{} if H\textasciicircum{}n(K)\allowbreak{}\allowbreak{}=0 for n<a,\allowbreak{} use the actual canonical truncation q:\allowbreak{}K \allowbreak{}-> tau\_\allowbreak{}\{>\allowbreak{}=a\}K. Its degree-a term is coker(d\_\allowbreak{}K\textasciicircum{}\{a-1\})\allowbreak{},\allowbreak{} its terms above a are K\textasciicircum{}n,\allowbreak{} and its terms below a are zero. The map q induces the identity identifications on H\^{}n for n>\allowbreak{}=a and is an isomorphism on the zero cohomology below a. It is therefore a quasi-isomorphism.

If the category has enough injectives,\allowbreak{} embed A into I\textasciicircum{}0,\allowbreak{} then embed I\textasciicircum{}0/\allowbreak{}A into I\textasciicircum{}1,\allowbreak{} then coker(d\_\allowbreak{}I\textasciicircum{}0)\allowbreak{} into I\textasciicircum{}2,\allowbreak{} and continue. The differential is the quotient map followed by the chosen monomorphism. Its kernel is the preceding image,\allowbreak{} including ker(d\_\allowbreak{}I\textasciicircum{}0)\allowbreak{}\allowbreak{}=A; successive differentials compose to zero. This proves the object resolution with all its original terms. DERIVED-RECON-036 reconciles P02-E0041 and French DERIVED-033 with MC-STK-ERR-0844:\allowbreak{} the quotient at original 6297 is I\textasciicircum{}0/\allowbreak{}A,\allowbreak{} not the undefined I\_\allowbreak{}0/\allowbreak{}A.

For a complex whose terms are zero below a,\allowbreak{} apply the source\textquotesingle s subcategory-right-resolution lemma,\allowbreak{} original 5386–\allowbreak{}5405,\allowbreak{} to the class of injective objects. It produces a termwise monomorphism K \allowbreak{}-> I,\allowbreak{} an isomorphism on cohomology,\allowbreak{} and I\textasciicircum{}n\allowbreak{}=0 for n\textless a. The construction and its exact pushout maps are proved in section 3 of PROOFS\_\allowbreak{}5091\_\allowbreak{}5900.md,\allowbreak{} by dualizing its fully written pullback induction. This application satisfies both hypotheses:\allowbreak{} zero is injective,\allowbreak{} and every object embeds into an injective. For K with only the cohomological bound,\allowbreak{} compose the quasi-isomorphism q above with the resolution of tau\_\allowbreak{}\{>\allowbreak{}=a\}K. That composite need not be a termwise monomorphism; no such assertion is made in the cohomologically bounded case.

\subsection{2. The null-homotopy induction and its corrected indices}\label{proofs-6209-7089-the-null-homotopy-induction-and-its-corrected-indices}

Let K be acyclic,\allowbreak{} let every I\^{}n be injective,\allowbreak{} and suppose I\textasciicircum{}n\allowbreak{}=0 for n\textless a. For a chain map alpha:\allowbreak{}K \allowbreak{}-> I set h\textasciicircum{}j\allowbreak{}=0 for j<\allowbreak{}=a. Suppose h has been defined through h\^{}n with

\[
 \alpha^j=d_I^{j-1}h^j+h^{j+1}d_K^j\quad(j<n).
\]

The next defect is b\textasciicircum{}n\allowbreak{}=alpha\textasciicircum{}n-d\_\allowbreak{}I\textasciicircum{}\{n-1\}h\textasciicircum{}n:\allowbreak{}K\textasciicircum{}n \allowbreak{}-> I\^{}n. With every domain explicit,\allowbreak{} the source calculation is

\[
\begin{aligned}
 b^n d_K^{n-1}
 &=\alpha^n d_K^{n-1}-d_I^{n-1}h^n d_K^{n-1}\\
 &=d_I^{n-1}\alpha^{n-1}-d_I^{n-1}h^n d_K^{n-1}\\
 &=d_I^{n-1}(d_I^{n-2}h^{n-1}+h^n d_K^{n-1})
                      -d_I^{n-1}h^n d_K^{n-1}=0.          
\end{aligned}
\tag{2.1}\]

Thus b\^{}n kills Im(d\_\allowbreak{}K\textasciicircum{}\{n-1\})\allowbreak{}\allowbreak{}=ker(d\_\allowbreak{}K\textasciicircum{}n)\allowbreak{}. The universal property of the cokernel identifies its induced map with a map Im(d\_\allowbreak{}K\textasciicircum{}n)\allowbreak{} \allowbreak{}-> I\^{}n. Extend it along Im(d\_\allowbreak{}K\textasciicircum{}n)\allowbreak{} \allowbreak{}-> K\textasciicircum{}\{n\allowbreak{}+1\},\allowbreak{} using injectivity of I\textasciicircum{}n,\allowbreak{} to h\textasciicircum{}\{n\allowbreak{}+1\}:\allowbreak{}K\textasciicircum{}\{n\allowbreak{}+1\} \allowbreak{}-> I\^{}n. The resulting equation is precisely alpha\textasciicircum{}n\allowbreak{}=d\_\allowbreak{}I\textasciicircum{}\{n-1\}h\textasciicircum{}n\allowbreak{}+h\textasciicircum{}\{n\allowbreak{}+1\}d\_\allowbreak{}K\textasciicircum{}n. Induction proves the null-homotopy in all degrees without an assumption on the lower terms of K.

DERIVED-RECON-037 reconciles P02-E0042 and French DERIVED-034 with all six operations in MC-STK-ERR-0845. Five occurrences of h\^{}\{n-1\} outside the previous-degree identity must be h\textasciicircum{}n,\allowbreak{} and the first alpha in original 6328 must be alpha\^{}n. The occurrence h\^{}\{n-1\} inside d\_\allowbreak{}I\textasciicircum{}\{n-2\}h\textasciicircum{}\{n-1\} in (2.1)\allowbreak{} remains unchanged. The current source implements this exact selective correction.

For a quasi-isomorphism alpha:\allowbreak{}K \allowbreak{}-> L,\allowbreak{} its original cone C(alpha)\allowbreak{} is acyclic. The distinguished homotopy-category triangle is (K,\allowbreak{}L,\allowbreak{}C(alpha)\allowbreak{},\allowbreak{}alpha,\allowbreak{}i,\allowbreak{}-p)\allowbreak{}. Applying Hom\_\allowbreak{}K(-,\allowbreak{}I)\allowbreak{} gives the exact segment

\[
 \operatorname{Hom}_K(C(\alpha),I)\longrightarrow
 \operatorname{Hom}_K(L,I)\xrightarrow{\alpha^*}
 \operatorname{Hom}_K(K,I)\longrightarrow
 \operatorname{Hom}_K(C(\alpha)[-1],I).
\]

Both outer groups vanish by the just proved induction,\allowbreak{} including the shifted acyclic source. Hence alpha\^{}* is an isomorphism. This proves existence and uniqueness of a lift up to homotopy independently of the direct proofs that follow.

\subsection{3. Strict lifting and the remaining sign error}\label{proofs-6209-7089-strict-lifting-and-the-remaining-sign-error}

Suppose alpha:\allowbreak{}K \allowbreak{}-> L is a termwise monomorphic quasi-isomorphism and gamma:\allowbreak{}K \allowbreak{}-> I is given. We construct beta:\allowbreak{}L \allowbreak{}-> I with beta alpha\allowbreak{}=gamma. Set beta\textasciicircum{}j\allowbreak{}=0 in degrees j\textless a. Assume the construction through degree n-1 is compatible with the differentials. To prescribe beta\^{}n on Im(d\_\allowbreak{}L\textasciicircum{}\{n-1\})\allowbreak{},\allowbreak{} first show that d\_\allowbreak{}I\textasciicircum{}\{n-1\}beta\textasciicircum{}\{n-1\} kills ker(d\_\allowbreak{}L\textasciicircum{}\{n-1\})\allowbreak{}.

It kills Im(d\_\allowbreak{}L\textasciicircum{}\{n-2\})\allowbreak{} by the already established chain equation. Its restriction to ker(d\_\allowbreak{}L\textasciicircum{}\{n-1\})\allowbreak{} therefore induces a map H\textasciicircum{}\{n-1\}(L)\allowbreak{} \allowbreak{}-> I\^{}n. Composing with H\textasciicircum{}\{n-1\}(alpha)\allowbreak{} gives zero:\allowbreak{} on K-cycles it is d\_\allowbreak{}I\textasciicircum{}\{n-1\}gamma\textasciicircum{}\{n-1\}\allowbreak{}=gamma\textasciicircum{}n d\_\allowbreak{}K\textasciicircum{}\{n-1\}\allowbreak{}=0. Since H\textasciicircum{}\{n-1\}(alpha)\allowbreak{} is an isomorphism,\allowbreak{} the induced map is zero. This proves the required factorization through Im(d\_\allowbreak{}L\textasciicircum{}\{n-1\})\allowbreak{}. At the starting degree it also follows immediately from I\textasciicircum{}\{a-1\}\allowbreak{}=0.

The prescription on alpha(K\textasciicircum{}n)\allowbreak{} is gamma\^{}n. The intersection identity needed to glue it to the prescription on Im(d\_\allowbreak{}L\textasciicircum{}\{n-1\})\allowbreak{} is

\[
 \alpha(\operatorname{Im}d_K^{n-1})
       =\alpha(K^n)\cap\operatorname{Im}d_L^{n-1}.         \tag{3.1}
\]

The forward inclusion follows from the chain equation. For the reverse,\allowbreak{} pull back the intersection along the monomorphism alpha\^{}n. Its objects lie in ker(d\_\allowbreak{}K\textasciicircum{}n)\allowbreak{},\allowbreak{} because alpha\textasciicircum{}\{n\allowbreak{}+1\} is monic; their images in H\textasciicircum{}n(L)\allowbreak{} vanish. Injectivity of H\textasciicircum{}n(alpha)\allowbreak{} forces their images in H\textasciicircum{}n(K)\allowbreak{} to vanish,\allowbreak{} so this pullback is contained in Im(d\_\allowbreak{}K\textasciicircum{}\{n-1\})\allowbreak{}. This proves the subobject identity (3.1)\allowbreak{} in the abelian category. On that intersection the two maps agree by beta\^{}\{n-1\}alpha\^{}\{n-1\} \allowbreak{}=gamma\textasciicircum{}\{n-1\} and the chain equation for gamma. They consequently give a map from the sum of the two subobjects to I\^{}n. Injectivity extends it to all of L\^{}n. Induction proves strict lifting. This also supplies the implicit well-definedness argument in the source\textquotesingle s direct proof; no additional source-body reconstruction is required.

For general alpha use exactly the source\textquotesingle s factorization at 2622–\allowbreak{}2724:\allowbreak{}

\[
 \widetilde L^n=L^n\oplus K^n\oplus K^{n+1},\qquad
 d_{\widetilde L}^n=
 \begin{pmatrix}d_L^n&0&0\\0&d_K^n&1\\0&0&-d_K^{n+1}\end{pmatrix},
\] \[
 \widetilde\alpha=(\alpha,1,0)^t,\qquad
 \pi=(1,0,0),\qquad s=(1,0,0)^t.                         \tag{3.2}
\]

Here pi tilde-alpha\allowbreak{}=alpha,\allowbreak{} pi s\allowbreak{}=1,\allowbreak{} and the original homotopy H(l,\allowbreak{}k,\allowbreak{}k\textasciicircum{}\allowbreak{}+)\allowbreak{}\allowbreak{}=(0,\allowbreak{}0,\allowbreak{}k)\allowbreak{} satisfies

\[
 (dH+Hd)(l,k,k^+)=(0,k,-dk)+(0,0,dk+k^+)=(0,k,k^+)
                        =(1-s\pi)(l,k,k^+).             \tag{3.3}
\]

Thus tilde-alpha is a split monomorphic quasi-isomorphism. A strict lift tilde-beta gives beta\allowbreak{}=tilde-beta s,\allowbreak{} with beta alpha homotopic to gamma by (3.3)\allowbreak{}. DERIVED-RECON-038 reconciles both occurrences,\allowbreak{} 6412 and 6502,\allowbreak{} in P02-E0043 and French DERIVED-035 with the two MC-STK-ERR-0846 operations restoring the complex symbol K\^{}bullet.

For the source\textquotesingle s direct uniqueness argument,\allowbreak{} after this replacement write L\textasciicircum{}n\allowbreak{}=K\textasciicircum{}n direct-sum M\textasciicircum{}n,\allowbreak{} where M is acyclic. The given homotopies satisfy

\[
 \gamma^n-\beta_i^n|_{K^n}
       =d_I^{n-1}h_i^n+h_i^{n+1}d_K^n.
\]

Subtracting the two displayed equations in the actual order gives

\[
 \beta_1^n|_{K^n}-\beta_2^n|_{K^n}
 =d_I^{n-1}(h_2^n-h_1^n)
                  +(h_2^{n+1}-h_1^{n+1})d_K^n.          \tag{3.4}
\]

\textbf{DERIVED-NEW-024:\allowbreak{}} original 6524 and 6527,\allowbreak{} still unchanged in the current chapter,\allowbreak{} use h\_\allowbreak{}1-h\_\allowbreak{}2 instead. Correct both to h\_\allowbreak{}2-h\_\allowbreak{}1. The original maps beta\_\allowbreak{}1,\allowbreak{}beta\_\allowbreak{}2,\allowbreak{}gamma and the original convention for each h\_\allowbreak{}i are retained; reversing a convention elsewhere would conceal the actual error.

Extend h\_\allowbreak{}2\textasciicircum{}n-h\_\allowbreak{}1\textasciicircum{}n by zero on the chosen M\^{}n summand and call it h\^{}n. Then

\[
 r^n=\beta_1^n-\beta_2^n
                -(d_I^{n-1}h^n+h^{n+1}d_L^n)            \tag{3.5}
\]

is a chain map and is zero on K. It factors through the quotient q:\allowbreak{}L \allowbreak{}-> M. Write r\allowbreak{}=tq. Section 2 gives t\allowbreak{}=d\_\allowbreak{}I k\allowbreak{}+k d\_\allowbreak{}M,\allowbreak{} so beta\_\allowbreak{}1-beta\_\allowbreak{}2\allowbreak{}=d\_\allowbreak{}I(h\allowbreak{}+kq)\allowbreak{}\allowbreak{}+(h\allowbreak{}+kq)\allowbreak{}d\_\allowbreak{}L. This is the required homotopy. DERIVED-RECON-039 confirms the separate MC-STK-ERR-0847 parenthesis repair reported by P02-E0044 and French DERIVED-036. It is retained in (3.5)\allowbreak{}; it did not correct the independent sign in (3.4)\allowbreak{}.

For a counterexample to the uncorrected intermediate formula,\allowbreak{} work in abelian groups and take K\allowbreak{}=L\allowbreak{}=I\allowbreak{}=(Q --1-\allowbreak{}-> Q)\allowbreak{} in degrees 0,\allowbreak{}1. All terms of I are injective. Put alpha\allowbreak{}=1,\allowbreak{} gamma\allowbreak{}=0,\allowbreak{} beta\_\allowbreak{}1\allowbreak{}=1,\allowbreak{} beta\_\allowbreak{}2\allowbreak{}=0,\allowbreak{} h\_\allowbreak{}1\textasciicircum{}1\allowbreak{}=-1 and h\_\allowbreak{}2\textasciicircum{}1\allowbreak{}=0,\allowbreak{} with all other homotopy components zero. These satisfy both given homotopy equations. The source\textquotesingle s h\_\allowbreak{}1-h\_\allowbreak{}2 has coboundary -1,\allowbreak{} whereas beta\_\allowbreak{}1-beta\_\allowbreak{}2\allowbreak{}=1. Its proposed residual is 2 times the identity and is not zero on K\allowbreak{}=L. The corrected h\_\allowbreak{}2-h\_\allowbreak{}1 has coboundary 1 and gives zero residual. Thus the formula needs correction,\allowbreak{} while the uniqueness theorem stays valid.

Finally,\allowbreak{} every derived arrow L \allowbreak{}-> I is a roof gamma alpha\^{}\{-1\}. Lifting gamma produces a chain map beta representing it. If the class of a chain map beta becomes zero,\allowbreak{} some quasi-isomorphism alpha makes beta alpha null-homotopic; the injectivity of alpha\^{}* proved in section 2 makes beta null-homotopic. Consequently

\[
 \operatorname{Hom}_{K(A)}(L,I)
       \xrightarrow{\sim}\operatorname{Hom}_{D(A)}(L,I).  \tag{3.6}
\]

\textbf{DERIVED-NEW-025:\allowbreak{}} the introductions of this lemma at original 6543 and its projective dual at 6824 each omit the article before ``bounded''. Insert ``a'' in those two sentences. This is a copyedit,\allowbreak{} with no changed bounds or mathematical hypotheses.

\subsection{\texorpdfstring{4. DERIVED-UNDER-009:\allowbreak{} degree windows and relative homotopies}{4.  degree windows and relative homotopies}}\label{proofs-6209-7089-derived-under-009-degree-windows-and-relative-homotopies}

Keep I\^{}n injective and zero for n\textless a. The proof of section 2 only uses exactness of the source at degrees n>\allowbreak{}=a. It therefore proves:\allowbreak{}

\[
 H^n(C)=0\ (n\ge a)
 \quad\Longrightarrow\quad
 \text{every chain map }C\longrightarrow I
                    \text{ is null-homotopic}.          \tag{4.1}
\]

Indeed,\allowbreak{} start h\textasciicircum{}j\allowbreak{}=0 for j<\allowbreak{}=a as before. At each subsequent step n,\allowbreak{} the only source exactness used is Im(d\_\allowbreak{}C\textasciicircum{}\{n-1\})\allowbreak{}\allowbreak{}=ker(d\_\allowbreak{}C\textasciicircum{}n)\allowbreak{}. No vanishing below a is used or claimed.

Now let alpha:\allowbreak{}K \allowbreak{}-> L induce cohomology isomorphisms only in degrees n>\allowbreak{}=a. Every chain map gamma:\allowbreak{}K \allowbreak{}-> I factors uniquely through the canonical map K \allowbreak{}-> tau\_\allowbreak{}\{>\allowbreak{}=a\}K:\allowbreak{} its degree-a component kills Im(d\_\allowbreak{}K\textasciicircum{}\{a-1\})\allowbreak{},\allowbreak{} and all higher terms of the truncation are unchanged. A homotopy h\textasciicircum{}n:\allowbreak{}K\textasciicircum{}n \allowbreak{}-> I\^{}\{n-1\} has h\textasciicircum{}n\allowbreak{}=0 for n<\allowbreak{}=a,\allowbreak{} and its remaining components also factor uniquely. The homotopy equation at degree a descends because d\_\allowbreak{}K\textasciicircum{}a kills Im(d\_\allowbreak{}K\textasciicircum{}\{a-1\})\allowbreak{}. These statements give canonical,\allowbreak{} functorial isomorphisms of homotopy-class groups

\[
 \operatorname{Hom}_K(\tau_{\ge a}K,I)
       \xrightarrow{\sim}\operatorname{Hom}_K(K,I),
 \qquad
 \operatorname{Hom}_K(\tau_{\ge a}L,I)
       \xrightarrow{\sim}\operatorname{Hom}_K(L,I).        \tag{4.2}
\]

The actual map tau\_\allowbreak{}\{>\allowbreak{}=a\}(alpha)\allowbreak{} is a quasi-isomorphism. Apply section 2 to it and use (4.2)\allowbreak{}. We obtain

\[
 \alpha^*:\operatorname{Hom}_K(L,I)
             \xrightarrow{\sim}\operatorname{Hom}_K(K,I)
 \quad\text{if }H^n(\alpha)\text{ is an isomorphism for }n\ge a.
                                                               \tag{4.3}
\]

If alpha is also termwise monomorphic,\allowbreak{} the truncated map is termwise monomorphic. In its bottom degree this follows exactly from (3.1)\allowbreak{},\allowbreak{} using the injectivity of H\textasciicircum{}a(alpha)\allowbreak{}; above a it is the original monomorphism. Applying strict lifting to the truncated map and composing with L \allowbreak{}-> tau\_\allowbreak{}\{>\allowbreak{}=a\}L therefore gives,\allowbreak{} for every gamma,\allowbreak{} a strict lift beta with beta alpha\allowbreak{}=gamma.

There is a further relative assertion for two such strict lifts. Their difference factors as L \allowbreak{}-> C\allowbreak{}=L/\allowbreak{}K \allowbreak{}-> I. The long exact sequence and the isomorphisms H\textasciicircum{}n(alpha)\allowbreak{},\allowbreak{}H\textasciicircum{}\{n\allowbreak{}+1\}(alpha)\allowbreak{} give H\textasciicircum{}n(C)\allowbreak{}\allowbreak{}=0 for n>\allowbreak{}=a. Apply (4.1)\allowbreak{} to the factored map C \allowbreak{}-> I and compose its homotopy with L \allowbreak{}-> C. The resulting homotopy h satisfies

\[
 \beta_1-\beta_2=d_Ih+h d_L,\qquad h^n\alpha^n=0
                                                    \text{ for every }n. \tag{4.4}
\]

Thus both strict lifts and their uniqueness can be taken relative to the entire original subcomplex K. The source theorem only states ordinary uniqueness up to homotopy.

For precision about duality,\allowbreak{} pass from an A-complex K to the A-opposite complex DK with

\[
 (DK)^n=K^{-n},\qquad d_{DK}^n=(d_K^{-n-1})^{op}.
\]

It reverses chain maps and satisfies D(K[1]\allowbreak{})\allowbreak{}\allowbreak{}=(DK)\allowbreak{}[-1]\allowbreak{},\allowbreak{} including the minus sign in both shifted differentials. A homotopy h dualizes with component (h\textasciicircum{}\{-m\allowbreak{}+1\})\allowbreak{}\textasciicircum{}\{op\} in degree m; reversing the two compositions in its equation gives the same sum of the two required homotopy terms. Applying (4.1)\allowbreak{}–\allowbreak{}(4.4)\allowbreak{} in A-opposite proves the exact dual:\allowbreak{} if P consists of projectives with P\textasciicircum{}n\allowbreak{}=0 for n\textgreater b and alpha:\allowbreak{}L \allowbreak{}-> K induces cohomology isomorphisms for n<\allowbreak{}=b,\allowbreak{} then

\[
 \operatorname{Hom}_K(P,L)\xrightarrow{\sim}\operatorname{Hom}_K(P,K).
                                                               \tag{4.5}
\]

For termwise epimorphic alpha,\allowbreak{} any map P \allowbreak{}-> K lifts strictly. Two strict lifts beta\_\allowbreak{}1,\allowbreak{}beta\_\allowbreak{}2 have a homotopy h with alpha\textasciicircum{}\{n-1\}h\textasciicircum{}n\allowbreak{}=0 in every degree. This dual statement retains the original complexes and upper bound; it does not replace them by unspecified projective models.

\subsection{5. Projective resolutions and the cohomology bounds in lifting}\label{proofs-6209-7089-projective-resolutions-and-the-cohomology-bounds-in-lifting}

Under this exact duality,\allowbreak{} sections 1–\allowbreak{}3 prove the stated projective resolution results:\allowbreak{} bounded-above resolutions,\allowbreak{} canonical truncation tau\_\allowbreak{}\{<\allowbreak{}=b\}K with degree-b term ker(d\_\allowbreak{}K\textasciicircum{}b)\allowbreak{},\allowbreak{} termwise epimorphic resolutions for complexes whose terms vanish above b,\allowbreak{} null-homotopy of maps from bounded-above projectives to acyclic complexes,\allowbreak{} strict lifts through termwise epimorphic quasi-isomorphisms,\allowbreak{} uniqueness,\allowbreak{} and the Hom\_\allowbreak{}K to Hom\_\allowbreak{}D isomorphism.

DERIVED-RECON-041 reconciles P02-E0046 and French DERIVED-038 with the two MC-STK-ERR-0849 replacements of Comp\textasciicircum{}\allowbreak{}+ by Comp\^{}- at original 6844–\allowbreak{}6845. The plus-bounded statement would be false:\allowbreak{} let C\textasciicircum{}n\allowbreak{}=Z for n>\allowbreak{}=0,\allowbreak{} C\textasciicircum{}n\allowbreak{}=0 otherwise,\allowbreak{} and all differentials zero. The sequence 0 \allowbreak{}-> 0 \allowbreak{}-> C \allowbreak{}-> C \allowbreak{}-> 0,\allowbreak{} with the last map the identity,\allowbreak{} is a short exact sequence of bounded-below complexes. But C has nonzero cohomology in every nonnegative degree and cannot have a bounded-above projective resolution. The corrected minus-bounded statement is precisely the dual of the injective construction.

For original lemma-precise-vanishing,\allowbreak{} retain the bounded projective complex P,\allowbreak{} the term bound P\textasciicircum{}i\allowbreak{}=0 for i<n,\allowbreak{} and the cohomology bound H\textasciicircum{}i(K)\allowbreak{}\allowbreak{}=0 for i>\allowbreak{}=n. The actual canonical map tau\_\allowbreak{}\{<\allowbreak{}=n-1\}K \allowbreak{}-> K is a quasi-isomorphism. Every chain map P \allowbreak{}-> tau\_\allowbreak{}\{<\allowbreak{}=n-1\}K is zero,\allowbreak{} because in each degree one of its source or target objects is zero. The projective Hom comparison consequently gives

\[
 \operatorname{Hom}_K(P,K)=\operatorname{Hom}_D(P,K)=0.    \tag{5.1}
\]

The source also argues with the canonical truncation triangle. That triangle is in D(A)\allowbreak{}; using the same projective Hom comparison on all its vertices justifies its displayed Hom\_\allowbreak{}K groups. No termwise splitting of the canonical truncation is presumed.

For lemma-lift-map,\allowbreak{} let alpha:\allowbreak{}E \allowbreak{}-> L induce cohomology isomorphisms above n and an epimorphism in degree n,\allowbreak{} and retain the same P with terms zero below n. In the cone long exact sequence,\allowbreak{} for every i>\allowbreak{}=n the map H\textasciicircum{}i(E)\allowbreak{} \allowbreak{}-> H\textasciicircum{}i(L)\allowbreak{} is epic and H\textasciicircum{}\{i\allowbreak{}+1\}(E)\allowbreak{} \allowbreak{}-> H\textasciicircum{}\{i\allowbreak{}+1\}(L)\allowbreak{} is monic. Therefore H\textasciicircum{}i(C(alpha)\allowbreak{})\allowbreak{}\allowbreak{}=0 for i>\allowbreak{}=n. Equation (5.1)\allowbreak{} makes i beta zero in Hom\_\allowbreak{}K(P,\allowbreak{}C(alpha)\allowbreak{})\allowbreak{}. The exact Hom sequence of the original cone triangle yields gamma:\allowbreak{}P \allowbreak{}-> E with alpha gamma homotopic to beta. This proves the claimed lift with precisely the source\textquotesingle s endpoint surjectivity,\allowbreak{} rather than a stronger assumed quasi-isomorphism.

\subsection{6. Short exact sequences and DERIVED-UNDER-010}\label{proofs-6209-7089-short-exact-sequences-and-derived-under-010}

Consider 0 \allowbreak{}-> A --j-\allowbreak{}-> B --q-\allowbreak{}-> C \allowbreak{}-> 0 in bounded-below complexes. The source construction first chooses,\allowbreak{} or retains,\allowbreak{} a resolution f\_\allowbreak{}1:\allowbreak{}A \allowbreak{}-> I\_\allowbreak{}1. Its pushout is

\[
 E=\operatorname{Coker}\big((f_1,-j):A\longrightarrow I_1\oplus B\big).
                                                               \tag{6.1}
\]

It has the induced short exact sequence 0 \allowbreak{}-> I\_\allowbreak{}1 \allowbreak{}-> E \allowbreak{}-> C \allowbreak{}-> 0 and a map of short exact sequences from the original one. The outer maps are f\_\allowbreak{}1 and the identity of C,\allowbreak{} so the long exact cohomology sequences give that B \allowbreak{}-> E is a quasi-isomorphism. Since each I\_\allowbreak{}1\textasciicircum{}n is injective,\allowbreak{} the lower sequence splits in every degree. Choose the splittings and write

\[
 E^n=I_1^n\oplus C^n,\qquad
 d_E^n=\begin{pmatrix}d_{I_1}^n&\delta^n\\0&d_C^n\end{pmatrix}.
                                                               \tag{6.2}
\]

The full off-diagonal identity is d\_\allowbreak{}\{I\_\allowbreak{}1\}\textasciicircum{}\{n\allowbreak{}+1\}delta\textasciicircum{}n\allowbreak{}+delta\textasciicircum{}\{n\allowbreak{}+1\}d\_\allowbreak{}C\textasciicircum{}n\allowbreak{}=0. Thus delta:\allowbreak{}C \allowbreak{}-> I\_\allowbreak{}1[1]\allowbreak{} is a chain map with the source\textquotesingle s positive upper-right entry. This is exactly the map pi\textasciicircum{}\{n\allowbreak{}+1\}d\_\allowbreak{}E\textasciicircum{}n s\^{}n in current Homology 3606–\allowbreak{}3632; its cohomology map is the connecting map by the subsequent lemma at 3636–\allowbreak{}3648.

For the source\textquotesingle s proof,\allowbreak{} choose a termwise monomorphic resolution f\_\allowbreak{}3:\allowbreak{}C \allowbreak{}-> I\_\allowbreak{}3. Strict lifting,\allowbreak{} with target I\_\allowbreak{}1[1]\allowbreak{},\allowbreak{} gives delta':\allowbreak{}I\_\allowbreak{}3 \allowbreak{}-> I\_\allowbreak{}1[1]\allowbreak{} satisfying delta'f\_\allowbreak{}3\allowbreak{}=delta. Put

\[
 I_2^n=I_1^n\oplus I_3^n,\qquad
 d_{I_2}^n=\begin{pmatrix}d_{I_1}^n&(\delta')^n\\0&d_{I_3}^n\end{pmatrix}.
                                                               \tag{6.3}
\]

Its squared differential is zero,\allowbreak{} because the off-diagonal entry is d\_\allowbreak{}\{I\_\allowbreak{}1\}\textasciicircum{}\{n\allowbreak{}+1\}(delta')\allowbreak{}\textasciicircum{}n\allowbreak{}+(delta')\allowbreak{}\textasciicircum{}\{n\allowbreak{}+1\}d\_\allowbreak{}\{I\_\allowbreak{}3\}\textasciicircum{}n\allowbreak{}=0. The diagonal map E \allowbreak{}-> I\_\allowbreak{}2 is a chain map by delta'f\_\allowbreak{}3\allowbreak{}=delta. It is a quasi-isomorphism by the two outer quasi-isomorphisms in the short exact sequences. Composing with B \allowbreak{}-> E gives the required diagram. Every I\_\allowbreak{}2\textasciicircum{}n is injective:\allowbreak{} extension of a map to its finite biproduct is obtained by extending the two components. All complexes retain any common lower term bound,\allowbreak{} so the three source clauses,\allowbreak{} including compatibility with a prescribed I\_\allowbreak{}1,\allowbreak{} hold.

DERIVED-RECON-040 reconciles P02-E0045 and French DERIVED-037 with MC-STK-ERR-0848. The literal extra item ``add more here'' supplies no mathematical statement. Its deletion leaves exactly the three proved clauses. The strengthening below is not inserted in place of that placeholder.

\textbf{DERIVED-UNDER-010.} Both outer resolutions can in fact be prescribed:\allowbreak{} take any given f\_\allowbreak{}1:\allowbreak{}A \allowbreak{}-> I\_\allowbreak{}1 and f\_\allowbreak{}3:\allowbreak{}C \allowbreak{}-> I\_\allowbreak{}3 satisfying the source definition,\allowbreak{} with no monomorphism requirement on f\_\allowbreak{}3. Once these resolutions exist,\allowbreak{} no additional enough-injectives hypothesis is needed. Keep (6.1)\allowbreak{}–\allowbreak{}(6.2)\allowbreak{}. Ordinary lifting up to homotopy gives delta':\allowbreak{}I\_\allowbreak{}3 \allowbreak{}-> I\_\allowbreak{}1[1]\allowbreak{} and maps h\textasciicircum{}n:\allowbreak{}C\textasciicircum{}n \allowbreak{}-> I\_\allowbreak{}1\textasciicircum{}n with the exact equation

\[
 (\delta')^n f_3^n-\delta^n
       =-d_{I_1}^n h^n+h^{n+1}d_C^n.                    \tag{6.4}
\]

The minus sign is the one in d\_\allowbreak{}\{I\_\allowbreak{}1[1]\allowbreak{}\}. Define I\_\allowbreak{}2 by (6.3)\allowbreak{},\allowbreak{} but now use the map

\[
 e^n:E^n\longrightarrow I_2^n,\qquad
 e^n=\begin{pmatrix}1&h^n\\0&f_3^n\end{pmatrix}.          \tag{6.5}
\]

The two possible upper-right entries in its chain-map equation are

\[
 (d_{I_2}^n e^n)_{12}=d_{I_1}^n h^n+(\delta')^n f_3^n,\qquad
 (e^{n+1}d_E^n)_{12}=\delta^n+h^{n+1}d_C^n.
\]

They are equal by (6.4)\allowbreak{}; the other three entries agree by identity and the chain equation for f\_\allowbreak{}3. Thus e is an actual chain map. Its restriction on I\_\allowbreak{}1 is the identity,\allowbreak{} and its map on the quotient is exactly the prescribed f\_\allowbreak{}3. The diagram

\[
\begin{array}{ccccccccc}
0&\longrightarrow&A&\longrightarrow&B&\longrightarrow&C&\longrightarrow&0\\
&&\downarrow f_1&&\downarrow&&\downarrow f_3&&\\
0&\longrightarrow&I_1&\longrightarrow&I_2&\longrightarrow&I_3&\longrightarrow&0
\end{array}                                                    \tag{6.6}
\]

therefore commutes strictly,\allowbreak{} its rows are exact,\allowbreak{} and the middle vertical map is a quasi-isomorphism. The lower sequence is termwise split. No term of I\_\allowbreak{}1 or I\_\allowbreak{}3 has been changed. A common lower bound on the original terms and the prescribed resolutions is retained.

For completeness the projective statement has an equally explicit construction. Prescribe f\_\allowbreak{}1:\allowbreak{}P\_\allowbreak{}1 \allowbreak{}-> A and f\_\allowbreak{}3:\allowbreak{}P\_\allowbreak{}3 \allowbreak{}-> C,\allowbreak{} and form E\allowbreak{}=B times\_\allowbreak{}C P\_\allowbreak{}3. The map E \allowbreak{}-> B is a quasi-isomorphism,\allowbreak{} and 0 \allowbreak{}-> A \allowbreak{}-> E \allowbreak{}-> P\_\allowbreak{}3 \allowbreak{}-> 0 is termwise split. Write d\_\allowbreak{}E\allowbreak{}=[[d\_\allowbreak{}A,\allowbreak{}delta]\allowbreak{},\allowbreak{}[0,\allowbreak{}d\_\allowbreak{}\{P\_\allowbreak{}3\}]\allowbreak{}]\allowbreak{},\allowbreak{} with delta:\allowbreak{}P\_\allowbreak{}3 \allowbreak{}-> A[1]\allowbreak{}. Projective lifting through f\_\allowbreak{}1[1]\allowbreak{} supplies delta':\allowbreak{}P\_\allowbreak{}3 \allowbreak{}-> P\_\allowbreak{}1[1]\allowbreak{} and h\textasciicircum{}n:\allowbreak{}P\_\allowbreak{}3\textasciicircum{}n \allowbreak{}-> A\^{}n such that

\[
 \delta^n-f_1^{n+1}(\delta')^n
              =-d_A^n h^n+h^{n+1}d_{P_3}^n.
\]

Set d\_\allowbreak{}\{P\_\allowbreak{}2\}\allowbreak{}=[[d\_\allowbreak{}\{P\_\allowbreak{}1\},\allowbreak{}delta']\allowbreak{},\allowbreak{}[0,\allowbreak{}d\_\allowbreak{}\{P\_\allowbreak{}3\}]\allowbreak{}]\allowbreak{} on P\_\allowbreak{}1 direct-sum P\_\allowbreak{}3 and use the map P\_\allowbreak{}2 \allowbreak{}-> E with matrix [[f\_\allowbreak{}1,\allowbreak{}h]\allowbreak{},\allowbreak{}[0,\allowbreak{}1]\allowbreak{}]\allowbreak{}. Its upper-right chain equation is d\_\allowbreak{}A h\allowbreak{}+delta\allowbreak{}=f\_\allowbreak{}1[1]\allowbreak{}delta'\allowbreak{}+h d\_\allowbreak{}\{P\_\allowbreak{}3\},\allowbreak{} exactly the displayed equation; the remaining entries are the original chain equations. Composing with E \allowbreak{}-> B gives a short exact sequence of projective resolutions with both outer resolutions fixed. Finite biproducts of projectives are projective,\allowbreak{} by lifting the two source components. Any common upper bound is retained. This also proves the projective receiving version without hiding the dual signs.

\subsection{7. Computing right derived functors and higher connecting maps}\label{proofs-6209-7089-computing-right-derived-functors-and-higher-connecting-maps}

For bounded-below injective I,\allowbreak{} let S\_\allowbreak{}I be the source\textquotesingle s index category I/\allowbreak{}Qis\textasciicircum{}\allowbreak{}+(A)\allowbreak{},\allowbreak{} formed in the homotopy category. Every object alpha:\allowbreak{}I \allowbreak{}-> K has a morphism to the identity object:\allowbreak{} a map beta:\allowbreak{}K \allowbreak{}-> I with beta alpha\allowbreak{}=1 in K(A)\allowbreak{},\allowbreak{} obtained in section 2. It is a quasi-isomorphism,\allowbreak{} since alpha is. This beta is unique in K(A)\allowbreak{},\allowbreak{} again by the isomorphism alpha\^{}* of section 2. Thus the identity index is terminal,\allowbreak{} and its singleton subcategory is cofinal. For any exact F:\allowbreak{}K\textasciicircum{}\allowbreak{}+(A)\allowbreak{} \allowbreak{}-> D the diagram of target values F(K)\allowbreak{} is consequently essentially constant with value F(I)\allowbreak{}. The canonical comparison at the identity is the identity of F(I)\allowbreak{}; this proves that I computes RF. The argument does not need enough injectives in A.

Taking I to be an injective object in degree zero,\allowbreak{} and taking the exact complex-level functor induced by an additive A \allowbreak{}-> B,\allowbreak{} proves that this object is right acyclic. The passage between bounded complex-level and abelian-category derived functors is the exact bounded comparison proved in section 2 of PROOFS\_\allowbreak{}5091\_\allowbreak{}5900.md. DERIVED-RECON-042 reconciles P02-E0047 and French DERIVED-039 with MC-STK-ERR-0850\textquotesingle s singular ``direct consequence'' at original 6962.

With enough injectives,\allowbreak{} every bounded-below complex has an injective resolution. Its computing object just proved gives existence of RF on every object of D\textasciicircum{}\allowbreak{}+(A)\allowbreak{},\allowbreak{} for any exact F from K\textasciicircum{}\allowbreak{}+(A)\allowbreak{}. For an additive A \allowbreak{}-> B this also follows from enough right acyclic objects. That latter route uses the repaired proposition-enough-acyclics:\allowbreak{} an acyclic complex is a zero object in the derived category,\allowbreak{} so RF is individually defined there,\allowbreak{} and the preceding Leray theorem proves its F-image acyclic. The counterexample to the old image induction and the noncircular replacement are fully proved in PROOFS\_\allowbreak{}5091\_\allowbreak{}5900.md,\allowbreak{} section 6. The present existence argument therefore does not reinstate the erroneous image-closure assertion.

The exactness of the everywhere defined derived functor,\allowbreak{} including its original shift comparison,\allowbreak{} was proved by the two-out-of-three construction in PROOFS\_\allowbreak{}4166\_\allowbreak{}5089.md. Compose it with the exact K\textasciicircum{}\allowbreak{}+(A)\allowbreak{} \allowbreak{}-> D\textasciicircum{}\allowbreak{}+(A)\allowbreak{} to obtain the stated exact functor on K\textasciicircum{}\allowbreak{}+(A)\allowbreak{}. For a short exact sequence of complexes,\allowbreak{} the canonical derived connecting triangle gives,\allowbreak{} after RF,\allowbreak{} a distinguished triangle; this is the asserted delta-functor on Comp\textasciicircum{}\allowbreak{}+(A)\allowbreak{}. Its restriction along A \allowbreak{}-> Comp\textasciicircum{}\allowbreak{}+(A)\allowbreak{},\allowbreak{} A mapped to A[0]\allowbreak{},\allowbreak{} is again a delta-functor. Alternatively,\allowbreak{} (6.6)\allowbreak{} computes that triangle by a termwise split sequence,\allowbreak{} whose F-image is still termwise split for any additive F.

For the source\textquotesingle s left exact F,\allowbreak{} the long exact cohomology sequence of this triangle is precisely the displayed sequence at 7057–\allowbreak{}7061. Negative R\^{}iF vanish by the lower-bound argument,\allowbreak{} and the canonical F \allowbreak{}-> R\^{}0F is an isomorphism by the kernel calculation in PROOFS\_\allowbreak{}5091\_\allowbreak{}5900.md,\allowbreak{} section 4. Positive R\^{}iF vanish on injectives because they are computed in degree zero. The actual connecting map R\textasciicircum{}iF(C)\allowbreak{} \allowbreak{}-> R\textasciicircum{}\{i\allowbreak{}+1\}F(A)\allowbreak{} is H\^{}i of RF of the canonical C \allowbreak{}-> A[1]\allowbreak{},\allowbreak{} followed by the inverse of the exact-functor shift comparison. This retains the original connecting-map sign.

Every A embeds into an injective,\allowbreak{} so all positive R\^{}iF are effaceable. The complete effacement-to-universality proof is in ../\allowbreak{}HOMOLOGY\_\allowbreak{}INTAKE\_\allowbreak{}20260930/\allowbreak{}PROOFS\_\allowbreak{}1901\_\allowbreak{}2765.md,\allowbreak{} lines 310–\allowbreak{}374,\allowbreak{} with its application to the present canonical connecting maps in PROOFS\_\allowbreak{}5091\_\allowbreak{}5900.md,\allowbreak{} section 4. Specifically R\textasciicircum{}1F(A)\allowbreak{} is the cokernel of F(I)\allowbreak{} \allowbreak{}-> F(I/\allowbreak{}A)\allowbreak{},\allowbreak{} and higher R\textasciicircum{}\{n\allowbreak{}+1\}F(A)\allowbreak{} identify with R\textasciicircum{}nF(I/\allowbreak{}A)\allowbreak{}; these maps construct and uniquely determine the extension of a degree-zero natural transformation. Comparing embeddings via their common injective direct sum gives independence,\allowbreak{} and naturality of the short exact sequences gives the delta-functor equations. All effacement hypotheses are thus realized by the actual injective embeddings,\allowbreak{} not assumed as unproved conditions.

\subsection{8. Receiving maps and propagation}\label{proofs-6209-7089-receiving-maps-and-propagation}

The correction NEW-024 preserves the original lifting and uniqueness theorems,\allowbreak{} (3.6)\allowbreak{},\allowbreak{} and their projective duals. Their cited uses therefore receive corrected proofs without changed theorem hypotheses. The following bounded checks identify the actual receiving maps.

In current Cohomology 2440–\allowbreak{}2464,\allowbreak{} G[0]\allowbreak{} \allowbreak{}-> J is an object injective resolution and hence a termwise monomorphism. The map G[0]\allowbreak{} \allowbreak{}-> f\_\allowbreak{}*I \allowbreak{}-> J\textquotesingle{} has the bounded-below injective target J\textquotesingle. Strict lifting gives the displayed beta:\allowbreak{}J \allowbreak{}-> J\textquotesingle. Formula (4.4)\allowbreak{} shows that two such strict choices are homotopic by a homotopy zero on G[0]\allowbreak{}. Applying the additive global-section functor preserves its two-term homotopy equation. The stated cohomology map is unchanged. In current Descent 3933–\allowbreak{}3956 the same construction,\allowbreak{} with source F[0]\allowbreak{} \allowbreak{}-> J|\_\allowbreak{}S and injective target I,\allowbreak{} gives the map J|\_\allowbreak{}S \allowbreak{}-> I used there. This check uses the exactness established in that receiving passage; it does not claim a new audit of its preceding vanishing theorems.

For the flat base-change constructions in current Cohomology 3368–\allowbreak{}3413 and Cohomology on Sites 2439–\allowbreak{}2484,\allowbreak{} the target (g')\allowbreak{}\_\allowbreak{}*J is a bounded-below complex of injectives. Indeed,\allowbreak{} for a monomorphism M \allowbreak{}-> N,\allowbreak{} exactness of (g')\allowbreak{}\textasciicircum{}* and injectivity of J\^{}n make Hom((g')\allowbreak{}\textasciicircum{}*N,\allowbreak{}J\textasciicircum{}n)\allowbreak{} \allowbreak{}-> Hom((g')\allowbreak{}\textasciicircum{}*M,\allowbreak{}J\textasciicircum{}n)\allowbreak{} epic; adjunction proves that (g')\allowbreak{}\_\allowbreak{}*J\textasciicircum{}n is injective. Sections 2–\allowbreak{}3 therefore construct the actual beta:\allowbreak{}I \allowbreak{}-> (g')\allowbreak{}\_\allowbreak{}*J and prove uniqueness up to homotopy. The map f\_\allowbreak{}*beta is carried by the original identification f\_\allowbreak{}*(g')\allowbreak{}\_\allowbreak{}* \allowbreak{}= g\_\allowbreak{}*(f')\allowbreak{}\_\allowbreak{}* and the adjunction to g\textasciicircum{}*f\_\allowbreak{}*I \allowbreak{}-> (f')\allowbreak{}\_\allowbreak{}*J. Additivity of f\_\allowbreak{}* preserves the homotopy equation,\allowbreak{} and the adjunction bijection in each degree preserves compositions with the differentials and addition. Thus the homotopy descends to the displayed base-change arrow. Comparing different injective resolutions by (3.6)\allowbreak{} and the same uniqueness argument gives the same derived arrow under their canonical resolution identifications. No new flatness or source comparison hypothesis is introduced.

UNDER-009 applies to these comparisons with the actual lower term bound of (g')\allowbreak{}\_\allowbreak{}*J:\allowbreak{} cohomological agreement in those degrees already suffices for the Hom comparison. This is recorded as a possible choice of comparison map justified by (4.3)\allowbreak{},\allowbreak{} while the source\textquotesingle s full quasi-isomorphisms remain in its proof.

For the Cech boundary comparison in current Cohomology 3150–\allowbreak{}3204,\allowbreak{} UNDER-010 allows any already chosen outer injective resolutions I\_\allowbreak{}1 and I\_\allowbreak{}3 to be retained in the displayed short exact sequence. The actual I\_\allowbreak{}2 of (6.3)\allowbreak{} and map (6.5)\allowbreak{} supply its middle term. Applying sections on every intersection of the original covering preserves the split sequence. Taking the stated products over intersection indices preserves this splitting as well,\allowbreak{} with the entire indexing family retained. The total-complex maps consequently form the same exact diagram. The connecting map is induced by the actual off-diagonal delta\textquotesingle{} from (6.3)\allowbreak{},\allowbreak{} with the sign verified in Homology 3606–\allowbreak{}3648. Naturality of the exact diagram under refinement proves the compatibility of the two boundary constructions as in the receiving proof. The cover and both outer resolutions are unchanged.

In current More Algebra 16590–\allowbreak{}16626,\allowbreak{} P is the bounded finite-free complex with P\textasciicircum{}i\allowbreak{}=0 for i<m\allowbreak{}+1,\allowbreak{} and E \allowbreak{}-> L is an isomorphism on cohomology above m and surjective in degree m. Represent the derived composite P \allowbreak{}-> K \allowbreak{}-> L by a chain map,\allowbreak{} using the projective Hom comparison. Section 5 applies with n\allowbreak{}=m\allowbreak{}+1 and gives the required gamma:\allowbreak{}P \allowbreak{}-> E. Its cone is the original complex of finite free modules,\allowbreak{} with terms P\textasciicircum{}\{i\allowbreak{}+1\} direct-sum E\^{}i and differential (-d\_\allowbreak{}P,\allowbreak{}0;gamma,\allowbreak{}d\_\allowbreak{}E)\allowbreak{}; it remains bounded.

For the cohomology bound on C(gamma)\allowbreak{} \allowbreak{}-> M,\allowbreak{} retain the two exact triangle sequences. In degree i>m\allowbreak{}+1,\allowbreak{} the adjacent maps for P \allowbreak{}-> K and E \allowbreak{}-> L are isomorphisms,\allowbreak{} so the cone map is an isomorphism. At i\allowbreak{}=m\allowbreak{}+1,\allowbreak{} P \allowbreak{}-> K is still surjective at the left endpoint and is an isomorphism in degree m\allowbreak{}+2,\allowbreak{} while E \allowbreak{}-> L is an isomorphism in both relevant degrees. A class in the kernel of the cone map has zero boundary in H\textasciicircum{}\{m\allowbreak{}+2\}(P)\allowbreak{},\allowbreak{} hence comes from H\textasciicircum{}\{m\allowbreak{}+1\}(E)\allowbreak{}; its image in H\textasciicircum{}\{m\allowbreak{}+1\}(L)\allowbreak{} comes from H\textasciicircum{}\{m\allowbreak{}+1\}(K)\allowbreak{},\allowbreak{} which lifts to H\textasciicircum{}\{m\allowbreak{}+1\}(P)\allowbreak{}. Subtracting that image makes the class zero,\allowbreak{} proving injectivity. Surjectivity follows by the same exact sequence,\allowbreak{} with the two next endpoint isomorphisms. At i\allowbreak{}=m,\allowbreak{} lift the boundary of a class of H\textasciicircum{}m(M)\allowbreak{} to H\textasciicircum{}\{m\allowbreak{}+1\}(P)\allowbreak{},\allowbreak{} using surjectivity there. Its image in H\textasciicircum{}\{m\allowbreak{}+1\}(E)\allowbreak{} is zero,\allowbreak{} since that group maps isomorphically to H\textasciicircum{}\{m\allowbreak{}+1\}(L)\allowbreak{}. Lift it to H\textasciicircum{}m(C(gamma)\allowbreak{})\allowbreak{}; the difference of its image from the original class comes from H\textasciicircum{}m(L)\allowbreak{},\allowbreak{} and can be removed using the epimorphism H\textasciicircum{}m(E)\allowbreak{} \allowbreak{}-> H\textasciicircum{}m(L)\allowbreak{}. Thus the cone map is surjective in degree m. This proves exactly the receiving pseudo-coherence estimate,\allowbreak{} including both endpoint distinctions.

These are dependency checks on the listed passages,\allowbreak{} not a review of the whole receiving chapters. No source operation in those chapters is needed for these results.
\clearpage
\section{\texorpdfstring{Cartan–\allowbreak{}Eilenberg resolutions,\allowbreak{} composition,\allowbreak{} and resolution functors}{   and resolution functors}}\label{proofs-7096-7578-cartaneilenberg-resolutions-composition-and-resolution-functors}

This is editorial review material,\allowbreak{} not a replacement exposition for the translated source. Source:\allowbreak{} the Stacks Project authors,\allowbreak{} Derived Categories,\allowbreak{} original commit a04446e5\allowbreak{}7ec1fbc2\allowbreak{}52a871af\allowbreak{}cec7752f\allowbreak{}b2807b14, lines 7096–\allowbreak{}7578. The frozen original and the current source have the identities recorded in SOURCE\_\allowbreak{}READING\_\allowbreak{}PART10.json. Original line numbers below refer to that original. The current Homology chapter supplies the commuting-square double-complex convention at 4207–\allowbreak{}4223 and the total differential at 4283–\allowbreak{}4300. Complete earlier arguments used here are retained in PROOFS\_\allowbreak{}6209\_\allowbreak{}7089.md,\allowbreak{} especially sections 2–\allowbreak{}7,\allowbreak{} and PROOFS\_\allowbreak{}5091\_\allowbreak{}5900.md with ADDENDUM\_\allowbreak{}01. No novelty is claimed.

\subsection{\texorpdfstring{1. Existence,\allowbreak{} the actual double-complex maps,\allowbreak{} and prescribed resolutions}{1.  the actual double-complex  and prescribed resolutions}}\label{proofs-7096-7578-existence-the-actual-double-complex-maps-and-prescribed-resolutions}

Keep the source\textquotesingle s complex K,\allowbreak{} with K\textasciicircum{}p\allowbreak{}=0 for p<n,\allowbreak{} and its actual objects Z\textasciicircum{}p\allowbreak{}=ker(d\_\allowbreak{}K\textasciicircum{}p)\allowbreak{},\allowbreak{} B\textasciicircum{}p\allowbreak{}=im(d\_\allowbreak{}K\textasciicircum{}\{p-1\})\allowbreak{},\allowbreak{} H\textasciicircum{}p\allowbreak{}=Z\textasciicircum{}p/\allowbreak{}B\textasciicircum{}p. Its two short exact sequences,\allowbreak{} in every degree p>\allowbreak{}=n,\allowbreak{} are

\[
 0\longrightarrow B^p\longrightarrow Z^p\longrightarrow H^p
       \longrightarrow0,\qquad
 0\longrightarrow Z^p\longrightarrow K^p\longrightarrow B^{p+1}
       \longrightarrow0.                                      \tag{1.1}
\]

Here B\textasciicircum{}n\allowbreak{}=0 and Z\textasciicircum{}n\allowbreak{}=H\textasciicircum{}n. The source starts by resolving Z\^{}n and applies the injective horseshoe construction successively along (1.1)\allowbreak{},\allowbreak{} retaining the first resolution at each step. This is the already proved construction in section 6 of PROOFS\_\allowbreak{}6209\_\allowbreak{}7089.md. It gives exact sequences of complexes

\[
 0\to J_B^{p,\bullet}\xrightarrow{b_p}J_Z^{p,\bullet}
       \xrightarrow{c_p}J_H^{p,\bullet}\to0,\qquad
 0\to J_Z^{p,\bullet}\xrightarrow{z_p}I^{p,\bullet}
       \xrightarrow{e_p}J_B^{p+1,\bullet}\to0.                   \tag{1.2}
\]

All the complexes in (1.2)\allowbreak{} have injective terms and vanish in vertical degrees q\textless0. Their augmentations form morphisms from the two source sequences. Put

\[
 d_1^{p,\bullet}=z_{p+1}b_{p+1}e_p,\qquad
 d_2^{p,q}=d_{I^{p,\bullet}}^q.                                \tag{1.3}
\]

These are the source\textquotesingle s maps,\allowbreak{} with their full factorization. Since e\_\allowbreak{}\{p\allowbreak{}+1\}z\_\allowbreak{}\{p\allowbreak{}+1\}\allowbreak{}=0,\allowbreak{} consecutive d\_\allowbreak{}1's compose to zero. Every factor in (1.3)\allowbreak{} is a vertical chain map,\allowbreak{} so d\_\allowbreak{}1d\_\allowbreak{}2\allowbreak{}=d\_\allowbreak{}2d\_\allowbreak{}1,\allowbreak{} with no inserted minus sign. Since e\_\allowbreak{}p is epic and z\_\allowbreak{}\{p\allowbreak{}+1\}b\_\allowbreak{}\{p\allowbreak{}+1\} is monic,\allowbreak{} its kernel is J\_\allowbreak{}Z\textasciicircum{}\{p,\allowbreak{}\textbackslash{}bullet\} and its image is J\_\allowbreak{}B\textasciicircum{}\{p\allowbreak{}+1,\allowbreak{}\textbackslash{}bullet\},\allowbreak{} with the actual inclusions from (1.2)\allowbreak{}. The horizontal cohomology in degree p is J\_\allowbreak{}Z\textasciicircum{}\{p,\allowbreak{}\textbackslash{}bullet\}/\allowbreak{}J\_\allowbreak{}B\textasciicircum{}\{p,\allowbreak{}\textbackslash{}bullet\}\allowbreak{}=J\_\allowbreak{}H\textasciicircum{}\{p,\allowbreak{}\textbackslash{}bullet\}. These identifications commute with d\_\allowbreak{}2. The augmentations therefore give every resolution required in original 7116–\allowbreak{}7122,\allowbreak{} including the image resolution indexed by p\allowbreak{}+1 rather than p in the B notation. Augmentation compatibility in (1.2)\allowbreak{} gives d\_\allowbreak{}1 epsilon\allowbreak{}=epsilon d\_\allowbreak{}K. Set every column with p\textless n equal to zero. This proves the stated existence theorem.

\textbf{DERIVED-UNDER-012:\allowbreak{} the boundary and cohomology resolutions may be prescribed simultaneously.} Fix injective resolutions J\_\allowbreak{}B\textasciicircum{}\{p,\allowbreak{}\textbackslash{}bullet\} of B\^{}p for every p>\allowbreak{}=n\allowbreak{}+1 and J\_\allowbreak{}H\textasciicircum{}\{p,\allowbreak{}\textbackslash{}bullet\} of H\^{}p for every p>\allowbreak{}=n,\allowbreak{} all vanishing for q\textless0. Set J\_\allowbreak{}B\textasciicircum{}\{p,\allowbreak{}\textbackslash{}bullet\}\allowbreak{}=0 for p<\allowbreak{}=n and J\_\allowbreak{}H\textasciicircum{}\{p,\allowbreak{}\textbackslash{}bullet\}\allowbreak{}=0 for p\textless n. Apply the both-prescribed-outer-resolution construction DERIVED-UNDER-010 first to the first sequence of (1.1)\allowbreak{},\allowbreak{} using the two given outer resolutions,\allowbreak{} to construct J\_\allowbreak{}Z\textasciicircum{}\{p,\allowbreak{}\textbackslash{}bullet\}. Apply it to the second sequence using that J\_\allowbreak{}Z\textasciicircum{}\{p,\allowbreak{}\textbackslash{}bullet\} and the already given J\_\allowbreak{}B\textasciicircum{}\{p\allowbreak{}+1,\allowbreak{}\textbackslash{}bullet\},\allowbreak{} to construct I\textasciicircum{}\{p,\allowbreak{}\textbackslash{}bullet\}. Thus both sequences (1.2)\allowbreak{} exist for each p with exactly the given outer complexes and augmentations. The proof of (1.3)\allowbreak{} then applies verbatim. In particular the resulting horizontal boundary and cohomology resolution complexes identify with the prescribed ones,\allowbreak{} not merely with unspecified quasi-isomorphic complexes.

The earlier horseshoe proof constructs a middle differential \[
 \begin{pmatrix}d_{J_1}&\delta'\\0&d_{J_3}\end{pmatrix}
\] and an augmentation through the actual matrix \(\left(\begin{smallmatrix}1&h\\0&f_3\end{smallmatrix}\right)\). The relation \(\delta'f_3-\delta=-d_{J_1}h+h\,d_C\) proves that augmentation is a chain map. We use that construction,\allowbreak{} including h; an arbitrary pair of chosen resolutions is not treated as if its connecting maps already agreed strictly. Both outer resolutions are zero below 0,\allowbreak{} so the constructed middle complexes are also zero below 0. Each step uses finite direct sums only. Once the prescribed resolutions exist,\allowbreak{} no further global enough-injectives hypothesis is needed for this version.

There is a necessary boundary restriction:\allowbreak{} a nonzero acyclic injective resolution of B\textasciicircum{}n\allowbreak{}=0 cannot be prescribed while requiring all columns p\textless n to be zero. The image of d\_\allowbreak{}1\textasciicircum{}\{n-1,\allowbreak{}\textbackslash{}bullet\} is then literally zero. Our statement explicitly sets that boundary complex to zero. It does not assert that arbitrary resolutions at that boundary work.

DERIVED-RECON-043 concerns only the article at 7113. Both reports P02-E0048 and DERIVED-040 describe the existing one-operation MC-STK-ERR-0851 correction from ``a i'' to ``an i''. The lower-bound condition itself is unchanged and valid.

\subsection{\texorpdfstring{2. Totalization,\allowbreak{} both spectral sequences,\allowbreak{} and the reported index error}{2.  both spectral  and the reported index error}}\label{proofs-7096-7578-totalization-both-spectral-sequences-and-the-reported-index-error}

Write T\allowbreak{}=Tot(I)\allowbreak{}. In total degree m the only possibly nonzero summands are I\textasciicircum{}\{p,\allowbreak{}m-p\} for n<\allowbreak{}=p<\allowbreak{}=m. Thus T is zero below n and has injective terms:\allowbreak{} a finite direct sum of injectives is injective,\allowbreak{} since extension of a map into it is extension of each of finitely many components. Its differential on I\textasciicircum{}\{p,\allowbreak{}q\} is exactly

\[
 d_T=d_1^{p,q}+(-1)^p d_2^{p,q}.                               \tag{2.1}
\]

The vertical augmentation equation d\_\allowbreak{}2\textasciicircum{}\{p,\allowbreak{}0\}epsilon\textasciicircum{}p\allowbreak{}=0 and the horizontal equation from section 1 prove that epsilon:\allowbreak{}K\allowbreak{}->T is a chain map. The first spectral sequence for I has E\_\allowbreak{}1\textasciicircum{}\{p,\allowbreak{}0\}\allowbreak{}=K\textasciicircum{}p,\allowbreak{} E\_\allowbreak{}1\textasciicircum{}\{p,\allowbreak{}q\}\allowbreak{}=0 for q!\allowbreak{}=0,\allowbreak{} and d\_\allowbreak{}1\allowbreak{}=d\_\allowbreak{}K. Finite diagonal support gives convergence. Consequently its E\_\allowbreak{}2 and E\_\allowbreak{}infinity terms are H\textasciicircum{}p(K)\allowbreak{} on q\allowbreak{}=0 and zero elsewhere. Apply the same construction to the double complex having K in row 0 and zero other rows. The augmentation is a map of double complexes and induces the identity under these E\_\allowbreak{}1 identifications. The associated cohomology filtrations have only one nonzero quotient in each total degree. Hence H\textasciicircum{}m(epsilon)\allowbreak{} itself is the resulting isomorphism. This verifies the map omitted at current Homology 6833–\allowbreak{}6835,\allowbreak{} rather than just an abstract isomorphism of objects. Therefore T is an injective resolution of K.

Let F:\allowbreak{}A\allowbreak{}->B be an additive functor of abelian categories. For the source theorem F is also left exact; we first keep that assumption and then track where it is used. The individual injective resolutions already constructed define RF at K\textasciicircum{}p[0]\allowbreak{},\allowbreak{} H\textasciicircum{}p(K)\allowbreak{}[0]\allowbreak{},\allowbreak{} and K,\allowbreak{} by the injective terminal-comparison argument proved in section 7 of the preceding note. No assumption that RF exists on every other object is needed. Finite additivity gives the canonical degreewise isomorphism F(Tot I)\allowbreak{}\allowbreak{}->Tot(F(I)\allowbreak{})\allowbreak{},\allowbreak{} preserving (2.1)\allowbreak{} and its sign. Put L\allowbreak{}=F(I)\allowbreak{}.

For the first filtration,\allowbreak{} current Homology 6704–\allowbreak{}6709 gives

\[
 {}'E_0^{p,q}=F(I^{p,q}),\quad {}'d_0=(-1)^p F(d_2),\quad
 {}'E_1^{p,q}=H^q(F(I^{p,\bullet}))=R^qF(K^p).                  \tag{2.2}
\]

Multiplication of a differential by -1 leaves its kernel and image unchanged; this identifies the displayed cohomology without losing the sign in the differential. Its d\_\allowbreak{}1 is the map R\textasciicircum{}qF(d\_\allowbreak{}K\textasciicircum{}p)\allowbreak{}. Indeed d\_\allowbreak{}1\textasciicircum{}\{p,\allowbreak{}\textbackslash{}bullet\} is a comparison of the injective resolutions of K\^{}p and K\textasciicircum{}\{p\allowbreak{}+1\} extending d\_\allowbreak{}K\textasciicircum{}p; the uniqueness on homotopy classes proved earlier identifies its induced derived map.

Fix a vertical degree p. In horizontal degree q the two exact sequences of objects are

\[
 0\to B_I^{q,p}\to Z_I^{q,p}\to H_I^{q,p}\to0,\qquad
 0\to Z_I^{q,p}\to I^{q,p}\to B_I^{q+1,p}\to0.                 \tag{2.3}
\]

They split because their first terms are injective. Applying F preserves those split sequences. Thus the kernel,\allowbreak{} image,\allowbreak{} and cohomology of F(d\_\allowbreak{}1)\allowbreak{} are respectively F(Z\_\allowbreak{}I)\allowbreak{},\allowbreak{} F(B\_\allowbreak{}I)\allowbreak{},\allowbreak{} and F(H\_\allowbreak{}I)\allowbreak{},\allowbreak{} via their actual induced inclusion and quotient maps. The cohomology identification is canonical:\allowbreak{} it is the map induced by F of the cycle inclusion and the quotient,\allowbreak{} whose kernel and image have just been proved; it does not depend on chosen splittings. Vertical chain maps preserve those inclusions and quotients. Hence

\[
 {}''E_1^{p,q}=F(H_I^q(I^{\bullet,p})),\qquad
 {}''d_1^{p,q}=(-1)^q F(d_{\,H_I^q(I^{\bullet,\bullet})}^{p}),
 \qquad
 {}''E_2^{p,q}=R^pF(H^q(K)).                                  \tag{2.4}
\]

This proves DERIVED-RECON-044 /\allowbreak{} P02-E0049. The two p\textquotesingle s at original 7248–\allowbreak{}7249 must be q:\allowbreak{} the resolution complex is H\_\allowbreak{}I\textasciicircum{}q(I)\allowbreak{},\allowbreak{} and the object it resolves is H\textasciicircum{}q(K)\allowbreak{}. The vertical cohomology degree p,\allowbreak{} the superscript p on E\_\allowbreak{}1,\allowbreak{} and the sign (-1)\allowbreak{}\textasciicircum{}q remain unchanged. The report\textquotesingle s phrase ``vertical homology degree'' and its final formula do not precisely describe both printed occurrences; the source and (2.4)\allowbreak{},\allowbreak{} not that wording,\allowbreak{} determine the two edits. No previous canonical correction covers either occurrence.

Both filtrations are finite in each degree and both spectral sequences converge to H\textasciicircum{}m(RF(K)\allowbreak{})\allowbreak{}. For a chosen column lower bound n,\allowbreak{} their supports are respectively p>\allowbreak{}=n,\allowbreak{}q>\allowbreak{}=0 and p>\allowbreak{}=0,\allowbreak{}q>\allowbreak{}=n. The differential has bidegree (r,\allowbreak{}1-r)\allowbreak{}. At a first-sequence term (p,\allowbreak{}q)\allowbreak{},\allowbreak{} an incoming differential is zero when r\textgreater p-n and an outgoing one is zero when r>q\allowbreak{}+1. At a second-sequence term they are zero when r\textgreater p and r>q-n\allowbreak{}+1 respectively. On total degree m>\allowbreak{}=n,\allowbreak{} every term has therefore stabilized by page r\allowbreak{}=m-n\allowbreak{}+2. Each resulting filtration has at most m-n\allowbreak{}+1 potentially nonzero graded pieces. All terms and the abutment vanish in total degrees m\textless n.

\textbf{DERIVED-UNDER-011,\allowbreak{} first part.} All arguments in this section used additivity,\allowbreak{} the given injective resolutions,\allowbreak{} and the split sequences (2.3)\allowbreak{},\allowbreak{} and never used left exactness of F. Thus (2.2)\allowbreak{}–\allowbreak{}(2.4)\allowbreak{},\allowbreak{} convergence,\allowbreak{} and the explicit bounds hold for every additive F with a given Cartan–\allowbreak{}Eilenberg resolution. In that statement R\textasciicircum{}0F(M)\allowbreak{} is not silently identified with F(M)\allowbreak{}.

For a strict example,\allowbreak{} fix a prime ell and let F(M)\allowbreak{}\allowbreak{}=M direct-sum (M/\allowbreak{}ell M)\allowbreak{} on abelian groups. It is additive. On multiplication by ell from Z to Z,\allowbreak{} its second component is zero on Z/\allowbreak{}ell Z,\allowbreak{} so F does not preserve that monomorphism and is not left exact. An injective group J is ell-divisible:\allowbreak{} extend the map ell Z\allowbreak{}->J sending ell to any given y across ell Z\allowbreak{}->Z. The image x of 1 satisfies ell x\allowbreak{}=y. Hence F(J)\allowbreak{}\allowbreak{}=J direct-sum 0. Applied to an injective resolution of M it gives that resolution direct-sum the zero complex. Its derived value has H\textasciicircum{}0\allowbreak{}=M direct-sum 0 and no positive cohomology,\allowbreak{} whereas F(M)\allowbreak{} can have a nonzero second summand. This preserves both summands and shows that the extension is strictly broader than the left-exact source statement.

\subsection{3. Functoriality with the exact page change}\label{proofs-7096-7578-functoriality-with-the-exact-page-change}

Original 7253–\allowbreak{}7262 points forward to the filtered construction. The current forward passage,\allowbreak{} 8835–\allowbreak{}8969,\allowbreak{} was read as a dependency,\allowbreak{} not as completed review of the intervening section. Here are the comparison maps that justify the use made here. The first spectral sequence is canonical starting at E\_\allowbreak{}1; the second is canonical starting at E\_\allowbreak{}2. An arbitrarily chosen injective resolution does not make its raw E\_\allowbreak{}0 terms canonical. The later source\textquotesingle s “r>\allowbreak{}=1” for the filtered construction makes this convention explicit.

Equip K first with S\textasciicircum{}aK\allowbreak{}=sigma\_\allowbreak{}\{>\allowbreak{}=a\}K. Give T the column filtration S\textasciicircum{}aT\allowbreak{}=Tot(I\textasciicircum{}\{p,\allowbreak{}q\}:\allowbreak{}p>\allowbreak{}=a)\allowbreak{}. Its associated graded in degree a is the vertical resolution I\textasciicircum{}\{a,\allowbreak{}\textbackslash{}bullet\}[-a]\allowbreak{},\allowbreak{} with vertical sign (-1)\allowbreak{}\textasciicircum{}a. The augmentation on that graded piece resolves K\textasciicircum{}a[-a]\allowbreak{}. Each T\^{}m has a finite split filtration whose graded objects are injective. Thus K\allowbreak{}->T is a filtered injective resolution for S.

For the second construction put G\textasciicircum{}sK\allowbreak{}=tau\_\allowbreak{}\{<\allowbreak{}=-s\}K,\allowbreak{} with its actual cycles in top degree. Put

\[
 G^sT=\operatorname{Tot}(\tau_{\le -s}^{\,\mathrm{horizontal}} I).
                                                                    \tag{3.1}
\]

This is a subcomplex:\allowbreak{} vertical differentials preserve horizontal cycles and the truncated horizontal differential vanishes on the last cycles. For k\allowbreak{}=-s,\allowbreak{} the associated graded horizontal complex has J\_\allowbreak{}B\textasciicircum{}\{k,\allowbreak{}\textbackslash{}bullet\} in horizontal degree k-1 and J\_\allowbreak{}Z\textasciicircum{}\{k,\allowbreak{}\textbackslash{}bullet\} in horizontal degree k,\allowbreak{} with the actual inclusion b\_\allowbreak{}k between them. Its total differential keeps the signs (-1)\allowbreak{}\textasciicircum{}\{k-1\} and (-1)\allowbreak{}\textasciicircum{}k in these two columns. Each associated graded object of T\^{}m is thus a finite direct sum of injectives. The induced augmentation from gr\_\allowbreak{}G\textasciicircum{}s K resolves its two terms B\^{}k and Z\^{}k by columns. The finite-diagonal comparison of section 2 proves it is a quasi-isomorphism on the total complexes. Moreover the quotient by the first column pair is J\_\allowbreak{}H\textasciicircum{}\{k,\allowbreak{}\textbackslash{}bullet\} in horizontal degree k:\allowbreak{} its kernel is the total complex of the identity on J\_\allowbreak{}B\textasciicircum{}\{k,\allowbreak{}\textbackslash{}bullet\},\allowbreak{} which is contractible. Explicitly,\allowbreak{} on the two copies in horizontal degrees k-1,\allowbreak{}k take the degree -1 map from the second copy to the first to be the identity and zero on the first; the vertical cross terms cancel by (-1)\allowbreak{}\textasciicircum{}\{k-1\}\allowbreak{}+(-1)\allowbreak{}\textasciicircum{}k\allowbreak{}=0. It follows that the graded total is quasi-isomorphic to J\_\allowbreak{}H\textasciicircum{}\{k,\allowbreak{}\textbackslash{}bullet\}[-k]\allowbreak{}. This also proves directly the asserted cohomology description of the corresponding graded source.

For completeness,\allowbreak{} the needed filtered comparison lemma follows from the same extension argument as the injective comparison already proved,\allowbreak{} with strictly exact filtered sequences. An object J with finite filtration and injective graded pieces is a split finite sum of its graded pieces:\allowbreak{} descending induction splits 0\allowbreak{}->F\textasciicircum{}\{a\allowbreak{}+1\}J\allowbreak{}->F\textasciicircum{}aJ\allowbreak{}->gr\textasciicircum{}a J\allowbreak{}->0 because F\textasciicircum{}\{a\allowbreak{}+1\}J is a finite sum of injectives. A filtered map A\allowbreak{}->J from a strict subobject A of B can be extended component by component. The component to the piece in filtration degree a factors through A/\allowbreak{}F\textasciicircum{}\{a\allowbreak{}+1\}A,\allowbreak{} which is a subobject of B/\allowbreak{}F\textasciicircum{}\{a\allowbreak{}+1\}B by strictness; use injectivity of that piece to extend it. The finitely many extensions assemble to a filtered extension.

An acyclic complex on associated graded objects is strictly acyclic with its induced cycle filtrations. To check this,\allowbreak{} lift a filtered cycle through the differential successively on the graded pieces; subtract the resulting boundaries and raise the filtration index at every step. The finite filtration in the degrees involved terminates the process. Its kernels and images are therefore the filtered cycles,\allowbreak{} and its cycle sequences are strict exact sequences. Given a chain map from such a complex to a bounded-below complex of filtered injectives,\allowbreak{} construct a null homotopy upwards from the first nonzero target degree:\allowbreak{} the residual map vanishes on cycles by the previously constructed homotopy equation,\allowbreak{} factors through the next cycles,\allowbreak{} and extends to the next entire term by the filtered extension property. This is the same equation f\textasciicircum{}m\allowbreak{}=d\_\allowbreak{}J\textasciicircum{}\{m-1\}h\textasciicircum{}m\allowbreak{}+h\textasciicircum{}\{m\allowbreak{}+1\}d\textasciicircum{}m proved in the preceding note,\allowbreak{} now with every map filtered. It proves that a filtered quasi-isomorphism induces a bijection on homotopy classes into these targets,\allowbreak{} by applying the Hom sequence to its strictly acyclic cone and its shift.

Consequently a map of source filtered complexes gives a unique homotopy class of comparison maps between their filtered injective resolutions. Identity and composition hold by that uniqueness. An additive F preserves their filtered homotopies:\allowbreak{} all term filtrations of the resolutions split,\allowbreak{} so their images define finite filtrations and F of every filtered component remains filtered. A filtered homotopy induces an ordinary homotopy on associated graded complexes and therefore equal maps on E\_\allowbreak{}1 and all later pages. This establishes choice independence and naturality for both filtered constructions,\allowbreak{} even when a given resolution is available without a global enough-injectives assumption.

Here is the exact relation of (3.1)\allowbreak{} to the second double-complex sequence,\allowbreak{} rather than just an agreement of named E\_\allowbreak{}2 objects. Let V\^{}aT be the vertical-column filtration,\allowbreak{} retaining summands I\textasciicircum{}\{p,\allowbreak{}q\} with q>\allowbreak{}=a. For a filtered complex (C,\allowbreak{}V)\allowbreak{} define

\[
 (\operatorname{Dec}V)^s C^m
   =V^{s+m}C^m\cap d^{-1}(V^{s+m+1}C^{m+1}).                  \tag{3.2}
\]

For C\allowbreak{}=T the first condition retains horizontal degrees p<\allowbreak{}=-s; the second requires that the component in horizontal degree -s be a horizontal cycle. All remaining components already satisfy it,\allowbreak{} including their vertical differentials. Hence Dec V\allowbreak{}=G with the actual subobjects (3.1)\allowbreak{}. After applying F the same equality holds:\allowbreak{} the required horizontal kernels are preserved by the split sequences (2.3)\allowbreak{}.

To verify pages and differentials,\allowbreak{} use \[
 Z_r^{a,m}(V)=V^a C^m\cap d^{-1}V^{a+r}C^{m+1},\qquad
 E_r^{a,m-a}(V)=
 \frac{Z_r^{a,m}(V)}
 {Z_{r-1}^{a+1,m}(V)+dZ_{r-1}^{a-r+1,m-1}(V)}
 \quad(r\ge1).                                               \tag{3.3}
\] Formula (3.2)\allowbreak{} and d\textasciicircum{}2\allowbreak{}=0 give Z\_\allowbreak{}r\textasciicircum{}\{s,\allowbreak{}m\}(Dec V)\allowbreak{}\allowbreak{}=Z\_\allowbreak{}\{r\allowbreak{}+1\}\textasciicircum{}\{s\allowbreak{}+m,\allowbreak{}m\}(V)\allowbreak{}. The two denominator terms become respectively Z\_\allowbreak{}r\textasciicircum{}\{s\allowbreak{}+m\allowbreak{}+1,\allowbreak{}m\}(V)\allowbreak{} and dZ\_\allowbreak{}r\textasciicircum{}\{s\allowbreak{}+m-r,\allowbreak{}m-1\}(V)\allowbreak{}. They are exactly the denominator for E\_\allowbreak{}\{r\allowbreak{}+1\}\textasciicircum{}\{s\allowbreak{}+m,\allowbreak{}-s\}(V)\allowbreak{}. The differential is induced by the same original d\_\allowbreak{}C on both quotients. Thus

\[
 E_r^{s,t}(G)= {}''E_{r+1}^{\,2s+t,-s},\qquad r\ge1,            \tag{3.4}
\]

compatibly with differentials and all maps. Its change of indices sends bidegree (r,\allowbreak{}1-r)\allowbreak{} to (r\allowbreak{}+1,\allowbreak{}-r)\allowbreak{},\allowbreak{} as required. At r\allowbreak{}=1,\allowbreak{} (3.4)\allowbreak{} gives precisely the source\textquotesingle s R\textasciicircum{}\{2s\allowbreak{}+t\}F(H\textasciicircum{}\{-s\}K)\allowbreak{}\allowbreak{}=\{\}''E\_\allowbreak{}2\textasciicircum{}\{2s\allowbreak{}+t,\allowbreak{}-s\}. This proves the promised functoriality from E\_\allowbreak{}1 and E\_\allowbreak{}2,\allowbreak{} respectively,\allowbreak{} without asserting it for arbitrary E\_\allowbreak{}0 choices. It also supplies the actual page-shifting comparison used in current Cohomology 7342–\allowbreak{}7367.

\subsection{\texorpdfstring{4. Composition:\allowbreak{} the evaluation object and the additive extension}{4.  the evaluation object and the additive extension}}\label{proofs-7096-7578-composition-the-evaluation-object-and-the-additive-extension}

Write F:\allowbreak{}A\allowbreak{}->B and G:\allowbreak{}B\allowbreak{}->C and assume A and B have enough injectives. Keep the canonical comparison t:\allowbreak{}R(GF)\allowbreak{}\allowbreak{}->RG RF constructed in section 1 of PROOFS\_\allowbreak{}5091\_\allowbreak{}5900.md. For X represented by A\^{}\textbackslash bullet and a resolution A\textasciicircum{}\textbackslash{}bullet\allowbreak{}->I\textasciicircum{}\textbackslash{}bullet,\allowbreak{} its component is exactly

\[
 t_X:\quad
 R(GF)(X)=G(F(I^\bullet))
       \xrightarrow{\eta_{G,F(I^\bullet)}}RG(F(I^\bullet))
       =RG(RF(X)).                                            \tag{4.1}
\]

The displayed equalities are the specified injective-resolution identifications,\allowbreak{} not an assertion that a nonexact F preserves arbitrary quasi-isomorphisms.

If every F(I)\allowbreak{},\allowbreak{} for I injective in A,\allowbreak{} is right G-acyclic,\allowbreak{} all terms of F(I\textasciicircum{}\textbackslash{}bullet)\allowbreak{} are G-acyclic. The already proved Leray acyclicity lemma makes the middle arrow of (4.1)\allowbreak{} an isomorphism. Conversely evaluate t at the object I[0]\allowbreak{} of D\textasciicircum{}\allowbreak{}+(A)\allowbreak{}. Its two sides identify with G(F(I)\allowbreak{})\allowbreak{}[0]\allowbreak{} and RG(F(I)\allowbreak{}[0]\allowbreak{})\allowbreak{},\allowbreak{} and its component is their canonical comparison. Thus an isomorphism t makes F(I)\allowbreak{} right G-acyclic.

\textbf{DERIVED-NEW-026.} Original 7311 instead says to evaluate t on RF(I)\allowbreak{}\allowbreak{}=F(I)\allowbreak{}. That is an object of D\textasciicircum{}\allowbreak{}+(B)\allowbreak{},\allowbreak{} whereas t has domain D\textasciicircum{}\allowbreak{}+(A)\allowbreak{}; for unrelated A and B this is not an allowable argument of t. Replace that clause by evaluation at I[0]\allowbreak{},\allowbreak{} using RF(I[0]\allowbreak{})\allowbreak{}\allowbreak{}=F(I)\allowbreak{}[0]\allowbreak{}. The argument above proves the conclusion with exactly that component. The theorem is valid; it is the evaluation instruction in the proof that is ill-typed.

\textbf{DERIVED-UNDER-011,\allowbreak{} composition part.} This equivalence also holds for additive F and G without left exactness. Enough injectives defines all three right derived functors; (4.1)\allowbreak{} and Leray acyclicity use additivity and the actual right-acyclic comparison. For this extension ``right G-acyclic'' means that G(M)\allowbreak{}[0]\allowbreak{}\allowbreak{}->RG(M[0]\allowbreak{})\allowbreak{} itself is an isomorphism. Vanishing R\textasciicircum{}iG(M)\allowbreak{} for i\textgreater0 alone is insufficient:\allowbreak{} it omits the degree-zero comparison,\allowbreak{} as the example in section 2 illustrates.

Under these equivalent conditions choose a Cartan–\allowbreak{}Eilenberg resolution of F(I\textasciicircum{}\textbackslash{}bullet)\allowbreak{},\allowbreak{} and apply section 2 with functor G. Its second sequence has \[
 E_2^{p,q}=R^pG(H^q(F(I^\bullet)))=R^pG(H^q(RF(X))).
                                                               \tag{4.2}
\] Its abutment is H\textasciicircum{}\{p\allowbreak{}+q\}(RG(RF(X)\allowbreak{})\allowbreak{})\allowbreak{},\allowbreak{} identified by H(t\_\allowbreak{}X)\allowbreak{} with H\textasciicircum{}\{p\allowbreak{}+q\}(R(GF)\allowbreak{}(X)\allowbreak{})\allowbreak{}. This proves the stated Grothendieck sequence with the exact comparison map. If I\^{}\textbackslash bullet is zero below n,\allowbreak{} F(I\textasciicircum{}\textbackslash{}bullet)\allowbreak{} is zero below n; take the Cartan–\allowbreak{}Eilenberg resolution with the same horizontal bound. Section 2 supplies page m-n\allowbreak{}+2 stabilization and at most m-n\allowbreak{}+1 graded pieces in each degree m>\allowbreak{}=n. Naturality from E\_\allowbreak{}2 follows from section 3 and (4.1)\allowbreak{}.

The report dispositions in this interval are specific:\allowbreak{} RECON-045 /\allowbreak{} P02-E0050 /\allowbreak{} DERIVED-041 is existing MC-STK-ERR-0852 (``be'' to ``are'' at 7278)\allowbreak{}. RECON-046 /\allowbreak{} P02-E0051 /\allowbreak{} DERIVED-042 is existing MC-STK-ERR-0853,\allowbreak{} both the premature period and missing article. RECON-047 /\allowbreak{} P02-E0052 /\allowbreak{} DERIVED-043 is existing MC-STK-ERR-0854,\allowbreak{} the extra parenthesis in (4.1)\allowbreak{}. RECON-048 /\allowbreak{} P02-E0053 /\allowbreak{} DERIVED-044 is existing MC-STK-ERR-0855. P02 proposes joining the fragment to the next sentence,\allowbreak{} whereas the existing two-operation edit makes it a complete assumption sentence. Both address the same defect; the existing edit retains exactly the original mathematical hypotheses and is not replaced solely to match P02\textquotesingle s wording.

\subsection{5. Finite-prefix propagation of DERIVED-UNDER-007}\label{proofs-7096-7578-finite-prefix-propagation-of-derived-under-007}

The stronger finite-prefix comparison already proved in PART07 has a precise receiver in (4.1)\allowbreak{}. Let F and G be additive,\allowbreak{} A and B have enough injectives,\allowbreak{} and fix X\allowbreak{}->I\textasciicircum{}\textbackslash{}bullet. It suffices here that the particular F(I\textasciicircum{}j)\allowbreak{} are right G-acyclic for j<\allowbreak{}=m; no assertion about all injectives is needed. Applying DERIVED-UNDER-007 to the actual complex F(I\textasciicircum{}\textbackslash{}bullet)\allowbreak{} and functor G gives \[
 H^j(t_X)\text{ is an isomorphism for }j<m,\qquad
 H^m(t_X)\text{ is monic}.                                    \tag{5.1}
\] Its boundary defect is exactly the intersection from the previous proof:\allowbreak{} with A\allowbreak{}=F(I\textasciicircum{}\textbackslash{}bullet)\allowbreak{},\allowbreak{} P\allowbreak{}=sigma\_\allowbreak{}\{<\allowbreak{}=m\}A and T\allowbreak{}=sigma\_\allowbreak{}\{>\allowbreak{}=m\allowbreak{}+1\}A,\allowbreak{} the cokernel of H\textasciicircum{}m(t\_\allowbreak{}X)\allowbreak{} is identified with im(delta\_\allowbreak{}G)\allowbreak{} intersect ker H\textasciicircum{}\{m\allowbreak{}+1\}(eta\_\allowbreak{}\{G,\allowbreak{}T\})\allowbreak{},\allowbreak{} where delta\_\allowbreak{}G:\allowbreak{}H\textasciicircum{}m(G(P)\allowbreak{})\allowbreak{}\allowbreak{}->H\textasciicircum{}\{m\allowbreak{}+1\}(G(T)\allowbreak{})\allowbreak{} is the boundary of the original termwise split triangle T\allowbreak{}->A\allowbreak{}->P\allowbreak{}->T[1]\allowbreak{}. This follows by the same comparison of its two exact cohomology sequences; every object and map is now substituted from (4.1)\allowbreak{}.

If G is also left exact,\allowbreak{} the complete PART07 addendum improves (5.1)\allowbreak{} to \[
 H^j(t_X)\text{ is an isomorphism for }j\le m+1,\qquad
 H^{m+2}(t_X)\text{ is monic}.                                \tag{5.2}
\] For clarity,\allowbreak{} the extra degrees come from the bottom of T:\allowbreak{} left exactness makes H\textasciicircum{}\{m\allowbreak{}+1\}(G(T)\allowbreak{})\allowbreak{}\allowbreak{}->H\textasciicircum{}\{m\allowbreak{}+1\}(RG(T)\allowbreak{})\allowbreak{} an isomorphism and the next comparison monic. Comparing cokernels of H\textasciicircum{}m(G(P)\allowbreak{})\allowbreak{}\allowbreak{}->H\textasciicircum{}\{m\allowbreak{}+1\}(G(T)\allowbreak{})\allowbreak{} proves the degree m\allowbreak{}+1 assertion,\allowbreak{} since P computes RG(P)\allowbreak{}. The adjacent cohomology objects of P vanish in the next degree,\allowbreak{} proving the last assertion. This retains all of A,\allowbreak{} including its terms above m. It is not a claim that the truncated term complex alone computes every displayed cohomology object.

The Grothendieck sequence without the all-terms acyclicity hypothesis still abuts to RG(RF(X)\allowbreak{})\allowbreak{}; (5.1)\allowbreak{} or (5.2)\allowbreak{} identifies only the proved degree range with R(GF)\allowbreak{}(X)\allowbreak{}. No global abutment identification is asserted from a finite prefix.

\subsection{6. The equivalence and the meaning of uniqueness}\label{proofs-7096-7578-the-equivalence-and-the-meaning-of-uniqueness}

The full subcategory of injectives is additive and closed under finite sums. Shifts and mapping cones of bounded-below complexes in it stay there. The canonical functor q:\allowbreak{}K\textasciicircum{}\allowbreak{}+(I)\allowbreak{}\allowbreak{}->D\textasciicircum{}\allowbreak{}+(A)\allowbreak{} therefore preserves distinguished cone triangles and shifts. Every object of D\textasciicircum{}\allowbreak{}+(A)\allowbreak{} has a bounded below representative and an injective resolution,\allowbreak{} giving essential surjectivity. For complexes U,\allowbreak{}V in K\textasciicircum{}\allowbreak{}+(I)\allowbreak{},\allowbreak{} the injective Hom comparison proves Hom\_\allowbreak{}\{K(A)\allowbreak{}\}(U,\allowbreak{}V)\allowbreak{}\allowbreak{}=Hom\_\allowbreak{}\{D(A)\allowbreak{}\}(U,\allowbreak{}V)\allowbreak{}. The bounded homotopy and derived subcategories are full on these objects,\allowbreak{} so this is the required full faithfulness. This proves original 7377–\allowbreak{}7397 with its actual functor.

RECON-049 /\allowbreak{} DERIVED-046 is the existing MC-STK-ERR-0857 ``Full faithfulness'' correction at 7396. This also closes the second locus of P02-E0039 /\allowbreak{} OCC-01511. Its first locus 6197–\allowbreak{}6198 was proved in section 4 of PROOFS\_\allowbreak{}5910\_\allowbreak{}6199.md and linked to RECON-035 /\allowbreak{} MC-STK-ERR-0840. The physical P02 report is counted once,\allowbreak{} in RECON-049,\allowbreak{} with an explicit link to that earlier component. The two French reports remain their distinct physical records.

A resolution functor was defined at 7407–\allowbreak{}7418 as the pair of choices (jK,\allowbreak{}i\_\allowbreak{}K:\allowbreak{}K\allowbreak{}->jK)\allowbreak{},\allowbreak{} not as j without augmentation. Precomposition with i\_\allowbreak{}K gives a bijection \[
 \operatorname{Hom}_{K(A)}(jK,jL)
       \longrightarrow\operatorname{Hom}_{K(A)}(K,jL).
                                                               \tag{6.1}
\] Let j(alpha)\allowbreak{} be the unique preimage of i\_\allowbreak{}L alpha. For the identity it is the identity; for beta alpha it is j(beta)\allowbreak{}j(alpha)\allowbreak{},\allowbreak{} since that composite has the required precomposition. Additivity follows from the same uniqueness and the additive map in (6.1)\allowbreak{}. Thus j is an additive functor and i a natural transformation in K\textasciicircum{}\allowbreak{}+(A)\allowbreak{}. Applying localization gives a natural isomorphism since each i\_\allowbreak{}K is a quasi-isomorphism.

For another pair (j',\allowbreak{}i')\allowbreak{},\allowbreak{} the same bijection gives a unique a\_\allowbreak{}K:\allowbreak{}jK\allowbreak{}->j'K with a\_\allowbreak{}K i\_\allowbreak{}K\allowbreak{}=i'\_\allowbreak{}K in K\textasciicircum{}\allowbreak{}+(A)\allowbreak{}. Construct b\_\allowbreak{}K in the opposite direction. The composites b\_\allowbreak{}Ka\_\allowbreak{}K and a\_\allowbreak{}Kb\_\allowbreak{}K are identities by (6.1)\allowbreak{}. Naturality is proved by precomposing the two candidate maps jK\allowbreak{}->j'L with i\_\allowbreak{}K. This proves existence and uniqueness of the augmentation-compatible natural isomorphism a. The existence statement for a set of objects follows by choosing one injective resolution for each such object.

The unqualified phrase at 7463–\allowbreak{}7464 is read using the defined structured object (j,\allowbreak{}i)\allowbreak{}. Without compatibility with i,\allowbreak{} uniqueness among all natural isomorphisms would be false. For example,\allowbreak{} over Q-vector spaces every object is injective and the identity functor on K\textasciicircum{}\allowbreak{}+ has distinct natural automorphisms 1 and -1; they differ on Q[0]\allowbreak{}. Only 1 is compatible with the identity augmentation. The source\textquotesingle s definition and its explicit comparison at 7571–\allowbreak{}7576 already provide this compatibility,\allowbreak{} so no additional source error is counted. This records the adverse interpretation and its exact counterexample without turning shorthand into an unsupported new defect.

RECON-050 /\allowbreak{} DERIVED-047 is existing MC-STK-ERR-0858,\allowbreak{} ``resolution functors'' at 7570. RECON-051 /\allowbreak{} P02-E0054 /\allowbreak{} DERIVED-045 is existing MC-STK-ERR-0856:\allowbreak{} with i:\allowbreak{}K\allowbreak{}->I,\allowbreak{} i':\allowbreak{}K\allowbreak{}->J and a:\allowbreak{}I\allowbreak{}->J,\allowbreak{} the equality is i'\allowbreak{}=a i in K\textasciicircum{}\allowbreak{}+(A)\allowbreak{}. The arrow a lies in K\textasciicircum{}\allowbreak{}+(I)\allowbreak{},\allowbreak{} but K need not. Both the composition-order and ambient-category repairs are needed and retained,\allowbreak{} even though P02 reports only the first.

\subsection{\texorpdfstring{7. Exactness,\allowbreak{} shifts,\allowbreak{} and the quasi-inverse}{7.   and the quasi-inverse}}\label{proofs-7096-7578-exactness-shifts-and-the-quasi-inverse}

There is a unique xi\_\allowbreak{}K:\allowbreak{}j(K[1]\allowbreak{})\allowbreak{}\allowbreak{}->j(K)\allowbreak{}[1]\allowbreak{} with \[
 \xi_K i_{K[1]}=i_K[1].
                                                               \tag{7.1}
\] Both arrows on the right and on the left\textquotesingle s domain are resolutions of K[1]\allowbreak{}. The comparison in section 6 makes xi\_\allowbreak{}K invertible. For alpha:\allowbreak{}K\allowbreak{}->L,\allowbreak{} both xi\_\allowbreak{}L j(alpha[1]\allowbreak{})\allowbreak{} and j(alpha)\allowbreak{}[1]\allowbreak{}xi\_\allowbreak{}K become i\_\allowbreak{}L[1]\allowbreak{}alpha[1]\allowbreak{} after precomposition with i\_\allowbreak{}\{K[1]\allowbreak{}\}; the same uniqueness proves their equality. This proves naturality,\allowbreak{} including the source\textquotesingle s shift convention d\_\allowbreak{}\{K[1]\allowbreak{}\}\textasciicircum{}m\allowbreak{}=-d\_\allowbreak{}K\textasciicircum{}\{m\allowbreak{}+1\}.

For a distinguished triangle (K,\allowbreak{}L,\allowbreak{}M,\allowbreak{}f,\allowbreak{}g,\allowbreak{}h)\allowbreak{},\allowbreak{} the triangle of chosen resolutions has final arrow xi\_\allowbreak{}K j(h)\allowbreak{}. Its image under q is isomorphic to the original localized triangle via the three i\textquotesingle s and (7.1)\allowbreak{}. It is consequently distinguished in D\textasciicircum{}\allowbreak{}+(A)\allowbreak{}. To see directly that q reflects this property,\allowbreak{} complete the first arrow in K\textasciicircum{}\allowbreak{}+(I)\allowbreak{} to a cone triangle. In D\textasciicircum{}\allowbreak{}+(A)\allowbreak{} the cone triangle and the proposed triangle have identical first arrows,\allowbreak{} and the triangle axiom gives a map on their third objects. It is an isomorphism:\allowbreak{} applying each representable cohomological functor to the two long exact sequences gives this by the five lemma. Full faithfulness of q lifts that map and its inverse to K\textasciicircum{}\allowbreak{}+(I)\allowbreak{},\allowbreak{} and lifts its commutative triangle equalities. Thus the proposed triangle is isomorphic to a cone triangle in K\textasciicircum{}\allowbreak{}+(I)\allowbreak{},\allowbreak{} proving that (j,\allowbreak{}xi)\allowbreak{} is exact.

\begingroup\raggedright
For a quasi-isomorphism alpha,\allowbreak{} naturality of i shows j(alpha)\allowbreak{}i\_\allowbreak{}K\allowbreak{}=i\_\allowbreak{}L alpha is a quasi-isomorphism. Since i\_\allowbreak{}K is one too,\allowbreak{} j(alpha)\allowbreak{} is a quasi-isomorphism. Between injective complexes it is invertible in K\textasciicircum{}\allowbreak{}+(I)\allowbreak{}:\allowbreak{} its image in D\textasciicircum{}\allowbreak{}+(A)\allowbreak{} is invertible and q is fully faithful. The localization universal property therefore factors j through Q:\allowbreak{}K\textasciicircum{}\allowbreak{}+(A)\allowbreak{}\allowbreak{}->D\textasciicircum{}\allowbreak{}+(A)\allowbreak{} as j'Q,\allowbreak{} uniquely with this specified equality. For U in K\textasciicircum{}\allowbreak{}+(I)\allowbreak{},\allowbreak{} i\_\allowbreak{}U:\allowbreak{}U\allowbreak{}->jU is itself a morphism in K\textasciicircum{}\allowbreak{}+(I)\allowbreak{} and is invertible; these maps give Id\allowbreak{}->j'q. For a representative K,\allowbreak{} Q(i\_\allowbreak{}K)\allowbreak{} gives QK\allowbreak{}->qjK. This transformation descends through Q:\allowbreak{} naturality holds on original maps by section 6 and on inverted quasi-isomorphisms by inverting that same naturality equation. It gives Id\allowbreak{}->qj' on D\textasciicircum{}\allowbreak{}+(A)\allowbreak{}. Both are natural isomorphisms,\allowbreak{} proving the quasi-inverse assertion omitted at 7560–\allowbreak{}7561.
\par\endgroup

The reference to a square in 7554–\allowbreak{}7556 uses the naturality square just proved; although its explicit alpha-labelled display occurs in the earlier resolution-functor lemma,\allowbreak{} that relation is also used in the exactness proof. This does not require a new mathematical correction. For a category with a proper class of objects the comparison proof still works for given choices; set-sized choice alone does not supply a class of such choices. The source explicitly defers that existence issue to its next section. This review does not assert that the issue has already been resolved.

\subsection{8. Receiving statements and propagation limits}\label{proofs-7096-7578-receiving-statements-and-propagation-limits}

Current Cohomology 2140–\allowbreak{}2193 applies composition to the actual pair f\_\allowbreak{}* and Gamma(Y,\allowbreak{}-)\allowbreak{},\allowbreak{} retaining the restriction-of-scalars functor along O\_\allowbreak{}Y(Y)\allowbreak{}\allowbreak{}->O\_\allowbreak{}X(X)\allowbreak{}. On a chosen injective resolution I\^{}\textbackslash bullet the comparison is \[
 \Gamma(Y,f_*I^\bullet)
     \longrightarrow R\Gamma(Y,f_*I^\bullet).
                                                               \tag{8.1}
\] Its left side is Gamma(X,\allowbreak{}I\textasciicircum{}\textbackslash{}bullet)\allowbreak{} with that specified scalar restriction. Section 4 evaluates acyclicity on I[0]\allowbreak{} in the category of O\_\allowbreak{}X-modules; it does not try to put f\_\allowbreak{}*I into that category. The receiver\textquotesingle s injective acyclicity hypothesis is exactly the one used in (4.1)\allowbreak{}. Its Leray sequence at 2261–\allowbreak{}2283 therefore retains E\_\allowbreak{}2\textasciicircum{}\{p,\allowbreak{}q\}\allowbreak{}=H\textasciicircum{}p(Y,\allowbreak{}R\textasciicircum{}qf\_\allowbreak{}*\textbackslash{}mathcal F\textasciicircum{}\textbackslash{}bullet)\allowbreak{} and its stated abutment via (8.1)\allowbreak{}. The correction of 7248–\allowbreak{}7249 changes no E\_\allowbreak{}2 statement here:\allowbreak{} it repairs the resolving complex in the proof.

For a chosen I\^{}\textbackslash bullet zero below n the spectral sequence has at most m-n\allowbreak{}+1 graded pieces in degree m and stabilizes there by page m-n\allowbreak{}+2. If only the prefix f\_\allowbreak{}*I\textasciicircum{}j,\allowbreak{} j<\allowbreak{}=m,\allowbreak{} is known Gamma(Y,\allowbreak{}-)\allowbreak{}-acyclic,\allowbreak{} (5.2)\allowbreak{} still proves that (8.1)\allowbreak{} is an isomorphism on cohomology through degree m\allowbreak{}+1 and monic in degree m\allowbreak{}+2. The existing receiver supplies acyclicity of every term,\allowbreak{} so this is an editorial strengthening,\allowbreak{} not a change to its hypotheses or statement.

Current Sites Cohomology 2304–\allowbreak{}2344 has the same exact receiver with Gamma(D,\allowbreak{}-)\allowbreak{},\allowbreak{} f\_\allowbreak{}* and Gamma(C,\allowbreak{}-)\allowbreak{}. Its cited total acyclicity of f\_\allowbreak{}*I supplies the full right-acyclic condition. Substitution into (4.1)\allowbreak{},\allowbreak{} (4.2)\allowbreak{},\allowbreak{} and (5.2)\allowbreak{} proves the identical comparison,\allowbreak{} page bounds,\allowbreak{} and finite-prefix refinement for this receiver,\allowbreak{} retaining its ringed topoi and module categories.

Current Cohomology 7342–\allowbreak{}7367 explicitly uses the reindexing p\allowbreak{}=-j,\allowbreak{}q\allowbreak{}=i\allowbreak{}+2j for the cohomology-of-cohomology sequence. Formula (3.4)\allowbreak{} supplies the page change and actual differential as well as that reindexing. Its first sequence at 7370–\allowbreak{}7391 uses the filtration S,\allowbreak{} whose graded injective resolution was checked in section 3. Thus the two identifications are compatible with maps from their canonical pages. The bounded-below cases here are proved; this note does not infer unbounded convergence from those bounds.

A root-chapter reference search also found other users of the two source lemmas. Those hits are dependency routes,\allowbreak{} not claims that their whole proofs were read. The source statements of the corrected lemmas are unchanged. The precise receiver calculations above and the forward filtered comparison are the propagation completed in this part. Joint research synthesis remains later than the correction and received-theorem queue.
\clearpage
\section{Functorial resolutions and finite filtered injectives}\label{proofs-7589-8350-functorial-resolutions-and-finite-filtered-injectives}

Editorial review of the Stacks Project authors\textquotesingle{} Derived Categories,\allowbreak{} original commit a04446e5\allowbreak{}7ec1fbc2\allowbreak{}52a871af\allowbreak{}cec7752f\allowbreak{}b2807b14, lines 7589–\allowbreak{}8350. The sources and actual reading ranges are bound in SOURCE\_\allowbreak{}READING\_\allowbreak{}PART11.json. The original formulas and objects remain the comparison authority. The stronger statements below are editorial consequences,\allowbreak{} with no novelty claim or replacement of source exposition. Earlier complete proofs are retained in PROOFS\_\allowbreak{}6209\_\allowbreak{}7089.md,\allowbreak{} PROOFS\_\allowbreak{}7096\_\allowbreak{}7578.md,\allowbreak{} and the earlier bounded-filtered-category review.

\subsection{\texorpdfstring{1. The embedding functor,\allowbreak{} its zero summand,\allowbreak{} and every differential}{1. The embedding  its zero  and every differential}}\label{proofs-7589-8350-the-embedding-functor-its-zero-summand-and-every-differential}

Current Homology 7232–\allowbreak{}7248 defines a functorial injective embedding as a functor to the arrow category. Equivalently it gives an ordinary functor J:\allowbreak{}A\allowbreak{}->A and natural monomorphisms eta\_\allowbreak{}A:\allowbreak{}A\allowbreak{}->J(A)\allowbreak{},\allowbreak{} with each J(A)\allowbreak{} injective. It does not require J to be additive.

Keep the original J and its zero-object contribution. Write s\_\allowbreak{}A\allowbreak{}=J(0\allowbreak{}->A)\allowbreak{}:\allowbreak{}J(0)\allowbreak{}\allowbreak{}->J(A)\allowbreak{} and p\_\allowbreak{}A\allowbreak{}=J(A\allowbreak{}->0)\allowbreak{}:\allowbreak{}J(A)\allowbreak{}\allowbreak{}->J(0)\allowbreak{}. Their composite p\_\allowbreak{}A s\_\allowbreak{}A is 1. Put J\_\allowbreak{}red(A)\allowbreak{}\allowbreak{}=ker p\_\allowbreak{}A,\allowbreak{} with inclusion k\_\allowbreak{}A. There is a unique r\_\allowbreak{}A:\allowbreak{}J(A)\allowbreak{}\allowbreak{}->J\_\allowbreak{}red(A)\allowbreak{} satisfying k\_\allowbreak{}A r\_\allowbreak{}A\allowbreak{}=1-s\_\allowbreak{}Ap\_\allowbreak{}A. The two inverse isomorphisms are

\[
 (s_A,k_A):J(0)\oplus J_{\rm red}(A)\longrightarrow J(A),
 \qquad
 \binom{p_A}{r_A}:J(A)\longrightarrow J(0)\oplus J_{\rm red}(A).
                                                               \tag{1.1}
\]

For f:\allowbreak{}A\allowbreak{}->B,\allowbreak{} p\_\allowbreak{}B J(f)\allowbreak{}\allowbreak{}=p\_\allowbreak{}A and J(f)\allowbreak{}s\_\allowbreak{}A\allowbreak{}=s\_\allowbreak{}B. Thus J(f)\allowbreak{} restricts to a map J\_\allowbreak{}red(f)\allowbreak{},\allowbreak{} and (1.1)\allowbreak{} identifies it with the matrix diag(1\_\allowbreak{}\{J(0)\allowbreak{}\},\allowbreak{}J\_\allowbreak{}red(f)\allowbreak{})\allowbreak{}. These restrictions preserve identity and composition. J\_\allowbreak{}red(0)\allowbreak{}\allowbreak{}=ker(1\_\allowbreak{}\{J(0)\allowbreak{}\})\allowbreak{}\allowbreak{}=0. Its values are injective direct summands:\allowbreak{} extend a map into J\_\allowbreak{}red(A)\allowbreak{} into J(A)\allowbreak{} and then compose with r\_\allowbreak{}A.

Naturality gives p\_\allowbreak{}A eta\_\allowbreak{}A\allowbreak{}=eta\_\allowbreak{}0(A\allowbreak{}->0)\allowbreak{}\allowbreak{}=0. Therefore eta\_\allowbreak{}A factors uniquely as k\_\allowbreak{}A eta\_\allowbreak{}red,\allowbreak{}A; the factor is monic. The factorization and its naturality are retained. We use J\_\allowbreak{}red for the source\textquotesingle s explicitly described replacement,\allowbreak{} without claiming that the original J(0)\allowbreak{} was absent. Because a zero morphism factors through 0,\allowbreak{} J\_\allowbreak{}red sends every zero morphism to zero. This is the property needed for the following complexes,\allowbreak{} rather than additivity.

Define functors Q\_\allowbreak{}0(A)\allowbreak{}\allowbreak{}=A and recursively \[
 E_q(A)=J_{\rm red}(Q_q(A)),\qquad
 Q_{q+1}(A)=\operatorname{coker}
       (\eta_{\rm red,Q_q(A)}:Q_q(A)\to E_q(A)).
                                                               \tag{1.2}
\] Let pi\_\allowbreak{}q:\allowbreak{}E\_\allowbreak{}q\allowbreak{}->Q\_\allowbreak{}\{q\allowbreak{}+1\} be the quotient and v\_\allowbreak{}q\allowbreak{}=eta\_\allowbreak{}red,\allowbreak{}Q\_\allowbreak{}\{q\allowbreak{}+1\} pi\_\allowbreak{}q:\allowbreak{}E\_\allowbreak{}q\allowbreak{}->E\_\allowbreak{}\{q\allowbreak{}+1\}. Cokernel maps are induced by the natural commutative squares,\allowbreak{} so Q\_\allowbreak{}q,\allowbreak{} E\_\allowbreak{}q,\allowbreak{} pi\_\allowbreak{}q,\allowbreak{} and v\_\allowbreak{}q are functorial. They preserve zero objects and zero maps. We have v\_\allowbreak{}\{q\allowbreak{}+1\}v\_\allowbreak{}q\allowbreak{}=0,\allowbreak{} ker v\_\allowbreak{}q\allowbreak{}=im eta\_\allowbreak{}red,\allowbreak{}Q\_\allowbreak{}q,\allowbreak{} and im v\_\allowbreak{}\{q-1\}\allowbreak{}=im eta\_\allowbreak{}red,\allowbreak{}Q\_\allowbreak{}q for q>\allowbreak{}=1. The initial kernel is the image of A. Thus \[
 0\to A\xrightarrow{\eta_{\rm red,A}}E_0(A)
     \xrightarrow{v_0}E_1(A)\xrightarrow{v_1}\cdots             \tag{1.3}
\] is an injective resolution,\allowbreak{} naturally in A.

For the original bounded-below K,\allowbreak{} with K\textasciicircum{}p\allowbreak{}=0 for p<N,\allowbreak{} take I\textasciicircum{}\{p,\allowbreak{}q\}\allowbreak{}=E\_\allowbreak{}q(K\textasciicircum{}p)\allowbreak{} for q>\allowbreak{}=0 and zero otherwise,\allowbreak{} d\_\allowbreak{}1\textasciicircum{}\{p,\allowbreak{}q\}\allowbreak{}=E\_\allowbreak{}q(d\_\allowbreak{}K\textasciicircum{}p)\allowbreak{},\allowbreak{} and d\_\allowbreak{}2\textasciicircum{}\{p,\allowbreak{}q\}\allowbreak{}=v\_\allowbreak{}q(K\textasciicircum{}p)\allowbreak{}. Functoriality and preservation of zero prove d\_\allowbreak{}1\textasciicircum{}2\allowbreak{}=0. Equation (1.3)\allowbreak{} proves d\_\allowbreak{}2\textasciicircum{}2\allowbreak{}=0; naturality of v\_\allowbreak{}q proves d\_\allowbreak{}1d\_\allowbreak{}2\allowbreak{}=d\_\allowbreak{}2d\_\allowbreak{}1. All columns p\textless N vanish. The augmentation epsilon\textasciicircum{}p\allowbreak{}=eta\_\allowbreak{}red,\allowbreak{}K\textasciicircum{}p is a horizontal chain map and is killed by d\_\allowbreak{}2. The total complex has the exact differential \[
 d_{\operatorname{Tot}}|_{I^{p,q}}=d_1^{p,q}+(-1)^p d_2^{p,q}.
                                                               \tag{1.4}
\] In degree m its only possible summands have N<\allowbreak{}=p<\allowbreak{}=m. They form a finite sum of injectives. Column exactness in (1.3)\allowbreak{},\allowbreak{} and the finite-diagonal spectral comparison proved in the preceding note,\allowbreak{} show that epsilon:\allowbreak{}K\allowbreak{}->Tot(I)\allowbreak{} is a quasi-isomorphism. Each epsilon\^{}m is also monic:\allowbreak{} it is the monomorphism into I\textasciicircum{}\{m,\allowbreak{}0\} followed by the inclusion into that finite direct sum.

For a chain map f:\allowbreak{}K\allowbreak{}->L,\allowbreak{} use E\_\allowbreak{}q(f\textasciicircum{}p)\allowbreak{} on each summand. It commutes with both differentials by functoriality and naturality,\allowbreak{} hence with (1.4)\allowbreak{},\allowbreak{} and makes the augmentation square commute strictly. Identity and composition hold before passage to homotopy. This supplies the omitted verification at 7675–\allowbreak{}7677 and the functor inj to the actual arrow category. Different chosen lower bounds N merely certify the support:\allowbreak{} the formula (1.2)\allowbreak{} itself makes no choice of a bound for the object or for a map.

On homotopy classes the functor is the previously defined resolution functor j. If f and g are homotopic,\allowbreak{} epsilon\_\allowbreak{}L f and epsilon\_\allowbreak{}L g are homotopic. Precomposition with epsilon\_\allowbreak{}K is a bijection into the bounded-below injective target jL,\allowbreak{} so inj(f)\allowbreak{} and inj(g)\allowbreak{} are homotopic. The same argument,\allowbreak{} applied to inj(f\allowbreak{}+g)\allowbreak{} and inj(f)\allowbreak{}\allowbreak{}+inj(g)\allowbreak{},\allowbreak{} proves additivity of j. The latter two maps agree strictly after epsilon\_\allowbreak{}K. Their difference factors through jK/\allowbreak{}K,\allowbreak{} which is acyclic because epsilon\_\allowbreak{}K is a monomorphic quasi-isomorphism. A null homotopy of that quotient-to-jL map,\allowbreak{} pulled back along jK\allowbreak{}->jK/\allowbreak{}K,\allowbreak{} gives a homotopy h with h epsilon\_\allowbreak{}K\allowbreak{}=0. Thus this additivity comparison can be made relative to the original K; we do not assert a canonical selection of those homotopies.

There is a concrete reason not to assume strict additivity. On Q-vector spaces let J(V)\allowbreak{}\allowbreak{}=V direct-sum (V tensor\_\allowbreak{}Q V)\allowbreak{},\allowbreak{} eta(v)\allowbreak{}\allowbreak{}=(v,\allowbreak{}0)\allowbreak{},\allowbreak{} and J(f)\allowbreak{}\allowbreak{}=f direct-sum (f tensor f)\allowbreak{}. This is a zero-preserving functorial injective embedding,\allowbreak{} since every vector space is injective (extend linear maps from a basis of a subspace to a basis of the whole space)\allowbreak{}. For K\allowbreak{}=Q[0]\allowbreak{},\allowbreak{} (1.2)\allowbreak{} gives I\textasciicircum{}q\allowbreak{}=Q direct-sum Q for all q>\allowbreak{}=0 and d\textasciicircum{}q(x,\allowbreak{}y)\allowbreak{}\allowbreak{}=(y,\allowbreak{}0)\allowbreak{}. On a scalar lambda its degree-q map is diag(lambda\textasciicircum{}\{2\textasciicircum{}q\},\allowbreak{}lambda\textasciicircum{}\{2\textasciicircum{}\{q\allowbreak{}+1\}\})\allowbreak{}. For f\allowbreak{}=g\allowbreak{}=1,\allowbreak{} the additivity defect is therefore \[
 \Delta^q=
 \begin{pmatrix}2^{2^q}-2&0\\0&2^{2^{q+1}}-2\end{pmatrix}.
                                                               \tag{1.5}
\] In degree 0 it has nonzero second component 2. Put h\textasciicircum{}0\allowbreak{}=0 and for q>\allowbreak{}=1 set \[
 h^q(x,y)=(0,(2^{2^q}-2)x):I^q\longrightarrow I^{q-1}.
                                                               \tag{1.6}
\] Then d\^{}\{q-1\}h\^{}q contributes the first diagonal entry of (1.5)\allowbreak{},\allowbreak{} and h\textasciicircum{}\{q\allowbreak{}+1\}d\textasciicircum{}q the second. Their sum is exactly Delta,\allowbreak{} including q\allowbreak{}=0. The homotopy vanishes on K. This example proves both the failure at the chain level and the exact homotopy correction,\allowbreak{} with all coefficients retained.

RECON-052 reconciles P02-E0056 and DERIVED-048 with existing MC-STK-ERR-0859 (“remains”)\allowbreak{}. RECON-053 reconciles P02-E0057 and DERIVED-050 with the two MC-STK-ERR-0861 diagram edits to j(K\textasciicircum{}bullet)\allowbreak{} and j(L\textasciicircum{}bullet)\allowbreak{}. P02 also cites 7695,\allowbreak{} where the notation is already correct; that line is context and receives no extra edit.

\subsection{2. Large categories and the exact lift of a derived functor}\label{proofs-7589-8350-large-categories-and-the-exact-lift-of-a-derived-functor}

The construction (1.2)\allowbreak{} consists of specified functors,\allowbreak{} natural maps,\allowbreak{} and their cokernels,\allowbreak{} so it does not require selecting a resolution independently from a proper class of objects. The source\textquotesingle s applications use functorial injective embeddings for modules and sheaves. The current Injectives theorem at 1003–\allowbreak{}1014 is exactly the required assertion for sheaves on a ringed site. Its preceding displayed construction is J(F)\allowbreak{}\allowbreak{}=(FreeFlat(F\textasciicircum{}vee)\allowbreak{})\allowbreak{}\textasciicircum{}vee with the natural evaluation embedding. Its injectivity calculation identifies Hom(G,\allowbreak{}J(F)\allowbreak{})\allowbreak{} with Hom(G tensor FreeFlat(F\textasciicircum{}vee)\allowbreak{},\allowbreak{}J\_\allowbreak{}ab)\allowbreak{}; flatness of the first factor and injectivity of J\_\allowbreak{}ab make this exact. These are the particular construction and theorem cited by the receiver; the whole Injectives chapter was not re-reviewed here.

For presheaves,\allowbreak{} the topology consisting of isomorphism coverings has every presheaf as a sheaf:\allowbreak{} matching data on a single isomorphism pulls back uniquely along its inverse,\allowbreak{} and the same is true after any base change and composition. The sheaf theorem therefore yields the asserted presheaf embedding. Modules over a ring are sheaves of modules on the one-object,\allowbreak{} identity-morphism site with that ring; the same construction applies. These are the three completed examples. RECON-054 /\allowbreak{} P02-E0058 /\allowbreak{} DERIVED-049 retains MC-STK-ERR-0860\textquotesingle s deletion of the literal ``Add more here as needed.'' item. No new example is substituted into that placeholder.

Let F:\allowbreak{}A\allowbreak{}->B now be the additive functor of original 7778–\allowbreak{}7811,\allowbreak{} and (j,\allowbreak{}i)\allowbreak{} the resolution functor,\allowbreak{} with quasi-inverse j':\allowbreak{}D\textasciicircum{}\allowbreak{}+(A)\allowbreak{}\allowbreak{}->K\textasciicircum{}\allowbreak{}+(I)\allowbreak{}. Additive F preserves chain homotopies,\allowbreak{} finite direct sums,\allowbreak{} the cone differential including its minus sign,\allowbreak{} and shifts with their differential signs. It therefore gives an exact functor K\textasciicircum{}\allowbreak{}+(I)\allowbreak{}\allowbreak{}->K\textasciicircum{}\allowbreak{}+(B)\allowbreak{}. For K\allowbreak{}->jK the canonical derived comparison identifies RF(K)\allowbreak{} with F(jK)\allowbreak{}. Naturality follows from j(alpha)\allowbreak{}i\_\allowbreak{}K\allowbreak{}=i\_\allowbreak{}L alpha and the unique injective comparison; the same equation after localization handles inverse quasi-isomorphisms. Thus the diagram in the source 2-commutes,\allowbreak{} and the composite Fj\textquotesingle{} is an actual exact lift to K\textasciicircum{}\allowbreak{}+(B)\allowbreak{}. It is stronger data than its image in D\textasciicircum{}\allowbreak{}+(B)\allowbreak{},\allowbreak{} but no strict chain-level inverse to localization is asserted.

RECON-055 /\allowbreak{} P02-E0059 /\allowbreak{} DERIVED-051 is MC-STK-ERR-0862,\allowbreak{} the terminal period at 7780. RECON-056 /\allowbreak{} DERIVED-052 is MC-STK-ERR-0863,\allowbreak{} the defining order (j,\allowbreak{}i)\allowbreak{}.

\subsection{\texorpdfstring{3. Finite filtrations,\allowbreak{} split injectives,\allowbreak{} and the reversed pushout}{3. Finite  split  and the reversed pushout}}\label{proofs-7589-8350-finite-filtrations-split-injectives-and-the-reversed-pushout}

All filtrations in this section are decreasing,\allowbreak{} finite,\allowbreak{} exhaustive and separated. A map is strict when its image of F\^{}p equals its full image intersected with F\^{}p of the target. For a complex A\allowbreak{}->B\allowbreak{}->C,\allowbreak{} exactness of the associated graded complex implies exactness of each filtered piece and each quotient. Here is the finite induction used by current Homology 5184–\allowbreak{}5233. Remove the top common filtration level n. The three quotient objects have shorter filtrations; their graded complex is exact. Their quotient complexes are exact by induction. The complex of top pieces is exact as a graded piece. The short exact sequence of these three-term complexes then makes the original complex exact at its middle term by the cohomology long exact sequence. It also proves the quotient assertions,\allowbreak{} and the reversed finite induction proves the filtered piece assertions. Exactness on pieces and on the underlying objects gives f(F\^{}p A)\allowbreak{}\allowbreak{}=ker g intersect F\^{}p B for f:\allowbreak{}A\allowbreak{}->B,\allowbreak{} so f is strict. For g,\allowbreak{} if g(b)\allowbreak{} lies in F\^{}p C,\allowbreak{} exactness on quotients puts the class of b in the image of A/\allowbreak{}F\textasciicircum{}p A; subtracting that image leaves a representative in F\^{}p B with the same g-image. This proves strictness of g as an equality of image subobjects. These arguments can be expressed using the corresponding pullbacks and images in an arbitrary abelian category; no element-lifting surjectivity is imposed.

If each gr\^{}p I is injective,\allowbreak{} split 0\allowbreak{}->F\textasciicircum{}\{p\allowbreak{}+1\}I\allowbreak{}->F\textasciicircum{}p I\allowbreak{}->gr\textasciicircum{}p I\allowbreak{}->0 downwards from the top level. The kernel at each step is a finite sum of injectives,\allowbreak{} hence injective,\allowbreak{} so that sequence splits. We obtain I\allowbreak{}=direct-sum I\_\allowbreak{}p with F\^{}a I\allowbreak{}=direct-sum\_\allowbreak{}\{p>\allowbreak{}=a\}I\_\allowbreak{}p. Conversely such a decomposition has the required injective graded pieces. This proves 7851–\allowbreak{}7864 without requiring injectivity of a quotient to split an epimorphism.

Let u:\allowbreak{}A\allowbreak{}->B be a strict monomorphism and f:\allowbreak{}A\allowbreak{}->I filtered,\allowbreak{} with that finite splitting of I. The component f\_\allowbreak{}p to I\_\allowbreak{}p vanishes on F\textasciicircum{}\{p\allowbreak{}+1\}A,\allowbreak{} so factors through A/\allowbreak{}F\textasciicircum{}\{p\allowbreak{}+1\}A. Strictness of u makes that quotient a subobject of B/\allowbreak{}F\textasciicircum{}\{p\allowbreak{}+1\}B. Extend f\_\allowbreak{}p there by injectivity of I\_\allowbreak{}p. Pull it back to B,\allowbreak{} and assemble the finitely many components. For x in F\^{}a B every component of degree p\textless a vanishes,\allowbreak{} so this is a filtered extension g with gu\allowbreak{}=f. Taking A\allowbreak{}=I and f\allowbreak{}=1 gives the strict-monomorphism splitting lemma as well. For I\allowbreak{}=0 its splitting is the unique map; the source\textquotesingle s top-index induction is used only for I!\allowbreak{}=0.

The pushout used at 7906 is \[
 P=I\amalg_A B
   =\operatorname{coker}((f,-u):A\to I\oplus B),\qquad
 f':I\longrightarrow P.                                      \tag{3.1}
\] It is f\textquotesingle{} that is the pushout of u along f. Its underlying map is monic because u is. The quotient filtration is the image of F\^{}p I direct-sum F\^{}p B. If the image of x in I lies in F\^{}p P,\allowbreak{} the quotient relation says x-f(a)\allowbreak{} lies in F\^{}p I for some a with u(a)\allowbreak{} in F\^{}p B. Strictness of u gives a in F\^{}p A,\allowbreak{} so f(a)\allowbreak{} lies in F\^{}p I. Thus x lies in F\^{}p I and f\textquotesingle{} is strict. This is the specific implication of Homology 4988–\allowbreak{}5013 used here. A filtered retraction r:\allowbreak{}P\allowbreak{}->I exists by the preceding splitting lemma. If v:\allowbreak{}B\allowbreak{}->P is the other map,\allowbreak{} g\allowbreak{}=rv satisfies gu\allowbreak{}=rvu\allowbreak{}=rf'f\allowbreak{}=f.

\textbf{DERIVED-NEW-027.} Original 7906 calls f\textquotesingle{} the pushout ``of f by u''. Those roles are reversed:\allowbreak{} that would give the other map v,\allowbreak{} whose domain is B. The strict monomorphism follows from u,\allowbreak{} not from arbitrary f. Replace that phrase by ``of u by f''. For example,\allowbreak{} over Q-vector spaces with their single filtration step,\allowbreak{} take A\allowbreak{}=B\allowbreak{}=Q,\allowbreak{} u\allowbreak{}=1,\allowbreak{} I\allowbreak{}=0,\allowbreak{} f\allowbreak{}=0. Then P\allowbreak{}=0 and f':\allowbreak{}0\allowbreak{}->0 is monic,\allowbreak{} whereas the pushout of f along u is Q\allowbreak{}->0,\allowbreak{} which is not. The theorem and its displayed arrow are valid; the phrase identifies the wrong arrow operation.

For the embedding at 7913,\allowbreak{} choose a<\allowbreak{}=b with gr\^{}p A\allowbreak{}=0 outside [a,\allowbreak{}b]\allowbreak{} and monomorphisms w\_\allowbreak{}p:\allowbreak{}A/\allowbreak{}F\textasciicircum{}\{p\allowbreak{}+1\}A\allowbreak{}->I\_\allowbreak{}p. Define I\allowbreak{}=direct-sum\_\allowbreak{}\{a<\allowbreak{}=p<\allowbreak{}=b\}I\_\allowbreak{}p with the stated step filtration,\allowbreak{} and u\allowbreak{}=(w\_\allowbreak{}p pi\_\allowbreak{}p)\allowbreak{}\_\allowbreak{}p:\allowbreak{}A\allowbreak{}->I,\allowbreak{} where pi\_\allowbreak{}p is the actual quotient map. The b component is a monomorphism because F\textasciicircum{}\{b\allowbreak{}+1\}A\allowbreak{}=0,\allowbreak{} so u is monic. For a<r<\allowbreak{}=b\allowbreak{}+1,\allowbreak{} the condition u(x)\allowbreak{} in F\^{}r I forces its r-1 component to vanish; monicity of w\_\allowbreak{}\{r-1\} gives x in F\^{}r A. The converse uses pi\_\allowbreak{}p(F\textasciicircum{}r A)\allowbreak{}\allowbreak{}=0 for p\textless r. At r<\allowbreak{}=a both sides are the whole object; above b\allowbreak{}+1 both are zero by monicity. This proves strictness with the source\textquotesingle s quotients and direct-sum components,\allowbreak{} including the implicit quotient in its phrase ``sum of the maps u\_\allowbreak{}n”.

\subsection{\texorpdfstring{4. Filtered resolutions of objects,\allowbreak{} maps,\allowbreak{} and short exact sequences}{4. Filtered resolutions of   and short exact sequences}}\label{proofs-7589-8350-filtered-resolutions-of-objects-maps-and-short-exact-sequences}

Successively embed A,\allowbreak{} then its filtered cokernel,\allowbreak{} then the next cokernel into filtered injectives by section 3. Strictness of each embedding makes 0\allowbreak{}->Q\_\allowbreak{}n\allowbreak{}->I\textasciicircum{}\{n\allowbreak{}+1\}\allowbreak{}->Q\_\allowbreak{}\{n\allowbreak{}+1\}\allowbreak{}->0 exact on every graded piece. The differential is quotient followed by the next embedding,\allowbreak{} so its square is zero and its graded kernel is the preceding graded image. Thus A[0]\allowbreak{}\allowbreak{}->I is a filtered quasi-isomorphism,\allowbreak{} with zero negative terms. RECON-057 /\allowbreak{} P02-E0060 /\allowbreak{} DERIVED-054 retains both MC-STK-ERR-0865 corrections:\allowbreak{} the filtered injective objects belong to Fil\textasciicircum{}f(A)\allowbreak{}.

For f:\allowbreak{}A\allowbreak{}->B and the given resolutions A\allowbreak{}->I,\allowbreak{} B\allowbreak{}->J,\allowbreak{} their augmented complexes are graded-exact. Section 3 therefore makes the augmentations strict monic and every differential strict. Extend bf across A\allowbreak{}->I\textasciicircum{}0 to f\textasciicircum{}0:\allowbreak{}I\textasciicircum{}0\allowbreak{}->J\textasciicircum{}0. The map d\_\allowbreak{}J\textasciicircum{}0 f\^{}0 vanishes on A,\allowbreak{} so factors through coker(A\allowbreak{}->I\textasciicircum{}0)\allowbreak{}\allowbreak{}=coim d\_\allowbreak{}I\textasciicircum{}0 \allowbreak{}=im d\_\allowbreak{}I\textasciicircum{}0,\allowbreak{} with its induced filtration. Extend it from that strict subobject of I\^{}1 to f\textasciicircum{}1:\allowbreak{}I\textasciicircum{}1\allowbreak{}->J\textasciicircum{}1. At each subsequent step d\_\allowbreak{}J f\^{}n vanishes on the preceding image because the earlier chain equation and d\_\allowbreak{}J\textasciicircum{}2\allowbreak{}=0 give that equality. The same factorization and extension continue. This proves the strict commuting diagram at 7960–\allowbreak{}7974. RECON-058 /\allowbreak{} P02-E0061 /\allowbreak{} DERIVED-053 is MC-STK-ERR-0864,\allowbreak{} replacing the undefined C[0]\allowbreak{} by the supplied B[0]\allowbreak{}.

Now keep the source\textquotesingle s short exact sequence 0\allowbreak{}->A --alpha-\allowbreak{}-> B --beta-\allowbreak{}-> C\allowbreak{}->0,\allowbreak{} with its prescribed resolutions A\allowbreak{}->I and C\allowbreak{}->J. Extend a:\allowbreak{}A\allowbreak{}->I\textasciicircum{}0 to b:\allowbreak{}B\allowbreak{}->I\textasciicircum{}0. Since d\_\allowbreak{}I\textasciicircum{}0 b kills alpha,\allowbreak{} there is a unique bar b:\allowbreak{}C\allowbreak{}->I\textasciicircum{}1 with bar b beta\allowbreak{}=d\_\allowbreak{}I\textasciicircum{}0 b. Extend bar b across c:\allowbreak{}C\allowbreak{}->J\textasciicircum{}0 to delta\textasciicircum{}0:\allowbreak{}J\textasciicircum{}0\allowbreak{}->I\textasciicircum{}1. The equality d\_\allowbreak{}I\textasciicircum{}1 bar b\allowbreak{}=0 follows after composition with the epimorphism beta. Thus d\_\allowbreak{}I\textasciicircum{}1 delta\^{}0 kills c(C)\allowbreak{} and factors through im d\_\allowbreak{}J\textasciicircum{}0. Extend to delta\textasciicircum{}1:\allowbreak{}J\textasciicircum{}1\allowbreak{}->I\textasciicircum{}2. Recursively,\allowbreak{} the preceding identity and d\_\allowbreak{}I\textasciicircum{}2\allowbreak{}=0 give \[
 d_I^{r+1}\delta^r=\delta^{r+1}d_J^r
 \quad(r\ge0).                                               \tag{4.1}
\] Each extension is filtered,\allowbreak{} since im d\_\allowbreak{}J\textasciicircum{}r is a strict subobject. Set M\textasciicircum{}r\allowbreak{}=I\textasciicircum{}r direct-sum J\^{}r and \[
 d_M^r=
 \begin{pmatrix}
 d_I^r&(-1)^{r+1}\delta^r\\
 0&d_J^r
 \end{pmatrix}.                                              \tag{4.2}
\] Its off-diagonal square is (-1)\allowbreak{}\textasciicircum{}\{r\allowbreak{}+1\}d\_\allowbreak{}I\textasciicircum{}\{r\allowbreak{}+1\}delta\textasciicircum{}r \allowbreak{}+(-1)\allowbreak{}\textasciicircum{}\{r\allowbreak{}+2\}delta\textasciicircum{}\{r\allowbreak{}+1\}d\_\allowbreak{}J\textasciicircum{}r\allowbreak{}=0 by (4.1)\allowbreak{}. The augmentation B\allowbreak{}->M\textasciicircum{}0 is (b,\allowbreak{}c beta)\allowbreak{}. Its differential has first component d\_\allowbreak{}I\textasciicircum{}0 b-delta\^{}0 c beta \allowbreak{}=bar b beta-bar b beta\allowbreak{}=0 and second component zero. The maps from I and to J are the actual inclusion and projection. They give a termwise filtered split exact sequence,\allowbreak{} compatible with the original short exact sequence. Applying gr\^{}p gives a map of short exact sequences of complexes with outer quasi-isomorphisms. The cohomology exact sequences prove that the middle augmentation is a quasi-isomorphism on every graded piece.

RECON-059 /\allowbreak{} P02-E0062 /\allowbreak{} DERIVED-063 is the existing MC-STK-ERR-0874 chapter-prefix repair at 8044. The unprefixed label actually names a different Derived result; it is not an absent label as P02 says. The required result is the strictness lemma in Homology proved above. RECON-060 /\allowbreak{} P02-E0063 /\allowbreak{} DERIVED-055 is MC-STK-ERR-0866,\allowbreak{} the source B of (b,\allowbreak{}c beta)\allowbreak{}. P02\textquotesingle s earlier diagram loci 8049–\allowbreak{}8050 and 8060 are supporting context; the correction is at 8092 alone.

\subsection{5. Arbitrary bounded-below filtered complexes}\label{proofs-7589-8350-arbitrary-bounded-below-filtered-complexes}

Keep K\textasciicircum{}p\allowbreak{}=0 for p\textless N in its original degrees. For each p>\allowbreak{}=N form the source\textquotesingle s strict short exact sequence \[
 0\to\ker d_K^p\to K^p
       \to\operatorname{coim}d_K^p\to0.                      \tag{5.1}
\] The kernel has the induced filtration and the coimage the quotient filtration. No strictness of d\_\allowbreak{}K\textasciicircum{}p itself is assumed. Since d\_\allowbreak{}K\textasciicircum{}\{p\allowbreak{}+1\}d\_\allowbreak{}K\textasciicircum{}p\allowbreak{}=0,\allowbreak{} there is a filtered map u\_\allowbreak{}p:\allowbreak{}coim d\_\allowbreak{}K\textasciicircum{}p\allowbreak{}->ker d\_\allowbreak{}K\textasciicircum{}\{p\allowbreak{}+1\} through which the original differential factors. The image of F\^{}r K\^{}p lies in the induced F\^{}r of that kernel,\allowbreak{} so the factor is a filtered map with the quotient filtration on its domain.

Resolve the two outer objects in (5.1)\allowbreak{},\allowbreak{} lift u\_\allowbreak{}p by section 4,\allowbreak{} and apply its filtered horseshoe with those outer resolutions fixed. Call the resulting maps alpha\_\allowbreak{}p:\allowbreak{}I\_\allowbreak{}ker,\allowbreak{}p\allowbreak{}->I\_\allowbreak{}p and beta\_\allowbreak{}p:\allowbreak{}I\_\allowbreak{}p\allowbreak{}->I\_\allowbreak{}coim,\allowbreak{}p. The horizontal differential is exactly \[
 d_p=\alpha_{p+1}u_p^\bullet\beta_p.                          \tag{5.2}
\] Its composite with d\_\allowbreak{}\{p-1\} is zero because beta\_\allowbreak{}p alpha\_\allowbreak{}p\allowbreak{}=0. All factors in (5.2)\allowbreak{} are vertical chain maps,\allowbreak{} so it commutes with the vertical differentials. Every column is a filtered injective resolution of K\textasciicircum{}p,\allowbreak{} zero in negative vertical degree. Totalize with d\_\allowbreak{}1\allowbreak{}+(-1)\allowbreak{}\textasciicircum{}p d\_\allowbreak{}2. In degree m there are only the columns N<\allowbreak{}=p<\allowbreak{}=m. Every term is therefore filtered injective and has finite filtration. The augmentation K\textasciicircum{}m\allowbreak{}->I\_\allowbreak{}m\textasciicircum{}0 is strict monic,\allowbreak{} since its associated graded maps are monic at the bottom of their resolutions. The inclusion of I\_\allowbreak{}m\textasciicircum{}0 into the total finite direct sum is strict monic too. Their composite is the required component.

The associated graded functor is additive,\allowbreak{} hence commutes with these finite total sums and with the differential including (-1)\allowbreak{}\textasciicircum{}p. Applying the column comparison to each gr\^{}r proves that the actual augmentation K\allowbreak{}->Tot I is a filtered quasi-isomorphism. It is zero below N. This proves all three clauses of 8095–\allowbreak{}8103 without a change of the original degree coordinates. The source\textquotesingle s shifted proof is compatible:\allowbreak{} undoing its initial shift on both the source and constructed resolution multiplies the differentials by the prescribed shift sign. No strict commutation of a nonadditive embedding functor with shifts is asserted. We have instead specified the construction directly in the original p.

RECON-061 /\allowbreak{} DERIVED-056 is MC-STK-ERR-0867,\allowbreak{} inserting ``by'' in the assignment of the total complex notation.

\subsection{\texorpdfstring{6. Filtered null homotopies,\allowbreak{} comparison,\allowbreak{} and equivalence}{6. Filtered null   and equivalence}}\label{proofs-7589-8350-filtered-null-homotopies-comparison-and-equivalence}

Let C be filtered acyclic and I a complex of filtered injectives,\allowbreak{} zero below N. For a filtered chain map f:\allowbreak{}C\allowbreak{}->I set h\textasciicircum{}r\allowbreak{}=0 for r<\allowbreak{}=N. Suppose h through degree r has been constructed,\allowbreak{} and put \[
 e^r=f^r-d_I^{r-1}h^r.
\] The earlier homotopy equation and the chain equation give e\^{}r d\_\allowbreak{}C\textasciicircum{}\{r-1\}\allowbreak{}=0. Filtered exactness of C identifies coker d\_\allowbreak{}C\textasciicircum{}\{r-1\} with im d\_\allowbreak{}C\textasciicircum{}r,\allowbreak{} with its induced filtration. Thus e\^{}r factors as a filtered map on im d\_\allowbreak{}C\textasciicircum{}r. Extend it across its strict inclusion into C\textasciicircum{}\{r\allowbreak{}+1\} to h\textasciicircum{}\{r\allowbreak{}+1\}:\allowbreak{}C\textasciicircum{}\{r\allowbreak{}+1\}\allowbreak{}->I\textasciicircum{}r. Then \[
 f^r=d_I^{r-1}h^r+h^{r+1}d_C^r.                              \tag{6.1}
\] At each fixed degree only this step and the preceding step contribute,\allowbreak{} so the infinite upward induction defines an ordinary homotopy,\allowbreak{} not an infinite sum in A. This proves 8196–\allowbreak{}8238 and the filtered null-homotopy fact used in PART10 section 3.

For a filtered quasi-isomorphism alpha:\allowbreak{}K\allowbreak{}->L,\allowbreak{} its cone and shifted cone are filtered acyclic. Apply Hom(-,\allowbreak{}I)\allowbreak{} to the source\textquotesingle s triangle (K,\allowbreak{}L,\allowbreak{}C(alpha)\allowbreak{},\allowbreak{}alpha,\allowbreak{}i,\allowbreak{}-p)\allowbreak{}. The outer Hom groups vanish by (6.1)\allowbreak{},\allowbreak{} so precomposition is a bijection \[
 \operatorname{Hom}_{K(\operatorname{Fil}^f A)}(L,I)
 \xrightarrow{\sim}
 \operatorname{Hom}_{K(\operatorname{Fil}^f A)}(K,I).           \tag{6.2}
\] A derived roof gamma alpha\^{}\{-1\} is represented by the unique homotopy class beta with beta alpha\allowbreak{}=gamma. If a homotopy class becomes zero after localization,\allowbreak{} a denominator alpha makes its precomposition zero,\allowbreak{} and (6.2)\allowbreak{} makes it zero. This proves both surjectivity and injectivity of the Hom comparison with DF(A)\allowbreak{}. RECON-062 /\allowbreak{} P02-E0064 /\allowbreak{} DERIVED-057 retains the two MC-STK-ERR-0868 repairs of the ambient Hom category. The assertion is (6.2)\allowbreak{}; it does not identify filtered and unfiltered homotopy classes.

Finite sums and shifts of filtered injectives remain filtered injective,\allowbreak{} so K\textasciicircum{}\allowbreak{}+(I\textasciicircum{}f)\allowbreak{} is closed under the actual cone triangles. Its functor to DF\textasciicircum{}\allowbreak{}+(A)\allowbreak{} is exact,\allowbreak{} and (6.2)\allowbreak{} makes it fully faithful. An object of DF\textasciicircum{}\allowbreak{}+(A)\allowbreak{} has a representative with gr cohomology zero below some N. The already proved bounded-representative lemma at original 4385–\allowbreak{}4431 replaces that representative by a filtered quasi-isomorphic complex zero below N:\allowbreak{} strictness of d\^{}\{N-1\} justifies the quotient K\textasciicircum{}N/\allowbreak{}im d\^{}\{N-1\} at its bottom degree. Applying section 5 then supplies a bounded-below filtered injective representative. Thus the functor is essentially surjective. Both gr\^{}p and forgetting the filtration send a filtered injective to an injective object,\allowbreak{} by its finite splitting. They preserve shifts,\allowbreak{} cones and filtered quasi-isomorphisms. The two squares in 8315–\allowbreak{}8335 therefore commute with their displayed functors.

\subsection{\texorpdfstring{7. DERIVED-UNDER-013:\allowbreak{} degree bounds and functorial fixed filtration windows}{7.  degree bounds and functorial fixed filtration windows}}\label{proofs-7589-8350-derived-under-013-degree-bounds-and-functorial-fixed-filtration-windows}

The construction in section 5 can retain more than just finiteness. Suppose K\textasciicircum{}p,\allowbreak{} p>\allowbreak{}=N,\allowbreak{} has filtration in the interval [a\_\allowbreak{}p,\allowbreak{}b\_\allowbreak{}p]\allowbreak{},\allowbreak{} meaning F\textasciicircum{}\{a\_\allowbreak{}p\}K\textasciicircum{}p\allowbreak{}=K\textasciicircum{}p and F\textasciicircum{}\{b\_\allowbreak{}p\allowbreak{}+1\}K\textasciicircum{}p\allowbreak{}=0. Its kernel and coimage in (5.1)\allowbreak{} have the same permitted bounds. Section 3\textquotesingle s embedding uses precisely the indices a\_\allowbreak{}p<\allowbreak{}=r<\allowbreak{}=b\_\allowbreak{}p,\allowbreak{} and its cokernel retains that interval. Iterating gives both outer resolutions in that same interval. Their horseshoe middle is their direct sum in each degree,\allowbreak{} so retains it too. Consequently the resulting total resolution has in degree m>\allowbreak{}=N the explicit interval \[
 \left[\min_{N\le p\le m}a_p,\
       \max_{N\le p\le m}b_p\right].                         \tag{7.1}
\] No term with p\textgreater m contributes. This is a finite interval for each m,\allowbreak{} even if the filtration widths of K have no common global bound. All terms below N remain zero.

If A has a functorial injective embedding,\allowbreak{} fix a single interval [a,\allowbreak{}b]\allowbreak{}. On the category of filtered objects supported in that interval define the functor \[
 \mathcal E(A)=
 \bigoplus_{r=a}^{b}J_{\rm red}(A/F^{r+1}A),\qquad
 F^p\mathcal E(A)=
 \bigoplus_{\substack{a\le r\le b\\r\ge p}}
              J_{\rm red}(A/F^{r+1}A).                       \tag{7.2}
\] The natural map is the tuple of quotient maps followed by eta\_\allowbreak{}red. The proof at the end of section 3 shows it is strict monic. Each graded piece is injective,\allowbreak{} so the target is filtered injective. Quotient maps of filtered morphisms give the functor on maps. It preserves zero because J\_\allowbreak{}red does,\allowbreak{} and the filtered cokernel of its natural embedding stays in [a,\allowbreak{}b]\allowbreak{}.

Repeat (1.2)\allowbreak{} in this filtered category,\allowbreak{} using (7.2)\allowbreak{} and these cokernels. Every augmented column is exact on each graded piece,\allowbreak{} because every embedding is strict. For a bounded-below complex all of whose terms have this fixed filtration interval,\allowbreak{} the horizontal maps are the images of its actual differentials under these zero-preserving functors. Their squares are zero,\allowbreak{} and naturality commutes the two differentials. The total construction (1.4)\allowbreak{} gives a functorial strict monomorphic filtered quasi-isomorphism,\allowbreak{} retains the original lower term bound N,\allowbreak{} and retains [a,\allowbreak{}b]\allowbreak{} in every degree. It uses the specified functors and no selection of resolutions over a proper class of objects.

The two constructions have different assertions:\allowbreak{} (7.1)\allowbreak{} covers arbitrary finite filtrations,\allowbreak{} with the source\textquotesingle s choices; (7.2)\allowbreak{} is functorial on a fixed filtration window. A window chosen separately for each object has not been shown to yield a functor on all finite filtered objects. We do not make that assertion.

The resulting homotopy functor is additive and exact by section 6. Its strict chain functor need not be additive,\allowbreak{} as (1.5)\allowbreak{} proves already with a single filtration step. The additivity defect has a filtered homotopy relative to the source:\allowbreak{} its factor through the cokernel of the strict augmentation has filtered acyclic domain,\allowbreak{} so (6.1)\allowbreak{} gives that homotopy. This adds an exact comparison to the stronger support and functorial assertions; no higher coherence of chosen homotopies is claimed.

\subsection{\texorpdfstring{8. DERIVED-UNDER-014:\allowbreak{} both prescribed resolutions for filtered complexes}{8.  both prescribed resolutions for filtered complexes}}\label{proofs-7589-8350-derived-under-014-both-prescribed-resolutions-for-filtered-complexes}

The source\textquotesingle s horseshoe in section 4 starts with objects in degree zero. It extends to any degreewise strict short exact sequence of bounded-below filtered complexes \[
 0\to A\xrightarrow{u}B\xrightarrow{v}C\to0,                 \tag{8.1}
\] with any two given filtered quasi-isomorphisms f\_\allowbreak{}1:\allowbreak{}A\allowbreak{}->I\_\allowbreak{}1 and f\_\allowbreak{}3:\allowbreak{}C\allowbreak{}->I\_\allowbreak{}3 into bounded-below filtered injective complexes. Neither given resolution map must be a termwise monomorphism.

Form the filtered-complex pushout \[
 E=\operatorname{coker}
       ((f_1,-u):A\to I_1\oplus B).                          \tag{8.2}
\] Section 3 makes 0\allowbreak{}->I\_\allowbreak{}1\allowbreak{}->E\allowbreak{}->C\allowbreak{}->0 degreewise strict exact. Apply gr\^{}p to the map from (8.1)\allowbreak{} into this sequence. Its outer maps are gr\^{}p f\_\allowbreak{}1 and the identity,\allowbreak{} so B\allowbreak{}->E is a filtered quasi-isomorphism. Split I\_\allowbreak{}1\textasciicircum{}n\allowbreak{}->E\textasciicircum{}n in the filtered category. This identifies \[
 E^n=I_1^n\oplus C^n,\qquad
 d_E^n=\begin{pmatrix}d_{I_1}^n&\delta^n\\0&d_C^n\end{pmatrix}.
                                                               \tag{8.3}
\] Every entry is filtered. The squared differential gives d\_\allowbreak{}\{I\_\allowbreak{}1\}delta\allowbreak{}+delta d\_\allowbreak{}C\allowbreak{}=0,\allowbreak{} so delta:\allowbreak{}C\allowbreak{}->I\_\allowbreak{}1[1]\allowbreak{} is a filtered chain map with the original positive off-diagonal entry. Apply (6.2)\allowbreak{} to f\_\allowbreak{}3 and target I\_\allowbreak{}1[1]\allowbreak{}. Choose its preimage delta':\allowbreak{}I\_\allowbreak{}3\allowbreak{}->I\_\allowbreak{}1[1]\allowbreak{} and a filtered homotopy h\textasciicircum{}n:\allowbreak{}C\textasciicircum{}n\allowbreak{}->I\_\allowbreak{}1\textasciicircum{}n with \[
 (\delta')^n f_3^n-\delta^n
      =-d_{I_1}^n h^n+h^{n+1}d_C^n.                         \tag{8.4}
\] The minus sign comes from d\_\allowbreak{}\{I\_\allowbreak{}1[1]\allowbreak{}\}. Set \[
 M^n=I_1^n\oplus I_3^n,\quad
 d_M^n=\begin{pmatrix}d_{I_1}^n&(\delta')^n\\0&d_{I_3}^n\end{pmatrix},
 \quad
 e^n=
 \begin{pmatrix}1&h^n\\0&f_3^n\end{pmatrix}:E^n\to M^n.        \tag{8.5}
\] The chain equation for delta\textquotesingle{} proves d\_\allowbreak{}M\textasciicircum{}2\allowbreak{}=0. The two off-diagonal entries in d\_\allowbreak{}M e and e d\_\allowbreak{}E are respectively d\_\allowbreak{}\{I\_\allowbreak{}1\}h\allowbreak{}+delta'f\_\allowbreak{}3 and delta\allowbreak{}+h d\_\allowbreak{}C; (8.4)\allowbreak{} makes them equal. Thus e is a filtered chain map. Its outer maps in the split sequences are 1 and f\_\allowbreak{}3,\allowbreak{} so gr\^{}p e is a quasi-isomorphism for every p by the cohomology exact sequences. Composing B\allowbreak{}->E\allowbreak{}->M gives the middle filtered resolution,\allowbreak{} compatible with both prescribed outer ones. Its terms are filtered injective,\allowbreak{} and it retains any common lower term bound and common filtration interval of I\_\allowbreak{}1,\allowbreak{}I\_\allowbreak{}3. No extra global enough-injectives hypothesis is used once these resolutions are supplied.

This is the exact filtered counterpart of DERIVED-UNDER-010,\allowbreak{} including the homotopy matrix h. Replacing e by the diagonal matrix would be unjustified unless delta'f\_\allowbreak{}3\allowbreak{}=delta strictly. The source\textquotesingle s degree-zero construction already obtains that strict equation in its own signed form (4.1)\allowbreak{}–\allowbreak{}(4.2)\allowbreak{}; it is retained,\allowbreak{} not replaced.

\subsection{9. Receiving calculations and limits}\label{proofs-7589-8350-receiving-calculations-and-limits}

Current Cohomology 104–\allowbreak{}145 and Sites Cohomology 145–\allowbreak{}182 construct RF using resolution functors on their actual module categories. Sections 1–\allowbreak{}2 justify the strict chain construction,\allowbreak{} its descent to additive homotopy functors,\allowbreak{} and its exact lift F j\textquotesingle{} in those categories. Their F is left exact and therefore additive; no extra condition is required. The typo on the next Sites Cohomology line 183 is a separate receiving-source observation,\allowbreak{} queued for that chapter\textquotesingle s report reconciliation. It changes none of the mathematical maps reviewed here.

Current Cohomology 7310–\allowbreak{}7339 uses a bounded-below filtered injective resolution with a finite filtration on every term. Formula (7.1)\allowbreak{} proves exactly that finiteness even when the original filtration intervals vary with degree. If the original terms share [a,\allowbreak{}b]\allowbreak{},\allowbreak{} (7.2)\allowbreak{} gives a functorial choice with that same interval. The filtered complex of global sections uses F\^{}p Gamma(X,\allowbreak{}I\textasciicircum{}m)\allowbreak{} \allowbreak{}=Gamma(X,\allowbreak{}F\textasciicircum{}p I\textasciicircum{}m)\allowbreak{}; the split filtration makes these literal subobjects and keeps the same bounds. Thus its associated graded complexes and their cohomology use the same columns and signs proved above.

PART10 section 3 used filtered comparison to prove Cartan–\allowbreak{}Eilenberg functoriality and the decalage formula. Section 6 now supplies the full source-ordered review of that comparison and its category correction. Section 5 supplies the general filtered resolution existence needed in that forward passage. The proof-specific pushout correction propagates through the extension,\allowbreak{} map,\allowbreak{} horseshoe and total resolution constructions in sections 3–\allowbreak{}5. No theorem statement is weakened or silently strengthened by that correction.

The subsequent large-category remark and general filtered derived functor discussion were read through original 8420 but are next in the semantic review. Original 8407 already announces the additive filtered derived-functor treatment. That is relevant prior source evidence for PART10\textquotesingle s additive consequence; it reinforces the absence of a novelty claim. No complete review of the unread continuation or of the entire Injectives chapter is asserted here.
\clearpage
\section{\texorpdfstring{Filtered derived functors,\allowbreak{} page bounds,\allowbreak{} and the first Ext comparison}{Filtered derived  page  and the first Ext comparison}}\label{proofs-8352-8772-filtered-derived-functors-page-bounds-and-the-first-ext-comparison}

Editorial review of the Stacks Project authors\textquotesingle{} Derived Categories,\allowbreak{} original commit a04446e5\allowbreak{}7ec1fbc2\allowbreak{}52a871af\allowbreak{}cec7752f\allowbreak{}b2807b14, lines 8352–\allowbreak{}8772. SOURCE\_\allowbreak{}READING\_\allowbreak{}PART12.json records the frozen identities and the actual reading ranges. Earlier proof files remain part of this review,\allowbreak{} in particular PROOFS\_\allowbreak{}7096\_\allowbreak{}7578.md and PROOFS\_\allowbreak{}7589\_\allowbreak{}8350.md. The extensions below are editorial consequences; neither novelty nor source-body replacement is claimed.

\subsection{1. The set of objects needed for a calculation}\label{proofs-8352-8772-the-set-of-objects-needed-for-a-calculation}

Original 8352–\allowbreak{}8366 invokes current Sets,\allowbreak{} lemma-abelian-injectives,\allowbreak{} 1135–\allowbreak{}1157. That lemma includes more than an exact embedding of a small abelian category:\allowbreak{} it requires its injective objects to be exactly the objects that are injective in the ambient category. This is the property needed to compare the resulting resolutions.

Here is the construction in the usual locally small big-category setting of that statement. Start with the given set of objects and,\allowbreak{} when a set of diagrams is specified,\allowbreak{} every object occurring in those diagrams. Inductively enlarge a set S\_\allowbreak{}n of objects to a set S\_\allowbreak{}\{n\allowbreak{}+1\} containing:\allowbreak{} the zero object and chosen finite biproducts of objects of S\_\allowbreak{}n; chosen kernels and cokernels of every morphism between objects of S\_\allowbreak{}n; a chosen monomorphism X\allowbreak{}->J\_\allowbreak{}X with J\_\allowbreak{}X injective for every X in S\_\allowbreak{}n; and witnesses to noninjectivity for every X in S\_\allowbreak{}n that is not injective in the ambient category. Such a witness consists of a monomorphism U\_\allowbreak{}X\allowbreak{}->V\_\allowbreak{}X and a map U\_\allowbreak{}X\allowbreak{}->X admitting no extension V\_\allowbreak{}X\allowbreak{}->X; include U\_\allowbreak{}X and V\_\allowbreak{}X. Each stage has a set of requests:\allowbreak{} the category is locally small,\allowbreak{} there are a set of pairs of existing objects,\allowbreak{} and each witness involves finitely many objects and maps. Set-indexed collection and choice make the enlargement a set. This does not select data simultaneously for every object of the proper class.

Let A\textquotesingle{} be the full subcategory on the union of the S\_\allowbreak{}n. It is additive and has kernels and cokernels computed in A,\allowbreak{} because every finite collection of its objects appears at one stage. The comparison from coimage to image is the ambient isomorphism,\allowbreak{} so A\textquotesingle{} is abelian and its inclusion is exact. Every object X embeds in its J\_\allowbreak{}X within A',\allowbreak{} hence A\textquotesingle{} has enough injectives. An ambient injective in A\textquotesingle{} is injective there:\allowbreak{} all monomorphisms of A\textquotesingle{} remain monic in A,\allowbreak{} and fullness includes the ambient extensions. Conversely a noninjective X in A\textquotesingle{} has its chosen witness in a later stage of A',\allowbreak{} so is not injective in A\textquotesingle. This proves the additional injectivity property rather than assuming that exactness of an embedding alone would imply it.

For a set of filtered complexes,\allowbreak{} include all their terms,\allowbreak{} all their filtration subobjects,\allowbreak{} and the objects in their specified diagrams in S\_\allowbreak{}0. There are only set many:\allowbreak{} the degrees and filtration indices are integers. Each such diagram is then a diagram in A\textquotesingle. A filtered injective resolution formed there is also one in A,\allowbreak{} since its graded injectives are ambient injectives. The homotopy Hom groups between these complexes agree by fullness. Computing bounded-below filtered derived Hom by these resolutions therefore proves that the inclusion on these derived objects is fully faithful and computes the same maps. The analogous argument without filtrations proves the bounded-below unfiltered comparison. Given two such choices of A',\allowbreak{} apply the same construction to the union of their object sets; the resulting comparison uses the same ambient Hom groups. This proves compatibility of the calculations and does not assert a quasi-inverse selected on the entire proper class.

The remark\textquotesingle s ``only if'' refers to its preceding use of ordinary choice to pick representatives for a set of objects. It is not used here as an impossibility theorem about every big category. For example,\allowbreak{} PART11 supplies an explicit functorial construction on a fixed filtration window when functorial embeddings are given. No new source defect is counted from that reading of the remark. RECON-063 /\allowbreak{} DERIVED-058 reconciles the existing MC-STK-ERR-0869 possessive correction at 8364.

\subsection{2. A canonical functor on split filtered objects}\label{proofs-8352-8772-a-canonical-functor-on-split-filtered-objects}

Let T:\allowbreak{}A\allowbreak{}->B be additive. For a filtered injective I,\allowbreak{} all the sequences \[
 0\to F^{p+1}I\longrightarrow F^p I
      \longrightarrow \operatorname{gr}^p I\to0              \tag{2.1}
\] and all inclusions F\^{}p I\allowbreak{}->I split,\allowbreak{} by PART11 section 3. Thus T(F\^{}p I)\allowbreak{}\allowbreak{}->T(I)\allowbreak{} is monic. Give T(I)\allowbreak{} the filtration consisting of those actual subobjects. It is finite and is functorial:\allowbreak{} a filtered map f:\allowbreak{}I\allowbreak{}->J gives the commutative square from F\^{}p I\allowbreak{}->I to F\^{}p J\allowbreak{}->J,\allowbreak{} whose T-image restricts T(f)\allowbreak{} to the required subobjects. Addition,\allowbreak{} identity and composition are those of T.

If a decomposition I\allowbreak{}=direct-sum I\_\allowbreak{}p is chosen as in the source,\allowbreak{} the finite direct-sum comparison identifies this filtered T(I)\allowbreak{} with direct-sum T(I\_\allowbreak{}p)\allowbreak{},\allowbreak{} with filtration direct-sum\_\allowbreak{}\{p>\allowbreak{}=a\}T(I\_\allowbreak{}p)\allowbreak{}. A filtered map in such coordinates has components f\_\allowbreak{}\{q p\}:\allowbreak{}I\_\allowbreak{}p\allowbreak{}->J\_\allowbreak{}q with f\_\allowbreak{}\{q p\}\allowbreak{}=0 for q\textless p. Apply T to each component. Composition of the matrices is preserved because all relevant sums are finite and T is additive. A different splitting changes coordinates by an invertible filtration-preserving matrix; applying T gives an invertible change of coordinates with the same composition identities. This supplies the omitted map construction and proves independence from the noncanonical splittings in 8415–\allowbreak{}8439.

Applying T to (2.1)\allowbreak{} preserves its split exactness. The resulting quotient map is therefore a canonical natural isomorphism \[
 \operatorname{gr}^p(T_{\rm ext}I)
       \xrightarrow{\sim}T(\operatorname{gr}^p I).            \tag{2.2}
\] It is induced by the original quotient in (2.1)\allowbreak{},\allowbreak{} and does not require choosing a splitting. Forgetting the filtration similarly identifies T\_\allowbreak{}ext(I)\allowbreak{} with T(I)\allowbreak{},\allowbreak{} naturally. These assertions hold for every finite split filtered object,\allowbreak{} not only for injectives.

They do not assert commutation on arbitrary filtered objects. For example fix a prime ell,\allowbreak{} let T(M)\allowbreak{}\allowbreak{}=M[ell]\allowbreak{},\allowbreak{} and filter Z by F\^{}0 Z\allowbreak{}=Z,\allowbreak{} F\^{}1 Z\allowbreak{}=ell Z,\allowbreak{} F\^{}2 Z\allowbreak{}=0. Then T(Z)\allowbreak{}\allowbreak{}=0,\allowbreak{} so gr\textasciicircum{}0(T(Z)\allowbreak{})\allowbreak{}\allowbreak{}=0,\allowbreak{} whereas T(gr\^{}0 Z)\allowbreak{}\allowbreak{}=T(Z/\allowbreak{}ell Z)\allowbreak{}\allowbreak{}=Z/\allowbreak{}ell Z. This T is left exact:\allowbreak{} an element of the kernel of a map on ell-torsion is exactly the ell-torsion of its original kernel. Thus even left exactness does not make arbitrary associated graded objects commute with T. Original 8379 already gives that warning. The commutation mentioned at 8448 is the restricted one (2.2)\allowbreak{},\allowbreak{} in the context of T\_\allowbreak{}ext on filtered injectives. No extra defect is counted from that contextual shorthand.

\subsection{\texorpdfstring{3. DERIVED-UNDER-015:\allowbreak{} the image filtration before deriving}{3.  the image filtration before deriving}}\label{proofs-8352-8772-derived-under-015-the-image-filtration-before-deriving}

There is also an additive extension to all finite filtered objects,\allowbreak{} without left exactness. Keep the original T(A)\allowbreak{} and define \[
 T_{\rm im}(A,F)=
 \left(T(A),\
 F^pT_{\rm im}(A)=
 \operatorname{im}\big(T(F^pA)\longrightarrow T(A)\big)\right).
                                                               \tag{3.1}
\] The map in (3.1)\allowbreak{} is T of the actual filtration inclusion. The images decrease with p since the inclusions factor. If F\^{}a A\allowbreak{}=A and F\textasciicircum{}\{b\allowbreak{}+1\}A\allowbreak{}=0,\allowbreak{} the corresponding maps are an identity and a zero map,\allowbreak{} respectively. Hence (3.1)\allowbreak{} has the same permitted finite bounds [a,\allowbreak{}b]\allowbreak{}.

For a filtered map f:\allowbreak{}A\allowbreak{}->B the T-image of its filtration square maps the subobject in (3.1)\allowbreak{} into the corresponding subobject for B. This defines T\_\allowbreak{}im(f)\allowbreak{} by restricting the original map T(f)\allowbreak{}. It preserves identities and composition because T does. It is additive on Hom groups:\allowbreak{} the sum of two maps preserving these fixed subobjects preserves them,\allowbreak{} and its underlying map is T(f)\allowbreak{}\allowbreak{}+T(g)\allowbreak{}\allowbreak{}=T(f\allowbreak{}+g)\allowbreak{}. Images of the finite direct-sum inclusions show that finite biproducts and their filtrations are preserved. This proves that T\_\allowbreak{}im is an additive functor Fil\textasciicircum{}f(A)\allowbreak{}\allowbreak{}->Fil\textasciicircum{}f(B)\allowbreak{}. If T is left exact,\allowbreak{} the filtration maps T(F\^{}p A)\allowbreak{}\allowbreak{}->T(A)\allowbreak{} are monic,\allowbreak{} and (3.1)\allowbreak{} is exactly the source\textquotesingle s extension at 8375–\allowbreak{}8378. On split objects it is the canonical T\_\allowbreak{}ext of section 2 for every additive T.

Applying T\_\allowbreak{}im to complexes preserves differentials,\allowbreak{} filtered homotopies,\allowbreak{} shifts and cone matrices because it is additive. It consequently gives an exact functor on the filtered homotopy categories. Assume A has enough injectives and let u:\allowbreak{}K\allowbreak{}->I be a bounded-below filtered injective resolution. For any filtered quasi-isomorphism s:\allowbreak{}K\allowbreak{}->L,\allowbreak{} the filtered Hom comparison gives a unique homotopy class v:\allowbreak{}L\allowbreak{}->I with vs\allowbreak{}=u. The map v is a filtered quasi-isomorphism by two-out-of-three. Between filtered injective resolutions these comparison maps are homotopy equivalences,\allowbreak{} with inverse and homotopy identities obtained by the same uniqueness.

Thus the denominator diagram defining the right derived functor of T\_\allowbreak{}im has the same essential-constancy witness as in the injective-computation proof already established:\allowbreak{} refine to I,\allowbreak{} where every further injective comparison is invertible up to filtered homotopy. The actual derived value is \[
 R(T_{\rm im})(K)=T_{\rm im}(I)=T_{\rm ext}(I)
       \quad\text{in }DF^+(B),                              \tag{3.2}
\] with comparison induced by T\_\allowbreak{}im(u)\allowbreak{}. The witness proves existence and the canonical derived comparison; (3.2)\allowbreak{} is not merely a coincidence of object values. Its naturality on roofs follows from the comparison equation vs\allowbreak{}=u and its uniqueness. Hence it agrees with the source\textquotesingle s RT\allowbreak{}=T\_\allowbreak{}ext j',\allowbreak{} including the implicit localization from K\textasciicircum{}\allowbreak{}+(Fil\textasciicircum{}f B)\allowbreak{} to DF\textasciicircum{}\allowbreak{}+(B)\allowbreak{}.

This extension does not assume that T preserves monomorphisms or exact sequences in A. For a strict example use T(M)\allowbreak{}\allowbreak{}=M direct-sum (M/\allowbreak{}ell M)\allowbreak{} from PART10 section 2. Applied to the monomorphism ell:\allowbreak{}Z\allowbreak{}->Z,\allowbreak{} its second component is zero on Z/\allowbreak{}ell Z,\allowbreak{} so T is not left exact. Formula (3.1)\allowbreak{} nevertheless supplies the filtration by the actual images,\allowbreak{} including that lost component,\allowbreak{} and (3.2)\allowbreak{} supplies its derived functor. No equality with T(gr A)\allowbreak{} on an arbitrary underived filtration is asserted.

\subsection{\texorpdfstring{4. Derived associated graded objects,\allowbreak{} convergence,\allowbreak{} and naturality}{4. Derived associated graded   and naturality}}\label{proofs-8352-8772-derived-associated-graded-objects-convergence-and-naturality}

Let K\allowbreak{}->I be a filtered injective resolution as above. Every gr\^{}p I and the underlying I are bounded-below complexes of injectives,\allowbreak{} and the corresponding maps from gr\^{}p K and underlying K are quasi-isomorphisms. Formula (2.2)\allowbreak{} is natural in the maps of filtered objects,\allowbreak{} hence commutes with the original differentials of I. It gives the natural isomorphisms of derived functors \[
 \operatorname{gr}^p RT(K)\simeq RT(\operatorname{gr}^p K),
 \qquad
 \operatorname{forget}RT(K)\simeq RT(\operatorname{forget}K).
                                                               \tag{4.1}
\] The RT on the right is the unfiltered one. The naturality on arbitrary derived maps is checked on the injective comparison maps; inverting a denominator inverts that same commuting square. Consequently the statement is independent of j\textquotesingle{} and of chosen splittings. RECON-064 /\allowbreak{} DERIVED-061 is the existing MC-STK-ERR-0872 singular ``Lemma'' at 8447.

Put C\allowbreak{}=T\_\allowbreak{}ext(I)\allowbreak{},\allowbreak{} with its actual finite term filtrations. Its filtered-complex spectral sequence has \[
 E_0^{p,q}=\operatorname{gr}^p C^{p+q},\qquad
 E_1^{p,q}=H^{p+q}(\operatorname{gr}^p C)
          =R^{p+q}T(\operatorname{gr}^p K),                  \tag{4.2}
\] and differential of bidegree (r,\allowbreak{}1-r)\allowbreak{}. The first equality is the definition of the filtered spectral sequence,\allowbreak{} while the last is (2.2)\allowbreak{} applied to the injective resolution of gr\^{}p K. The boundedness,\allowbreak{} convergence and finite filtrations follow directly from the finite term filtrations,\allowbreak{} with the complete explicit calculation in section 5 below. The abutment is H\textasciicircum{}*(C)\allowbreak{}\allowbreak{}=R\textasciicircum{}*T(K)\allowbreak{}; each degree H\textasciicircum{}m(C)\allowbreak{} has its own finite filtration. RECON-065 /\allowbreak{} DERIVED-062 reconciles both MC-STK-ERR-0873 changes at 8530–\allowbreak{}8531. They state the graded abutment and its per-degree filtrations precisely; no different cohomology object is introduced.

For a filtered chain map K\allowbreak{}->L,\allowbreak{} the filtered Hom comparison provides a unique homotopy class I\allowbreak{}->J between their filtered injective resolutions,\allowbreak{} compatible with the augmentations. T\_\allowbreak{}ext,\allowbreak{} equivalently T\_\allowbreak{}im on these objects,\allowbreak{} preserves filtered homotopies. Such a homotopy induces a homotopy on the associated graded complexes,\allowbreak{} so equal maps on E\_\allowbreak{}1. Equality on each later page follows by passing to cohomology of the preceding page. The same uniqueness proves identity and composition. Thus the sequence is functorial and independent of choices from E\_\allowbreak{}1,\allowbreak{} exactly as stated at 8501. It need not have a choice-independent E\_\allowbreak{}0. For example,\allowbreak{} resolve Q[0]\allowbreak{} either by Q[0]\allowbreak{} or by Q[0]\allowbreak{} plus the contractible complex Q\allowbreak{}->Q in degrees 1,\allowbreak{}2,\allowbreak{} with a single filtration step. The E\_\allowbreak{}0 complexes differ whereas their E\_\allowbreak{}1 and later pages agree.

All of this argument uses additivity and split filtrations on I,\allowbreak{} and therefore also proves the additive version of the filtered spectral- sequence lemma. The source itself already announces additive RT at 8407. This propagates DERIVED-UNDER-011 through its actual filtered construction,\allowbreak{} with no suggestion of novelty.

\subsection{\texorpdfstring{5. DERIVED-UNDER-016:\allowbreak{} explicit local stabilization pages}{5.  explicit local stabilization pages}}\label{proofs-8352-8772-derived-under-016-explicit-local-stabilization-pages}

Let (C,\allowbreak{}F)\allowbreak{} be any complex in an abelian category,\allowbreak{} with a finite filtration on every term. For each m fix integers a\_\allowbreak{}m<\allowbreak{}=b\_\allowbreak{}m satisfying F\textasciicircum{}\{a\_\allowbreak{}m\}C\textasciicircum{}m\allowbreak{}=C\textasciicircum{}m and F\textasciicircum{}\{b\_\allowbreak{}m\allowbreak{}+1\}C\textasciicircum{}m\allowbreak{}=0. Zero terms may be given any such bounds. No lower or upper bound on the complex itself is required for this calculation.

For r>\allowbreak{}=1 use the representatives \[
 \mathcal Z_r^{p,m}
     =F^pC^m\cap(d^m)^{-1}(F^{p+r}C^{m+1}).
                                                               \tag{5.1}
\] The source\textquotesingle s quotient formula is equivalently \[
 E_r^{p,m-p}=
 \frac{\mathcal Z_r^{p,m}}
 {\mathcal Z_{r-1}^{p+1,m}
       +d^{m-1}\mathcal Z_{r-1}^{p-r+1,m-1}}.                 \tag{5.2}
\] Both terms in its denominator lie in the numerator:\allowbreak{} the first by its differential condition,\allowbreak{} and the second because it lies in F\^{}p and is a boundary. This formula is obtained from the source\textquotesingle s Z\_\allowbreak{}r and B\_\allowbreak{}r inside F\textasciicircum{}p/\allowbreak{}F\textasciicircum{}\{p\allowbreak{}+1\} by taking representatives and then intersecting the added F\textasciicircum{}\{p\allowbreak{}+1\} with the numerator. No quotient component is discarded.

If \[
 r\ge R(p,m):=
 \max\{1,\ b_{m+1}-p+1,\ p-a_{m-1}+1\},                    \tag{5.3}
\] then F\textasciicircum{}\{p\allowbreak{}+r\}C\textasciicircum{}\{m\allowbreak{}+1\}\allowbreak{}=0 and F\textasciicircum{}\{p-r\allowbreak{}+1\}C\textasciicircum{}\{m-1\}\allowbreak{}=C\textasciicircum{}\{m-1\}. Write Z\textasciicircum{}m\allowbreak{}=ker d\^{}m and B\textasciicircum{}m\allowbreak{}=im d\^{}\{m-1\}. The numerator and denominator in (5.2)\allowbreak{} become exactly \[
 E_r^{p,m-p}=
 \frac{F^pC^m\cap Z^m}
 {(F^{p+1}C^m\cap Z^m)+(F^pC^m\cap B^m)}.                  \tag{5.4}
\] For the second denominator term,\allowbreak{} the image of the inverse image of F\^{}p under d\^{}\{m-1\} is F\^{}p intersect B\^{}m. This identity follows by pulling back the epimorphism C\textasciicircum{}\{m-1\}\allowbreak{}->B\textasciicircum{}m along that intersection,\allowbreak{} so is valid in any abelian category.

The induced filtration on H\textasciicircum{}m(C)\allowbreak{}\allowbreak{}=Z\textasciicircum{}m/\allowbreak{}B\textasciicircum{}m is the image of F\^{}pC\^{}m intersect Z\^{}m. Its associated graded object is precisely (5.4)\allowbreak{},\allowbreak{} by the subobject-quotient isomorphisms. The numerator and both denominator terms remain identical for every larger r. Thus (5.3)\allowbreak{} is an explicit stabilization page,\allowbreak{} with the canonical identification E\_\allowbreak{}r\textasciicircum{}\{p,\allowbreak{}m-p\}\allowbreak{}=gr\textasciicircum{}p H\textasciicircum{}m(C)\allowbreak{}. It proves the convergence comparison itself,\allowbreak{} not only equality of dimensions.

Only a\_\allowbreak{}m<\allowbreak{}=p<\allowbreak{}=b\_\allowbreak{}m can contribute in degree m. Taking a maximum of (5.3)\allowbreak{} over that finite interval gives the single sufficient page \[
 R_m=\max\{1,\ b_{m+1}-a_m+1,\
                   b_m-a_{m-1}+1\}.                         \tag{5.5}
\] The filtration on H\textasciicircum{}m(C)\allowbreak{} has at most b\_\allowbreak{}m-a\_\allowbreak{}m\allowbreak{}+1 potentially nonzero graded pieces:\allowbreak{} it is all of H\^{}m at a\_\allowbreak{}m and zero at b\_\allowbreak{}m\allowbreak{}+1,\allowbreak{} as follows directly from its cycle representatives. It is consequently exhaustive,\allowbreak{} separated and complete; the inverse system of its quotients is eventually H\^{}m itself. Together with (5.3)\allowbreak{} this proves convergence in the source\textquotesingle s exact sense at Homology 6406–\allowbreak{}6421. Boundedness there means finitely many E\_\allowbreak{}0 terms on each total-degree diagonal,\allowbreak{} which is also given by [a\_\allowbreak{}m,\allowbreak{}b\_\allowbreak{}m]\allowbreak{}.

If every term has a common interval [a,\allowbreak{}b]\allowbreak{},\allowbreak{} all degrees stabilize by page b-a\allowbreak{}+1 and each cohomology filtration has at most b-a\allowbreak{}+1 graded pieces. No complex-degree bound is needed for this uniform-window assertion. For varying intervals,\allowbreak{} the adjacent degrees m-1 and m\allowbreak{}+1 in (5.3)\allowbreak{} and (5.5)\allowbreak{} are essential to the calculation and have been retained.

For the filtered derived sequence in section 4,\allowbreak{} PART11 supplies a resolution zero below N whose term in degree m>\allowbreak{}=N has interval \[
 A_m=\min_{N\le j\le m}a_j,\qquad
 B_m=\max_{N\le j\le m}b_j
                                                               \tag{5.6}
\] when [a\_\allowbreak{}j,\allowbreak{}b\_\allowbreak{}j]\allowbreak{} are the input term intervals. T\_\allowbreak{}ext preserves those intervals. Apply (5.3)\allowbreak{} and (5.5)\allowbreak{} with A\_\allowbreak{}m,\allowbreak{}B\_\allowbreak{}m. For the zero terms below N choose A\_\allowbreak{}m\allowbreak{}=A\_\allowbreak{}N and B\_\allowbreak{}m\allowbreak{}=B\_\allowbreak{}N,\allowbreak{} which are valid bounds on zero. This gives an explicit page and filtration length from the input data,\allowbreak{} with their actual integer indices intact.

\subsection{\texorpdfstring{6. The two Cartan–\allowbreak{}Eilenberg sequences and their page shift}{6. The two  sequences and their page shift}}\label{proofs-8352-8772-the-two-cartaneilenberg-sequences-and-their-page-shift}

Original 8558–\allowbreak{}8617 uses the stupid filtration S\^{}pK \allowbreak{}=sigma\_\allowbreak{}\{>\allowbreak{}=p\}K and the canonical filtration G\textasciicircum{}sK\allowbreak{}=tau\_\allowbreak{}\{<\allowbreak{}=-s\}K. The complete comparison was proved in PART10 section 3,\allowbreak{} including the injective filtrations on the actual total complex and the formula \[
 E_r^{s,t}(G)=
 {}''E_{r+1}^{2s+t,-s}\quad(r\ge1).                          \tag{6.1}
\] The new review in sections 2–\allowbreak{}4 confirms that the filtered derived functor used there has the precise associated-graded and forgetful comparisons (4.1)\allowbreak{}. For S,\allowbreak{} the graded piece is K\textasciicircum{}p[-p]\allowbreak{},\allowbreak{} so its derived E\_\allowbreak{}1 is R\textasciicircum{}qT(K\textasciicircum{}p)\allowbreak{}. For G,\allowbreak{} the graded piece is quasi-isomorphic to H\textasciicircum{}\{-s\}(K)\allowbreak{}[s]\allowbreak{},\allowbreak{} so its derived E\_\allowbreak{}1 is R\textasciicircum{}\{2s\allowbreak{}+t\}T(H\textasciicircum{}\{-s\}(K)\allowbreak{})\allowbreak{}. The exact cycle and boundary quotient calculation establishing (6.1)\allowbreak{} retains the original total differential d\_\allowbreak{}1\allowbreak{}+(-1)\allowbreak{}\textasciicircum{}p d\_\allowbreak{}2. It identifies the differentials as well:\allowbreak{} bidegree (r,\allowbreak{}1-r)\allowbreak{} becomes (r\allowbreak{}+1,\allowbreak{}-r)\allowbreak{}. Thus the source\textquotesingle s two promised sequences agree from their canonical E\_\allowbreak{}1 and E\_\allowbreak{}2 pages,\allowbreak{} respectively.

The split horizontal kernel sequences in the Cartan–\allowbreak{}Eilenberg resolution make this argument valid for additive T as proved in PART10. We do not assume that arbitrary underived kernels commute with T. RECON-066 /\allowbreak{} P02-E0065 /\allowbreak{} DERIVED-059 is the existing MC-STK-ERR-0870 insertion of ``is''. RECON-067 /\allowbreak{} DERIVED-060 is MC-STK-ERR-0871,\allowbreak{} the existing idiom correction at 8603.

\subsection{7. Ext groups and their exact boundary maps}\label{proofs-8352-8772-ext-groups-and-their-exact-boundary-maps}

Keep the source\textquotesingle s definition Ext\textasciicircum{}i\_\allowbreak{}A(X,\allowbreak{}Y)\allowbreak{}\allowbreak{}=Hom\_\allowbreak{}D(A)\allowbreak{}(X,\allowbreak{}Y[i]\allowbreak{})\allowbreak{}. The equality with Hom\_\allowbreak{}D(A)\allowbreak{}(X[-i]\allowbreak{},\allowbreak{}Y)\allowbreak{} is the natural bijection given by the inverse shift autoequivalence. For a triangle Y\allowbreak{}->Y'\allowbreak{}->Y'' --h-\allowbreak{}->Y[1]\allowbreak{},\allowbreak{} the connecting map on Ext sends f:\allowbreak{}X\allowbreak{}->Y''[i]\allowbreak{} to h[i]\allowbreak{}f:\allowbreak{}X\allowbreak{}->Y[i\allowbreak{}+1]\allowbreak{}. The long exact Hom sequence for that triangle proves exactness with the displayed objects. For a triangle X\allowbreak{}->X'\allowbreak{}->X'' --k-\allowbreak{}->X[1]\allowbreak{},\allowbreak{} the connecting map sends f:\allowbreak{}X\allowbreak{}->Y[i]\allowbreak{} to f[1]\allowbreak{}k:\allowbreak{}X''\allowbreak{}->Y[i\allowbreak{}+1]\allowbreak{}. This is the shifted contravariant Hom boundary. These are the actual maps in 8654–\allowbreak{}8666. Fullness of the bounded derived subcategories allows the same computation when both inputs belong to one of them,\allowbreak{} since their shifts stay there.

If Y\allowbreak{}->I is a bounded-below injective resolution,\allowbreak{} the target I[i]\allowbreak{} remains such an injective complex,\allowbreak{} with lower bound shifted by -i and differential (-1)\allowbreak{}\textasciicircum{}i d\_\allowbreak{}I. The injective Hom comparison identifies Ext\textasciicircum{}i(X,\allowbreak{}Y)\allowbreak{} with Hom\_\allowbreak{}K(A)\allowbreak{}(X,\allowbreak{}I[i]\allowbreak{})\allowbreak{}. The source X need not have any term bound. Dually,\allowbreak{} for a bounded-above projective resolution P\allowbreak{}->X,\allowbreak{} P[-i]\allowbreak{} has upper bound shifted by i and its full shift differential. The projective Hom comparison gives Ext\textasciicircum{}i(X,\allowbreak{}Y)\allowbreak{}\allowbreak{}= Hom\_\allowbreak{}K(A)\allowbreak{}(P[-i]\allowbreak{},\allowbreak{}Y)\allowbreak{} without a bound on Y. This proves 8681–\allowbreak{}8706 with the source\textquotesingle s opposite bounds.

\subsection{8. The exact lowest Ext degree}\label{proofs-8352-8772-the-exact-lowest-ext-degree}

Assume H\textasciicircum{}j(X)\allowbreak{}\allowbreak{}=0 for j\textgreater a and H\textasciicircum{}j(Y)\allowbreak{}\allowbreak{}=0 for j\textless b. Use the original canonical truncations:\allowbreak{} Y\allowbreak{}->Y\_\allowbreak{}b\allowbreak{}=tau\_\allowbreak{}\{>\allowbreak{}=b\}Y is a quasi-isomorphism; a denominator L\allowbreak{}->X may be replaced by L\_\allowbreak{}a\allowbreak{}=tau\_\allowbreak{}\{<\allowbreak{}=a\}L\allowbreak{}->X. Thus every derived map X\allowbreak{}->Y[n]\allowbreak{} has a roof whose numerator is a chain map L\_\allowbreak{}a\allowbreak{}->Y\_\allowbreak{}b[n]\allowbreak{}. Here L\_\allowbreak{}a\textasciicircum{}j\allowbreak{}=0 for j>a,\allowbreak{} and Y\_\allowbreak{}b\textasciicircum{}j\allowbreak{}=0 for j\textless b. For n\textless b-a and j<\allowbreak{}=a,\allowbreak{} j\allowbreak{}+n<b,\allowbreak{} so the target component is zero. For j\textgreater a the source component is zero. All components of the numerator vanish,\allowbreak{} proving Ext\textasciicircum{}n(X,\allowbreak{}Y)\allowbreak{}\allowbreak{}=0.

At n\allowbreak{}=b-a only the component in degree a can be nonzero. The chain equations make it kill im d\_\allowbreak{}\{L\_\allowbreak{}a\}\textasciicircum{}\{a-1\} and land in ker d\_\allowbreak{}\{Y\_\allowbreak{}b\}\textasciicircum{}b. The target shift multiplies that last differential by (-1)\allowbreak{}\textasciicircum{}\{b-a\},\allowbreak{} which has the same kernel; the sign has not been removed from the shifted complex. Therefore the map factors uniquely as \[
 L_a^a\longrightarrow H^a(L_a)
 \xrightarrow{\bar f} H^b(Y_b)
 \longrightarrow Y_b^b.                                    \tag{8.1}
\] There are no nonzero homotopies between such maps:\allowbreak{} a homotopy component in degree j has source L\_\allowbreak{}a\textasciicircum{}j and target Y\_\allowbreak{}b\textasciicircum{}\{j-1\allowbreak{}+b-a\}; for j<\allowbreak{}=a the target is below b,\allowbreak{} and for j\textgreater a the source is zero.

The denominator induces H\textasciicircum{}a(L\_\allowbreak{}a)\allowbreak{}\textasciitilde{}\allowbreak{}=H\textasciicircum{}a(X)\allowbreak{}. Consequently a roof gives the map bar f H\textasciicircum{}a(s)\allowbreak{}\textasciicircum{}\{-1\} from H\textasciicircum{}a(X)\allowbreak{} to H\textasciicircum{}b(Y)\allowbreak{}. Refining the roof preserves that map,\allowbreak{} by the cohomology equation of the refining denominator. Conversely,\allowbreak{} given H\textasciicircum{}a(X)\allowbreak{}\allowbreak{}->H\textasciicircum{}b(Y)\allowbreak{},\allowbreak{} take tau\_\allowbreak{}\{<\allowbreak{}=a\}X,\allowbreak{} map its top term to H\textasciicircum{}a(X)\allowbreak{},\allowbreak{} follow the given map and the cycle inclusion into tau\_\allowbreak{}\{>\allowbreak{}=b\}Y[b-a]\allowbreak{},\allowbreak{} and use zero maps in every other degree. It is a chain map by (8.1)\allowbreak{} and defines a roof with precisely the given cohomology map. A roof and this constructed one agree after a common denominator,\allowbreak{} because their only component is the unique factorization (8.1)\allowbreak{}. This proves the natural isomorphism \[
 \operatorname{Ext}^{b-a}_A(X,Y)
   \xrightarrow{\sim}
 \operatorname{Hom}_A(H^a(X),H^b(Y)).                         \tag{8.2}
\] No injectives,\allowbreak{} projectives,\allowbreak{} or replacement objects with altered support are assumed.

RECON-068 /\allowbreak{} P02-E0067 /\allowbreak{} DERIVED-064 retains MC-STK-ERR-0875:\allowbreak{} after tau\_\allowbreak{}\{<\allowbreak{}=a\},\allowbreak{} the vanishing is L\textasciicircum{}j\allowbreak{}=0 for j>a,\allowbreak{} not j\textless a. The report\textquotesingle s mention of cohomology is supporting intuition; the corrected printed statement is about the terms,\allowbreak{} exactly as needed in this proof. The case a\allowbreak{}=b\allowbreak{}=0 gives Ext\textasciicircum{}n(B,\allowbreak{}A)\allowbreak{}\allowbreak{}=0 for n\textless0 and Ext\textasciicircum{}0(B,\allowbreak{}A)\allowbreak{}\allowbreak{}=Hom(B,\allowbreak{}A)\allowbreak{}.

Finally a short exact sequence of objects gives its canonical triangle in D(A)\allowbreak{}. Applying the two Hom sequences of section 7 and this negative-degree vanishing produces the two displayed long exact sequences at 8753–\allowbreak{}8765,\allowbreak{} including their initial zero terms. They therefore exist without enough injectives or enough projectives. The Yoneda construction following 8772 is the next substantive review,\allowbreak{} not assumed as proof of the assertions above.

\subsection{9. Receiving use and the next review}\label{proofs-8352-8772-receiving-use-and-the-next-review}

For the bounded-below ringed-space construction in current Cohomology 7307–\allowbreak{}7339,\allowbreak{} take its actual injective resolution I and T\allowbreak{}=Gamma(X,\allowbreak{}-)\allowbreak{}. The split filtration on each I\^{}m makes Gamma(X,\allowbreak{}F\textasciicircum{}p I\textasciicircum{}m)\allowbreak{} the subobject specified in (3.1)\allowbreak{}. If the input intervals are [a\_\allowbreak{}j,\allowbreak{}b\_\allowbreak{}j]\allowbreak{},\allowbreak{} PART11\textquotesingle s construction and (5.6)\allowbreak{} give explicit intervals [A\_\allowbreak{}m,\allowbreak{}B\_\allowbreak{}m]\allowbreak{} on Gamma(X,\allowbreak{}I\textasciicircum{}m)\allowbreak{}. Equation (5.5)\allowbreak{} gives its stabilization page,\allowbreak{} and (5.4)\allowbreak{} identifies its graded abutment. For a common input interval [a,\allowbreak{}b]\allowbreak{} the page is b-a\allowbreak{}+1. This strengthens the receiver\textquotesingle s qualitative finite- filtration conclusion. Its unbounded K-injective variant at 7286–\allowbreak{}7305 has separate hypotheses; this review does not assert that its filtration intervals satisfy the bounded-below construction (5.6)\allowbreak{}.

The original continuation through 8930 was read but its Yoneda proofs are not yet adjudicated. In particular,\allowbreak{} the extension-composition display at 8913–\allowbreak{}8929 appears to reverse the endpoint variance relative to the following extensions. The next review must check the whole composition passage and its sign and roof conventions before recording an operation. This is a routed observation,\allowbreak{} not a counted correction or an accepted Yoneda comparison.
\clearpage
\section{\texorpdfstring{Yoneda extensions,\allowbreak{} composition signs,\allowbreak{} splitting,\allowbreak{} and K-groups}{Yoneda  composition   and K-groups}}\label{proofs-8774-9196-yoneda-extensions-composition-signs-splitting-and-k-groups}

This is a private source-bound review of original Derived Categories lines 8774–\allowbreak{}9196,\allowbreak{} with a receiving correction in Cohomology. It continues PART12. The authority is Stacks commit a04446e5\allowbreak{}7ec1fbc2\allowbreak{}52a871af\allowbreak{}cec7752f\allowbreak{}b2807b14; the receiving public edition is 95a51593\allowbreak{}aceb247a\allowbreak{}c7a18311\allowbreak{}1678c3a0\allowbreak{}f8b8f4b5. The accompanying reading and operation records bind the exact bytes. This note proves the arguments used for review; it does not replace the translated source by its editorial strengthenings. No novelty is claimed.

Human source:\allowbreak{} \href{https://github.com/stacks/stacks-project/blob/a04446e57ec1fbc252a871afcec7752fb2807b14/derived.tex\#L8774}{The Stacks Project,\allowbreak{} Derived Categories,\allowbreak{} original source}. The \href{https://stacks.math.columbia.edu/tag/06XP}{current online Ext section} was checked as a secondary comparison on 2026-09-30; it does not replace the frozen authority. The arguments below use the original TeX,\allowbreak{} not a PDF.

\subsection{1. The exact pushout construction and common refinement}\label{proofs-8774-9196-the-exact-pushout-construction-and-common-refinement}

For a degree n extension E of B by A,\allowbreak{} put A in degree -n,\allowbreak{} its successive middle objects in degrees -n\allowbreak{}+1 through 0,\allowbreak{} and use its actual arrows as the differentials. Write C\_\allowbreak{}E for this complex,\allowbreak{} s\_\allowbreak{}E:\allowbreak{}C\_\allowbreak{}E\allowbreak{}→B[0]\allowbreak{} for the augmentation,\allowbreak{} and f\_\allowbreak{}E:\allowbreak{}C\_\allowbreak{}E\allowbreak{}→A[n]\allowbreak{} for the projection with component id\_\allowbreak{}A in degree -n. Thus δ(E)\allowbreak{}\allowbreak{}=f\_\allowbreak{}E s\_\allowbreak{}E\textasciicircum{}\{-1\}. The identity component is the natural projection in the source\textquotesingle s display; retaining it matters for Section 3.

Suppose s:\allowbreak{}L\allowbreak{}→B[0]\allowbreak{} is a quasi-isomorphism and f:\allowbreak{}L\allowbreak{}→A[n]\allowbreak{} is a chain map. Replace L by τ\_\allowbreak{}\{\textbackslash{}leq0\}L,\allowbreak{} restricting both maps. This keeps the represented morphism and makes L\textasciicircum{}r\allowbreak{}=0 for r\textgreater0. In particular H\textasciicircum{}r(L)\allowbreak{}\allowbreak{}=0 for r<0,\allowbreak{} and L\textasciicircum{}0\allowbreak{}→B is an epimorphism with kernel im(d\_\allowbreak{}L\textasciicircum{}\{-1\})\allowbreak{}. The component f\^{}\{-n\} kills im(d\_\allowbreak{}L\textasciicircum{}\{-n-1\})\allowbreak{}. Therefore it factors through Q\allowbreak{}=L\textasciicircum{}\{-n\}/\allowbreak{}im(d\_\allowbreak{}L\textasciicircum{}\{-n-1\})\allowbreak{}. Exactness at degree -n identifies Q with im(d\_\allowbreak{}L\textasciicircum{}\{-n\})\allowbreak{} and gives a monomorphism j:\allowbreak{}Q\allowbreak{}→L\textasciicircum{}\{-n\allowbreak{}+1\}.

Form the pushout

P\allowbreak{}=coker((d\_\allowbreak{}L\textasciicircum{}\{-n\},\allowbreak{}-f\textasciicircum{}\{-n\})\allowbreak{}:\allowbreak{}L\textasciicircum{}\{-n\}\allowbreak{}→L\textasciicircum{}\{-n\allowbreak{}+1\}\allowbreak{}⊕A)\allowbreak{}.

Equivalently,\allowbreak{} push j out along Q\allowbreak{}→A. The resulting map A\allowbreak{}→P,\allowbreak{} a\allowbreak{}↦[0,\allowbreak{}a]\allowbreak{},\allowbreak{} is a monomorphism,\allowbreak{} and coker(A\allowbreak{}→P)\allowbreak{}\allowbreak{}=coker(j)\allowbreak{}. This follows directly from the pushout of the exact sequence 0\allowbreak{}→Q\allowbreak{}→L\textasciicircum{}\{-n\allowbreak{}+1\}\allowbreak{}→coker(j)\allowbreak{}\allowbreak{}→0 in an abelian category. For n\textgreater1 the next map is [x,\allowbreak{}a]\allowbreak{}\allowbreak{}↦d\_\allowbreak{}L\textasciicircum{}\{-n\allowbreak{}+1\}x; it is well defined and its kernel is the image of A. The remaining terms are L\textasciicircum{}\{-n\allowbreak{}+2\},\allowbreak{}…,\allowbreak{}L\textasciicircum{}0 followed by B. For n\allowbreak{}=1 the next map is [x,\allowbreak{}a]\allowbreak{}\allowbreak{}↦s\textasciicircum{}0(x)\allowbreak{},\allowbreak{} and the same pushout sequence identifies its cokernel as zero and its kernel as A. Thus both cases give an exact degree n extension E(f,\allowbreak{}s)\allowbreak{}.

There is a chain map u:\allowbreak{}L\allowbreak{}→C\_\allowbreak{}\{E(f,\allowbreak{}s)\allowbreak{}\}. Its components are zero below -n,\allowbreak{} f\^{}\{-n\} in degree -n,\allowbreak{} x\allowbreak{}↦[x,\allowbreak{}0]\allowbreak{} in degree -n\allowbreak{}+1,\allowbreak{} and identities above that. The essential equation is {[}d x,\allowbreak{}0]\allowbreak{}\allowbreak{}=[0,\allowbreak{}f x]\allowbreak{},\allowbreak{} which is precisely the relation in the pushout. All other chain equations follow from f d\allowbreak{}=0 and d²\allowbreak{}=0. The equalities s\_\allowbreak{}E u\allowbreak{}=s and f\_\allowbreak{}E u\allowbreak{}=f are strict. On H\^{}0 the first equality shows that u is an isomorphism,\allowbreak{} and all other cohomology objects on either side vanish. Hence u is a quasi-isomorphism and δ(E(f,\allowbreak{}s)\allowbreak{})\allowbreak{}\allowbreak{}=f s\^{}\{-1\}. This proves the surjectivity in the source\textquotesingle s Yoneda classification.

For injectivity,\allowbreak{} let E,\allowbreak{}E' have equal classes. The common-denominator description of the localization of K(A)\allowbreak{} gives a complex L and quasi-isomorphisms t:\allowbreak{}L\allowbreak{}→C\_\allowbreak{}E and t':\allowbreak{}L\allowbreak{}→C\_\allowbreak{}\{E'\} such that s\_\allowbreak{}E t\allowbreak{}=s\_\allowbreak{}\{E'\}t' and f\_\allowbreak{}E t\allowbreak{}=f\_\allowbreak{}\{E'\}t' in K(A)\allowbreak{}. Replace L by τ\_\allowbreak{}\{\textbackslash{}leq0\}L. A homotopy from L to B[0]\allowbreak{} could contribute to degree 0 only through a map L\textasciicircum{}1\allowbreak{}→B. That term is zero. Consequently the first equality is already a strict equality of chain maps.

Set α\allowbreak{}=t\textasciicircum{}\{-n\}:\allowbreak{}L\textasciicircum{}\{-n\}\allowbreak{}→A and α'\allowbreak{}=t'\textasciicircum{}\{-n\}. The second equality says α-α'\allowbreak{}=h d\_\allowbreak{}L\textasciicircum{}\{-n\} for a map h:\allowbreak{}L\textasciicircum{}\{-n\allowbreak{}+1\}\allowbreak{}→A. There are no other possibly nonzero components in a homotopy to A[n]\allowbreak{}. Let ι':\allowbreak{}A\allowbreak{}→Z'\_\allowbreak{}\{n-1\} denote the first arrow of E\textquotesingle. Modify t\textquotesingle{} by the homotopy whose only component is h:\allowbreak{}

t̃'\textasciicircum{}\{-n\}\allowbreak{}=α,\allowbreak{} t̃'\textasciicircum{}\{-n\allowbreak{}+1\}\allowbreak{}=t'\textasciicircum{}\{-n\allowbreak{}+1\}\allowbreak{}+ι'h,\allowbreak{}

and keep every other component equal to t\textquotesingle. These formulas are t̃'\allowbreak{}=t'\allowbreak{}+d h\allowbreak{}+h d. They remain valid for n\allowbreak{}=1:\allowbreak{} the augmentation kills ι'h,\allowbreak{} so it is unchanged. Thus t̃\textquotesingle{} is a chain map,\allowbreak{} is homotopic to t',\allowbreak{} and has the same augmentation as t. Both maps now have the same top component α strictly.

Apply the pushout construction to (α,\allowbreak{}s\_\allowbreak{}E t)\allowbreak{}. The map from its first middle term to Z\_\allowbreak{}\{n-1\} is [x,\allowbreak{}a]\allowbreak{}\allowbreak{}↦t\textasciicircum{}\{-n\allowbreak{}+1\}x\allowbreak{}+ι a. It kills the relation (d x,\allowbreak{}-α x)\allowbreak{} by the chain equation for t. The analogous formula using t̃\textquotesingle{} maps to Z'\_\allowbreak{}\{n-1\}. The later components are those of t and t̃',\allowbreak{} and both endpoint maps are identities. This constructs the actual common middle extension in the definition,\allowbreak{} including the homotopy adjustment that the source leaves to the reader. Conversely,\allowbreak{} a map of extensions fixing A,\allowbreak{}B commutes with s and f and is a quasi-isomorphism,\allowbreak{} so it preserves δ. The classification follows,\allowbreak{} and the stated common-refinement relation is reflexive,\allowbreak{} symmetric and transitive because equality of δ is.

The omitted adjustment is supplied here as an editorial proof. The source explicitly omits details,\allowbreak{} so omission alone is not counted as another mathematical error.

\subsection{2. Naturality and the actual group law in degree one}\label{proofs-8774-9196-naturality-and-the-actual-group-law-in-degree-one}

For a map of extensions from E of B by A to F of B\textquotesingle{} by A\textquotesingle{} with endpoint maps β:\allowbreak{}B\allowbreak{}→B' and α:\allowbreak{}A\allowbreak{}→A',\allowbreak{} the induced chain map v satisfies s\_\allowbreak{}F v\allowbreak{}=βs\_\allowbreak{}E and f\_\allowbreak{}F v\allowbreak{}=α[n]\allowbreak{}f\_\allowbreak{}E. In D(A)\allowbreak{} this proves

δ(F)\allowbreak{}β\allowbreak{}=α[n]\allowbreak{}δ(E)\allowbreak{}.

Pushout of the first monomorphism along α and pullback of the last epimorphism along β exist in an abelian category and preserve the exact extension. The displayed identity proves their usual covariance and contravariance,\allowbreak{} with no assumption about enough injectives or projectives. Direct sums give δ(E\allowbreak{}⊕F)\allowbreak{}\allowbreak{}=δ(E)\allowbreak{}\allowbreak{}⊕δ(F)\allowbreak{}.

In degree one,\allowbreak{} morphisms of extensions fixing the endpoints are isomorphisms:\allowbreak{} apply the kernel and cokernel exact sequences to their commutative diagram with exact rows,\allowbreak{} or the short five lemma. Thus the common-refinement classification is the isomorphism classification in Homology\textquotesingle s definition. Its Baer sum is the direct sum followed by the pushout along Σ:\allowbreak{}A\allowbreak{}⊕A\allowbreak{}→A and pullback along Δ:\allowbreak{}B\allowbreak{}→B\allowbreak{}⊕B. Naturality gives

δ(E\allowbreak{}+F)\allowbreak{}\allowbreak{}=Σ\href{δ(E)⊕δ(F)}{1}Δ\allowbreak{}=δ(E)\allowbreak{}\allowbreak{}+δ(F)\allowbreak{}.

For the split extension A\allowbreak{}→A\allowbreak{}⊕B\allowbreak{}→B,\allowbreak{} projection A\allowbreak{}⊕B\allowbreak{}→A is a homotopy from f\_\allowbreak{}E to zero on its two-term complex,\allowbreak{} so δ(E)\allowbreak{}\allowbreak{}=0. Pushout by -id\_\allowbreak{}A gives additive inverses. This proves the group identification,\allowbreak{} not merely the set bijection. These are details supporting the original lemma.

One sign from an earlier source section must remain explicit. For a short exact sequence 0\allowbreak{}→A\allowbreak{}→Z\allowbreak{}→B\allowbreak{}→0,\allowbreak{} C\_\allowbreak{}E is the source cone of A[0]\allowbreak{}\allowbreak{}→Z[0]\allowbreak{}. Its cone differential is the positive inclusion A\allowbreak{}→Z. The source\textquotesingle s distinguished cone triangle has last map -p,\allowbreak{} so its connecting morphism is

∂\_\allowbreak{}E\allowbreak{}=-p s\_\allowbreak{}E\textasciicircum{}\{-1\}\allowbreak{}=-δ(E)\allowbreak{}.

This is the convention in current derived.tex 4286–\allowbreak{}4326. Both maps have been retained; the proof below never silently identifies them.

\subsection{3. Endpoint variance and the full signed splicing formula}\label{proofs-8774-9196-endpoint-variance-and-the-full-signed-splicing-formula}

Let E be a degree i extension of B by A,\allowbreak{} and E\textquotesingle{} a degree j extension of C by B. Their classes belong to Ext\textasciicircum{}i(B,\allowbreak{}A)\allowbreak{} and Ext\textasciicircum{}j(C,\allowbreak{}B)\allowbreak{}. The appropriate composition is δ(E)\allowbreak{}[j]\allowbreak{}δ(E')\allowbreak{}:\allowbreak{}C[0]\allowbreak{}\allowbreak{}→A[i\allowbreak{}+j]\allowbreak{}. Therefore the source display at 8913–\allowbreak{}8915 must have domain Ext\textasciicircum{}i(B,\allowbreak{}A)\allowbreak{}×Ext\textasciicircum{}j(C,\allowbreak{}B)\allowbreak{} and codomain Ext\textasciicircum{}\{i\allowbreak{}+j\}(C,\allowbreak{}A)\allowbreak{}. This is DERIVED-NEW-028. The preceding display with the general objects X,\allowbreak{}Y,\allowbreak{}Z is correct and is retained.

Let S be the splice with all the original arrows:\allowbreak{}

0\allowbreak{}→A\allowbreak{}→Z\_\allowbreak{}\{i-1\}\allowbreak{}→…\allowbreak{}→Z\_\allowbreak{}0\allowbreak{}→Z'\_\allowbreak{}\{j-1\}\allowbreak{}→…\allowbreak{}→Z'\_\allowbreak{}0\allowbreak{}→C\allowbreak{}→0.

The joining arrow is Z\_\allowbreak{}0\allowbreak{}→B\allowbreak{}→Z'\_\allowbreak{}\{j-1\}. Its kernel and image are exactly the required adjacent image and kernel,\allowbreak{} so S is exact. Let J\allowbreak{}=C\_\allowbreak{}S,\allowbreak{} K\allowbreak{}=C\_\allowbreak{}E,\allowbreak{} and L\allowbreak{}=C\_\allowbreak{}\{E'\}. Define a:\allowbreak{}J\allowbreak{}→L as follows:\allowbreak{} it is the identity on the Z\textquotesingle{} terms in degrees -j\allowbreak{}+1,\allowbreak{}…,\allowbreak{}0,\allowbreak{} the augmentation Z\_\allowbreak{}0\allowbreak{}→B in degree -j,\allowbreak{} and zero in all lower degrees. It is a chain map,\allowbreak{} its augmentation is s\_\allowbreak{}S,\allowbreak{} and it induces the identity on the only nonzero cohomology,\allowbreak{} C in degree 0. Hence a is a quasi-isomorphism.

The shift convention is d\_\allowbreak{}\{K[j]\allowbreak{}\}\textasciicircum{}r\allowbreak{}=(-1)\allowbreak{}\textasciicircum{}j d\_\allowbreak{}K\textasciicircum{}\{r\allowbreak{}+j\}. Define b:\allowbreak{}J\allowbreak{}→K[j]\allowbreak{} by

b\textasciicircum{}r\allowbreak{}=(-1)\allowbreak{}\textasciicircum{}\{j(r\allowbreak{}+j)\allowbreak{}\} id\_\allowbreak{}\{K\textasciicircum{}\{r\allowbreak{}+j\}\} (-i-j≤r≤-j)\allowbreak{},\allowbreak{} b\textasciicircum{}r\allowbreak{}=0 (r>-j)\allowbreak{}.

Outside the displayed support its components are zero. For r\textless-j the chain equation is the scalar identity

(-1)\allowbreak{}\textasciicircum{}j(-1)\allowbreak{}\textasciicircum{}\{j(r\allowbreak{}+j)\allowbreak{}\}\allowbreak{}=(-1)\allowbreak{}\textasciicircum{}\{j(r\allowbreak{}+1\allowbreak{}+j)\allowbreak{}\}.

At r\allowbreak{}=-j both sides of the chain equation are zero. Thus b is a chain map. Its degree -j component is id\_\allowbreak{}\{Z\_\allowbreak{}0\},\allowbreak{} and consequently

s\_\allowbreak{}E[j]\allowbreak{} b\allowbreak{}=f\_\allowbreak{}\{E'\}a.

In the derived category this equality and the invertibility of s\_\allowbreak{}E[j]\allowbreak{} compute the desired composition without omitting any shift:\allowbreak{}

δ(E)\allowbreak{}[j]\allowbreak{}δ(E')\allowbreak{} \allowbreak{}= f\_\allowbreak{}E[j]\allowbreak{} b(s\_\allowbreak{}\{E'\}a)\allowbreak{}\textasciicircum{}\{-1\} \allowbreak{}= (-1)\allowbreak{}\textasciicircum{}\{ij\} f\_\allowbreak{}S s\_\allowbreak{}S\textasciicircum{}\{-1\} \allowbreak{}= (-1)\allowbreak{}\textasciicircum{}\{ij\}δ(S)\allowbreak{}. (3.1)\allowbreak{}

The sign comes from b\textasciicircum{}\{-i-j\}\allowbreak{}=(-1)\allowbreak{}\textasciicircum{}\{ij\}id\_\allowbreak{}A. In particular,\allowbreak{} the natural projection convention in Section 1 does not send unsigned splicing to composition. DERIVED-NEW-029 records this sign and its propagation into the higher-vanishing proof. The correction makes f\_\allowbreak{}E\textasciicircum{}\{-i\}\allowbreak{}=id\_\allowbreak{}A explicit; it does not change the differentials or rescale the existing class. Because the source leaves the component of f implicit,\allowbreak{} the finding is also a convention-clarity issue:\allowbreak{} a different,\allowbreak{} degree-dependent choice of f would be a different convention and cannot silently repair this calculation.

Here is a counterexample to dropping the sign. Put R\allowbreak{}=Q[ε]\allowbreak{}/\allowbreak{}(ε²)\allowbreak{} and k\allowbreak{}=R/\allowbreak{}(ε)\allowbreak{}. The sequence E:\allowbreak{}0\allowbreak{}→k\allowbreak{}→R\allowbreak{}→k\allowbreak{}→0 has injection 1\allowbreak{}↦ε and quotient q. The complex P with P\textasciicircum{}r\allowbreak{}=R for r≤0,\allowbreak{} P\textasciicircum{}r\allowbreak{}=0 for r>0,\allowbreak{} differential ε in every negative degree,\allowbreak{} and augmentation q is a projective resolution of k:\allowbreak{} ker(ε)\allowbreak{}\allowbreak{}=εR\allowbreak{}=im(ε)\allowbreak{},\allowbreak{} and coker(ε)\allowbreak{}\allowbreak{}=k. The chain map P\allowbreak{}→C\_\allowbreak{}E has components id\_\allowbreak{}R in degree 0,\allowbreak{} q in degree -1,\allowbreak{} and zero below. It represents δ(E)\allowbreak{} by q:\allowbreak{}P\textasciicircum{}\{-1\}\allowbreak{}→k.

A lift g:\allowbreak{}P\allowbreak{}→P[1]\allowbreak{} of that class has g\textasciicircum{}r\allowbreak{}=(-1)\allowbreak{}\textasciicircum{}\{r\allowbreak{}+1\}id\_\allowbreak{}R for r≤-1 and zero in degree 0. The chain equation -ε g\textasciicircum{}r\allowbreak{}=g\textasciicircum{}\{r\allowbreak{}+1\}ε holds. Then δ(E)\allowbreak{}[1]\allowbreak{}δ(E)\allowbreak{} is represented in degree -2 by q g\textasciicircum{}\{-2\}\allowbreak{}=-q. The splice S has middle arrow ε:\allowbreak{}R\allowbreak{}→R; a map P\allowbreak{}→C\_\allowbreak{}S has components id\_\allowbreak{}R in degrees 0 and -1 and q in degree -2,\allowbreak{} so δ(S)\allowbreak{} is represented by q. Every differential in Hom\_\allowbreak{}R(P,\allowbreak{}k)\allowbreak{} is zero because ε acts by zero on k. Thus Ext\textasciicircum{}2\_\allowbreak{}R(k,\allowbreak{}k)\allowbreak{}\allowbreak{}=k,\allowbreak{} and q and -q are distinct. This verifies an actual failure of unsigned composition,\allowbreak{} not just a formal discrepancy of chain maps.

The signs are associative:\allowbreak{} either bracketing of three splices contributes ij\allowbreak{}+ik\allowbreak{}+jk modulo 2. In particular no associativity axiom of the graded derived composition has been changed.

\subsection{\texorpdfstring{4. Zero products,\allowbreak{} the extension diagram,\allowbreak{} and the set of lifts}{4. Zero  the extension  and the set of lifts}}\label{proofs-8774-9196-zero-products-the-extension-diagram-and-the-set-of-lifts}

Write E:\allowbreak{}0\allowbreak{}→A\allowbreak{}→Z\allowbreak{}→B\allowbreak{}→0 and E':\allowbreak{}0\allowbreak{}→B\allowbreak{}→Z'\allowbreak{}→C\allowbreak{}→0,\allowbreak{} with inclusion u:\allowbreak{}B\allowbreak{}→Z' and quotient v:\allowbreak{}Z'\allowbreak{}→C. Applying Hom\_\allowbreak{}D(-,\allowbreak{}A[1]\allowbreak{})\allowbreak{} to the source\textquotesingle s triangle for E\textquotesingle{} gives the exact segment

Hom(B,\allowbreak{}A)\allowbreak{} --∂-\allowbreak{}-> Ext\textasciicircum{}1(C,\allowbreak{}A)\allowbreak{} --v*-\allowbreak{}-> Ext\textasciicircum{}1(Z',\allowbreak{}A)\allowbreak{} --u*-\allowbreak{}-> Ext\textasciicircum{}1(B,\allowbreak{}A)\allowbreak{} --boundary-\allowbreak{}-> Ext\textasciicircum{}2(C,\allowbreak{}A)\allowbreak{}.

For e\allowbreak{}=δ(E)\allowbreak{} the last boundary is e[1]\allowbreak{}∂\_\allowbreak{}\{E'\}\allowbreak{}=-e[1]\allowbreak{}δ(E')\allowbreak{}. Its vanishing is therefore equivalent to vanishing of the product in the original lemma,\allowbreak{} even though the sign must be kept in the exact formula. By exactness this holds precisely when there is w\allowbreak{}∈Ext\textasciicircum{}1(Z',\allowbreak{}A)\allowbreak{} with u*w\allowbreak{}=e. Section 2 represents w by 0\allowbreak{}→A\allowbreak{}→W\allowbreak{}→Z'\allowbreak{}→0. Its pullback to B is isomorphic,\allowbreak{} fixing A and B,\allowbreak{} to E. Choose that isomorphism and compose it with W×\_\allowbreak{}\{Z'\}B\allowbreak{}→W to obtain Z\allowbreak{}→W.

The induced square is a pullback,\allowbreak{} so Z\allowbreak{}→W is a monomorphism,\allowbreak{} and its quotient is C:\allowbreak{} compose W\allowbreak{}→Z'\allowbreak{}→C,\allowbreak{} whose kernel is the inverse image of B. The result is precisely the diagram in original 8941–\allowbreak{}8966,\allowbreak{} with all rows and columns exact. Conversely such a diagram identifies Z with W×\_\allowbreak{}\{Z'\}B by the short five lemma,\allowbreak{} so its middle row gives a lift of e and the product vanishes. This proves both directions of the omitted proof.

DERIVED-UNDER-017 records the extra information supplied by this argument. The set

F\_\allowbreak{}e\allowbreak{}=\{w\allowbreak{}∈Ext\textasciicircum{}1(Z',\allowbreak{}A)\allowbreak{}:\allowbreak{}u*w\allowbreak{}=e\}

is empty if the product is nonzero. If it is nonempty,\allowbreak{} addition makes it a torsor under

im(v*)\allowbreak{} ≅ Ext\textasciicircum{}1(C,\allowbreak{}A)\allowbreak{}/\allowbreak{}im(∂:\allowbreak{}Hom(B,\allowbreak{}A)\allowbreak{}\allowbreak{}→Ext\textasciicircum{}1(C,\allowbreak{}A)\allowbreak{})\allowbreak{}.

Indeed the difference of two lifts lies in ker(u*)\allowbreak{}\allowbreak{}=im(v*)\allowbreak{},\allowbreak{} every such element can be added to a lift,\allowbreak{} and the resulting action of im(v*)\allowbreak{} is free. Exactness identifies its kernel description as the stated quotient. These are isomorphism classes of middle-row extensions admitting a compatible top row; a choice of the isomorphism to that top row is additional data. The assertion is deliberately not a classification of diagrams with that choice fixed.

The earlier boundary is ∂(a)\allowbreak{}\allowbreak{}=-a[1]\allowbreak{}δ(E')\allowbreak{}; in particular its image equals the image of the pushout classes a[1]\allowbreak{}δ(E')\allowbreak{}. This keeps the sign while identifying the same quotient group. If that quotient is zero,\allowbreak{} an existing middle-row class is unique. This consequence is editorial,\allowbreak{} without a replacement of the source lemma.

\subsection{5. The receiving ideal-filtration diagram}\label{proofs-8774-9196-the-receiving-ideal-filtration-diagram}

Current cohomology.tex 14916–\allowbreak{}14934 cites the zero-product criterion. For every index i for which these ideal powers are defined,\allowbreak{} retain exactly

A\allowbreak{}=I\textasciicircum{}\{i\allowbreak{}+2\}/\allowbreak{}I\textasciicircum{}\{i\allowbreak{}+3\},\allowbreak{} B\allowbreak{}=I\textasciicircum{}\{i\allowbreak{}+1\}/\allowbreak{}I\textasciicircum{}\{i\allowbreak{}+2\},\allowbreak{} C\allowbreak{}=I\textasciicircum{}i/\allowbreak{}I\textasciicircum{}\{i\allowbreak{}+1\},\allowbreak{} Z\allowbreak{}=I\textasciicircum{}\{i\allowbreak{}+1\}/\allowbreak{}I\textasciicircum{}\{i\allowbreak{}+3\},\allowbreak{} Z'\allowbreak{}=I\textasciicircum{}i/\allowbreak{}I\textasciicircum{}\{i\allowbreak{}+2\},\allowbreak{} W\allowbreak{}=I\textasciicircum{}i/\allowbreak{}I\textasciicircum{}\{i\allowbreak{}+3\}.

All maps are induced by inclusions or quotient maps of the original ideal powers. The needed rows and columns are

0\allowbreak{}→A\allowbreak{}→Z\allowbreak{}→B\allowbreak{}→0,\allowbreak{} 0\allowbreak{}→A\allowbreak{}→W\allowbreak{}→Z'\allowbreak{}→0,\allowbreak{} 0\allowbreak{}→Z\allowbreak{}→W\allowbreak{}→C\allowbreak{}→0,\allowbreak{} 0\allowbreak{}→B\allowbreak{}→Z'\allowbreak{}→C\allowbreak{}→0.

They are exact because for subobjects J⊆K⊆L the sequence 0\allowbreak{}→K/\allowbreak{}J\allowbreak{}→L/\allowbreak{}J\allowbreak{}→L/\allowbreak{}K\allowbreak{}→0 is exact. Applied to each consecutive triple of the actual ideal powers,\allowbreak{} this also proves commutation of every square. Thus Section 4 proves the product is zero. Both source connecting maps b\^{}i and b\textasciicircum{}\{i\allowbreak{}+1\} are negatives of the corresponding degree-one Yoneda classes,\allowbreak{} so the two minus signs cancel in b\textasciicircum{}\{i\allowbreak{}+1\}[1]\allowbreak{}b\textasciicircum{}i. Applying M⊗\^{}L- to this zero composite still gives zero; applying cohomology therefore gives the zero composite of the Bockstein maps.

Original cohomology.tex 14425–\allowbreak{}14426 describes W as an extension in the opposite order from its quotient sequence. DERIVED-NEW-031 corrects that receiving sentence by displaying the actual middle-row sequence 0\allowbreak{}→I\textasciicircum{}\{i\allowbreak{}+2\}/\allowbreak{}I\textasciicircum{}\{i\allowbreak{}+3\}\allowbreak{}→I\textasciicircum{}i/\allowbreak{}I\textasciicircum{}\{i\allowbreak{}+3\}\allowbreak{}→I\textasciicircum{}i/\allowbreak{}I\textasciicircum{}\{i\allowbreak{}+2\}\allowbreak{}→0,\allowbreak{} whose pullback along B\allowbreak{}→Z' is the top row. This both corrects the reversed description and supplies exactly the hypothesis of the cited criterion. No ideal factor,\allowbreak{} exponent,\allowbreak{} or connecting map has been discarded.

The zero-class use in current adequate.tex 1216–\allowbreak{}1232 also survives (3.1)\allowbreak{}:\allowbreak{} multiplication of a class by a sign does not change whether it is zero. That bounded receiver check concerns the invocation of Yoneda classification,\allowbreak{} not acceptance of the remainder of that chapter\textquotesingle s proof.

\subsection{\texorpdfstring{6. Higher vanishing,\allowbreak{} including p\allowbreak{}=0 and p\allowbreak{}=1}{6. Higher  including  and }}\label{proofs-8774-9196-higher-vanishing-including-p0-and-p1}

For the original higher-vanishing lemma,\allowbreak{} if p\allowbreak{}=0 then Hom(A,\allowbreak{}A)\allowbreak{}\allowbreak{}=0 for every A,\allowbreak{} whence id\_\allowbreak{}A\allowbreak{}=0. Every A is a zero object and so is every complex up to its unique maps. All claimed Ext groups vanish. This case cannot use a positive-degree Yoneda extension of degree p.

Now let p≥1 and i\textgreater p. Represent ξ by the exact extension E as in Section 1 and set C\allowbreak{}=im(Z\_\allowbreak{}p\allowbreak{}→Z\_\allowbreak{}\{p-1\})\allowbreak{}. This is defined also when p\allowbreak{}=1. There are the two exact extensions printed by the source,\allowbreak{} E\textquotesingle{} of B by C of degree p and E\textquotesingle\textquotesingle{} of C by A of degree i-p. Equation (3.1)\allowbreak{} gives the complete equality

δ(E)\allowbreak{}\allowbreak{}=(-1)\allowbreak{}\textasciicircum{}\{p(i-p)\allowbreak{}\}δ(E'')\allowbreak{}[p]\allowbreak{}δ(E')\allowbreak{}.

The assumption makes δ(E')\allowbreak{} zero,\allowbreak{} so ξ\allowbreak{}=0. The case i\allowbreak{}=p is the assumption itself. DERIVED-NEW-030 inserts the p\allowbreak{}=0 argument and uses the image formula for C:\allowbreak{} the source\textquotesingle s additional kernel formula contains Z\_\allowbreak{}\{-1\} when p\allowbreak{}=1 and neither formula is defined as written when p\allowbreak{}=0. The signed equality is the propagation of NEW-029,\allowbreak{} not another counted finding. The theorem itself remains valid with its original hypotheses.

\subsection{7. A smaller splitting hypothesis and all choices of splitting}\label{proofs-8774-9196-a-smaller-splitting-hypothesis-and-all-choices-of-splitting}

Let K\allowbreak{}∈D\textasciicircum{}b(A)\allowbreak{},\allowbreak{} and write H\_\allowbreak{}r\allowbreak{}=H\textasciicircum{}r(K)\allowbreak{}. Assume only

Ext\textasciicircum{}\{r-s\allowbreak{}+1\}(H\_\allowbreak{}r,\allowbreak{}H\_\allowbreak{}s)\allowbreak{}\allowbreak{}=0 for every r\textgreater s. (7.1)\allowbreak{}

This is weaker than requiring every degree at least two for every such pair,\allowbreak{} as the original lemma does. Choose a≤b containing the cohomology support. The canonical truncation triangle at r is

τ\_\allowbreak{}\{\textbackslash{}leq r-1\}K\allowbreak{}→τ\_\allowbreak{}\{\textbackslash{}leq r\}K\allowbreak{}→H\_\allowbreak{}r[-r]\allowbreak{}\allowbreak{}→(τ\_\allowbreak{}\{\textbackslash{}leq r-1\}K)\allowbreak{}[1]\allowbreak{}.

Inductively the lower truncation has an isomorphism τ\_\allowbreak{}\{\textbackslash{}leq r-1\}K≅\allowbreak{}⊕\_\allowbreak{}\{a≤s<r\}H\_\allowbreak{}s[-s]\allowbreak{} inducing identity on its cohomology. Consequently the group containing the last arrow is

\allowbreak{}⊕\_\allowbreak{}\{a≤s<r\} Hom\_\allowbreak{}D(H\_\allowbreak{}r[-r]\allowbreak{},\allowbreak{}H\_\allowbreak{}s[-s\allowbreak{}+1]\allowbreak{})\allowbreak{} \allowbreak{}= \allowbreak{}⊕\_\allowbreak{}\{a≤s<r\} Ext\textasciicircum{}\{r-s\allowbreak{}+1\}(H\_\allowbreak{}r,\allowbreak{}H\_\allowbreak{}s)\allowbreak{}\allowbreak{}=0.

The triangle splits. Choose the splitting that is compatible with its first and second maps; combined with the previous identity on cohomology,\allowbreak{} it gives τ\_\allowbreak{}\{\textbackslash{}leq r\}K≅\allowbreak{}⊕\_\allowbreak{}\{a≤s≤r\}H\_\allowbreak{}s[-s]\allowbreak{},\allowbreak{} still inducing the identity. The one-degree and zero-cohomology cases follow from canonical truncation quasi-isomorphisms. This completes the finite induction and proves (7.1)\allowbreak{} is sufficient. Without the group vanishing,\allowbreak{} the actual last arrow is the exact obstruction:\allowbreak{} its being zero is necessary and sufficient for this stage to split,\allowbreak{} whether or not its ambient group is zero.

Set T\allowbreak{}=\allowbreak{}⊕\_\allowbreak{}\{a≤r≤b\}H\_\allowbreak{}r[-r]\allowbreak{}. A matrix entry from H\_\allowbreak{}r[-r]\allowbreak{} to H\_\allowbreak{}s[-s]\allowbreak{} is Ext\textasciicircum{}\{r-s\}(H\_\allowbreak{}r,\allowbreak{}H\_\allowbreak{}s)\allowbreak{}. It vanishes for r\textless s by negative Ext vanishing. For r\allowbreak{}=s it is Hom(H\_\allowbreak{}r,\allowbreak{}H\_\allowbreak{}r)\allowbreak{},\allowbreak{} and induces precisely that map on H\^{}r. Thus endomorphisms of T inducing zero on every cohomology object form the strictly triangular ideal

N\allowbreak{}=\allowbreak{}⊕\_\allowbreak{}\{a≤s<r≤b\}Ext\textasciicircum{}\{r-s\}(H\_\allowbreak{}r,\allowbreak{}H\_\allowbreak{}s)\allowbreak{} ⊆ End\_\allowbreak{}D(T)\allowbreak{}.

Its multiplication is actual shifted composition. A product of m entries requires a strictly descending chain of m cohomological degrees; hence N\textasciicircum{}\{b-a\allowbreak{}+1\}\allowbreak{}=0. Every 1\allowbreak{}+n is invertible,\allowbreak{} with inverse 1-n\allowbreak{}+n²-…\allowbreak{}+(-1)\allowbreak{}\textasciicircum{}\{b-a\}n\textasciicircum{}\{b-a\}. Conversely an automorphism inducing identity on cohomology has the form 1\allowbreak{}+n. Once an identity-on-cohomology splitting T\allowbreak{}→K exists,\allowbreak{} all such splittings form a torsor under 1\allowbreak{}+N by precomposition. The action is free and transitive by composing one splitting with the inverse of another.

Therefore the extra vanishing Ext\textasciicircum{}\{r-s\}(H\_\allowbreak{}r,\allowbreak{}H\_\allowbreak{}s)\allowbreak{}\allowbreak{}=0 for r\textgreater s makes the identity-on-cohomology splitting unique. Without it there is no general claim of a canonical splitting. DERIVED-UNDER-018 comprises (7.1)\allowbreak{},\allowbreak{} the precise obstruction,\allowbreak{} and this description of the choices. All are proved consequences of the source argument,\allowbreak{} kept in editorial material.

The source\textquotesingle s global Ext²\allowbreak{}=0 conclusion follows from Section 6 and then (7.1)\allowbreak{}. The two checked More on Algebra receivers,\allowbreak{} current 21522–\allowbreak{}21531 and 28714–\allowbreak{}28725,\allowbreak{} assume projective or free cohomology. Their positive Ext groups vanish,\allowbreak{} so (7.1)\allowbreak{} and the uniqueness condition both hold. Their stated existence conclusions remain valid; no receiving body replacement is needed for the stronger editorial result.

\subsection{8. K-groups and all omitted verifications in this interval}\label{proofs-8774-9196-k-groups-and-all-omitted-verifications-in-this-interval}

Use the source\textquotesingle s convention that object collections are sets (or work in the chosen universe)\allowbreak{}. In a triangulated category,\allowbreak{} the triangle of zeros implies [0]\allowbreak{}\allowbreak{}=0. An isomorphism X\allowbreak{}→Y has a triangle with third object zero,\allowbreak{} so [X]\allowbreak{}\allowbreak{}=[Y]\allowbreak{}. The triangle X\allowbreak{}→0\allowbreak{}→X[1]\allowbreak{} gives [X[1]\allowbreak{}]\allowbreak{}\allowbreak{}=-[X]\allowbreak{}; induction in both directions gives [X[n]\allowbreak{}]\allowbreak{}\allowbreak{}=(-1)\allowbreak{}\textasciicircum{}n[X]\allowbreak{} for every integer n. Thus the definition using all objects has the required invariance under isomorphism. The report P02-E0068/\allowbreak{}OCC-01538 corrects the preposition in original 9067. There is no prior admitted correction of this locus.

For a finite exact sequence 0\allowbreak{}→M\_\allowbreak{}a\allowbreak{}→M\_\allowbreak{}\{a\allowbreak{}+1\}\allowbreak{}→…\allowbreak{}→M\_\allowbreak{}b\allowbreak{}→0 in an abelian category,\allowbreak{} let B\_\allowbreak{}r be the image entering M\_\allowbreak{}r. The short exact sequences 0\allowbreak{}→B\_\allowbreak{}r\allowbreak{}→M\_\allowbreak{}r\allowbreak{}→B\_\allowbreak{}\{r\allowbreak{}+1\}\allowbreak{}→0 give [M\_\allowbreak{}r]\allowbreak{}\allowbreak{}=[B\_\allowbreak{}r]\allowbreak{}\allowbreak{}+[B\_\allowbreak{}\{r\allowbreak{}+1\}]\allowbreak{}. The alternating sum telescopes,\allowbreak{} since B\_\allowbreak{}a\allowbreak{}=B\_\allowbreak{}\{b\allowbreak{}+1\}\allowbreak{}=0. Consequently any bounded long exact cohomology sequence of a triangle X\allowbreak{}→Y\allowbreak{}→Z gives χ(Y)\allowbreak{}\allowbreak{}=χ(X)\allowbreak{}\allowbreak{}+χ(Z)\allowbreak{},\allowbreak{} with χ(X)\allowbreak{}\allowbreak{}=Σ\_\allowbreak{}r(-1)\allowbreak{}\textasciicircum{}r[H\textasciicircum{}r(X)\allowbreak{}]\allowbreak{}. The signs follow by arranging its terms H\textasciicircum{}r(X)\allowbreak{},\allowbreak{}H\textasciicircum{}r(Y)\allowbreak{},\allowbreak{}H\textasciicircum{}r(Z)\allowbreak{},\allowbreak{}H\textasciicircum{}\{r\allowbreak{}+1\}(X)\allowbreak{},\allowbreak{}… in their actual sequence order:\allowbreak{} up to an overall sign the three coefficients in degree r are (-1)\allowbreak{}\textasciicircum{}r,\allowbreak{}-(-1)\allowbreak{}\textasciicircum{}r,\allowbreak{}(-1)\allowbreak{}\textasciicircum{}r.

The assignment A\allowbreak{}↦A[0]\allowbreak{} sends each short exact sequence to the source\textquotesingle s distinguished triangle and so induces K\_\allowbreak{}0(A)\allowbreak{}\allowbreak{}→K\_\allowbreak{}0(D\textasciicircum{}b(A)\allowbreak{})\allowbreak{}. The Euler assignment above descends in the reverse direction and their composition on [A]\allowbreak{} is the identity. For any K\allowbreak{}∈D\textasciicircum{}b(A)\allowbreak{},\allowbreak{} the finite canonical truncation triangles give

[K]\allowbreak{}\allowbreak{}=Σ\_\allowbreak{}r[H\textasciicircum{}r(K)\allowbreak{}[-r]\allowbreak{}]\allowbreak{}\allowbreak{}=Σ\_\allowbreak{}r(-1)\allowbreak{}\textasciicircum{}r[H\textasciicircum{}r(K)\allowbreak{}[0]\allowbreak{}]\allowbreak{}.

The finite cohomological support makes this induction terminate at zero. This proves the other composition is the identity,\allowbreak{} without assuming the initial representative of K is termwise bounded. If needed,\allowbreak{} a bounded representative is obtained by the actual canonical truncations at the two cohomological endpoints. The source\textquotesingle s alternative finite induction on stupid truncations of that representative is valid.

If B⊆A is weak Serre,\allowbreak{} it is a full abelian subcategory whose inclusion is exact. In the long cohomology sequence of a triangle in D\textasciicircum{}b\_\allowbreak{}B(A)\allowbreak{},\allowbreak{} all terms lie in B. Kernels and cokernels of its maps belong to B,\allowbreak{} so all intervening images belong to B as well. The preceding telescoping argument therefore takes place in K\_\allowbreak{}0(B)\allowbreak{},\allowbreak{} not just K\_\allowbreak{}0(A)\allowbreak{}. The canonical truncations stay in D\textasciicircum{}b\_\allowbreak{}B(A)\allowbreak{},\allowbreak{} since their cohomology objects are retained or replaced by zero. The same formula gives mutually inverse maps K\_\allowbreak{}0(B)\allowbreak{}⇄K\_\allowbreak{}0(D\textasciicircum{}b\_\allowbreak{}B(A)\allowbreak{})\allowbreak{}. This fills the proof at original 9164–\allowbreak{}9176. It does not assume that every such complex has a representative with all its terms in B.

An exact functor of triangulated categories sends each triangle relation to a triangle relation,\allowbreak{} so its assignment on objects descends to K\_\allowbreak{}0. For the source\textquotesingle s homological functor H with finite support on every shift orbit,\allowbreak{} apply its long exact sequence to a triangle. The union of the three finite supports is finite,\allowbreak{} and the same Euler calculation proves that [X]\allowbreak{}\allowbreak{}↦Σ\_\allowbreak{}r(-1)\allowbreak{}\textasciicircum{}r[H(X[r]\allowbreak{})\allowbreak{}]\allowbreak{} respects every relation. Finally,\allowbreak{} a bifunctor exact in each variable respects the triangle relations separately in each input. Its map on pairs of generators therefore descends to the claimed bilinear map on the two K-groups. These arguments complete the omitted verifications through original 9196.

\subsection{\texorpdfstring{9. Dispositions,\allowbreak{} propagation,\allowbreak{} and limits}{9.   and limits}}\label{proofs-8774-9196-dispositions-propagation-and-limits}

One additional physical report is complete:\allowbreak{} P02-E0068/\allowbreak{}OCC-01538,\allowbreak{} DERIVED-RECON-069,\allowbreak{} missing copyedit. Cumulatively that is 117 of 176 reports,\allowbreak{} 69 groups,\allowbreak{} 108 already corrected,\allowbreak{} eight missing copyedit/\allowbreak{}notation reports and one missing mathematical-index report.

New unreported findings are NEW-028 (Ext endpoints)\allowbreak{},\allowbreak{} NEW-029 (signed splicing with its actual projection convention)\allowbreak{},\allowbreak{} NEW-030 (p\allowbreak{}=0/\allowbreak{}p\allowbreak{}=1 proof endpoints)\allowbreak{},\allowbreak{} and NEW-031 (receiving Cohomology extension sequence)\allowbreak{}. There are eight new proposed operations:\allowbreak{} seven in Derived Categories and one in Cohomology. The first seven include the received copyedit. The two new editorial consequences are UNDER-017 and UNDER-018.

Propagation has been proved into the higher-vanishing lemma,\allowbreak{} the zero-product criterion,\allowbreak{} the ideal-filtration Bockstein argument,\allowbreak{} the bounded splitting lemmas,\allowbreak{} and the bounded receiving uses specified above. The common-refinement and Baer-sum arguments support correct source claims and are not counted as errors merely because their source proofs are short. The K-group section requires only its received preposition correction.

Semantic review is complete through 9196; sequential reading reaches 9232. The next section is Unbounded complexes,\allowbreak{} beginning at 9206; the first incomplete source lemma begins at 9214. New unread source starts at 9233. No claim is made about that lemma\textquotesingle s proof yet. The Illusie I.3.0–\allowbreak{}I.3.1.1 producer relay received during this review is queued separately,\allowbreak{} without mathematical acceptance or interruption of the correction priority. No new sessions,\allowbreak{} live source composition,\allowbreak{} builds,\allowbreak{} Lean runs,\allowbreak{} publication,\allowbreak{} or cleanup were performed.
\clearpage
\section{\texorpdfstring{Unbounded resolutions,\allowbreak{} derived adjunction,\allowbreak{} and products of K-injectives}{Unbounded  derived  and products of K-injectives}}\label{proofs-9206-9744-unbounded-resolutions-derived-adjunction-and-products-of-k-injectives}

Private source-order review of original derived.tex 9206–\allowbreak{}9744. Authority:\allowbreak{} Stacks commit a04446e5\allowbreak{}7ec1fbc2\allowbreak{}52a871af\allowbreak{}cec7752f\allowbreak{}b2807b14. Receiving public edition:\allowbreak{} 95a51593\allowbreak{}aceb247a\allowbreak{}c7a18311\allowbreak{}1678c3a0\allowbreak{}f8b8f4b5. The recorder binds the source files,\allowbreak{} original reports,\allowbreak{} operations,\allowbreak{} and exact proof identity. Editorial consequences remain separate from the source body. No novelty is claimed.

Human sources are \href{https://github.com/stacks/stacks-project/blob/a04446e57ec1fbc252a871afcec7752fb2807b14/derived.tex\#L9206}{the original Derived Categories TeX} and the referenced Categories and Homology chapters at the same authority. Secondary online comparisons located \href{https://stacks.math.columbia.edu/tag/0794}{the unbounded existence proposition} and \href{https://stacks.math.columbia.edu/tag/0BK6}{the K-injective product lemma}; the frozen TeX,\allowbreak{} rather than the mutable online edition,\allowbreak{} determines the corrections.

\subsection{1. A strict resolution system with its actual upper bounds}\label{proofs-9206-9744-a-strict-resolution-system-with-its-actual-upper-bounds}

Let A be abelian and let P contain zero,\allowbreak{} be closed under finite direct sums,\allowbreak{} and provide an epimorphism from an object of P onto each object of A. For an arbitrary complex K write T\_\allowbreak{}n\allowbreak{}=τ\_\allowbreak{}\{\textbackslash{}leq n\}K,\allowbreak{} n≥1. We construct termwise epimorphic quasi-isomorphisms p\_\allowbreak{}n:\allowbreak{}P\_\allowbreak{}n\allowbreak{}→T\_\allowbreak{}n with

P\_\allowbreak{}n\textasciicircum{}r\allowbreak{}=0 for r>n,\allowbreak{}

strictly compatible with termwise split monomorphisms P\_\allowbreak{}n\allowbreak{}→P\_\allowbreak{}\{n\allowbreak{}+1\},\allowbreak{} whose termwise cokernels belong to P.

The bounded resolution lemma,\allowbreak{} current derived.tex 5654–\allowbreak{}5737,\allowbreak{} gives P\_\allowbreak{}1\allowbreak{}→T\_\allowbreak{}1 with upper bound 1 and with all the required termwise properties. Its descending construction is worth specifying:\allowbreak{} to extend a partial resolution in degrees ≥r,\allowbreak{} choose an epimorphism from P\^{}\{r-1\} onto K\textasciicircum{}\{r-1\}×\_\allowbreak{}\{ker(d\_\allowbreak{}K\textasciicircum{}r)\allowbreak{}\}ker(d\_\allowbreak{}P\textasciicircum{}r)\allowbreak{}. Projection to the two factors gives the augmentation and differential. Surjectivity onto cycles at degree r and exactness on the previously established cohomology follow from this pullback. Starting with zero terms above the bound produces the stated resolution. This is the bounded construction used below.

Suppose P\_\allowbreak{}n and p\_\allowbreak{}n are given. Let Y\allowbreak{}=T\_\allowbreak{}\{n\allowbreak{}+1\} and let a:\allowbreak{}P\_\allowbreak{}n\allowbreak{}→Y be the composite with the canonical truncation inclusion. Use the actual source cone

C(a)\allowbreak{}\textasciicircum{}r\allowbreak{}=Y\textasciicircum{}r\allowbreak{}⊕P\_\allowbreak{}n\textasciicircum{}\{r\allowbreak{}+1\},\allowbreak{} d\_\allowbreak{}C\textasciicircum{}r\allowbreak{}=[[d\_\allowbreak{}Y\textasciicircum{}r,\allowbreak{}a\textasciicircum{}\{r\allowbreak{}+1\}]\allowbreak{},\allowbreak{}[0,\allowbreak{}-d\_\allowbreak{}\{P\_\allowbreak{}n\}\textasciicircum{}\{r\allowbreak{}+1\}]\allowbreak{}]\allowbreak{}.

It vanishes above n\allowbreak{}+1. Choose a termwise epimorphic quasi-isomorphism q:\allowbreak{}Q\allowbreak{}→C(a)\allowbreak{} from a complex with terms in P and Q\textasciicircum{}r\allowbreak{}=0 for r>n\allowbreak{}+1. Write q\textasciicircum{}r\allowbreak{}=(u\textasciicircum{}r,\allowbreak{}v\textasciicircum{}r)\allowbreak{},\allowbreak{} so

d\_\allowbreak{}Y u\allowbreak{}+a v\allowbreak{}=u d\_\allowbreak{}Q,\allowbreak{} -d\_\allowbreak{}\{P\_\allowbreak{}n\}v\allowbreak{}=v d\_\allowbreak{}Q. (1.1)\allowbreak{}

Define the next complex and map by

P\_\allowbreak{}\{n\allowbreak{}+1\}\textasciicircum{}r\allowbreak{}=P\_\allowbreak{}n\textasciicircum{}r\allowbreak{}⊕Q\textasciicircum{}r,\allowbreak{} d\_\allowbreak{}\{n\allowbreak{}+1\}\textasciicircum{}r\allowbreak{}=[[d\_\allowbreak{}\{P\_\allowbreak{}n\}\textasciicircum{}r,\allowbreak{}-v\textasciicircum{}r]\allowbreak{},\allowbreak{}[0,\allowbreak{}d\_\allowbreak{}Q\textasciicircum{}r]\allowbreak{}]\allowbreak{},\allowbreak{} p\_\allowbreak{}\{n\allowbreak{}+1\}\textasciicircum{}r\allowbreak{}=(a\textasciicircum{}r,\allowbreak{}u\textasciicircum{}r)\allowbreak{}:\allowbreak{}P\_\allowbreak{}n\textasciicircum{}r\allowbreak{}⊕Q\textasciicircum{}r\allowbreak{}→Y\textasciicircum{}r. (1.2)\allowbreak{}

Equation (1.1)\allowbreak{} gives d\_\allowbreak{}\{n\allowbreak{}+1\}²\allowbreak{}=0 and p\_\allowbreak{}\{n\allowbreak{}+1\}d\_\allowbreak{}\{n\allowbreak{}+1\}\allowbreak{}=d\_\allowbreak{}Y p\_\allowbreak{}\{n\allowbreak{}+1\}. The inclusion of P\_\allowbreak{}n is a chain map,\allowbreak{} and its composite with p\_\allowbreak{}\{n\allowbreak{}+1\} equals a strictly. The quotient complex is Q. Thus the inclusion is termwise split,\allowbreak{} its cokernel terms are in P,\allowbreak{} and the terms of P\_\allowbreak{}\{n\allowbreak{}+1\} belong to P. Its upper bound is n\allowbreak{}+1.

The map u\^{}r is the composite of the epimorphism q\^{}r with the biproduct projection C(a)\allowbreak{}\textasciicircum{}r\allowbreak{}→Y\textasciicircum{}r; hence u\textasciicircum{}r,\allowbreak{} and therefore p\_\allowbreak{}\{n\allowbreak{}+1\}\textasciicircum{}r,\allowbreak{} are epimorphisms. Finally,\allowbreak{} both C(p\_\allowbreak{}\{n\allowbreak{}+1\})\allowbreak{}\textasciicircum{}r and C(q)\allowbreak{}\textasciicircum{}r are,\allowbreak{} in the same ordered coordinates,\allowbreak{}

Y\textasciicircum{}r\allowbreak{}⊕P\_\allowbreak{}n\textasciicircum{}\{r\allowbreak{}+1\}\allowbreak{}⊕Q\textasciicircum{}\{r\allowbreak{}+1\},\allowbreak{}

and both differentials have the full matrix

[[d\_\allowbreak{}Y\textasciicircum{}r,\allowbreak{} a\textasciicircum{}\{r\allowbreak{}+1\},\allowbreak{} u\textasciicircum{}\{r\allowbreak{}+1\}]\allowbreak{},\allowbreak{} [0,\allowbreak{} -d\_\allowbreak{}\{P\_\allowbreak{}n\}\textasciicircum{}\{r\allowbreak{}+1\},\allowbreak{} v\textasciicircum{}\{r\allowbreak{}+1\}]\allowbreak{},\allowbreak{} [0,\allowbreak{} 0,\allowbreak{} -d\_\allowbreak{}Q\textasciicircum{}\{r\allowbreak{}+1\}]\allowbreak{}]\allowbreak{}.

The associativity map of these biproducts is an isomorphism of complexes. Since q is a quasi-isomorphism,\allowbreak{} its cone is acyclic; the same is true of C(p\_\allowbreak{}\{n\allowbreak{}+1\})\allowbreak{}. Thus p\_\allowbreak{}\{n\allowbreak{}+1\} is a quasi-isomorphism. This proves the induction. No compatibility merely in the homotopy category has been substituted for compatibility of the actual resolution maps.

DERIVED-UNDER-019 records the term bounds P\_\allowbreak{}n\textasciicircum{}r\allowbreak{}=0 for r\textgreater n and this explicit strict construction,\allowbreak{} as well as its dual in Section 4. The source states only boundedness above,\allowbreak{} although its hypotheses permit these sharper bounds. The proof is editorial; the existing source proof is also valid as checked next.

\subsection{\texorpdfstring{2. The original homotopy adjustment,\allowbreak{} disk correction,\allowbreak{} and colimit map}{2. The original homotopy  disk  and colimit map}}\label{proofs-9206-9744-the-original-homotopy-adjustment-disk-correction-and-colimit-map}

The source first constructs a middle complex P\textquotesingle{} in a triangle and a map f:\allowbreak{}P'\allowbreak{}→Y with f i-a\allowbreak{}=d h\allowbreak{}+h d on P\_\allowbreak{}n. Here i:\allowbreak{}P\_\allowbreak{}n\allowbreak{}→P' is termwise split. Choose graded retractions r\textasciicircum{}m:\allowbreak{}P'\textasciicircum{}m\allowbreak{}→P\_\allowbreak{}n\textasciicircum{}m and extend the homotopy by h̃\textasciicircum{}m\allowbreak{}=h\textasciicircum{}m r\^{}m. Then

f̃\allowbreak{}=f-d\_\allowbreak{}Y h̃-h̃ d\_\allowbreak{}\{P'\}

is a chain map homotopic to f,\allowbreak{} and f̃ i\allowbreak{}=a strictly. The equality holds because i is a chain map and r i\allowbreak{}=id; r need not be a chain map. This verifies the source\textquotesingle s omitted adjustment,\allowbreak{} including its sign.

After that adjustment,\allowbreak{} the source adds the contractible disk D with copies of an object P in degrees n,\allowbreak{}n\allowbreak{}+1,\allowbreak{} differential id\_\allowbreak{}P,\allowbreak{} and map to T\_\allowbreak{}\{n\allowbreak{}+1\} given by q:\allowbreak{}P\allowbreak{}→K\textasciicircum{}n and d\_\allowbreak{}K\textasciicircum{}n q:\allowbreak{}P\allowbreak{}→ker(d\_\allowbreak{}K\textasciicircum{}\{n\allowbreak{}+1\})\allowbreak{}. Here q is epic. The map is a chain map,\allowbreak{} the disk is contractible,\allowbreak{} and adding it preserves the quasi-isomorphism and the split inclusion of P\_\allowbreak{}n. Surjectivity in degree n is immediate. In degree n\allowbreak{}+1,\allowbreak{} the disk image is im(d\_\allowbreak{}K\textasciicircum{}n)\allowbreak{}. The map H\textasciicircum{}\{n\allowbreak{}+1\}(f̃)\allowbreak{} is epic; hence the image of f̃ in ker(d\_\allowbreak{}K\textasciicircum{}\{n\allowbreak{}+1\})\allowbreak{},\allowbreak{} together with im(d\_\allowbreak{}K\textasciicircum{}n)\allowbreak{},\allowbreak{} is all of ker(d\_\allowbreak{}K\textasciicircum{}\{n\allowbreak{}+1\})\allowbreak{}. This proves termwise surjectivity in the remaining degree. At degrees below n it already follows from the previous epimorphic augmentation. Thus the disk argument is correct; Section 1 gives a choice for which this final disk is unnecessary.

Assume A has exact sequential colimits. Put P\_\allowbreak{}∞\allowbreak{}=colim\_\allowbreak{}n P\_\allowbreak{}n termwise. The compatible p\_\allowbreak{}n give a chain map p:\allowbreak{}P\_\allowbreak{}∞\allowbreak{}→K because colim\_\allowbreak{}n T\_\allowbreak{}n\allowbreak{}=K. Indeed in each fixed degree r the term and adjacent differential of T\_\allowbreak{}n are those of K for all sufficiently large n. Exactness of colimits,\allowbreak{} applied to kernels,\allowbreak{} images and their quotient sequences,\allowbreak{} gives

H\textasciicircum{}r(P\_\allowbreak{}∞)\allowbreak{}\allowbreak{}=colim\_\allowbreak{}n H\textasciicircum{}r(P\_\allowbreak{}n)\allowbreak{}\allowbreak{}=H\textasciicircum{}r(K)\allowbreak{}.

For n≥max(1,\allowbreak{}r)\allowbreak{},\allowbreak{} the maps identify H\textasciicircum{}r(P\_\allowbreak{}n)\allowbreak{} with H\textasciicircum{}r(K)\allowbreak{} and are identities under these identifications. Hence p is a quasi-isomorphism. It is termwise epic too:\allowbreak{} choose n>r,\allowbreak{} so P\_\allowbreak{}n\textasciicircum{}r\allowbreak{}→K\textasciicircum{}r is already epic,\allowbreak{} and factor it through P\_\allowbreak{}∞\textasciicircum{}r.

The upper bound from Section 1 means that the map P\_\allowbreak{}n\allowbreak{}→P\_\allowbreak{}∞ factors strictly through τ\_\allowbreak{}\{\textbackslash{}leq n\}P\_\allowbreak{}∞. At degree n its image lies in the kernel of d\_\allowbreak{}\{P\_\allowbreak{}∞\}\textasciicircum{}n because P\_\allowbreak{}n\textasciicircum{}\{n\allowbreak{}+1\}\allowbreak{}=0. This factorization is a quasi-isomorphism,\allowbreak{} by the preceding cohomology calculation. This supplies exactly the stage maps used in the source\textquotesingle s comparison proof; it does not assume that a cohomological upper bound alone gives a factorization of arbitrary chain maps.

For each fixed degree r,\allowbreak{} construction (1.2)\allowbreak{} gives the compatible decompositions

P\_\allowbreak{}n\textasciicircum{}r\allowbreak{}=P\_\allowbreak{}1\textasciicircum{}r\allowbreak{}⊕Q\_\allowbreak{}1\textasciicircum{}r\allowbreak{}⊕…\allowbreak{}⊕Q\_\allowbreak{}\{n-1\}\textasciicircum{}r,\allowbreak{} P\_\allowbreak{}∞\textasciicircum{}r\allowbreak{}=P\_\allowbreak{}1\textasciicircum{}r\allowbreak{}⊕\allowbreak{}⊕\_\allowbreak{}\{n≥1\}Q\_\allowbreak{}n\textasciicircum{}r. (2.1)\allowbreak{}

The last isomorphism follows from the coproduct universal property:\allowbreak{} a compatible cocone on the finite partial sums is exactly one map out of each listed summand. This observation concerns the terms; the complex differential retains all the off-diagonal maps -v\_\allowbreak{}n from (1.2)\allowbreak{}.

\subsection{3. Unbounded left derivation and the weaker additive hypothesis}\label{proofs-9206-9744-unbounded-left-derivation-and-the-weaker-additive-hypothesis}

Here is the complete strengthened assertion,\allowbreak{} DERIVED-UNDER-020. Let A,\allowbreak{}B be abelian categories with exact sequential colimits. Let P satisfy Section 1\textquotesingle s hypotheses,\allowbreak{} and let F:\allowbreak{}A\allowbreak{}→B be additive. Suppose:\allowbreak{}

* F takes every bounded above acyclic complex with terms in P to an acyclic complex; * F commutes with countable coproducts.

Then the termwise functor has a left derived functor on all of D(A)\allowbreak{}. For the system of Section 1 it is computed by F(P\_\allowbreak{}∞)\allowbreak{}. The original proposition follows:\allowbreak{} its right exact functor is additive,\allowbreak{} and preservation of sequential colimits implies preservation of countable coproducts by expressing them as colimits of finite partial sums. Right exactness is not used in the argument below.

First consider bounded above complexes. Each has a resolution by a bounded above complex with terms in P. A quasi-isomorphism between two such resolutions has an acyclic cone whose terms are finite sums of P objects. F preserves the actual cone matrix because it is additive. The first assumption therefore makes its image a quasi-isomorphism. The left version of the source\textquotesingle s existence-and-computation criterion (current 5444–\allowbreak{}5470)\allowbreak{} produces LF\^{}- on D\textasciicircum{}-(A)\allowbreak{},\allowbreak{} computed by these resolutions. Canonical upper truncation supplies bounded representatives of cohomologically bounded-above objects; all comparisons are in D\textasciicircum{}-(A)\allowbreak{}.

Let S be the class of complexes that admit a presentation as a colimit of Section 1\textquotesingle s systems,\allowbreak{} with the stage maps to their own truncations as in Section 2. Every complex receives a quasi-isomorphism from an object of S. For P,\allowbreak{}Q in S choose witnesses P\_\allowbreak{}n,\allowbreak{}Q\_\allowbreak{}n and let α:\allowbreak{}P\allowbreak{}→Q be any quasi-isomorphism. It need not preserve these witnesses.

Write e\_\allowbreak{}\{P,\allowbreak{}n\}:\allowbreak{}F(P\_\allowbreak{}n)\allowbreak{}\allowbreak{}→LF\textasciicircum{}-(τ\_\allowbreak{}\{\textbackslash{}leq n\}P)\allowbreak{} and e\_\allowbreak{}\{Q,\allowbreak{}n\}:\allowbreak{}F(Q\_\allowbreak{}n)\allowbreak{}\allowbreak{}→LF\textasciicircum{}-(τ\_\allowbreak{}\{\textbackslash{}leq n\}Q)\allowbreak{} for the isomorphisms induced by their resolution maps. The maps

θ\_\allowbreak{}n\allowbreak{}=e\_\allowbreak{}\{Q,\allowbreak{}n\}\textasciicircum{}\{-1\} LF\textasciicircum{}-(τ\_\allowbreak{}\{\textbackslash{}leq n\}α)\allowbreak{}e\_\allowbreak{}\{P,\allowbreak{}n\} :\allowbreak{}F(P\_\allowbreak{}n)\allowbreak{}\allowbreak{}→F(Q\_\allowbreak{}n)\allowbreak{} (3.1)\allowbreak{}

are isomorphisms in D(B)\allowbreak{}. They are compatible with the actual maps P\_\allowbreak{}n\allowbreak{}→P\_\allowbreak{}\{n\allowbreak{}+1\},\allowbreak{} Q\_\allowbreak{}n\allowbreak{}→Q\_\allowbreak{}\{n\allowbreak{}+1\}:\allowbreak{} the truncation inclusions commute with τ\_\allowbreak{}\{\textbackslash{}leq n\}α,\allowbreak{} LF\^{}- is a functor,\allowbreak{} and the identifications e respect maps between the P resolutions. This proves stage compatibility,\allowbreak{} not just existence of unrelated isomorphisms.

It remains to identify their colimit with the map induced by F(α)\allowbreak{}. In K\textasciicircum{}-(A)\allowbreak{},\allowbreak{} choose a common roof R\allowbreak{}→P\_\allowbreak{}n and R\allowbreak{}→Q\_\allowbreak{}n over τ\_\allowbreak{}\{\textbackslash{}leq n\}α,\allowbreak{} commuting up to homotopy after the augmentation maps. One obtains such a roof by the calculus of quasi-isomorphisms,\allowbreak{} then by truncating its numerator above n and resolving it by a bounded above complex of P objects. Both R\allowbreak{}→P\_\allowbreak{}n and R\allowbreak{}→Q\_\allowbreak{}n are quasi-isomorphisms. Denote them a and b. The bounded computation says θ\_\allowbreak{}n\allowbreak{}=F(b)\allowbreak{}F(a)\allowbreak{}\textasciicircum{}\{-1\}. After composing to P and Q,\allowbreak{} the roof identity gives the homotopy α i\_\allowbreak{}\{P,\allowbreak{}n\}a∼i\_\allowbreak{}\{Q,\allowbreak{}n\}b. Additivity of F preserves this homotopy,\allowbreak{} so in D(B)\allowbreak{}

F(α)\allowbreak{} F(i\_\allowbreak{}\{P,\allowbreak{}n\})\allowbreak{}\allowbreak{}=F(i\_\allowbreak{}\{Q,\allowbreak{}n\})\allowbreak{}θ\_\allowbreak{}n. (3.2)\allowbreak{}

The definition (3.1)\allowbreak{} proves independence of all roof choices. It also proves that equality of alternative roofs,\allowbreak{} which can be checked after a common refinement,\allowbreak{} is carried to equality,\allowbreak{} not only to an unspecified isomorphism.

By (2.1)\allowbreak{} and preservation of countable coproducts,\allowbreak{} F(P)\allowbreak{}\allowbreak{}=colim\_\allowbreak{}n F(P\_\allowbreak{}n)\allowbreak{} and F(Q)\allowbreak{}\allowbreak{}=colim\_\allowbreak{}n F(Q\_\allowbreak{}n)\allowbreak{},\allowbreak{} canonically degree by degree. The comparison is a chain map because it commutes with the maps from each finite stage and those maps determine a map out of the colimit. Exact sequential colimits in B therefore give

H\textasciicircum{}r(F(P)\allowbreak{})\allowbreak{}\allowbreak{}=colim\_\allowbreak{}n H\textasciicircum{}r(F(P\_\allowbreak{}n)\allowbreak{})\allowbreak{},\allowbreak{} H\textasciicircum{}r(F(Q)\allowbreak{})\allowbreak{}\allowbreak{}=colim\_\allowbreak{}n H\textasciicircum{}r(F(Q\_\allowbreak{}n)\allowbreak{})\allowbreak{}.

Equations (3.1)\allowbreak{} and (3.2)\allowbreak{} identify H\textasciicircum{}r(F(α)\allowbreak{})\allowbreak{} with the colimit of the compatible isomorphisms H\textasciicircum{}r(θ\_\allowbreak{}n)\allowbreak{},\allowbreak{} so it is an isomorphism for every r. Thus F sends every quasi-isomorphism between members of S to a quasi-isomorphism. The same existence-and-computation criterion,\allowbreak{} now in the full homotopy category,\allowbreak{} constructs LF everywhere,\allowbreak{} with every member of S computing it. This proves the strengthened assertion.

We have retained exact sequential colimits in both categories. No claim that exact coproducts alone make all sequential colimits exact is needed. The weaker preservation condition on F suffices because the systems being used are split in each degree,\allowbreak{} not because arbitrary colimits were ignored.

For a concrete instance beyond right exact functors,\allowbreak{} take abelian groups,\allowbreak{} P the free groups,\allowbreak{} and F(M)\allowbreak{}\allowbreak{}=M[ℓ]\allowbreak{}\allowbreak{}=ker(ℓ:\allowbreak{}M\allowbreak{}→M)\allowbreak{},\allowbreak{} for a prime ℓ. F is additive and commutes with countable coproducts,\allowbreak{} since an element of a direct sum is killed by ℓ exactly when each of its finitely many nonzero components is. It is zero on free groups,\allowbreak{} so it sends all the required bounded above P complexes to the zero complex. It is not right exact:\allowbreak{} the epimorphism Z\allowbreak{}→Z/\allowbreak{}ℓ induces 0\allowbreak{}→Z/\allowbreak{}ℓ. Nevertheless the preceding proof constructs LF\allowbreak{}=0 on all of D(Ab)\allowbreak{},\allowbreak{} since every term in (2.1)\allowbreak{} is free. This is its left derived functor,\allowbreak{} not its right derived functor. Without right exactness there is no claim that H\textasciicircum{}0(LF(M[0]\allowbreak{})\allowbreak{})\allowbreak{} equals F(M)\allowbreak{}.

The received direction corrections are now semantically reconciled. MC-STK-ERR-0876 makes α run from P to Q,\allowbreak{} as (3.1)\allowbreak{} and (3.2)\allowbreak{} require,\allowbreak{} and writes its resulting map as H\textasciicircum{}r(F(α)\allowbreak{})\allowbreak{}. P02-E0070 proposes reversing the arrow rather than exchanging the displayed objects; both yield the same typed map. The three admitted operations are sufficient. P02-E0071 and the French DERIVED-065 report are completed with this entire argument,\allowbreak{} including compatibility in n.

\subsection{4. The dual system with the exact reversed indexing}\label{proofs-9206-9744-the-dual-system-with-the-exact-reversed-indexing}

Apply Section 1 in A\^{}op using the complex D(K)\allowbreak{}\textasciicircum{}r\allowbreak{}=(K\textasciicircum{}\{-r\})\allowbreak{}\textasciicircum{}op and d\_\allowbreak{}\{D(K)\allowbreak{}\}\textasciicircum{}r\allowbreak{}=(d\_\allowbreak{}K\textasciicircum{}\{-r-1\})\allowbreak{}\textasciicircum{}op. This is a complex because its two consecutive differentials compose to the opposite of d\_\allowbreak{}K². It sends a chain map K\allowbreak{}→L to D(L)\allowbreak{}\allowbreak{}→D(K)\allowbreak{},\allowbreak{} and H\textasciicircum{}r(D(K)\allowbreak{})\allowbreak{}\allowbreak{}=H\textasciicircum{}\{-r\}(K)\allowbreak{}\textasciicircum{}op:\allowbreak{} kernels and cokernels exchange under passage to the opposite category. Direct sums become finite products,\allowbreak{} and epimorphisms become monomorphisms. The actual canonical truncations satisfy

D(τ\_\allowbreak{}\{\textbackslash{}geq-n\}K)\allowbreak{}\allowbreak{}=τ\_\allowbreak{}\{\textbackslash{}leq n\}D(K)\allowbreak{}.

At the endpoint this equality identifies the kernel in A\^{}op with the cokernel of d\_\allowbreak{}K\textasciicircum{}\{-n-1\}; away from the endpoint it is the identity. Apply D back to the constructed diagram. It gives τ\_\allowbreak{}\{\textbackslash{}geq-n\}K\allowbreak{}→I\_\allowbreak{}n with termwise monomorphic quasi-isomorphisms,\allowbreak{} I\_\allowbreak{}n\textasciicircum{}r\allowbreak{}=0 for r<-n,\allowbreak{} and compatible termwise split epimorphisms I\_\allowbreak{}\{n\allowbreak{}+1\}\allowbreak{}→I\_\allowbreak{}n whose kernel terms belong to the given class I. All squares commute as chain maps. This proves the original inverse-system lemma and the lower-bound part of UNDER-019. No exact inverse-limit hypothesis has been asserted,\allowbreak{} and no conclusion about K\allowbreak{}→lim I\_\allowbreak{}n is deduced from this construction alone.

\subsection{5. Derived adjunction with explicit representatives and naturality}\label{proofs-9206-9744-derived-adjunction-with-explicit-representatives-and-naturality}

Retain the source\textquotesingle s exact functors F:\allowbreak{}D\allowbreak{}→D' and G:\allowbreak{}D'\allowbreak{}→D with G left adjoint to F,\allowbreak{} and the multiplicative systems S,\allowbreak{}S'. Fix K and M at which RF and LG respectively exist. Let s:\allowbreak{}K\allowbreak{}→I range over S,\allowbreak{} and let t:\allowbreak{}P\allowbreak{}→M range over S\textquotesingle. The values F(I)\allowbreak{} are an essentially constant ind-system in (S')\allowbreak{}\textasciicircum{}\{-1\}D' with value RF(K)\allowbreak{},\allowbreak{} and G(P)\allowbreak{} form an essentially constant pro-system in S\^{}\{-1\}D with value LG(M)\allowbreak{}.

Write c\_\allowbreak{}s:\allowbreak{}F(I)\allowbreak{}\allowbreak{}→RF(K)\allowbreak{} for the ind-system structure maps and d\_\allowbreak{}t:\allowbreak{}LG(M)\allowbreak{}\allowbreak{}→G(P)\allowbreak{} for the pro-system structure maps,\allowbreak{} in the respective localized categories. The Hom formulas for essentially constant systems and for the calculus of fractions identify each side of

Hom\_\allowbreak{}\{(S')\allowbreak{}\textasciicircum{}\{-1\}D'\}(M,\allowbreak{}RF(K)\allowbreak{})\allowbreak{} ≅ Hom\_\allowbreak{}\{S\textasciicircum{}\{-1\}D\}(LG(M)\allowbreak{},\allowbreak{}K)\allowbreak{} (5.1)\allowbreak{}

with the filtered colimit,\allowbreak{} in the two indices s and t,\allowbreak{} of Hom\_\allowbreak{}\{D'\}(P,\allowbreak{}F(I)\allowbreak{})\allowbreak{}≅Hom\_\allowbreak{}D(G(P)\allowbreak{},\allowbreak{}I)\allowbreak{}. The first index is K/\allowbreak{}S and the second is (S'/\allowbreak{}M)\allowbreak{}\textasciicircum{}op. Thus they are independent filtered indices; changing the order of their colimits is justified by the universal property of a cocone on their product,\allowbreak{} not by an exchange of a limit with a colimit.

More concretely,\allowbreak{} if h:\allowbreak{}P\allowbreak{}→F(I)\allowbreak{} and h\textasciicircum{}{\Rsixtyonesymbolfont ♭}:\allowbreak{}G(P)\allowbreak{}\allowbreak{}→I are adjoint,\allowbreak{} then (5.1)\allowbreak{} sends

c\_\allowbreak{}s h t\^{}\{-1\} to s\textasciicircum{}\{-1\}h\textasciicircum{}{\Rsixtyonesymbolfont ♭}d\_\allowbreak{}t. (5.2)\allowbreak{}

Changing s or t along a refinement gives the same localized maps:\allowbreak{} the naturality identities of the original adjunction identify the refined h\textasciicircum{}{\Rsixtyonesymbolfont ♭},\allowbreak{} and the structure maps c\_\allowbreak{}s,\allowbreak{}d\_\allowbreak{}t are compatible with those refinements. Conversely any arrow on the right is represented,\allowbreak{} by essential constancy of the pro-system,\allowbreak{} by some map G(P)\allowbreak{}\allowbreak{}→K in the localization. A left fraction writes it as s\^{}\{-1\}b with b:\allowbreak{}G(P)\allowbreak{}\allowbreak{}→I. Adjoint b back to h:\allowbreak{}P\allowbreak{}→F(I)\allowbreak{} and use the left expression in (5.2)\allowbreak{}. Equality of fractions or of representatives is exactly equality after a refinement,\allowbreak{} so these constructions are inverse.

These formulas prove naturality for actual arrows in K and M. A localized arrow is a composite of such arrows and inverses of arrows in S or S\textquotesingle. Naturality for those invertible arrows implies naturality for their inverses by multiplying the naturality square by the inverse isomorphisms. Hence (5.1)\allowbreak{} is functorial on the full domains where the two derived functors are defined,\allowbreak{} as stated by the source.

MC-STK-ERR-0877 correctly replaces the two undefined abelian derived categories in the general proof by (S')\allowbreak{}\textasciicircum{}\{-1\}D' and S\^{}\{-1\}D. The two P02 reports and French DERIVED-066 are reconciled by (5.1)\allowbreak{}–\allowbreak{}(5.2)\allowbreak{}.

For adjoint functors between abelian categories,\allowbreak{} G preserves zero objects and finite coproducts,\allowbreak{} while F preserves zero objects and finite products. Since the latter are biproducts,\allowbreak{} both functors are additive. Their degreewise adjunction restricts to chain maps:\allowbreak{} naturality in each variable equates the two chain-map equations. It also equates the homotopy relations,\allowbreak{} by additivity and the same naturality for their degree -1 components. Thus it descends to an adjunction on homotopy categories. Both induced functors preserve the actual shift and cone formulas,\allowbreak{} so they are exact there. This proves the abelian-category specialization.

If RF and LG are everywhere defined,\allowbreak{} define the unit η\_\allowbreak{}M:\allowbreak{}M\allowbreak{}→RF LG(M)\allowbreak{} as the preimage of id\_\allowbreak{}\{LG(M)\allowbreak{}\} under (5.1)\allowbreak{},\allowbreak{} and the counit ε\_\allowbreak{}K:\allowbreak{}LG RF(K)\allowbreak{}\allowbreak{}→K as the image of id\_\allowbreak{}\{RF(K)\allowbreak{}\}. Naturality in M gives ε\_\allowbreak{}\{LG(M)\allowbreak{}\} LG(η\_\allowbreak{}M)\allowbreak{}\allowbreak{}=id\_\allowbreak{}\{LG(M)\allowbreak{}\}. For the other identity,\allowbreak{} naturality in K sends the image of η\_\allowbreak{}\{RF(K)\allowbreak{}\} under (5.1)\allowbreak{} to ε\_\allowbreak{}K. But the image of id\_\allowbreak{}\{RF(K)\allowbreak{}\} is also ε\_\allowbreak{}K,\allowbreak{} and (5.1)\allowbreak{} is injective; therefore RF(ε\_\allowbreak{}K)\allowbreak{}η\_\allowbreak{}\{RF(K)\allowbreak{}\}\allowbreak{}=id\_\allowbreak{}\{RF(K)\allowbreak{}\}. This proves the actual adjunction,\allowbreak{} including its two triangle identities. No such composition is asserted in a partial domain unless its terms exist.

\subsection{\texorpdfstring{6. K-injectivity,\allowbreak{} localization,\allowbreak{} and bounded below complexes}{6.   and bounded below complexes}}\label{proofs-9206-9744-k-injectivity-localization-and-bounded-below-complexes}

The defining vanishing for maps from acyclic complexes also holds for all their shifts. If I is K-injective and u:\allowbreak{}M\allowbreak{}→N is a quasi-isomorphism,\allowbreak{} apply Hom\_\allowbreak{}K(-,\allowbreak{}I)\allowbreak{} to the cone triangle. The neighboring Hom groups of C(u)\allowbreak{} and its shifts vanish,\allowbreak{} so precomposition by u is bijective. Conversely,\allowbreak{} this bijectivity identifies a fraction f s\textasciicircum{}\{-1\}:\allowbreak{}N\allowbreak{}→I with the unique g:\allowbreak{}N\allowbreak{}→I in K satisfying g s\allowbreak{}=f. It is independent of the fraction because a common refinement identifies both such g. It is also injective:\allowbreak{} if g maps to zero in D,\allowbreak{} some quasi-isomorphism s makes g s zero in K,\allowbreak{} whence g\allowbreak{}=0 by that same bijectivity. Thus Hom\_\allowbreak{}K(N,\allowbreak{}I)\allowbreak{}\allowbreak{}→Hom\_\allowbreak{}D(N,\allowbreak{}I)\allowbreak{} is an isomorphism. If this last condition holds and N is acyclic,\allowbreak{} N is zero in D,\allowbreak{} so its Hom to I is zero in K. This proves all three equivalences.

For a distinguished triangle in K,\allowbreak{} apply Hom\_\allowbreak{}K(N,\allowbreak{}-)\allowbreak{} for acyclic N and each of its shifts. Its long exact sequence implies that if two terms are K-injective,\allowbreak{} the third is as well. Shifts of a K-injective complex are K-injective because Hom\_\allowbreak{}K(N,\allowbreak{}I[r]\allowbreak{})\allowbreak{}\allowbreak{}=Hom\_\allowbreak{}K(N[-r]\allowbreak{},\allowbreak{}I)\allowbreak{}.

For completeness let I have injective terms and vanish below a. Given f:\allowbreak{}N\allowbreak{}→I from an acyclic N,\allowbreak{} construct h\textasciicircum{}r:\allowbreak{}N\textasciicircum{}r\allowbreak{}→I\textasciicircum{}\{r-1\} with f\allowbreak{}=d\_\allowbreak{}I h\allowbreak{}+h d\_\allowbreak{}N. Set h\textasciicircum{}r\allowbreak{}=0 for r≤a. At r\allowbreak{}=a,\allowbreak{} f\^{}a kills im(d\_\allowbreak{}N\textasciicircum{}\{a-1\})\allowbreak{}\allowbreak{}=ker(d\_\allowbreak{}N\textasciicircum{}a)\allowbreak{},\allowbreak{} so it factors through im(d\_\allowbreak{}N\textasciicircum{}a)\allowbreak{}⊆N\textasciicircum{}\{a\allowbreak{}+1\}; injectivity of I\^{}a extends this factor to h\textasciicircum{}\{a\allowbreak{}+1\}. Inductively f\textasciicircum{}r-d\_\allowbreak{}I\textasciicircum{}\{r-1\}h\textasciicircum{}r kills ker(d\_\allowbreak{}N\textasciicircum{}r)\allowbreak{}:\allowbreak{} on the image of d\_\allowbreak{}N\textasciicircum{}\{r-1\} this follows from the chain equation for f and the preceding homotopy equation,\allowbreak{} with d\_\allowbreak{}I²\allowbreak{}=0. Extend again to obtain h\textasciicircum{}\{r\allowbreak{}+1\}. The construction proves every component of the homotopy equation. Hence f is zero in K and I is K-injective.

\subsection{7. The Hom-complex index correction and the product theorem}\label{proofs-9206-9744-the-hom-complex-index-correction-and-the-product-theorem}

For arbitrary complexes K,\allowbreak{}I define the full Hom complex of abelian groups by

C\textasciicircum{}r\allowbreak{}=∏\_\allowbreak{}\{b\allowbreak{}∈Z\}Hom\_\allowbreak{}A(K\textasciicircum{}b,\allowbreak{}I\textasciicircum{}\{b\allowbreak{}+r\})\allowbreak{},\allowbreak{} (d\_\allowbreak{}C\textasciicircum{}r f)\allowbreak{}\textasciicircum{}b\allowbreak{}=d\_\allowbreak{}I\textasciicircum{}\{b\allowbreak{}+r\}f\textasciicircum{}b-(-1)\allowbreak{}\textasciicircum{}r f\textasciicircum{}\{b\allowbreak{}+1\}d\_\allowbreak{}K\textasciicircum{}b. (7.1)\allowbreak{}

Expanding d\_\allowbreak{}C\textasciicircum{}\{r\allowbreak{}+1\}d\_\allowbreak{}C\textasciicircum{}r gives the two terms containing d\_\allowbreak{}I² and d\_\allowbreak{}K²,\allowbreak{} which vanish,\allowbreak{} and two terms containing d\_\allowbreak{}I f d\_\allowbreak{}K whose coefficients -(-1)\allowbreak{}\textasciicircum{}r and -(-1)\allowbreak{}\textasciicircum{}\{r\allowbreak{}+1\} cancel. Thus it is a complex. A degree r cocycle is exactly a chain map K\allowbreak{}→I[r]\allowbreak{},\allowbreak{} since the target differential is (-1)\allowbreak{}\textasciicircum{}r d\_\allowbreak{}I. If f\allowbreak{}=d\_\allowbreak{}C\textasciicircum{}\{r-1\}h,\allowbreak{} the corresponding chain homotopy has components (-1)\allowbreak{}\textasciicircum{}r h\textasciicircum{}b:\allowbreak{}K\textasciicircum{}b\allowbreak{}→I[r]\allowbreak{}\textasciicircum{}\{b-1\}. Indeed its boundary is

(-1)\allowbreak{}\textasciicircum{}r d\_\allowbreak{}I((-1)\allowbreak{}\textasciicircum{}r h)\allowbreak{}\allowbreak{}+((-1)\allowbreak{}\textasciicircum{}r h)\allowbreak{}d\_\allowbreak{}K \allowbreak{}= d\_\allowbreak{}I h\allowbreak{}+(-1)\allowbreak{}\textasciicircum{}r h d\_\allowbreak{}K\allowbreak{}=d\_\allowbreak{}C\textasciicircum{}\{r-1\}h.

Conversely any such homotopy arises this way. Therefore H\textasciicircum{}r(C)\allowbreak{}\allowbreak{}=Hom\_\allowbreak{}K(K,\allowbreak{}I[r]\allowbreak{})\allowbreak{},\allowbreak{} including H\textasciicircum{}0(C)\allowbreak{}\allowbreak{}=Hom\_\allowbreak{}K(K,\allowbreak{}I)\allowbreak{}.

The six factors in original 9712–\allowbreak{}9714 and 9720–\allowbreak{}9722 have K\^{}\{-b\} where (7.1)\allowbreak{} requires K\^{}b. DERIVED-NEW-032 corrects them and makes the full Hom complex explicit. This is not an innocuous dummy-index choice:\allowbreak{} changing b to -b in all factors would also change the exponents of I. For example,\allowbreak{} take K\allowbreak{}=I\allowbreak{}=Q[-1]\allowbreak{} in complexes of Q-vector spaces,\allowbreak{} so both are Q in degree 1. Their identity gives Hom\_\allowbreak{}K(K,\allowbreak{}I)\allowbreak{}\allowbreak{}=Q. In the printed middle product,\allowbreak{} nonzero K\^{}\{-b\} requires b\allowbreak{}=-1,\allowbreak{} whereas nonzero I\^{}b requires b\allowbreak{}=1. The product is zero. In this example the two adjacent printed groups are zero too,\allowbreak{} so no differential could repair the asserted middle cohomology. I is a bounded below complex of injectives,\allowbreak{} so this example is within the stated lemma\textquotesingle s setting.

Now let \{I\_\allowbreak{}t\} be the given set-indexed family,\allowbreak{} with the termwise products I\textasciicircum{}m\allowbreak{}=∏\_\allowbreak{}t I\_\allowbreak{}t\textasciicircum{}m existing in A. The product universal property defines d\_\allowbreak{}I by its components d\_\allowbreak{}\{I\_\allowbreak{}t\}; its square is zero because all its projections are zero. Formula (7.1)\allowbreak{} gives the actual isomorphism of complexes

Hom\textasciicircum{}•(K,\allowbreak{}I)\allowbreak{} ≅ ∏\_\allowbreak{}t Hom\textasciicircum{}•(K,\allowbreak{}I\_\allowbreak{}t)\allowbreak{}.

In degree r this is the canonical rearrangement ∏\_\allowbreak{}b Hom(K\textasciicircum{}b,\allowbreak{}∏\_\allowbreak{}t I\_\allowbreak{}t\textasciicircum{}\{b\allowbreak{}+r\})\allowbreak{}≅∏\_\allowbreak{}t∏\_\allowbreak{}b Hom(K\textasciicircum{}b,\allowbreak{}I\_\allowbreak{}t\textasciicircum{}\{b\allowbreak{}+r\})\allowbreak{}. Both terms in the differential agree after each projection,\allowbreak{} so this is an isomorphism of complexes,\allowbreak{} not just of their graded objects. Products of abelian groups are exact:\allowbreak{} kernels are coordinatewise,\allowbreak{} and a family of surjections is surjective by choosing a preimage in each coordinate. Hence cohomology commutes with this product. In particular,\allowbreak{} for every K and r,\allowbreak{}

Hom\_\allowbreak{}K(K,\allowbreak{}I[r]\allowbreak{})\allowbreak{} ≅ ∏\_\allowbreak{}t Hom\_\allowbreak{}K(K,\allowbreak{}I\_\allowbreak{}t[r]\allowbreak{})\allowbreak{}. (7.2)\allowbreak{}

For acyclic K each right-hand group vanishes because I\_\allowbreak{}t is K-injective and K[-r]\allowbreak{} is acyclic. Thus the full Hom complexes are acyclic,\allowbreak{} which justifies the source\textquotesingle s acyclicity step at every degree,\allowbreak{} not only at zero. Equation (7.2)\allowbreak{} proves I is K-injective. Applying Section 6 to I and the I\_\allowbreak{}t,\allowbreak{} now for arbitrary K,\allowbreak{} yields

Hom\_\allowbreak{}D(K,\allowbreak{}I)\allowbreak{} ≅ ∏\_\allowbreak{}t Hom\_\allowbreak{}D(K,\allowbreak{}I\_\allowbreak{}t)\allowbreak{}.

These isomorphisms are induced by the actual product projections and are natural in K,\allowbreak{} so they prove the universal property of the product in D(A)\allowbreak{}. Only the displayed term products in A were required. No exactness of products in A was assumed:\allowbreak{} the exact products used above are products of abelian groups. The theorem remains true.

The received article correction at 9709,\allowbreak{} MC-STK-ERR-0883,\allowbreak{} is already present. P02-E0074 and French DERIVED-072 are reconciled separately from the new Hom-index finding; neither report previously states it.

\subsection{8. Receiving statements and the review boundary}\label{proofs-9206-9744-receiving-statements-and-the-review-boundary}

The special direct-system uses in current more-algebra.tex 14989–\allowbreak{}15023 and cohomology.tex 6465–\allowbreak{}6495 require precisely the strict compatible epimorphic quasi-isomorphisms and split term inclusions proved in Sections 1–\allowbreak{}2. The improved bounds and the same colimit map apply to those choices of P. Their subsequent K-flatness arguments are separate uses of their cited lemmas; this bounded check does not replace those proofs.

The inverse system used in current cohomology.tex 10284–\allowbreak{}10310 has the strict maps and lower bounds of Section 4. Its later inverse-limit claim is not certified by the direct-system colimit calculation. The missing word before ``produce a diagram'' is routed as a reading observation for that chapter\textquotesingle s own report join,\allowbreak{} not counted as an additional finding here.

Current sites-cohomology.tex 9337–\allowbreak{}9350 invokes the unbounded criterion for a left adjoint g\_\allowbreak{}!,\allowbreak{} which is right exact and preserves all colimits. The new additive criterion therefore preserves those verified hypotheses; the later verification of its acyclicity condition is outside the excerpt. The further cited reductions at 10439–\allowbreak{}10450 and 10836–\allowbreak{}10850 motivate the explicit compatible-stage proof of Section 3. These excerpts do not certify the entirety of either receiving theorem.

Current injectives.tex 2038–\allowbreak{}2042 constructs a derived product from K-injective resolutions. Section 7 supplies the corrected Hom calculation for that exact invocation. The product projections and derived universal property are unchanged. No new receiving source operation is needed.

Eight more physical reports are complete in groups DERIVED-RECON-070,\allowbreak{} 071 and 072. All eight correspond to the three existing units MC-STK-ERR-0876,\allowbreak{} 0877 and 0883,\allowbreak{} comprising six already applied operations. Totals:\allowbreak{} 125/\allowbreak{}176 reports,\allowbreak{} 72 groups,\allowbreak{} 116 already corrected reports,\allowbreak{} eight missing copyedit/\allowbreak{}notation reports and one missing mathematical-index report; 67 prior Derived units and 84 prior operations reconciled. NEW-032 adds four proposed operations,\allowbreak{} giving 61 cumulative operations. UNDER-019 and UNDER-020 bring the editorial consequence count to twenty. Semantic review ends at 9744. Sequential reading has reached 9760; the next incomplete lemma starts at 9746 and unread source begins at 9761. The overall goal and received-work priority remain active.
\clearpage
\section{K-injective limits and finite cohomological dimension}\label{proofs-9746-10132-k-injective-limits-and-finite-cohomological-dimension}

Private source-order review of original derived.tex 9746–\allowbreak{}10132. Authority:\allowbreak{} Stacks a04446e5\allowbreak{}7ec1fbc2\allowbreak{}52a871af\allowbreak{}cec7752f\allowbreak{}b2807b14. Receiving public edition:\allowbreak{} 95a51593\allowbreak{}aceb247a\allowbreak{}c7a18311\allowbreak{}1678c3a0\allowbreak{}f8b8f4b5. The accompanying source-reading and operation records bind the files. Editorial consequences are additional statements with complete arguments; they are not substituted for the original author\textquotesingle s hypotheses. No novelty is claimed.

Human source:\allowbreak{} \href{https://github.com/stacks/stacks-project/blob/a04446e57ec1fbc252a871afcec7752fb2807b14/derived.tex\#L9746}{The Stacks Project,\allowbreak{} original Derived Categories TeX},\allowbreak{} with its cited Homology results. This review continues the full Hom-complex and shift calculation in PROOFS\_\allowbreak{}9206\_\allowbreak{}9744.md,\allowbreak{} Section 7.

\subsection{1. Why a K-injective complex computes every right derived functor}\label{proofs-9746-10132-why-a-k-injective-complex-computes-every-right-derived-functor}

Let I be K-injective and let s:\allowbreak{}I\allowbreak{}→N be a quasi-isomorphism in K(A)\allowbreak{}. Precomposition gives a bijection Hom\_\allowbreak{}K(N,\allowbreak{}I)\allowbreak{}\allowbreak{}→Hom\_\allowbreak{}K(I,\allowbreak{}I)\allowbreak{}. Thus there exists a unique r:\allowbreak{}N\allowbreak{}→I with r s\allowbreak{}=id\_\allowbreak{}I in K(A)\allowbreak{}. It is also a quasi-isomorphism:\allowbreak{} H(r)\allowbreak{} is inverse to H(s)\allowbreak{}. Therefore id\_\allowbreak{}I is a final object of I/\allowbreak{}Qis,\allowbreak{} including the uniqueness required for finality,\allowbreak{} not only the existence of a retraction.

For an arbitrary functor F on K(A)\allowbreak{},\allowbreak{} the system sending s:\allowbreak{}I\allowbreak{}→N to F(N)\allowbreak{} has value F(I)\allowbreak{} at that final object and is an essentially constant ind-object. Indeed the unique arrows to the final object provide its cocone,\allowbreak{} and the Hom colimit from any object is the Hom to F(I)\allowbreak{}. In the source\textquotesingle s exact-functor setting this proves I computes RF. If K is isomorphic in D(A)\allowbreak{} to I,\allowbreak{} the isomorphism K\allowbreak{}→I is represented by an actual map in K(A)\allowbreak{},\allowbreak{} since maps to a K-injective complex agree in K and D. Its induced cohomology maps are isomorphisms. Hence the computation transfers to K by that quasi-isomorphism.

If every complex maps by a quasi-isomorphism to a K-injective complex,\allowbreak{} the computation criterion gives RF everywhere. A quasi-isomorphism s:\allowbreak{}I\allowbreak{}→J between K-injectives is an isomorphism already in K:\allowbreak{} the unique retraction r satisfies r s\allowbreak{}=id\_\allowbreak{}I,\allowbreak{} and precomposition by s is injective on Hom\_\allowbreak{}K(J,\allowbreak{}J)\allowbreak{},\allowbreak{} so (s r)\allowbreak{}s\allowbreak{}=s\allowbreak{}=id\_\allowbreak{}J s implies s r\allowbreak{}=id\_\allowbreak{}J. Every functor F therefore takes this map to an isomorphism. This proves the two source lemmas without imposing injectivity on individual terms.

\subsection{2. The split inverse system and compatible primitives}\label{proofs-9746-10132-the-split-inverse-system-and-compatible-primitives}

Let …\allowbreak{}→I\_\allowbreak{}2\allowbreak{}→I\_\allowbreak{}1 be the source tower,\allowbreak{} with K-injective stages,\allowbreak{} split surjections in each degree,\allowbreak{} and existing termwise limit I. For an acyclic M let D\_\allowbreak{}n\allowbreak{}=Hom\textasciicircum{}•(M,\allowbreak{}I\_\allowbreak{}n)\allowbreak{},\allowbreak{} using the full differential

D\_\allowbreak{}n\textasciicircum{}r\allowbreak{}=∏\_\allowbreak{}a Hom(M\textasciicircum{}a,\allowbreak{}I\_\allowbreak{}n\textasciicircum{}\{a\allowbreak{}+r\})\allowbreak{},\allowbreak{} (d f)\allowbreak{}\textasciicircum{}a\allowbreak{}=d\_\allowbreak{}\{I\_\allowbreak{}n\}\textasciicircum{}\{a\allowbreak{}+r\}f\textasciicircum{}a-(-1)\allowbreak{}\textasciicircum{}r f\textasciicircum{}\{a\allowbreak{}+1\}d\_\allowbreak{}M\textasciicircum{}a.

Every D\_\allowbreak{}n is acyclic by the all-shift K-injectivity calculation in PART14. The transition maps on D\_\allowbreak{}n\textasciicircum{}r are surjective:\allowbreak{} lift each component through the chosen split epimorphism of I terms. Representable Hom and products preserve limits,\allowbreak{} and the differentials agree after each projection. Consequently Hom\textasciicircum{}•(M,\allowbreak{}I)\allowbreak{}\allowbreak{}=lim\_\allowbreak{}n D\_\allowbreak{}n as actual complexes of abelian groups.

Here is an explicit proof of its degree-zero exactness. A compatible cycle z\allowbreak{}=(z\_\allowbreak{}n)\allowbreak{} admits primitives h\_\allowbreak{}n\allowbreak{}∈D\_\allowbreak{}n\textasciicircum{}\{-1\}. Choose h\_\allowbreak{}1 first. Given a compatible primitive through stage n,\allowbreak{} choose any k\allowbreak{}∈D\_\allowbreak{}\{n\allowbreak{}+1\}\textasciicircum{}\{-1\} with d k\allowbreak{}=z\_\allowbreak{}\{n\allowbreak{}+1\}. Its discrepancy t\_\allowbreak{}n(k)\allowbreak{}-h\_\allowbreak{}n is a cycle in D\_\allowbreak{}n\textasciicircum{}\{-1\},\allowbreak{} where t\_\allowbreak{}n is the transition. Acyclicity gives y\_\allowbreak{}n\allowbreak{}∈D\_\allowbreak{}n\textasciicircum{}\{-2\} with d y\_\allowbreak{}n\allowbreak{}=t\_\allowbreak{}n(k)\allowbreak{}-h\_\allowbreak{}n. Lift y\_\allowbreak{}n to y\allowbreak{}∈D\_\allowbreak{}\{n\allowbreak{}+1\}\textasciicircum{}\{-2\} by surjectivity and put h\_\allowbreak{}\{n\allowbreak{}+1\}\allowbreak{}=k-d y. Then d h\_\allowbreak{}\{n\allowbreak{}+1\}\allowbreak{}=z\_\allowbreak{}\{n\allowbreak{}+1\} and t\_\allowbreak{}n(h\_\allowbreak{}\{n\allowbreak{}+1\})\allowbreak{}\allowbreak{}=h\_\allowbreak{}n. The resulting compatible primitive proves H\^{}0(lim D\_\allowbreak{}n)\allowbreak{}\allowbreak{}=0,\allowbreak{} and hence I is K-injective. Shifting this argument proves exactness in every degree.

This is exactly the use of the degree -2 system in the source\textquotesingle s invocation of Homology\textquotesingle s Mittag-Leffler lemma. It does not assume that inverse limits are exact in A. In particular it proves I is K-injective,\allowbreak{} but does not prove that an unrelated compatible map K\allowbreak{}→I is a quasi-isomorphism. The source explicitly keeps that latter issue open for a later lemma.

\subsection{3. The exact obstruction for a tower without split maps}\label{proofs-9746-10132-the-exact-obstruction-for-a-tower-without-split-maps}

DERIVED-UNDER-021 records a stronger conclusion of the preceding argument. Retain K-injective stages and existing termwise limits,\allowbreak{} but allow arbitrary transition maps. For each acyclic M use D\_\allowbreak{}n as above and put Z\_\allowbreak{}n\textasciicircum{}r\allowbreak{}=ker(d:\allowbreak{}D\_\allowbreak{}n\textasciicircum{}r\allowbreak{}→D\_\allowbreak{}n\textasciicircum{}\{r\allowbreak{}+1\})\allowbreak{}. For a tower A\_\allowbreak{}n of abelian groups write

Δ\_\allowbreak{}A:\allowbreak{}∏\_\allowbreak{}n A\_\allowbreak{}n\allowbreak{}→∏\_\allowbreak{}n A\_\allowbreak{}n,\allowbreak{} (x\_\allowbreak{}n)\allowbreak{}\allowbreak{}↦(x\_\allowbreak{}n-t\_\allowbreak{}n x\_\allowbreak{}\{n\allowbreak{}+1\})\allowbreak{},\allowbreak{} C\textasciicircum{}1(A\_\allowbreak{}n)\allowbreak{}\allowbreak{}=coker(Δ\_\allowbreak{}A)\allowbreak{}.

This defines the obstruction groups directly; no exactness of a limit functor is assumed. The exact sequences 0\allowbreak{}→Z\_\allowbreak{}n\textasciicircum{}\{-1\}\allowbreak{}→D\_\allowbreak{}n\textasciicircum{}\{-1\}\allowbreak{}→Z\_\allowbreak{}n\textasciicircum{}0\allowbreak{}→0 and exactness of products of abelian groups give,\allowbreak{} by the kernel-cokernel sequence for the vertical maps Δ,\allowbreak{} a natural identification

Hom\_\allowbreak{}K(M,\allowbreak{}I)\allowbreak{} ≅ ker(C\textasciicircum{}1(Z\_\allowbreak{}n\textasciicircum{}\{-1\})\allowbreak{}\allowbreak{}→C\textasciicircum{}1(D\_\allowbreak{}n\textasciicircum{}\{-1\})\allowbreak{})\allowbreak{}. (3.1)\allowbreak{}

Explicitly,\allowbreak{} for z\allowbreak{}=(z\_\allowbreak{}n)\allowbreak{}\allowbreak{}∈lim Z\_\allowbreak{}n\textasciicircum{}0 choose primitives h\_\allowbreak{}n with d h\_\allowbreak{}n\allowbreak{}=z\_\allowbreak{}n. Their discrepancies e\_\allowbreak{}n\allowbreak{}=h\_\allowbreak{}n-t\_\allowbreak{}n h\_\allowbreak{}\{n\allowbreak{}+1\} lie in Z\_\allowbreak{}n\textasciicircum{}\{-1\}. Changing the primitives by cycles changes (e\_\allowbreak{}n)\allowbreak{} by Δ\_\allowbreak{}Z of those cycles,\allowbreak{} and changing z by the boundary of a compatible h changes its class by zero. Moreover e\allowbreak{}=Δ\_\allowbreak{}D h,\allowbreak{} so its class maps to zero in C\textasciicircum{}1(D\_\allowbreak{}n\textasciicircum{}\{-1\})\allowbreak{}. Conversely,\allowbreak{} if a cycle family e represents an element of that kernel,\allowbreak{} write e\allowbreak{}=Δ\_\allowbreak{}D h. Then (d h\_\allowbreak{}n)\allowbreak{} is a compatible cycle family. If the class of e is zero in C\textasciicircum{}1(Z\_\allowbreak{}n\textasciicircum{}\{-1\})\allowbreak{},\allowbreak{} subtract the corresponding cycle family from h to make it compatible,\allowbreak{} so z is a boundary. These operations prove both the bijection and its naturality,\allowbreak{} including the exact map.

Thus I is K-injective precisely when the right side of (3.1)\allowbreak{} vanishes for every acyclic M. This is an explicit obstruction,\allowbreak{} not an assertion that arbitrary towers work. A sufficient condition is that for every acyclic M the tower D\_\allowbreak{}n\textasciicircum{}\{-2\} is Mittag-Leffler. Its quotient tower Z\_\allowbreak{}n\textasciicircum{}\{-1\}\allowbreak{}=im(D\_\allowbreak{}n\textasciicircum{}\{-2\}\allowbreak{}→D\_\allowbreak{}n\textasciicircum{}\{-1\})\allowbreak{} is then Mittag-Leffler too.

For completeness,\allowbreak{} a countable Mittag-Leffler tower A\_\allowbreak{}n has C\textasciicircum{}1(A\_\allowbreak{}n)\allowbreak{}\allowbreak{}=0. Let B\_\allowbreak{}n be the eventual image in A\_\allowbreak{}n; Mittag-Leffler makes it an attained image,\allowbreak{} and B\_\allowbreak{}\{n\allowbreak{}+1\}\allowbreak{}→B\_\allowbreak{}n is surjective by choosing a stage large enough for both images to have stabilized. The quotient tower A\_\allowbreak{}n/\allowbreak{}B\_\allowbreak{}n is pro-zero. On a pro-zero tower Q,\allowbreak{} the expression y\_\allowbreak{}n\allowbreak{}=Σ\_\allowbreak{}\{k≥n\}t\_\allowbreak{}\{n,\allowbreak{}k\}x\_\allowbreak{}k is a finite sum at each fixed n and satisfies Δ\_\allowbreak{}Q y\allowbreak{}=x. On a surjective tower B,\allowbreak{} solve Δ\_\allowbreak{}B y\allowbreak{}=x recursively,\allowbreak{} starting with any y\_\allowbreak{}1 and lifting y\_\allowbreak{}n-x\_\allowbreak{}n to y\_\allowbreak{}\{n\allowbreak{}+1\}. Given x in ∏A\_\allowbreak{}n,\allowbreak{} solve first in the quotient,\allowbreak{} lift a solution to ∏A\_\allowbreak{}n,\allowbreak{} then correct its remaining error in ∏B\_\allowbreak{}n. This proves Δ\_\allowbreak{}A is surjective. Applied to Z\_\allowbreak{}n\textasciicircum{}\{-1\},\allowbreak{} it makes the right side of (3.1)\allowbreak{} zero.

This sufficient condition includes some nonsplit towers. For K-injective I\_\allowbreak{}0,\allowbreak{}J let every stage be I\_\allowbreak{}0\allowbreak{}⊕J and every transition diag(id,\allowbreak{}0)\allowbreak{}. It is not termwise surjective where J is nonzero,\allowbreak{} but every Hom tower has stable image the I\_\allowbreak{}0 summand,\allowbreak{} and the termwise limit is I\_\allowbreak{}0. The actual projection and inclusion identify that limit and its K-injectivity. No novelty is claimed for this criterion or obstruction.

The source\textquotesingle s exact-adjoint preservation lemma is also valid. For v left adjoint to u with v exact,\allowbreak{} the degreewise adjunction carries chain maps and homotopies to one another,\allowbreak{} so Hom\_\allowbreak{}K(B)\allowbreak{}(M,\allowbreak{}uI)\allowbreak{}\allowbreak{}=Hom\_\allowbreak{}K(A)\allowbreak{}(vM,\allowbreak{}I)\allowbreak{}. An acyclic M remains acyclic under v:\allowbreak{} exactness preserves the kernel-image short exact sequences defining its cohomology. Therefore the last Hom group is zero. The existing article correction MC-STK-ERR-0884 is complete.

\subsection{4. Repairing the initial cutoff in the replacement lemma}\label{proofs-9746-10132-repairing-the-initial-cutoff-in-the-replacement-lemma}

The dimension function d is permitted to have value infinity. The limit condition n\allowbreak{}+d(K\textasciicircum{}n)\allowbreak{}\allowbreak{}→-infinity as n\allowbreak{}→-infinity implies finiteness far to the left,\allowbreak{} but permits finitely many infinite values in any remaining bounded range. Replacing only σ\_\allowbreak{}\{\textbackslash{}geq0\}K does not necessarily remove them. Choose instead an integer a such that d(K\textasciicircum{}n)\allowbreak{}<infinity for every n\textless a. Resolve σ\_\allowbreak{}\{\textbackslash{}geq a\}K by the source\textquotesingle s bounded below construction to a termwise monomorphic quasi-isomorphism into M,\allowbreak{} with M\textasciicircum{}n\allowbreak{}=0 for n\textless a and d(M\textasciicircum{}n)\allowbreak{}\allowbreak{}=0 for n≥a.

The class of d-zero objects contains zero. To check this from the stated axioms,\allowbreak{} choose A with d(A)\allowbreak{}\allowbreak{}=0 by axiom (1)\allowbreak{} and apply axiom (3)\allowbreak{} to 0\allowbreak{}→A\allowbreak{}→A\allowbreak{}→0\allowbreak{}→0; it gives d(0)\allowbreak{}≤max(-1,\allowbreak{}0)\allowbreak{}\allowbreak{}=0. Thus the cited bounded resolution lemma is applicable.

Splice the initial map into

…\allowbreak{}→K\textasciicircum{}\{a-2\}\allowbreak{}→K\textasciicircum{}\{a-1\}\allowbreak{}→M\textasciicircum{}a\allowbreak{}→M\textasciicircum{}\{a\allowbreak{}+1\}\allowbreak{}→…,\allowbreak{}

using the composite K\textasciicircum{}\{a-1\}\allowbreak{}→K\textasciicircum{}a\allowbreak{}→M\textasciicircum{}a at the join. Its square with the next differential is zero by the chain-map equation. The map from K is the identity below a and the resolution map at and above a. The short exact sequences with subcomplex σ\_\allowbreak{}\{\textbackslash{}geq a\} and common quotient σ\_\allowbreak{}\{\textbackslash{}leq a-1\} show it is a quasi-isomorphism. Equivalently the quotient of the monomorphic map is the acyclic quotient of that bounded resolution. All its d-values are now finite,\allowbreak{} d is zero for n≥a,\allowbreak{} and the far-left limit condition is unchanged.

DERIVED-NEW-033 changes the source\textquotesingle s fixed cutoff 0 to this cutoff a. Here is a complete example of the defect in the original finite-maximum argument. In abelian groups set d(G)\allowbreak{}\allowbreak{}=0 for divisible G and infinity otherwise. Finite sums and quotients of divisible groups are divisible. Consequently axioms (2)\allowbreak{} and (3)\allowbreak{} hold:\allowbreak{} the only nonautomatic case of (3)\allowbreak{} has both its first and middle groups divisible,\allowbreak{} hence its quotient is. Every group embeds into a divisible group. One direct construction adjoins roots to all its elements by a pushout of the monomorphism \allowbreak{}⊕\_\allowbreak{}\{(g,\allowbreak{}m)\allowbreak{}\}Z\allowbreak{}→\allowbreak{}⊕\_\allowbreak{}\{(g,\allowbreak{}m)\allowbreak{}\}Z whose components are multiplication by m≥1,\allowbreak{} with the map from the first sum to G sending its generator to g. Pushout preserves that monomorphism,\allowbreak{} so G embeds in the resulting group,\allowbreak{} where every original g has an m-th root. Iterating countably and taking the union embeds G in a divisible group. Thus axiom (1)\allowbreak{} holds as well. Let K be Z in degree -1 and zero elsewhere. It meets the far-left condition,\allowbreak{} but the source\textquotesingle s resolution of σ\_\allowbreak{}\{\textbackslash{}geq0\}K leaves it unchanged,\allowbreak{} so the printed maximum ξ is infinite. The corrected cutoff removes this term before the finite-maximum argument. The lemma remains true.

\subsection{5. Elementary replacements and a quantitative stabilization bound}\label{proofs-9746-10132-elementary-replacements-and-a-quantitative-stabilization-bound}

After Section 4 assume all d-values are finite and zero in sufficiently high degrees. For n with d(K\textasciicircum{}n)\allowbreak{}>0 choose a monomorphism j:\allowbreak{}K\textasciicircum{}n\allowbreak{}→M with d(M)\allowbreak{}\allowbreak{}=0. Use the pushout

M'\allowbreak{}=coker((j,\allowbreak{}-d\_\allowbreak{}K\textasciicircum{}n)\allowbreak{}:\allowbreak{}K\textasciicircum{}n\allowbreak{}→M\allowbreak{}⊕K\textasciicircum{}\{n\allowbreak{}+1\})\allowbreak{}.

The new differential into M is j d\_\allowbreak{}K\textasciicircum{}\{n-1\}; the next differential is m\allowbreak{}↦[m,\allowbreak{}0]\allowbreak{}; the differential out of M\textquotesingle{} is [m,\allowbreak{}x]\allowbreak{}\allowbreak{}↦d\_\allowbreak{}K\textasciicircum{}\{n\allowbreak{}+1\}x. The vertical map K\textasciicircum{}\{n\allowbreak{}+1\}\allowbreak{}→M' is x\allowbreak{}↦[0,\allowbreak{}x]\allowbreak{}. The identity {[}j y,\allowbreak{}0]\allowbreak{}\allowbreak{}=[0,\allowbreak{}d\_\allowbreak{}K\textasciicircum{}n y]\allowbreak{} proves the chain-map equation,\allowbreak{} and the relation is killed by the outgoing differential because d\_\allowbreak{}K²\allowbreak{}=0. All other components are unchanged.

Both changed vertical components are monomorphisms. The quotient complex is M/\allowbreak{}j(K\textasciicircum{}n)\allowbreak{} in degrees n and n\allowbreak{}+1,\allowbreak{} with identity differential:\allowbreak{} the second quotient M'/\allowbreak{}K\textasciicircum{}\{n\allowbreak{}+1\} is canonically M/\allowbreak{}j(K\textasciicircum{}n)\allowbreak{}. It is contractible. Hence this elementary replacement is a quasi-isomorphism. The defining short exact sequence for M\textquotesingle{} and the d axioms give exactly

d((K')\allowbreak{}\textasciicircum{}n)\allowbreak{}\allowbreak{}=0,\allowbreak{} d((K')\allowbreak{}\textasciicircum{}\{n\allowbreak{}+1\})\allowbreak{}≤max(d(K\textasciicircum{}n)\allowbreak{}-1,\allowbreak{}d(K\textasciicircum{}\{n\allowbreak{}+1\})\allowbreak{})\allowbreak{}.

If there are no positive d-values,\allowbreak{} the construction stops. Otherwise ξ\allowbreak{}=max\{n\allowbreak{}+d(K\textasciicircum{}n)\allowbreak{}:\allowbreak{}d(K\textasciicircum{}n)\allowbreak{}>0\} is a finite integer. Its maximum is attained at finitely many indices because weights tend to -infinity to the left and d is zero to the right. An elementary step cannot increase ξ. At a maximal-weight index n,\allowbreak{} any new positive weight at n\allowbreak{}+1 is at most ξ,\allowbreak{} and the old positive value at n disappears. With ξ fixed,\allowbreak{} the positive integer Σ\_\allowbreak{}\{n\allowbreak{}+d(K\textasciicircum{}n)\allowbreak{}\allowbreak{}=ξ\}(ξ-n)\allowbreak{} strictly decreases on every such step:\allowbreak{} a new maximal index can only replace n by n\allowbreak{}+1,\allowbreak{} or coincide with an already maximal index. Thus finitely many steps decrease ξ. Repeating either terminates or makes ξ tend to -infinity.

DERIVED-UNDER-022 gives a bound that also directly proves stabilization. Let d\_\allowbreak{}j be the finite dimensions immediately after the initial splice. For each integer m put

N\_\allowbreak{}m\allowbreak{}=Σ\_\allowbreak{}\{j:\allowbreak{}d\_\allowbreak{}j>0,\allowbreak{} j\allowbreak{}+d\_\allowbreak{}j≥m\}d\_\allowbreak{}j. (5.1)\allowbreak{}

This is a finite sum. During the algorithm let the same expression,\allowbreak{} formed from the current dimensions,\allowbreak{} be its current mass above level m. A replacement with n\allowbreak{}+d(K\textasciicircum{}n)\allowbreak{}≥m decreases that mass by at least one. Indeed it removes d(K\textasciicircum{}n)\allowbreak{}; if its recipient was already counted,\allowbreak{} the recipient increases by at most d(K\textasciicircum{}n)\allowbreak{}-1,\allowbreak{} and if it was not counted,\allowbreak{} its new counted value is at most d(K\textasciicircum{}n)\allowbreak{}-1. A replacement of lower weight cannot increase the mass:\allowbreak{} it cannot create a weight ≥m,\allowbreak{} and cannot increase a recipient already of weight ≥m.

A step changing a degree at least m has n\allowbreak{}+1≥m and d(K\textasciicircum{}n)\allowbreak{}≥1,\allowbreak{} hence n\allowbreak{}+d(K\textasciicircum{}n)\allowbreak{}≥m. Therefore at most N\_\allowbreak{}m elementary replacements affect the entire tail in degrees ≥m. This proves the source\textquotesingle s required local stabilization with an explicit bound.

Define L degree by degree by the eventual unchanged terms and differentials. Every component of K\allowbreak{}→L is the finite composite of the actual replacement maps affecting that component. Compatibility holds because the chain equations hold at every finite stage and eventually stabilize. Cohomology at degree r depends only on degrees r-1,\allowbreak{}r,\allowbreak{}r\allowbreak{}+1,\allowbreak{} which stabilize after finitely many steps; the finite composite to that stage is a quasi-isomorphism,\allowbreak{} so K\allowbreak{}→L induces an isomorphism on H\^{}r. Every surviving d-value is zero,\allowbreak{} since an eventual positive value would give a fixed weight while ξ tends to -infinity. No infinite colimit or exactness of an infinite limit is used. P02-E0076 also identifies the missing article at original 9968; its copyedit is separate from the cutoff repair.

\subsection{6. Finite cohomological dimension and the cone bound}\label{proofs-9746-10132-finite-cohomological-dimension-and-the-cone-bound}

Let F be left exact,\allowbreak{} with enough right acyclic objects,\allowbreak{} and suppose R\textasciicircum{}nF\allowbreak{}=0. If n\allowbreak{}=0,\allowbreak{} left exactness identifies R\^{}0F with F,\allowbreak{} so F is the zero functor. Its derived functor is zero on every complex,\allowbreak{} and all conclusions follow. In the remainder take n≥1.

The bounded below derived functor exists by enough acyclics. For a monomorphism A\allowbreak{}→A' with A\textquotesingle{} right acyclic and quotient C,\allowbreak{} the long exact sequence gives R\textasciicircum{}\{r\allowbreak{}+1\}F(A)\allowbreak{}≅R\textasciicircum{}rF(C)\allowbreak{} for r≥1. Starting with r\allowbreak{}=n and iterating shows R\textasciicircum{}rF\allowbreak{}=0 for all r≥n. Define

d(A)\allowbreak{}\allowbreak{}=max(\{0\}∪\{r:\allowbreak{}R\textasciicircum{}rF(A)\allowbreak{}≠0\})\allowbreak{}.

Then 0≤d(A)\allowbreak{}≤n-1. Additivity gives the direct-sum axiom for d. For 0\allowbreak{}→A\allowbreak{}→B\allowbreak{}→C\allowbreak{}→0 and r>max(d(A)\allowbreak{}-1,\allowbreak{}d(B)\allowbreak{})\allowbreak{},\allowbreak{} the two adjacent terms R\textasciicircum{}rF(B)\allowbreak{} and R\textasciicircum{}\{r\allowbreak{}+1\}F(A)\allowbreak{} in the long exact sequence vanish,\allowbreak{} so R\textasciicircum{}rF(C)\allowbreak{}\allowbreak{}=0. This proves the third d axiom. The enough-acyclics assumption supplies the first. The corrected replacement lemma therefore maps every complex to a complex of right acyclic objects. The admitted parentheses in MC-STK-ERR-0879 give this actual maximum.

Let L\allowbreak{}→M be a quasi-isomorphism between complexes of right acyclic objects. Fix a cohomological degree i and put c\allowbreak{}=i-n-1. Let Q\allowbreak{}=C(σ\_\allowbreak{}\{\textbackslash{}geq c\}L\allowbreak{}→σ\_\allowbreak{}\{\textbackslash{}geq c\}M)\allowbreak{}. The source uses stupid truncations. They agree with L,\allowbreak{}M in degrees at least c,\allowbreak{} so their cohomology agrees in degrees greater than c and vanishes below c. The long exact cone sequence therefore makes H\textasciicircum{}q(Q)\allowbreak{} zero except possibly for q\allowbreak{}=c-1,\allowbreak{}c.

Apply the bounded below hypercohomology spectral sequence to Q. Its second page has terms R\textasciicircum{}pF(H\textasciicircum{}q(Q)\allowbreak{})\allowbreak{},\allowbreak{} with 0≤p≤n-1 and q\allowbreak{}∈\{c-1,\allowbreak{}c\}. Every possible total degree p\allowbreak{}+q lies in

[c-1,\allowbreak{}c\allowbreak{}+n-1]\allowbreak{}\allowbreak{}=[i-n-2,\allowbreak{}i-2]\allowbreak{}. (6.1)\allowbreak{}

The filtration is bounded in each degree,\allowbreak{} so the same support holds for RF(Q)\allowbreak{}. In particular both H\textasciicircum{}\{i-1\}(RF(Q)\allowbreak{})\allowbreak{} and H\textasciicircum{}i(RF(Q)\allowbreak{})\allowbreak{} vanish. The triangle now proves an isomorphism on H\^{}i from RF(σ\_\allowbreak{}\{\textbackslash{}geq c\}L)\allowbreak{} to RF(σ\_\allowbreak{}\{\textbackslash{}geq c\}M)\allowbreak{}. Both are computed termwise by F because their terms are right acyclic and they are bounded below. Since i>c,\allowbreak{} their H\^{}i is that of F(L)\allowbreak{} and F(M)\allowbreak{}. Thus F(L)\allowbreak{}\allowbreak{}→F(M)\allowbreak{} is a quasi-isomorphism. The computation criterion constructs RF on all of D(A)\allowbreak{},\allowbreak{} and every complex of right acyclic objects computes it.

DERIVED-NEW-034 tightens the printed upper endpoint i-1 to i-2 in original 10048. The printed interval is a true but insufficient estimate:\allowbreak{} it does not exclude H\textasciicircum{}\{i-1\}(RF(Q)\allowbreak{})\allowbreak{},\allowbreak{} whose vanishing is needed for the claimed isomorphism. Formula (6.1)\allowbreak{} proves the missing vanishing rather than replacing it by an assumption.

\subsection{\texorpdfstring{7. Canonical truncation maps,\allowbreak{} their endpoints,\allowbreak{} and finite windows}{7. Canonical truncation  their  and finite windows}}\label{proofs-9746-10132-canonical-truncation-maps-their-endpoints-and-finite-windows}

For E\allowbreak{}∈D(A)\allowbreak{},\allowbreak{} the functor just constructed agrees with bounded below derivation on D\textasciicircum{}\allowbreak{}+(A)\allowbreak{}. In particular it sends D\^{}\{≥a\} into D\textasciicircum{}\{≥a\}:\allowbreak{} represent the input by a complex zero below a and use the bounded below acyclic resolution with the same lower bound. Apply RF to τ\_\allowbreak{}\{\textbackslash{}leq a\}E\allowbreak{}→E\allowbreak{}→τ\_\allowbreak{}\{\textbackslash{}geq a\allowbreak{}+1\}E. The long exact sequence proves H\textasciicircum{}i(RF(τ\_\allowbreak{}\{\textbackslash{}leq a\}E)\allowbreak{})\allowbreak{}\allowbreak{}→H\textasciicircum{}i(RF(E)\allowbreak{})\allowbreak{} is an isomorphism for i≤a. It is also a monomorphism for i\allowbreak{}=a\allowbreak{}+1. This proves the statement with the existing parentheses correction MC-STK-ERR-0878.

For the second source comparison write E as a complex L of right acyclic objects. For any integer c,\allowbreak{} RF(σ\_\allowbreak{}\{\textbackslash{}geq c\}L)\allowbreak{} is computed by σ\_\allowbreak{}\{\textbackslash{}geq c\}F(L)\allowbreak{}. The quantifier on c makes the existing unintroduced cutoff explicit; P02-E0079 is repaired this way,\allowbreak{} preserving the true general statement instead of replacing c throughout by one special value.

Set q\allowbreak{}=b-n. There are the two source sequences

H\textasciicircum{}q(L)\allowbreak{}[-q]\allowbreak{}\allowbreak{}→τ\_\allowbreak{}\{\textbackslash{}geq q\}L\allowbreak{}→τ\_\allowbreak{}\{\textbackslash{}geq q\allowbreak{}+1\}L,\allowbreak{} 0\allowbreak{}→im(d\_\allowbreak{}L\textasciicircum{}\{q-1\})\allowbreak{}[-q]\allowbreak{}\allowbreak{}→σ\_\allowbreak{}\{\textbackslash{}geq q\}L\allowbreak{}→τ\_\allowbreak{}\{\textbackslash{}geq q\}L\allowbreak{}→0.

For either left term T[-q]\allowbreak{},\allowbreak{} RF(T[-q]\allowbreak{})\allowbreak{} has cohomology only in degrees [q,\allowbreak{}q\allowbreak{}+n-1]\allowbreak{}\allowbreak{}=[b-n,\allowbreak{}b-1]\allowbreak{}. Their long exact sequences show that both maps RF(σ\_\allowbreak{}\{\textbackslash{}geq q\}L)\allowbreak{}\allowbreak{}→RF(τ\_\allowbreak{}\{\textbackslash{}geq q\}L)\allowbreak{} and RF(τ\_\allowbreak{}\{\textbackslash{}geq q\}L)\allowbreak{}\allowbreak{}→RF(τ\_\allowbreak{}\{\textbackslash{}geq q\allowbreak{}+1\}L)\allowbreak{} induce isomorphisms in degrees ≥b. The map RF(σ\_\allowbreak{}\{\textbackslash{}geq q\}L)\allowbreak{}\allowbreak{}→RF(L)\allowbreak{} is the actual inclusion σ\_\allowbreak{}\{\textbackslash{}geq q\}F(L)\allowbreak{}\allowbreak{}→F(L)\allowbreak{},\allowbreak{} also an isomorphism in degrees ≥b since n≥1. All these maps commute with L\allowbreak{}→τ\_\allowbreak{}\{\textbackslash{}geq q\allowbreak{}+1\}L. Thus that canonical map induces the asserted isomorphism

H\textasciicircum{}i(RF(E)\allowbreak{})\allowbreak{}\allowbreak{}→H\textasciicircum{}i(RF(τ\_\allowbreak{}\{\textbackslash{}geq b-n\allowbreak{}+1\}E)\allowbreak{})\allowbreak{},\allowbreak{} i≥b.

If the cohomology support of E lies in [a,\allowbreak{}b]\allowbreak{},\allowbreak{} the first comparison gives no RF cohomology below a and the second gives none at or above b\allowbreak{}+n. This proves support in [a,\allowbreak{}b\allowbreak{}+n-1]\allowbreak{},\allowbreak{} including unbounded endpoints.

DERIVED-UNDER-023 records the finite-window consequence. For every i,\allowbreak{} the canonical zigzag

E\allowbreak{}→τ\_\allowbreak{}\{\textbackslash{}geq i-n\allowbreak{}+1\}E←τ\_\allowbreak{}\{\textbackslash{}leq i\}τ\_\allowbreak{}\{\textbackslash{}geq i-n\allowbreak{}+1\}E

becomes isomorphisms on H\^{}i after applying RF. Hence H\textasciicircum{}i(RF(E)\allowbreak{})\allowbreak{} is canonically determined by the actual truncation window [i-n\allowbreak{}+1,\allowbreak{}i]\allowbreak{},\allowbreak{} retaining all its extension data,\allowbreak{} not merely the list of cohomology objects. The upper-truncation comparison is monomorphic at its next degree as already shown. The lower-truncation comparison E\allowbreak{}→τ\_\allowbreak{}\{\textbackslash{}geq c\}E is epimorphic on H\textasciicircum{}\{c\allowbreak{}+n-2\}(RF)\allowbreak{}:\allowbreak{} its triangle has first term RF(τ\_\allowbreak{}\{\textbackslash{}leq c-1\}E)\allowbreak{},\allowbreak{} supported at most in degree c\allowbreak{}+n-2,\allowbreak{} so the next cohomology of that term is zero. It is an isomorphism in degrees ≥c\allowbreak{}+n-1.

The left-derived assertions are obtained with the exact contravariant complex transformation D(K)\allowbreak{}\textasciicircum{}r\allowbreak{}=(K\textasciicircum{}\{-r\})\allowbreak{}\textasciicircum{}op and d\_\allowbreak{}D\textasciicircum{}r\allowbreak{}=(d\_\allowbreak{}K\textasciicircum{}\{-r-1\})\allowbreak{}\textasciicircum{}op used in PART14. It interchanges the canonical upper and lower truncations,\allowbreak{} sends H\^{}r to H\textasciicircum{}\{-r\},\allowbreak{} reverses arrows,\allowbreak{} and changes the interval [a,\allowbreak{}b]\allowbreak{} to [-b,\allowbreak{}-a]\allowbreak{}. The positive left-derived index n in this dual assertion means the vanishing in degree -n of LF(A[0]\allowbreak{})\allowbreak{},\allowbreak{} equivalently R\textasciicircum{}n(F\textasciicircum{}op)\allowbreak{}\allowbreak{}=0. The right interval [-b,\allowbreak{}-a\allowbreak{}+n-1]\allowbreak{} therefore returns exactly [a-n\allowbreak{}+1,\allowbreak{}b]\allowbreak{}. The canonical comparisons become

H\textasciicircum{}i(LF(τ\_\allowbreak{}\{\textbackslash{}leq a\allowbreak{}+n-1\}E)\allowbreak{})\allowbreak{}\allowbreak{}→H\textasciicircum{}i(LF(E)\allowbreak{})\allowbreak{},\allowbreak{} i≤a,\allowbreak{} H\textasciicircum{}i(LF(E)\allowbreak{})\allowbreak{}\allowbreak{}→H\textasciicircum{}i(LF(τ\_\allowbreak{}\{\textbackslash{}geq b\}E)\allowbreak{})\allowbreak{},\allowbreak{} i≥b.

These are the source\textquotesingle s maps with MC-STK-ERR-0880\textquotesingle s restored parenthesis. The dual finite window is [i,\allowbreak{}i\allowbreak{}+n-1]\allowbreak{}. At the first additional degree the upper-cut comparison is monomorphic; the lower-cut comparison is epimorphic at degree b-1. For n\allowbreak{}=0 all these functors are zero,\allowbreak{} as in the right-derived case. No replacement by a differently indexed derived functor has been made.

\subsection{\texorpdfstring{8. Receivers,\allowbreak{} dispositions,\allowbreak{} and the next section}{8.   and the next section}}\label{proofs-9746-10132-receivers-dispositions-and-the-next-section}

Current cohomology.tex 10300–\allowbreak{}10317 distinguishes K-injectivity of the termwise limit from its being a resolution. Sections 2–\allowbreak{}3 prove its K-injectivity invocation and preserve that distinction. Current injectives.tex 2078–\allowbreak{}2082 uses the existence of RF from enough K-injectives,\allowbreak{} exactly as Section 1 proves. Current perfect.tex 1643–\allowbreak{}1649 applies the exact-left-adjoint criterion,\allowbreak{} verified in Section 3. These are bounded checks of the stated invocations.

Current more-algebra.tex 23649–\allowbreak{}23678 applies the finite cohomological dimension theorem to inverse limits of abelian-group towers. The correction (6.1)\allowbreak{} preserves that theorem and its acyclic-complex computation. The exact obstruction groups in Section 3 were constructed directly and do not depend circularly on this receiving Rlim theorem. Current cohomology.tex 14218–\allowbreak{}14228 also uses that any complex of acyclics computes RF. That invocation is preserved. The preceding displayed identification of an ordinary Hom with RΓ at current 14213 is routed as a type-check observation for its own chapter review,\allowbreak{} not adjudicated here.

Ten more physical reports are complete:\allowbreak{} groups DERIVED-RECON-073 through 078. Eight are covered by existing units 0884,\allowbreak{} 0878,\allowbreak{} 0879,\allowbreak{} 0880; two need copyedit/\allowbreak{}notation operations (P02-E0076 and P02-E0079)\allowbreak{}. Totals are 135/\allowbreak{}176 reports,\allowbreak{} 78 groups,\allowbreak{} 124 already corrected,\allowbreak{} ten missing copyedit/\allowbreak{}notation reports and one missing mathematical-index report. There are 71 prior Derived units and 88 operations reconciled. NEW-033 and NEW-034 add two unreported findings; together with the two received copyedits there are four new operations,\allowbreak{} giving 65 cumulative operations. UNDER-021,\allowbreak{} 022 and 023 bring the editorial consequence total to 23.

Review is complete through original 10132; reading has reached 10310. Next is Derived colimits at 10139; unread source starts at 10311. The read-ahead functoriality and subsequence formulas require their own full diagram and report comparison before further operations are proposed. No new sessions,\allowbreak{} source composition,\allowbreak{} build,\allowbreak{} Lean run,\allowbreak{} cleanup,\allowbreak{} publication,\allowbreak{} or producer mathematical acceptance occurred.
\clearpage
\section{\texorpdfstring{Derived colimits:\allowbreak{} telescope maps and comparison defects}{Derived  telescope maps and comparison defects}}\label{proofs-10139-10427-derived-colimits-telescope-maps-and-comparison-defects}

Private source-order review of original derived.tex 10139–\allowbreak{}10427. Authority:\allowbreak{} Stacks a04446e5\allowbreak{}7ec1fbc2\allowbreak{}52a871af\allowbreak{}cec7752f\allowbreak{}b2807b14. Receiving public edition:\allowbreak{} 95a51593\allowbreak{}aceb247a\allowbreak{}c7a18311\allowbreak{}1678c3a0\allowbreak{}f8b8f4b5. The operation and source-reading records bind the actual files. The strengthened comparison statement below is editorial material,\allowbreak{} not a replacement for the source\textquotesingle s compactness hypothesis. No novelty is claimed.

Human source:\allowbreak{} \href{https://github.com/stacks/stacks-project/blob/a04446e57ec1fbc252a871afcec7752fb2807b14/derived.tex\#L10139}{The Stacks Project,\allowbreak{} original Derived Categories TeX}. The Hom exact sequence also occurs in \href{https://github.com/stacks/stacks-project/blob/a04446e57ec1fbc252a871afcec7752fb2807b14/more-algebra.tex}{More on Algebra,\allowbreak{} original TeX},\allowbreak{} Lemma named lemma-map-from-hocolim. The source labels,\allowbreak{} rather than current line offsets,\allowbreak{} identify that receiving result.

\subsection{1. The defining triangle and actual comparison maps}\label{proofs-10139-10427-the-defining-triangle-and-actual-comparison-maps}

For a sequence K\_\allowbreak{}1\allowbreak{}→K\_\allowbreak{}2\allowbreak{}→… with maps f\_\allowbreak{}n put S\allowbreak{}=\allowbreak{}⊕\_\allowbreak{}n K\_\allowbreak{}n and let ι\_\allowbreak{}n:\allowbreak{}K\_\allowbreak{}n\allowbreak{}→S be its coprojections. Define u:\allowbreak{}S\allowbreak{}→S by the exact formulas

u ι\_\allowbreak{}n\allowbreak{}=ι\_\allowbreak{}n-ι\_\allowbreak{}\{n\allowbreak{}+1\}f\_\allowbreak{}n.

Each column has two components,\allowbreak{} so the universal property of the coproduct defines u even if the individual K\_\allowbreak{}n are not compact. A homotopy colimit is a choice of distinguished triangle

S --u-\allowbreak{}-> S --v-\allowbreak{}-> C --w-\allowbreak{}-> S[1]\allowbreak{}. (1.1)\allowbreak{}

Its structural maps are i\_\allowbreak{}n\allowbreak{}=vι\_\allowbreak{}n,\allowbreak{} and v u\allowbreak{}=0 gives i\_\allowbreak{}\{n\allowbreak{}+1\}f\_\allowbreak{}n\allowbreak{}=i\_\allowbreak{}n. TR1 gives existence as soon as S exists. Two choices with this same u admit a map of triangles whose first two maps are identities,\allowbreak{} by TR3. Applying representable Hom gives maps of long exact sequences; the five-lemma argument shows that the third map induces bijections on every representable Hom,\allowbreak{} so is an isomorphism. This is uniqueness up to a possibly nonunique isomorphism compatible with both v and w,\allowbreak{} as stated in the source.

For a morphism of sequences a\_\allowbreak{}n:\allowbreak{}K\_\allowbreak{}n\allowbreak{}→L\_\allowbreak{}n satisfying a\_\allowbreak{}\{n\allowbreak{}+1\}f\_\allowbreak{}n\allowbreak{}=g\_\allowbreak{}n a\_\allowbreak{}n,\allowbreak{} let A\allowbreak{}=\allowbreak{}⊕a\_\allowbreak{}n:\allowbreak{}S\_\allowbreak{}K\allowbreak{}→S\_\allowbreak{}L. The two maps in the square are both A,\allowbreak{} and the column equation gives A u\_\allowbreak{}K\allowbreak{}=u\_\allowbreak{}L A. TR3 supplies a map a:\allowbreak{}C\_\allowbreak{}K\allowbreak{}→C\_\allowbreak{}L with

a v\_\allowbreak{}K\allowbreak{}=v\_\allowbreak{}L A,\allowbreak{} w\_\allowbreak{}L a\allowbreak{}=A[1]\allowbreak{}w\_\allowbreak{}K.

After precomposition with ι\_\allowbreak{}n this says a i\_\allowbreak{}n\allowbreak{}=j\_\allowbreak{}n a\_\allowbreak{}n. The source\textquotesingle s printed a i\_\allowbreak{}n\allowbreak{}=j\_\allowbreak{}n has different domains on its two sides. Existing unit MC-STK-ERR-0881 supplies precisely the missing a\_\allowbreak{}n; no new operation is needed for that reported defect.

For completeness,\allowbreak{} differences of two such completions ε:\allowbreak{}C\_\allowbreak{}K\allowbreak{}→C\_\allowbreak{}L satisfy ε v\_\allowbreak{}K\allowbreak{}=0 and w\_\allowbreak{}L ε\allowbreak{}=0. Exactness of Hom applied to (1.1)\allowbreak{} gives ε\allowbreak{}=t w\_\allowbreak{}K for some t:\allowbreak{}S\_\allowbreak{}K[1]\allowbreak{}\allowbreak{}→C\_\allowbreak{}L. If η:\allowbreak{}C\_\allowbreak{}L\allowbreak{}→C\_\allowbreak{}M is another such difference,\allowbreak{} write η\allowbreak{}=s w\_\allowbreak{}L. Then ηε\allowbreak{}=s w\_\allowbreak{}L ε\allowbreak{}=0. Thus these differences form the square-zero ideal already proved as DERIVED-UNDER-001. This is a receiving application of that result,\allowbreak{} not a new underclaim. It does not assert a canonical choice of a.

\subsection{2. Maps out of a telescope and their ambiguity}\label{proofs-10139-10427-maps-out-of-a-telescope-and-their-ambiguity}

Apply Hom(-,\allowbreak{}L)\allowbreak{} to (1.1)\allowbreak{}. The relevant exact segment is

Hom(S[1]\allowbreak{},\allowbreak{}L)\allowbreak{} --u[1]\allowbreak{}\textasciicircum{}*-\allowbreak{}-> Hom(S[1]\allowbreak{},\allowbreak{}L)\allowbreak{} --w\textasciicircum{}*-\allowbreak{}-> Hom(C,\allowbreak{}L)\allowbreak{} --v\textasciicircum{}*-\allowbreak{}-> Hom(S,\allowbreak{}L)\allowbreak{} --u\textasciicircum{}*-\allowbreak{}-> Hom(S,\allowbreak{}L)\allowbreak{}.

The coproduct property gives Hom(S,\allowbreak{}L)\allowbreak{}\allowbreak{}=∏\_\allowbreak{}n Hom(K\_\allowbreak{}n,\allowbreak{}L)\allowbreak{}. Under this exact identification,\allowbreak{}

u\textasciicircum{}*((x\_\allowbreak{}n)\allowbreak{})\allowbreak{}\allowbreak{}=(x\_\allowbreak{}n-x\_\allowbreak{}\{n\allowbreak{}+1\}f\_\allowbreak{}n)\allowbreak{}.

The same formula on shifted objects identifies u[1]\allowbreak{}\textasciicircum{}* with the difference map for the tower Hom(K\_\allowbreak{}n,\allowbreak{}L[-1]\allowbreak{})\allowbreak{},\allowbreak{} using the exact shift adjunction Hom(K\_\allowbreak{}n[1]\allowbreak{},\allowbreak{}L)\allowbreak{}\allowbreak{}=Hom(K\_\allowbreak{}n,\allowbreak{}L[-1]\allowbreak{})\allowbreak{}. If Δ denotes that product difference map,\allowbreak{} the exact segment becomes

0\allowbreak{}→coker Δ\_\allowbreak{}\{Hom(K\_\allowbreak{}n,\allowbreak{}L[-1]\allowbreak{})\allowbreak{}\}\allowbreak{}→Hom(C,\allowbreak{}L)\allowbreak{}\allowbreak{}→lim Hom(K\_\allowbreak{}n,\allowbreak{}L)\allowbreak{}\allowbreak{}→0. (2.1)\allowbreak{}

The left injection sends a family z\_\allowbreak{}n:\allowbreak{}K\_\allowbreak{}n[1]\allowbreak{}\allowbreak{}→L to z w,\allowbreak{} where z:\allowbreak{}S[1]\allowbreak{}\allowbreak{}→L has those components. A change by an element of im Δ does not change z w because u[1]\allowbreak{}w\allowbreak{}=0. The right map sends b:\allowbreak{}C\allowbreak{}→L to (b i\_\allowbreak{}n)\allowbreak{}. Thus every compatible family lifts,\allowbreak{} and its lifts form a torsor under the displayed cokernel; no functorial choice of a lift follows. The source\textquotesingle s R\^{}1 lim notation denotes this cokernel for abelian-group towers. This proves the exact maps in the displayed source sequence and the receiving More on Algebra lemma without a new count.

\subsection{\texorpdfstring{3. Passing to a subsequence:\allowbreak{} all four maps}{3. Passing to a  all four maps}}\label{proofs-10139-10427-passing-to-a-subsequence-all-four-maps}

Let n\_\allowbreak{}1<n\_\allowbreak{}2<… be the given subsequence. Write

t\_\allowbreak{}\{q,\allowbreak{}p\}\allowbreak{}=f\_\allowbreak{}\{q-1\}⋯f\_\allowbreak{}p:\allowbreak{}K\_\allowbreak{}p\allowbreak{}→K\_\allowbreak{}q for q>p,\allowbreak{} t\_\allowbreak{}\{p,\allowbreak{}p\}\allowbreak{}=id,\allowbreak{} T\allowbreak{}=\allowbreak{}⊕\_\allowbreak{}i K\_\allowbreak{}\{n\_\allowbreak{}i\},\allowbreak{} g\_\allowbreak{}i\allowbreak{}=t\_\allowbreak{}\{n\_\allowbreak{}\{i\allowbreak{}+1\},\allowbreak{}n\_\allowbreak{}i\},\allowbreak{} u\_\allowbreak{}T κ\_\allowbreak{}i\allowbreak{}=κ\_\allowbreak{}i-κ\_\allowbreak{}\{i\allowbreak{}+1\}g\_\allowbreak{}i,\allowbreak{}

where κ\_\allowbreak{}i:\allowbreak{}K\_\allowbreak{}\{n\_\allowbreak{}i\}\allowbreak{}→T. Retain S,\allowbreak{}u from Section 1. The two telescope triangles and forward comparison have this form:\allowbreak{}

\[
\begin{array}{ccccccc}
T&\xrightarrow{u_T}&T&\xrightarrow{v_T}&C_T&\xrightarrow{w_T}&T[1]\\
\scriptstyle b\downarrow&&\scriptstyle a\downarrow&&\scriptstyle φ\downarrow&&\scriptstyle b[1]\downarrow\\
S&\xrightarrow{u}&S&\xrightarrow{v}&C&\xrightarrow{w}&S[1].
\end{array}
\]

The reverse comparison has first map d,\allowbreak{} second map c and third map ψ. Define all components,\allowbreak{} including the initial segment before n\_\allowbreak{}1,\allowbreak{} by

a κ\_\allowbreak{}i\allowbreak{}=ι\_\allowbreak{}\{n\_\allowbreak{}i\},\allowbreak{} b κ\_\allowbreak{}i\allowbreak{}=Σ\_\allowbreak{}\{n\_\allowbreak{}i≤q<n\_\allowbreak{}\{i\allowbreak{}+1\}\}ι\_\allowbreak{}q t\_\allowbreak{}\{q,\allowbreak{}n\_\allowbreak{}i\},\allowbreak{} c ι\_\allowbreak{}j\allowbreak{}=κ\_\allowbreak{}i t\_\allowbreak{}\{n\_\allowbreak{}i,\allowbreak{}j\},\allowbreak{} where i is the least index with j≤n\_\allowbreak{}i,\allowbreak{} d ι\_\allowbreak{}j\allowbreak{}=κ\_\allowbreak{}i if j\allowbreak{}=n\_\allowbreak{}i,\allowbreak{} and zero otherwise. (3.1)\allowbreak{}

Each expression defines a map from a coproduct by finitely many components on each source summand. The least index in c exists since the strictly increasing natural-number subsequence is unbounded. At j\allowbreak{}=n\_\allowbreak{}i the transition composite in c is the identity.

Telescoping each finite sum proves

u b κ\_\allowbreak{}i\allowbreak{}=ι\_\allowbreak{}\{n\_\allowbreak{}i\}-ι\_\allowbreak{}\{n\_\allowbreak{}\{i\allowbreak{}+1\}\}t\_\allowbreak{}\{n\_\allowbreak{}\{i\allowbreak{}+1\},\allowbreak{}n\_\allowbreak{}i\} \allowbreak{}=a u\_\allowbreak{}T κ\_\allowbreak{}i.

For the reverse square,\allowbreak{} if j is not a selected index,\allowbreak{} j and j\allowbreak{}+1 have the same least selected upper endpoint and c u ι\_\allowbreak{}j\allowbreak{}=0\allowbreak{}=u\_\allowbreak{}T dι\_\allowbreak{}j. If j\allowbreak{}=n\_\allowbreak{}i,\allowbreak{} then c u ι\_\allowbreak{}j\allowbreak{}=κ\_\allowbreak{}i-κ\_\allowbreak{}\{i\allowbreak{}+1\}g\_\allowbreak{}i\allowbreak{}=u\_\allowbreak{}T dι\_\allowbreak{}j. Hence u b\allowbreak{}=a u\_\allowbreak{}T and c u\allowbreak{}=u\_\allowbreak{}T d,\allowbreak{} exactly the two required squares. TR3 gives φ and ψ with all triangle compatibilities.

The identities c a\allowbreak{}=id\_\allowbreak{}T and d b\allowbreak{}=id\_\allowbreak{}T hold on every κ\_\allowbreak{}i:\allowbreak{} the finite interval [n\_\allowbreak{}i,\allowbreak{}n\_\allowbreak{}\{i\allowbreak{}+1\})\allowbreak{} contains exactly one selected index,\allowbreak{} namely n\_\allowbreak{}i. These are the first two maps of the composite of triangles from the T triangle to itself. Consequently ψφ is an isomorphism. The printed φψ at original 10278 reverses these domains.

For the other composite,\allowbreak{} define h:\allowbreak{}S\allowbreak{}→S on each summand as follows. Let i be the least index with j≤n\_\allowbreak{}i. Set

h ι\_\allowbreak{}j\allowbreak{}=Σ\_\allowbreak{}\{j≤q<n\_\allowbreak{}i\}ι\_\allowbreak{}q t\_\allowbreak{}\{q,\allowbreak{}j\}. (3.2)\allowbreak{}

The sum is empty,\allowbreak{} hence zero,\allowbreak{} when j\allowbreak{}=n\_\allowbreak{}i. This includes every positive j before n\_\allowbreak{}1. The last summand is K\_\allowbreak{}\{n\_\allowbreak{}i-1\},\allowbreak{} with transition f\_\allowbreak{}\{n\_\allowbreak{}i-2\}⋯f\_\allowbreak{}j; there is no K\_\allowbreak{}\{n\_\allowbreak{}i\} summand.

The first exact identity is obtained by telescoping:\allowbreak{}

u h ι\_\allowbreak{}j\allowbreak{}=ι\_\allowbreak{}j-ι\_\allowbreak{}\{n\_\allowbreak{}i\}t\_\allowbreak{}\{n\_\allowbreak{}i,\allowbreak{}j\}\allowbreak{}=(id\_\allowbreak{}S-a c)\allowbreak{}ι\_\allowbreak{}j.

For the second,\allowbreak{} when j<n\_\allowbreak{}i the difference of (3.2)\allowbreak{} at j and j\allowbreak{}+1 leaves exactly ι\_\allowbreak{}j,\allowbreak{} so h u ι\_\allowbreak{}j\allowbreak{}=ι\_\allowbreak{}j\allowbreak{}=(id\_\allowbreak{}S-b d)\allowbreak{}ι\_\allowbreak{}j. When j\allowbreak{}=n\_\allowbreak{}i,\allowbreak{} the first sum is zero and the second has upper endpoint n\_\allowbreak{}\{i\allowbreak{}+1\}; therefore

h u ι\_\allowbreak{}\{n\_\allowbreak{}i\} \allowbreak{}=-Σ\_\allowbreak{}\{n\_\allowbreak{}i<q<n\_\allowbreak{}\{i\allowbreak{}+1\}\}ι\_\allowbreak{}q t\_\allowbreak{}\{q,\allowbreak{}n\_\allowbreak{}i\} \allowbreak{}=ι\_\allowbreak{}\{n\_\allowbreak{}i\}-bκ\_\allowbreak{}i \allowbreak{}=(id\_\allowbreak{}S-b d)\allowbreak{}ι\_\allowbreak{}\{n\_\allowbreak{}i\}.

Thus the two source homotopy equations hold exactly:\allowbreak{}

u h\allowbreak{}=id\_\allowbreak{}S-a c,\allowbreak{} h u\allowbreak{}=id\_\allowbreak{}S-b d. (3.3)\allowbreak{}

The printed h\_\allowbreak{}j includes its upper endpoint. This is a real defect. For example,\allowbreak{} work in D(Ab)\allowbreak{},\allowbreak{} let every K\_\allowbreak{}j\allowbreak{}=Z[0]\allowbreak{} and every f\_\allowbreak{}j\allowbreak{}=id,\allowbreak{} and choose n\_\allowbreak{}i\allowbreak{}=2i. At j\allowbreak{}=3 the printed h has column ι\_\allowbreak{}3\allowbreak{}+ι\_\allowbreak{}4. Its image under u is ι\_\allowbreak{}3-ι\_\allowbreak{}5,\allowbreak{} whereas (id\_\allowbreak{}S-a c)\allowbreak{}ι\_\allowbreak{}3\allowbreak{}=ι\_\allowbreak{}3-ι\_\allowbreak{}4. These maps of degree-zero free abelian groups are different. The corrected formula gives hι\_\allowbreak{}3\allowbreak{}=ι\_\allowbreak{}3 and satisfies (3.3)\allowbreak{}.

Now θ\allowbreak{}=φψ is the endomorphism on C. Its triangle\textquotesingle s first two components are b d and a c. Let e\allowbreak{}=id\_\allowbreak{}C-θ. Then

e v\allowbreak{}=v(id\_\allowbreak{}S-a c)\allowbreak{}\allowbreak{}=v u h\allowbreak{}=0,\allowbreak{} w e\allowbreak{}=(id\_\allowbreak{}S-b d)\allowbreak{}[1]\allowbreak{}w\allowbreak{}=h[1]\allowbreak{}u[1]\allowbreak{}w\allowbreak{}=0.

Exactness gives e\allowbreak{}=t w for some t:\allowbreak{}S[1]\allowbreak{}\allowbreak{}→C; hence e²\allowbreak{}=t w e\allowbreak{}=0. It follows that θ\allowbreak{}=id\_\allowbreak{}C-e is invertible,\allowbreak{} with inverse id\_\allowbreak{}C\allowbreak{}+e. This supplies the source\textquotesingle s omitted small argument and shows why the printed ψφ in original 10292 and 10294 must instead be φψ. DERIVED-NEW-035 records these three reversed composite expressions. DERIVED-NEW-036 records the wrong homotopy endpoint,\allowbreak{} with its zero case and initial-segment indexing made explicit in the corrected formula.

Finally both ψφ and φψ are invertible. The two candidates (ψφ)\allowbreak{}\textasciicircum{}(-1)\allowbreak{}ψ and ψ(φψ)\allowbreak{}\textasciicircum{}(-1)\allowbreak{} for an inverse to φ coincide:\allowbreak{} multiply (ψφ)\allowbreak{}ψ\allowbreak{}=ψ(φψ)\allowbreak{} on the left and right by the respective inverses. They are respectively a left and a right inverse,\allowbreak{} so φ is an isomorphism. Its component equation φ v\_\allowbreak{}T\allowbreak{}=v a gives φ v\_\allowbreak{}T κ\_\allowbreak{}i\allowbreak{}=vι\_\allowbreak{}\{n\_\allowbreak{}i\},\allowbreak{} which is precisely the source\textquotesingle s required diagram. No compactness,\allowbreak{} enhancement,\allowbreak{} infinite sum in an individual column,\allowbreak{} or uniqueness of a TR3 completion was assumed.

\subsection{\texorpdfstring{4. Direct sums in the derived category,\allowbreak{} including uniqueness}{4. Direct sums in the derived  including uniqueness}}\label{proofs-10139-10427-direct-sums-in-the-derived-category-including-uniqueness}

Assume A has exact countable direct sums. For complexes K\_\allowbreak{}i let S be their termwise direct sum. Exactness gives H\textasciicircum{}r(S)\allowbreak{}\allowbreak{}=\allowbreak{}⊕\_\allowbreak{}i H\textasciicircum{}r(K\_\allowbreak{}i)\allowbreak{}:\allowbreak{} apply direct sums to the short exact sequences defining cycles,\allowbreak{} boundaries and cohomology. In particular a direct sum of quasi-isomorphisms is a quasi-isomorphism.

For maps α\_\allowbreak{}i:\allowbreak{}K\_\allowbreak{}i\allowbreak{}→L in D(A)\allowbreak{},\allowbreak{} choose roofs K\_\allowbreak{}i←s\_\allowbreak{}i M\_\allowbreak{}i\allowbreak{}→f\_\allowbreak{}i L with s\_\allowbreak{}i a quasi-isomorphism. Their coproduct s:\allowbreak{}\allowbreak{}⊕M\_\allowbreak{}i\allowbreak{}→S is a quasi-isomorphism,\allowbreak{} and the maps f\_\allowbreak{}i define f:\allowbreak{}\allowbreak{}⊕M\_\allowbreak{}i\allowbreak{}→L. The roof f s\textasciicircum{}(-1)\allowbreak{} restricts to every α\_\allowbreak{}i,\allowbreak{} which proves existence of the required coproduct map.

Here is the uniqueness omitted in the source. Suppose γ:\allowbreak{}S\allowbreak{}→L restricts to zero on every K\_\allowbreak{}i. Represent it by f s\textasciicircum{}(-1)\allowbreak{} with s:\allowbreak{}M\allowbreak{}→S a quasi-isomorphism. For each coprojection K\_\allowbreak{}i\allowbreak{}→S,\allowbreak{} the fraction calculus supplies a quasi-isomorphism t\_\allowbreak{}i:\allowbreak{}N\_\allowbreak{}i\allowbreak{}→K\_\allowbreak{}i and u\_\allowbreak{}i:\allowbreak{}N\_\allowbreak{}i\allowbreak{}→M with s u\_\allowbreak{}i\allowbreak{}=ι\_\allowbreak{}i t\_\allowbreak{}i in K(A)\allowbreak{}. The zero restriction implies f u\_\allowbreak{}i\allowbreak{}=0 in D(A)\allowbreak{}. The zero criterion for localization supplies a further quasi-isomorphism q\_\allowbreak{}i:\allowbreak{}N'\_\allowbreak{}i\allowbreak{}→N\_\allowbreak{}i with f u\_\allowbreak{}i q\_\allowbreak{}i\allowbreak{}=0 in K(A)\allowbreak{}. Replace t\_\allowbreak{}i,\allowbreak{}u\_\allowbreak{}i by their composites with q\_\allowbreak{}i.

Let U:\allowbreak{}\allowbreak{}⊕N'\_\allowbreak{}i\allowbreak{}→M have components u\_\allowbreak{}i and let T\allowbreak{}=\allowbreak{}⊕t\_\allowbreak{}i:\allowbreak{}\allowbreak{}⊕N'\_\allowbreak{}i\allowbreak{}→S. The component homotopies combine to give s U\allowbreak{}=T in K(A)\allowbreak{},\allowbreak{} because maps and homotopies out of a termwise coproduct are specified by maps and homotopies on its summands. T and s are quasi-isomorphisms,\allowbreak{} so U is one. Likewise the component nullhomotopies combine to f U\allowbreak{}=0 in K(A)\allowbreak{}. Consequently γ\allowbreak{}=f U(s U)\allowbreak{}\textasciicircum{}(-1)\allowbreak{}\allowbreak{}=0 in D(A)\allowbreak{}. This proves the full coproduct property without presuming that derived Hom out of a coproduct is a product.

The same argument applies to any indexing set for which the corresponding direct sums exist and are exact; the countable source statement remains the one under review. This routine index-set observation is not counted as another editorial result.

\subsection{5. Exact sequential colimits and the telescope sequence}\label{proofs-10139-10427-exact-sequential-colimits-and-the-telescope-sequence}

Assume sequential colimits exist and are exact in A. The colimit of the finite partial sums gives a countable coproduct:\allowbreak{} maps from the colimit to an object are precisely compatible maps from the finite sums,\allowbreak{} equivalently a family of maps from their summands. Finite sums are exact,\allowbreak{} and taking the assumed exact sequential colimit proves countable sums are exact. The existing correction 0885 changes only the source\textquotesingle s plural verb.

For a sequence A\_\allowbreak{}1\allowbreak{}→A\_\allowbreak{}2\allowbreak{}→… set B\_\allowbreak{}n\allowbreak{}=\allowbreak{}⊕\_\allowbreak{}\{1≤j≤n\}A\_\allowbreak{}j and B'\_\allowbreak{}n\allowbreak{}=\allowbreak{}⊕\_\allowbreak{}\{1≤j<n\}A\_\allowbreak{}j. Define β\_\allowbreak{}n:\allowbreak{}B'\_\allowbreak{}n\allowbreak{}→B\_\allowbreak{}n by columns ι\_\allowbreak{}j-ι\_\allowbreak{}\{j\allowbreak{}+1\}f\_\allowbreak{}j and q\_\allowbreak{}n:\allowbreak{}B\_\allowbreak{}n\allowbreak{}→A\_\allowbreak{}n by columns t\_\allowbreak{}\{n,\allowbreak{}j\}. Then q\_\allowbreak{}n β\_\allowbreak{}n\allowbreak{}=0,\allowbreak{} and q\_\allowbreak{}n is split epic using the last summand. The sequence 0\allowbreak{}→B'\_\allowbreak{}n\allowbreak{}→B\_\allowbreak{}n\allowbreak{}→A\_\allowbreak{}n\allowbreak{}→0 is exact. For a completely explicit verification,\allowbreak{} β\_\allowbreak{}n is injective by successive coordinates:\allowbreak{} its first coordinate is x\_\allowbreak{}1 and its j-th coordinate is x\_\allowbreak{}j-f\_\allowbreak{}\{j-1\}x\_\allowbreak{}\{j-1\}. An element y of ker q\_\allowbreak{}n has unique preimage with x\_\allowbreak{}1\allowbreak{}=y\_\allowbreak{}1 and x\_\allowbreak{}j\allowbreak{}=y\_\allowbreak{}j\allowbreak{}+f\_\allowbreak{}\{j-1\}x\_\allowbreak{}\{j-1\} for 1\textless j\textless n. The last-coordinate equation is precisely q\_\allowbreak{}n(y)\allowbreak{}\allowbreak{}=0. These calculations are matrix identities of finite biproduct maps,\allowbreak{} so hold in any abelian category,\allowbreak{} without an element interpretation.

The transition maps B'\_\allowbreak{}n\allowbreak{}→B'\_\allowbreak{}\{n\allowbreak{}+1\} and B\_\allowbreak{}n\allowbreak{}→B\_\allowbreak{}\{n\allowbreak{}+1\} are inclusions; the right transition is f\_\allowbreak{}n. The column equations show these form a morphism of short exact sequences. Their colimit is therefore

0\allowbreak{}→\allowbreak{}⊕\_\allowbreak{}n A\_\allowbreak{}n --(1-f)\allowbreak{}-\allowbreak{}-> \allowbreak{}⊕\_\allowbreak{}n A\_\allowbreak{}n\allowbreak{}→colim\_\allowbreak{}n A\_\allowbreak{}n\allowbreak{}→0. (5.1)\allowbreak{}

For a sequence of complexes,\allowbreak{} apply (5.1)\allowbreak{} in every degree. The differentials commute with every map,\allowbreak{} so this is a short exact sequence of complexes. Its distinguished triangle in D(A)\allowbreak{},\allowbreak{} together with Section 4,\allowbreak{} identifies the actual termwise colimit and its canonical stage maps as a homotopy colimit. The source\textquotesingle s exact sequential-colimit assumption is retained.

\subsection{6. Homological functors and compact-source maps}\label{proofs-10139-10427-homological-functors-and-compact-source-maps}

Let H:\allowbreak{}D\allowbreak{}→A be homological,\allowbreak{} commute with countable coproducts,\allowbreak{} and let A have exact sequential colimits. Apply H to (1.1)\allowbreak{}. Shift is an equivalence,\allowbreak{} hence preserves the existing coproduct S. The exact segment is

\allowbreak{}⊕H(K\_\allowbreak{}n)\allowbreak{} --(1-H(f)\allowbreak{})\allowbreak{}-\allowbreak{}-> \allowbreak{}⊕H(K\_\allowbreak{}n)\allowbreak{}\allowbreak{}→H(C)\allowbreak{}\allowbreak{}→ \allowbreak{}⊕H(K\_\allowbreak{}n[1]\allowbreak{})\allowbreak{} --(1-H(f[1]\allowbreak{})\allowbreak{})\allowbreak{}-\allowbreak{}-> \allowbreak{}⊕H(K\_\allowbreak{}n[1]\allowbreak{})\allowbreak{}.

The last map is monic by (5.1)\allowbreak{}; its first map\textquotesingle s cokernel is colim H(K\_\allowbreak{}n)\allowbreak{}. Thus the induced map colim H(K\_\allowbreak{}n)\allowbreak{}\allowbreak{}→H(C)\allowbreak{} is an isomorphism. Its components are H(i\_\allowbreak{}n)\allowbreak{},\allowbreak{} so the asserted identification is the actual canonical one.

Take A\allowbreak{}=Ab and H\allowbreak{}=Hom\_\allowbreak{}D(P,\allowbreak{}-)\allowbreak{}. It is homological by the triangle axioms. If Hom\_\allowbreak{}D(P,\allowbreak{}-)\allowbreak{} commutes with countable sums,\allowbreak{} the preceding argument proves the source comparison

λ:\allowbreak{}colim\_\allowbreak{}n Hom\_\allowbreak{}D(P,\allowbreak{}K\_\allowbreak{}n)\allowbreak{}\allowbreak{}→Hom\_\allowbreak{}D(P,\allowbreak{}C)\allowbreak{}

is bijective. It follows explicitly that every map P\allowbreak{}→C factors through some i\_\allowbreak{}n and that two stage maps giving the same map to C become equal at a later common stage. Both assertions are the element description of a sequential colimit of abelian groups. No countability assertion about the individual Hom groups is needed.

\subsection{7. The exact comparison defect without compactness}\label{proofs-10139-10427-the-exact-comparison-defect-without-compactness}

DERIVED-UNDER-024 gives a criterion for this one sequence and object P. Do not assume P compact or that Hom(P,\allowbreak{}-)\allowbreak{} preserves any infinite sum. Retain the actual triangle (1.1)\allowbreak{}. For r\allowbreak{}=0,\allowbreak{}1 define abelian groups

A\_\allowbreak{}r\allowbreak{}=\allowbreak{}⊕\_\allowbreak{}n Hom\_\allowbreak{}D(P,\allowbreak{}K\_\allowbreak{}n[r]\allowbreak{})\allowbreak{},\allowbreak{} B\_\allowbreak{}r\allowbreak{}=Hom\_\allowbreak{}D(P,\allowbreak{}S[r]\allowbreak{})\allowbreak{},\allowbreak{} Q\_\allowbreak{}r\allowbreak{}=B\_\allowbreak{}r/\allowbreak{}A\_\allowbreak{}r.

The canonical map A\_\allowbreak{}r\allowbreak{}→B\_\allowbreak{}r is injective:\allowbreak{} if a finite sum of stage maps gives zero in B\_\allowbreak{}r,\allowbreak{} compose with each coordinate projection S[r]\allowbreak{}\allowbreak{}→K\_\allowbreak{}n[r]\allowbreak{} to show that every component is zero. Let d\_\allowbreak{}r:\allowbreak{}B\_\allowbreak{}r\allowbreak{}→B\_\allowbreak{}r be postcomposition by u[r]\allowbreak{}. Its restriction δ\_\allowbreak{}r to A\_\allowbreak{}r has the original finite-support columns 1-(f\_\allowbreak{}n[r]\allowbreak{})\allowbreak{}\_\allowbreak{}*; it is injective,\allowbreak{} by successive components. It induces an endomorphism e\_\allowbreak{}r:\allowbreak{}Q\_\allowbreak{}r\allowbreak{}→Q\_\allowbreak{}r. These definitions retain all the actual maps,\allowbreak{} not only their kernels.

The triangle yields a short exact sequence

0\allowbreak{}→coker d\_\allowbreak{}0 --v\_\allowbreak{}*-\allowbreak{}-> Hom\_\allowbreak{}D(P,\allowbreak{}C)\allowbreak{} --w\_\allowbreak{}*-\allowbreak{}-> ker d\_\allowbreak{}1\allowbreak{}→0. (7.1)\allowbreak{}

The square of exact sequences 0\allowbreak{}→A\_\allowbreak{}0\allowbreak{}→B\_\allowbreak{}0\allowbreak{}→Q\_\allowbreak{}0\allowbreak{}→0 with vertical maps δ\_\allowbreak{}0,\allowbreak{}d\_\allowbreak{}0,\allowbreak{}e\_\allowbreak{}0 has kernel-cokernel exact sequence

0\allowbreak{}→ker d\_\allowbreak{}0\allowbreak{}→ker e\_\allowbreak{}0 --∂-\allowbreak{}-> coker δ\_\allowbreak{}0 \allowbreak{}→coker d\_\allowbreak{}0\allowbreak{}→coker e\_\allowbreak{}0\allowbreak{}→0. (7.2)\allowbreak{}

Here ∂ has an exact representative description. If q\allowbreak{}∈Q\_\allowbreak{}0 is killed by e\_\allowbreak{}0,\allowbreak{} choose b\allowbreak{}∈B\_\allowbreak{}0 lifting q. Then d\_\allowbreak{}0(b)\allowbreak{} lies in A\_\allowbreak{}0,\allowbreak{} and ∂q is its class in coker δ\_\allowbreak{}0. Changing b by an element of A\_\allowbreak{}0 changes this representative by im δ\_\allowbreak{}0. This proves well-definedness,\allowbreak{} and the kernel and image assertions follow by lifting and subtracting representatives in the two rows.

Since coker δ\_\allowbreak{}0\allowbreak{}=colim Hom\_\allowbreak{}D(P,\allowbreak{}K\_\allowbreak{}n)\allowbreak{},\allowbreak{} its composite into (7.1)\allowbreak{} is exactly λ,\allowbreak{} not a different comparison. Equations (7.1)\allowbreak{}–\allowbreak{}(7.2)\allowbreak{} give

ker λ ≅ coker(ker d\_\allowbreak{}0\allowbreak{}→ker e\_\allowbreak{}0)\allowbreak{},\allowbreak{} (7.3)\allowbreak{} 0\allowbreak{}→coker e\_\allowbreak{}0\allowbreak{}→coker λ\allowbreak{}→ker d\_\allowbreak{}1\allowbreak{}→0. (7.4)\allowbreak{}

In (7.3)\allowbreak{},\allowbreak{} the isomorphism is induced by ∂; in (7.4)\allowbreak{},\allowbreak{} the first injection comes from the quotient of coker d\_\allowbreak{}0 by im(coker δ\_\allowbreak{}0)\allowbreak{},\allowbreak{} and the last map is induced by w\_\allowbreak{}*. Thus λ is bijective exactly when

ker d\_\allowbreak{}0\allowbreak{}→ker e\_\allowbreak{}0 is surjective,\allowbreak{} coker e\_\allowbreak{}0\allowbreak{}=0,\allowbreak{} ker d\_\allowbreak{}1\allowbreak{}=0. (7.5)\allowbreak{}

In particular it suffices that e\_\allowbreak{}0 is invertible and d\_\allowbreak{}1 is injective. It also suffices that just the two comparison maps A\_\allowbreak{}r\allowbreak{}→B\_\allowbreak{}r for r\allowbreak{}=0,\allowbreak{}1 are surjective:\allowbreak{} then Q\_\allowbreak{}0\allowbreak{}=Q\_\allowbreak{}1\allowbreak{}=0 and d\_\allowbreak{}1\allowbreak{}=δ\_\allowbreak{}1 is injective. This is weaker than requiring preservation of every countable sum. The source\textquotesingle s compactness assumption implies these conditions,\allowbreak{} and its theorem remains unchanged.

The distinction is substantive even for familiar categories. For a constant sequence K\_\allowbreak{}n\allowbreak{}=X and f\_\allowbreak{}n\allowbreak{}=id,\allowbreak{} define ε:\allowbreak{}S\allowbreak{}→X by ε ι\_\allowbreak{}n\allowbreak{}=id and σ:\allowbreak{}S\allowbreak{}→S by

σ ι\_\allowbreak{}n\allowbreak{}=-Σ\_\allowbreak{}\{1≤j<n\}ι\_\allowbreak{}j.

Finite telescoping gives σ u\allowbreak{}=id\_\allowbreak{}S and u σ\allowbreak{}=id\_\allowbreak{}S-ι\_\allowbreak{}1 ε,\allowbreak{} while εu\allowbreak{}=0 and ει\_\allowbreak{}1\allowbreak{}=id\_\allowbreak{}X. The maps (σ,\allowbreak{}ε)\allowbreak{}:\allowbreak{}S\allowbreak{}→S\allowbreak{}⊕X and (u,\allowbreak{}ι\_\allowbreak{}1)\allowbreak{}:\allowbreak{}S\allowbreak{}⊕X\allowbreak{}→S are inverse:\allowbreak{} the two displayed identities give their products,\allowbreak{} with σι\_\allowbreak{}1\allowbreak{}=0 and εu\allowbreak{}=0 supplying the off-diagonal zeros. Therefore u is identified with the split inclusion S\allowbreak{}→S\allowbreak{}⊕X,\allowbreak{} whose distinguished cone triangle has third object X. The comparison from (1.1)\allowbreak{} gives C≅X with every i\_\allowbreak{}n\allowbreak{}=id\_\allowbreak{}X. Hence λ is bijective for every P.

This example also directly satisfies the weaker criterion:\allowbreak{} σ preserves the finite-support subgroup A\_\allowbreak{}r,\allowbreak{} and on Q\_\allowbreak{}r its products with e\_\allowbreak{}r are identities because ι\_\allowbreak{}1 ε factors through one summand. Thus e\_\allowbreak{}0 is invertible; d\_\allowbreak{}1 is split monic with left inverse induced by σ[1]\allowbreak{}. For a concrete noncompact P take \allowbreak{}⊕\_\allowbreak{}m Z[0]\allowbreak{} in D(Ab)\allowbreak{}. The identity map P\allowbreak{}→\allowbreak{}⊕\_\allowbreak{}m Z[0]\allowbreak{} has infinitely many nonzero target coordinates and therefore is not in \allowbreak{}⊕\_\allowbreak{}m Hom(P,\allowbreak{}Z[0]\allowbreak{})\allowbreak{}. Nevertheless the constant-sequence comparison above holds for it. This proves an actual application beyond the global compactness hypothesis,\allowbreak{} without claiming a new result in the literature.

\subsection{\texorpdfstring{8. Propagation,\allowbreak{} report reconciliation,\allowbreak{} and continuation}{8.  report  and continuation}}\label{proofs-10139-10427-propagation-report-reconciliation-and-continuation}

More on Algebra current 23899–\allowbreak{}23917 states the Hom sequence proved with its actual maps in Section 2. Its current 21011–\allowbreak{}21054 specializes Section 6 to D(R)\allowbreak{} and an actual termwise colimit,\allowbreak{} justified by Section 5. Sections 6–\allowbreak{}7 retain that statement and give a separate exact description of its possible failure if one removes its compactness assumption. The earlier square-zero ideal DERIVED-UNDER-001 is propagated to the telescope completions and corrected subsequence proof in Sections 1–\allowbreak{}3.

Perfect Complexes current 4393–\allowbreak{}4401 uses only the factorization of a compact-source map through a stage. Section 6 proves this precise invocation; the later generator and factor-through claims have not been adjudicated here. Differential Graded Algebra current 7124–\allowbreak{}7127 uses cohomology commuting with the homotopy colimit,\allowbreak{} as in Section 6; the remainder of its countability equivalence is outside this check.

Duality current 798–\allowbreak{}801 invokes countable coproducts for the family M\_\allowbreak{}n[n]\allowbreak{}. In each complex degree there is only one nonzero term,\allowbreak{} so its termwise sum and product coincide. Section 4 proves the coproduct side. The product lemma belongs to the next review section. The printed period in Ext\textasciicircum{}n\_\allowbreak{}A(B. M\_\allowbreak{}n)\allowbreak{} at current 805 is routed to that chapter\textquotesingle s report comparison; it is not counted here.

Two physical French reports,\allowbreak{} DERIVED-070 and DERIVED-074,\allowbreak{} are semantically reconciled with existing units 0881 and 0885. Groups DERIVED-RECON-079 and 080 bring the totals to 137/\allowbreak{}176 reports,\allowbreak{} 80 groups,\allowbreak{} 126 already corrected reports,\allowbreak{} ten missing copyedit/\allowbreak{}notation reports and one missing mathematical-index report. There are 73 prior Derived units and 90 operations reconciled. NEW-035 and NEW-036 bring the unreported finding count to 36; their four operations bring the cumulative operation count to 69. UNDER-024 brings the editorial consequence count to 24.

Review is complete through original 10427. Sequential reading has reached 10630,\allowbreak{} including read-ahead into Derived limits; this is not a claim of semantic acceptance there. Continue semantic review at 10435,\allowbreak{} with unread source starting 10631. The producer handoffs remain queued and no new sessions,\allowbreak{} live source composition,\allowbreak{} build,\allowbreak{} Lean run,\allowbreak{} cleanup,\allowbreak{} publication or producer mathematical acceptance occurred in this pass.
\clearpage
\section{Derived limits and the exact truncation-recovery obstruction}\label{proofs-10435-10686-derived-limits-and-the-exact-truncation-recovery-obstruction}

Private source-order review of original derived.tex 10435–\allowbreak{}10686. Authority:\allowbreak{} Stacks a04446e5\allowbreak{}7ec1fbc2\allowbreak{}52a871af\allowbreak{}cec7752f\allowbreak{}b2807b14. Receiving public edition:\allowbreak{} 95a51593\allowbreak{}aceb247a\allowbreak{}c7a18311\allowbreak{}1678c3a0\allowbreak{}f8b8f4b5. The operation and source-reading records bind the files. The two consequences in Sections 6–\allowbreak{}7 are separate editorial results; the original hypotheses and theorem statements are retained. No novelty is claimed.

Human source:\allowbreak{} \href{https://github.com/stacks/stacks-project/blob/a04446e57ec1fbc252a871afcec7752fb2807b14/derived.tex\#L10435}{The Stacks Project,\allowbreak{} original Derived Categories TeX}. The receiving source is \href{https://github.com/stacks/stacks-project/blob/a04446e57ec1fbc252a871afcec7752fb2807b14/more-algebra.tex}{More on Algebra,\allowbreak{} original TeX},\allowbreak{} lemmas named lemma-pro-isomorphism-iso-Rlim,\allowbreak{} lemma-map-into-Rlim,\allowbreak{} and lemma-pro-isomorphism-Rlim.

\subsection{1. Products and the defining triangle}\label{proofs-10435-10686-products-and-the-defining-triangle}

For a tower with transition maps f\_\allowbreak{}\{n\allowbreak{}+1\}:\allowbreak{}K\_\allowbreak{}\{n\allowbreak{}+1\}\allowbreak{}→K\_\allowbreak{}n,\allowbreak{} let P\allowbreak{}=∏\_\allowbreak{}n K\_\allowbreak{}n with projections π\_\allowbreak{}n and define Δ:\allowbreak{}P\allowbreak{}→P by

π\_\allowbreak{}n Δ\allowbreak{}=π\_\allowbreak{}n-f\_\allowbreak{}\{n\allowbreak{}+1\}π\_\allowbreak{}\{n\allowbreak{}+1\}.

This is a map specified by its product components,\allowbreak{} not an assumption that the K\_\allowbreak{}n have elements. A homotopy limit is a distinguished triangle

C --j-\allowbreak{}-> P --Δ-\allowbreak{}-> P --w-\allowbreak{}-> C[1]\allowbreak{}. (1.1)\allowbreak{}

Existence follows by completing Δ to a triangle and rotating. Two such triangles have an isomorphism respecting j,\allowbreak{} Δ and w,\allowbreak{} by TR3 and the two-isomorphisms criterion for maps of triangles. This need not choose a unique map between their first objects.

Suppose A is abelian with exact countable products. The termwise product of complexes has H\textasciicircum{}p(∏K\_\allowbreak{}n)\allowbreak{}\allowbreak{}=∏H\textasciicircum{}p(K\_\allowbreak{}n)\allowbreak{}:\allowbreak{} apply exact products to the short exact sequences for kernels,\allowbreak{} images and cohomology. Thus products preserve quasi-isomorphisms. For maps L\allowbreak{}→K\_\allowbreak{}n in D(A)\allowbreak{},\allowbreak{} choose right roofs L\allowbreak{}→f\_\allowbreak{}n M\_\allowbreak{}n←s\_\allowbreak{}n K\_\allowbreak{}n,\allowbreak{} where s\_\allowbreak{}n is a quasi-isomorphism. Then s\allowbreak{}=∏s\_\allowbreak{}n is a quasi-isomorphism,\allowbreak{} and the maps f\_\allowbreak{}n give f:\allowbreak{}L\allowbreak{}→∏M\_\allowbreak{}n. The right roof s\textasciicircum{}(-1)\allowbreak{}f gives the desired product map.

Here is the uniqueness omitted in the source. Put P\allowbreak{}=∏K\_\allowbreak{}n termwise,\allowbreak{} and suppose γ:\allowbreak{}L\allowbreak{}→P has zero composite with every projection. Write γ\allowbreak{}=s\textasciicircum{}(-1)\allowbreak{}f with s:\allowbreak{}P\allowbreak{}→M a quasi-isomorphism. For each projection p\_\allowbreak{}n:\allowbreak{}P\allowbreak{}→K\_\allowbreak{}n,\allowbreak{} right fraction calculus gives u\_\allowbreak{}n:\allowbreak{}M\allowbreak{}→N\_\allowbreak{}n and a quasi-isomorphism t\_\allowbreak{}n:\allowbreak{}K\_\allowbreak{}n\allowbreak{}→N\_\allowbreak{}n with u\_\allowbreak{}n s\allowbreak{}=t\_\allowbreak{}n p\_\allowbreak{}n in K(A)\allowbreak{}. Since u\_\allowbreak{}n f is zero in D(A)\allowbreak{},\allowbreak{} the zero criterion for localization allows postcomposition by a further quasi-isomorphism N\_\allowbreak{}n\allowbreak{}→N'\_\allowbreak{}n to make u\_\allowbreak{}n f zero in K(A)\allowbreak{}. Use these refined u\_\allowbreak{}n and t\_\allowbreak{}n.

Form U:\allowbreak{}M\allowbreak{}→∏N'\_\allowbreak{}n and T\allowbreak{}=∏t\_\allowbreak{}n:\allowbreak{}P\allowbreak{}→∏N'\_\allowbreak{}n. Coordinate homotopies combine to U s\allowbreak{}=T and U f\allowbreak{}=0 in K(A)\allowbreak{},\allowbreak{} by the degreewise product property. T and s are quasi-isomorphisms,\allowbreak{} so U is one. It follows that γ\allowbreak{}=(U s)\allowbreak{}\textasciicircum{}(-1)\allowbreak{}U f\allowbreak{}=0. This proves the product property in D(A)\allowbreak{} without assuming it in the uniqueness argument.

With enough injectives and merely existing countable products,\allowbreak{} a tower whose individual objects lie in D\textasciicircum{}\allowbreak{}+(A)\allowbreak{} also has a derived product:\allowbreak{} choose a bounded below injective representative I\_\allowbreak{}n for each object. Every I\_\allowbreak{}n is K-injective,\allowbreak{} and the full Hom-complex product argument in PART14 proves ∏I\_\allowbreak{}n is K-injective and represents the product in D(A)\allowbreak{}. No common lower bound is required. Its difference map gives (1.1)\allowbreak{}. This proves the source\textquotesingle s second existence criterion without adding exactness of products.

\subsection{2. Maps into a homotopy limit}\label{proofs-10435-10686-maps-into-a-homotopy-limit}

Use the shifted boundary k\allowbreak{}=-w[-1]\allowbreak{}:\allowbreak{}P[-1]\allowbreak{}\allowbreak{}→C. With the source rotation convention this is the first arrow of the distinguished triangle P[-1]\allowbreak{}\allowbreak{}→C\allowbreak{}→P\allowbreak{}→P. Applying Hom(T,\allowbreak{}-)\allowbreak{} gives the exact segment

Hom(T,\allowbreak{}P[-1]\allowbreak{})\allowbreak{} --Δ[-1]\allowbreak{}\_\allowbreak{}*-\allowbreak{}-> Hom(T,\allowbreak{}P[-1]\allowbreak{})\allowbreak{} --k\_\allowbreak{}*-\allowbreak{}-> Hom(T,\allowbreak{}C)\allowbreak{} --j\_\allowbreak{}*-\allowbreak{}-> Hom(T,\allowbreak{}P)\allowbreak{} --Δ\_\allowbreak{}*-\allowbreak{}-> Hom(T,\allowbreak{}P)\allowbreak{}.

Representable Hom preserves products. Thus Δ\_\allowbreak{}* is precisely the difference map on ∏Hom(T,\allowbreak{}K\_\allowbreak{}n)\allowbreak{},\allowbreak{} and the preceding map is the difference map on ∏Hom(T,\allowbreak{}K\_\allowbreak{}n[-1]\allowbreak{})\allowbreak{}. Consequently

0\allowbreak{}→coker Δ\_\allowbreak{}\{Hom(T,\allowbreak{}K\_\allowbreak{}n[-1]\allowbreak{})\allowbreak{}\}\allowbreak{}→Hom(T,\allowbreak{}C)\allowbreak{} \allowbreak{}→lim\_\allowbreak{}n Hom(T,\allowbreak{}K\_\allowbreak{}n)\allowbreak{}\allowbreak{}→0. (2.1)\allowbreak{}

The injection sends a family b:\allowbreak{}T\allowbreak{}→P[-1]\allowbreak{} to k b\allowbreak{}=-w[-1]\allowbreak{}b,\allowbreak{} with its indicated sign. The right map is given by all π\_\allowbreak{}n j. The cokernel is the R\^{}1 lim of this abelian-group tower,\allowbreak{} by the explicit two-term complex from More on Algebra. In particular compatible maps T\allowbreak{}→K\_\allowbreak{}n lift to T\allowbreak{}→C,\allowbreak{} and their possible lifts form a torsor under the left group. Nothing here requires exact products in the underlying abelian category when the needed product already exists in D.

\subsection{\texorpdfstring{3. Pro-morphisms:\allowbreak{} exact finite sums and independence}{3.  exact finite sums and independence}}\label{proofs-10435-10686-pro-morphisms-exact-finite-sums-and-independence}

Write t\_\allowbreak{}\{r,\allowbreak{}s\}:\allowbreak{}K\_\allowbreak{}s\allowbreak{}→K\_\allowbreak{}r for the original transition composite,\allowbreak{} with t\_\allowbreak{}\{r,\allowbreak{}r\}\allowbreak{}=id. Let g denote transitions in a second tower L. A pro-morphism may be represented by strictly increasing indices m\_\allowbreak{}n and maps a\_\allowbreak{}n:\allowbreak{}K\_\allowbreak{}\{m\_\allowbreak{}n\}\allowbreak{}→L\_\allowbreak{}n satisfying g\_\allowbreak{}\{n\allowbreak{}+1\}a\_\allowbreak{}\{n\allowbreak{}+1\}\allowbreak{}=a\_\allowbreak{}n t\_\allowbreak{}\{m\_\allowbreak{}n,\allowbreak{}m\_\allowbreak{}\{n\allowbreak{}+1\}\}. To obtain such a representation,\allowbreak{} choose successive representatives of the inverse-limit-of-colimits definition of a pro-morphism. At stage n\allowbreak{}+1,\allowbreak{} choose an index larger than m\_\allowbreak{}n and large enough that the required equality with the n-th representative holds. That equality is available after a sufficiently late transition by the definition of equality in the filtered colimit. Earlier equalities then follow by composition. This also permits m\_\allowbreak{}n≥n.

The source\textquotesingle s two product maps have exact components

π\_\allowbreak{}n a'\allowbreak{}=a\_\allowbreak{}n π\_\allowbreak{}\{m\_\allowbreak{}n\},\allowbreak{} π\_\allowbreak{}n a''\allowbreak{}=Σ\_\allowbreak{}\{m\_\allowbreak{}n≤r<m\_\allowbreak{}\{n\allowbreak{}+1\}\}a\_\allowbreak{}n t\_\allowbreak{}\{m\_\allowbreak{}n,\allowbreak{}r\}π\_\allowbreak{}r. (3.1)\allowbreak{}

The sum in each component is finite. Telescoping proves

π\_\allowbreak{}n a''Δ\_\allowbreak{}K \allowbreak{}=a\_\allowbreak{}n π\_\allowbreak{}\{m\_\allowbreak{}n\}-a\_\allowbreak{}n t\_\allowbreak{}\{m\_\allowbreak{}n,\allowbreak{}m\_\allowbreak{}\{n\allowbreak{}+1\}\}π\_\allowbreak{}\{m\_\allowbreak{}\{n\allowbreak{}+1\}\} \allowbreak{}=π\_\allowbreak{}n Δ\_\allowbreak{}L a\textquotesingle.

Thus Δ\_\allowbreak{}L a'\allowbreak{}=a''Δ\_\allowbreak{}K. TR3,\allowbreak{} after rotation,\allowbreak{} gives a map C\_\allowbreak{}K\allowbreak{}→C\_\allowbreak{}L and a map of triangles with product components a\textquotesingle{} and a\textquotesingle\textquotesingle. This verifies every transition factor in original 10539. The spelling correction in unit 0886 changes none of them.

The construction induces well-defined maps on kernel and cokernel of the abelian-group difference maps. Here are the full comparison homotopies,\allowbreak{} including the omitted representation-independence step. If m\_\allowbreak{}n is enlarged to M\_\allowbreak{}n≥m\_\allowbreak{}n and b\_\allowbreak{}n\allowbreak{}=a\_\allowbreak{}n t\_\allowbreak{}\{m\_\allowbreak{}n,\allowbreak{}M\_\allowbreak{}n\},\allowbreak{} define a map H:\allowbreak{}∏K\_\allowbreak{}r\allowbreak{}→∏L\_\allowbreak{}n by

π\_\allowbreak{}n H\allowbreak{}=Σ\_\allowbreak{}\{m\_\allowbreak{}n≤r<M\_\allowbreak{}n\}a\_\allowbreak{}n t\_\allowbreak{}\{m\_\allowbreak{}n,\allowbreak{}r\}π\_\allowbreak{}r.

The same finite telescoping gives

a'-b'\allowbreak{}=HΔ\_\allowbreak{}K,\allowbreak{} a''-b''\allowbreak{}=Δ\_\allowbreak{}L H. (3.2)\allowbreak{}

For the second equality,\allowbreak{} expand π\_\allowbreak{}n H-g\_\allowbreak{}\{n\allowbreak{}+1\}π\_\allowbreak{}\{n\allowbreak{}+1\}H. Compatibility converts the second sum into one with coefficients a\_\allowbreak{}n t\_\allowbreak{}\{m\_\allowbreak{}n,\allowbreak{}r\}; subtracting the two intervals leaves the interval [m\_\allowbreak{}n,\allowbreak{}m\_\allowbreak{}\{n\allowbreak{}+1\})\allowbreak{} minus [M\_\allowbreak{}n,\allowbreak{}M\_\allowbreak{}\{n\allowbreak{}+1\})\allowbreak{},\allowbreak{} exactly a\textquotesingle\textquotesingle-b\textquotesingle\textquotesingle. Thus enlargement does not change the maps on ker Δ or coker Δ. Two representatives of the same pro-map have a common increasing enlargement where their component maps agree; (3.2)\allowbreak{} proves their maps on these groups agree.

For a second pro-map b\_\allowbreak{}n:\allowbreak{}L\_\allowbreak{}\{ℓ\_\allowbreak{}n\}\allowbreak{}→N\_\allowbreak{}n,\allowbreak{} the composite uses indices m\_\allowbreak{}\{ℓ\_\allowbreak{}n\} and components b\_\allowbreak{}n a\_\allowbreak{}\{ℓ\_\allowbreak{}n\}. The degree-zero product map is b\textquotesingle a\textquotesingle. In degree one,\allowbreak{} expand b''a'':\allowbreak{} its outer interval ℓ\_\allowbreak{}n≤q<ℓ\_\allowbreak{}\{n\allowbreak{}+1\} and inner intervals m\_\allowbreak{}q≤r<m\_\allowbreak{}\{q\allowbreak{}+1\} partition [m\_\allowbreak{}\{ℓ\_\allowbreak{}n\},\allowbreak{}m\_\allowbreak{}\{ℓ\_\allowbreak{}\{n\allowbreak{}+1\}\})\allowbreak{}. Compatibility converts each coefficient to b\_\allowbreak{}n a\_\allowbreak{}\{ℓ\_\allowbreak{}n\}t\_\allowbreak{}\{m\_\allowbreak{}\{ℓ\_\allowbreak{}n\},\allowbreak{}r\}. Hence its degree-one map is exactly the composite\textquotesingle s map. For an identity represented by transitions t\_\allowbreak{}\{n,\allowbreak{}m\_\allowbreak{}n\},\allowbreak{} compare to m\_\allowbreak{}n\allowbreak{}=n in (3.2)\allowbreak{}. The induced maps on kernels and cokernels are identities. This proves functoriality and invariance under a pro-isomorphism.

Apply this to the towers Hom(T,\allowbreak{}K\_\allowbreak{}n)\allowbreak{} and Hom(T,\allowbreak{}K\_\allowbreak{}n[-1]\allowbreak{})\allowbreak{}. A pro-isomorphism gives isomorphisms on both outer terms of (2.1)\allowbreak{}. The map of triangles supplies a commuting map of its short exact sequences,\allowbreak{} so the map Hom(T,\allowbreak{}C\_\allowbreak{}K)\allowbreak{}\allowbreak{}→Hom(T,\allowbreak{}C\_\allowbreak{}L)\allowbreak{} is an isomorphism for every T. Yoneda then makes C\_\allowbreak{}K\allowbreak{}→C\_\allowbreak{}L an isomorphism. This proves the asserted receiving More on Algebra result using the same a',\allowbreak{}a'' as the source,\allowbreak{} including any TR3 completion. It does not assert that the resulting map itself is uniquely chosen.

\subsection{4. The truncation tower and choice of its comparison}\label{proofs-10435-10686-the-truncation-tower-and-choice-of-its-comparison}

Let T\_\allowbreak{}n\allowbreak{}=τ\_\allowbreak{}\{\textbackslash{}geq -n\}K and assume its product in D(A)\allowbreak{} exists. Let C be its homotopy limit. The canonical maps c\_\allowbreak{}n:\allowbreak{}K\allowbreak{}→T\_\allowbreak{}n are compatible,\allowbreak{} so (2.1)\allowbreak{} supplies c:\allowbreak{}K\allowbreak{}→C with π\_\allowbreak{}n j c\allowbreak{}=c\_\allowbreak{}n. For each fixed p,\allowbreak{} put N\_\allowbreak{}p\allowbreak{}=max(1,\allowbreak{}-p)\allowbreak{}. The cohomology tower has H\textasciicircum{}p(T\_\allowbreak{}n)\allowbreak{}\allowbreak{}=0 for n<N\_\allowbreak{}p and H\textasciicircum{}p(T\_\allowbreak{}n)\allowbreak{}\allowbreak{}=H\textasciicircum{}p(K)\allowbreak{} for n≥N\_\allowbreak{}p,\allowbreak{} with identity transitions on the latter factors.

For m≥N\_\allowbreak{}p,\allowbreak{} the composite H\textasciicircum{}p(K)\allowbreak{}\allowbreak{}→H\textasciicircum{}p(C)\allowbreak{}\allowbreak{}→H\textasciicircum{}p(T\_\allowbreak{}m)\allowbreak{}\allowbreak{}=H\textasciicircum{}p(K)\allowbreak{} is the identity. Thus H\textasciicircum{}p(c)\allowbreak{} is an isomorphism if and only if the second map is an isomorphism:\allowbreak{} when either is invertible the other is its inverse. The second map is independent of c. Any isomorphism of choices of the triangle preserves these projections,\allowbreak{} so the criterion is also independent of the chosen homotopy limit. This proves original 10574–\allowbreak{}10585. Existing unit 0887 removes the redundant connector without changing this argument.

The comparison c is one possible lift,\allowbreak{} not a canonical choice deduced from compatible stage maps. In particular,\allowbreak{} showing that a specific chain map γ supplies one such lift proves the criterion for every lift,\allowbreak{} by the preceding argument; it does not prove that all lifts are literally equal.

\subsection{5. The split injective tower and the actual chain map}\label{proofs-10435-10686-the-split-injective-tower-and-the-actual-chain-map}

Assume countable products and enough injectives in A. Use the strictly commuting diagram of original Lemma lemma-special-inverse-system:\allowbreak{} maps τ\_\allowbreak{}\{\textbackslash{}geq -n\}K\allowbreak{}→I\_\allowbreak{}n are quasi-isomorphisms and the maps f\_\allowbreak{}\{n\allowbreak{}+1\}:\allowbreak{}I\_\allowbreak{}\{n\allowbreak{}+1\}\allowbreak{}→I\_\allowbreak{}n are termwise split epimorphisms. Each I\_\allowbreak{}n is bounded below with injective terms,\allowbreak{} hence K-injective.

Fix a degree p and choose sections s\_\allowbreak{}\{n\allowbreak{}+1\}\textasciicircum{}p of f\_\allowbreak{}\{n\allowbreak{}+1\}\textasciicircum{}p. Put C\_\allowbreak{}1\textasciicircum{}p\allowbreak{}=I\_\allowbreak{}1\textasciicircum{}p and C\_\allowbreak{}n\textasciicircum{}p\allowbreak{}=ker f\_\allowbreak{}n\textasciicircum{}p for n≥2. The successive chosen decompositions identify I\_\allowbreak{}n\textasciicircum{}p with the finite sum of C\_\allowbreak{}1\textasciicircum{}p,\allowbreak{}…,\allowbreak{}C\_\allowbreak{}n\textasciicircum{}p,\allowbreak{} keeping the transition map as the projection dropping its last component. The termwise limit is therefore I\textasciicircum{}p\allowbreak{}=∏\_\allowbreak{}n C\_\allowbreak{}n\textasciicircum{}p,\allowbreak{} with the actual projections λ\_\allowbreak{}n\textasciicircum{}p:\allowbreak{}I\textasciicircum{}p\allowbreak{}→I\_\allowbreak{}n\textasciicircum{}p from those decompositions. The differentials on the I\_\allowbreak{}n induce a unique differential on I commuting with all λ\_\allowbreak{}n; its square is zero after every projection,\allowbreak{} so is zero. Thus this is a complex and the termwise limit.

Let P\textasciicircum{}p\allowbreak{}=∏\_\allowbreak{}n I\_\allowbreak{}n\textasciicircum{}p and let Δ\^{}p have components x\_\allowbreak{}n-f\_\allowbreak{}\{n\allowbreak{}+1\}\textasciicircum{}p x\_\allowbreak{}\{n\allowbreak{}+1\}. Define a section σ\textasciicircum{}p:\allowbreak{}P\textasciicircum{}p\allowbreak{}→P\textasciicircum{}p recursively by

(σ\^{}p y)\allowbreak{}\_\allowbreak{}1\allowbreak{}=0,\allowbreak{} (σ\^{}p y)\allowbreak{}\_\allowbreak{}\{n\allowbreak{}+1\}\allowbreak{}=s\_\allowbreak{}\{n\allowbreak{}+1\}\textasciicircum{}p((σ\textasciicircum{}p y)\allowbreak{}\_\allowbreak{}n-y\_\allowbreak{}n)\allowbreak{}. (5.1)\allowbreak{}

Every component is a finite expression in the projections of y,\allowbreak{} so (5.1)\allowbreak{} defines a morphism in A by the product property. It satisfies Δ\^{}p σ\textasciicircum{}p\allowbreak{}=id exactly. The compatible sequences are exactly im λ\textasciicircum{}p,\allowbreak{} which proves

0\allowbreak{}→I\textasciicircum{}p --λ\textasciicircum{}p-\allowbreak{}-> P\^{}p --Δ\textasciicircum{}p-\allowbreak{}-> P\textasciicircum{}p\allowbreak{}→0 (5.2)\allowbreak{}

is split exact. To identify the other splitting mentioned in the source,\allowbreak{} let ρ\textasciicircum{}p:\allowbreak{}P\textasciicircum{}p\allowbreak{}→∏C\_\allowbreak{}n\textasciicircum{}p take from x\_\allowbreak{}n its C\_\allowbreak{}n\textasciicircum{}p component. Then ρ\^{}p λ\textasciicircum{}p\allowbreak{}=id,\allowbreak{} ρ\^{}p σ\textasciicircum{}p\allowbreak{}=0,\allowbreak{} and λ\textasciicircum{}pρ\textasciicircum{}p\allowbreak{}+σ\textasciicircum{}pΔ\textasciicircum{}p\allowbreak{}=id:\allowbreak{} subtract σ\^{}pΔ\^{}p x from x to get a compatible sequence,\allowbreak{} whose C\_\allowbreak{}n coordinates are precisely ρ\^{}p x.

The splittings need not be chain maps,\allowbreak{} and are not asserted to be. The λ and Δ are chain maps and give a termwise split exact sequence of complexes. Its triangle in K(A)\allowbreak{} and D(A)\allowbreak{} is I\allowbreak{}→P\allowbreak{}→P\allowbreak{}→I[1]\allowbreak{}. The product-of-K-injectives argument from PART14 identifies P as the derived product of the I\_\allowbreak{}n and makes it K-injective. Closure under triangles then makes I K-injective. Equivalently PART15\textquotesingle s split inverse-system argument proves this. The displayed triangle identifies I with holim I\_\allowbreak{}n and,\allowbreak{} through the stage quasi-isomorphisms,\allowbreak{} with holim τ\_\allowbreak{}\{\textbackslash{}geq -n\}K.

The strictly compatible composite maps K\allowbreak{}→τ\_\allowbreak{}\{\textbackslash{}geq -n\}K\allowbreak{}→I\_\allowbreak{}n give a unique actual chain map γ:\allowbreak{}K\allowbreak{}→I by the termwise limit property. It represents a lift c\_\allowbreak{}γ:\allowbreak{}K\allowbreak{}→holim τ\_\allowbreak{}\{\textbackslash{}geq -n\}K with the prescribed stage maps. Hence γ is a quasi-isomorphism exactly when c\_\allowbreak{}γ is an isomorphism. Section 4 makes this equivalent to any chosen comparison c being an isomorphism. This justifies the source\textquotesingle s conclusion without identifying all possible c with a single predetermined map in its displayed diagram. There is no new source-defect count for that choice convention.

No assertion that γ is a quasi-isomorphism has yet been made under countable products alone. Its obstruction is computed next.

\subsection{6. The exact recovery obstruction without exact products}\label{proofs-10435-10686-the-exact-recovery-obstruction-without-exact-products}

DERIVED-UNDER-025 concerns the actual truncation tower T\_\allowbreak{}n,\allowbreak{} assuming only that its product P and homotopy limit C exist in D(A)\allowbreak{}. Exactness of products in A and enough injectives are not assumed. Write A\_\allowbreak{}p\allowbreak{}=H\textasciicircum{}p(P)\allowbreak{} and δ\_\allowbreak{}p\allowbreak{}=H\textasciicircum{}p(Δ)\allowbreak{}. The triangle gives a short exact sequence in A

0\allowbreak{}→C\_\allowbreak{}p\allowbreak{}→H\textasciicircum{}p(C)\allowbreak{} --z\_\allowbreak{}p-\allowbreak{}-> Z\_\allowbreak{}p\allowbreak{}→0,\allowbreak{} C\_\allowbreak{}p\allowbreak{}=coker δ\_\allowbreak{}\{p-1\},\allowbreak{} Z\_\allowbreak{}p\allowbreak{}=ker δ\_\allowbreak{}p. (6.1)\allowbreak{}

The first injection is induced by the shifted boundary of (1.1)\allowbreak{}; thus its sign is the one specified in Section 2. All these kernels and cokernels exist since A is abelian. There is no replacement of A\_\allowbreak{}p by a product of cohomology objects.

For m≥N\_\allowbreak{}p the projection P\allowbreak{}→T\_\allowbreak{}m induces q\_\allowbreak{}p:\allowbreak{}Z\_\allowbreak{}p\allowbreak{}→H\textasciicircum{}p(K)\allowbreak{}. This map is independent of m:\allowbreak{} on Z\_\allowbreak{}p the equations δ\_\allowbreak{}p\allowbreak{}=0 say that adjacent projections are related by the cohomology transition maps,\allowbreak{} which are identities after N\_\allowbreak{}p. The canonical product map κ:\allowbreak{}K\allowbreak{}→P has components c\_\allowbreak{}n. The equality Δκ\allowbreak{}=0 lets H\textasciicircum{}p(κ)\allowbreak{} factor as s\_\allowbreak{}p:\allowbreak{}H\textasciicircum{}p(K)\allowbreak{}\allowbreak{}→Z\_\allowbreak{}p,\allowbreak{} and q\_\allowbreak{}p s\_\allowbreak{}p\allowbreak{}=id. Define B\_\allowbreak{}p\allowbreak{}=ker q\_\allowbreak{}p and

E\_\allowbreak{}p\allowbreak{}=ker(q\_\allowbreak{}p z\_\allowbreak{}p:\allowbreak{}H\textasciicircum{}p(C)\allowbreak{}\allowbreak{}→H\textasciicircum{}p(K)\allowbreak{})\allowbreak{}.

Pull back the epimorphism z\_\allowbreak{}p along B\_\allowbreak{}p\allowbreak{}→Z\_\allowbreak{}p in (6.1)\allowbreak{}. Its kernel remains C\_\allowbreak{}p and the pullback object is E\_\allowbreak{}p,\allowbreak{} so

0\allowbreak{}→C\_\allowbreak{}p\allowbreak{}→E\_\allowbreak{}p\allowbreak{}→B\_\allowbreak{}p\allowbreak{}→0 (6.2)\allowbreak{}

is exact. These are explicit objects with the indicated maps,\allowbreak{} independent of the choice of c. For any lift c with the required stage maps,\allowbreak{} (q\_\allowbreak{}p z\_\allowbreak{}p)\allowbreak{}H\textasciicircum{}p(c)\allowbreak{}\allowbreak{}=id. Consequently

H\textasciicircum{}p(K)\allowbreak{}\allowbreak{}⊕E\_\allowbreak{}p\allowbreak{}→H\textasciicircum{}p(C)\allowbreak{},\allowbreak{} (x,\allowbreak{}e)\allowbreak{}\allowbreak{}↦H\textasciicircum{}p(c)\allowbreak{}x\allowbreak{}+e (6.3)\allowbreak{}

is an isomorphism. Its inverse sends y to ((q\_\allowbreak{}p z\_\allowbreak{}p)\allowbreak{}y,\allowbreak{} y-H\textasciicircum{}p(c)\allowbreak{}(q\_\allowbreak{}p z\_\allowbreak{}p)\allowbreak{}y)\allowbreak{}. This writes out the actual splitting without suppressing the defect. It depends on the chosen c; E\_\allowbreak{}p itself does not.

It follows that every such c is an isomorphism in D(A)\allowbreak{} exactly when

coker δ\_\allowbreak{}\{p-1\}\allowbreak{}=0 and ker(q\_\allowbreak{}p:\allowbreak{}ker δ\_\allowbreak{}p\allowbreak{}→H\textasciicircum{}p(K)\allowbreak{})\allowbreak{}\allowbreak{}=0 for every integer p. (6.4)\allowbreak{}

Indeed (6.2)\allowbreak{} makes these conditions equivalent to E\_\allowbreak{}p\allowbreak{}=0,\allowbreak{} and (6.3)\allowbreak{} then makes every H\textasciicircum{}p(c)\allowbreak{} invertible. Conversely invertibility of c makes the kernel in (6.3)\allowbreak{} zero,\allowbreak{} forcing both outer objects in (6.2)\allowbreak{} to vanish. This is a proved necessary-and-sufficient criterion,\allowbreak{} not an assumption that the missing product exactness holds. Under Section 5\textquotesingle s hypotheses it is also the exact criterion for the specific γ:\allowbreak{}K\allowbreak{}→I to be a quasi-isomorphism.

\subsection{\texorpdfstring{7. Exact products,\allowbreak{} stabilization,\allowbreak{} and removal of an unused hypothesis}{7. Exact   and removal of an unused hypothesis}}\label{proofs-10435-10686-exact-products-stabilization-and-removal-of-an-unused-hypothesis}

DERIVED-UNDER-026 assumes exact countable products in A,\allowbreak{} but does not assume enough injectives. Section 1 supplies products and homotopy limits for every countable tower in D(A)\allowbreak{}. For any such tower K\_\allowbreak{}n its triangle has the cohomology exact sequence

\begingroup\raggedright
0\allowbreak{}→coker(Δ:\allowbreak{}∏H\textasciicircum{}\{p-1\}(K\_\allowbreak{}n)\allowbreak{}\allowbreak{}→∏H\textasciicircum{}\{p-1\}(K\_\allowbreak{}n)\allowbreak{})\allowbreak{} \allowbreak{}→H\textasciicircum{}p(holim K\_\allowbreak{}n)\allowbreak{} \allowbreak{}→ker(Δ:\allowbreak{}∏H\textasciicircum{}p(K\_\allowbreak{}n)\allowbreak{}\allowbreak{}→∏H\textasciicircum{}p(K\_\allowbreak{}n)\allowbreak{})\allowbreak{}\allowbreak{}→0. (7.1)\allowbreak{}
\par\endgroup

All maps are induced by the original projections,\allowbreak{} transitions and boundary of the chosen triangle. This follows from (1.1)\allowbreak{} and H\^{}p of the actual product,\allowbreak{} as proved in Section 1.

Here is the needed vanishing for any eventually isomorphic tower B\_\allowbreak{}n in A,\allowbreak{} with f\_\allowbreak{}\{n\allowbreak{}+1\} isomorphisms for n≥N. Write t\_\allowbreak{}\{r,\allowbreak{}s\}:\allowbreak{}B\_\allowbreak{}s\allowbreak{}→B\_\allowbreak{}r for its unchanged transition composites. Define σ:\allowbreak{}∏B\_\allowbreak{}n\allowbreak{}→∏B\_\allowbreak{}n by finite component formulas

(σ y)\allowbreak{}\_\allowbreak{}N\allowbreak{}=0,\allowbreak{} (σ y)\allowbreak{}\_\allowbreak{}n\allowbreak{}=Σ\_\allowbreak{}\{n≤j<N\}t\_\allowbreak{}\{n,\allowbreak{}j\}y\_\allowbreak{}j,\allowbreak{} n<N,\allowbreak{} (σ y)\allowbreak{}\_\allowbreak{}n\allowbreak{}=-Σ\_\allowbreak{}\{N≤j<n\}(t\_\allowbreak{}\{j,\allowbreak{}n\})\allowbreak{}\textasciicircum{}(-1)\allowbreak{}y\_\allowbreak{}j,\allowbreak{} n\textgreater N. (7.2)\allowbreak{}

Every inverse in the last formula is of a transition within the isomorphism range. Direct cancellation gives (σy)\allowbreak{}\_\allowbreak{}n-f\_\allowbreak{}\{n\allowbreak{}+1\}(σy)\allowbreak{}\_\allowbreak{}\{n\allowbreak{}+1\}\allowbreak{}=y\_\allowbreak{}n,\allowbreak{} also at n\allowbreak{}=N. Thus Δσ\allowbreak{}=id. Its kernel is B\_\allowbreak{}N by the map with components t\_\allowbreak{}\{n,\allowbreak{}N\} for n\textless N and (t\_\allowbreak{}\{N,\allowbreak{}n\})\allowbreak{}\textasciicircum{}(-1)\allowbreak{} for n≥N. The inverse is projection at N. These exact formulas retain the initial finite part and all transition factors.

Apply this to the truncation tower\textquotesingle s H\^{}p. Its nonzero tail starts at N\_\allowbreak{}p\allowbreak{}=max(1,\allowbreak{}-p)\allowbreak{}. The preceding calculation makes C\_\allowbreak{}p\allowbreak{}=0 in (6.1)\allowbreak{} and identifies Z\_\allowbreak{}p with H\textasciicircum{}p(K)\allowbreak{} through the very map q\_\allowbreak{}p,\allowbreak{} so B\_\allowbreak{}p\allowbreak{}=0 as well. Section 6 therefore proves every comparison

K\allowbreak{}→holim\_\allowbreak{}n τ\_\allowbreak{}\{\textbackslash{}geq -n\}K

is an isomorphism,\allowbreak{} without enough injectives. This is the asserted strengthening of the intermediate recovery claim in original 10670–\allowbreak{}10685. Enough injectives remain necessary as an assumption of the source\textquotesingle s separate construction of the I\_\allowbreak{}n; retaining them,\allowbreak{} Section 5 now gives the stated K-injective resolution. The source theorem is not replaced by a stronger existence claim.

The same proof gives a more general convergence assertion. If for a fixed p the towers H\textasciicircum{}\{p-1\}(K\_\allowbreak{}n)\allowbreak{} and H\textasciicircum{}p(K\_\allowbreak{}n)\allowbreak{} are eventually isomorphic,\allowbreak{} (7.1)\allowbreak{} identifies H\^{}p(holim K\_\allowbreak{}n)\allowbreak{} with the eventual H\textasciicircum{}p(K\_\allowbreak{}n)\allowbreak{},\allowbreak{} by the actual stable projection. In fact for this one p it suffices that the first difference map in (7.1)\allowbreak{} is epic and that the kernel of the second projects isomorphically to the chosen stable value. If E\allowbreak{}→K\_\allowbreak{}n are compatible and induce isomorphisms on H\^{}p for all sufficiently large n for each p,\allowbreak{} then any lift E\allowbreak{}→holim K\_\allowbreak{}n is an isomorphism:\allowbreak{} the compatibility makes the cohomology transitions eventually isomorphic in every degree,\allowbreak{} and the stable projection composite is exactly the given isomorphism. The stabilization index may depend on p; no uniform bound is used.

Existing unit MC-STK-ERR-0882 makes n the explicit product index in original 10681 while keeping the condition p≥-n. For fixed p it is precisely the product over n≥N\_\allowbreak{}p used here. No bound on the fixed degree p is substituted for the tower index.

\subsection{8. The source\textquotesingle s omitted dual statements}\label{proofs-10435-10686-the-sources-omitted-dual-statements}

The source says that the duals of its three colimit lemmas could be stated here. They follow with the following exact hypotheses. These supplementary proofs are not counted as source errors.

First,\allowbreak{} assume all countable inverse limits in A exist and the inverse-limit functor on towers is exact. The inverse limit of the finite products ∏\_\allowbreak{}\{j≤n\}A\_\allowbreak{}j,\allowbreak{} with the projection transitions,\allowbreak{} is the countable product by its universal property. Finite products are exact,\allowbreak{} so exact inverse limits make countable products exact. For a tower A\_\allowbreak{}n,\allowbreak{} consider the finite split exact sequence

0\allowbreak{}→A\_\allowbreak{}n\allowbreak{}→∏\_\allowbreak{}\{j≤n\}A\_\allowbreak{}j\allowbreak{}→∏\_\allowbreak{}\{j<n\}A\_\allowbreak{}j\allowbreak{}→0.

The first map has components t\_\allowbreak{}\{j,\allowbreak{}n\}; the second has components x\_\allowbreak{}j-f\_\allowbreak{}\{j\allowbreak{}+1\}x\_\allowbreak{}\{j\allowbreak{}+1\} for j\textless n. Its kernel is the image of the first map,\allowbreak{} and it is split epic:\allowbreak{} prescribe x\_\allowbreak{}n\allowbreak{}=0 and recursively solve x\_\allowbreak{}j\allowbreak{}=y\_\allowbreak{}j\allowbreak{}+f\_\allowbreak{}\{j\allowbreak{}+1\}x\_\allowbreak{}\{j\allowbreak{}+1\} backwards. The sequence maps to the preceding one by the projections on the products and the transition A\_\allowbreak{}n\allowbreak{}→A\_\allowbreak{}\{n-1\}. Its exact inverse limit is

0\allowbreak{}→lim A\_\allowbreak{}n\allowbreak{}→∏A\_\allowbreak{}n --Δ-\allowbreak{}-> ∏A\_\allowbreak{}n\allowbreak{}→0. (8.1)\allowbreak{}

For a tower of complexes,\allowbreak{} apply (8.1)\allowbreak{} termwise. The short exact sequence yields a distinguished triangle in D(A)\allowbreak{},\allowbreak{} and Section 1 identifies its two products as derived products. Thus the termwise inverse limit is a homotopy limit. Exact countable products alone are not substituted for exactness of the inverse-limit functor in this statement.

Finally let D have countable products and let Q satisfy \allowbreak{}⊕Hom\_\allowbreak{}D(E\_\allowbreak{}n,\allowbreak{}Q)\allowbreak{}≅Hom\_\allowbreak{}D(∏E\_\allowbreak{}n,\allowbreak{}Q)\allowbreak{} for every countable family. For any inverse tower with limit triangle (1.1)\allowbreak{},\allowbreak{} applying Hom\_\allowbreak{}D(-,\allowbreak{}Q)\allowbreak{} gives

Hom(P,\allowbreak{}Q)\allowbreak{} --Δ\textasciicircum{}*-\allowbreak{}-> Hom(P,\allowbreak{}Q)\allowbreak{}\allowbreak{}→Hom(C,\allowbreak{}Q)\allowbreak{}\allowbreak{}→ Hom(P[-1]\allowbreak{},\allowbreak{}Q)\allowbreak{} --Δ[-1]\allowbreak{}\textasciicircum{}*-\allowbreak{}-> Hom(P[-1]\allowbreak{},\allowbreak{}Q)\allowbreak{}.

Under the stipulated direct-sum identifications,\allowbreak{} both difference maps have the finite-support columns 1-f\_\allowbreak{}\{n\allowbreak{}+1\}\textasciicircum{}* and are injective. Their cokernel in the first row is colim Hom(K\_\allowbreak{}n,\allowbreak{}Q)\allowbreak{}. Hence the canonical map colim Hom(K\_\allowbreak{}n,\allowbreak{}Q)\allowbreak{}\allowbreak{}→Hom(C,\allowbreak{}Q)\allowbreak{} is an isomorphism. The injection proof is by successive finite-support coordinates,\allowbreak{} exactly as in PART16. This proves the final dual statement without invoking any hypothetical exact inverse-limit functor on D.

\subsection{9. Receiving checks and disposition totals}\label{proofs-10435-10686-receiving-checks-and-disposition-totals}

More on Algebra current 23752–\allowbreak{}23824 and 23828–\allowbreak{}23846 receive the explicit pro-map homotopies and Hom exact sequence in Sections 2–\allowbreak{}3. Its later pro-isomorphism consequence is thereby proved with the actual maps a',\allowbreak{}a'' and compatible boundaries. The repeated spelling at current 23772 and the abbreviated transition notation at current 23820 are routed for that chapter\textquotesingle s own report comparison,\allowbreak{} not newly adjudicated here.

More on Algebra current 14522–\allowbreak{}14533 has both exact products and enough injectives for modules,\allowbreak{} so the K-injective existence and derived-limit invocations follow exactly from Sections 1,\allowbreak{}5,\allowbreak{}7. The ring-module hypotheses were not weakened in this receiving check.

Cohomology current 10338–\allowbreak{}10341 and Cohomology on Sites current 5604–\allowbreak{}5607 use the equivalence between the particular injective tower\textquotesingle s quasi-isomorphism and truncation recovery,\allowbreak{} as proved in Sections 4–\allowbreak{}5. Perfect Complexes current 138–\allowbreak{}143 uses that equivalence in its footnote. These uses are valid; their separate geometric recovery hypotheses are not replaced by an assertion that products of arbitrary sheaves are exact. Cohomology current 10087–\allowbreak{}10116 explicitly verifies local recovery through stable cohomology maps. Section 6 describes the exact global defect those hypotheses must kill,\allowbreak{} while preserving that chapter\textquotesingle s geometric proof for its own review. The product side of Duality current 798–\allowbreak{}801,\allowbreak{} queued in PART16,\allowbreak{} is now justified by Section 1:\allowbreak{} for the shifts M\_\allowbreak{}n[n]\allowbreak{} there is only one nonzero term in each degree,\allowbreak{} and modules have exact products.

Five physical reports are complete in new groups DERIVED-RECON-081,\allowbreak{} 082 and 083:\allowbreak{} two reports each for units 0886 and 0887,\allowbreak{} and one for 0882. Totals are 142/\allowbreak{}176 reports,\allowbreak{} 83 groups,\allowbreak{} 131 already corrected,\allowbreak{} ten missing copyedit/\allowbreak{}notation reports and one missing mathematical-index report. There are 76 prior Derived units and 93 prior operations reconciled. This pass adds no new correction operation or unreported defect. The cumulative 69 proposed operations and 36 unreported findings remain unchanged. UNDER-025 and UNDER-026 bring the editorial consequence count to 26.

Review is complete through original 10686. Reading has reached 10840; the following operations-on-subcategories section is read ahead only. Next semantic review starts at 10691; unread source starts at 10841. The correction queue remains first,\allowbreak{} and no new session,\allowbreak{} live source composition,\allowbreak{} TeX or Lean run,\allowbreak{} cleanup,\allowbreak{} publication or producer mathematical acceptance occurred.
\clearpage
\section{\texorpdfstring{Derived Categories:\allowbreak{} full subcategories and generators}{Derived  full subcategories and generators}}\label{proofs-10691-11108-derived-categories-full-subcategories-and-generators}

Private correction review,\allowbreak{} 2026-09-30. Authority:\allowbreak{} The Stacks Project,\allowbreak{} commit a04446e5\allowbreak{}7ec1fbc2\allowbreak{}52a871af\allowbreak{}cec7752f\allowbreak{}b2807b14, derived.tex 10691–\allowbreak{}11108,\allowbreak{} with its exact-functor definition at 190–\allowbreak{}203. The receiver equiv.tex is frozen separately in authority-equiv-part18.tex and current-equiv-part18.tex. Those two files are identical. Source identities,\allowbreak{} original reports,\allowbreak{} exact replacements and actual reading coverage are bound by PART18\_\allowbreak{}CHECKPOINT.json and its associated ledgers.

The original objects,\allowbreak{} shifts,\allowbreak{} extension triangles and tensor factors are retained below. Sections 6 and 7 are editorial consequences,\allowbreak{} not replacement source statements or novelty claims. The corrections in Sections 2 and 8 are distinct from the seven earlier corrections reconciled here.

\subsection{\texorpdfstring{1. The operations,\allowbreak{} their exact witnesses and finite closure}{1. The  their exact witnesses and finite closure}}\label{proofs-10691-11108-the-operations-their-exact-witnesses-and-finite-closure}

Write T for the given triangulated category. A full subcategory A is an actual subset of Ob(T)\allowbreak{},\allowbreak{} as specified at original 10695–\allowbreak{}10696. Thus A[a,\allowbreak{}b]\allowbreak{} contains precisely the objects A[-i]\allowbreak{} for A in A and integer a <\allowbreak{}= i <\allowbreak{}= b; it is not implicitly closed under arbitrary isomorphism. In contrast,\allowbreak{} smd(A)\allowbreak{} and add(A)\allowbreak{} include objects isomorphic to actual direct summands and finite direct sums,\allowbreak{} respectively. The empty finite sum is allowed. A star B consists of the middle objects of distinguished triangles A \allowbreak{}-> X \allowbreak{}-> B \allowbreak{}-> A[1]\allowbreak{}.

For associativity,\allowbreak{} start with triangles

\[
 A\xrightarrow f X\xrightarrow p B\xrightarrow h A[1],
 \qquad X\xrightarrow g Y\xrightarrow q C\xrightarrow k X[1].
\]

Complete gf to A \allowbreak{}-> Y \allowbreak{}-> Z \allowbreak{}-> A[1]\allowbreak{}. TR4 gives a distinguished triangle B \allowbreak{}-> Z \allowbreak{}-> C \allowbreak{}-> B[1]\allowbreak{},\allowbreak{} whose last map is p[1]\allowbreak{}k,\allowbreak{} and maps of triangles having first two components (id\_\allowbreak{}A,\allowbreak{}g)\allowbreak{} and (f,\allowbreak{}id\_\allowbreak{}Y)\allowbreak{}. Therefore Y belongs to A star (B star C)\allowbreak{}.

Conversely take triangles

\[
 A\xrightarrow f Y\xrightarrow g Z\xrightarrow h A[1],
 \qquad B\xrightarrow u Z\xrightarrow v C\xrightarrow w B[1].
\]

Apply TR4 to Y --g-\allowbreak{}-> Z --v-\allowbreak{}-> C,\allowbreak{} using the rotated triangles (Y,\allowbreak{}Z,\allowbreak{}A[1]\allowbreak{},\allowbreak{}g,\allowbreak{}h,\allowbreak{}-f[1]\allowbreak{})\allowbreak{} and (Z,\allowbreak{}C,\allowbreak{}B[1]\allowbreak{},\allowbreak{}v,\allowbreak{}w,\allowbreak{}-u[1]\allowbreak{})\allowbreak{}. Choose a triangle (Y,\allowbreak{}C,\allowbreak{}W,\allowbreak{}vg,\allowbreak{}p,\allowbreak{}d)\allowbreak{}. Its cone triangle is

\[
 A[1]\xrightarrow\alpha W\xrightarrow\beta B[1]
       \xrightarrow{-(hu)[1]} A[2].
\]

Put X\allowbreak{}=W[-1]\allowbreak{}. Inverse rotation gives X --(-d[-1]\allowbreak{})\allowbreak{}-\allowbreak{}-> Y --vg-\allowbreak{}-> C --p-\allowbreak{}-> X[1]\allowbreak{}. Inverse shift of the cone triangle,\allowbreak{} with its triangle boundary sign,\allowbreak{} gives A --alpha[-1]\allowbreak{}-\allowbreak{}-> X --beta[-1]\allowbreak{}-\allowbreak{}-> B --hu-\allowbreak{}-> A[1]\allowbreak{}. Here maps are transported along the chosen inverse-shift identifications; no literal equality of composed shift functors is assumed. This proves the other inclusion with the original A,\allowbreak{} B,\allowbreak{} C and Y. It supplies the direction left implicit at 10733–\allowbreak{}10735.

Next let A1 \allowbreak{}-> X \allowbreak{}-> B1 \allowbreak{}-> A1[1]\allowbreak{},\allowbreak{} with maps f,\allowbreak{}g,\allowbreak{}h,\allowbreak{} be distinguished. Choose the actual complement witnesses for membership in smd(A)\allowbreak{} and smd(B)\allowbreak{}:\allowbreak{} isomorphisms A1\allowbreak{}+ A2 \allowbreak{}-> A0 in A and B1\allowbreak{}+ B2 \allowbreak{}-> B0 in B. The sum of this triangle with A2 --id-\allowbreak{}-> A2 \allowbreak{}-> 0 and 0 \allowbreak{}-> B2 --id-\allowbreak{}-> B2 gives,\allowbreak{} after permuting the middle summands,\allowbreak{} the distinguished triangle

\[
 A_1\oplus A_2
 \xrightarrow{(a_1,a_2)\mapsto(a_2,f a_1,0)}
 A_2\oplus X\oplus B_2
 \xrightarrow{(a_2,x,b_2)\mapsto(gx,b_2)}
 B_1\oplus B_2
 \xrightarrow{(b_1,b_2)\mapsto(hb_1,0)}
 (A_1\oplus A_2)[1].
\]

Transport its endpoints to A0,\allowbreak{}B0. Its middle object has X as a specified summand,\allowbreak{} proving smd(A)\allowbreak{} star smd(B)\allowbreak{} subset smd(A star B)\allowbreak{}. Actual complement witnesses compose:\allowbreak{} if Y\allowbreak{}=X\allowbreak{}+X' and Z\allowbreak{}=Y\allowbreak{}+Y',\allowbreak{} then Z\allowbreak{}=X\allowbreak{}+(X'\allowbreak{}+Y')\allowbreak{}. Hence smd(smd(A)\allowbreak{})\allowbreak{}\allowbreak{}=smd(A)\allowbreak{}. The reverse inclusion in the asserted equality after another smd follows from A subset smd(A)\allowbreak{} and B subset smd(B)\allowbreak{}. No assumption that every idempotent in T splits has entered this proof.

Finite sums of distinguished triangles are distinguished. Combining the finite lists of summands proves add(A)\allowbreak{} star add(B)\allowbreak{} is closed under finite sums; combining complement witnesses proves the same for smd(add(A)\allowbreak{})\allowbreak{}. The assertions about direct sums at 10763 and 10783 are used in this finite sense,\allowbreak{} as in their displayed binary-sum proof. They do not establish closure under infinite coproducts. For example,\allowbreak{} every object in smd(add(\{Z\})\allowbreak{})\allowbreak{} in D(Ab)\allowbreak{} has finitely generated H\textasciicircum{}0,\allowbreak{} whereas the countable direct sum of Z does not.

Set C\_\allowbreak{}n(A)\allowbreak{}\allowbreak{}=smd(add(A)\allowbreak{}\textasciicircum{}\{star n\})\allowbreak{},\allowbreak{} n>\allowbreak{}=1. The preceding arguments prove that C\_\allowbreak{}n(A)\allowbreak{} is strictly full,\allowbreak{} closed under finite sums and actual summands. Associativity and the smd equality give

\[
 C_{n+m}(A)=smd(C_n(A)\star C_m(A)).
\]

Since 0 is an empty finite sum,\allowbreak{} adding a zero factor through X --id-\allowbreak{}-> X \allowbreak{}-> 0 proves C\_\allowbreak{}n(A)\allowbreak{} subset C\_\allowbreak{}\{n\allowbreak{}+1\}(A)\allowbreak{}. All these operations are monotone in A.

For an increasing sequence A\_\allowbreak{}j,\allowbreak{} every operation in 10823–\allowbreak{}10841 commutes with union:\allowbreak{} a shift or summand has one witness in some A\_\allowbreak{}j; a finite sum has finitely many such witnesses; a fixed n-fold extension has finitely many leaves in its extension tree. Their indices have a common upper bound. The same triangles,\allowbreak{} maps,\allowbreak{} complements and sums then witness membership using that A\_\allowbreak{}j. The reverse inclusions follow from monotonicity. This is an actual finite-witness proof and makes no assertion for infinite sums.

\subsection{2. Exact functors and literal object sets --- DERIVED-NEW-037}\label{proofs-10691-11108-exact-functors-and-literal-object-sets-derived-new-037}

The source definition equips an exact functor F with natural isomorphisms xi\_\allowbreak{}X:\allowbreak{}F(X[1]\allowbreak{})\allowbreak{}\allowbreak{}->F(X)\allowbreak{}[1]\allowbreak{}. Iterating xi,\allowbreak{} its inverse and the chosen inverse-shift identifications gives,\allowbreak{} for every integer i,\allowbreak{}

\[
 \xi_X^{(-i)}:F(X[-i])\xrightarrow{\sim}F(X)[-i].
\]

Consequently F(A[a,\allowbreak{}b]\allowbreak{})\allowbreak{} and F(A)\allowbreak{}[a,\allowbreak{}b]\allowbreak{} have the same isomorphism closure. This is sufficient for every subsequent occurrence inside add,\allowbreak{} smd or C\_\allowbreak{}n. It does not give the literal object-set equality printed at 10800.

Here is a counterexample in the source\textquotesingle s actual category conventions. Let Q be the complex Z --id-\allowbreak{}-> Z in degrees 0 and 1 in D(Ab)\allowbreak{}. Its contracting homotopy has h\textasciicircum{}1\allowbreak{}=id\_\allowbreak{}Z and all other components zero,\allowbreak{} so id\_\allowbreak{}Q\allowbreak{}=0 in D(Ab)\allowbreak{}. Define F on every object to be Q and on every morphism to be the unique morphism Q\allowbreak{}->Q,\allowbreak{} which is simultaneously zero and id\_\allowbreak{}Q. This respects identities and composition and is additive. The unique map Q\allowbreak{}->Q[1]\allowbreak{} is an isomorphism,\allowbreak{} since both are zero objects in D(Ab)\allowbreak{}; it defines xi. Every image triangle is isomorphic to the zero triangle,\allowbreak{} so F is exact.

Take A to have the single object the zero complex,\allowbreak{} and a\allowbreak{}=b\allowbreak{}=1. F(A[1,\allowbreak{}1]\allowbreak{})\allowbreak{} has the sole object Q. F(A)\allowbreak{}[1,\allowbreak{}1]\allowbreak{} has the sole object Q[-1]\allowbreak{}. These are different complexes,\allowbreak{} with different nonzero terms,\allowbreak{} and thus different objects of D(Ab)\allowbreak{},\allowbreak{} although isomorphic there. The literal assertion is false. The repair states equality of isomorphism closures without changing either F or A.

The other four inclusions need no repair. An additive functor sends an actual decomposition X\allowbreak{}+Y \allowbreak{}-> A0 to a decomposition F(X)\allowbreak{}\allowbreak{}+F(Y)\allowbreak{} \allowbreak{}-> F(A0)\allowbreak{},\allowbreak{} and sends a finite sum to its canonically isomorphic finite sum of images. For an extension it sends the distinguished triangle to

\[
 F(A)\xrightarrow{F(f)}F(X)\xrightarrow{F(g)}F(B)
       \xrightarrow{\xi_A F(h)}F(A)[1].
\]

Thus its actual middle object belongs to F(A)\allowbreak{} star F(B)\allowbreak{}. Induction gives the n-fold version. At n\allowbreak{}=1 this is the identity inclusion,\allowbreak{} not an appeal to isomorphism closure of A.

MC-STK-ERR-0888 changes ``objects consists'' to ``objects consist''. It corrects the plural verb and is separate from the formula issue. The public source comparison is \href{https://stacks.math.columbia.edu/tag/0FX5}{the operations remark} and \href{https://stacks.math.columbia.edu/tag/05QK}{the exact-functor definition}. The frozen original TeX,\allowbreak{} not that mutable comparison,\allowbreak{} is the authority.

\subsection{3. Cohomology layers and actual complex terms}\label{proofs-10691-11108-cohomology-layers-and-actual-complex-terms}

Let A be an abelian category and let E be the given set of objects of A,\allowbreak{} embedded in D(A)\allowbreak{} in degree zero. Suppose H\textasciicircum{}i(K)\allowbreak{}\allowbreak{}=0 outside [a,\allowbreak{}b]\allowbreak{}. If a\allowbreak{}=b,\allowbreak{} the canonical truncation maps identify K with H\textasciicircum{}a(K)\allowbreak{}[-a]\allowbreak{}. They induce the identity on its sole cohomology object and are quasi-isomorphisms. This proves the repair in MC-STK-ERR-0889:\allowbreak{} the index is a,\allowbreak{} not an unbound i.

For b\textgreater a the canonical truncation triangle is

\[
 \tau_{\leq b-1}K\longrightarrow K\longrightarrow H^b(K)[-b]
                 \longrightarrow(\tau_{\leq b-1}K)[1].
\]

More precisely,\allowbreak{} its middle term is first tau\_\allowbreak{}\{\textbackslash{}leq b\}K,\allowbreak{} whose canonical map to K is a quasi-isomorphism under the stated upper bound,\allowbreak{} and the displayed triangle is its transport. Its third term is cohomology,\allowbreak{} not the last term of an arbitrary representative. For example,\allowbreak{} for the contractible complex M --id-\allowbreak{}-> M in degrees b-1,\allowbreak{}b,\allowbreak{} both K and its canonical lower truncation are zero in D(A)\allowbreak{},\allowbreak{} but K\textasciicircum{}b[-b]\allowbreak{}\allowbreak{}=M[-b]\allowbreak{} is nonzero when M is nonzero. It cannot be the third term of that triangle. This verifies MC-STK-ERR-0890.

If each H\textasciicircum{}i(K)\allowbreak{} belongs to smd(add(E)\allowbreak{})\allowbreak{},\allowbreak{} then H\textasciicircum{}i(K)\allowbreak{}[-i]\allowbreak{} belongs to C\_\allowbreak{}1(E[a,\allowbreak{}b]\allowbreak{})\allowbreak{},\allowbreak{} by shifting its actual finite-sum and complement witnesses. Induction on b-a,\allowbreak{} using Section 1,\allowbreak{} gives K in C\_\allowbreak{}\{b-a\allowbreak{}+1\}(E[a,\allowbreak{}b]\allowbreak{})\allowbreak{}. Membership H\textasciicircum{}i(K)\allowbreak{} in E is the special case used in part (1)\allowbreak{}. Zero intermediate layers cause no problem because 0 belongs to C\_\allowbreak{}1 regardless of whether E contains zero.

The omitted termwise proof uses a different triangle. If K\^{}bullet has actual terms zero outside [a,\allowbreak{}b]\allowbreak{},\allowbreak{} its top term is a subcomplex,\allowbreak{} and the degreewise split exact sequence of complexes is

\[
 0\longrightarrow K^b[-b]\longrightarrow K^\bullet
   \longrightarrow \sigma_{\leq b-1}K^\bullet\longrightarrow0.
\]

Here sigma is the stupid quotient truncation,\allowbreak{} with last term K\^{}\{b-1\} and outgoing differential zero; it is not canonical cohomological truncation. The connecting map of this sequence retains d\textasciicircum{}\{b-1\}:\allowbreak{}K\textasciicircum{}\{b-1\}\allowbreak{}->K\textasciicircum{}b with the shift sign of its triangle. Removing top terms successively proves parts (3)\allowbreak{} and (4)\allowbreak{},\allowbreak{} in descending term order,\allowbreak{} using K\textasciicircum{}i[-i]\allowbreak{} in C\_\allowbreak{}1(E[a,\allowbreak{}b]\allowbreak{})\allowbreak{}. The entire differential information remains in the extension maps. No direct-sum splitting of K into its terms or its cohomology is asserted.

\subsection{4. Degree support is independent of extension length}\label{proofs-10691-11108-degree-support-is-independent-of-extension-length}

For the homological functor in the source let H\textasciicircum{}q(X)\allowbreak{}\allowbreak{}=H(X[q]\allowbreak{})\allowbreak{}. Shift comparisons give H\textasciicircum{}q(E[-j]\allowbreak{})\allowbreak{}\allowbreak{}=H\textasciicircum{}\{q-j\}(E)\allowbreak{},\allowbreak{} through the specified natural identifications. Thus if E has support in [a,\allowbreak{}b]\allowbreak{},\allowbreak{} E[-j]\allowbreak{} has support in [a\allowbreak{}+j,\allowbreak{}b\allowbreak{}+j]\allowbreak{}. For -m<\allowbreak{}=j<\allowbreak{}=m the union of these intervals is contained in [a-m,\allowbreak{}b\allowbreak{}+m]\allowbreak{}.

Additivity preserves vanishing under finite sums and summands. In a distinguished triangle A\allowbreak{}->X\allowbreak{}->B\allowbreak{}->A[1]\allowbreak{},\allowbreak{} the exact sequence H\textasciicircum{}q(A)\allowbreak{}\allowbreak{}->H\textasciicircum{}q(X)\allowbreak{}\allowbreak{}->H\textasciicircum{}q(B)\allowbreak{} proves H\textasciicircum{}q(X)\allowbreak{}\allowbreak{}=0 whenever its two neighbors vanish. Induction proves that every X in C\_\allowbreak{}n(E[-m,\allowbreak{}m]\allowbreak{})\allowbreak{} has support in [a-m,\allowbreak{}b\allowbreak{}+m]\allowbreak{},\allowbreak{} for every n>\allowbreak{}=1. This verifies MC-STK-ERR-0891 without an unstated boundedness assumption on a chosen representative.

The former interval [-m\allowbreak{}+na,\allowbreak{}m\allowbreak{}+nb]\allowbreak{} is false:\allowbreak{} in D(k)\allowbreak{},\allowbreak{} take the singleton E\allowbreak{}=\{k[-2]\allowbreak{}\},\allowbreak{} a\allowbreak{}=b\allowbreak{}=2,\allowbreak{} m\allowbreak{}=1 and n\allowbreak{}=2. The split triangle E --id-\allowbreak{}-> E \allowbreak{}-> 0 and 0 in add(E[-1,\allowbreak{}1]\allowbreak{})\allowbreak{} put E in C\_\allowbreak{}2(E[-1,\allowbreak{}1]\allowbreak{})\allowbreak{}. Its H\^{}2 is k,\allowbreak{} outside the formerly claimed interval [3,\allowbreak{}5]\allowbreak{}. The argument also proves the more precise support statement:\allowbreak{} for any set S of degrees supporting all E and any set J of permitted shift indices,\allowbreak{} all finite extensions,\allowbreak{} sums and summands of E[-j]\allowbreak{} have support contained in \{s\allowbreak{}+j:\allowbreak{}s in S,\allowbreak{}j in J\}. Extension length does not multiply or fill gaps in that set.

\subsection{\texorpdfstring{5. Generator levels,\allowbreak{} minimality and detection}{5. Generator  minimality and detection}}\label{proofs-10691-11108-generator-levels-minimality-and-detection}

Let A\allowbreak{}=E[-infinity,\allowbreak{}infinity]\allowbreak{}. The source definition gives <E>\_\allowbreak{}1\allowbreak{}=smd(add(A)\allowbreak{})\allowbreak{}. Induction with Section 1 gives

\[
 \langle E\rangle_n=C_n(A)
 =smd(\langle E\rangle_1^{\star n})
 =\bigcup_{m\geq1}C_n(E[-m,m]).
\]

The last equality follows from the finite-witness union argument. In particular every individual membership has a finite shift window,\allowbreak{} even though no common window for all members is claimed. We also have <E>\_\allowbreak{}\{n\allowbreak{}+n'\}\allowbreak{}=smd(<E>\_\allowbreak{}n star <E>\_\allowbreak{}\{n'\})\allowbreak{}. MC-STK-ERR-0892 supplies the missing ``by'' in ``Denote by''. MC-STK-ERR-0893 removes the unmatched bracket in the formula for <E>\_\allowbreak{}1; the balanced formula above is the source\textquotesingle s definition.

The increasing union \textless E\textgreater{} of these levels is closed under finite sums and summands. A shift of a finite witness is again a finite witness with all its indices shifted,\allowbreak{} so it is closed under both shifts. If X is in level r and Y in level s,\allowbreak{} a cone triangle X\allowbreak{}->Y\allowbreak{}->Z\allowbreak{}->X[1]\allowbreak{} puts Z in level s\allowbreak{}+r. Thus this strictly full subcategory is triangulated and saturated.

If D\textquotesingle{} is another strictly full saturated triangulated subcategory containing E,\allowbreak{} it is stable under shifts,\allowbreak{} finite sums,\allowbreak{} actual summands and extensions. Each witness for every level therefore lies in D\textquotesingle. This proves minimality. The correct closure assertion is add(D')\allowbreak{} subset D',\allowbreak{} as in MC-STK-ERR-0894. The former add(D)\allowbreak{} subset D\textquotesingle{} is false:\allowbreak{} take D\allowbreak{}=D(k)\allowbreak{},\allowbreak{} E\allowbreak{}=0,\allowbreak{} and D\textquotesingle{} the strictly full subcategory of zero objects. It contains E but not k. The P02-E0085 explanation concerning a functor landing in D\textquotesingle{} does not describe this proof; the replacement is justified instead by this exact closure argument.

For the right orthogonal,\allowbreak{} assume Hom(E,\allowbreak{}K[i]\allowbreak{})\allowbreak{}\allowbreak{}=0 for every i. Shifting the source and target shows Hom(E[j]\allowbreak{},\allowbreak{}K)\allowbreak{}\allowbreak{}=0 for every j. Finite sums and summands preserve this vanishing. For A\allowbreak{}->X\allowbreak{}->B in an extension witness,\allowbreak{} exactness of Hom(B,\allowbreak{}K)\allowbreak{}\allowbreak{}->Hom(X,\allowbreak{}K)\allowbreak{}\allowbreak{}->Hom(A,\allowbreak{}K)\allowbreak{} proves the same vanishing for X. Induction over the levels proves Hom(X,\allowbreak{}K)\allowbreak{}\allowbreak{}=0 for all X in \textless E\textgreater. Conversely choose X\allowbreak{}=E[-i]\allowbreak{}. If E is classical,\allowbreak{} <E>\allowbreak{}=D,\allowbreak{} so this vanishing includes id\_\allowbreak{}K\allowbreak{}=0. The unique zero maps between K and a zero object are then inverse,\allowbreak{} proving K is zero. Hence classical generators are weak generators.

The source already proves the change-of-generator level bound:\allowbreak{} if E\textquotesingle{} is in <E>\_\allowbreak{}r,\allowbreak{} every shift,\allowbreak{} finite sum and summand of E\textquotesingle{} is in <E>\_\allowbreak{}r,\allowbreak{} so <E'>\_\allowbreak{}1 subset <E>\_\allowbreak{}r and <E'>\_\allowbreak{}n subset <E>\_\allowbreak{}\{rn\}. Thus a classical generator is strong whenever D has a strong generator. This is not counted as a new editorial consequence.

Finally consider the object property T in 11094–\allowbreak{}11108. From T(E)\allowbreak{} and the triangle E --id-\allowbreak{}-> E \allowbreak{}-> 0 one obtains T(0)\allowbreak{}. If X is isomorphic to Y,\allowbreak{} the triangle X\allowbreak{}->Y\allowbreak{}->0 shows T(X)\allowbreak{} implies T(Y)\allowbreak{}. Thus isomorphism invariance follows from the actual assumptions; it need not be added. The triangle X\allowbreak{}->0\allowbreak{}->X[1]\allowbreak{} and its inverse rotation give shift closure. The stated sum,\allowbreak{} triangle and summand assumptions now show that the objects satisfying T form a strictly full saturated triangulated subcategory containing E,\allowbreak{} so minimality proves the remark.

\subsection{6. Sparse and weighted generation --- DERIVED-UNDER-027}\label{proofs-10691-11108-sparse-and-weighted-generation-derived-under-027}

Keep the actual K and its canonical truncation maps. Let p1\textless...\textless pt be its nonzero cohomology degrees,\allowbreak{} a finite nonempty list. Suppose E is a full subcategory of D(A)\allowbreak{} and,\allowbreak{} for each pj,\allowbreak{}

\[
 H^{p_j}(K)\in C_{r_j}(E),\qquad r_j\geq1.
\]

Write a\allowbreak{}=p1 and b\allowbreak{}=pt. Then

\[
 K\in C_{r_1+\cdots+r_t}(E[a,b]).
\]

Indeed,\allowbreak{} shifting each witness gives H\textasciicircum{}\{pj\}(K)\allowbreak{}[-pj]\allowbreak{} in C\_\allowbreak{}\{rj\}(E[a,\allowbreak{}b]\allowbreak{})\allowbreak{}. Between two consecutive nonzero cohomology degrees,\allowbreak{} the canonical truncation maps are quasi-isomorphisms,\allowbreak{} since they induce identity on every retained cohomology object and zero-to-zero elsewhere. Compose these actual maps across each gap. The remaining tower has distinguished triangles

\[
 \tau_{\leq p_{j-1}}K\longrightarrow\tau_{\leq p_j}K
 \longrightarrow H^{p_j}(K)[-p_j]
 \longrightarrow(\tau_{\leq p_{j-1}}K)[1]
\]

for j>\allowbreak{}=2,\allowbreak{} after transport along those quasi-isomorphisms. Its first stage is H\textasciicircum{}\{p1\}(K)\allowbreak{}[-p1]\allowbreak{},\allowbreak{} and its last stage is canonically isomorphic to K. Apply the level-addition identity at every displayed triangle. This proves the stated sum of costs. When all cohomology is zero,\allowbreak{} K is isomorphic to zero and belongs to C\_\allowbreak{}1(E[a,\allowbreak{}b]\allowbreak{})\allowbreak{} for any chosen nonempty interval; no level zero is introduced into the source notation.

There is an equally exact termwise statement. For a given bounded complex with nonzero terms in degrees q1<...<qu,\allowbreak{} assume K\^{}\{qj\} in C\_\allowbreak{}\{sj\}(E)\allowbreak{}. Peel off its actual top term using the degreewise split sequence in Section 3,\allowbreak{} then repeat with the quotient. The resulting extension order is qu,\allowbreak{}...,\allowbreak{}q1,\allowbreak{} and its cost is s1\allowbreak{}+...\allowbreak{}+su in the shift window [q1,\allowbreak{}qu]\allowbreak{}. All intervening differentials are the original differentials in those triangles.

For the source hypothesis rj\allowbreak{}=1,\allowbreak{} this reduces the displayed width b-a\allowbreak{}+1 to the number of nonzero layers. Weighted costs allow different finite constructions for different layers. Neither bound is asserted to be minimal,\allowbreak{} and no splitting of the actual extension data is used. This consequence belongs in editorial material.

\subsection{7. Change of generator with the full shift window --- DERIVED-UNDER-028}\label{proofs-10691-11108-change-of-generator-with-the-full-shift-window-derived-under-028}

Let G,\allowbreak{}E be objects of T,\allowbreak{} let a<\allowbreak{}=b and c<\allowbreak{}=d be finite integers,\allowbreak{} and let r,\allowbreak{}n>\allowbreak{}=1. For the original witnesses satisfying

\[
 E\in C_r(G[c,d]),\qquad X\in C_n(E[a,b]),
\]

one has

\[
 X\in C_{rn}(G[a+c,b+d]).
\]

To prove this,\allowbreak{} fix i in [a,\allowbreak{}b]\allowbreak{}. Shift every triangle,\allowbreak{} finite sum,\allowbreak{} complement and isomorphism in the witness for E by [-i]\allowbreak{}. Its leaves become G[-j]\allowbreak{}[-i]\allowbreak{},\allowbreak{} j in [c,\allowbreak{}d]\allowbreak{}. The specified composition isomorphisms identify these with G[-(i\allowbreak{}+j)\allowbreak{}]\allowbreak{}. Therefore E[-i]\allowbreak{} belongs to C\_\allowbreak{}r(G[c\allowbreak{}+i,\allowbreak{}d\allowbreak{}+i]\allowbreak{})\allowbreak{},\allowbreak{} which is contained in C\_\allowbreak{}r(G[a\allowbreak{}+c,\allowbreak{}b\allowbreak{}+d]\allowbreak{})\allowbreak{}. This receiving category is strictly full,\allowbreak{} closed under finite sums and actual summands,\allowbreak{} so the same statement holds for every object of smd(add(E[a,\allowbreak{}b]\allowbreak{})\allowbreak{})\allowbreak{}.

For each of the n extension factors witnessing X,\allowbreak{} substitute that membership. Apply Section 1 successively to the same extension triangles and the final summand:\allowbreak{} their costs add to rn. Their common window remains [a\allowbreak{}+c,\allowbreak{}b\allowbreak{}+d]\allowbreak{}. It is not multiplied by n or r. Forgetting the finite windows gives the source\textquotesingle s already established unbounded rn rule; retaining the exact windows is the refinement recorded here. It does not change the source definitions.

\subsection{8. The tensor factor in the geometric receiver --- DERIVED-NEW-038}\label{proofs-10691-11108-the-tensor-factor-in-the-geometric-receiver-derived-new-038}

In equiv.tex 1784–\allowbreak{}1842,\allowbreak{} X is proper and smooth over the field k,\allowbreak{} E and G are the fixed finite locally free sheaves chosen from the diagonal construction,\allowbreak{} and

\[
 \mathcal O_\Delta\in C_n((E\boxtimes G)[-m,m]).
\]

For a perfect K,\allowbreak{} the exact functor

\[
 \Phi_K(M)=R{\rm pr}_{2,*}(L{\rm pr}_1^*K
                   \otimes_{\mathcal O_{X\times X}}^{\mathbf L}M)
\]

sends O\_\allowbreak{}Delta to K and E boxtimes G to

\[
 V_K\otimes_k G,\qquad
 V_K=R\Gamma(X,K\otimes_{\mathcal O_X}^{\mathbf L}E).
\]

The first identification is restriction to the diagonal followed by its inverse projection. For the second,\allowbreak{} flat base change from Spec(k)\allowbreak{} identifies Rpr\_\allowbreak{}\{2,\allowbreak{}*\}Lpr\_\allowbreak{}1\textasciicircum{}*(K tensor\^{}L E)\allowbreak{} with RΓ(X,\allowbreak{}K tensor\^{}L E)\allowbreak{} tensor\_\allowbreak{}k O\_\allowbreak{}X,\allowbreak{} and the projection formula tensors it with G. All schemes here are quasi-compact and separated,\allowbreak{} and all complexes have quasi-coherent cohomology; over the field k the base change is flat. The projection formula used is perfect.tex,\allowbreak{} lemma-cohomology-base-change,\allowbreak{} current 5051–\allowbreak{}5094. These are isomorphisms applied inside the strictly full C\_\allowbreak{}n,\allowbreak{} so the repair in Section 2 preserves the entire argument.

The complex K tensor\^{}L E is perfect:\allowbreak{} on each open with a bounded finite locally free representative for K,\allowbreak{} tensor it with the finite locally free E. Properness gives finitely many such opens,\allowbreak{} hence a bounded coherent cohomology complex globally. Applying perfect.tex,\allowbreak{} lemma-direct-image-coherent,\allowbreak{} current 2634–\allowbreak{}2660,\allowbreak{} to X\allowbreak{}->Spec(k)\allowbreak{} proves V\_\allowbreak{}K has finite-dimensional cohomology in a finite interval [a,\allowbreak{}b]\allowbreak{}. Its hypotheses hold:\allowbreak{} Spec(k)\allowbreak{} is Noetherian,\allowbreak{} X is of finite type,\allowbreak{} and every support is closed in proper X,\allowbreak{} hence proper. The bounds must therefore refer to H\textasciicircum{}i(X,\allowbreak{}K tensor\^{}L E)\allowbreak{},\allowbreak{} not H\textasciicircum{}i(X,\allowbreak{}K)\allowbreak{}.

Here is the precise use of the field. For any complex V of k-vector spaces write Z\textasciicircum{}q\allowbreak{}=ker d\^{}q and B\textasciicircum{}q\allowbreak{}=im d\^{}\{q-1\}. Choose vector-space complements Z\textasciicircum{}q\allowbreak{}=B\textasciicircum{}q\allowbreak{}+T\textasciicircum{}q and V\textasciicircum{}q\allowbreak{}=Z\textasciicircum{}q\allowbreak{}+S\textasciicircum{}q. The map d\^{}q restricts to an isomorphism S\textasciicircum{}q\allowbreak{}->B\textasciicircum{}\{q\allowbreak{}+1\}; T\^{}q maps isomorphically to H\textasciicircum{}q(V)\allowbreak{}. Thus V is the direct sum of its zero-differential T\textasciicircum{}q[-q]\allowbreak{} pieces and the two-term disks S\^{}q --d\textasciicircum{}q-\allowbreak{}-> B\textasciicircum{}\{q\allowbreak{}+1\}. Every degree has just the stated finite number of types of pieces,\allowbreak{} so this decomposition is well-defined even when V itself is unbounded. On each disk the inverse of d\^{}q is a contraction. The projection onto the T pieces and their inclusions give an explicit homotopy equivalence,\allowbreak{} with the contractions accounting for every discarded disk.

Apply this to V\_\allowbreak{}K. Its nonzero T pieces are finite-dimensional and confined to [a,\allowbreak{}b]\allowbreak{},\allowbreak{} so V\_\allowbreak{}K tensor\_\allowbreak{}k G is isomorphic to the finite sum of H\textasciicircum{}q(V\_\allowbreak{}K)\allowbreak{} tensor\_\allowbreak{}k G[-q]\allowbreak{},\allowbreak{} a member of add(G[a,\allowbreak{}b]\allowbreak{})\allowbreak{}. Tensoring preserves the displayed contraction. The same shift calculation as Section 7 gives

\[
 C_n((V_K\otimes_k G)[-m,m])
       \subset C_n(G[a-m,b+m]).
\]

This proves the corrected inference at 1836–\allowbreak{}1840 with its original n,\allowbreak{}m,\allowbreak{}K,\allowbreak{}E,\allowbreak{}G,\allowbreak{} and completes the strong-generator argument.

For the next lemma,\allowbreak{} replace K by a coherent sheaf F. Smooth X is regular of finite dimension,\allowbreak{} so F is perfect locally; moreover F tensor\^{}L E\allowbreak{}=F tensor E because E is locally free. This is a coherent sheaf in degree zero. The bounds a\allowbreak{}=0,\allowbreak{}b\allowbreak{}=dim(X)\allowbreak{} apply to its cohomology,\allowbreak{} by the Noetherian dimension vanishing theorem (cohomology.tex,\allowbreak{} proposition-vanishing-Noetherian,\allowbreak{} current statement 4083–\allowbreak{}4086)\allowbreak{}. Thus the source\textquotesingle s m\allowbreak{}=max(m',\allowbreak{}m'\allowbreak{}+dim(X)\allowbreak{})\allowbreak{} still works. The corrected arbitrary-window assertion at 1862 must name H\textasciicircum{}i(X,\allowbreak{}F tensor E)\allowbreak{},\allowbreak{} with E the fixed sheaf from the preceding proof. Neither lemma\textquotesingle s final conclusion or dimension constant needs weakening.

The missing tensor factor is demonstrably substantive. On X\allowbreak{}=P\textasciicircum{}1\_\allowbreak{}k use homogeneous coordinates [X0:\allowbreak{}X1]\allowbreak{},\allowbreak{} t\allowbreak{}=X1/\allowbreak{}X0 on U0,\allowbreak{} and s\allowbreak{}=X0/\allowbreak{}X1\allowbreak{}=t\textasciicircum{}\{-1\} on U1. The frames X0\^{}e and X1\^{}e of O(e)\allowbreak{} obey X1\textasciicircum{}e\allowbreak{}=t\textasciicircum{}e X0\^{}e. The two-affine-cover Cech complex is

\[
 k[t]\oplus k[t^{-1}]
 \xrightarrow{(f,g)\mapsto t^e g-f}
 k[t,t^{-1}].
\]

All opens and their intersection are affine,\allowbreak{} so quasi-coherent affine acyclicity identifies its kernel and cokernel with H\^{}0 and H\textasciicircum{}1,\allowbreak{} with higher groups zero. For e\allowbreak{}=-1 the two sets of exponents are >\allowbreak{}=0 and <\allowbreak{}=-1. They are disjoint and their union is Z,\allowbreak{} so both kernel and cokernel vanish. For e\allowbreak{}=-2 their union omits exactly -1; the kernel is zero and the cokernel is k times t\^{}\{-1\}. Therefore

\[
 H^*(X,O(-1))=0,\qquad H^1(X,O(-2))=k.
\]

Take K\allowbreak{}=O(-1)\allowbreak{} and E\allowbreak{}=G\allowbreak{}=O\allowbreak{}+O(-1)\allowbreak{}. These E,\allowbreak{}G meet the actual diagonal condition:\allowbreak{} the determinant X0Y1-X1Y0 defines the diagonal and gives the exact resolution

\[
 0\to O(-1,-1)\xrightarrow{X_0Y_1-X_1Y_0}
 O\to O_\Delta\to0.
\]

On the coordinate charts the diagonal equation is the coordinate difference,\allowbreak{} or its invertible-coordinate counterpart,\allowbreak{} a non-zero-divisor; its quotient is precisely the diagonal chart. This proves injectivity,\allowbreak{} identifies the quotient and establishes the sequence globally. Both vector bundles are summands of E boxtimes G. The associated triangle O \allowbreak{}-> O\_\allowbreak{}Delta \allowbreak{}-> O(-1,\allowbreak{}-1)\allowbreak{}[1]\allowbreak{} puts O\_\allowbreak{}Delta in C\_\allowbreak{}2((E boxtimes G)\allowbreak{}[-1,\allowbreak{}1]\allowbreak{})\allowbreak{}.

Nevertheless K tensor E\allowbreak{}=O(-1)\allowbreak{}\allowbreak{}+O(-2)\allowbreak{} has nonzero H\textasciicircum{}1,\allowbreak{} while K has zero global cohomology in every degree. For example [0,\allowbreak{}0]\allowbreak{} bounds H\textasciicircum{}*(X,\allowbreak{}K)\allowbreak{} but does not bound V\_\allowbreak{}K. The stated inference from that range is therefore invalid.

The unqualified ``for any a<\allowbreak{}=b” claim in the next proof also fails. For any fixed m',\allowbreak{} choose F\allowbreak{}=O(-1)\allowbreak{} and a\allowbreak{}=b\allowbreak{}=m'\allowbreak{}+1. Its global cohomology is zero and satisfies that printed condition. Yet G is a sheaf in degree zero and the alleged generating window becomes [1,\allowbreak{}2m'\allowbreak{}+1]\allowbreak{}. Section 4 shows that every object of C\_\allowbreak{}n(G[1,\allowbreak{}2m'\allowbreak{}+1]\allowbreak{})\allowbreak{} has zero cohomology sheaf in degree zero. F has nonzero degree-zero cohomology sheaf O(-1)\allowbreak{},\allowbreak{} a contradiction. Using F tensor E in the premise excludes this false range and proves the intended assertion.

These two loci are recorded as one propagated missing-factor finding,\allowbreak{} with two exact operations. The 101 equiv.tex report routing records have no unlocated record and none overlaps 1836–\allowbreak{}1838 or 1862–\allowbreak{}1863. This bounded inventory check is not a claim to have reviewed the remaining Equivalences reports.

\subsection{\texorpdfstring{9. Propagation,\allowbreak{} scope and remaining work}{9.  scope and remaining work}}\label{proofs-10691-11108-propagation-scope-and-remaining-work}

The downstream boundedness lemma equiv.tex 1871–\allowbreak{}1888 remains valid. Choose the uniform m,\allowbreak{}n,\allowbreak{}G from the repaired diagonal lemma. For its homological functor choose the finite [a,\allowbreak{}b]\allowbreak{} supporting H\textasciicircum{}*(G)\allowbreak{},\allowbreak{} as allowed by that lemma\textquotesingle s assumptions. Section 4 gives the uniform range [a-m,\allowbreak{}b\allowbreak{}+m]\allowbreak{} for every coherent F. Enlarge it to [-M,\allowbreak{}M]\allowbreak{} with M\allowbreak{}=max(1,\allowbreak{}m-a,\allowbreak{}m\allowbreak{}+b)\allowbreak{}. This explicit M does not depend on the extension count n. For a cohomological functor apply the same reasoning with its reversed exact Hom sequence and its corresponding shift indices; a symmetric interval is unchanged by reversing signs.

The call to the classical-to-strong-generator lemma in more-algebra.tex 21561–\allowbreak{}21578 is preserved:\allowbreak{} if the cited earlier ring-specific lemmas give a classical generator R and any strong generator,\allowbreak{} Section 5 supplies the precise finite rn comparison. Only this categorical implication is checked here; this pass does not certify the unreviewed ring-specific premises.

All eleven physical reports in this span concern seven previously corrected units 0888–\allowbreak{}0894. P02-E0086 reports the unbound base-case index,\allowbreak{} although its range also touches the different truncation repair; its disposition is joined to 0889 by its actual text. P02-E0085 has the correct replacement but the inaccurate functor rationale discussed in Section 5. Neither range overlap nor an earlier admission is substituted for these semantic checks.

New source operations are limited to the shift-isomorphism qualification and the two missing-tensor-factor replacements. The sparse costs and bounded-window generator transfer remain separate editorial results with complete derivations above. This pass completes semantic review through 11108 and leaves Compact objects,\allowbreak{} starting 11121,\allowbreak{} next. Reading through 11170 is only read-ahead. No whole-chapter Equivalences or More on Algebra review,\allowbreak{} live-source composition,\allowbreak{} build or publication is claimed.
\clearpage
\section{Compact objects and Brown representability}\label{proofs-11121-11450-compact-objects-and-brown-representability}

Private source review,\allowbreak{} 2026-09-30. Authority is the frozen Stacks Project commit a04446e5\allowbreak{}7ec1fbc2\allowbreak{}52a871af\allowbreak{}cec7752f\allowbreak{}b2807b14, derived.tex 11121–\allowbreak{}11450. The inherited five chapter preimages are unchanged. This note reconciles the received reports,\allowbreak{} proves the exact maps behind two omitted shifts,\allowbreak{} and checks the geometric receiver in perfect.tex. The finite-stage strengthening in Section 3 is editorial material. Its applications below are explicit; no novelty is claimed.

\subsection{1. Closure of compact objects}\label{proofs-11121-11450-closure-of-compact-objects}

For an object P and a family (M\_\allowbreak{}j)\allowbreak{},\allowbreak{} denote the canonical map by

\[
 \theta_{P,(M_j)}:\bigoplus_j\operatorname{Hom}(P,M_j)
       \longrightarrow\operatorname{Hom}(P,\bigoplus_jM_j).
\]

Its restriction to the j-th summand composes with the actual coproduct inclusion. Compactness is bijectivity of this map for every family. Shifts preserve coproducts because they are equivalences; the shift isomorphisms identify theta for P[r]\allowbreak{} and (M\_\allowbreak{}j)\allowbreak{} with theta for P and (M\_\allowbreak{}j[-r]\allowbreak{})\allowbreak{}. Thus shifts preserve compactness. Isomorphisms do also,\allowbreak{} by precomposition.

For a distinguished triangle X\allowbreak{}->Y\allowbreak{}->Z\allowbreak{}->X[1]\allowbreak{},\allowbreak{} apply Hom(-,\allowbreak{}M\_\allowbreak{}j)\allowbreak{} and take the direct sum of its exact sequences of abelian groups. Compare with Hom(-,\allowbreak{}direct sum M\_\allowbreak{}j)\allowbreak{}. The five adjacent terms

\[
 \operatorname{Hom}(Y[1],-)\longrightarrow
 \operatorname{Hom}(X[1],-)\longrightarrow
 \operatorname{Hom}(Z,-)\longrightarrow
 \operatorname{Hom}(Y,-)\longrightarrow
 \operatorname{Hom}(X,-)
\]

are exact with the triangle\textquotesingle s given connecting signs. If X and Y are compact,\allowbreak{} the four comparison maps surrounding the Z term are isomorphisms. The five lemma makes the Z comparison an isomorphism. Rotations handle every other pair of compact terms. Finite sums preserve compactness,\allowbreak{} since Hom out of a finite biproduct is a finite product and finite products of abelian groups commute with direct sums.

If P is a retract of a compact Q,\allowbreak{} choose i:\allowbreak{}P\allowbreak{}->Q and r:\allowbreak{}Q\allowbreak{}->P with ri\allowbreak{}=id\_\allowbreak{}P. Precomposition with r and i makes the comparison map for P a retract,\allowbreak{} in the category of homomorphisms,\allowbreak{} of that for Q. Its inverse is the corresponding composite of the inverse for Q with these inclusion and retraction maps; the equations ri\allowbreak{}=id\_\allowbreak{}P prove both inverse identities. Thus P is compact.

The previously proved original lemma-projectors-have-images-triangulated,\allowbreak{} 805–\allowbreak{}822,\allowbreak{} makes the ambient pre-triangulated category Karoubian when it has countable coproducts:\allowbreak{} split monomorphisms have cokernels by their split triangles,\allowbreak{} and the countable-coproduct idempotent splitting criterion applies. Its actual summands are compact by the preceding retract calculation. This proves every assertion at 11146–\allowbreak{}11165,\allowbreak{} including the pre-triangulated version,\allowbreak{} without using TR4 in the compactness comparison.

\subsection{2. The cellular tower and the missing inverse shift --- DERIVED-NEW-039}\label{proofs-11121-11450-the-cellular-tower-and-the-missing-inverse-shift-derived-new-039}

Let (E\_\allowbreak{}i)\allowbreak{} be the source\textquotesingle s compact family and assume their coproduct is a weak generator. For a fixed X put

\[
 X_1=\bigoplus_{(i,m,\varphi:E_i[m]\to X)}E_i[m],
 \qquad p_1|_{(i,m,\varphi)}=\varphi.
\]

This is a set-indexed coproduct:\allowbreak{} I and Z are sets and each Hom collection is a set. Given p\_\allowbreak{}n:\allowbreak{}X\_\allowbreak{}n\allowbreak{}->X let Y\_\allowbreak{}n be the coproduct indexed by maps phi:\allowbreak{}E\_\allowbreak{}i[m]\allowbreak{}\allowbreak{}->X\_\allowbreak{}n with p\_\allowbreak{}n phi\allowbreak{}=0. Its component map is v\_\allowbreak{}n:\allowbreak{}Y\_\allowbreak{}n\allowbreak{}->X\_\allowbreak{}n. Complete it to the actual triangle

\[
 Y_n\xrightarrow{v_n}X_n\xrightarrow{u_n}X_{n+1}
          \xrightarrow{d_n}Y_n[1].
\]

The equality p\_\allowbreak{}n v\_\allowbreak{}n\allowbreak{}=0 follows on every coproduct component. Exactness of Hom(-,\allowbreak{}X)\allowbreak{} supplies p\_\allowbreak{}\{n\allowbreak{}+1\} with p\_\allowbreak{}\{n\allowbreak{}+1\}u\_\allowbreak{}n\allowbreak{}=p\_\allowbreak{}n. There is no uniqueness requirement on this choice.

Put S\allowbreak{}=direct sum\_\allowbreak{}n X\_\allowbreak{}n and t i\_\allowbreak{}n\allowbreak{}=i\_\allowbreak{}n-i\_\allowbreak{}\{n\allowbreak{}+1\}u\_\allowbreak{}n,\allowbreak{} with i\_\allowbreak{}n the coproduct inclusions. Choose the telescope triangle

\[
 S\xrightarrow t S\xrightarrow q H\xrightarrow w S[1].
\]

The component maps p\_\allowbreak{}n define p:\allowbreak{}S\allowbreak{}->X with pt\allowbreak{}=0,\allowbreak{} so exactness gives f:\allowbreak{}H\allowbreak{}->X satisfying fq\allowbreak{}=p. Set q\_\allowbreak{}n\allowbreak{}=q i\_\allowbreak{}n. Complete f to C --a-\allowbreak{}-> H --f-\allowbreak{}-> X --b-\allowbreak{}-> C[1]\allowbreak{}. For any P\allowbreak{}=E\_\allowbreak{}i[m]\allowbreak{},\allowbreak{} compactness and the previously proved telescope comparison give

\[
 \operatorname{Hom}(P,H)=
 \mathop{\rm colim}_n\operatorname{Hom}(P,X_n).
\]

If z:\allowbreak{}P\allowbreak{}->C,\allowbreak{} write az\allowbreak{}=q\_\allowbreak{}n h for some h:\allowbreak{}P\allowbreak{}->X\_\allowbreak{}n. Since faz\allowbreak{}=0 and fq\_\allowbreak{}n\allowbreak{}=p\_\allowbreak{}n,\allowbreak{} we have p\_\allowbreak{}n h\allowbreak{}=0. The summand labelled h in Y\_\allowbreak{}n shows u\_\allowbreak{}n h\allowbreak{}=0. Therefore az\allowbreak{}=q\_\allowbreak{}\{n\allowbreak{}+1\}u\_\allowbreak{}n h\allowbreak{}=0. The inverse-rotated triangle is

\[
 X[-1]\xrightarrow{-b[-1]}C\xrightarrow a H\xrightarrow f X.
\]

Exactness supplies y:\allowbreak{}P\allowbreak{}->X[-1]\allowbreak{} with z\allowbreak{}=(-b[-1]\allowbreak{})\allowbreak{}y. Shift y to y[1]\allowbreak{}:\allowbreak{}P[1]\allowbreak{}\allowbreak{}->X. Since P[1]\allowbreak{} is a shift of E\_\allowbreak{}i,\allowbreak{} it labels a summand of X\_\allowbreak{}1. Write e:\allowbreak{}P[1]\allowbreak{}\allowbreak{}->X\_\allowbreak{}1 for its inclusion. Then p\_\allowbreak{}1 e\allowbreak{}=y[1]\allowbreak{}. The shifted map

\[
 P\xrightarrow{e[-1]}X_1[-1]\xrightarrow{q_1[-1]}H[-1]
\]

composed with f[-1]\allowbreak{} is y,\allowbreak{} using the specified shift identifications. Since (-b[-1]\allowbreak{})\allowbreak{}f[-1]\allowbreak{}\allowbreak{}=0,\allowbreak{} z\allowbreak{}=0. Hence Hom(E\_\allowbreak{}i[m]\allowbreak{},\allowbreak{}C)\allowbreak{}\allowbreak{}=0 for every i,\allowbreak{}m. The coproduct universal property gives Hom(direct sum\_\allowbreak{}i E\_\allowbreak{}i,\allowbreak{}C[-m]\allowbreak{})\allowbreak{}\allowbreak{}=0 for every m,\allowbreak{} so generation forces C to be zero and f to be an isomorphism.

The map just constructed lands in H[-1]\allowbreak{},\allowbreak{} not H. The wording at 11209 omits [-1]\allowbreak{} from the displayed homotopy colimit. The corrected object is (hocolim X\_\allowbreak{}n)\allowbreak{}[-1]\allowbreak{}. This is a type correction to the proof,\allowbreak{} not a change in the conclusion or a choice of an unspecified map H\allowbreak{}->X[-1]\allowbreak{}. All choices of lifts p\_\allowbreak{}\{n\allowbreak{}+1\} and f satisfy the argument.

\subsection{3. Exact finite-stage factorization and its level --- DERIVED-UNDER-029}\label{proofs-11121-11450-exact-finite-stage-factorization-and-its-level-derived-under-029}

Here a stronger hypothesis accounting is useful. Let (E\_\allowbreak{}i)\allowbreak{} be any family of compact objects in a triangulated category with coproducts. A weak-generation assumption is not needed in this section. Suppose X\_\allowbreak{}1 is a coproduct of shifts of this family and there are triangles

\[
 Y_{n-1}\xrightarrow u X_{n-1}\xrightarrow v X_n
        \xrightarrow d Y_{n-1}[1]
\]

with each Y\_\allowbreak{}\{n-1\} a coproduct of such shifts. For any compact C and map h:\allowbreak{}C\allowbreak{}->X\_\allowbreak{}n there exist a finite J subset I,\allowbreak{} an object Q in \textless direct sum\_\allowbreak{}\{j in J\} E\_\allowbreak{}j>\_\allowbreak{}n and actual maps

\[
 C\xrightarrow s Q\xrightarrow r X_n,\qquad rs=h.
\]

The subscript n is a proved extension-cost bound,\allowbreak{} not a claim of minimality. It strengthens the source\textquotesingle s unbounded membership in \textless direct sum E\_\allowbreak{}j>. Generation is used in Section 2 to obtain the isomorphism from the telescope,\allowbreak{} but not for this finite-stage factorization.

For n\allowbreak{}=1 compactness factors h through a finite subcoproduct Q. Its objects are shifts of a finite family,\allowbreak{} so Q is in level one. Suppose the statement holds for n-1. Compactness factors dh through a finite subcoproduct of Y\_\allowbreak{}\{n-1\}[1]\allowbreak{}. Write it as

\[
 C\xrightarrow\delta E'[1]\xrightarrow{j[1]}Y_{n-1}[1],
 \qquad j[1]\delta=dh,
\]

where E\textquotesingle{} is a finite direct sum of shifts of the E\_\allowbreak{}i with indices I\textquotesingle. Complete delta to a triangle E\textquotesingle{} --a-\allowbreak{}-> C\textquotesingle{} --b-\allowbreak{}-> C --delta-\allowbreak{}-> E'[1]\allowbreak{}. TR3 applied to the rotated triangles gives a map t:\allowbreak{}C'\allowbreak{}->X\_\allowbreak{}\{n-1\} such that (j,\allowbreak{}t,\allowbreak{}h)\allowbreak{} is a morphism to the preceding displayed triangle:\allowbreak{} ta\allowbreak{}=uj,\allowbreak{} v t\allowbreak{}=h b,\allowbreak{} and j[1]\allowbreak{}delta\allowbreak{}=d h. In particular C\textquotesingle{} is compact by Section 1.

Induction factors t as C\textquotesingle{} --alpha-\allowbreak{}-> E\textquotesingle\textquotesingle{} --beta-\allowbreak{}-> X\_\allowbreak{}\{n-1\} with E\textquotesingle\textquotesingle{} in level n-1 for a finite family I\textquotesingle\textquotesingle. Complete alpha a to

\[
 E'\xrightarrow{\alpha a}E''\xrightarrow l E
                     \xrightarrow e E'[1].
\]

TR3 gives p:\allowbreak{}C\allowbreak{}->E and q:\allowbreak{}E\allowbreak{}->X\_\allowbreak{}n making (id,\allowbreak{}alpha,\allowbreak{}p)\allowbreak{} and (j,\allowbreak{}beta,\allowbreak{}q)\allowbreak{} morphisms of triangles. Thus e p\allowbreak{}=delta and d q\allowbreak{}=j[1]\allowbreak{}e. It follows exactly that

\[
 d(h-qp)=dh-j[1]ep=dh-j[1]\delta=0.
\]

The common target in this comparison is Y\_\allowbreak{}\{n-1\}[1]\allowbreak{}. The source at 11255 prints Y\_\allowbreak{}\{n-1\},\allowbreak{} missing [1]\allowbreak{}. This is DERIVED-NEW-040. Exactness of Hom(C,\allowbreak{}-)\allowbreak{} now provides z:\allowbreak{}C\allowbreak{}->X\_\allowbreak{}\{n-1\} with vz\allowbreak{}=h-qp. Inductively factor z as C --gamma-\allowbreak{}-> E\textquotesingle\textquotesingle\textquotesingle{} --eta-\allowbreak{}-> X\_\allowbreak{}\{n-1\},\allowbreak{} using a finite family I\textquotesingle\textquotesingle\textquotesingle{} and level n-1. The complete factorization is

\[
 C\xrightarrow{\binom p\gamma}E\oplus E'''
       \xrightarrow{(q,\ v\eta)}X_n,
 \qquad (q,\ v\eta)\binom p\gamma=qp+v\eta\gamma=h.
\]

It preserves the original h,\allowbreak{} not merely its image under d. The objects E\textquotesingle\textquotesingle{} and E\textquotesingle\textquotesingle\textquotesingle{} are different induction outputs. This verifies the already applied MC-STK-ERR-0895 replacement of E\textquotesingle\textquotesingle{} by E\textquotesingle\textquotesingle\textquotesingle{} at 11258 and MC-STK-ERR-0896 replacement of I\textquotesingle{} by I\textquotesingle\textquotesingle\textquotesingle{} at 11259. In particular I\textquotesingle\textquotesingle\textquotesingle{} is a finite index set,\allowbreak{} not an ``injective object'' as the P02-E0089 explanation calls it.

To prove the cost bound,\allowbreak{} put J\allowbreak{}=I' union I\textquotesingle\textquotesingle{} union I\textquotesingle\textquotesingle\textquotesingle{} and G\_\allowbreak{}J\allowbreak{}=direct sum\_\allowbreak{}\{j in J\}E\_\allowbreak{}j. Enlarging a finite family preserves membership in each level:\allowbreak{} the smaller direct sum is a specified summand of G\_\allowbreak{}J,\allowbreak{} so replace its leaf witnesses by those summands. We have E\textquotesingle{} in level one,\allowbreak{} E'',\allowbreak{}E''' in level n-1. Rotation of the E triangle gives E''\allowbreak{}->E\allowbreak{}->E'[1]\allowbreak{},\allowbreak{} hence E is in level n by the level-addition formula proved in the preceding checkpoint. Levels are increasing and closed under finite sums,\allowbreak{} so E\allowbreak{}+E''' is in level n as required. If J is empty,\allowbreak{} all these objects are zero and the empty finite sum is the generator 0. Every membership also has a finite shift window by the finite witness argument of PROOFS\_\allowbreak{}10691\_\allowbreak{}11108.md,\allowbreak{} Section 5.

\subsection{4. Compact generators and the family conclusion}\label{proofs-11121-11450-compact-generators-and-the-family-conclusion}

If compact E classically generates D\_\allowbreak{}c and D is compactly generated by a family (F\_\allowbreak{}i)\allowbreak{},\allowbreak{} then Hom(E,\allowbreak{}K[r]\allowbreak{})\allowbreak{}\allowbreak{}=0 for all r implies Hom(F\_\allowbreak{}i,\allowbreak{}K[r]\allowbreak{})\allowbreak{}\allowbreak{}=0 for all i,\allowbreak{}r. Indeed F\_\allowbreak{}i belongs to <E>,\allowbreak{} and the right-orthogonal calculation already proved at original 11033–\allowbreak{}11050 applies to every shift of K. Generation by direct sum F\_\allowbreak{}i forces K\allowbreak{}=0. This proves (1)\allowbreak{} implies (2)\allowbreak{}.

Conversely a compact weak generator E gives the source\textquotesingle s one-element compact generating set. For compact X use the tower of Section 2 and its isomorphism f:\allowbreak{}H\allowbreak{}->X. Compactness factors f\textasciicircum{}\{-1\}:\allowbreak{}X\allowbreak{}->H through a stage:\allowbreak{} f\textasciicircum{}\{-1\}\allowbreak{}=q\_\allowbreak{}n z. Put r\allowbreak{}=f q\_\allowbreak{}n. Then rz\allowbreak{}=id\_\allowbreak{}X. Apply Section 3 to z to factor it through a finite object Q in <E>\_\allowbreak{}n. Thus X is an actual retract of Q and hence is in <E>\_\allowbreak{}n,\allowbreak{} which is closed under summands. The complement exists in the Karoubian ambient category; its compactness also follows from Section 1. The reverse inclusion \textless E\textgreater{} subset D\_\allowbreak{}c follows from compact closure. Therefore D\_\allowbreak{}c\allowbreak{}=<E>,\allowbreak{} proving the proposition.

Exactly the same argument for a family yields

\[
 D_c=\bigcup_{\substack{J\subset I\ {\rm finite}\\ n\geq1}}
          \left\langle\bigoplus_{j\in J}E_j\right\rangle_n .
\]

The actual retract uses f,\allowbreak{}q\_\allowbreak{}n,\allowbreak{}z and the Section 3 factorization. This is recorded with UNDER-029 as its family consequence,\allowbreak{} not as an independent novelty claim. MC-STK-ERR-0897 correctly changes ``the lemma is proven'' at 11309 to ``the proposition is proven''. Its earlier registry description gives an inaccurate label; the actual source label is proposition-generator-versus-classical-generator.

\subsection{5. Brown\textquotesingle s representing object and its exact transformation}\label{proofs-11121-11450-browns-representing-object-and-its-exact-transformation}

Now let H:\allowbreak{}D\textasciicircum{}op\allowbreak{}->Ab be cohomological and take arbitrary coproducts to products,\allowbreak{} with D compactly generated. Enlarge the set of compact generators by all integer shifts,\allowbreak{} retaining the actual objects and their shift isomorphisms; use that indexed set below. For each pair (i,\allowbreak{}a)\allowbreak{} with a in H(E\_\allowbreak{}i)\allowbreak{},\allowbreak{} take a copy of E\_\allowbreak{}i:\allowbreak{}

\[
 X_1=\bigoplus_{(i,a)}E_i,\qquad
 a_1\in H(X_1)\cong\prod_{(i,a)}H(E_i),\quad
 a_1=(a)_{(i,a)}.
\]

The product comparison is the one induced by all inclusions. The associated Yoneda transformation is explicitly theta\_\allowbreak{}\{1,\allowbreak{}W\}(f)\allowbreak{}\allowbreak{}=H(f)\allowbreak{}(a\_\allowbreak{}1)\allowbreak{} for f:\allowbreak{}W\allowbreak{}->X\_\allowbreak{}1.

Suppose (X\_\allowbreak{}n,\allowbreak{}a\_\allowbreak{}n)\allowbreak{} has been chosen. For all pairs (i,\allowbreak{}k)\allowbreak{} with k:\allowbreak{}E\_\allowbreak{}i\allowbreak{}->X\_\allowbreak{}n and H(k)\allowbreak{}a\_\allowbreak{}n\allowbreak{}=0,\allowbreak{} let K\_\allowbreak{}\{n\allowbreak{}+1\} be their coproduct and g\_\allowbreak{}n:\allowbreak{}K\_\allowbreak{}\{n\allowbreak{}+1\}\allowbreak{}->X\_\allowbreak{}n their component map. The product comparison gives H(g\_\allowbreak{}n)\allowbreak{}a\_\allowbreak{}n\allowbreak{}=0 because each coordinate is zero. Complete g\_\allowbreak{}n to

\[
 K_{n+1}\xrightarrow{g_n}X_n\xrightarrow{u_n}X_{n+1}
                       \longrightarrow K_{n+1}[1].
\]

Cohomological exactness gives a\_\allowbreak{}\{n\allowbreak{}+1\} with H(u\_\allowbreak{}n)\allowbreak{}a\_\allowbreak{}\{n\allowbreak{}+1\}\allowbreak{}=a\_\allowbreak{}n. This justifies the source\textquotesingle s use of the kernel functor without assuming that this kernel is cohomological. Its needed product property follows because kernels commute with products,\allowbreak{} or directly from the coordinate calculation just given.

Form S\allowbreak{}=direct sum X\_\allowbreak{}n and the telescope triangle S --t-\allowbreak{}-> S --q-\allowbreak{}-> X --w-\allowbreak{}-> S[1]\allowbreak{},\allowbreak{} with t i\_\allowbreak{}n\allowbreak{}=i\_\allowbreak{}n-i\_\allowbreak{}\{n\allowbreak{}+1\}u\_\allowbreak{}n. Under H(S)\allowbreak{}\allowbreak{}=product H(X\_\allowbreak{}n)\allowbreak{},\allowbreak{} H(t)\allowbreak{} is the homomorphism

\[
 (b_n)_n\longmapsto
           (b_n-H(u_n)b_{n+1})_n .
\]

The tuple (a\_\allowbreak{}n)\allowbreak{} is in its kernel. Exactness supplies a in H(X)\allowbreak{} with H(q)\allowbreak{}a\allowbreak{}=(a\_\allowbreak{}n)\allowbreak{}. Put theta\_\allowbreak{}W(f)\allowbreak{}\allowbreak{}=H(f)\allowbreak{}a for f:\allowbreak{}W\allowbreak{}->X. Naturality follows from H(fg)\allowbreak{}\allowbreak{}=H(g)\allowbreak{}H(f)\allowbreak{},\allowbreak{} and theta\_\allowbreak{}W(q\_\allowbreak{}n f)\allowbreak{}\allowbreak{}=H(f)\allowbreak{}a\_\allowbreak{}n. This is the exact transformation asserted at 11385–\allowbreak{}11387. The requested copyedit inserts the finite verb in ``Hence there is a natural transformation''.

For each shifted generator E\_\allowbreak{}i,\allowbreak{} compactness identifies Hom(E\_\allowbreak{}i,\allowbreak{}X)\allowbreak{} with colim Hom(E\_\allowbreak{}i,\allowbreak{}X\_\allowbreak{}n)\allowbreak{}. The transformation is surjective there because the pair (i,\allowbreak{}b)\allowbreak{} for every b in H(E\_\allowbreak{}i)\allowbreak{} occurs in X\_\allowbreak{}1. If the image of a class represented by f:\allowbreak{}E\_\allowbreak{}i\allowbreak{}->X\_\allowbreak{}n is zero,\allowbreak{} H(f)\allowbreak{}a\_\allowbreak{}n\allowbreak{}=0. It is a summand label of K\_\allowbreak{}\{n\allowbreak{}+1\},\allowbreak{} so u\_\allowbreak{}n f\allowbreak{}=0. Its colimit class is zero. Hence theta is also injective on every shifted generator.

For completeness,\allowbreak{} let D\textquotesingle{} consist of Y for which theta\_\allowbreak{}\{Y[r]\allowbreak{}\} is an isomorphism for every integer r. Isomorphism invariance and shift closure are immediate from this definition. Two objects of a distinguished triangle in D\textquotesingle{} force the third into D\textquotesingle{} by the five adjacent terms of the contravariant exact sequences and the five lemma. A summand is handled by the same explicit retraction of comparison maps as in Section 1. Both Hom(-,\allowbreak{}X)\allowbreak{} and H send coproducts to products; a product of the isomorphisms theta\_\allowbreak{}\{Y\_\allowbreak{}j[r]\allowbreak{}\} is an isomorphism. Thus D\textquotesingle{} is closed under coproducts. In particular the telescope triangle of any sequence of its objects has all three terms in D\textquotesingle.

Every generator is in D',\allowbreak{} hence all X\_\allowbreak{}1,\allowbreak{}Y\_\allowbreak{}n and X\_\allowbreak{}n in the cellular presentation of every object are in D',\allowbreak{} by induction. Section 2 gives every such object as that telescope. Hence D'\allowbreak{}=D and theta is an isomorphism on every object,\allowbreak{} as required. No assertion that the original generators classically generate all of D is made.

Different choices of representing data give canonically isomorphic representing pairs:\allowbreak{} if theta:\allowbreak{}Hom(-,\allowbreak{}X)\allowbreak{}\allowbreak{}->H and theta':\allowbreak{}Hom(-,\allowbreak{}X')\allowbreak{}\allowbreak{}->H are the isomorphisms,\allowbreak{} Yoneda turns (theta')\allowbreak{}\textasciicircum{}\{-1\}theta into a unique u:\allowbreak{}X\allowbreak{}->X'. Its inverse comes from theta\textasciicircum{}\{-1\}theta',\allowbreak{} so both composites are identities. This accounts for the nonunique choices made in the construction.

\subsection{6. The exact right adjoint}\label{proofs-11121-11450-the-exact-right-adjoint}

For an exact coproduct-preserving F:\allowbreak{}D\allowbreak{}->D',\allowbreak{} fix Y in D\textquotesingle{} and apply Section 5 to H\_\allowbreak{}Y(W)\allowbreak{}\allowbreak{}=Hom\_\allowbreak{}\{D'\}(F(W)\allowbreak{},\allowbreak{}Y)\allowbreak{}. Exactness of F and the contravariant representable Hom sequence make this cohomological. The coproduct universal property and preservation by F make H\_\allowbreak{}Y(direct sum W\_\allowbreak{}j)\allowbreak{}\allowbreak{}=product H\_\allowbreak{}Y(W\_\allowbreak{}j)\allowbreak{}.

Choose a representing object G(Y)\allowbreak{} and the natural bijections

\[
 \phi_{W,Y}:\operatorname{Hom}_D(W,G(Y))
        \xrightarrow{\sim}\operatorname{Hom}_{D'}(F(W),Y).
\]

For a:\allowbreak{}Y\allowbreak{}->Y',\allowbreak{} postcomposition by a and these bijections give a natural transformation Hom(-,\allowbreak{}G(Y)\allowbreak{})\allowbreak{}\allowbreak{}->Hom(-,\allowbreak{}G(Y')\allowbreak{})\allowbreak{}. Yoneda defines G(a)\allowbreak{}. The identities and composites are preserved because their natural transformations are,\allowbreak{} proving that G is a functor and that phi is natural in both variables. This proves the actual adjunction,\allowbreak{} rather than objectwise representability alone.

The shift comparison and preservation of triangles are proved in PROOFS\_\allowbreak{}1844\_\allowbreak{}2401.md,\allowbreak{} Section 3,\allowbreak{} with the original 2290–\allowbreak{}2370 definition and its corrected Hom domains. Explicitly the inverse of G(Y)\allowbreak{}[1]\allowbreak{}\allowbreak{}->G(Y[1]\allowbreak{})\allowbreak{} is G\textquotesingle s exact-functor comparison; the forward map is adjoint to

\[
 F(G(Y)[1])\xrightarrow{\xi_{G(Y)}}F(G(Y))[1]
                         \xrightarrow{\epsilon_Y[1]}Y[1].
\]

Its naturality and bijectivity are the Yoneda calculation in that proof. Applying F to a cone of G(A)\allowbreak{}\allowbreak{}->G(B)\allowbreak{} and using the counit gives the comparison cone\allowbreak{}->G(C)\allowbreak{}; testing with Hom\_\allowbreak{}D(W,\allowbreak{}-)\allowbreak{},\allowbreak{} using adjunction and the two exact sequences,\allowbreak{} makes that comparison an isomorphism. The domain is D,\allowbreak{} not D',\allowbreak{} as reconciled there. Thus the already verified exact-adjoint theorem applies without any new operation.

\subsection{7. Geometric receiving category and three received reports}\label{proofs-11121-11450-geometric-receiving-category-and-three-received-reports}

The citation in perfect.tex 4323–\allowbreak{}4326 must be read in its actual category. Proposition compact-is-perfect at 4246–\allowbreak{}4251 concerns D\_\allowbreak{}QCoh(O\_\allowbreak{}X)\allowbreak{},\allowbreak{} for quasi-compact quasi-separated X. Its proof at 4253–\allowbreak{}4265 uses

\[
 \operatorname{Hom}_{D(O_X)}(P,M)
    =H^0R\Gamma(X,P^\vee\otimes^{\mathbf L}M)
\]

and the preservation of coproducts by RΓ on D\_\allowbreak{}QCoh. The underlying morphisms may be computed in D(O\_\allowbreak{}X)\allowbreak{} because the inclusion is fully faithful. The compactness quantifier,\allowbreak{} however,\allowbreak{} ranges over families in D\_\allowbreak{}QCoh. It does not thereby range over arbitrary objects of D(O\_\allowbreak{}X)\allowbreak{}.

The exact sufficient comparison is the inclusion

\[
 D_{\mathrm{QCoh},T}(O_X)
    \xrightarrow{\ j\ }D_{\mathrm{QCoh}}(O_X)
    \xrightarrow{\ k\ }D(O_X).
\]

Both are full triangulated inclusions and preserve the indicated coproducts. For a perfect supported P and a family M\_\allowbreak{}l in the supported category,\allowbreak{} its compactness comparison there is,\allowbreak{} under full faithfulness,\allowbreak{} precisely its comparison for the same family in D\_\allowbreak{}QCoh. Thus it is an isomorphism. This proves the forward implication of compact-is-perfect-with-support in the needed category; Section 4 supplies the categorical part of its converse with the supported perfect generator.

This validates P05-EN-CH36-0050 /\allowbreak{} OCC-05095:\allowbreak{} replace the ambient D(O\_\allowbreak{}X)\allowbreak{} in line 4324 by D\_\allowbreak{}QCoh(O\_\allowbreak{}X)\allowbreak{}. It is an incorrectly widened cited conclusion; no universal assertion that O\_\allowbreak{}X fails compactness in D(O\_\allowbreak{}X)\allowbreak{} is needed or claimed here. P05-EN-CH36-0051 /\allowbreak{} OCC-05096 is the distinct missing ``that'' at 4326. Their ranges overlap,\allowbreak{} but their semantic dispositions and operations remain separate.

For P05-EN-CH36-0052 /\allowbreak{} OCC-05097,\allowbreak{} the object K\allowbreak{}=RSheafHom(P,\allowbreak{}E)\allowbreak{} at 4379–\allowbreak{}4385 is a sheaf complex. The source writes H\textasciicircum{}0(K)\allowbreak{}\allowbreak{}=Hom\_\allowbreak{}\{D(O\_\allowbreak{}X)\allowbreak{}\}(O\_\allowbreak{}X[0]\allowbreak{},\allowbreak{}K)\allowbreak{}. The former is a sheaf; the latter is an abelian group. The precise equality is

\[
 H^0(X,K):=H^0R\Gamma(X,K)
      =\operatorname{Hom}_{D(O_X)}(O_X[0],K).
\]

To see the map,\allowbreak{} represent K by a K-injective sheaf complex J. Morphisms O\_\allowbreak{}X[0]\allowbreak{}\allowbreak{}->J in the homotopy category are degree-zero cycles of Γ(X,\allowbreak{}J)\allowbreak{} modulo boundaries,\allowbreak{} hence H\textasciicircum{}0Γ(X,\allowbreak{}J)\allowbreak{}. That is the defining computation of H\textasciicircum{}0RΓ(X,\allowbreak{}K)\allowbreak{}. The derived tensor-Hom adjunction sends the original alpha:\allowbreak{}P\allowbreak{}->E to O\_\allowbreak{}X\allowbreak{}->RSheafHom(P,\allowbreak{}E)\allowbreak{},\allowbreak{} so it defines exactly this global class. The correction supplies the missing argument X in H\textasciicircum{}0(X,\allowbreak{}K)\allowbreak{}; it does not replace the original internal Hom object.

The ambient-category repair propagates to 4396–\allowbreak{}4404. Use A\allowbreak{}=D\_\allowbreak{}QCoh(O\_\allowbreak{}X)\allowbreak{} as ambient category and D\allowbreak{}=D\_\allowbreak{}\{QCoh,\allowbreak{}T\}(O\_\allowbreak{}X)\allowbreak{} for the supported tower. The inclusion D\allowbreak{}->A is exact and preserves coproducts,\allowbreak{} so the actual telescope K\allowbreak{}=hocolim K\_\allowbreak{}n holds in A as well as after applying k. The perfect object O\_\allowbreak{}X is compact in A by the cited proposition,\allowbreak{} and alpha factors through K\_\allowbreak{}n there.

Now apply the finite-stage construction of Section 3 in A:\allowbreak{} the tower is made from the perfect supported generator E,\allowbreak{} which is compact in A,\allowbreak{} while its compact domain O\_\allowbreak{}X need not be supported on T. Crucially,\allowbreak{} that finite-stage proof does not require E to generate the ambient A. We obtain O\_\allowbreak{}X\allowbreak{}->Q\allowbreak{}->K\_\allowbreak{}n with Q in <E>\_\allowbreak{}n,\allowbreak{} hence Q perfect and supported on T. Complete O\_\allowbreak{}X\allowbreak{}->Q to I\allowbreak{}->O\_\allowbreak{}X\allowbreak{}->Q\allowbreak{}->I[1]\allowbreak{}. Then I is perfect; restricting to U\allowbreak{}=X\textbackslash{}T makes I\allowbreak{}->O\_\allowbreak{}X an isomorphism since Q|\_\allowbreak{}U\allowbreak{}=0. The composite I\allowbreak{}->O\_\allowbreak{}X\allowbreak{}->K is zero by that triangle. Under the tensor-Hom adjunction,\allowbreak{} this is precisely the required vanishing of I tensor\^{}L P\allowbreak{}->E. Thus the receiving conclusion and its original objects are preserved.

The preview implements three reported repairs (4324,\allowbreak{} 4326,\allowbreak{} 4385)\allowbreak{} and one propagated block repair (4396–\allowbreak{}4404)\allowbreak{},\allowbreak{} under DERIVED-PERFECT-RECEIVER-001/\allowbreak{}002/\allowbreak{}003. These are not added to the count of previously unreported Derived findings. Exact original report records are frozen in PERFECT\_\allowbreak{}RECEIVER\_\allowbreak{}ORIGINAL\_\allowbreak{}REPORTS\_\allowbreak{}PART19.json.

The related reported use at perfect.tex 4500–\allowbreak{}4505,\allowbreak{} P05-EN-CH36-0053,\allowbreak{} was encountered in the recovered report record but its full DGA argument is not read or adjudicated in this pass. It remains explicitly routed to the later Perfect chapter review. This bounded propagation does not certify every occurrence of ambient compactness in that chapter.

\subsection{\texorpdfstring{8. Other receivers,\allowbreak{} copyedits and review limits}{8. Other  copyedits and review limits}}\label{proofs-11121-11450-other-receivers-copyedits-and-review-limits}

At perfect.tex 4904–\allowbreak{}4915 the Brown argument has source D\_\allowbreak{}QCoh(O\_\allowbreak{}X)\allowbreak{} and target D(O\_\allowbreak{}X)\allowbreak{}. Compact generation is needed only in the source. Its cited perfect generator and compactness in D\_\allowbreak{}QCoh therefore give the correct hypothesis for Section 6; the target need not be compactly generated. This pass checks that categorical use,\allowbreak{} not a fresh proof of the cited geometric existence theorem.

At duality.tex 867–\allowbreak{}885 the class under consideration is closed under coproducts,\allowbreak{} shifts and triangles,\allowbreak{} contains the supported perfect generator E,\allowbreak{} and hence contains the first stage Q\_\allowbreak{}1 and every attaching coproduct in Section 2. Induction and the actual telescope triangle then put every supported Q in that class. The omitted [-1]\allowbreak{} in the Derived proof creates no extra hypothesis in this receiving argument.

At injectives.tex 2765–\allowbreak{}2773 the source is D(Mod\_\allowbreak{}R)\allowbreak{}. The free module R in degree zero is compact because Hom(R,\allowbreak{}M)\allowbreak{}\allowbreak{}=H\textasciicircum{}0(M)\allowbreak{},\allowbreak{} and module coproducts are exact and computed termwise in the derived category. Its shifts detect every nonzero cohomology module,\allowbreak{} hence it generates. The exact functor F in that source is left adjoint and sends coproducts to coproducts,\allowbreak{} so H composed with F satisfies Section 5 and is representable. The remainder of that Grothendieck-category argument is outside this bounded check.

P02-E0087 joins the introductory fragment ``With assumptions and notation as in Lemma ...'' to the following if/\allowbreak{}then sentence. The if/\allowbreak{}then sentence itself is complete; the producer calls the wrong part a fragment. The proposed comma and lowercase ``if'' fix the actual introductory fragment. This is a copyedit,\allowbreak{} not a mathematical repair. P02-E0092 adds ``there is'' to the existence sentence in Brown\textquotesingle s proof. Its later repeated wording at 11536 is a different report,\allowbreak{} pending with the next section.

The two omitted shifts remain separate from all these reported corrections. The public comparisons at \href{https://stacks.math.columbia.edu/tag/09SN}{09SN} and \href{https://stacks.math.columbia.edu/tag/09SP}{09SP} were inspected only as secondary comparisons; the frozen original TeX is the authority. No PDF was downloaded or used as authority.

Semantic review ends at 11450. Brown representability,\allowbreak{} bis,\allowbreak{} starting 11460,\allowbreak{} has been read only through 11560 and is next. In particular its argument about maps out of the weaker generating family has not been certified by the compact argument above. Source reports beyond 11450 remain pending.
\clearpage
\section{\texorpdfstring{Brown representability,\allowbreak{} bis:\allowbreak{} complete correction and receiving arguments}{Brown   complete correction and receiving arguments}}\label{proofs-11460-11655-brown-representability-bis-complete-correction-and-receiving-arguments}

Authority:\allowbreak{} derived.tex at a04446e5\allowbreak{}7ec1fbc2\allowbreak{}52a871af\allowbreak{}cec7752f\allowbreak{}b2807b14, lines 11460–\allowbreak{}11655. The frozen authority and current public preimage are bound in INPUT\_\allowbreak{}BINDING.json and CURRENT\_\allowbreak{}PUBLIC\_\allowbreak{}BINDING.json. This note supports PART20; it does not certify subsequent source sections.

Human source:\allowbreak{} Henning Krause,\allowbreak{} *Derived categories,\allowbreak{} resolutions,\allowbreak{} and Brown representability*,\allowbreak{} arXiv:\allowbreak{}math/\allowbreak{}0511047v3,\allowbreak{} author source dated 17 September 2006,\allowbreak{} §4,\allowbreak{} especially §4.5,\allowbreak{} source lines 1165–\allowbreak{}1318. The original compressed source and decompressed TeX are retained under literature/\allowbreak{}KRAUSE\_\allowbreak{}MATH\_\allowbreak{}0511047V3. The source-reading ledger records actual coverage. Krause identifies the proof as following his *A Brown representability theorem via coherent functors*,\allowbreak{} Topology 41 (2002)\allowbreak{},\allowbreak{} 853–\allowbreak{}861,\allowbreak{} the existing Stacks bibliography entry Krause. The 2002 article\textquotesingle s TeX has not been acquired; no PDF substitutes for it. The proof below supplies the restricted-functor construction explicitly and does not assume an unproved coproduct comparison.

\subsection{\texorpdfstring{1. The original tower,\allowbreak{} its shifts,\allowbreak{} and the invalid assertion}{1. The original  its  and the invalid assertion}}\label{proofs-11460-11655-the-original-tower-its-shifts-and-the-invalid-assertion}

Let D be the source triangulated category with arbitrary coproducts. Let E be its given set of objects. Replacing E by all its integral shifts preserves both assumptions. Detection is preserved because the old set is included. To verify the factorization assumption for A[i]\allowbreak{} with A in the old set,\allowbreak{} shift a map A[i]\allowbreak{}\allowbreak{}→\allowbreak{}⊕Y\_\allowbreak{}n by −i,\allowbreak{} apply the old assumption to Y\_\allowbreak{}n[−i]\allowbreak{},\allowbreak{} and shift the resulting entire diagram by i. Its intermediate objects are E\_\allowbreak{}n[i]\allowbreak{} in the enlarged set. The shift equivalences preserve coproducts by their universal property.

Keep exactly the source functor H:\allowbreak{}D\textasciicircum{}op\allowbreak{}→Ab and elements a\_\allowbreak{}n. Put

\[
 X_1=\bigoplus_{(E,a),\,E\in\mathcal E,\ a\in H(E)}E,\qquad
 a_1=(a)_{(E,a)}\in H(X_1).
\]

At stage n define θ\_\allowbreak{}n(f)\allowbreak{}\allowbreak{}=H(f)\allowbreak{}(a\_\allowbreak{}n)\allowbreak{}. Form the source coproduct K\_\allowbreak{}\{n\allowbreak{}+1\} of all pairs (E,\allowbreak{}k)\allowbreak{} with k:\allowbreak{}E\allowbreak{}→X\_\allowbreak{}n and θ\_\allowbreak{}n(k)\allowbreak{}\allowbreak{}=0,\allowbreak{} its evaluation e\_\allowbreak{}n:\allowbreak{}K\_\allowbreak{}\{n\allowbreak{}+1\}\allowbreak{}→X\_\allowbreak{}n,\allowbreak{} and a distinguished triangle

\[
 K_{n+1}\xrightarrow{e_n}X_n\xrightarrow{u_n}X_{n+1}
 \longrightarrow K_{n+1}[1].
\]

Every coordinate of H(e\_\allowbreak{}n)\allowbreak{}(a\_\allowbreak{}n)\allowbreak{} is zero. The comparison from H on the coproduct to the product is an isomorphism,\allowbreak{} so H(e\_\allowbreak{}n)\allowbreak{}(a\_\allowbreak{}n)\allowbreak{}\allowbreak{}=0. The exact H sequence provides a\_\allowbreak{}\{n\allowbreak{}+1\} with H(u\_\allowbreak{}n)\allowbreak{}(a\_\allowbreak{}\{n\allowbreak{}+1\})\allowbreak{}\allowbreak{}=a\_\allowbreak{}n. In particular,\allowbreak{} θ\_\allowbreak{}\{n\allowbreak{}+1\}u\_\allowbreak{}\{n*\}\allowbreak{}=θ\_\allowbreak{}n.

For every E in the generating set,\allowbreak{} θ\_\allowbreak{}n:\allowbreak{}Hom(E,\allowbreak{}X\_\allowbreak{}n)\allowbreak{}\allowbreak{}→H(E)\allowbreak{} is surjective:\allowbreak{} start with the summand of X\_\allowbreak{}1 indexed by (E,\allowbreak{}a)\allowbreak{},\allowbreak{} then compose with the transition maps. Also u\_\allowbreak{}\{n*\} kills ker θ\_\allowbreak{}n,\allowbreak{} because each such map is an evaluation summand of e\_\allowbreak{}n. Conversely u\_\allowbreak{}n k\allowbreak{}=0 implies θ\_\allowbreak{}n(k)\allowbreak{}\allowbreak{}=0 by compatibility. Thus

\[
 \ker(u_{n*})=\ker(\theta_n).
\]

The source at 11554 claims u\_\allowbreak{}nβ\_\allowbreak{}n\allowbreak{}=0 for the arbitrary factors supplied by assumption (2)\allowbreak{}. This is false:\allowbreak{} those factors need not belong to ker θ\_\allowbreak{}n. The surrounding assertion tα\allowbreak{}=α for every α is false as well.

Here is a counterexample to those two assertions within the hypotheses,\allowbreak{} without claiming the theorem itself false. Take D\allowbreak{}=D(k)\allowbreak{} for a field k,\allowbreak{} E\allowbreak{}=\{k[j]\allowbreak{}:\allowbreak{}j\allowbreak{}∈Z\},\allowbreak{} and H\allowbreak{}=Hom\_\allowbreak{}D(−,\allowbreak{}k)\allowbreak{}. The family detects nonzero cohomology,\allowbreak{} and each k[j]\allowbreak{} is compact,\allowbreak{} since Hom\_\allowbreak{}D(k[j]\allowbreak{},\allowbreak{}Y)\allowbreak{}\allowbreak{}=H\textasciicircum{}\{−j\}(Y)\allowbreak{}. Consequently assumption (2)\allowbreak{} holds:\allowbreak{} a map from k[j]\allowbreak{} into a coproduct has finite support; take every intermediate object to be k[j]\allowbreak{},\allowbreak{} take β\_\allowbreak{}n to be its component,\allowbreak{} and let γ have the identity in its finitely many supported positions. Set β\_\allowbreak{}1:\allowbreak{}k\allowbreak{}→X\_\allowbreak{}1 equal to the summand indexed by id\_\allowbreak{}k\allowbreak{}∈H(k)\allowbreak{}. Then θ\_\allowbreak{}1(β\_\allowbreak{}1)\allowbreak{}\allowbreak{}=id\_\allowbreak{}k,\allowbreak{} and θ\_\allowbreak{}2(u\_\allowbreak{}1β\_\allowbreak{}1)\allowbreak{}\allowbreak{}=id\_\allowbreak{}k,\allowbreak{} so u\_\allowbreak{}1β\_\allowbreak{}1≠0. Take β\_\allowbreak{}n\allowbreak{}=0 for n\textgreater1 and let γ be the first-summand inclusion. For S\allowbreak{}=\allowbreak{}⊕X\_\allowbreak{}n,\allowbreak{} write ι\_\allowbreak{}n:\allowbreak{}X\_\allowbreak{}n\allowbreak{}→S and sι\_\allowbreak{}n\allowbreak{}=ι\_\allowbreak{}\{n\allowbreak{}+1\}u\_\allowbreak{}n. Then α\allowbreak{}=ι\_\allowbreak{}1β\_\allowbreak{}1 satisfies

\[
 (1-s)\alpha-\alpha=-\iota_2u_1\beta_1\ne0,
\]

as composition with the second projection shows.

DERIVED-NEW-041 is therefore a proof gap with explicit adverse evidence. Sections 2–\allowbreak{}4 repair it without altering either source hypothesis or replacing the specified tower.

\subsection{2. The exact restricted-functor category and its countable coproducts}\label{proofs-11460-11655-the-exact-restricted-functor-category-and-its-countable-coproducts}

Let P be the full subcategory of objects isomorphic to coproducts of objects in E. For Y\allowbreak{}∈D define hY\allowbreak{}=Hom\_\allowbreak{}D(−,\allowbreak{}Y)\allowbreak{}|\_\allowbreak{}P. Let A be the category of additive functors P\textasciicircum{}op\allowbreak{}→Ab that take all coproducts in P to products,\allowbreak{} with natural transformations. These are genuine products of abelian groups; coproducts in A will not be identified with pointwise sums.

A is locally small. A natural transformation is determined by its components at the set E:\allowbreak{} for any decomposition P\allowbreak{}=\allowbreak{}⊕E\_\allowbreak{}i,\allowbreak{} naturality for the inclusions identifies its component at P with the product of its components at E\_\allowbreak{}i. Hence the class of natural transformations injects into the set ∏\_\allowbreak{}\{E\allowbreak{}∈E\}Hom\_\allowbreak{}Ab(F(E)\allowbreak{},\allowbreak{}G(E)\allowbreak{})\allowbreak{}.

Kernels and cokernels of a natural transformation in A,\allowbreak{} calculated pointwise,\allowbreak{} still take coproducts to products. For kernels this follows coordinate by coordinate. For cokernels the needed identity is

\[
 \operatorname{coker}\left(\prod_i F(P_i)\to\prod_i G(P_i)\right)
   \cong\prod_i\operatorname{coker}(F(P_i)\to G(P_i)).
\]

Indeed the product of the quotient maps is onto by choosing a representative in every coordinate,\allowbreak{} and its kernel consists exactly of coordinatewise images,\allowbreak{} which are the image of the product map. The same coordinate argument identifies images and coimages. Thus A is abelian,\allowbreak{} its exact sequences are pointwise exact,\allowbreak{} and evaluation at any P\allowbreak{}∈P is exact. In particular,\allowbreak{} idempotents split in A,\allowbreak{} with image and kernel calculated pointwise.

For each Y choose

\[
 P_Y=\bigoplus_{(E,f),\,f:E\to Y}E,\qquad p_Y:P_Y\to Y
\]

by evaluation. Hom(P,\allowbreak{}p\_\allowbreak{}Y)\allowbreak{} is onto for every P\allowbreak{}∈P:\allowbreak{} lift each component from E and combine them by the coproduct universal property. Complete p\_\allowbreak{}Y to a triangle C\_\allowbreak{}Y\allowbreak{}→P\_\allowbreak{}Y\allowbreak{}→Y\allowbreak{}→C\_\allowbreak{}Y[1]\allowbreak{},\allowbreak{} choose an evaluation map q\_\allowbreak{}Y:\allowbreak{}Q\_\allowbreak{}Y\allowbreak{}→C\_\allowbreak{}Y with Q\_\allowbreak{}Y\allowbreak{}∈P in the same way,\allowbreak{} and compose it with C\_\allowbreak{}Y\allowbreak{}→P\_\allowbreak{}Y. The Hom exact sequence gives the pointwise exact presentation

\[
 hQ_Y\longrightarrow hP_Y\longrightarrow hY\longrightarrow0. \tag{2.1}
\]

All three functors are in A. For P\_\allowbreak{}0\allowbreak{}∈P,\allowbreak{} hP\_\allowbreak{}0 is representable on P,\allowbreak{} so Yoneda identifies Nat(hP\_\allowbreak{}0,\allowbreak{}F)\allowbreak{} with F(P\_\allowbreak{}0)\allowbreak{}.

We now prove the precise coproduct comparison used in the repair:\allowbreak{}

\[
 h\left(\bigoplus_{n\ge1}Y_n\right)
       \ \cong\ \coprod_{n\ge1}^{A}hY_n. \tag{2.2}
\]

First,\allowbreak{} assumption (2)\allowbreak{} implies that any countable family of maps p\_\allowbreak{}n:\allowbreak{}P\_\allowbreak{}n\allowbreak{}→Y\_\allowbreak{}n which is onto on Hom(E,\allowbreak{}−)\allowbreak{} for every E\allowbreak{}∈E has a coproduct with the same property. For α:\allowbreak{}E\allowbreak{}→\allowbreak{}⊕Y\_\allowbreak{}n write α\allowbreak{}=(\allowbreak{}⊕β\_\allowbreak{}n)\allowbreak{}γ as in the assumption. Lift β\_\allowbreak{}n:\allowbreak{}E\_\allowbreak{}n\allowbreak{}→Y\_\allowbreak{}n through p\_\allowbreak{}n,\allowbreak{} say to \(\widetilde\beta_n:E_n\to P_n\). Then (\allowbreak{}⊕\(\widetilde\beta_n\))\allowbreak{}γ is a lift of α. For a domain P\allowbreak{}∈P do this on each coproduct component.

Apply this to the evaluation maps p\_\allowbreak{}\{Y\_\allowbreak{}n\} and also to q\_\allowbreak{}\{Y\_\allowbreak{}n\}. Coproducts of distinguished triangles are distinguished in D. Here is the needed argument without imposing an additional exactness axiom on coproducts. For triangles A\_\allowbreak{}i\allowbreak{}→B\_\allowbreak{}i\allowbreak{}→C\_\allowbreak{}i\allowbreak{}→A\_\allowbreak{}i[1]\allowbreak{},\allowbreak{} choose a distinguished triangle on \allowbreak{}⊕A\_\allowbreak{}i\allowbreak{}→\allowbreak{}⊕B\_\allowbreak{}i with third object Z. The triangle morphism axiom applied to the inclusions A\_\allowbreak{}i\allowbreak{}→\allowbreak{}⊕A\_\allowbreak{}i and B\_\allowbreak{}i\allowbreak{}→\allowbreak{}⊕B\_\allowbreak{}i gives maps C\_\allowbreak{}i\allowbreak{}→Z. They combine to c:\allowbreak{}\allowbreak{}⊕C\_\allowbreak{}i\allowbreak{}→Z,\allowbreak{} and all three squares commute on the actual inclusions. Applying Hom(−,\allowbreak{}W)\allowbreak{},\allowbreak{} the candidate coproduct triangle has an exact sequence because it is the product of the exact sequences for the individual triangles and products of abelian groups are exact. The comparison with the distinguished triangle on \allowbreak{}⊕A\_\allowbreak{}i\allowbreak{}→\allowbreak{}⊕B\_\allowbreak{}i has identities on the adjacent A and B terms. The five-lemma argument proves Hom(c,\allowbreak{}W)\allowbreak{} is an isomorphism for every W. Yoneda implies c is an isomorphism,\allowbreak{} proving the assertion. Thus the coproduct of the triangles C\_\allowbreak{}\{Y\_\allowbreak{}n\}\allowbreak{}→P\_\allowbreak{}\{Y\_\allowbreak{}n\}\allowbreak{}→Y\_\allowbreak{}n\allowbreak{}→C\_\allowbreak{}\{Y\_\allowbreak{}n\}[1]\allowbreak{},\allowbreak{} together with these two surjectivities,\allowbreak{} gives the exact presentation

\[
 h(\bigoplus Q_{Y_n})\longrightarrow
 h(\bigoplus P_{Y_n})\longrightarrow
 h(\bigoplus Y_n)\longrightarrow0. \tag{2.3}
\]

For an arbitrary F\allowbreak{}∈A,\allowbreak{} applying Nat(−,\allowbreak{}F)\allowbreak{} to the cokernel presentation (2.3)\allowbreak{},\allowbreak{} then Yoneda and the product property of F,\allowbreak{} gives

\[
\begin{aligned}
 \operatorname{Nat}(h(\bigoplus Y_n),F)
  &=\ker\left(F(\bigoplus P_{Y_n})\to F(\bigoplus Q_{Y_n})\right)\\
  &=\ker\left(\prod_n F(P_{Y_n})\to\prod_n F(Q_{Y_n})\right)\\
  &=\prod_n\ker\left(F(P_{Y_n})\to F(Q_{Y_n})\right)\\
  &=\prod_n\operatorname{Nat}(hY_n,F).
\end{aligned}
\]

The maps are restriction along the actual inclusions Y\_\allowbreak{}n\allowbreak{}→\allowbreak{}⊕Y\_\allowbreak{}n,\allowbreak{} since both presentation diagrams commute with the inclusions of P\_\allowbreak{}\{Y\_\allowbreak{}n\} and Q\_\allowbreak{}\{Y\_\allowbreak{}n\}. This proves the coproduct universal property,\allowbreak{} including its canonical structure maps. It is not a compactness assertion:\allowbreak{} evaluation at E need not preserve these coproducts.

The same proof works under the weaker condition that countable coproducts preserve all maps which are onto on Hom(E,\allowbreak{}−)\allowbreak{} for every E\allowbreak{}∈E. Conversely,\allowbreak{} (2.2)\allowbreak{} implies this condition. If all h(f\_\allowbreak{}n)\allowbreak{} are epimorphisms in A,\allowbreak{} their coproduct is an epimorphism:\allowbreak{} two maps agreeing after it agree after every epimorphism h(f\_\allowbreak{}n)\allowbreak{},\allowbreak{} hence on all summands. By (2.2)\allowbreak{} this is h(\allowbreak{}⊕f\_\allowbreak{}n)\allowbreak{},\allowbreak{} and its evaluation at E is onto because epimorphisms in A are pointwise onto. This proves the exact equivalence between that weaker condition and (2.2)\allowbreak{}.

\subsection{3. The split telescope sequence with all maps specified}\label{proofs-11460-11655-the-split-telescope-sequence-with-all-maps-specified}

Set T\allowbreak{}=H|\_\allowbreak{}P and M\_\allowbreak{}n\allowbreak{}=hX\_\allowbreak{}n. All θ\_\allowbreak{}n:\allowbreak{}M\_\allowbreak{}n\allowbreak{}→T are epimorphisms in A. To see this directly at P\allowbreak{}=\allowbreak{}⊕E\_\allowbreak{}i,\allowbreak{} use Hom(P,\allowbreak{}X\_\allowbreak{}n)\allowbreak{}\allowbreak{}=∏Hom(E\_\allowbreak{}i,\allowbreak{}X\_\allowbreak{}n)\allowbreak{},\allowbreak{} H(P)\allowbreak{}\allowbreak{}=∏H(E\_\allowbreak{}i)\allowbreak{},\allowbreak{} and the surjectivity established in §1. The transition v\_\allowbreak{}n\allowbreak{}=h(u\_\allowbreak{}n)\allowbreak{} kills ker θ\_\allowbreak{}n pointwise at E\_\allowbreak{}i,\allowbreak{} and therefore at P. Since θ\_\allowbreak{}1 is a pointwise epimorphism,\allowbreak{} v\_\allowbreak{}1 factors uniquely as

\[
 M_1\xrightarrow{\theta_1}T\xrightarrow{\sigma_2}M_2.
\]

Compatibility and surjectivity of θ\_\allowbreak{}1 give θ\_\allowbreak{}2σ\_\allowbreak{}2\allowbreak{}=id\_\allowbreak{}T. For n≥2 set σ\_\allowbreak{}\{n\allowbreak{}+1\}\allowbreak{}=v\_\allowbreak{}nσ\_\allowbreak{}n. Induction gives

\[
 \theta_n\sigma_n=1_T,\qquad v_n\sigma_n=\sigma_{n+1}.
\]

Write L\_\allowbreak{}n\allowbreak{}=ker θ\_\allowbreak{}n for n≥2. The isomorphism T\allowbreak{}⊕L\_\allowbreak{}n\allowbreak{}→M\_\allowbreak{}n sends (x,\allowbreak{}z)\allowbreak{} to σ\_\allowbreak{}nx\allowbreak{}+z; its inverse sends m to (θ\_\allowbreak{}nm,\allowbreak{}m−σ\_\allowbreak{}nθ\_\allowbreak{}nm)\allowbreak{}. Since v\_\allowbreak{}n kills L\_\allowbreak{}n,\allowbreak{} in these actual coordinates the transition is

\[
 T\oplus L_n\longrightarrow T\oplus L_{n+1},
       \qquad (x,z)\longmapsto(x,0). \tag{3.1}
\]

Let V\allowbreak{}=\allowbreak{}⊕\textasciicircum{}A\_\allowbreak{}\{n≥2\}M\_\allowbreak{}n,\allowbreak{} which exists by (2.2)\allowbreak{}. The idempotent on V induced on its n-th summand by σ\_\allowbreak{}nθ\_\allowbreak{}n has image C and complementary image L. The coproduct universal property and the component retractions show

\[
 C=\coprod_{n\ge2}^{A}T,\quad
 L=\coprod_{n\ge2}^{A}L_n,\quad V=C\oplus L.
\]

For example,\allowbreak{} maps from C to any F correspond exactly to families of maps M\_\allowbreak{}n\allowbreak{}→F that vanish on L\_\allowbreak{}n,\allowbreak{} hence to families of maps T\allowbreak{}→F. This also proves that the displayed coproducts exist; their existence is not assumed from pointwise sums.

On C let j\_\allowbreak{}n:\allowbreak{}T\allowbreak{}→C denote its n-th coproduct inclusion,\allowbreak{} let s\_\allowbreak{}Cj\_\allowbreak{}n\allowbreak{}=j\_\allowbreak{}\{n\allowbreak{}+1\},\allowbreak{} and let ε:\allowbreak{}C\allowbreak{}→T be the fold,\allowbreak{} εj\_\allowbreak{}n\allowbreak{}=id\_\allowbreak{}T. Define \[
 \ell_Cj_n=-\sum_{j=2}^{n-1}j_j,\qquad
 \ell_Cj_2=0.
\] Each sum is finite,\allowbreak{} so the coproduct universal property defines an actual morphism ℓ\_\allowbreak{}C:\allowbreak{}C\allowbreak{}→C. Direct calculation on each inclusion gives

\[
 \ell_C(1-s_C)j_n=j_n,\qquad
 (1-s_C)\ell_Cj_n=j_n-j_2,\qquad \varepsilon j_2=1_T.
\]

Consequently \[
 \ell_C(1-s_C)=1_C,\quad
 (1-s_C)\ell_C=1_C-j_2\varepsilon,\quad
 \varepsilon(1-s_C)=0.
\] These identities prove the split exact sequence 0\allowbreak{}→C\allowbreak{}→\textasciicircum{}\{1−s\_\allowbreak{}C\}C\allowbreak{}→\textasciicircum{}εT\allowbreak{}→0. In (3.1)\allowbreak{} the transition is zero on L,\allowbreak{} so the tail telescope map t\_\allowbreak{}V:\allowbreak{}V\allowbreak{}→V is (1−s\_\allowbreak{}C)\allowbreak{}\allowbreak{}⊕id\_\allowbreak{}L. Its left inverse is ℓ\_\allowbreak{}V\allowbreak{}=ℓ\_\allowbreak{}C\allowbreak{}⊕id\_\allowbreak{}L,\allowbreak{} and its cokernel is T with map induced by the θ\_\allowbreak{}n.

Retain the initial object M\_\allowbreak{}1,\allowbreak{} which need not split as T\allowbreak{}⊕ker θ\_\allowbreak{}1. Put b\allowbreak{}=j\_\allowbreak{}2\textasciicircum{}Vv\_\allowbreak{}1:\allowbreak{}M\_\allowbreak{}1\allowbreak{}→V. Relative to M\allowbreak{}=M\_\allowbreak{}1\allowbreak{}⊕V\allowbreak{}=hS,\allowbreak{} where S\allowbreak{}=\allowbreak{}⊕\_\allowbreak{}\{n≥1\}X\_\allowbreak{}n,\allowbreak{} the actual telescope map t\allowbreak{}=h(1−s)\allowbreak{} is

\[
 t(x,y)=(x,t_Vy-bx).
\]

The target automorphism A\_\allowbreak{}0(x,\allowbreak{}y)\allowbreak{}\allowbreak{}=(x,\allowbreak{}y\allowbreak{}+bx)\allowbreak{},\allowbreak{} whose inverse is (x,\allowbreak{}y)\allowbreak{}\allowbreak{}↦(x,\allowbreak{}y−bx)\allowbreak{},\allowbreak{} satisfies \[
 A_0t(x,y)=(x,t_Vy).
\] Thus \[
 \ell(x,y)=(x,\ell_V(y+bx))
\] is a left inverse of t. Its cokernel is T. In fact the cokernel map Θ:\allowbreak{}M\allowbreak{}→T is induced by the original θ\_\allowbreak{}n; in these coordinates Θ(x,\allowbreak{}y)\allowbreak{}\allowbreak{}=θ\_\allowbreak{}1x\allowbreak{}+Θ\_\allowbreak{}Vy,\allowbreak{} because Θ\_\allowbreak{}Vb\allowbreak{}=θ\_\allowbreak{}2v\_\allowbreak{}1\allowbreak{}=θ\_\allowbreak{}1. A section is x\allowbreak{}↦(0,\allowbreak{}j\_\allowbreak{}2\textasciicircum{}Vσ\_\allowbreak{}2x)\allowbreak{}. We have proved the split exact sequence in A

\[
 0\longrightarrow hS\xrightarrow{h(1-s)}hS
          \xrightarrow{\Theta}H|_P\longrightarrow0. \tag{3.2}
\]

Evaluating its left inverse at E proves exactly the source injectivity Hom(E,\allowbreak{}S)\allowbreak{}\allowbreak{}→\textasciicircum{}\{(1−s)\allowbreak{}\_\allowbreak{}*\}Hom(E,\allowbreak{}S)\allowbreak{}. It does not imply that 1−s:\allowbreak{}S\allowbreak{}→S has a left inverse in D:\allowbreak{} the natural transformation ℓ is only defined on the restriction to P. Since E[−1]\allowbreak{} is also in the enlarged set,\allowbreak{} the natural shift comparison gives injectivity on Hom(E,\allowbreak{}S[1]\allowbreak{})\allowbreak{} as well.

\subsection{\texorpdfstring{4. The homotopy-colimit map,\allowbreak{} bijectivity,\allowbreak{} and full representability}{4. The homotopy-colimit   and full representability}}\label{proofs-11460-11655-the-homotopy-colimit-map-bijectivity-and-full-representability}

Choose exactly the source telescope triangle

\[
 S\xrightarrow{t_D=1-s}S\xrightarrow{q}X\xrightarrow{w}S[1].
\]

Let a\_\allowbreak{}S\allowbreak{}∈H(S)\allowbreak{} correspond to (a\_\allowbreak{}n)\allowbreak{}. On its n-th coordinate,\allowbreak{} H(t\_\allowbreak{}D)\allowbreak{}(a\_\allowbreak{}S)\allowbreak{} is a\_\allowbreak{}n−H(u\_\allowbreak{}n)\allowbreak{}(a\_\allowbreak{}\{n\allowbreak{}+1\})\allowbreak{}\allowbreak{}=0. Thus the H exact sequence gives a\allowbreak{}∈H(X)\allowbreak{} with H(q)\allowbreak{}(a)\allowbreak{}\allowbreak{}=a\_\allowbreak{}S. Let θ:\allowbreak{}h\_\allowbreak{}X\allowbreak{}→H be the Yoneda transformation f\allowbreak{}↦H(f)\allowbreak{}(a)\allowbreak{}. Its composite with every q\_\allowbreak{}n\allowbreak{}=qι\_\allowbreak{}n is θ\_\allowbreak{}n.

The injectivity proved in §3,\allowbreak{} also for the shifted S,\allowbreak{} and the Hom sequence of the telescope triangle give

\[
 \operatorname{Hom}(E,S)\xrightarrow{t_{D*}}
 \operatorname{Hom}(E,S)\xrightarrow{q_*}
 \operatorname{Hom}(E,X)\longrightarrow0. \tag{4.1}
\]

Comparison with (3.2)\allowbreak{} identifies θ\_\allowbreak{}E:\allowbreak{}Hom(E,\allowbreak{}X)\allowbreak{}\allowbreak{}→H(E)\allowbreak{} as an isomorphism. This argument also proves the stronger natural statement hX≅H|\_\allowbreak{}P,\allowbreak{} with the isomorphism induced by the actual element a.

For completeness,\allowbreak{} the original next paragraph,\allowbreak{} 11571–\allowbreak{}11599,\allowbreak{} now works with its original maps. Lift f:\allowbreak{}E\allowbreak{}→X to α:\allowbreak{}E\allowbreak{}→S using (4.1)\allowbreak{},\allowbreak{} and factor α\allowbreak{}=(\allowbreak{}⊕β\_\allowbreak{}n)\allowbreak{}γ by assumption (2)\allowbreak{}. Choose δ\_\allowbreak{}n:\allowbreak{}E\_\allowbreak{}n\allowbreak{}→X\_\allowbreak{}1 with θ\_\allowbreak{}1(δ\_\allowbreak{}n)\allowbreak{}\allowbreak{}=θ\_\allowbreak{}n(β\_\allowbreak{}n)\allowbreak{},\allowbreak{} and let u\_\allowbreak{}\{n,\allowbreak{}1\}:\allowbreak{}X\_\allowbreak{}1\allowbreak{}→X\_\allowbreak{}n be the transition composite (the identity when n\allowbreak{}=1)\allowbreak{}. Then β\_\allowbreak{}n−u\_\allowbreak{}\{n,\allowbreak{}1\}δ\_\allowbreak{}n is in ker θ\_\allowbreak{}n,\allowbreak{} so u\_\allowbreak{}nβ\_\allowbreak{}n\allowbreak{}=u\_\allowbreak{}nu\_\allowbreak{}\{n,\allowbreak{}1\}δ\_\allowbreak{}n. Since q\_\allowbreak{}\{n\allowbreak{}+1\}u\_\allowbreak{}n\allowbreak{}=q\_\allowbreak{}n,\allowbreak{}

\[
 q_n\beta_n=q_1\delta_n.
\]

Maps out of \allowbreak{}⊕E\_\allowbreak{}n are determined by their components. Therefore qα\allowbreak{}=q\_\allowbreak{}1dγ,\allowbreak{} where d:\allowbreak{}\allowbreak{}⊕E\_\allowbreak{}n\allowbreak{}→X\_\allowbreak{}1 has components δ\_\allowbreak{}n. This proves Hom(E,\allowbreak{}X\_\allowbreak{}1)\allowbreak{}\allowbreak{}→Hom(E,\allowbreak{}X)\allowbreak{} is onto. If θ\_\allowbreak{}1(g)\allowbreak{}\allowbreak{}=0 then u\_\allowbreak{}1g\allowbreak{}=0,\allowbreak{} so q\_\allowbreak{}1g\allowbreak{}=0. Since θ\_\allowbreak{}1 is onto,\allowbreak{} θ\_\allowbreak{}E is onto; if θ\_\allowbreak{}E(f)\allowbreak{}\allowbreak{}=0,\allowbreak{} choose f\allowbreak{}=q\_\allowbreak{}1g,\allowbreak{} obtain θ\_\allowbreak{}1(g)\allowbreak{}\allowbreak{}=0 and f\allowbreak{}=0. These are the exact surjectivity and injectivity claims in the source paragraph.

Define D\textquotesingle{} exactly as in the source:\allowbreak{} Y\allowbreak{}∈D' if θ\_\allowbreak{}\{Y[j]\allowbreak{}\} is an isomorphism for every j\allowbreak{}∈Z. This class is strictly full and shift stable. For a distinguished triangle,\allowbreak{} the two Hom/\allowbreak{}H long exact sequences form a commutative diagram; if θ is an isomorphism at the two other objects in all shifts,\allowbreak{} the five-lemma argument on five successive terms gives it at the third. For a retraction Y\allowbreak{}→Z\allowbreak{}→Y,\allowbreak{} the naturality squares make θ\_\allowbreak{}Y a retract of θ\_\allowbreak{}Z,\allowbreak{} and the inverse is obtained by composing the inclusion,\allowbreak{} θ\_\allowbreak{}Z\textasciicircum{}\{-1\},\allowbreak{} and projection. Thus D\textquotesingle{} is saturated. On a coproduct of objects in D',\allowbreak{} the map θ is the product of their isomorphisms,\allowbreak{} since both functors take coproducts to products. Hence D\textquotesingle{} is closed under coproducts,\allowbreak{} and the telescope triangle makes it closed under these homotopy colimits. The preceding calculation puts E in D\textquotesingle.

For arbitrary Y\allowbreak{}∈D,\allowbreak{} repeat the already established construction for H\_\allowbreak{}Y\allowbreak{}=Hom\_\allowbreak{}D(−,\allowbreak{}Y)\allowbreak{},\allowbreak{} naming the resulting tower X\_\allowbreak{}n\textasciicircum{}Y and telescope X\^{}Y. The construction uses sums of E and their cones,\allowbreak{} so all X\_\allowbreak{}n\textasciicircum{}Y and X\^{}Y belong to the original D\textquotesingle. Yoneda identifies the constructed element with a map f\_\allowbreak{}Y:\allowbreak{}X\textasciicircum{}Y\allowbreak{}→Y,\allowbreak{} and the preceding argument proves Hom(E,\allowbreak{}f\_\allowbreak{}Y)\allowbreak{} is bijective for every E\allowbreak{}∈E. Complete f\_\allowbreak{}Y to a triangle X\textasciicircum{}Y\allowbreak{}→Y\allowbreak{}→C\allowbreak{}→X\textasciicircum{}Y[1]\allowbreak{}. Since E is shift stable,\allowbreak{} the maps on both adjacent Hom terms are isomorphisms,\allowbreak{} so Hom(E,\allowbreak{}C)\allowbreak{}\allowbreak{}=0 for every E. Assumption (1)\allowbreak{} implies C\allowbreak{}=0. Hence f\_\allowbreak{}Y is an isomorphism and Y\allowbreak{}∈D'. Thus D'\allowbreak{}=D,\allowbreak{} and θ:\allowbreak{}h\_\allowbreak{}X\allowbreak{}→H is the required natural isomorphism.

This also records exactly where detection is used:\allowbreak{} only the final cone argument. Sections 1–\allowbreak{}3 and the restricted conclusion hX≅H|\_\allowbreak{}P do not require detection.

\subsection{\texorpdfstring{5. Underclaim,\allowbreak{} weaker hypothesis,\allowbreak{} and the exact right adjoint}{5.  weaker  and the exact right adjoint}}\label{proofs-11460-11655-underclaim-weaker-hypothesis-and-the-exact-right-adjoint}

DERIVED-UNDER-030 records the split sequence (3.2)\allowbreak{},\allowbreak{} its explicit retraction,\allowbreak{} and the natural restricted representation hX≅H|\_\allowbreak{}P. It is a proved editorial strengthening of the source\textquotesingle s stated injectivity. Its hypotheses are the shift-stable set,\allowbreak{} the specified kernel-killing tower for the cohomological coproduct-to-product functor,\allowbreak{} and countable preservation of E-epimorphisms proved in §2. Detection is not needed for that restricted conclusion.

Under the weaker countable preservation condition itself,\allowbreak{} sections 2–\allowbreak{}4 up to hX≅H|\_\allowbreak{}P remain unchanged. The alternative source factorization paragraph is no longer necessary:\allowbreak{} (3.2)\allowbreak{} and the telescope exact sequence already prove θ\_\allowbreak{}E is an isomorphism. With detection,\allowbreak{} the final cone argument proves full Brown representability under that weaker condition. This is the known perfect-generator mechanism of Krause,\allowbreak{} not a novelty claim. The source\textquotesingle s original assumption (2)\allowbreak{} and its original formulation remain in the corrected source body.

For the proposition at 11633–\allowbreak{}11655 let F:\allowbreak{}D\allowbreak{}→D\_\allowbreak{}1 be exact and preserve coproducts. For Y\allowbreak{}∈D\_\allowbreak{}1 define H\_\allowbreak{}Y(W)\allowbreak{}\allowbreak{}=Hom\_\allowbreak{}\{D\_\allowbreak{}1\}(F(W)\allowbreak{},\allowbreak{}Y)\allowbreak{}. Applying Hom(−,\allowbreak{}Y)\allowbreak{} to the distinguished triangle F sends a triangle to proves cohomological exactness,\allowbreak{} including the specified shift isomorphism of F. The coproduct universal property and preservation by F give H\_\allowbreak{}Y(\allowbreak{}⊕W\_\allowbreak{}i)\allowbreak{}\allowbreak{}=∏H\_\allowbreak{}Y(W\_\allowbreak{}i)\allowbreak{}. The repaired theorem supplies GY and an isomorphism η\_\allowbreak{}\{W,\allowbreak{}Y\}:\allowbreak{}Hom\_\allowbreak{}D(W,\allowbreak{}GY)\allowbreak{}\allowbreak{}→Hom\_\allowbreak{}\{D\_\allowbreak{}1\}(F(W)\allowbreak{},\allowbreak{}Y)\allowbreak{}.

For a:\allowbreak{}Y\allowbreak{}→Y',\allowbreak{} define G(a)\allowbreak{}:\allowbreak{}GY\allowbreak{}→GY' as the unique map corresponding by Yoneda to the natural transformation η\_\allowbreak{}\{−,\allowbreak{}Y'\}\textasciicircum{}\{-1\}∘(a∘−)\allowbreak{}∘η\_\allowbreak{}\{−,\allowbreak{}Y\}. The uniqueness proves identity and composition laws. Thus η is natural in both variables and is an adjunction. The unit at W is η\_\allowbreak{}\{W,\allowbreak{}FW\}\textasciicircum{}\{−1\}(id\_\allowbreak{}\{FW\})\allowbreak{}; the counit at Y is η\_\allowbreak{}\{GY,\allowbreak{}Y\}(id\_\allowbreak{}\{GY\})\allowbreak{}. Naturality and the two corresponding identity maps give the two triangle identities.

Exactness is the already proved exact-adjoint lemma,\allowbreak{} with full shift comparison and morphism-of-triangles argument retained in PROOFS\_\allowbreak{}1844\_\allowbreak{}2401.md §3. Concretely its shift isomorphism is the Yoneda representative of Hom\_\allowbreak{}D(W,\allowbreak{}G(Y[1]\allowbreak{})\allowbreak{})\allowbreak{}≅Hom\_\allowbreak{}\{D\_\allowbreak{}1\}(F(W)\allowbreak{},\allowbreak{}Y[1]\allowbreak{})\allowbreak{} ≅Hom\_\allowbreak{}\{D\_\allowbreak{}1\}(F(W[−1]\allowbreak{})\allowbreak{},\allowbreak{}Y)\allowbreak{}≅Hom\_\allowbreak{}D(W,\allowbreak{}GY[1]\allowbreak{})\allowbreak{}. For a triangle in D\_\allowbreak{}1,\allowbreak{} choose a cone of its first image map in D,\allowbreak{} construct the comparison to the third image using the adjunction,\allowbreak{} and apply Hom\_\allowbreak{}D(W,\allowbreak{}−)\allowbreak{}. The five-term exact comparison proves that cone comparison is an isomorphism for every W. This is exactly the source lemma\textquotesingle s argument already checked,\allowbreak{} with no additional compactness assumption introduced here.

\subsection{6. Existing reports and exactly-once scope accounting}\label{proofs-11460-11655-existing-reports-and-exactly-once-scope-accounting}

P02-E0093 /\allowbreak{} OCC-01561 reports the missing finite verb at 11536. ``Hence there is a natural transformation'' repairs this sentence. The identical earlier phrase was separately corrected in PART19; the operation here is located by the authority/\allowbreak{}current line alignment,\allowbreak{} not by choosing the first matching text.

P02-E0094 /\allowbreak{} OCC-01562 originally mentioned both 11522 and 11624–\allowbreak{}11625. The authority at 11522 already says hocolim. Its own scope resolution P02-ER0003 /\allowbreak{} OCC-02894,\allowbreak{} original producer ledger line 1430,\allowbreak{} explicitly withdraws the first locus. That resolution has no numeric routing range and is metadata,\allowbreak{} not another defect. Preserve its original empty range,\allowbreak{} record the effective linked locus,\allowbreak{} and count its physical report exactly once with metadata-only disposition.

At 11625 the authority says ordinary colim,\allowbreak{} while the construction is the telescope defined at 11522–\allowbreak{}11527. The current source already contains the exact admitted operation MC-STK-ERR-0898-OP1 replacing this by hocolim. DERIVED-087 /\allowbreak{} OCC-12612 reports the same defect. The two defect reports and the metadata resolution share one reconciliation group,\allowbreak{} with the metadata\textquotesingle s individual disposition explicit. They add no new correction operation at 11625 and no operation at 11522.

Four physical reports are completed in PART20:\allowbreak{} one missing copyedit,\allowbreak{} two already-corrected defect reports,\allowbreak{} and one metadata resolution. There are two new semantic reconciliation groups,\allowbreak{} not three new mathematical defects. The unreported Brown proof gap is separately identified as DERIVED-NEW-041.

\subsection{7. Downstream scope and receiving maps}\label{proofs-11460-11655-downstream-scope-and-receiving-maps}

A bounded search of root TeX chapter files for the Brown-bis section,\allowbreak{} lemma and proposition labels found the internal proposition,\allowbreak{} sites-cohomology.tex 11481/\allowbreak{}11484,\allowbreak{} and the literature pointer in injectives.tex 2755. The search does not certify files outside that scope or invisible semantic dependencies.

The Injectives passage at 2743–\allowbreak{}2779 directs the reader to this more general theorem. Its actual displayed argument at 2765–\allowbreak{}2773 uses compact Brown on D(Mod\_\allowbreak{}R)\allowbreak{},\allowbreak{} as already checked in PART19. The repair does not change that argument or its hypotheses.

The receiving proposition in sites-cohomology.tex 11458–\allowbreak{}11486 is about the full category QC(O)\allowbreak{} defined at 11098–\allowbreak{}11111 by the actual base-change morphisms

\[
 R\Gamma(V,K)\otimes_{\mathcal O(V)}^{\mathbf L}\mathcal O(U)
       \longrightarrow R\Gamma(U,K).
\]

The cited lemma at 11113–\allowbreak{}11134 supplies its triangulated structure and coproducts. On this chaotic site,\allowbreak{} evaluation is exact,\allowbreak{} so the first factor is computed by the original complex of sections. Evaluation,\allowbreak{} derived tensor product and their comparison are exact and commute with coproducts. The objects on which the comparison is an isomorphism are therefore closed under cones,\allowbreak{} shifts,\allowbreak{} retracts and coproducts,\allowbreak{} by the same exact-sequence and retraction arguments used in §4. This proves the categorical closure required by Brown.

The source cardinal lemma at 11296–\allowbreak{}11456 supplies the two original Brown hypotheses. To spell out the receiving set,\allowbreak{} take representatives of QC(O)\allowbreak{} objects with a complex of cardinality at most κ. This is a set:\allowbreak{} the disjoint union of all degree/\allowbreak{}object components has size at most κ,\allowbreak{} so representatives of its underlying sets and all module,\allowbreak{} restriction and differential structure maps range over sets. Replacing each small object in the cardinal lemma by its chosen representative transports every map by an isomorphism. Its two conclusions become exactly Brown (1)\allowbreak{} and (2)\allowbreak{}. Shifts preserve the cardinality,\allowbreak{} since they reindex the same degree/\allowbreak{}object disjoint union. No compactness of these small objects is asserted or needed.

The countable factorization in that lemma is compatible with the actual maps:\allowbreak{} a chain map from a κ-small complex into \allowbreak{}⊕K\_\allowbreak{}n has images in the selected small Cartesian subcomplexes E\_\allowbreak{}n⊂K\_\allowbreak{}n. Each section already has finite support in the original coproduct; passing to these subcomplexes keeps the support and yields the chain map to \allowbreak{}⊕E\_\allowbreak{}n. Its composition with the inclusions is the original map. This observation uses the cited cardinal enlargement lemma; the complete independent review of that chapter and of its More on Algebra enlargement dependency remains in their own source queues. Reading this receiving argument does not mark those chapters reviewed.

Therefore the repaired Brown theorem applies with exactly the receiver\textquotesingle s existing hypotheses. Its part (2)\allowbreak{} retains representability,\allowbreak{} part (3)\allowbreak{} retains the exact right adjoint,\allowbreak{} and part (4)\allowbreak{} uses the actual full inclusion QC(O)\allowbreak{}\allowbreak{}→D(O)\allowbreak{},\allowbreak{} which is exact and preserves coproducts by the closure argument. No receiving source emendation is required for DERIVED-NEW-041. These propagation conclusions concern this Brown dependency,\allowbreak{} not blanket approval of the whole receiving chapter.

\subsection{8. Rejected shortcuts and retained scope}\label{proofs-11460-11655-rejected-shortcuts-and-retained-scope}

Coordinate projections of a map E\allowbreak{}→\allowbreak{}⊕X\_\allowbreak{}n may all vanish without the map vanishing when E is not compact. The attempted argument from projections alone therefore does not prove the required injectivity. Nor does the factorization hypothesis put each arbitrary β\_\allowbreak{}n into ker θ\_\allowbreak{}n. Neither shortcut is used in §§2–\allowbreak{}4.

For reference,\allowbreak{} the direct finite-sum calculation from an arbitrary factorization does prove a cokernel statement. Choose δ\_\allowbreak{}n as in §4,\allowbreak{} put B\allowbreak{}=\allowbreak{}⊕E\_\allowbreak{}n and β|\_\allowbreak{}\{E\_\allowbreak{}n\}\allowbreak{}=ι\_\allowbreak{}nβ\_\allowbreak{}n,\allowbreak{} \(\widetilde\beta|_{E_n}=\iota_nu_{n,1}\delta_n\),\allowbreak{} and k\allowbreak{}=β−\(\widetilde\beta\). Then sk\allowbreak{}=0 and t\_\allowbreak{}Dk\allowbreak{}=k. Define

\[
 g|_{E_n}=\sum_{j=1}^{n-1}\iota_ju_{j,1}\delta_n,\qquad
 d|_{E_n}=\delta_n.
\]

Every defining sum is finite,\allowbreak{} and direct telescoping gives t\_\allowbreak{}Dg\allowbreak{}=ι\_\allowbreak{}1d−\(\widetilde\beta\),\allowbreak{} whence β\allowbreak{}=ι\_\allowbreak{}1d\allowbreak{}+t\_\allowbreak{}D(k−g)\allowbreak{}. After composition with γ,\allowbreak{} every α:\allowbreak{}E\allowbreak{}→S therefore has the form ι\_\allowbreak{}1z\allowbreak{}+t\_\allowbreak{}Dη. This calculation alone neither proves that t\_\allowbreak{}D is injective on Hom(E,\allowbreak{}S)\allowbreak{} nor that every E\allowbreak{}→X lifts to S. The restricted-functor argument supplies exactly these missing facts.

The theorem,\allowbreak{} tower,\allowbreak{} original generating hypotheses,\allowbreak{} all transition maps and the receiving formulations are retained. The stronger split statement and weaker perfectness formulation are editorial consequences with human-source attribution; neither is a claimed new discovery or an undisclosed replacement of the source theorem.
\clearpage
\section{\texorpdfstring{Admissible subcategories:\allowbreak{} hypotheses,\allowbreak{} projections,\allowbreak{} and receiving criteria}{Admissible    and receiving criteria}}\label{proofs-11665-12003-admissible-subcategories-hypotheses-projections-and-receiving-criteria}

Authority:\allowbreak{} derived.tex at a04446e5\allowbreak{}7ec1fbc2\allowbreak{}52a871af\allowbreak{}cec7752f\allowbreak{}b2807b14, original lines 11665–\allowbreak{}12003. Current-source comparisons use the frozen public preimages at 95a51593\allowbreak{}aceb247a\allowbreak{}c7a18311\allowbreak{}1678c3a0\allowbreak{}f8b8f4b5. This note supports PART21. The source cites Bondal–\allowbreak{}Kapranov,\allowbreak{} §1; the arguments below are proved from the specified triangulated-category axioms and the earlier source lemmas,\allowbreak{} without claiming to have read that reference in this pass or claiming novelty.

\subsection{1. Orthogonals and the two exact Hom criteria}\label{proofs-11665-12003-orthogonals-and-the-two-exact-hom-criteria}

For a full subcategory A of an additive category D,\allowbreak{} the original right orthogonal A\^{}⊥ has objects Y for which Hom\_\allowbreak{}D(A,\allowbreak{}Y)\allowbreak{}\allowbreak{}=0 for every A\allowbreak{}∈A. The original left orthogonal \^{}⊥A has objects Y for which Hom\_\allowbreak{}D(Y,\allowbreak{}A)\allowbreak{}\allowbreak{}=0 for every A\allowbreak{}∈A. They are full subcategories by definition.

Suppose D is triangulated and A is invariant under all shifts. In the triangle \[
 X\xrightarrow fY\xrightarrow gZ\xrightarrow hX[1]
\] the natural exact Hom sequence contains \[
 \operatorname{Hom}(A[1],Z)\longrightarrow\operatorname{Hom}(A,X)
 \xrightarrow{f_*}\operatorname{Hom}(A,Y)
 \longrightarrow\operatorname{Hom}(A,Z).
\] If Z\allowbreak{}∈A\textasciicircum{}⊥,\allowbreak{} both outer groups vanish. Conversely,\allowbreak{} suppose f\_\allowbreak{}* is an isomorphism for every A\allowbreak{}∈A. The exact sequence \[
 \operatorname{Hom}(A,X)\xrightarrow{f_*}\operatorname{Hom}(A,Y)
 \longrightarrow\operatorname{Hom}(A,Z)
 \longrightarrow\operatorname{Hom}(A[-1],X)
 \xrightarrow{f_*}\operatorname{Hom}(A[-1],Y)
\] has its first and last arrows isomorphisms,\allowbreak{} hence Hom(A,\allowbreak{}Z)\allowbreak{}\allowbreak{}=0. This proves 11682–\allowbreak{}11715. The equality notation in its statement means the isomorphism induced by f,\allowbreak{} not an arbitrary abstract group isomorphism.

For the dual statement fix a shift-invariant B. If X\allowbreak{}∈\textasciicircum{}⊥B,\allowbreak{} the exact sequence \[
 \operatorname{Hom}(X[1],B)\longrightarrow\operatorname{Hom}(Z,B)
 \xrightarrow{g^*}\operatorname{Hom}(Y,B)
 \longrightarrow\operatorname{Hom}(X,B)
\] gives the required isomorphism. Conversely,\allowbreak{} apply the assumption to B and B[1]\allowbreak{} in \[
 \operatorname{Hom}(Z,B)\longrightarrow\operatorname{Hom}(Y,B)
 \longrightarrow\operatorname{Hom}(X,B)
 \longrightarrow\operatorname{Hom}(Z[-1],B)
 \longrightarrow\operatorname{Hom}(Y[-1],B).
\] The last arrow is the assumed isomorphism for B[1]\allowbreak{},\allowbreak{} under the shift identifications Hom(Z[−1]\allowbreak{},\allowbreak{}B)\allowbreak{}\allowbreak{}=Hom(Z,\allowbreak{}B[1]\allowbreak{})\allowbreak{}. Thus Hom(X,\allowbreak{}B)\allowbreak{}\allowbreak{}=0. This explicitly proves the omitted dual at 11717–\allowbreak{}11735.

Both orthogonals are strictly full because composition with an isomorphism identifies the relevant Hom groups. They contain zero and are closed under finite biproducts,\allowbreak{} by additivity of Hom. For a retraction Y\allowbreak{}→Z\allowbreak{}→Y,\allowbreak{} the identity composite makes Hom(A,\allowbreak{}Y)\allowbreak{} a retract of Hom(A,\allowbreak{}Z)\allowbreak{},\allowbreak{} and similarly for Hom(Y,\allowbreak{}A)\allowbreak{}. Thus a retract of an orthogonal object remains orthogonal. Shifts preserve orthogonality because Hom(A,\allowbreak{}Y[j]\allowbreak{})\allowbreak{}\allowbreak{}=Hom(A[−j]\allowbreak{},\allowbreak{}Y)\allowbreak{} and the dual formula retain the actual shift isomorphisms. The two exact-sequence criteria show closure under cones; rotation gives two-out-of-three closure. The full triangulated-subcategory criterion at 860–\allowbreak{}880 then supplies the inherited triangulated structure. This proves 11737–\allowbreak{}11758,\allowbreak{} including saturation in the source definition at 1758–\allowbreak{}1763.

The direct sums needed in that proof are finite biproducts. Arbitrary coproduct closure of A\^{}⊥ is not used or inferred. In contrast,\allowbreak{} \^{}⊥A is closed under any existing coproducts by Hom(\allowbreak{}⊕Y\_\allowbreak{}i,\allowbreak{}A)\allowbreak{}\allowbreak{}=∏Hom(Y\_\allowbreak{}i,\allowbreak{}A)\allowbreak{}; A\^{}⊥ is closed under any existing products by Hom(A,\allowbreak{}∏Y\_\allowbreak{}i)\allowbreak{}\allowbreak{}=∏Hom(A,\allowbreak{}Y\_\allowbreak{}i)\allowbreak{}. If every testing object A is compact,\allowbreak{} A\^{}⊥ is also closed under existing coproducts by the defining compact comparison. These precise universal-property statements do not strengthen the source hypotheses silently.

\subsection{2. Existence of approximation triangles is stable under triangles}\label{proofs-11665-12003-existence-of-approximation-triangles-is-stable-under-triangles}

Let A be the source full triangulated subcategory,\allowbreak{} not necessarily strictly full. Suppose X\_\allowbreak{}i has a triangle \[
 A_i\xrightarrow{a_i}X_i\xrightarrow{b_i}B_i
 \xrightarrow{\delta_i}A_i[1],\qquad A_i\in A,\quad B_i\in A^\perp
\] for i\allowbreak{}=1,\allowbreak{}2,\allowbreak{} and let f:\allowbreak{}X\_\allowbreak{}1\allowbreak{}→X\_\allowbreak{}2 be the first map of the given distinguished triangle on X\_\allowbreak{}1,\allowbreak{}X\_\allowbreak{}2,\allowbreak{}X\_\allowbreak{}3. By §1 the map Hom(A\_\allowbreak{}1,\allowbreak{}A\_\allowbreak{}2)\allowbreak{}\allowbreak{}→Hom(A\_\allowbreak{}1,\allowbreak{}X\_\allowbreak{}2)\allowbreak{},\allowbreak{} α\allowbreak{}↦a\_\allowbreak{}2α,\allowbreak{} is an isomorphism. Define α by a\_\allowbreak{}2α\allowbreak{}=fa\_\allowbreak{}1. This specifies the exact commutative square used in the source.

Apply the previously proved nine-object proposition,\allowbreak{} original 1031–\allowbreak{}1052 and PROOFS\_\allowbreak{}0001\_\allowbreak{}1140.md,\allowbreak{} to that square. Choose the middle column to be the specified triangle X\_\allowbreak{}1\allowbreak{}→X\_\allowbreak{}2\allowbreak{}→X\_\allowbreak{}3\allowbreak{}→X\_\allowbreak{}1[1]\allowbreak{},\allowbreak{} and complete the first column on α. The resulting rows are A\_\allowbreak{}i\allowbreak{}→X\_\allowbreak{}i\allowbreak{}→Q\_\allowbreak{}i\allowbreak{}→A\_\allowbreak{}i[1]\allowbreak{},\allowbreak{} and the third column is the distinguished triangle Q\_\allowbreak{}1\allowbreak{}→Q\_\allowbreak{}2\allowbreak{}→Q\_\allowbreak{}3\allowbreak{}→Q\_\allowbreak{}1[1]\allowbreak{}. The lower-right square has the anticommuting sign required by the nine-object proposition; it is not replaced by a commuting square. All other squares commute.

The first two rows are triangles on the original a\_\allowbreak{}i. The triangle morphism axiom gives (id\_\allowbreak{}\{A\_\allowbreak{}i\},\allowbreak{}id\_\allowbreak{}\{X\_\allowbreak{}i\},\allowbreak{}q\_\allowbreak{}i)\allowbreak{} to the originally chosen triangle,\allowbreak{} and the Hom exact sequences show q\_\allowbreak{}i:\allowbreak{}Q\_\allowbreak{}i\allowbreak{}→B\_\allowbreak{}i is an isomorphism. Therefore Q\_\allowbreak{}1,\allowbreak{}Q\_\allowbreak{}2 belong to the strictly full orthogonal,\allowbreak{} and its cone closure gives Q\_\allowbreak{}3\allowbreak{}∈A\textasciicircum{}⊥. The object A\_\allowbreak{}3 is the cone of α. Since A is a full triangulated subcategory,\allowbreak{} there is an object A\_\allowbreak{}3'\allowbreak{}∈A and an isomorphism t:\allowbreak{}A\_\allowbreak{}3'\allowbreak{}→A\_\allowbreak{}3. Transport the third-row triangle along t:\allowbreak{} its first map becomes A\_\allowbreak{}3'\allowbreak{}→\textasciicircum{}tA\_\allowbreak{}3\allowbreak{}→X\_\allowbreak{}3,\allowbreak{} and its last map becomes Q\_\allowbreak{}3\allowbreak{}→A\_\allowbreak{}3[1]\allowbreak{}\allowbreak{}→\textasciicircum{}\{t[1]\allowbreak{}\textasciicircum{}\{-1\}\}A\_\allowbreak{}3'[1]\allowbreak{}. This is precisely a triangle witnessing the source property P(X\_\allowbreak{}3)\allowbreak{},\allowbreak{} without treating A as replete.

The property is invariant under shifts,\allowbreak{} using representatives in A where necessary and the corresponding shift of its triangle with the distinguished-triangle sign convention. Rotating the given triangle therefore proves the other two cases of two-out-of-three. The split triangle X\_\allowbreak{}1\allowbreak{}→X\_\allowbreak{}1\allowbreak{}⊕X\_\allowbreak{}2\allowbreak{}→X\_\allowbreak{}2\allowbreak{}→\textasciicircum{}0X\_\allowbreak{}1[1]\allowbreak{} proves the biproduct case.

For the dual property,\allowbreak{} start with triangles A\_\allowbreak{}i\allowbreak{}→X\_\allowbreak{}i\allowbreak{}→B\_\allowbreak{}i\allowbreak{}→A\_\allowbreak{}i[1]\allowbreak{},\allowbreak{} where B\_\allowbreak{}i\allowbreak{}∈B and A\_\allowbreak{}i\allowbreak{}∈\textasciicircum{}⊥B. The isomorphism Hom(B\_\allowbreak{}1,\allowbreak{}B\_\allowbreak{}2)\allowbreak{}\allowbreak{}→Hom(X\_\allowbreak{}1,\allowbreak{}B\_\allowbreak{}2)\allowbreak{},\allowbreak{} β\allowbreak{}↦βb\_\allowbreak{}1,\allowbreak{} specifies the unique β with βb\_\allowbreak{}1\allowbreak{}=b\_\allowbreak{}2f. Apply the same nine-object construction to the square on X\_\allowbreak{}i\allowbreak{}→B\_\allowbreak{}i; equivalently use the opposite triangulated category,\allowbreak{} whose translation is [−1]\allowbreak{}. Its cone column lies in \^{}⊥B by §1,\allowbreak{} and the B cone is isomorphic to an object of B. Transport along that isomorphism as above. Rotation and the split triangle give all remaining cases. This proves 11813–\allowbreak{}11829 with the actual dual maps.

\subsection{\texorpdfstring{3. Adjoints,\allowbreak{} saturation,\allowbreak{} and double orthogonals}{3.   and double orthogonals}}\label{proofs-11665-12003-adjoints-saturation-and-double-orthogonals}

Let i:\allowbreak{}A\allowbreak{}→D be a full triangulated inclusion. If it has right adjoint v,\allowbreak{} the counit ε\_\allowbreak{}X:\allowbreak{}i vX\allowbreak{}→X induces \[
 \operatorname{Hom}_A(A',vX)
   \xrightarrow{\ \varepsilon_X\circ i(-)\ }
 \operatorname{Hom}_D(iA',X)
\] as the adjunction bijection. Completing ε\_\allowbreak{}X to a triangle and applying §1 puts its cone in A\^{}⊥.

Conversely choose,\allowbreak{} for every X,\allowbreak{} a triangle A\_\allowbreak{}X\allowbreak{}→\textasciicircum{}\{a\_\allowbreak{}X\}X\allowbreak{}→B\_\allowbreak{}X\allowbreak{}→A\_\allowbreak{}X[1]\allowbreak{} as in the source condition (2)\allowbreak{}. Set vX\allowbreak{}=A\_\allowbreak{}X. For f:\allowbreak{}X\allowbreak{}→Y let vf be the unique α with a\_\allowbreak{}Yα\allowbreak{}=fa\_\allowbreak{}X. This uniqueness proves v(id)\allowbreak{}\allowbreak{}=id and v(gf)\allowbreak{}\allowbreak{}=vg vf. The bijections Hom\_\allowbreak{}A(A,\allowbreak{}A\_\allowbreak{}X)\allowbreak{}\allowbreak{}→Hom\_\allowbreak{}D(iA,\allowbreak{}X)\allowbreak{} are induced by a\_\allowbreak{}X,\allowbreak{} and are natural in both A and X by this defining identity. They therefore give the required right adjunction,\allowbreak{} with a\_\allowbreak{}X as counit. The inclusion is fully faithful,\allowbreak{} so the unit A\allowbreak{}→v iA is an isomorphism. Indeed its image under the adjunction is id\_\allowbreak{}\{iA\}; the triangle identities and faithfulness give its two-sided inverse.

If A is strictly full,\allowbreak{} then A⊂\textasciicircum{}⊥(A\textasciicircum{}⊥)\allowbreak{}. For X\allowbreak{}∈\textasciicircum{}⊥(A\textasciicircum{}⊥)\allowbreak{},\allowbreak{} its approximation triangle has b\_\allowbreak{}X\allowbreak{}=0. The split-triangle lemma,\allowbreak{} applied after rotation,\allowbreak{} gives A\_\allowbreak{}X≅X\allowbreak{}⊕B\_\allowbreak{}X[−1]\allowbreak{}. Since B\_\allowbreak{}X[−1]\allowbreak{}\allowbreak{}∈A\textasciicircum{}⊥,\allowbreak{} the projection from A\_\allowbreak{}X to that summand is zero; composing it with the inclusion gives id\_\allowbreak{}\{B\_\allowbreak{}X[−1]\allowbreak{}\}\allowbreak{}=0. Hence B\_\allowbreak{}X\allowbreak{}=0 and a\_\allowbreak{}X is an isomorphism,\allowbreak{} so X\allowbreak{}∈A. Thus A\allowbreak{}=\textasciicircum{}⊥(A\textasciicircum{}⊥)\allowbreak{}. Without strict fullness the same proof gives equality with the strictly full isomorphism closure of A. Saturation in the source\textquotesingle s sense follows:\allowbreak{} any summand X of an object isomorphic to an A object is in \textasciicircum{}⊥(A\textasciicircum{}⊥)\allowbreak{},\allowbreak{} so is isomorphic to an object of A. The other summand is treated by the same argument. No splitting of arbitrary idempotents in D was assumed.

For j:\allowbreak{}B\allowbreak{}→D and a left adjoint w the unit η\_\allowbreak{}X:\allowbreak{}X\allowbreak{}→j wX induces Hom\_\allowbreak{}B(wX,\allowbreak{}B)\allowbreak{}\allowbreak{}→Hom\_\allowbreak{}D(X,\allowbreak{}jB)\allowbreak{},\allowbreak{} β\allowbreak{}↦j(β)\allowbreak{}η\_\allowbreak{}X. Completing η\_\allowbreak{}X backwards gives a triangle A\_\allowbreak{}X\allowbreak{}→X\allowbreak{}→j wX\allowbreak{}→A\_\allowbreak{}X[1]\allowbreak{},\allowbreak{} and §1 places A\_\allowbreak{}X in \^{}⊥B. Conversely,\allowbreak{} such triangles define wX\allowbreak{}=B\_\allowbreak{}X and wf as the unique β with βb\_\allowbreak{}X\allowbreak{}=b\_\allowbreak{}Yf. Uniqueness proves the functor laws and the adjunction naturality. If X\allowbreak{}∈(\textasciicircum{}⊥B)\allowbreak{}\textasciicircum{}⊥,\allowbreak{} then a\_\allowbreak{}X\allowbreak{}=0; the triangle splits as B\_\allowbreak{}X≅X\allowbreak{}⊕A\_\allowbreak{}X[1]\allowbreak{}. The inclusion of this last summand vanishes because Hom(A\_\allowbreak{}X[1]\allowbreak{},\allowbreak{}B\_\allowbreak{}X)\allowbreak{}\allowbreak{}=0,\allowbreak{} so A\_\allowbreak{}X\allowbreak{}=0. Strict fullness gives B\allowbreak{}=(\textasciicircum{}⊥B)\allowbreak{}\textasciicircum{}⊥,\allowbreak{} and without it the equality holds after isomorphism closure. This also proves saturation and the entire omitted dual at 11900–\allowbreak{}11924.

Both adjoints are exact by the full exact-adjoint proof in PROOFS\_\allowbreak{}1844\_\allowbreak{}2401.md §3. Its shift isomorphisms are induced by the actual Hom adjunctions; no assertion of literal equality of chosen objects replaces them.

\subsection{4. The missing hypotheses in the summary proposition}\label{proofs-11665-12003-the-missing-hypotheses-in-the-summary-proposition}

DERIVED-NEW-042 concerns 11954–\allowbreak{}11963. As printed,\allowbreak{} A and B are merely subcategories. Condition (3)\allowbreak{},\allowbreak{} the degree-zero orthogonality and decomposition triangle for every object,\allowbreak{} does not force them to be admissible triangulated subcategories.

An explicit counterexample has D\allowbreak{}=D\textasciicircum{}b(k)\allowbreak{} for a field k,\allowbreak{} A\allowbreak{}=D\textasciicircum{}\{≤0\}(k)\allowbreak{},\allowbreak{} and B\allowbreak{}=D\textasciicircum{}\{≥1\}(k)\allowbreak{},\allowbreak{} both strictly full. Hom\_\allowbreak{}D(A,\allowbreak{}B)\allowbreak{}\allowbreak{}=0:\allowbreak{} use the natural truncation quasi-isomorphisms to a bounded complex supported in degrees ≤0 for A and one supported in degrees ≥1 for B. A bounded complex of vector spaces is a complex of projectives that computes these derived Hom groups; its ordinary chain maps to the second complex are zero by the disjoint degree supports. Transporting through the stated quasi-isomorphisms gives the asserted Hom vanishing for the original A and B.

For the decomposition of an arbitrary bounded complex X,\allowbreak{} use the actual subcomplex A\_\allowbreak{}X with terms X\^{}n for n<0,\allowbreak{} ker(d\_\allowbreak{}X\textasciicircum{}0)\allowbreak{} for n\allowbreak{}=0,\allowbreak{} and zero for n\textgreater0. Its inclusion a\_\allowbreak{}X:\allowbreak{}A\_\allowbreak{}X\allowbreak{}→X has quotient complex B\_\allowbreak{}X\allowbreak{}=X/\allowbreak{}A\_\allowbreak{}X. The quotient has zero terms below zero,\allowbreak{} term X\textasciicircum{}0/\allowbreak{}ker(d\_\allowbreak{}X\textasciicircum{}0)\allowbreak{} in degree zero,\allowbreak{} and terms X\^{}n in degrees n\textgreater0. The degree-zero differential of B\_\allowbreak{}X is injective,\allowbreak{} so H\textasciicircum{}n(B\_\allowbreak{}X)\allowbreak{}\allowbreak{}=0 for n≤0,\allowbreak{} while H\textasciicircum{}n(A\_\allowbreak{}X)\allowbreak{}\allowbreak{}=0 for n\textgreater0. The actual short exact sequence of these complexes gives a distinguished triangle A\_\allowbreak{}X\allowbreak{}→X\allowbreak{}→B\_\allowbreak{}X\allowbreak{}→A\_\allowbreak{}X[1]\allowbreak{}. Thus printed condition (3)\allowbreak{} holds.

But k\allowbreak{}∈A while k[−1]\allowbreak{}\allowbreak{}∉A,\allowbreak{} and k[−1]\allowbreak{}\allowbreak{}∈B while k[−1]\allowbreak{}[1]\allowbreak{}\allowbreak{}=k\allowbreak{}∉B. Neither subcategory is invariant under all shifts,\allowbreak{} so neither is triangulated and (1)\allowbreak{},\allowbreak{}(2)\allowbreak{} both fail. The failure is mathematical,\allowbreak{} not just an omitted word in an otherwise valid proof.

Strict fullness is also not implied by printed (3)\allowbreak{},\allowbreak{} which tests ambient Hom groups and membership of selected objects,\allowbreak{} not the subcategory\textquotesingle s complete morphism sets or isomorphism closure. For example,\allowbreak{} in a nonzero derived category take the subcategory with all objects but only identity and zero maps,\allowbreak{} and let B be its full zero subcategory. The identity triangle X\allowbreak{}→X\allowbreak{}→0 verifies (3)\allowbreak{},\allowbreak{} yet the first subcategory is not full. Hence the conclusion cannot silently restore fullness.

The corrected premise is:\allowbreak{} A and B are \textbf{strictly full subcategories invariant under all shifts}. It is unnecessary to assume their cone closure separately. Here is the complete proof of the missing implication under precisely this premise.

Assume (3)\allowbreak{}. We already have B⊂A\^{}⊥ and A⊂\^{}⊥B. If X\allowbreak{}∈A\textasciicircum{}⊥,\allowbreak{} the first map A\_\allowbreak{}X\allowbreak{}→X of its decomposition triangle vanishes. The split triangle gives B\_\allowbreak{}X≅X\allowbreak{}⊕A\_\allowbreak{}X[1]\allowbreak{}. By shift invariance and (3)\allowbreak{},\allowbreak{} Hom(A\_\allowbreak{}X[1]\allowbreak{},\allowbreak{}B\_\allowbreak{}X)\allowbreak{}\allowbreak{}=0; composing the inclusion of A\_\allowbreak{}X[1]\allowbreak{} with its projection forces its identity to vanish. Thus A\_\allowbreak{}X\allowbreak{}=0 and X≅B\_\allowbreak{}X\allowbreak{}∈B. Since B is strictly full,\allowbreak{} X\allowbreak{}∈B,\allowbreak{} proving B\allowbreak{}=A\textasciicircum{}⊥.

If X\allowbreak{}∈\textasciicircum{}⊥B,\allowbreak{} its map X\allowbreak{}→B\_\allowbreak{}X vanishes,\allowbreak{} hence A\_\allowbreak{}X≅X\allowbreak{}⊕B\_\allowbreak{}X[−1]\allowbreak{}. Shift invariance gives Hom(A\_\allowbreak{}X,\allowbreak{}B\_\allowbreak{}X[−1]\allowbreak{})\allowbreak{}\allowbreak{}=0,\allowbreak{} so the projection to B\_\allowbreak{}X[−1]\allowbreak{} is zero and B\_\allowbreak{}X\allowbreak{}=0. Strict fullness gives X\allowbreak{}∈A. Thus A\allowbreak{}=\textasciicircum{}⊥B. Section 1 now shows that both categories are saturated and triangulated. Section 3 applied to the original decomposition triangles proves both (1)\allowbreak{} and (2)\allowbreak{}. Conversely (1)\allowbreak{} supplies (3)\allowbreak{} by the right-adjoint criterion,\allowbreak{} and (2)\allowbreak{} supplies it by the left-adjoint criterion.

The source patch adds exactly the strict-fullness and shift hypotheses and supplies this implication. It does not assume a missing conclusion in place of proving it.

\subsection{5. Functorial triangles and the exact morphism obstruction}\label{proofs-11665-12003-functorial-triangles-and-the-exact-morphism-obstruction}

DERIVED-UNDER-031 strengthens the statement at 11938–\allowbreak{}11950. For an admissible pair let \[
 A_X\xrightarrow{a_X}X\xrightarrow{b_X}B_X
 \xrightarrow{\delta_X}A_X[1]
\] be the selected triangles,\allowbreak{} with A\_\allowbreak{}X\allowbreak{}∈A,\allowbreak{} B\_\allowbreak{}X\allowbreak{}∈B\allowbreak{}=A\textasciicircum{}⊥. For any f:\allowbreak{}X\allowbreak{}→Y,\allowbreak{} the adjunction gives a unique α\_\allowbreak{}f:\allowbreak{}A\_\allowbreak{}X\allowbreak{}→A\_\allowbreak{}Y with a\_\allowbreak{}Yα\_\allowbreak{}f\allowbreak{}=fa\_\allowbreak{}X. There is also a unique β\_\allowbreak{}f:\allowbreak{}B\_\allowbreak{}X\allowbreak{}→B\_\allowbreak{}Y with β\_\allowbreak{}fb\_\allowbreak{}X\allowbreak{}=b\_\allowbreak{}Yf:\allowbreak{} the relevant exact sequence is \[
 \operatorname{Hom}(A_X[1],B_Y)\longrightarrow
 \operatorname{Hom}(B_X,B_Y)\xrightarrow{b_X^*}
 \operatorname{Hom}(X,B_Y)\longrightarrow\operatorname{Hom}(A_X,B_Y),
\] whose outer terms are zero. The triangle morphism axiom gives a completion of the first square,\allowbreak{} and uniqueness identifies its third map with β\_\allowbreak{}f. Consequently \[
 \delta_Y\beta_f=\alpha_f[1]\delta_X.
\] The unique choices respect identities and composition. In particular,\allowbreak{} two selected triangles on the same X have a \textbf{unique} isomorphism inducing id\_\allowbreak{}X:\allowbreak{} the unique maps in both directions compose to identities by the same uniqueness. These constructions give a natural triangle of functors i v\allowbreak{}→Id\_\allowbreak{}D\allowbreak{}→j w\allowbreak{}→(i v)\allowbreak{}[1]\allowbreak{},\allowbreak{} with the actual exact-adjoint shift identifications. The conclusion is stronger than an unspecified isomorphism between the third objects of two cones.

There is further exact information about arbitrary maps X\allowbreak{}→Y. Using the original a,\allowbreak{}b,\allowbreak{}δ,\allowbreak{} define \[
\begin{aligned}
 d_{-1}(k,l)&=k\delta_X-\delta_Y[-1]\,l,\\
 \phi(h)&=a_Yh b_X,\\
 \rho(f)&=(\alpha_f,\beta_f),\\
 d_0(\alpha,\beta)&=\alpha[1]\delta_X-\delta_Y\beta.
\end{aligned}
\] The domains and codomains form the following sequence,\allowbreak{} exact at each of its three middle terms:\allowbreak{} \[
\begin{gathered}
 \operatorname{Hom}(A_X[1],A_Y)\oplus
 \operatorname{Hom}(B_X,B_Y[-1])
 \xrightarrow{d_{-1}}\operatorname{Hom}(B_X,A_Y)
 \xrightarrow{\phi}\operatorname{Hom}(X,Y)\\
 \xrightarrow{\rho}
 \operatorname{Hom}(A_X,A_Y)\oplus\operatorname{Hom}(B_X,B_Y)
 \xrightarrow{d_0}\operatorname{Hom}(B_X,A_Y[1]).
 \tag{5.1}
\end{gathered}
\] No surjectivity onto the last group is asserted.

For exactness at the pair of Hom groups,\allowbreak{} the displayed triangle morphism identity gives d\_\allowbreak{}0ρ\allowbreak{}=0. Conversely if α[1]\allowbreak{}δ\_\allowbreak{}X\allowbreak{}=δ\_\allowbreak{}Yβ,\allowbreak{} rotate each triangle twice,\allowbreak{} obtaining B\_\allowbreak{}X\allowbreak{}→\textasciicircum{}\{δ\_\allowbreak{}X\}A\_\allowbreak{}X[1]\allowbreak{}\allowbreak{}→\textasciicircum{}\{−a\_\allowbreak{}X[1]\allowbreak{}\}X[1]\allowbreak{}\allowbreak{}→\textasciicircum{}\{−b\_\allowbreak{}X[1]\allowbreak{}\}B\_\allowbreak{}X[1]\allowbreak{} and its Y counterpart. The first square,\allowbreak{} with maps β and α[1]\allowbreak{},\allowbreak{} commutes. The triangle morphism axiom gives a map f[1]\allowbreak{} on the third objects. Its two squares contain the same minus signs on both sides,\allowbreak{} so shifting back gives fa\_\allowbreak{}X\allowbreak{}=a\_\allowbreak{}Yα and b\_\allowbreak{}Yf\allowbreak{}=βb\_\allowbreak{}X. The uniqueness above identifies ρ(f)\allowbreak{}\allowbreak{}=(α,\allowbreak{}β)\allowbreak{}.

If ρ(f)\allowbreak{}\allowbreak{}=0,\allowbreak{} then fa\_\allowbreak{}X\allowbreak{}=0. Exactness of Hom(−,\allowbreak{}Y)\allowbreak{} gives g:\allowbreak{}B\_\allowbreak{}X\allowbreak{}→Y with f\allowbreak{}=g b\_\allowbreak{}X. Since b\_\allowbreak{}Yf\allowbreak{}=0 and Hom(B\_\allowbreak{}X,\allowbreak{}B\_\allowbreak{}Y)\allowbreak{}\allowbreak{}→Hom(X,\allowbreak{}B\_\allowbreak{}Y)\allowbreak{} is injective,\allowbreak{} b\_\allowbreak{}Yg\allowbreak{}=0. The Hom(B\_\allowbreak{}X,\allowbreak{}−)\allowbreak{} sequence for Y gives h:\allowbreak{}B\_\allowbreak{}X\allowbreak{}→A\_\allowbreak{}Y with g\allowbreak{}=a\_\allowbreak{}Yh. Thus f\allowbreak{}=φ(h)\allowbreak{}. Conversely b\_\allowbreak{}Xa\_\allowbreak{}X\allowbreak{}=0 and b\_\allowbreak{}Ya\_\allowbreak{}Y\allowbreak{}=0 show ρφ\allowbreak{}=0 by the defining uniqueness of α and β.

Finally φ(h)\allowbreak{}\allowbreak{}=0 means a\_\allowbreak{}Yh b\_\allowbreak{}X\allowbreak{}=0. The Hom(−,\allowbreak{}Y)\allowbreak{} sequence gives r:\allowbreak{}A\_\allowbreak{}X[1]\allowbreak{}\allowbreak{}→Y with a\_\allowbreak{}Yh\allowbreak{}=rδ\_\allowbreak{}X. Since A\_\allowbreak{}X[1]\allowbreak{}\allowbreak{}∈A,\allowbreak{} its adjunction bijection writes r uniquely as a\_\allowbreak{}Yk for k:\allowbreak{}A\_\allowbreak{}X[1]\allowbreak{}\allowbreak{}→A\_\allowbreak{}Y. Hence a\_\allowbreak{}Y(h−kδ\_\allowbreak{}X)\allowbreak{}\allowbreak{}=0. The inverse-rotated triangle B\_\allowbreak{}Y[−1]\allowbreak{}\allowbreak{}→\textasciicircum{}\{−δ\_\allowbreak{}Y[-1]\allowbreak{}\}A\_\allowbreak{}Y\allowbreak{}→\textasciicircum{}\{a\_\allowbreak{}Y\}Y\allowbreak{}→B\_\allowbreak{}Y gives l:\allowbreak{}B\_\allowbreak{}X\allowbreak{}→B\_\allowbreak{}Y[−1]\allowbreak{} with h−kδ\_\allowbreak{}X\allowbreak{}=−δ\_\allowbreak{}Y[−1]\allowbreak{}l. Therefore h\allowbreak{}=d\_\allowbreak{}\{−1\}(k,\allowbreak{}l)\allowbreak{}. Conversely the triangle composites δ\_\allowbreak{}Xb\_\allowbreak{}X\allowbreak{}=0 and a\_\allowbreak{}Yδ\_\allowbreak{}Y[−1]\allowbreak{}\allowbreak{}=0 give φd\_\allowbreak{}\{−1\}\allowbreak{}=0. This proves every asserted exactness and keeps the rotation sign.

Consequently a compatible pair (α,\allowbreak{}β)\allowbreak{},\allowbreak{} namely d\_\allowbreak{}0(α,\allowbreak{}β)\allowbreak{}\allowbreak{}=0,\allowbreak{} lifts to a map X\allowbreak{}→Y. The set of its lifts is a torsor under the explicitly defined quotient group \[
 \operatorname{Hom}(B_X,A_Y)/\operatorname{Im}(d_{-1}),
\] acting by adding a\_\allowbreak{}Yh b\_\allowbreak{}X. Surjectivity of the action is exactness at Hom(X,\allowbreak{}Y)\allowbreak{},\allowbreak{} and freeness is exactness at Hom(B\_\allowbreak{}X,\allowbreak{}A\_\allowbreak{}Y)\allowbreak{}. Thus the connecting maps δ\_\allowbreak{}X,\allowbreak{}δ\_\allowbreak{}Y determine both the obstruction and the exact ambiguity; one may not replace the category by a direct product merely because each object has two projections. This is an editorial consequence of the original triangles,\allowbreak{} with no novelty claim and no change to their objects or maps.

\subsection{6. The quotient equivalences and their natural isomorphisms}\label{proofs-11665-12003-the-quotient-equivalences-and-their-natural-isomorphisms}

Keep i:\allowbreak{}A\allowbreak{}→D and its exact right adjoint v,\allowbreak{} and j:\allowbreak{}B\allowbreak{}→D and its exact left adjoint w. Then ker v\allowbreak{}=B:\allowbreak{} if vX\allowbreak{}=0 the adjunction makes Hom(A,\allowbreak{}X)\allowbreak{}\allowbreak{}=0 for all A; if X\allowbreak{}∈B then Hom\_\allowbreak{}A(A,\allowbreak{}vX)\allowbreak{}\allowbreak{}=0 for all A,\allowbreak{} so taking A\allowbreak{}=vX shows id\_\allowbreak{}\{vX\}\allowbreak{}=0. Dually ker w\allowbreak{}=A.

Let q\_\allowbreak{}B:\allowbreak{}D\allowbreak{}→D/\allowbreak{}B be the actual Verdier quotient. Since v is exact and vanishes on B,\allowbreak{} its value on every morphism whose cone lies in B is an isomorphism,\allowbreak{} by its triangle exactness. Thus the quotient universal property defines an exact functor \(\bar v:D/B\to A\) with \(\bar v q_B\cong v\). The adjunction unit gives \(\bar v q_B i\cong v i\cong Id_A\). The counit ε\_\allowbreak{}X\allowbreak{}=i vX\allowbreak{}→X has cone B\_\allowbreak{}X\allowbreak{}∈B,\allowbreak{} so q\_\allowbreak{}B(ε\_\allowbreak{}X)\allowbreak{}:\allowbreak{}q\_\allowbreak{}B i vX\allowbreak{}→q\_\allowbreak{}BX is an isomorphism. These isomorphisms are natural for the original arrows by the counit identity; they are natural for inverses of quotient denominators by multiplying that identity by their inverses. They are therefore natural for every roof in D/\allowbreak{}B. This supplies \(q_B i\bar v\cong Id_{D/B}\),\allowbreak{} the other quasi-inverse comparison.

For q\_\allowbreak{}A:\allowbreak{}D\allowbreak{}→D/\allowbreak{}A,\allowbreak{} the exact left adjoint gives \(\bar w:D/A\to B\) with \(\bar w q_A\cong w\). Its counit gives \(\bar w q_A j\cong wj\cong Id_B\). The original unit b\_\allowbreak{}X:\allowbreak{}X\allowbreak{}→j wX has cone A\_\allowbreak{}X[1]\allowbreak{}\allowbreak{}∈A,\allowbreak{} so its quotient isomorphism gives \(Id_{D/A}\cong q_A j\bar w\). Naturality for roofs follows by the same inverse calculation. Thus both quotient equivalences in 11965–\allowbreak{}12003,\allowbreak{} and both descriptions of the adjoints as quotient composites,\allowbreak{} are proved with their actual comparison maps. Object choices give natural isomorphisms,\allowbreak{} not an unproved assertion that independently chosen objects are literally equal.

\subsection{7. Existing reports and corrected receiving quantifiers}\label{proofs-11665-12003-existing-reports-and-corrected-receiving-quantifiers}

P02-E0095 /\allowbreak{} DERIVED-088 report the spelling at original 11744. MC-STK-ERR-0899-OP1 already makes exactly that correction. P02-E0096 /\allowbreak{} DERIVED-089 report the spelling at original 11890; MC-STK-ERR-0900-OP1 already corrects it. Their four physical reports form two reconciliation groups. No additional operation is needed for those spellings.

The source-order receiver search also found the following original reports,\allowbreak{} retained in ADJOINT\_\allowbreak{}RECEIVER\_\allowbreak{}ORIGINAL\_\allowbreak{}REPORTS\_\allowbreak{}PART21.json:\allowbreak{}

\begin{itemize}
\tightlist
\item
  P05-EN-CH36-0056,\allowbreak{} OCC-05101,\allowbreak{} perfect.tex 4848–\allowbreak{}4856.
\item
  P09-ERR-0448,\allowbreak{} OCC-08627,\allowbreak{} spaces-perfect.tex 4459–\allowbreak{}4467.
\item
  P09-ERR-1013,\allowbreak{} OCC-09189,\allowbreak{} spaces-perfect.tex 4459–\allowbreak{}4466.
\end{itemize}

The last two are distinct physical reports of the same source defect. Their original producer/\allowbreak{}ledger identities are preserved,\allowbreak{} and they share one receiving correction group. These reports are separate from the 176-report Derived denominator.

Both receiving passages assert a triangle \[
 E'\longrightarrow E\longrightarrow K\longrightarrow E'[1],
 \qquad E'\in D_{\rm QCoh}(\mathcal O_X),\quad
 \operatorname{Hom}(M,K)=0\quad(M\in D_{\rm QCoh}(\mathcal O_X)).
\] Their universal quantifier incorrectly says ``for every object K''. The arbitrary input in §3 is the middle object,\allowbreak{} so it must say ``for every object E''. K is the cone that is required to lie in the right orthogonal. The subsequent source argument already uses E as its input,\allowbreak{} so changing precisely that quantified letter preserves all following formulas.

The printed statement is demonstrably false as a criterion. Take X\allowbreak{}=Spec(k)\allowbreak{},\allowbreak{} as either a scheme or an algebraic space. The inclusion D\_\allowbreak{}QCoh(O\_\allowbreak{}X)\allowbreak{}\allowbreak{}→D(O\_\allowbreak{}X)\allowbreak{} is the identity and has a right adjoint. But demanding the displayed orthogonality for every K includes K\allowbreak{}=k and M\allowbreak{}=k,\allowbreak{} for which Hom(k,\allowbreak{}k)\allowbreak{} contains id\_\allowbreak{}k≠0. This validates the producer reports mathematically,\allowbreak{} rather than matching their loci only.

For the corrected scheme argument,\allowbreak{} the source\textquotesingle s Mayer–\allowbreak{}Vietoris triangle has middle-input objects Rj\_\allowbreak{}\{U,\allowbreak{}*\}E|\_\allowbreak{}U,\allowbreak{} Rj\_\allowbreak{}\{V,\allowbreak{}*\}E|\_\allowbreak{}V,\allowbreak{} Rj\_\allowbreak{}\{U∩V,\allowbreak{}*\}E|\_\allowbreak{}\{U∩V\}. The preparation lemma in §2 propagates approximation existence through that triangle and its finite sum. For one of these open maps,\allowbreak{} start with the triangle L'\allowbreak{}→L\allowbreak{}→Q\allowbreak{}→L'[1]\allowbreak{} on U. Applying the actual exact Rj\_\allowbreak{}* gives Rj\_\allowbreak{}*L'\allowbreak{}→Rj\_\allowbreak{}*L\allowbreak{}→Rj\_\allowbreak{}*Q\allowbreak{}→Rj\_\allowbreak{}*L'[1]\allowbreak{}. The cited quasi-coherent pushforward result puts Rj\_\allowbreak{}*L' in the required subcategory. Its cone lies in the orthogonal because \[
 \operatorname{Hom}_X(M,Rj_*Q)
   \cong\operatorname{Hom}_U(Lj^*M,Q)=0,
\] with the actual pullback preserving the quasi-coherent category. The algebraic-space passage uses the same maps for the stated étale morphisms and the elementary distinguished-square triangle; its pullback–\allowbreak{}pushforward adjunction is the same precise Hom comparison. The quantifier correction therefore repairs both receiving uses. Their separate geometric dependencies remain in their own chapter reviews; this bounded check does not certify those entire chapters.

\subsection{8. Other receiving uses and limits of this review}\label{proofs-11665-12003-other-receiving-uses-and-limits-of-this-review}

In stacks-perfect.tex 861–\allowbreak{}934 the kernel category \[
 B=D_{\textit{Parasitic}\cap\textit{LQCoh}^{fbc}}(\mathcal O_X)
 \subset D_{\textit{LQCoh}^{fbc}}(\mathcal O_X)
\] is the stated full derived subcategory and is invariant under all shifts. Its left orthogonal A is strictly full and shift invariant by §1. Thus the pair in the source remark satisfies the repaired premise of DERIVED-NEW-042. The preceding lemma provides E'\allowbreak{}=Lg\_\allowbreak{}!g\textasciicircum{}*E\allowbreak{}→E with cone P\allowbreak{}∈B and the actual comparison \[
 \operatorname{Hom}(Lg_!g^*E,P')
  \cong\operatorname{Hom}(g^*E,g^*P')=0.
\] It supplies the original decomposition and orthogonality,\allowbreak{} so the admissibility and quotient conclusions follow from §§3,\allowbreak{}4,\allowbreak{}6 with no new change to this receiver. Section 5 further proves that its canonical triangle is unique up to a unique isomorphism fixing E and is functorial in E. That strengthening is recorded editorially. The complete independent review of the preceding stack/\allowbreak{}topos comparison lemmas remains in their own chapter queues.

In equiv.tex 2821–\allowbreak{}2864 the essential image A of the fully faithful exact F is taken with all its ambient morphisms and all objects isomorphic to its image; it is strictly full. Exactness preserves shifts. Full faithfulness lifts every map between image objects to the source; applying F to its source cone proves cone closure up to the prescribed isomorphism. Thus A is triangulated. If F\_\allowbreak{}r is its right adjoint as cited there,\allowbreak{} the counit F F\_\allowbreak{}rN\allowbreak{}→N gives the right approximation:\allowbreak{} for a source object M,\allowbreak{} adjunction and the fully faithful unit identify Hom(FM,\allowbreak{}F F\_\allowbreak{}rN)\allowbreak{}\allowbreak{}→Hom(FM,\allowbreak{}N)\allowbreak{} with an isomorphism. Section 3 gives the right admissibility and A\allowbreak{}=\textasciicircum{}⊥(A\textasciicircum{}⊥)\allowbreak{} used by the receiver. Its later point-sheaf detection argument is not a new condition needed for this categorical step. It remains subject to its own source review.

The root-chapter reference search was bounded to explicit references to this section\textquotesingle s labels; it is not a claim of exhaustive semantic dependency discovery. Later Postnikov source was read ahead through 12150 only,\allowbreak{} and is not covered by this completed review.
\clearpage
\section{\texorpdfstring{Derived Categories:\allowbreak{} Postnikov systems,\allowbreak{} original lines 12008–\allowbreak{}12370}{Derived  Postnikov  original lines }}\label{proofs-12008-12370-derived-categories-postnikov-systems-original-lines-1200812370}

This is a private mathematical review and editorial proof supplement. Its source is the Stacks Project at commit a04446e5\allowbreak{}7ec1fbc2\allowbreak{}52a871af\allowbreak{}cec7752f\allowbreak{}b2807b14, original derived.tex,\allowbreak{} SHA-256 EFDA4CB3\allowbreak{}61DD4090\allowbreak{}9D199116\allowbreak{}1EAE716E\allowbreak{}58B6CF9E\allowbreak{}79380971\allowbreak{}502E1DA1\allowbreak{}2396D402. The current preimage is the separately frozen R60 chapter. The original source statement remains identifiable. No new theorem is attributed to the original author,\allowbreak{} and no novelty is claimed. Section 4 proves DERIVED-NEW-043; Section 5 proves DERIVED-UNDER-032. Section 6 adjudicates the three physical reports in this part.

\subsection{\texorpdfstring{1. Original objects,\allowbreak{} triangle maps,\allowbreak{} and the complex example}{1. Original  triangle  and the complex example}}\label{proofs-12008-12370-original-objects-triangle-maps-and-the-complex-example}

Write the given differential as \(d_j:X_j\to X_{j-1}\). A system has an isomorphism \(p_0:Y_0\to X_0\) and distinguished triangles \[
Y_j\xrightarrow{p_j}X_j\xrightarrow{t_j}Y_{j-1}
 \xrightarrow{\delta_j}Y_j[1],\qquad p_{j-1}t_j=d_j. \tag{1.1}
\] A map over \(f_j:X_j\to X'_j\) is a family \(\phi_j:Y_j\to Y'_j\) with \[
p'_j\phi_j=f_jp_j,\qquad
\phi_{j-1}t_j=t'_jf_j,\qquad
\phi_j[1]\delta_j=\delta'_j\phi_{j-1}, \tag{1.2}
\] and \(p'_0\phi_0=f_0p_0\). The third equation,\allowbreak{} including its shift,\allowbreak{} is essential.

Retain the original cochain convention \(M[s]^q=M^{q+s}\),\allowbreak{} \(d_{M[s]}=(-1)^s d_M\),\allowbreak{} and the source rotation convention \((a,b,c)\mapsto(b,c,-a[1])\). Put \(C_j^{-i}=A_i\) for \(0\leq i\leq j\),\allowbreak{} zero otherwise,\allowbreak{} with the original \(d_i:A_i\to A_{i-1}\). Thus the original objects are \(X_j=A_j\) and \(Y_j=C_j[-j]\),\allowbreak{} with differential \((-1)^j d_i\); none of these differentials is discarded.

Let \(\iota_j:C_{j-1}\to C_j\) be the literal inclusion and \(q_j:C_j\to A_j[j]\) the literal projection. The degreewise split exact sequence \[
0\longrightarrow C_{j-1}\xrightarrow{\iota_j}C_j
 \xrightarrow{q_j}A_j[j]\longrightarrow0
\] has connecting map \(e_j:A_j[j]\to C_{j-1}[1]\) whose only nonzero component is \(d_j\) in degree \(-j\). Indeed,\allowbreak{} the source definition of this connecting map is \(\pi^{q+1}d^q s^q\) (original derived.tex 2819–\allowbreak{}2848; Homology\textquotesingle s termwise-split cochain construction)\allowbreak{}. After shifting by \(-j\),\allowbreak{} the connecting map is \((-1)^j e_j[-j]\). Rotating once gives the distinguished triangle with arrows \[
Y_j\xrightarrow{q_j[-j]}A_j
 \xrightarrow{(-1)^j d_j}Y_{j-1}
 \xrightarrow{-\iota_j[-j+1]}Y_j[1]. \tag{1.3}
\] Here a displayed map \(d_j:A_j\to Y_{j-1}\) has degree-zero component \(d_j\) and all other components zero. It is a chain map because \(d_{j-1}d_j=0\).

For explicit coherent choices with literal inclusion maps in the tower,\allowbreak{} set \[
c_j=(-1)^{j(j-1)/2},\qquad
p_j=c_j q_j[-j],\qquad t_j=c_{j-1}d_j,\qquad
\delta_j=\iota_j[-j+1].
\] Take \(p_0=1\). Since \(c_j=(-1)^{j+1}c_{j-1}\),\allowbreak{} conjugating (1.3)\allowbreak{} by the three isomorphisms \(1,c_j,-1\) gives exactly these three arrows. Thus each triangle is distinguished. Also \(p_{j-1}t_j=c_{j-1}^2d_j=d_j\),\allowbreak{} and \(\delta_j[j-1]=\iota_j\). This proves the example with all its original objects and differentials. The phrase ``evident canonical maps'' needs a sign convention; the displayed construction supplies one. An alternative retaining every literal projection \(p_j\) has \(t_j=d_j\) and \(\delta_j=(-1)^{j+1}\iota_j[-j+1]\). These two conventions must not be mixed. This is recorded as an explicit clarification,\allowbreak{} not an additional false theorem.

For any system,\allowbreak{} put \(Z_j=Y_j[j]\),\allowbreak{} \(u_j=\delta_j[j-1]\). Shifting (1.1)\allowbreak{},\allowbreak{} rotating twice and conjugating by the object signs gives the distinguished triangle \[
Z_{j-1}\xrightarrow{u_j}Z_j\xrightarrow{p_j[j]}X_j[j]
 \xrightarrow{(-1)^{j-1}t_j[j]}Z_{j-1}[1]. \tag{1.4}
\] For example,\allowbreak{} before conjugation the three arrows are \((-1)^{j-1}\delta_j[j-1],-p_j[j],-t_j[j]\). Conjugation on its three objects by \((-1)^{j-1},1,-1\) gives (1.4)\allowbreak{}. This retains the signs in every use of the tower below.

Under the source\textquotesingle s exactness and existence assumption for colimits over \(\mathbf N\),\allowbreak{} the colimit of the literal complexes \(C_j\) is the original complex with \(A_i\) in degree \(-i\). The sequential-colimit/\allowbreak{}homotopy-colimit comparison proved in PROOFS\_\allowbreak{}10139\_\allowbreak{}10427.md applies to the actual inclusion maps \(\iota_j\). It therefore identifies \(\operatorname{hocolim}Y_j[j]\) with that complex. No assertion that \(\operatorname{hocolim}Y_j\) has the same value is made.

If every \(d_j=0\),\allowbreak{} the split construction is \(Y_j=\bigoplus_{i=0}^jX_i[i-j]\). The map \(p_j\) is the projection to \(X_j\),\allowbreak{} \(t_j=0\),\allowbreak{} and \(\delta_j\) is the inclusion of the other summands into \(Y_j[1]\),\allowbreak{} with the sign of that summand chosen to make the triangle distinguished. Then \(Y_j[j]=\bigoplus_{i=0}^jX_i[i]\). This proves precisely the assertion ``we can take'' in original 12117–\allowbreak{}12118; it does not say every system for the zero complex splits in this way.

\subsection{\texorpdfstring{2. Small lengths,\allowbreak{} with actual failure examples}{2. Small  with actual failure examples}}\label{proofs-12008-12370-small-lengths-with-actual-failure-examples}

For length zero choose \(Y_0=X_0\). For two choices the forced comparison is \((p'_0)^{-1}f_0p_0\). For length one put \(t_1=p_0^{-1}d_1\),\allowbreak{} complete it to a triangle and rotate backwards. TR3,\allowbreak{} applied to the commuting square on \(X_1,Y_0\),\allowbreak{} supplies a comparison; rotation gives (1.2)\allowbreak{}.

Nonuniqueness already occurs at length one. In \(D^b(k)\),\allowbreak{} let \(K=k[0]\),\allowbreak{} \(V=K\oplus K\),\allowbreak{} \(X_0=Y_0=K[1]\),\allowbreak{} \(X_1=K\),\allowbreak{} \(d_1=0\),\allowbreak{} \(Y_1=V\). Use \[
V\xrightarrow{\operatorname{pr}_1}K\xrightarrow{0}K[1]
 \xrightarrow{\operatorname{inc}_2[1]}V[1]. \tag{2.1}
\] This is distinguished:\allowbreak{} rotate the split triangle \(K\xrightarrow{\operatorname{inc}_2}V\xrightarrow{\operatorname{pr}_1}K\to K[1]\),\allowbreak{} and then change the sign of its third object. If \(h(x,y)=(0,x)\),\allowbreak{} then \(\operatorname{pr}_1h=0\) and \(h[1]\operatorname{inc}_2[1]=0\). Thus both \((\phi_0,\phi_1)=(0,0)\) and \((0,h)\) are maps over the zero map of the underlying complex. The map \(h\) is nonzero in the derived category because it is nonzero on degree-zero cohomology.

For length two,\allowbreak{} given the first triangle,\allowbreak{} exactness of \[
\operatorname{Hom}(X_2,Y_1)\longrightarrow
\operatorname{Hom}(X_2,X_1)\longrightarrow
\operatorname{Hom}(X_2,Y_0)
\] lifts \(d_2\) to \(t_2\),\allowbreak{} since \(p_0t_1d_2=d_1d_2=0\) and \(p_0\) is invertible. Complete \(t_2\) to a triangle. This proves existence for all such complexes.

Comparisons need not exist at length two. In \(D^b(k)\),\allowbreak{} use the common complex \(K\to0\to K[1]\),\allowbreak{} with zero arrows,\allowbreak{} and \(Y_0=K[1]\),\allowbreak{} \(Y_1=K\),\allowbreak{} with first triangle \(K\to0\to K[1]\xrightarrow{1}K[1]\). This triangle is distinguished by a sign change in a rotation of an identity triangle. For one system choose \(t_2=0:K\to K\),\allowbreak{} with fiber object \(Y_2=K\oplus K[-1]\). For the other choose \(t'_2=1_K\),\allowbreak{} with \(Y'_2=0\). Any comparison over the identity of the complex would have \(\phi_0=1\),\allowbreak{} and its compatibility with the first connecting map forces \(\phi_1=1\). The stage-two equation would then read \(1\circ0=1\circ1\),\allowbreak{} a contradiction.

There is an actual nonexistence example at length three,\allowbreak{} not merely a failure of the displayed construction. In \(D^b(\mathbf{Ab})\),\allowbreak{} use the nonsplit exact sequence \[
0\to A=\mathbf Z\xrightarrow{2}B=\mathbf Z
 \xrightarrow{q}C=\mathbf Z/2\to0
\] with connecting map \(\epsilon:C\to A[1]\). Consider \[
X_3=A\xrightarrow{2}X_2=B\xrightarrow{q}X_1=C
 \xrightarrow{\epsilon}X_0=A[1]. \tag{2.2}
\] Consecutive composites vanish because these are consecutive arrows of a distinguished triangle. Normalize \(p_0\) only by the explicit isomorphism \(Y_0\to X_0\),\allowbreak{} so no underlying object or map is replaced without comparison. Every first-stage triangle for \(\epsilon\) is isomorphic,\allowbreak{} identically on \(C,A[1]\),\allowbreak{} to \[
B\xrightarrow{q}C\xrightarrow{\epsilon}A[1]
 \xrightarrow{-2[1]}B[1].
\] This follows from TR3 and the two-isomorphisms property,\allowbreak{} applied after rotation. Transport along that isomorphism identifies any second-stage lift with an endomorphism \(a:B\to B\) satisfying \(qa=q\). Since \(\operatorname{End}_{D^b(\mathbf{Ab})}(\mathbf Z)=\mathbf Z\),\allowbreak{} this is multiplication by an odd integer. The third-stage obstruction is \(a\circ2\),\allowbreak{} which is nonzero. Exactness of Hom applied to the second-stage triangle says a third-stage lift would force it to be zero. Thus no choices give a system. For \(n>3\) prepend zero objects to (2.2)\allowbreak{}; a system would restrict to the impossible length-three system.

In the printed proof at original 12160 the relevant composite is consequently \(X_n\to X_{n-1}\to Y_{n-2}\). The last arrow is \(t_{n-1}\),\allowbreak{} not an arrow to \(Y_{n-1}\). This confirms P02-E0097,\allowbreak{} while correcting the report\textquotesingle s imprecise description of which triangle is being used.

\subsection{3. The two vanishing hypotheses and their exact scope}\label{proofs-12008-12370-the-two-vanishing-hypotheses-and-their-exact-scope}

Write the map-extension condition as \[
(P)\qquad \operatorname{Hom}(X_i[i-j-1],X'_j)=0
 \quad(i>j+1).
\] For any target system it implies \(\operatorname{Hom}(X_i[i-j-1],Y'_j)=0\). Induct on \(j\). The case zero uses \(p'_0\). For \(j>0\),\allowbreak{} the inverse rotation of its triangle gives the exact segment with outer groups \[
\operatorname{Hom}(X_i[i-j-1],Y'_{j-1}[-1])
 =\operatorname{Hom}(X_i[i-j],Y'_{j-1}),\quad
\operatorname{Hom}(X_i[i-j-1],X'_j).
\] The first vanishes by induction (the required inequality \(i>j\) follows)\allowbreak{},\allowbreak{} and the second by \(P\). Hence the middle group vanishes.

Given comparisons through \(n-1\),\allowbreak{} define \(\alpha=t'_nf_n-\phi_{n-1}t_n:X_n\to Y'_{n-1}\). The chain-map and earlier-square equations give \[
p'_{n-1}\alpha=d'_nf_n-f_{n-1}d_n=0.
\] Therefore \(\alpha=-\delta'_{n-1}[-1]z\) for some \(z:X_n\to Y'_{n-2}[-1]\). Its domain group is \(\operatorname{Hom}(X_n[1],Y'_{n-2})=0\) by the preceding conclusion. Thus \(\alpha=0\),\allowbreak{} and TR3 extends the comparison. The cases zero and one are Section 2. Notice that this argument extends any already chosen partial comparison:\allowbreak{} induction gives a compatible infinite comparison too,\allowbreak{} rather than a disconnected family of finite comparisons.

For existence,\allowbreak{} the source assumes the different condition \[
(P')\qquad \operatorname{Hom}(X_i[i-j-2],X_j)=0
 \quad(i>j+2).
\] Induction on \(j\),\allowbreak{} using the same exact sequence and shifted source,\allowbreak{} gives \(\operatorname{Hom}(X_i[i-j-2],Y_j)=0\) for \(i>j+2\). For \(n>2\),\allowbreak{} the obstruction \(o=t_{n-1}d_n:X_n\to Y_{n-2}\) satisfies \(p_{n-2}o=d_{n-1}d_n=0\). It factors through \(Y_{n-3}[-1]\),\allowbreak{} but \(\operatorname{Hom}(X_n,Y_{n-3}[-1])
 =\operatorname{Hom}(X_n[1],Y_{n-3})=0\) by the just proved conclusion. Thus \(o=0\),\allowbreak{} and the exact sequence lifts \(d_n\). Lengths zero through two are already proved. This induction also extends every partial system and hence constructs an infinite system stage by stage when \(P'\) holds for every index.

When \(P\) holds for a complex compared with itself,\allowbreak{} any two systems have a comparison over the identity. Its degree-zero map is an isomorphism. At every subsequent stage two of the three component maps are isomorphisms; exact Hom and the five lemma show the third is an isomorphism. Their inverses satisfy the inverse commuting squares. Thus the systems are isomorphic. No uniqueness of this isomorphism has been asserted.

The last statement must not be read as saying \(P\) alone implies existence. Here is a counterexample to that reading. In the abelian category of representations of the quiver \(1\to2\) over \(k\),\allowbreak{} let \[
A=(0\to k),\quad B=(k\xrightarrow{1}k),\quad C=(k\to0).
\] Let \(f:A\to B\) be inclusion at vertex two and \(g:B\to C\) projection at vertex one. They give a nonsplit exact sequence. Indeed,\allowbreak{} a section \(C\to B\) would be identity at vertex one and zero at vertex two,\allowbreak{} contrary to the arrow-commuting equation. Let \(\epsilon:C\to A[1]\) be its connecting map,\allowbreak{} and use \(A\to B\to C\to A[1]\) as in (2.2)\allowbreak{}. Every map \(B\to A\) is zero:\allowbreak{} its first component is zero,\allowbreak{} and commutativity forces its second component to be zero. Thus 
\begingroup\fontsize{9.5}{11.5}\selectfont
\[
\operatorname{Hom}(X_2[1],X_0)=\operatorname{Hom}(B,A)=0,\quad
\operatorname{Hom}(X_3[2],X_0)=\operatorname{Ext}^{-1}(A,A)=0,\quad
\operatorname{Hom}(X_3[1],X_1)=\operatorname{Ext}^{-1}(A,C)=0.
\]
\endgroup
 These are all the groups in \(P\); negative Ext between degree-zero objects is zero by the standard derived t-structure. After comparison of the first triangle,\allowbreak{} a second-stage lift \(a:B\to B\) satisfies \(ga=g\). Exact Hom gives \(a-1=fv\) for \(v:B\to A\),\allowbreak{} so \(a=1\). The next obstruction is \(af=f\ne0\). No system exists. In contrast,\allowbreak{} the group in \(P'\) at \((i,j)=(3,0)\) is \(\operatorname{Hom}(A[1],A[1])=k\),\allowbreak{} so the actual existence hypothesis fails. The source\textquotesingle s two separate assertions survive exactly as written.

\subsection{\texorpdfstring{4. DERIVED-NEW-043:\allowbreak{} terminal uniqueness versus uniqueness of all stages}{4.  terminal uniqueness versus uniqueness of all stages}}\label{proofs-12008-12370-derived-new-043-terminal-uniqueness-versus-uniqueness-of-all-stages}

The source lemma at original 12242–\allowbreak{}12317 says any of its three conditions gives at most one map of systems. Conditions (1)\allowbreak{} and (2)\allowbreak{} give only uniqueness of the terminal map. Condition (3)\allowbreak{} does give uniqueness of the entire map of systems.

For (1)\allowbreak{},\allowbreak{} let \(T=Z'_n\). By (1.4)\allowbreak{},\allowbreak{} the zero group \(\operatorname{Hom}(X_j[j],T)\) makes \(\operatorname{Hom}(Z_j,T)\to\operatorname{Hom}(Z_{j-1},T)\) injective for every \(1\leq j\leq n\). The restriction of a terminal comparison to \(Z_0=Y_0\) is the fixed composite \[
Y_0\xrightarrow{p_0}X_0\xrightarrow{f_0}X'_0
 \xrightarrow{(p'_0)^{-1}}Y'_0\longrightarrow Z'_n.
\] Successive injectivity determines its terminal map. It determines the composites of intermediate maps with \(Z'_j\to Z'_n\),\allowbreak{} not the intermediate maps themselves.

For (2)\allowbreak{},\allowbreak{} set \(S=Z_n\). Its hypotheses are \(\operatorname{Hom}(S,X'_j[j])=0\) for \(0\leq j<n\). Starting with \(Z'_0\cong X'_0\),\allowbreak{} induction on the target tower and exact Hom show \(\operatorname{Hom}(S,Z'_{n-1})=0\). The last target triangle makes \[
\operatorname{Hom}(S,Z'_n)\longrightarrow\operatorname{Hom}(S,X'_n[n])
\] injective. Its value on a terminal comparison is \(f_n[n]p_n[n]\),\allowbreak{} which is fixed. This again proves terminal uniqueness.

For (3)\allowbreak{},\allowbreak{} induct on length. Assume the maps below stage \(n\) have been proved unique. Inverse rotation gives triangles with objects \((Y_{n-1}[-1],Y_n,X_n)\) and their primed counterparts,\allowbreak{} with first arrow \(-\delta_n[-1]\). Part (5)\allowbreak{} of the source uniqueness-of-third-arrow lemma gives uniqueness of the middle map if \[
\operatorname{Hom}(Y_{n-1},X'_n)=0,\qquad
\operatorname{Hom}(X_n,Y'_{n-1}[-1])=0. \tag{4.1}
\] For completeness,\allowbreak{} that part follows as follows. If two middle maps differ by \(b\),\allowbreak{} then the two outside commuting squares imply \(bf=0\) and \(g'b=0\) for triangles \(X\xrightarrow fY\xrightarrow gZ\to X[1]\),\allowbreak{} \(X'\xrightarrow {f'}Y'\xrightarrow {g'}Z'\to X'[1]\). Exactness gives \(b=eg\). Then \(g'eg=0\) implies \(g'e=vh\),\allowbreak{} where \(h:Z\to X[1]\) and \(v:X[1]\to Z'\). If \(\operatorname{Hom}(X[1],Z')=0\),\allowbreak{} then \(g'e=0\),\allowbreak{} so \(e=f'w\) with \(w:Z\to X'\). If also \(\operatorname{Hom}(Z,X')=0\),\allowbreak{} then \(e=0\),\allowbreak{} and \(b=0\). These two groups become exactly (4.1)\allowbreak{}. The source and target filtrations express their outer objects by successive triangle layers \(X_{n-i}[-i+1]\),\allowbreak{} respectively \(X'_{n-i}[-i]\),\allowbreak{} for \(1\leq i\leq n\). Exact Hom and induction from \(Y_0\cong X_0\) imply (4.1)\allowbreak{} from the two groups in source condition (3)\allowbreak{}. This proves uniqueness at stage \(n\) and hence of the whole map. It does not prove existence without the separate extension argument.

Here is a counterexample to whole-system uniqueness satisfying BOTH (1)\allowbreak{} and (2)\allowbreak{}. Use \(K,V,h\) and the common first stage (2.1)\allowbreak{}. Set the higher stages as follows; all underlying maps \(f_j\) in the comparison are zero.

{ % Captionless table using the bundled longtable counter support.
\begin{longtable}[]{@{}
  >{\raggedright\arraybackslash}p{(\linewidth - 8\tabcolsep) * \real{0.2000}}
  >{\raggedright\arraybackslash}p{(\linewidth - 8\tabcolsep) * \real{0.2000}}
  >{\raggedright\arraybackslash}p{(\linewidth - 8\tabcolsep) * \real{0.2000}}
  >{\raggedright\arraybackslash}p{(\linewidth - 8\tabcolsep) * \real{0.2000}}
  >{\raggedright\arraybackslash}p{(\linewidth - 8\tabcolsep) * \real{0.2000}}@{}}
\toprule\noalign{}
\begin{minipage}[b]{\linewidth}\raggedright
Stage
\end{minipage} & \begin{minipage}[b]{\linewidth}\raggedright
Source \(X_j\)
\end{minipage} & \begin{minipage}[b]{\linewidth}\raggedright
Source \(Y_j\)
\end{minipage} & \begin{minipage}[b]{\linewidth}\raggedright
Target \(X'_j\)
\end{minipage} & \begin{minipage}[b]{\linewidth}\raggedright
Target \(Y'_j\)
\end{minipage} \\
\midrule\noalign{}
\endhead
\bottomrule\noalign{}
\endlastfoot
0 & \(K[1]\) & \(K[1]\) & \(K[1]\) & \(K[1]\) \\
1 & \(K\) & \(V\) & \(K\) & \(V\) \\
2 & \(0\) & \(V[-1]\) & \(V\) & \(0\) \\
3 & \(V[-1]\) & \(0\) & \(0\) & \(0\) \\
\end{longtable}
}

At source stage two choose \[
V[-1]\xrightarrow0 0\xrightarrow0 V\xrightarrow{1}V;
\] at target stage two choose \[
0\longrightarrow V\xrightarrow{1}V\longrightarrow0.
\] At source stage three choose \(0\to V[-1]\xrightarrow{1}V[-1]\to0\); at target stage three take the zero triangle. All are distinguished by rotations and sign changes of identity triangles. The underlying source complex has all arrows zero. The target complex has \(d'_2=\operatorname{pr}_1:V\to K\),\allowbreak{} with other arrows zero,\allowbreak{} so it too is a complex.

One comparison is identically zero. Another is \[
\phi_0=0,\qquad \phi_1=h,\qquad \phi_2=0,\qquad\phi_3=0.
\] Stage one satisfies (1.2)\allowbreak{} by the two explicit equations for \(h\). At stage two the equation involving \(t_2\) has source \(X_2=0\),\allowbreak{} and the connecting-map equation has target connecting map zero and \(\phi_2=0\); both vanish. The projection equation also vanishes. Stage three has zero target objects and maps. Thus all required squares commute. Since \(h\ne0\),\allowbreak{} the two system maps differ. However \(Y_3=Y'_3=0\),\allowbreak{} so every group in both conditions (1)\allowbreak{} and (2)\allowbreak{} vanishes. This disproves the original conclusion under exactly its printed hypotheses.

The minimal conclusion replacement is:\allowbreak{} any two maps of systems induce the same map \(Y_n\to Y'_n\); in case (3)\allowbreak{},\allowbreak{} the map of systems itself is unique. This retains every valid part of the source and introduces no new existence assertion.

\subsection{\texorpdfstring{5. DERIVED-UNDER-032:\allowbreak{} exact obstruction and groups of choices}{5.  exact obstruction and groups of choices}}\label{proofs-12008-12370-derived-under-032-exact-obstruction-and-groups-of-choices}

This consequence retains the actual chosen partial system; it does not replace a nonzero obstruction by an assumption. Fix the triangle \[
Y_{n-1}\xrightarrow{p}X_{n-1}\xrightarrow{t}Y_{n-2}
 \xrightarrow{\delta}Y_{n-1}[1].
\] For the given \(d_n\),\allowbreak{} the complete extension obstruction is the concrete element \[
o=t d_n\in\operatorname{Hom}(X_n,Y_{n-2}). \tag{5.1}
\] The relevant exact sequence is 
\begingroup\fontsize{9.5}{11.5}\selectfont
\[
\operatorname{Hom}(X_n,X_{n-1}[-1])
 \xrightarrow{t[-1]_*}\operatorname{Hom}(X_n,Y_{n-2}[-1])
 \xrightarrow{(-\delta[-1])_*}\operatorname{Hom}(X_n,Y_{n-1})
 \xrightarrow{p_*}\operatorname{Hom}(X_n,X_{n-1})
 \xrightarrow{t_*}\operatorname{Hom}(X_n,Y_{n-2}). \tag{5.2}
\]
\endgroup
 Exactness proves that a lift exists exactly when the computed element (5.1)\allowbreak{} is zero. In Sections 2 and 3 this element was actually computed or annihilated using the stated hypotheses,\allowbreak{} including an example where every possible preceding choice gives nonzero obstruction.

When a lift \(a\) exists,\allowbreak{} the additive group \[
Q_n=
\operatorname{Hom}(X_n,Y_{n-2}[-1])/
 t[-1]_*\operatorname{Hom}(X_n,X_{n-1}[-1]) \tag{5.3}
\] acts on all lifts by \([z]\cdot a=a-\delta[-1]z\). Exactness at the second group proves this is independent of the representative and free. Exactness at the third group proves transitivity,\allowbreak{} since the difference of any two lifts lies in \(\ker p_*\). Thus (5.3)\allowbreak{} is the exact group of lift choices; the inverse-rotation minus sign is retained. Completion of a fixed lift to a triangle exists by TR1 and gives fiber objects isomorphic over the two fixed adjacent objects by TR3. Such a triangle isomorphism need not be unique,\allowbreak{} as Section 4 shows.

For fixed source and target systems and a fixed comparison through stage \(n-1\),\allowbreak{} the complete next-map obstruction is the actual element \[
\alpha=t'_nf_n-\phi_{n-1}t_n\in\operatorname{Hom}(X_n,Y'_{n-1}). \tag{5.4}
\] A completion requires \(\alpha=0\) by its middle square,\allowbreak{} and TR3 proves sufficiency. The calculation in Section 3 places it in the image of \((-\delta'_{n-1}[-1])_*\operatorname{Hom}(X_n,Y'_{n-2}[-1])\); this records its more precise location even when that group does not vanish. If one completion \(\phi_n\) exists,\allowbreak{} all completions form a torsor under \[
\ker\!\left[
\operatorname{Hom}(Y_n,Y'_n)\longrightarrow
\operatorname{Hom}(Y_{n-1}[-1],Y'_n)\oplus\operatorname{Hom}(Y_n,X'_n)
\right],\qquad
h\longmapsto\big(h(-\delta_n[-1]),p'_nh\big). \tag{5.5}
\] Indeed,\allowbreak{} subtract the three equations (1.2)\allowbreak{}. The middle equation involves no \(\phi_n\),\allowbreak{} while the other two are precisely the two displayed zero conditions after shifting the last equation back. Conversely any such \(h\) added to \(\phi_n\) preserves all three squares. Addition is free and transitive. Formula (5.5)\allowbreak{} exactly describes the intermediate ambiguity missed by the printed conclusion.

These are proved exact descriptions,\allowbreak{} not prospective existence theorems with the missing mathematics assumed. The obstruction depends on the specified partial system or partial comparison. No choice-independent global obstruction class is asserted. The results are formal consequences of the source exact Hom sequences and carry no novelty claim.

\subsection{6. Physical reports and prior correction identity}\label{proofs-12008-12370-physical-reports-and-prior-correction-identity}

P02-E0097,\allowbreak{} original producer line 101,\allowbreak{} spans original 12148–\allowbreak{}12160. Its exact minimal replacement \(Y_{n-1}\mapsto Y_{n-2}\) at the final target is correct by (1.1)\allowbreak{},\allowbreak{} (2.2)\allowbreak{},\allowbreak{} and (5.1)\allowbreak{}. The explanation is restated in Section 2 because the report\textquotesingle s wording about the triangle itself is imprecise. This is one missing mathematical-index report,\allowbreak{} semantic group DERIVED-RECON-100,\allowbreak{} and one new operation.

P02-E0098,\allowbreak{} producer line 102,\allowbreak{} and French DERIVED-090,\allowbreak{} producer line 967,\allowbreak{} are two physical reports of one already corrected defect,\allowbreak{} semantic group DERIVED-RECON-101. In original 12278 the first composition must end at \(X'_n[n]\),\allowbreak{} since its penultimate object is \(Y'_n[n]\) and its final map is \(p'_n[n]\). The next line is \(f_n[n]p_n[n]\),\allowbreak{} with the same codomain. This is exactly prior MC-STK-ERR-0901,\allowbreak{} which the recorder replays against authority and current bytes. The operation is not repeated. Its correctness does not establish the stronger uniqueness claim:\allowbreak{} Section 4 independently tests and corrects that claim.

The part therefore completes three additional physical Derived reports,\allowbreak{} bringing the chapter to 172/\allowbreak{}176 in 101 groups. The remaining four reports lie in the final section after this one.

\subsection{7. Receiving proofs and retained pending chapter work}\label{proofs-12008-12370-receiving-proofs-and-retained-pending-chapter-work}

A complete literal-reference search of the root chapter TeX files finds no use of the uniqueness-of-maps lemma outside its definition. This is a bounded repository-reference statement,\allowbreak{} not a claim about all mathematics or all external uses. Existing candidate/\allowbreak{}archive copies are not separate live receiving chapters.

The references in equiv.tex use the complex example or the separate existence/\allowbreak{}isomorphism lemma,\allowbreak{} not the false whole-system uniqueness conclusion.

At equiv.tex 1990–\allowbreak{}2033,\allowbreak{} the two exact functors send a chosen complex Postnikov system to systems over naturally isomorphic complexes. Under the displayed negative-Ext vanishing,\allowbreak{} \(P\) holds because \(i-j-1>0\); equivalently the relevant Ext degree is negative. Section 3 supplies an isomorphism,\allowbreak{} hence the terminal objects are isomorphic after the exact functors\textquotesingle{} shift comparison. This proves the precise use in the siblings argument without claiming a unique isomorphism or a natural isomorphism of the two functors on all objects.

At 3008–\allowbreak{}3113,\allowbreak{} the original resolution objects \(M_n\) admit the signed construction of Section 1. The printed Künneth and full-faithfulness computations give the displayed negative Ext vanishing on the new terms. Both \(P'\) and \(P\) follow from that vanishing. Stagewise existence constructs the infinite system. At 3116–\allowbreak{}3178,\allowbreak{} the two exact functors yield systems over the stated common complex; the same negative Ext hypothesis gives extensions of each already selected partial comparison. Thus a compatible infinite family of isomorphisms exists. This checks the categorical use and its signs; it does not independently adjudicate the entire Fourier–\allowbreak{}Mukai proof or its earlier geometric input.

At spaces-simplicial.tex 2957–\allowbreak{}3120,\allowbreak{} the source\textquotesingle s finite-sum adjunction calculation reduces \(P'\) to the assumed positive-source-shift vanishing on \(K_i\),\allowbreak{} and its comparison calculation similarly gives \(P\) at each \(K_m\). Section 3 justifies extending a chosen partial system and a chosen partial comparison,\allowbreak{} including their infinite compatibility. The exact functors in 2784–\allowbreak{}2826 and 3230–\allowbreak{}3273 transport the maps of Section 1 along their specified shift isomorphisms.

The receiver contains additional local issues which must not be hidden by saying this chapter is approved:\allowbreak{} 3016 calls \(g_m^{-1}Y_n[n]\) itself a Postnikov system,\allowbreak{} although its stages are \(g_m^{-1}Y_n\); 3049 and 3068 omit the \([n]\) on the target tower in the homotopy colimit; and 3095–\allowbreak{}3110 interchange the indices of the explicitly defined \(Y'_{m,n}\) and \(\gamma_{m,n}\). The correct comparison supplied by the present proof is exactly \[
\operatorname{hocolim}_n g_m^{-1}Y_n[n]
 \xrightarrow{\operatorname{hocolim}_n\gamma_{m,n}[n]}
\operatorname{hocolim}_n Y'_{m,n}[n]\ \cong\ K_m. \tag{7.1}
\] Here ``hocolim of the comparison'' means a TR3 completion of the square on the two telescope sums; it is an isomorphism by the two-isomorphisms property. No canonical choice is asserted. The stage-zero labels in the later diagram are \(Y'_{m,0},Y'_{n,0}\),\allowbreak{} and its comparisons are \(\gamma_{m,0},\gamma_{n,0}\). The formula \(Y'_{m,n}=B_{m,\leq n}[-n]\otimes^{\mathbf L}K_m\) makes these indices and shifts explicit. These receiver observations are durably routed for its own complete report reconciliation; they are not silently counted as completed Derived reports or as an approved receiving chapter. Also original receiver reports P10-E0193/\allowbreak{}E0194 concern the module statement and a reference; their exact original evidence remains to be read in that chapter\textquotesingle s queue.

In particular,\allowbreak{} the corrected uniqueness result does not invalidate these uses of existence and isomorphism. The receiver-specific notation issues remain visible rather than being conflated with DERIVED-NEW-043.

\subsection{8. Scope and continuation}\label{proofs-12008-12370-scope-and-continuation}

Semantic review of this source chapter is complete through original 12370. Original 12380–\allowbreak{}12430 has been read ahead only. Continue with Essentially constant systems at 12380,\allowbreak{} next unread line 12431. The exact obstruction descriptions remain editorial additions. Later synthesis and novelty review stay after the core correction and received-theorem queues. Pursuing Stacks and Illusie already have readers; the hub must not duplicate those assignments.
\clearpage
\section{\texorpdfstring{Derived Categories:\allowbreak{} essentially constant systems,\allowbreak{} original 12380–\allowbreak{}12604}{Derived  essentially constant  original }}\label{proofs-12380-12604-derived-categories-essentially-constant-systems-original-1238012604}

Authority:\allowbreak{} Stacks Project commit a04446e5\allowbreak{}7ec1fbc2\allowbreak{}52a871af\allowbreak{}cec7752f\allowbreak{}b2807b14, derived.tex,\allowbreak{} SHA-256 EFDA4CB3\allowbreak{}61DD4090\allowbreak{}9D199116\allowbreak{}1EAE716E\allowbreak{}58B6CF9E\allowbreak{}79380971\allowbreak{}502E1DA1\allowbreak{}2396D402. This is a complete editorial proof supplement for the terminal mathematical section. DERIVED-NEW-044 repairs a finite-stage proof gap without changing the theorem. DERIVED-UNDER-033 records the finite cohomological-degree bound proved below; no novelty is claimed.

\subsection{\texorpdfstring{1. Essential constancy,\allowbreak{} its splitting,\allowbreak{} and the actual pro-category maps}{1. Essential  its  and the actual pro-category maps}}\label{proofs-12380-12604-essential-constancy-its-splitting-and-the-actual-pro-category-maps}

For a directed inverse system write \(t_{ji}:A_j\to A_i\) for \(j\geq i\). The source definition of essential constancy gives an object \(A\),\allowbreak{} a cone \(q_i:A\to A_i\),\allowbreak{} and \(p:A_{i_0}\to A\) with \(pq_{i_0}=1_A\),\allowbreak{} such that for each \(i\) some \(j\geq i,i_0\) satisfies \[
t_{ji}=q_i p t_{j i_0}. \tag{1.1}
\] This is the inverse-system specialization of Categories 2945–\allowbreak{}2955. It is stronger than mere existence of a limit.

For \(i\geq i_0\),\allowbreak{} put \(p_i=p t_{i i_0}\); then \(p_iq_i=1_A\). Complete \(p_i\) to a triangle \[
Z_i\xrightarrow{z_i}A_i\xrightarrow{p_i}A\xrightarrow{\epsilon_i}Z_i[1].
\] Consecutive composites vanish,\allowbreak{} so \(\epsilon_i=\epsilon_ip_iq_i=0\). The split-triangle proof (original Derived 699–\allowbreak{}726)\allowbreak{} makes \(z_i\oplus q_i:Z_i\oplus A\to A_i\) an isomorphism. For any \(T\),\allowbreak{} exact Hom gives \(0\to\operatorname{Hom}(T,Z_i)\to\operatorname{Hom}(T,A_i)
 \to\operatorname{Hom}(T,A)\to0\); the initial zero follows from \(\epsilon_i[-1]=0\). Thus \(z_i\) is an actual categorical kernel of \(p_i\). The identities \(p_i t_{ji}=p_j\) and \(t_{ji}q_j=q_i\) show the transition matrix is \(\operatorname{diag}(1_A,z_{ji})\) in the ordered decomposition \(A\oplus Z_i\). Equation (1.1)\allowbreak{} says \(z_{ji}=0\) for a sufficiently late \(j\) at each fixed \(i\). This proves the source\textquotesingle s forward implication without assuming arbitrary idempotents split in the triangulated category:\allowbreak{} the exhibited retraction has a complement by its triangle.

Conversely,\allowbreak{} suppose the displayed decompositions and transition matrices of the source lemma exist on a tail. The inclusions of \(A\) give a cone there,\allowbreak{} and hence on the whole system by composing to each earlier index; independence follows by taking a common later index. The projection at the first index has composite identity on \(A\). For each \(i\),\allowbreak{} choose \(j\) killing \(Z_j\to Z_i\); then the transition factors through \(A\) exactly as in (1.1)\allowbreak{}. This is essential constancy. For a compatible cone from any \(T\),\allowbreak{} its component into \(Z_i\) is zero,\allowbreak{} because it factors through a transition \(Z_j\to Z_i\) which is zero. Its components into \(A\) are equal. Consequently the actual limit is \(A\),\allowbreak{} with the specified cone maps.

For clarity about the subsequent pro-category argument,\allowbreak{} let \[
F_X(T)=\operatorname{colim}_i\operatorname{Hom}(X_i,T).
\] A pro-map \(a:X\to Y\) consists of compatible classes \([a_j]\in\operatorname{colim}_i\operatorname{Hom}(X_i,Y_j)\). It induces \(F_Y\to F_X\):\allowbreak{} the class of \(u:Y_j\to T\) is sent to the class of \(u a_j\). Refining indices proves this is independent of both representatives; compatibility in \(j\) proves it respects the colimit. It is natural in \(T\) and respects composition. Conversely,\allowbreak{} evaluate a natural transformation \(F_Y\to F_X\) at the class of \(1_{Y_j}\); the resulting classes give a pro-map,\allowbreak{} and naturality recovers every class \(u\). These constructions are inverse. Equivalently,\allowbreak{} \[
\operatorname{Nat}(F_Y,F_X)
 =\lim_j\operatorname{colim}_i\operatorname{Hom}(X_i,Y_j)
 =\operatorname{Hom}_{\operatorname{Pro}(\mathcal D)}(X,Y). \tag{1.2}
\] This proves the exact contravariant fully faithful comparison cited by the source,\allowbreak{} rather than using only the name of the copresheaf construction.

An isomorphism from \(X\) to a constant pro-object \(A\) gives an actual cone \(A\to X_i\) and a map \(X_{i_0}\to A\) representing its inverse. The constant composite is literally \(1_A\); equality of the other composite with the pro-identity is precisely eventual equality (1.1)\allowbreak{} at each index. Thus pro-isomorphism to a constant system and the source\textquotesingle s essential constancy are equivalent.

\subsection{\texorpdfstring{2. DERIVED-NEW-044:\allowbreak{} the needed additional tail in the triangle comparison}{2.  the needed additional tail in the triangle comparison}}\label{proofs-12380-12604-derived-new-044-the-needed-additional-tail-in-the-triangle-comparison}

Suppose \[
A_n\xrightarrow{a_n}B_n\xrightarrow{b_n}C_n
 \xrightarrow{d_n}A_n[1] \tag{2.1}
\] is an inverse sequence of triangles and the outer systems are essentially constant. After taking a common tail,\allowbreak{} Section 1 gives decompositions \(A_n=A\oplus A'_n\),\allowbreak{} \(C_n=C\oplus C'_n\),\allowbreak{} with diagonal transitions and pro-zero primed systems. Write the actual connecting map in blocks:\allowbreak{} \[
d_n=
\begin{pmatrix}
\delta_n&\beta_n\\
\gamma_n&\eta_n
\end{pmatrix}
:C\oplus C'_n\longrightarrow A[1]\oplus A'_n[1]. \tag{2.2}
\] For \(m\geq n\),\allowbreak{} compatibility gives \[
\delta_m=\delta_n,\qquad
\gamma_n=a'_{mn}[1]\gamma_m,\qquad
\beta_m=\beta_n c'_{mn},\qquad
a'_{mn}[1]\eta_m=\eta_n c'_{mn}. \tag{2.3}
\] At each fixed \(n\),\allowbreak{} choose \(m\) killing \(a'_{mn}\); then \(\gamma_n=0\). Thus the assertion in original 12450–\allowbreak{}12451 that \(C\) maps into \(A[1]\) by a fixed map \(\delta\) is correct.

However,\allowbreak{} the projections \(C_n\to C\) and \(A_n[1]\to A[1]\) commute with the connecting maps only when \(\beta_n=0\). This does not follow at the first index merely from \(\gamma_n=0\). Fix the first remaining index \(n_0\),\allowbreak{} and choose \(N\geq n_0\) with \(c'_{N n_0}=0\). Equation (2.3)\allowbreak{} gives \(\beta_N=0\),\allowbreak{} and then \(\beta_n=0\) for every \(n\geq N\). Passing to this additional tail makes the needed square commute. This is the missing finite-stage justification before the TR3 step at original 12456–\allowbreak{}12461.

The failure at an unadjusted first index is explicit. In \(D^b(k)\),\allowbreak{} put \(K=k[0]\). Let \(A_n=K\) with identity transitions; let \(C_1=K[1]\) and \(C_n=0\) for \(n\geq2\); let \(B_1=0\) and \(B_n=K\) for \(n\geq2\). At the first index use \[
K\longrightarrow0\longrightarrow K[1]\xrightarrow{1}K[1],
\] and at every later index use \(K\xrightarrow{1}K\to0\to K[1]\). The transition from index two to one has components \(1_K,0,0\); every square commutes,\allowbreak{} including the shifted last square. The outer systems have values \(K\) and \(0\). Their constant triangle is \(K\xrightarrow{1}K\to0\to K[1]\). A comparison from the first triangle to this one inducing the identity on \(A_1=K\) would force \(0=1_K\) in the first square. There is no such map. At index two it does exist. Thus the proof needs the additional tail; the theorem,\allowbreak{} which asks for some index,\allowbreak{} is unchanged.

After the adjustment,\allowbreak{} choose a distinguished triangle \[
A\xrightarrow{a}B\xrightarrow{b}C\xrightarrow{\delta}A[1].
\] Apply TR3 to the square on \(C_N,A_N[1]\),\allowbreak{} then rotate back,\allowbreak{} obtaining a triangle map from (2.1)\allowbreak{} at \(N\) to this fixed triangle. For \(n\geq N\),\allowbreak{} compose with the given transition triangle maps to obtain compatible maps \(p^B_n:B_n\to B\).

For every \(T\),\allowbreak{} the induced long exact Hom sequences give a natural diagram \[
\operatorname{Hom}(C[s],T)\to\operatorname{Hom}(B[s],T)
 \to\operatorname{Hom}(A[s],T)\to\operatorname{Hom}(C[s-1],T)
\] over the corresponding filtered colimits of the groups from (2.1)\allowbreak{},\allowbreak{} for every integer \(s\). The colimit sequence is exact:\allowbreak{} for an element whose image is zero,\allowbreak{} represent it at one stage,\allowbreak{} advance to a stage where its image is literally zero,\allowbreak{} and use exactness at that stage. The outer comparisons are isomorphisms because the primed systems are pro-zero. Applying the five lemma to the five consecutive terms surrounding the middle term makes the middle comparison an isomorphism. These are maps of functors in \(T\),\allowbreak{} not just group isomorphisms selected separately. Section 1 identifies the resulting comparison as a pro-isomorphism \(B_\bullet\to B\),\allowbreak{} hence as essential constancy with value \(B\).

For completeness,\allowbreak{} a pro-isomorphism from a system \(B_n\) to a constant \(B\),\allowbreak{} represented on a tail by the compatible maps \(p^B_n\),\allowbreak{} induces the required isomorphism on its actual limit. Let \(q_n:B\to B_n\) be the inverse cone. Then \(p^B_nq_n=1_B\) at every index where \(p^B_n\) represents the pro-map,\allowbreak{} and \(q_n p^B_m=t_{mn}\) after a sufficiently late \(m\),\allowbreak{} for each fixed \(n\). These are the identities from essential constancy and show its limit cone is \(q_n\). The same reasoning applies to \(A,C\). Thus the final limit clause of the source lemma is proved as well.

The minimal correction adds the calculation killing the components \(C'_n\to A[1]\) and one more increase of the starting index. The original statement,\allowbreak{} its triangle,\allowbreak{} and all receiving hypotheses are preserved.

\subsection{3. Essentially constant bounded cohomology}\label{proofs-12380-12604-essentially-constant-bounded-cohomology}

Suppose every \(E_n\) has cohomology in the same original interval \([a,b]\),\allowbreak{} and every system \(H^i(E_n)\) is essentially constant. If \(a=b\),\allowbreak{} the natural truncation comparisons identify \(E_n\) with \(H^a(E_n)[-a]\),\allowbreak{} compatibly with its transition maps. Thus the system is essentially constant.

For \(a<b\),\allowbreak{} use the actual functorial triangles \[
\tau_{\leq a}E_n\longrightarrow E_n\longrightarrow
\tau_{\geq a+1}E_n\longrightarrow(\tau_{\leq a}E_n)[1].
\] The first system is naturally isomorphic to \(H^a(E_n)[-a]\). The third has cohomology in \([a+1,b]\) and meets the induction hypothesis. Section 2 proves essential constancy of the middle system. This induction ends after the finite interval; no exactness of products,\allowbreak{} enough injectives,\allowbreak{} or existence of an arbitrary derived inverse limit is used.

Let its value be \(E\). Applying any functor \(G\) to the explicit identities (1.1)\allowbreak{} and \(pq=1\) shows that \(G(E_n)\) is essentially constant with value \(G(E)\). In particular,\allowbreak{} \(H^i(E)\) is the original value of \(H^i(E_n)\). This is the omitted cohomology assertion,\allowbreak{} and it does not assume cohomology commutes with arbitrary inverse limits.

\subsection{4. Pro-isomorphisms and the actual cone sequence}\label{proofs-12380-12604-pro-isomorphisms-and-the-actual-cone-sequence}

If \(A_n\to B_n\to C_n\to A_n[1]\) is an inverse sequence of triangles with \(C_n\) pro-zero,\allowbreak{} then both \(\operatorname{colim}_n\operatorname{Hom}(C_n,T)\) and \(\operatorname{colim}_n\operatorname{Hom}(C_n[-1],T)\) are zero. Indeed,\allowbreak{} a representative map out of \(C_n\) becomes zero after precomposition with a zero transition. Exactness gives a natural isomorphism \[
\operatorname{colim}_n\operatorname{Hom}(B_n,T)
 \longrightarrow\operatorname{colim}_n\operatorname{Hom}(A_n,T).
\] The explicit comparison (1.2)\allowbreak{} proves that the original maps \(A_n\to B_n\) give a pro-isomorphism.

Now suppose a strict sequence of maps \(f_n:A_n\to B_n\) has both systems in the same interval \([a,b]\),\allowbreak{} and every \(H^p(f_n)\) is a pro-isomorphism. Choose a cone triangle for each \(f_n\). For each adjacent transition,\allowbreak{} TR3 supplies a third component; define transitions over larger gaps by composition. Because the index is the natural numbers,\allowbreak{} this constructs a coherent inverse sequence of triangles. It does not claim an arbitrary diagram of cones has automatically coherent choices.

The long exact cohomology sequence shows \[
H^p(C_n)=0\quad\hbox{if }p<a-1\hbox{ or }p>b. \tag{4.1}
\] The lower endpoint is \(a-1\),\allowbreak{} not \(a\). For each \(p\),\allowbreak{} put \[
Q_n^p=\operatorname{coker}H^p(f_n),\qquad
J_n^{p+1}=\ker H^{p+1}(f_n).
\] There is a compatible short exact sequence \[
0\longrightarrow Q_n^p\longrightarrow H^p(C_n)
 \longrightarrow J_n^{p+1}\longrightarrow0. \tag{4.2}
\]

Here is the exact reason these two outside systems are pro-zero. For a strict map \(u:X_\bullet\to Y_\bullet\) with pro-inverse,\allowbreak{} the equality \(u u^{-1}=1_Y\) implies that,\allowbreak{} at every target index \(n\),\allowbreak{} some transition \(Y_m\to Y_n\) factors through \(u_n\). It therefore kills the cokernel system at \(n\). Similarly \(u^{-1}u=1_X\) implies that after advancing the source index some transition \(X_k\to X_n\) has the form \(v u_k\) for an actual map \(v:Y_k\to X_n\),\allowbreak{} possibly after refining the representative of the inverse. It kills \(\ker u_k\). These assertions follow directly from eventual equality in the filtered colimit defining each component of a pro-map; they do not use an unproved abelian structure on a pro-category.

Fix \(n\). Choose \(m\geq n\) killing \(Q_m^p\to Q_n^p\),\allowbreak{} and then \(k\geq m\) killing \(J_k^{p+1}\to J_m^{p+1}\). By (4.2)\allowbreak{},\allowbreak{} the image of \(H^p(C_k)\to H^p(C_m)\) lies in \(Q_m^p\); its image at \(n\) is therefore zero. Thus \(H^p(C_n)\) is pro-zero. Together with (4.1)\allowbreak{},\allowbreak{} Section 3 makes \(C_n\) essentially constant with an object \(C\) whose cohomology is zero in every degree. Such an object is zero in the derived category,\allowbreak{} by its definition as the localization at quasi-isomorphisms. Hence \(C_n\) is pro-zero. The preceding paragraph proves the required pro-isomorphism \(A_\bullet\to B_\bullet\).

The converse also holds for these maps:\allowbreak{} a pro-isomorphism remains one after the functor \(H^p\),\allowbreak{} because functors preserve the component maps,\allowbreak{} refinements,\allowbreak{} and inverse equations. Thus,\allowbreak{} in the common bounded interval,\allowbreak{} cohomological pro-isomorphism is an exact detection criterion.

\subsection{\texorpdfstring{5. DERIVED-UNDER-033:\allowbreak{} a finite-degree annihilation bound}{5.  a finite-degree annihilation bound}}\label{proofs-12380-12604-derived-under-033-a-finite-degree-annihilation-bound}

A morphism \(f:E\to F\) in \(D(\mathcal A)\) will be called cohomologically zero when \(H^q(f)=0\) for every integer \(q\). Retain the original degree labels and canonical truncation morphisms.

For each \(q\),\allowbreak{} the map \(\tau_{\leq q}f\) factors through \(\tau_{\leq q-1}F\). To prove this,\allowbreak{} compose it with \(\tau_{\leq q}F\to H^q(F)[-q]\). The restriction of this composite to \(\tau_{\leq q-1}E\) is zero by \(\operatorname{Hom}(D^{\leq q-1},D^{\geq q})=0\). Exact Hom for the truncation triangle of \(E\) makes the composite factor through \(H^q(E)[-q]\). Its induced map on degree-\(q\) cohomology is \(H^q(f)=0\),\allowbreak{} so that intervening morphism between shifted heart objects is zero. Thus the whole composite is zero. Exact Hom for the truncation triangle of \(F\) now gives the asserted factorization. This constructs the actual morphism needed to pass down one filtration step; it is not assumed to be unique.

Let \(S=\{s_1<\cdots<s_r\}\) be a finite set of integers and suppose every object in a chain \[
E_0\xrightarrow{f_1}E_1\xrightarrow{f_2}\cdots
 \xrightarrow{f_r}E_r
\] has cohomology zero outside these exact degrees. If every \(f_j\) is cohomologically zero,\allowbreak{} their composite is zero. Indeed \(E_0\cong\tau_{\leq s_r}E_0\). The factorization just proved makes the first map land in \(\tau_{\leq s_r-1}E_1\). Because every cohomology group in the gap \((s_{r-1},s_r)\) is zero,\allowbreak{} the canonical map \(\tau_{\leq s_{r-1}}E_1\to\tau_{\leq s_r-1}E_1\) is a quasi-isomorphism; use its inverse in the derived category to factor through the former. Naturality of truncation allows composition with \(\tau_{\leq s_{r-1}}f_2\),\allowbreak{} which then factors through \(\tau_{\leq s_{r-1}-1}E_2\). Repeating gives a factorization of the complete composite through \(\tau_{\leq s_1-1}E_r=0\). For \(r=1\) this uses the first factorization directly; if \(S\) is empty,\allowbreak{} every object is already zero. Every skipped degree is retained as an explicitly vanishing cohomology group,\allowbreak{} not removed by a coordinate change.

In particular,\allowbreak{} for a common interval \([a,b]\),\allowbreak{} any composite of \(b-a+1\) cohomologically zero maps is zero. A cohomologically zero endomorphism has \(f^{b-a+1}=0\). More precisely,\allowbreak{} the number of possibly nonzero original cohomological degrees gives the bound \(r\).

For an inverse system supported in \(S\) whose cohomology systems are pro-zero,\allowbreak{} fix an index \(n_0\). At each step choose \(n_{j+1}\geq n_j\) so that the transition kills all \(H^{s_\ell}\),\allowbreak{} \(1\leq\ell\leq r\),\allowbreak{} simultaneously. Such an index exists because finitely many indices have a common upper bound. All other cohomology groups vanish from the outset. The transition \[
E_{n_r}\longrightarrow E_{n_{r-1}}\longrightarrow\cdots
 \longrightarrow E_{n_0}
\] is a composite of \(r\) cohomologically zero maps and hence is zero. This supplies a direct finite procedure proving pro-zero,\allowbreak{} not just an existence statement obtained from an abstract constant value.

Applied to the cones in Section 4,\allowbreak{} there are at most \(b-a+2\) possible degrees,\allowbreak{} namely \([a-1,b]\). At each target index the two-step construction from (4.2)\allowbreak{} can be performed for all these finitely many degrees simultaneously; repeating this for \(b-a+2\) transitions makes the cone transition itself zero. This quantitatively propagates the consequence into the source\textquotesingle s pro-isomorphism criterion.

A finite common range cannot be dropped without a replacement hypothesis. In \(D(k)\),\allowbreak{} let \(E_n\) be the actual complex with \(k\) in every degree \(q\geq n\),\allowbreak{} zero in smaller degrees,\allowbreak{} and all differentials zero. Use the literal inclusions \(E_{n+1}\to E_n\). For every fixed \(q\),\allowbreak{} the system \(H^q(E_n)\) is eventually zero. Yet every transition \(E_m\to E_n\),\allowbreak{} \(m\geq n\),\allowbreak{} is nonzero:\allowbreak{} its cohomology map in degree \(m\) is \(1_k\). Therefore the system is not pro-zero. If it were essentially constant with value \(E\),\allowbreak{} Section 3\textquotesingle s functorial argument would give \(H^q(E)=0\) for all \(q\),\allowbreak{} hence \(E=0\),\allowbreak{} a contradiction. This exhibits the exact limitation of the source result while the finite-degree result proves what does work.

These bounds are editorial consequences of truncation orthogonality,\allowbreak{} with no novelty claim and no hypothesis removed from the source body.

\subsection{6. Remaining physical reports}\label{proofs-12380-12604-remaining-physical-reports}

P02-E0100 (producer line 103)\allowbreak{} and French DERIVED-091 (968)\allowbreak{} both concern original 12551–\allowbreak{}12555. Insert ``Let'' before the displayed inverse system to supply the missing introductory verb. This is exactly existing MC-STK-ERR-0902.

P02-E0101 (104)\allowbreak{} and French DERIVED-092 (969)\allowbreak{} both concern original 12559–\allowbreak{}12560. The head noun ``system'' is singular,\allowbreak{} so ``define'' becomes ``defines''. The p02 report inaccurately names its subject ``This''; the original sentence and French report identify the actual subject. The replacement itself is correct and is exactly existing MC-STK-ERR-0903. Both units are replayed against authority and current bytes; no duplicate operation is added.

These four physical reports form groups DERIVED-RECON-102 and DERIVED-RECON-103. All 176 chapter reports have now been individually accounted for:\allowbreak{} 160 already corrected,\allowbreak{} thirteen missing copyedit/\allowbreak{}notation reports,\allowbreak{} two missing mathematical-index reports,\allowbreak{} and one metadata-only resolution. The independent proof gap in Section 2 is not counted as a received report.

\subsection{\texorpdfstring{7. Receiving uses,\allowbreak{} completed scope,\allowbreak{} and next work}{7. Receiving  completed  and next work}}\label{proofs-12380-12604-receiving-uses-completed-scope-and-next-work}

The literal-reference search of current root chapter sources routes the two-out-of-three theorem to duality.tex 8222 and 8360,\allowbreak{} and the cohomology/\allowbreak{}pro-isomorphism results to more-algebra.tex 27867,\allowbreak{} 27872,\allowbreak{} 29719,\allowbreak{} 30341,\allowbreak{} 30346 and 30389.

In duality.tex 8218–\allowbreak{}8224 the inputs are inverse sequences of restricted triangles with essentially constant outer systems. Taking the extra tail in Section 2 preserves their pro-objects,\allowbreak{} so the same middle value exists. At 8348–\allowbreak{}8362 the cited restriction triangle can be rotated,\allowbreak{} with its prescribed rotation sign,\allowbreak{} so any pair of essentially constant terms gives the third; exact restriction/\allowbreak{}direct-image functors preserve the identities (1.1)\allowbreak{}. Thus the categorical inference is unchanged. Earlier geometric and sheaf-theoretic dependencies remain in that chapter\textquotesingle s review.

In more-algebra.tex 27859–\allowbreak{}27874,\allowbreak{} the negative truncations of the \(r\)-term Koszul complexes have cohomology only in the \(r\) original degrees \([-r,-1]\); once their individual cohomology transitions are killed by the subsequent argument,\allowbreak{} \(r\) successive simultaneous choices kill the derived transition. The augmentation triangle then gives the same pro-isomorphism. The zero-generator case has zero negative truncation.

At 29704–\allowbreak{}29724 and 30366–\allowbreak{}30389,\allowbreak{} the receiving proof first replaces an unbounded derived tensor presentation by the pro-isomorphic Koszul presentation with a common finite cohomology range. The criterion is applied to that bounded presentation,\allowbreak{} and pro-isomorphism is then transported back by the actual tensor functor. No boundedness of the original derived tensor objects is inferred. At 30333–\allowbreak{}30352,\allowbreak{} the cone sequence has the finite range coming from the displayed finite complex; its known pro-zero cohomology makes it pro-zero. These are the exact categorical uses checked here. They do not certify the unreviewed Artin–\allowbreak{}Rees inputs or prove the entire receiving chapter.

All mathematical text of the Derived authority chapter has now been read and semantically reviewed through 12590; terminal apparatus through 12604 has been read. The final bibliography/\allowbreak{}input commands contain no additional substantive mathematical result. The next step is to consolidate the chapter\textquotesingle s complete evidence and proposed operations,\allowbreak{} retaining all inherited corrections,\allowbreak{} conceptual errata,\allowbreak{} editorial consequences,\allowbreak{} and receiver followups before admission/\allowbreak{}composition and the required build/\allowbreak{}review/\allowbreak{}publication workflow. Completion of this chapter\textquotesingle s source review does not complete the larger correction queue or the active programme.
\clearpage
\end{document}
